% !TeX program = xelatex
% Standalone English volume. Use XeLaTeX until the references stabilize.
\documentclass[a4paper,11pt]{article}
\usepackage[T1]{fontenc}
\usepackage{lmodern}
\usepackage[english]{babel}
\usepackage[left=26mm,right=26mm,top=25mm,bottom=25mm,headheight=15pt,headsep=7mm,footskip=11mm]{geometry}
\usepackage{amsmath,amssymb,mathtools,amsthm}
\newtheorem{proposition}{Proposition}[section]
\newtheorem{lemma}[proposition]{Lemma}
\newtheorem{theorem}[proposition]{Theorem}
\newtheorem{corollary}[proposition]{Corollary}
\newtheorem{maintheorem}[proposition]{Main Theorem}
\theoremstyle{definition}
\newtheorem{definition}[proposition]{Definition}
\theoremstyle{remark}
\newtheorem{remark}[proposition]{Remark}
\makeatletter
\renewcommand*\l@section{\@dottedtocline{1}{0em}{2.4em}}
\renewcommand*\l@subsection{\@dottedtocline{2}{1.45em}{3.45em}}
\renewcommand*\l@subsubsection{\@dottedtocline{3}{3.8em}{4.55em}}
\makeatother
\usepackage{xcolor,fancyhdr,titlesec,needspace,xurl}
\usepackage{longtable,multicol}
\usepackage{hyperref}
\definecolor{ink}{HTML}{1A3542}
\definecolor{muted}{HTML}{4F606A}
\hypersetup{unicode=true,colorlinks=true,linkcolor=ink,filecolor=ink,urlcolor=ink,citecolor=ink,bookmarksnumbered=true,
 pdfdisplaydoctitle=true,pdflang={en},
 pdftitle={The generalized Riemann hypothesis for Dirichlet L-functions. Edition 930, Volume I},
 pdfsubject={A natural-language reconstruction of the proposed proof},
 pdfauthor={A mathematical reconstruction}}
\titleformat{\section}{\Large\bfseries\color{ink}}{\thesection}{.65em}{}
\titlespacing*{\section}{0pt}{.8em}{.65em}
\setlength{\parindent}{1.2em}
\linespread{1.045}
\setlength{\parskip}{.22em plus .05em minus .03em}
\setlength{\emergencystretch}{2em}
\setcounter{topnumber}{2}
\setcounter{tocdepth}{3}\setcounter{secnumdepth}{3}
\widowpenalty=4500
\clubpenalty=4500
\pagestyle{fancy}\fancyhf{}
\fancyhead[L]{\small\color{muted}Dirichlet GRH \textperiodcentered\ Volume I}
\fancyhead[R]{\small\color{muted}Dirichlet zeros and power intervals}
\fancyfoot[C]{\small\thepage}
\renewcommand{\headrulewidth}{.3pt}
\newcommand{\lead}[1]{\par\smallskip\noindent{\bfseries\color{ink}#1.}\enspace}
\newcommand{\sour}[2]{\unskip\penalty150\hspace{.35em}{\small\hyperlink{ref#1}{[#2]}}}
\titleformat{\subsection}{\large\bfseries}{\thesubsection}{.6em}{}
\titlespacing*{\subsection}{0pt}{1.15em plus .15em}{.4em}
\setlength{\abovedisplayskip}{9pt plus 2pt minus 3pt}
\setlength{\belowdisplayskip}{9pt plus 2pt minus 3pt}
\newcommand{\RR}{\mathbb R}\newcommand{\CC}{\mathbb C}\newcommand{\QQ}{\mathbb Q}
\newcommand{\ZZ}{\mathbb Z}\newcommand{\ord}{\operatorname{ord}}
\newcommand{\Frob}{\operatorname{Frob}}\newcommand{\Gal}{\operatorname{Gal}}
\newcommand{\len}{\operatorname{length}}\newcommand{\Res}{\operatorname*{Res}}
\newcommand{\wmean}{\operatorname{av}}\newcommand{\Ker}{\operatorname{ker}}
\newcommand{\Lh}{\mathcal L}\newcommand{\Gg}{\mathcal G}
\newcommand{\real}{\operatorname{Re}}\newcommand{\imag}{\operatorname{Im}}
\newcommand{\abs}[1]{\left|#1\right|}\newcommand{\eps}{\varepsilon}
\newcommand{\anchoritem}[2]{\par\Needspace{2.8\baselineskip}\hypertarget{ref#1}{}\noindent{\bfseries\color{ink}#2.}\enspace}
\newcommand{\srcloc}[2]{{\small\nolinkurl{#1.lean}}, line~#2}
\numberwithin{equation}{section}
\raggedbottom

\newcommand{\Sel}{\mathcal S}
\newcommand{\Cl}{\operatorname{Cl}}
\newcommand{\Aut}{\operatorname{Aut}}
\newcommand{\GL}{\operatorname{GL}}
\newcommand{\OO}{\mathcal O}
\hypersetup{bookmarksdepth=3}


\addto\captionsenglish{\renewcommand{\partname}{Volume}}


\makeatletter
\expandafter\def\csname r@cor:actualdenom803\endcsname{{52.21}{80}{}{proposition.52.21}{Dirichlet-GRH-v930-Volume-X.pdf}}
\expandafter\def\csname r@cor:actualdyadic525\endcsname{{57.23}{84}{}{proposition.57.23}{Dirichlet-GRH-v930-Volume-XI.pdf}}
\expandafter\def\csname r@cor:actualgram535\endcsname{{113.8}{89}{}{proposition.113.8}{Dirichlet-GRH-v930-Volume-XXI.pdf}}
\expandafter\def\csname r@cor:actualidentity929\endcsname{{139.36}{138}{}{proposition.139.36}{Dirichlet-GRH-v930-Volume-XXVI.pdf}}
\expandafter\def\csname r@cor:actualmarks714\endcsname{{12.10}{38}{}{proposition.12.10}{Dirichlet-GRH-v930-Volume-II.pdf}}
\expandafter\def\csname r@cor:actualperiodmultiples768\endcsname{{16.28}{59}{}{proposition.16.28}{Dirichlet-GRH-v930-Volume-III.pdf}}
\expandafter\def\csname r@cor:actualsource520\endcsname{{146.14}{103}{}{proposition.146.14}{Dirichlet-GRH-v930-Volume-XXVII.pdf}}
\expandafter\def\csname r@cor:adelicfirst816\endcsname{{77.19}{22}{}{proposition.77.19}{Dirichlet-GRH-v930-Volume-XV.pdf}}
\expandafter\def\csname r@cor:adelicfirst915\endcsname{{139.12}{104}{}{proposition.139.12}{Dirichlet-GRH-v930-Volume-XXVI.pdf}}
\expandafter\def\csname r@cor:alg97\endcsname{{87.3}{34}{}{proposition.87.3}{Dirichlet-GRH-v930-Volume-XVII.pdf}}
\expandafter\def\csname r@cor:alignedrows638\endcsname{{117.72}{106}{}{proposition.117.72}{Dirichlet-GRH-v930-Volume-XXII.pdf}}
\expandafter\def\csname r@cor:allcost454\endcsname{{141.22}{22}{}{proposition.141.22}{Dirichlet-GRH-v930-Volume-XXVII.pdf}}
\expandafter\def\csname r@cor:allcritical733\endcsname{{13.11}{70}{}{proposition.13.11}{Dirichlet-GRH-v930-Volume-II.pdf}}
\expandafter\def\csname r@cor:allmarks888\endcsname{{136.37}{39}{}{proposition.136.37}{Dirichlet-GRH-v930-Volume-XXVI.pdf}}
\expandafter\def\csname r@cor:allmoving663\endcsname{{26.29}{69}{}{proposition.26.29}{Dirichlet-GRH-v930-Volume-V.pdf}}
\expandafter\def\csname r@cor:allprimeperiod862\endcsname{{84.3}{78}{}{proposition.84.3}{Dirichlet-GRH-v930-Volume-XVI.pdf}}
\expandafter\def\csname r@cor:allprimes514\endcsname{{22.20}{78}{}{proposition.22.20}{Dirichlet-GRH-v930-Volume-IV.pdf}}
\expandafter\def\csname r@cor:allprimesstrict601\endcsname{{22.35}{90}{}{proposition.22.35}{Dirichlet-GRH-v930-Volume-IV.pdf}}
\expandafter\def\csname r@cor:allrational645\endcsname{{117.84}{118}{}{proposition.117.84}{Dirichlet-GRH-v930-Volume-XXII.pdf}}
\expandafter\def\csname r@cor:allsample341\endcsname{{22.15}{74}{}{proposition.22.15}{Dirichlet-GRH-v930-Volume-IV.pdf}}
\expandafter\def\csname r@cor:allscalartests770\endcsname{{16.30}{62}{}{proposition.16.30}{Dirichlet-GRH-v930-Volume-III.pdf}}
\expandafter\def\csname r@cor:allterminalcoeff641\endcsname{{117.81}{115}{}{proposition.117.81}{Dirichlet-GRH-v930-Volume-XXII.pdf}}
\expandafter\def\csname r@cor:alltwistcost738\endcsname{{50.10}{10}{}{proposition.50.10}{Dirichlet-GRH-v930-Volume-X.pdf}}
\expandafter\def\csname r@cor:allweightlayer529\endcsname{{57.28}{87}{}{proposition.57.28}{Dirichlet-GRH-v930-Volume-XI.pdf}}
\expandafter\def\csname r@cor:annram405\endcsname{{41.24}{52}{}{proposition.41.24}{Dirichlet-GRH-v930-Volume-VIII.pdf}}
\expandafter\def\csname r@cor:apartmentideals528\endcsname{{96.11}{80}{}{proposition.96.11}{Dirichlet-GRH-v930-Volume-XVIII.pdf}}
\expandafter\def\csname r@cor:arbitraryorders488\endcsname{{68.31}{107}{}{proposition.68.31}{Dirichlet-GRH-v930-Volume-XIII.pdf}}
\expandafter\def\csname r@cor:armdefect879\endcsname{{136.15}{17}{}{proposition.136.15}{Dirichlet-GRH-v930-Volume-XXVI.pdf}}
\expandafter\def\csname r@cor:augdet720\endcsname{{45.12}{49}{}{proposition.45.12}{Dirichlet-GRH-v930-Volume-IX.pdf}}
\expandafter\def\csname r@cor:augheight501\endcsname{{144.15}{72}{}{proposition.144.15}{Dirichlet-GRH-v930-Volume-XXVII.pdf}}
\expandafter\def\csname r@cor:avoidrange343\endcsname{{126.16}{25}{}{proposition.126.16}{Dirichlet-GRH-v930-Volume-XXIV.pdf}}
\expandafter\def\csname r@cor:avoidspace706\endcsname{{133.21}{79}{}{proposition.133.21}{Dirichlet-GRH-v930-Volume-XXV.pdf}}
\expandafter\def\csname r@cor:baseseparation723\endcsname{{45.21}{58}{}{proposition.45.21}{Dirichlet-GRH-v930-Volume-IX.pdf}}
\expandafter\def\csname r@cor:bothaux572\endcsname{{60.6}{28}{}{proposition.60.6}{Dirichlet-GRH-v930-Volume-XII.pdf}}
\expandafter\def\csname r@cor:bothinitial579\endcsname{{60.24}{44}{}{proposition.60.24}{Dirichlet-GRH-v930-Volume-XII.pdf}}
\expandafter\def\csname r@cor:centeredterm496\endcsname{{144.3}{63}{}{proposition.144.3}{Dirichlet-GRH-v930-Volume-XXVII.pdf}}
\expandafter\def\csname r@cor:classfiber372\endcsname{{37.37}{70}{}{proposition.37.37}{Dirichlet-GRH-v930-Volume-VII.pdf}}
\expandafter\def\csname r@cor:classmass503\endcsname{{56.21}{66}{}{proposition.56.21}{Dirichlet-GRH-v930-Volume-XI.pdf}}
\expandafter\def\csname r@cor:closedmoment693\endcsname{{131.35}{47}{}{proposition.131.35}{Dirichlet-GRH-v930-Volume-XXV.pdf}}
\expandafter\def\csname r@cor:cmcomp508\endcsname{{144.21}{76}{}{proposition.144.21}{Dirichlet-GRH-v930-Volume-XXVII.pdf}}
\expandafter\def\csname r@cor:coarsemixed773\endcsname{{51.14}{49}{}{proposition.51.14}{Dirichlet-GRH-v930-Volume-X.pdf}}
\expandafter\def\csname r@cor:codimmarked624\endcsname{{96.33}{97}{}{proposition.96.33}{Dirichlet-GRH-v930-Volume-XVIII.pdf}}
\expandafter\def\csname r@cor:collisionloss545\endcsname{{48.3}{106}{}{proposition.48.3}{Dirichlet-GRH-v930-Volume-IX.pdf}}
\expandafter\def\csname r@cor:column276\endcsname{{83.11}{60}{}{proposition.83.11}{Dirichlet-GRH-v930-Volume-XVI.pdf}}
\expandafter\def\csname r@cor:common240\endcsname{{21.15}{53}{}{proposition.21.15}{Dirichlet-GRH-v930-Volume-IV.pdf}}
\expandafter\def\csname r@cor:commoncapacity489\endcsname{{143.12}{59}{}{proposition.143.12}{Dirichlet-GRH-v930-Volume-XXVII.pdf}}
\expandafter\def\csname r@cor:commoncurves597\endcsname{{62.10}{78}{}{proposition.62.10}{Dirichlet-GRH-v930-Volume-XII.pdf}}
\expandafter\def\csname r@cor:commonminors510\endcsname{{145.4}{82}{}{proposition.145.4}{Dirichlet-GRH-v930-Volume-XXVII.pdf}}
\expandafter\def\csname r@cor:commonwitness575\endcsname{{60.15}{34}{}{proposition.60.15}{Dirichlet-GRH-v930-Volume-XII.pdf}}
\expandafter\def\csname r@cor:complementreproduce812\endcsname{{77.4}{6}{}{proposition.77.4}{Dirichlet-GRH-v930-Volume-XV.pdf}}
\expandafter\def\csname r@cor:complexallweights542\endcsname{{57.32}{90}{}{proposition.57.32}{Dirichlet-GRH-v930-Volume-XI.pdf}}
\expandafter\def\csname r@cor:componentdefect375\endcsname{{38.12}{78}{}{proposition.38.12}{Dirichlet-GRH-v930-Volume-VII.pdf}}
\expandafter\def\csname r@cor:componentgraph419\endcsname{{44.20}{19}{}{proposition.44.20}{Dirichlet-GRH-v930-Volume-IX.pdf}}
\expandafter\def\csname r@cor:conditionedwave917\endcsname{{139.16}{109}{}{proposition.139.16}{Dirichlet-GRH-v930-Volume-XXVI.pdf}}
\expandafter\def\csname r@cor:confluent486\endcsname{{68.23}{102}{}{proposition.68.23}{Dirichlet-GRH-v930-Volume-XIII.pdf}}
\expandafter\def\csname r@cor:coupledglobal386\endcsname{{127.16}{43}{}{proposition.127.16}{Dirichlet-GRH-v930-Volume-XXIV.pdf}}
\expandafter\def\csname r@cor:crossangles664\endcsname{{117.99}{128}{}{proposition.117.99}{Dirichlet-GRH-v930-Volume-XXII.pdf}}
\expandafter\def\csname r@cor:cubecutoff584\endcsname{{61.6}{52}{}{proposition.61.6}{Dirichlet-GRH-v930-Volume-XII.pdf}}
\expandafter\def\csname r@cor:curvesupp371\endcsname{{37.33}{66}{}{proposition.37.33}{Dirichlet-GRH-v930-Volume-VII.pdf}}
\expandafter\def\csname r@cor:cut232\endcsname{{20.20}{39}{}{proposition.20.20}{Dirichlet-GRH-v930-Volume-IV.pdf}}
\expandafter\def\csname r@cor:cyclicuniform582\endcsname{{116.4}{28}{}{proposition.116.4}{Dirichlet-GRH-v930-Volume-XXII.pdf}}
\expandafter\def\csname r@cor:cyclomonodromy476\endcsname{{47.24}{93}{}{proposition.47.24}{Dirichlet-GRH-v930-Volume-IX.pdf}}
\expandafter\def\csname r@cor:denomperiod562\endcsname{{115.18}{17}{}{proposition.115.18}{Dirichlet-GRH-v930-Volume-XXII.pdf}}
\expandafter\def\csname r@cor:detectingsplit650\endcsname{{22.54}{105}{}{proposition.22.54}{Dirichlet-GRH-v930-Volume-IV.pdf}}
\expandafter\def\csname r@cor:detheight810\endcsname{{73.23}{95}{}{proposition.73.23}{Dirichlet-GRH-v930-Volume-XIV.pdf}}
\expandafter\def\csname r@cor:dual68\endcsname{{32.9}{33}{}{proposition.32.9}{Dirichlet-GRH-v930-Volume-VI.pdf}}
\expandafter\def\csname r@cor:dyadicchoice462\endcsname{{141.25}{24}{}{proposition.141.25}{Dirichlet-GRH-v930-Volume-XXVII.pdf}}
\expandafter\def\csname r@cor:dyadiccyclic661\endcsname{{22.60}{110}{}{proposition.22.60}{Dirichlet-GRH-v930-Volume-IV.pdf}}
\expandafter\def\csname r@cor:dyadicdepth458\endcsname{{56.3}{52}{}{proposition.56.3}{Dirichlet-GRH-v930-Volume-XI.pdf}}
\expandafter\def\csname r@cor:dyadicorbit459\endcsname{{56.6}{53}{}{proposition.56.6}{Dirichlet-GRH-v930-Volume-XI.pdf}}
\expandafter\def\csname r@cor:eliminateupper472\endcsname{{57.3}{68}{}{proposition.57.3}{Dirichlet-GRH-v930-Volume-XI.pdf}}
\expandafter\def\csname r@cor:exactdiagonal638\endcsname{{117.71}{106}{}{proposition.117.71}{Dirichlet-GRH-v930-Volume-XXII.pdf}}
\expandafter\def\csname r@cor:exactphase695\endcsname{{132.3}{52}{}{proposition.132.3}{Dirichlet-GRH-v930-Volume-XXV.pdf}}
\expandafter\def\csname r@cor:exceptionalweights642\endcsname{{57.45}{101}{}{proposition.57.45}{Dirichlet-GRH-v930-Volume-XI.pdf}}
\expandafter\def\csname r@cor:family30\endcsname{{24.14}{13}{}{proposition.24.14}{Dirichlet-GRH-v930-Volume-V.pdf}}
\expandafter\def\csname r@cor:fieldrelations712\endcsname{{12.4}{32}{}{proposition.12.4}{Dirichlet-GRH-v930-Volume-II.pdf}}
\expandafter\def\csname r@cor:finitecurve473\endcsname{{47.15}{87}{}{proposition.47.15}{Dirichlet-GRH-v930-Volume-IX.pdf}}
\expandafter\def\csname r@cor:finitedet709\endcsname{{142.10}{44}{}{proposition.142.10}{Dirichlet-GRH-v930-Volume-XXVII.pdf}}
\expandafter\def\csname r@cor:finiteoddtotal675\endcsname{{11.26}{23}{}{proposition.11.26}{Dirichlet-GRH-v930-Volume-II.pdf}}
\expandafter\def\csname r@cor:finitepositive749\endcsname{{15.20}{24}{}{proposition.15.20}{Dirichlet-GRH-v930-Volume-III.pdf}}
\expandafter\def\csname r@cor:finmark557\endcsname{{48.19}{121}{}{proposition.48.19}{Dirichlet-GRH-v930-Volume-IX.pdf}}
\expandafter\def\csname r@cor:firstaverage395\endcsname{{128.18}{71}{}{proposition.128.18}{Dirichlet-GRH-v930-Volume-XXIV.pdf}}
\expandafter\def\csname r@cor:firstglobal909\endcsname{{139.2}{94}{}{proposition.139.2}{Dirichlet-GRH-v930-Volume-XXVI.pdf}}
\expandafter\def\csname r@cor:firstliftrank711\endcsname{{142.16}{49}{}{proposition.142.16}{Dirichlet-GRH-v930-Volume-XXVII.pdf}}
\expandafter\def\csname r@cor:firstprojection708\endcsname{{142.6}{41}{}{proposition.142.6}{Dirichlet-GRH-v930-Volume-XXVII.pdf}}
\expandafter\def\csname r@cor:firstsample898\endcsname{{137.19}{62}{}{proposition.137.19}{Dirichlet-GRH-v930-Volume-XXVI.pdf}}
\expandafter\def\csname r@cor:firsttorus469\endcsname{{47.6}{81}{}{proposition.47.6}{Dirichlet-GRH-v930-Volume-IX.pdf}}
\expandafter\def\csname r@cor:fivepowers674\endcsname{{145.15}{92}{}{proposition.145.15}{Dirichlet-GRH-v930-Volume-XXVII.pdf}}
\expandafter\def\csname r@cor:fixedcollision457\endcsname{{55.35}{50}{}{proposition.55.35}{Dirichlet-GRH-v930-Volume-XI.pdf}}
\expandafter\def\csname r@cor:fixedhankel482\endcsname{{68.13}{95}{}{proposition.68.13}{Dirichlet-GRH-v930-Volume-XIII.pdf}}
\expandafter\def\csname r@cor:fixedheight479\endcsname{{143.9}{56}{}{proposition.143.9}{Dirichlet-GRH-v930-Volume-XXVII.pdf}}
\expandafter\def\csname r@cor:fixedimaginary539\endcsname{{22.29}{86}{}{proposition.22.29}{Dirichlet-GRH-v930-Volume-IV.pdf}}
\expandafter\def\csname r@cor:fixedinterval732\endcsname{{13.7}{66}{}{proposition.13.7}{Dirichlet-GRH-v930-Volume-II.pdf}}
\expandafter\def\csname r@cor:fixedminor486\endcsname{{68.24}{102}{}{proposition.68.24}{Dirichlet-GRH-v930-Volume-XIII.pdf}}
\expandafter\def\csname r@cor:fixedproductjets482\endcsname{{68.12}{95}{}{proposition.68.12}{Dirichlet-GRH-v930-Volume-XIII.pdf}}
\expandafter\def\csname r@cor:flagcost383\endcsname{{127.13}{39}{}{proposition.127.13}{Dirichlet-GRH-v930-Volume-XXIV.pdf}}
\expandafter\def\csname r@cor:flagcotangent438\endcsname{{55.17}{32}{}{proposition.55.17}{Dirichlet-GRH-v930-Volume-XI.pdf}}
\expandafter\def\csname r@cor:flagfinitered439\endcsname{{55.21}{37}{}{proposition.55.21}{Dirichlet-GRH-v930-Volume-XI.pdf}}
\expandafter\def\csname r@cor:flagmoments653\endcsname{{117.87}{120}{}{proposition.117.87}{Dirichlet-GRH-v930-Volume-XXII.pdf}}
\expandafter\def\csname r@cor:flagshifts520\endcsname{{146.13}{102}{}{proposition.146.13}{Dirichlet-GRH-v930-Volume-XXVII.pdf}}
\expandafter\def\csname r@cor:flatmean853\endcsname{{82.4}{29}{}{proposition.82.4}{Dirichlet-GRH-v930-Volume-XVI.pdf}}
\expandafter\def\csname r@cor:fourcoeff385\endcsname{{40.28}{26}{}{proposition.40.28}{Dirichlet-GRH-v930-Volume-VIII.pdf}}
\expandafter\def\csname r@cor:fouriermean354\endcsname{{120.9}{20}{}{proposition.120.9}{Dirichlet-GRH-v930-Volume-XXIII.pdf}}
\expandafter\def\csname r@cor:frameheight509\endcsname{{144.24}{79}{}{proposition.144.24}{Dirichlet-GRH-v930-Volume-XXVII.pdf}}
\expandafter\def\csname r@cor:frobeniusstrict399\endcsname{{41.6}{34}{}{proposition.41.6}{Dirichlet-GRH-v930-Volume-VIII.pdf}}
\expandafter\def\csname r@cor:frobideal470\endcsname{{47.9}{83}{}{proposition.47.9}{Dirichlet-GRH-v930-Volume-IX.pdf}}
\expandafter\def\csname r@cor:fullcenter452\endcsname{{141.16}{17}{}{proposition.141.16}{Dirichlet-GRH-v930-Volume-XXVII.pdf}}
\expandafter\def\csname r@cor:fullcollision502\endcsname{{56.18}{63}{}{proposition.56.18}{Dirichlet-GRH-v930-Volume-XI.pdf}}
\expandafter\def\csname r@cor:fullinput632\endcsname{{117.60}{96}{}{proposition.117.60}{Dirichlet-GRH-v930-Volume-XXII.pdf}}
\expandafter\def\csname r@cor:fullnonsplit466\endcsname{{56.12}{59}{}{proposition.56.12}{Dirichlet-GRH-v930-Volume-XI.pdf}}
\expandafter\def\csname r@cor:fullperiodmoments400\endcsname{{128.26}{77}{}{proposition.128.26}{Dirichlet-GRH-v930-Volume-XXIV.pdf}}
\expandafter\def\csname r@cor:fullreal529\endcsname{{57.29}{87}{}{proposition.57.29}{Dirichlet-GRH-v930-Volume-XI.pdf}}
\expandafter\def\csname r@cor:gausssum402\endcsname{{41.12}{40}{}{proposition.41.12}{Dirichlet-GRH-v930-Volume-VIII.pdf}}
\expandafter\def\csname r@cor:globalreducedheight628\endcsname{{96.43}{106}{}{proposition.96.43}{Dirichlet-GRH-v930-Volume-XVIII.pdf}}
\expandafter\def\csname r@cor:gluedperiod811\endcsname{{73.25}{97}{}{proposition.73.25}{Dirichlet-GRH-v930-Volume-XIV.pdf}}
\expandafter\def\csname r@cor:goodtest648\endcsname{{57.60}{114}{}{proposition.57.60}{Dirichlet-GRH-v930-Volume-XI.pdf}}
\expandafter\def\csname r@cor:gramerror622\endcsname{{117.44}{83}{}{proposition.117.44}{Dirichlet-GRH-v930-Volume-XXII.pdf}}
\expandafter\def\csname r@cor:gramsignal450\endcsname{{141.13}{14}{}{proposition.141.13}{Dirichlet-GRH-v930-Volume-XXVII.pdf}}
\expandafter\def\csname r@cor:halftail647\endcsname{{57.57}{111}{}{proposition.57.57}{Dirichlet-GRH-v930-Volume-XI.pdf}}
\expandafter\def\csname r@cor:heckebudget345\endcsname{{36.5}{4}{}{proposition.36.5}{Dirichlet-GRH-v930-Volume-VII.pdf}}
\expandafter\def\csname r@cor:heightfloor463\endcsname{{141.28}{26}{}{proposition.141.28}{Dirichlet-GRH-v930-Volume-XXVII.pdf}}
\expandafter\def\csname r@cor:heightoneplanes610\endcsname{{117.20}{57}{}{proposition.117.20}{Dirichlet-GRH-v930-Volume-XXII.pdf}}
\expandafter\def\csname r@cor:homogeneousrank498\endcsname{{144.9}{67}{}{proposition.144.9}{Dirichlet-GRH-v930-Volume-XXVII.pdf}}
\expandafter\def\csname r@cor:ideal37\endcsname{{64.3}{6}{}{proposition.64.3}{Dirichlet-GRH-v930-Volume-XIII.pdf}}
\expandafter\def\csname r@cor:identityblock703\endcsname{{133.11}{71}{}{proposition.133.11}{Dirichlet-GRH-v930-Volume-XXV.pdf}}
\expandafter\def\csname r@cor:identitynorm863\endcsname{{84.6}{81}{}{proposition.84.6}{Dirichlet-GRH-v930-Volume-XVI.pdf}}
\expandafter\def\csname r@cor:indepder494\endcsname{{68.40}{114}{}{proposition.68.40}{Dirichlet-GRH-v930-Volume-XIII.pdf}}
\expandafter\def\csname r@cor:indexperiod576\endcsname{{115.26}{23}{}{proposition.115.26}{Dirichlet-GRH-v930-Volume-XXII.pdf}}
\expandafter\def\csname r@cor:inertcorrection461\endcsname{{68.3}{89}{}{proposition.68.3}{Dirichlet-GRH-v930-Volume-XIII.pdf}}
\expandafter\def\csname r@cor:inertproduct905\endcsname{{138.17}{83}{}{proposition.138.17}{Dirichlet-GRH-v930-Volume-XXVI.pdf}}
\expandafter\def\csname r@cor:inertsum902\endcsname{{138.9}{74}{}{proposition.138.9}{Dirichlet-GRH-v930-Volume-XXVI.pdf}}
\expandafter\def\csname r@cor:inertterminal903\endcsname{{138.11}{76}{}{proposition.138.11}{Dirichlet-GRH-v930-Volume-XXVI.pdf}}
\expandafter\def\csname r@cor:integralinduction585\endcsname{{61.9}{55}{}{proposition.61.9}{Dirichlet-GRH-v930-Volume-XII.pdf}}
\expandafter\def\csname r@cor:integralmixed627\endcsname{{96.40}{104}{}{proposition.96.40}{Dirichlet-GRH-v930-Volume-XVIII.pdf}}
\expandafter\def\csname r@cor:integralranks623\endcsname{{117.47}{85}{}{proposition.117.47}{Dirichlet-GRH-v930-Volume-XXII.pdf}}
\expandafter\def\csname r@cor:integratedmark873\endcsname{{134.13}{98}{}{proposition.134.13}{Dirichlet-GRH-v930-Volume-XXV.pdf}}
\expandafter\def\csname r@cor:interval\endcsname{{30.13}{11}{}{proposition.30.13}{Dirichlet-GRH-v930-Volume-VI.pdf}}
\expandafter\def\csname r@cor:intervalbase416\endcsname{{44.8}{8}{}{proposition.44.8}{Dirichlet-GRH-v930-Volume-IX.pdf}}
\expandafter\def\csname r@cor:intervalintegers668\endcsname{{141.31}{28}{}{proposition.141.31}{Dirichlet-GRH-v930-Volume-XXVII.pdf}}
\expandafter\def\csname r@cor:intervalroots523\endcsname{{57.15}{78}{}{proposition.57.15}{Dirichlet-GRH-v930-Volume-XI.pdf}}
\expandafter\def\csname r@cor:intervalzero634\endcsname{{57.42}{99}{}{proposition.57.42}{Dirichlet-GRH-v930-Volume-XI.pdf}}
\expandafter\def\csname r@cor:irreduciblelabel374\endcsname{{38.8}{75}{}{proposition.38.8}{Dirichlet-GRH-v930-Volume-VII.pdf}}
\expandafter\def\csname r@cor:irreducibleweight423\endcsname{{44.40}{33}{}{proposition.44.40}{Dirichlet-GRH-v930-Volume-IX.pdf}}
\expandafter\def\csname r@cor:isolatedcost353\endcsname{{99.21}{61}{}{proposition.99.21}{Dirichlet-GRH-v930-Volume-XIX.pdf}}
\expandafter\def\csname r@cor:jet49\endcsname{{64.13}{14}{}{proposition.64.13}{Dirichlet-GRH-v930-Volume-XIII.pdf}}
\expandafter\def\csname r@cor:jetslength490\endcsname{{47.31}{97}{}{proposition.47.31}{Dirichlet-GRH-v930-Volume-IX.pdf}}
\expandafter\def\csname r@cor:jointimage355\endcsname{{127.4}{29}{}{proposition.127.4}{Dirichlet-GRH-v930-Volume-XXIV.pdf}}
\expandafter\def\csname r@cor:labelledtrace735\endcsname{{13.16}{75}{}{proposition.13.16}{Dirichlet-GRH-v930-Volume-II.pdf}}
\expandafter\def\csname r@cor:largeprimetable857\endcsname{{82.12}{40}{}{proposition.82.12}{Dirichlet-GRH-v930-Volume-XVI.pdf}}
\expandafter\def\csname r@cor:lastblock393\endcsname{{128.12}{65}{}{proposition.128.12}{Dirichlet-GRH-v930-Volume-XXIV.pdf}}
\expandafter\def\csname r@cor:levelcost725\endcsname{{45.24}{61}{}{proposition.45.24}{Dirichlet-GRH-v930-Volume-IX.pdf}}
\expandafter\def\csname r@cor:levelcurve370\endcsname{{37.29}{62}{}{proposition.37.29}{Dirichlet-GRH-v930-Volume-VII.pdf}}
\expandafter\def\csname r@cor:leveltwentyfive467\endcsname{{56.15}{60}{}{proposition.56.15}{Dirichlet-GRH-v930-Volume-XI.pdf}}
\expandafter\def\csname r@cor:localjetnorm487\endcsname{{68.28}{105}{}{proposition.68.28}{Dirichlet-GRH-v930-Volume-XIII.pdf}}
\expandafter\def\csname r@cor:maincost730\endcsname{{12.33}{59}{}{proposition.12.33}{Dirichlet-GRH-v930-Volume-II.pdf}}
\expandafter\def\csname r@cor:markbudget785\endcsname{{52.6}{63}{}{proposition.52.6}{Dirichlet-GRH-v930-Volume-X.pdf}}
\expandafter\def\csname r@cor:markedbessel679\endcsname{{130.4}{7}{}{proposition.130.4}{Dirichlet-GRH-v930-Volume-XXV.pdf}}
\expandafter\def\csname r@cor:markeddet715\endcsname{{47.37}{102}{}{proposition.47.37}{Dirichlet-GRH-v930-Volume-IX.pdf}}
\expandafter\def\csname r@cor:markedheight477\endcsname{{143.3}{51}{}{proposition.143.3}{Dirichlet-GRH-v930-Volume-XXVII.pdf}}
\expandafter\def\csname r@cor:markedlocal801\endcsname{{17.34}{120}{}{proposition.17.34}{Dirichlet-GRH-v930-Volume-III.pdf}}
\expandafter\def\csname r@cor:markednearisometry625\endcsname{{96.35}{99}{}{proposition.96.35}{Dirichlet-GRH-v930-Volume-XVIII.pdf}}
\expandafter\def\csname r@cor:markedpaths394\endcsname{{128.16}{69}{}{proposition.128.16}{Dirichlet-GRH-v930-Volume-XXIV.pdf}}
\expandafter\def\csname r@cor:markedperiod564\endcsname{{115.20}{19}{}{proposition.115.20}{Dirichlet-GRH-v930-Volume-XXII.pdf}}
\expandafter\def\csname r@cor:markedtests341\endcsname{{71.28}{43}{}{proposition.71.28}{Dirichlet-GRH-v930-Volume-XIV.pdf}}
\expandafter\def\csname r@cor:matchedlength617\endcsname{{22.45}{98}{}{proposition.22.45}{Dirichlet-GRH-v930-Volume-IV.pdf}}
\expandafter\def\csname r@cor:middleperiod570\endcsname{{115.23}{21}{}{proposition.115.23}{Dirichlet-GRH-v930-Volume-XXII.pdf}}
\expandafter\def\csname r@cor:midtangent333\endcsname{{34.26}{90}{}{proposition.34.26}{Dirichlet-GRH-v930-Volume-VI.pdf}}
\expandafter\def\csname r@cor:minor67\endcsname{{66.19}{60}{}{proposition.66.19}{Dirichlet-GRH-v930-Volume-XIII.pdf}}
\expandafter\def\csname r@cor:minor77\endcsname{{66.25}{65}{}{proposition.66.25}{Dirichlet-GRH-v930-Volume-XIII.pdf}}
\expandafter\def\csname r@cor:mixedchain418\endcsname{{44.16}{16}{}{proposition.44.16}{Dirichlet-GRH-v930-Volume-IX.pdf}}
\expandafter\def\csname r@cor:mixedhalf633\endcsname{{57.38}{96}{}{proposition.57.38}{Dirichlet-GRH-v930-Volume-XI.pdf}}
\expandafter\def\csname r@cor:mixedhalfexponent643\endcsname{{57.49}{104}{}{proposition.57.49}{Dirichlet-GRH-v930-Volume-XI.pdf}}
\expandafter\def\csname r@cor:mixedmatrix541\endcsname{{113.23}{102}{}{proposition.113.23}{Dirichlet-GRH-v930-Volume-XXI.pdf}}
\expandafter\def\csname r@cor:mixedphase441\endcsname{{128.33}{83}{}{proposition.128.33}{Dirichlet-GRH-v930-Volume-XXIV.pdf}}
\expandafter\def\csname r@cor:mixedrestrictions626\endcsname{{96.38}{101}{}{proposition.96.38}{Dirichlet-GRH-v930-Volume-XVIII.pdf}}
\expandafter\def\csname r@cor:movingcubic665\endcsname{{26.32}{71}{}{proposition.26.32}{Dirichlet-GRH-v930-Volume-V.pdf}}
\expandafter\def\csname r@cor:multipleaugmentations519\endcsname{{146.10}{100}{}{proposition.146.10}{Dirichlet-GRH-v930-Volume-XXVII.pdf}}
\expandafter\def\csname r@cor:negativejets797\endcsname{{17.26}{109}{}{proposition.17.26}{Dirichlet-GRH-v930-Volume-III.pdf}}
\expandafter\def\csname r@cor:negwedge398\endcsname{{41.3}{31}{}{proposition.41.3}{Dirichlet-GRH-v930-Volume-VIII.pdf}}
\expandafter\def\csname r@cor:nonzerotable856\endcsname{{82.9}{37}{}{proposition.82.9}{Dirichlet-GRH-v930-Volume-XVI.pdf}}
\expandafter\def\csname r@cor:normaldegree347\endcsname{{127.2}{27}{}{proposition.127.2}{Dirichlet-GRH-v930-Volume-XXIV.pdf}}
\expandafter\def\csname r@cor:normalizedtransfer722\endcsname{{45.17}{55}{}{proposition.45.17}{Dirichlet-GRH-v930-Volume-IX.pdf}}
\expandafter\def\csname r@cor:normalrange694\endcsname{{131.38}{50}{}{proposition.131.38}{Dirichlet-GRH-v930-Volume-XXV.pdf}}
\expandafter\def\csname r@cor:normaltranslate388\endcsname{{127.23}{52}{}{proposition.127.23}{Dirichlet-GRH-v930-Volume-XXIV.pdf}}
\expandafter\def\csname r@cor:octicspectrum543\endcsname{{57.34}{92}{}{proposition.57.34}{Dirichlet-GRH-v930-Volume-XI.pdf}}
\expandafter\def\csname r@cor:oddwindow677\endcsname{{11.30}{27}{}{proposition.11.30}{Dirichlet-GRH-v930-Volume-II.pdf}}
\expandafter\def\csname r@cor:oldlogperiod748\endcsname{{15.16}{21}{}{proposition.15.16}{Dirichlet-GRH-v930-Volume-III.pdf}}
\expandafter\def\csname r@cor:oldperiod752\endcsname{{15.29}{32}{}{proposition.15.29}{Dirichlet-GRH-v930-Volume-III.pdf}}
\expandafter\def\csname r@cor:onedeepweight464\endcsname{{56.9}{56}{}{proposition.56.9}{Dirichlet-GRH-v930-Volume-XI.pdf}}
\expandafter\def\csname r@cor:oneindexperiod603\endcsname{{117.4}{42}{}{proposition.117.4}{Dirichlet-GRH-v930-Volume-XXII.pdf}}
\expandafter\def\csname r@cor:openprofiles615\endcsname{{117.32}{70}{}{proposition.117.32}{Dirichlet-GRH-v930-Volume-XXII.pdf}}
\expandafter\def\csname r@cor:orbitconstants764\endcsname{{16.19}{51}{}{proposition.16.19}{Dirichlet-GRH-v930-Volume-III.pdf}}
\expandafter\def\csname r@cor:orbitideal606\endcsname{{117.12}{51}{}{proposition.117.12}{Dirichlet-GRH-v930-Volume-XXII.pdf}}
\expandafter\def\csname r@cor:orbitkernel500\endcsname{{144.12}{70}{}{proposition.144.12}{Dirichlet-GRH-v930-Volume-XXVII.pdf}}
\expandafter\def\csname r@cor:ordinarycurve598\endcsname{{62.13}{81}{}{proposition.62.13}{Dirichlet-GRH-v930-Volume-XII.pdf}}
\expandafter\def\csname r@cor:originalinterval600\endcsname{{22.32}{88}{}{proposition.22.32}{Dirichlet-GRH-v930-Volume-IV.pdf}}
\expandafter\def\csname r@cor:orth\endcsname{{70.9}{10}{}{proposition.70.9}{Dirichlet-GRH-v930-Volume-XIV.pdf}}
\expandafter\def\csname r@cor:orthometric491\endcsname{{143.15}{61}{}{proposition.143.15}{Dirichlet-GRH-v930-Volume-XXVII.pdf}}
\expandafter\def\csname r@cor:outsidecenter453\endcsname{{141.19}{19}{}{proposition.141.19}{Dirichlet-GRH-v930-Volume-XXVII.pdf}}
\expandafter\def\csname r@cor:packetsave504\endcsname{{144.18}{74}{}{proposition.144.18}{Dirichlet-GRH-v930-Volume-XXVII.pdf}}
\expandafter\def\csname r@cor:paddedminor814\endcsname{{77.12}{13}{}{proposition.77.12}{Dirichlet-GRH-v930-Volume-XV.pdf}}
\expandafter\def\csname r@cor:parabolicmass701\endcsname{{133.3}{66}{}{proposition.133.3}{Dirichlet-GRH-v930-Volume-XXV.pdf}}
\expandafter\def\csname r@cor:parabolicslice680\endcsname{{130.8}{11}{}{proposition.130.8}{Dirichlet-GRH-v930-Volume-XXV.pdf}}
\expandafter\def\csname r@cor:parallelind578\endcsname{{60.21}{40}{}{proposition.60.21}{Dirichlet-GRH-v930-Volume-XII.pdf}}
\expandafter\def\csname r@cor:parity43\endcsname{{64.9}{11}{}{proposition.64.9}{Dirichlet-GRH-v930-Volume-XIII.pdf}}
\expandafter\def\csname r@cor:pathperiod757\endcsname{{16.3}{37}{}{proposition.16.3}{Dirichlet-GRH-v930-Volume-III.pdf}}
\expandafter\def\csname r@cor:periodblock346\endcsname{{28.21}{104}{}{proposition.28.21}{Dirichlet-GRH-v930-Volume-V.pdf}}
\expandafter\def\csname r@cor:periodfieldweights644\endcsname{{57.52}{107}{}{proposition.57.52}{Dirichlet-GRH-v930-Volume-XI.pdf}}
\expandafter\def\csname r@cor:periodhalf529\endcsname{{57.27}{87}{}{proposition.57.27}{Dirichlet-GRH-v930-Volume-XI.pdf}}
\expandafter\def\csname r@cor:periodlarge425\endcsname{{54.4}{5}{}{proposition.54.4}{Dirichlet-GRH-v930-Volume-XI.pdf}}
\expandafter\def\csname r@cor:periodlifts759\endcsname{{16.6}{40}{}{proposition.16.6}{Dirichlet-GRH-v930-Volume-III.pdf}}
\expandafter\def\csname r@cor:periodrank707\endcsname{{142.3}{38}{}{proposition.142.3}{Dirichlet-GRH-v930-Volume-XXVII.pdf}}
\expandafter\def\csname r@cor:phasevaluation605\endcsname{{117.9}{48}{}{proposition.117.9}{Dirichlet-GRH-v930-Volume-XXII.pdf}}
\expandafter\def\csname r@cor:pivotoperator875\endcsname{{136.5}{6}{}{proposition.136.5}{Dirichlet-GRH-v930-Volume-XXVI.pdf}}
\expandafter\def\csname r@cor:planeinterval565\endcsname{{59.12}{13}{}{proposition.59.12}{Dirichlet-GRH-v930-Volume-XII.pdf}}
\expandafter\def\csname r@cor:planeweight566\endcsname{{59.15}{15}{}{proposition.59.15}{Dirichlet-GRH-v930-Volume-XII.pdf}}
\expandafter\def\csname r@cor:pointalternative550\endcsname{{48.15}{117}{}{proposition.48.15}{Dirichlet-GRH-v930-Volume-IX.pdf}}
\expandafter\def\csname r@cor:polarsplit593\endcsname{{61.19}{63}{}{proposition.61.19}{Dirichlet-GRH-v930-Volume-XII.pdf}}
\expandafter\def\csname r@cor:primeproduct906\endcsname{{138.20}{86}{}{proposition.138.20}{Dirichlet-GRH-v930-Volume-XXVI.pdf}}
\expandafter\def\csname r@cor:productjets481\endcsname{{68.9}{93}{}{proposition.68.9}{Dirichlet-GRH-v930-Volume-XIII.pdf}}
\expandafter\def\csname r@cor:projecteddiagonal540\endcsname{{113.20}{99}{}{proposition.113.20}{Dirichlet-GRH-v930-Volume-XXI.pdf}}
\expandafter\def\csname r@cor:projectedminors485\endcsname{{68.20}{100}{}{proposition.68.20}{Dirichlet-GRH-v930-Volume-XIII.pdf}}
\expandafter\def\csname r@cor:projectsymbol483\endcsname{{68.17}{97}{}{proposition.68.17}{Dirichlet-GRH-v930-Volume-XIII.pdf}}
\expandafter\def\csname r@cor:properperiod602\endcsname{{116.18}{39}{}{proposition.116.18}{Dirichlet-GRH-v930-Volume-XXII.pdf}}
\expandafter\def\csname r@cor:quadraticcomponent583\endcsname{{61.3}{50}{}{proposition.61.3}{Dirichlet-GRH-v930-Volume-XII.pdf}}
\expandafter\def\csname r@cor:quadraticcutoff587\endcsname{{61.14}{59}{}{proposition.61.14}{Dirichlet-GRH-v930-Volume-XII.pdf}}
\expandafter\def\csname r@cor:quadraticrange507\endcsname{{57.12}{75}{}{proposition.57.12}{Dirichlet-GRH-v930-Volume-XI.pdf}}
\expandafter\def\csname r@cor:quadraticrange609\endcsname{{22.42}{96}{}{proposition.22.42}{Dirichlet-GRH-v930-Volume-IV.pdf}}
\expandafter\def\csname r@cor:quotientchainheight734\endcsname{{13.14}{73}{}{proposition.13.14}{Dirichlet-GRH-v930-Volume-II.pdf}}
\expandafter\def\csname r@cor:rankobstruction337\endcsname{{83.22}{70}{}{proposition.83.22}{Dirichlet-GRH-v930-Volume-XVI.pdf}}
\expandafter\def\csname r@cor:rankonedual599\endcsname{{62.16}{84}{}{proposition.62.16}{Dirichlet-GRH-v930-Volume-XII.pdf}}
\expandafter\def\csname r@cor:rankpacking603\endcsname{{117.3}{42}{}{proposition.117.3}{Dirichlet-GRH-v930-Volume-XXII.pdf}}
\expandafter\def\csname r@cor:rationaldepth750\endcsname{{15.23}{26}{}{proposition.15.23}{Dirichlet-GRH-v930-Volume-III.pdf}}
\expandafter\def\csname r@cor:rationalfixed629\endcsname{{117.51}{89}{}{proposition.117.51}{Dirichlet-GRH-v930-Volume-XXII.pdf}}
\expandafter\def\csname r@cor:rationalperiod538\endcsname{{22.26}{83}{}{proposition.22.26}{Dirichlet-GRH-v930-Volume-IV.pdf}}
\expandafter\def\csname r@cor:rawfailure672\endcsname{{141.38}{35}{}{proposition.141.38}{Dirichlet-GRH-v930-Volume-XXVII.pdf}}
\expandafter\def\csname r@cor:rayinterval444\endcsname{{55.30}{46}{}{proposition.55.30}{Dirichlet-GRH-v930-Volume-XI.pdf}}
\expandafter\def\csname r@cor:raystable656\endcsname{{22.56}{107}{}{proposition.22.56}{Dirichlet-GRH-v930-Volume-IV.pdf}}
\expandafter\def\csname r@cor:raytamagawa726\endcsname{{12.21}{46}{}{proposition.12.21}{Dirichlet-GRH-v930-Volume-II.pdf}}
\expandafter\def\csname r@cor:real91\endcsname{{94.10}{45}{}{proposition.94.10}{Dirichlet-GRH-v930-Volume-XVIII.pdf}}
\expandafter\def\csname r@cor:realgramposition911\endcsname{{139.4}{96}{}{proposition.139.4}{Dirichlet-GRH-v930-Volume-XXVI.pdf}}
\expandafter\def\csname r@cor:realhighprime642\endcsname{{57.46}{102}{}{proposition.57.46}{Dirichlet-GRH-v930-Volume-XI.pdf}}
\expandafter\def\csname r@cor:realmarked682\endcsname{{131.2}{16}{}{proposition.131.2}{Dirichlet-GRH-v930-Volume-XXV.pdf}}
\expandafter\def\csname r@cor:realsystem927\endcsname{{139.32}{133}{}{proposition.139.32}{Dirichlet-GRH-v930-Volume-XXVI.pdf}}
\expandafter\def\csname r@cor:remaininggram807\endcsname{{73.12}{84}{}{proposition.73.12}{Dirichlet-GRH-v930-Volume-XIV.pdf}}
\expandafter\def\csname r@cor:remainingresidues426\endcsname{{54.8}{9}{}{proposition.54.8}{Dirichlet-GRH-v930-Volume-XI.pdf}}
\expandafter\def\csname r@cor:res28\endcsname{{24.10}{10}{}{proposition.24.10}{Dirichlet-GRH-v930-Volume-V.pdf}}
\expandafter\def\csname r@cor:residuecolumn492\endcsname{{68.34}{110}{}{proposition.68.34}{Dirichlet-GRH-v930-Volume-XIII.pdf}}
\expandafter\def\csname r@cor:rest104\endcsname{{19.17}{14}{}{proposition.19.17}{Dirichlet-GRH-v930-Volume-IV.pdf}}
\expandafter\def\csname r@cor:retainedtraces635\endcsname{{117.63}{99}{}{proposition.117.63}{Dirichlet-GRH-v930-Volume-XXII.pdf}}
\expandafter\def\csname r@cor:rotationroots612\endcsname{{117.27}{65}{}{proposition.117.27}{Dirichlet-GRH-v930-Volume-XXII.pdf}}
\expandafter\def\csname r@cor:rowhaar380\endcsname{{127.9}{34}{}{proposition.127.9}{Dirichlet-GRH-v930-Volume-XXIV.pdf}}
\expandafter\def\csname r@cor:rowperiod607\endcsname{{117.16}{54}{}{proposition.117.16}{Dirichlet-GRH-v930-Volume-XXII.pdf}}
\expandafter\def\csname r@cor:saturatedcontent448\endcsname{{141.10}{11}{}{proposition.141.10}{Dirichlet-GRH-v930-Volume-XXVII.pdf}}
\expandafter\def\csname r@cor:saturatedcotangent434\endcsname{{46.21}{76}{}{proposition.46.21}{Dirichlet-GRH-v930-Volume-IX.pdf}}
\expandafter\def\csname r@cor:scalarfl724\endcsname{{12.16}{42}{}{proposition.12.16}{Dirichlet-GRH-v930-Volume-II.pdf}}
\expandafter\def\csname r@cor:scalarfrob713\endcsname{{12.7}{35}{}{proposition.12.7}{Dirichlet-GRH-v930-Volume-II.pdf}}
\expandafter\def\csname r@cor:scalarpower806\endcsname{{73.9}{82}{}{proposition.73.9}{Dirichlet-GRH-v930-Volume-XIV.pdf}}
\expandafter\def\csname r@cor:scalarsum717\endcsname{{12.13}{40}{}{proposition.12.13}{Dirichlet-GRH-v930-Volume-II.pdf}}
\expandafter\def\csname r@cor:secondlatticedepth767\endcsname{{16.26}{57}{}{proposition.16.26}{Dirichlet-GRH-v930-Volume-III.pdf}}
\expandafter\def\csname r@cor:seedcomponent437\endcsname{{55.13}{29}{}{proposition.55.13}{Dirichlet-GRH-v930-Volume-XI.pdf}}
\expandafter\def\csname r@cor:seedcomponents595\endcsname{{62.5}{72}{}{proposition.62.5}{Dirichlet-GRH-v930-Volume-XII.pdf}}
\expandafter\def\csname r@cor:selmercongruence614\endcsname{{62.20}{87}{}{proposition.62.20}{Dirichlet-GRH-v930-Volume-XII.pdf}}
\expandafter\def\csname r@cor:selmergenerators410\endcsname{{42.11}{72}{}{proposition.42.11}{Dirichlet-GRH-v930-Volume-VIII.pdf}}
\expandafter\def\csname r@cor:severalflags526\endcsname{{96.4}{75}{}{proposition.96.4}{Dirichlet-GRH-v930-Volume-XVIII.pdf}}
\expandafter\def\csname r@cor:sexticmoving669\endcsname{{26.35}{74}{}{proposition.26.35}{Dirichlet-GRH-v930-Volume-V.pdf}}
\expandafter\def\csname r@cor:shortscalar763\endcsname{{51.3}{40}{}{proposition.51.3}{Dirichlet-GRH-v930-Volume-X.pdf}}
\expandafter\def\csname r@cor:siegeltrace431\endcsname{{46.13}{70}{}{proposition.46.13}{Dirichlet-GRH-v930-Volume-IX.pdf}}
\expandafter\def\csname r@cor:signedfactor513\endcsname{{145.13}{90}{}{proposition.145.13}{Dirichlet-GRH-v930-Volume-XXVII.pdf}}
\expandafter\def\csname r@cor:signedfirst828\endcsname{{78.5}{59}{}{proposition.78.5}{Dirichlet-GRH-v930-Volume-XV.pdf}}
\expandafter\def\csname r@cor:signedprecision516\endcsname{{146.3}{94}{}{proposition.146.3}{Dirichlet-GRH-v930-Volume-XXVII.pdf}}
\expandafter\def\csname r@cor:signedslice604\endcsname{{117.7}{45}{}{proposition.117.7}{Dirichlet-GRH-v930-Volume-XXII.pdf}}
\expandafter\def\csname r@cor:signedspan478\endcsname{{143.6}{54}{}{proposition.143.6}{Dirichlet-GRH-v930-Volume-XXVII.pdf}}
\expandafter\def\csname r@cor:signedterminal390\endcsname{{128.3}{58}{}{proposition.128.3}{Dirichlet-GRH-v930-Volume-XXIV.pdf}}
\expandafter\def\csname r@cor:singularselected874\endcsname{{136.2}{4}{}{proposition.136.2}{Dirichlet-GRH-v930-Volume-XXVI.pdf}}
\expandafter\def\csname r@cor:skewsize363\endcsname{{36.27}{31}{}{proposition.36.27}{Dirichlet-GRH-v930-Volume-VII.pdf}}
\expandafter\def\csname r@cor:smallgroups590\endcsname{{116.10}{32}{}{proposition.116.10}{Dirichlet-GRH-v930-Volume-XXII.pdf}}
\expandafter\def\csname r@cor:smallorder427\endcsname{{54.14}{14}{}{proposition.54.14}{Dirichlet-GRH-v930-Volume-XI.pdf}}
\expandafter\def\csname r@cor:smallresidual515\endcsname{{22.23}{81}{}{proposition.22.23}{Dirichlet-GRH-v930-Volume-IV.pdf}}
\expandafter\def\csname r@cor:smalltransverse589\endcsname{{116.7}{30}{}{proposition.116.7}{Dirichlet-GRH-v930-Volume-XXII.pdf}}
\expandafter\def\csname r@cor:sphericalcupbound762\endcsname{{16.16}{48}{}{proposition.16.16}{Dirichlet-GRH-v930-Volume-III.pdf}}
\expandafter\def\csname r@cor:splitcost456\endcsname{{120.18}{29}{}{proposition.120.18}{Dirichlet-GRH-v930-Volume-XXIII.pdf}}
\expandafter\def\csname r@cor:splitdet475\endcsname{{47.21}{91}{}{proposition.47.21}{Dirichlet-GRH-v930-Volume-IX.pdf}}
\expandafter\def\csname r@cor:splitlayers646\endcsname{{57.54}{109}{}{proposition.57.54}{Dirichlet-GRH-v930-Volume-XI.pdf}}
\expandafter\def\csname r@cor:splitloss574\endcsname{{60.12}{32}{}{proposition.60.12}{Dirichlet-GRH-v930-Volume-XII.pdf}}
\expandafter\def\csname r@cor:splitperiod858\endcsname{{82.15}{42}{}{proposition.82.15}{Dirichlet-GRH-v930-Volume-XVI.pdf}}
\expandafter\def\csname r@cor:splitrank480\endcsname{{68.6}{91}{}{proposition.68.6}{Dirichlet-GRH-v930-Volume-XIII.pdf}}
\expandafter\def\csname r@cor:splitray663\endcsname{{26.28}{68}{}{proposition.26.28}{Dirichlet-GRH-v930-Volume-V.pdf}}
\expandafter\def\csname r@cor:standardall908\endcsname{{138.25}{92}{}{proposition.138.25}{Dirichlet-GRH-v930-Volume-XXVI.pdf}}
\expandafter\def\csname r@cor:standardeuler899\endcsname{{137.22}{66}{}{proposition.137.22}{Dirichlet-GRH-v930-Volume-XXVI.pdf}}
\expandafter\def\csname r@cor:strictdet471\endcsname{{47.12}{85}{}{proposition.47.12}{Dirichlet-GRH-v930-Volume-IX.pdf}}
\expandafter\def\csname r@cor:strictmindet484\endcsname{{47.28}{95}{}{proposition.47.28}{Dirichlet-GRH-v930-Volume-IX.pdf}}
\expandafter\def\csname r@cor:strictminors415\endcsname{{44.4}{5}{}{proposition.44.4}{Dirichlet-GRH-v930-Volume-IX.pdf}}
\expandafter\def\csname r@cor:strictvertical499\endcsname{{47.34}{100}{}{proposition.47.34}{Dirichlet-GRH-v930-Volume-IX.pdf}}
\expandafter\def\csname r@cor:sum31\endcsname{{30.2}{4}{}{proposition.30.2}{Dirichlet-GRH-v930-Volume-VI.pdf}}
\expandafter\def\csname r@cor:sumconductor655\endcsname{{62.29}{94}{}{proposition.62.29}{Dirichlet-GRH-v930-Volume-XII.pdf}}
\expandafter\def\csname r@cor:sympcomponent468\endcsname{{47.3}{78}{}{proposition.47.3}{Dirichlet-GRH-v930-Volume-IX.pdf}}
\expandafter\def\csname r@cor:tamebudget800\endcsname{{17.31}{117}{}{proposition.17.31}{Dirichlet-GRH-v930-Volume-III.pdf}}
\expandafter\def\csname r@cor:tamelayers618\endcsname{{22.49}{101}{}{proposition.22.49}{Dirichlet-GRH-v930-Volume-IV.pdf}}
\expandafter\def\csname r@cor:terminalamplitudes640\endcsname{{117.78}{112}{}{proposition.117.78}{Dirichlet-GRH-v930-Volume-XXII.pdf}}
\expandafter\def\csname r@cor:terminalcomparison851\endcsname{{81.21}{24}{}{proposition.81.21}{Dirichlet-GRH-v930-Volume-XVI.pdf}}
\expandafter\def\csname r@cor:terminalfixer878\endcsname{{136.13}{14}{}{proposition.136.13}{Dirichlet-GRH-v930-Volume-XXVI.pdf}}
\expandafter\def\csname r@cor:terminalgram518\endcsname{{146.6}{97}{}{proposition.146.6}{Dirichlet-GRH-v930-Volume-XXVII.pdf}}
\expandafter\def\csname r@cor:terminalheight497\endcsname{{144.6}{66}{}{proposition.144.6}{Dirichlet-GRH-v930-Volume-XXVII.pdf}}
\expandafter\def\csname r@cor:terminalidentity639\endcsname{{117.75}{109}{}{proposition.117.75}{Dirichlet-GRH-v930-Volume-XXII.pdf}}
\expandafter\def\csname r@cor:terminalprefix582\endcsname{{116.3}{28}{}{proposition.116.3}{Dirichlet-GRH-v930-Volume-XXII.pdf}}
\expandafter\def\csname r@cor:terminalrow455\endcsname{{120.15}{26}{}{proposition.120.15}{Dirichlet-GRH-v930-Volume-XXIII.pdf}}
\expandafter\def\csname r@cor:thirdnormratio840\endcsname{{79.14}{92}{}{proposition.79.14}{Dirichlet-GRH-v930-Volume-XV.pdf}}
\expandafter\def\csname r@cor:thirdrealmatrix841\endcsname{{79.16}{94}{}{proposition.79.16}{Dirichlet-GRH-v930-Volume-XV.pdf}}
\expandafter\def\csname r@cor:thirdseed838\endcsname{{79.10}{86}{}{proposition.79.10}{Dirichlet-GRH-v930-Volume-XV.pdf}}
\expandafter\def\csname r@cor:threeextensions761\endcsname{{16.12}{44}{}{proposition.16.12}{Dirichlet-GRH-v930-Volume-III.pdf}}
\expandafter\def\csname r@cor:threeminor447\endcsname{{141.7}{8}{}{proposition.141.7}{Dirichlet-GRH-v930-Volume-XXVII.pdf}}
\expandafter\def\csname r@cor:torsionimage378\endcsname{{40.8}{7}{}{proposition.40.8}{Dirichlet-GRH-v930-Volume-VIII.pdf}}
\expandafter\def\csname r@cor:torustests741\endcsname{{15.6}{9}{}{proposition.15.6}{Dirichlet-GRH-v930-Volume-III.pdf}}
\expandafter\def\csname r@cor:torusvaluation852\endcsname{{82.2}{26}{}{proposition.82.2}{Dirichlet-GRH-v930-Volume-XVI.pdf}}
\expandafter\def\csname r@cor:tower27\endcsname{{19.11}{9}{}{proposition.19.11}{Dirichlet-GRH-v930-Volume-IV.pdf}}
\expandafter\def\csname r@cor:towerdescent571\endcsname{{60.3}{26}{}{proposition.60.3}{Dirichlet-GRH-v930-Volume-XII.pdf}}
\expandafter\def\csname r@cor:transfer100\endcsname{{11.18}{17}{}{proposition.11.18}{Dirichlet-GRH-v930-Volume-II.pdf}}
\expandafter\def\csname r@cor:transverseplane567\endcsname{{59.18}{18}{}{proposition.59.18}{Dirichlet-GRH-v930-Volume-XII.pdf}}
\expandafter\def\csname r@cor:transversetrace581\endcsname{{115.29}{26}{}{proposition.115.29}{Dirichlet-GRH-v930-Volume-XXII.pdf}}
\expandafter\def\csname r@cor:twointegral377\endcsname{{40.4}{4}{}{proposition.40.4}{Dirichlet-GRH-v930-Volume-VIII.pdf}}
\expandafter\def\csname r@cor:twoperiod754\endcsname{{50.22}{28}{}{proposition.50.22}{Dirichlet-GRH-v930-Volume-X.pdf}}
\expandafter\def\csname r@cor:twoplane474\endcsname{{47.18}{89}{}{proposition.47.18}{Dirichlet-GRH-v930-Volume-IX.pdf}}
\expandafter\def\csname r@cor:typedbudget799\endcsname{{52.14}{72}{}{proposition.52.14}{Dirichlet-GRH-v930-Volume-X.pdf}}
\expandafter\def\csname r@cor:unboundedprimes428\endcsname{{54.19}{18}{}{proposition.54.19}{Dirichlet-GRH-v930-Volume-XI.pdf}}
\expandafter\def\csname r@cor:undistinguished435\endcsname{{55.4}{22}{}{proposition.55.4}{Dirichlet-GRH-v930-Volume-XI.pdf}}
\expandafter\def\csname r@cor:uniformstrict561\endcsname{{59.6}{7}{}{proposition.59.6}{Dirichlet-GRH-v930-Volume-XII.pdf}}
\expandafter\def\csname r@cor:unitcap505\endcsname{{57.6}{70}{}{proposition.57.6}{Dirichlet-GRH-v930-Volume-XI.pdf}}
\expandafter\def\csname r@cor:unitrows662\endcsname{{26.25}{66}{}{proposition.26.25}{Dirichlet-GRH-v930-Volume-V.pdf}}
\expandafter\def\csname r@cor:unorderedcutoff588\endcsname{{61.17}{61}{}{proposition.61.17}{Dirichlet-GRH-v930-Volume-XII.pdf}}
\expandafter\def\csname r@cor:unusedposition885\endcsname{{136.29}{31}{}{proposition.136.29}{Dirichlet-GRH-v930-Volume-XXVI.pdf}}
\expandafter\def\csname r@cor:verticaltrace654\endcsname{{62.25}{91}{}{proposition.62.25}{Dirichlet-GRH-v930-Volume-XII.pdf}}
\expandafter\def\csname r@cor:weightedkernel389\endcsname{{127.26}{54}{}{proposition.127.26}{Dirichlet-GRH-v930-Volume-XXIV.pdf}}
\expandafter\def\csname r@cor:whittakerterminal527\endcsname{{96.8}{78}{}{proposition.96.8}{Dirichlet-GRH-v930-Volume-XVIII.pdf}}
\expandafter\def\csname r@cor:zeromass506\endcsname{{57.8}{72}{}{proposition.57.8}{Dirichlet-GRH-v930-Volume-XI.pdf}}
\expandafter\def\csname r@eq:A100\endcsname{{11.3}{16}{}{equation.11.3}{Dirichlet-GRH-v930-Volume-II.pdf}}
\expandafter\def\csname r@eq:cp73\endcsname{{66.1}{46}{}{equation.66.1}{Dirichlet-GRH-v930-Volume-XIII.pdf}}
\expandafter\def\csname r@eq:degrees73\endcsname{{66.2}{47}{}{equation.66.2}{Dirichlet-GRH-v930-Volume-XIII.pdf}}
\expandafter\def\csname r@eq:mu106\endcsname{{86.15}{22}{}{equation.86.15}{Dirichlet-GRH-v930-Volume-XVII.pdf}}
\expandafter\def\csname r@eq:pub104\endcsname{{19.9}{13}{}{equation.19.9}{Dirichlet-GRH-v930-Volume-IV.pdf}}
\expandafter\def\csname r@eq:region198\endcsname{{20.12}{33}{}{equation.20.12}{Dirichlet-GRH-v930-Volume-IV.pdf}}
\expandafter\def\csname r@eq:rest104\endcsname{{19.11}{14}{}{equation.19.11}{Dirichlet-GRH-v930-Volume-IV.pdf}}
\expandafter\def\csname r@eq:src109\endcsname{{87.4}{37}{}{equation.87.4}{Dirichlet-GRH-v930-Volume-XVII.pdf}}
\expandafter\def\csname r@eq:strictdigits329\endcsname{{22.1}{66}{}{equation.22.1}{Dirichlet-GRH-v930-Volume-IV.pdf}}
\expandafter\def\csname r@eq:sub81\endcsname{{66.22}{67}{}{equation.66.22}{Dirichlet-GRH-v930-Volume-XIII.pdf}}
\expandafter\def\csname r@eq:vor104\endcsname{{19.10}{13}{}{equation.19.10}{Dirichlet-GRH-v930-Volume-IV.pdf}}
\expandafter\def\csname r@lem:basis63\endcsname{{65.14}{34}{}{proposition.65.14}{Dirichlet-GRH-v930-Volume-XIII.pdf}}
\expandafter\def\csname r@lem:bilinearsource710\endcsname{{142.11}{44}{}{proposition.142.11}{Dirichlet-GRH-v930-Volume-XXVII.pdf}}
\expandafter\def\csname r@lem:cofactormark581\endcsname{{115.27}{24}{}{proposition.115.27}{Dirichlet-GRH-v930-Volume-XXII.pdf}}
\expandafter\def\csname r@lem:commondirection913\endcsname{{139.8}{100}{}{proposition.139.8}{Dirichlet-GRH-v930-Volume-XXVI.pdf}}
\expandafter\def\csname r@lem:des28\endcsname{{24.8}{8}{}{proposition.24.8}{Dirichlet-GRH-v930-Volume-V.pdf}}
\expandafter\def\csname r@lem:dual241\endcsname{{36.2}{2}{}{proposition.36.2}{Dirichlet-GRH-v930-Volume-VII.pdf}}
\expandafter\def\csname r@lem:excessmass671\endcsname{{141.35}{32}{}{proposition.141.35}{Dirichlet-GRH-v930-Volume-XXVII.pdf}}
\expandafter\def\csname r@lem:factorial41\endcsname{{64.4}{7}{}{proposition.64.4}{Dirichlet-GRH-v930-Volume-XIII.pdf}}
\expandafter\def\csname r@lem:finiteunitdensity930\endcsname{{139.37}{139}{}{proposition.139.37}{Dirichlet-GRH-v930-Volume-XXVI.pdf}}
\expandafter\def\csname r@lem:fit13\endcsname{{11.2}{3}{}{proposition.11.2}{Dirichlet-GRH-v930-Volume-II.pdf}}
\expandafter\def\csname r@lem:flagsign658\endcsname{{117.88}{120}{}{proposition.117.88}{Dirichlet-GRH-v930-Volume-XXII.pdf}}
\expandafter\def\csname r@lem:frameend512\endcsname{{145.8}{86}{}{proposition.145.8}{Dirichlet-GRH-v930-Volume-XXVII.pdf}}
\expandafter\def\csname r@lem:inertmiddle902\endcsname{{138.6}{72}{}{proposition.138.6}{Dirichlet-GRH-v930-Volume-XXVI.pdf}}
\expandafter\def\csname r@lem:int13\endcsname{{11.3}{4}{}{proposition.11.3}{Dirichlet-GRH-v930-Volume-II.pdf}}
\expandafter\def\csname r@lem:mass209\endcsname{{93.7}{26}{}{proposition.93.7}{Dirichlet-GRH-v930-Volume-XVIII.pdf}}
\expandafter\def\csname r@lem:primitiveclass502\endcsname{{56.16}{61}{}{proposition.56.16}{Dirichlet-GRH-v930-Volume-XI.pdf}}
\expandafter\def\csname r@lem:scalarden562\endcsname{{115.16}{16}{}{proposition.115.16}{Dirichlet-GRH-v930-Volume-XXII.pdf}}
\expandafter\def\csname r@lem:signalsupport452\endcsname{{141.14}{15}{}{proposition.141.14}{Dirichlet-GRH-v930-Volume-XXVII.pdf}}
\expandafter\def\csname r@lem:size97\endcsname{{87.1}{33}{}{proposition.87.1}{Dirichlet-GRH-v930-Volume-XVII.pdf}}
\expandafter\def\csname r@lem:terminalblock496\endcsname{{144.1}{61}{}{proposition.144.1}{Dirichlet-GRH-v930-Volume-XXVII.pdf}}
\expandafter\def\csname r@lem:terminalintegral867\endcsname{{84.13}{90}{}{proposition.84.13}{Dirichlet-GRH-v930-Volume-XVI.pdf}}
\expandafter\def\csname r@lem:terminalrotation694\endcsname{{131.36}{48}{}{proposition.131.36}{Dirichlet-GRH-v930-Volume-XXV.pdf}}
\expandafter\def\csname r@lem:unit41\endcsname{{64.5}{7}{}{proposition.64.5}{Dirichlet-GRH-v930-Volume-XIII.pdf}}
\expandafter\def\csname r@lem:very43\endcsname{{64.7}{10}{}{proposition.64.7}{Dirichlet-GRH-v930-Volume-XIII.pdf}}
\expandafter\def\csname r@nav:I-37\endcsname{{70.1}{3}{}{subsection.70.1}{Dirichlet-GRH-v930-Volume-XIV.pdf}}
\expandafter\def\csname r@nav:I-39\endcsname{{70.3}{6}{}{subsection.70.3}{Dirichlet-GRH-v930-Volume-XIV.pdf}}
\expandafter\def\csname r@nav:II-15\endcsname{{19.3}{6}{}{subsection.19.3}{Dirichlet-GRH-v930-Volume-IV.pdf}}
\expandafter\def\csname r@nav:IV-09\endcsname{{32.1}{22}{}{subsection.32.1}{Dirichlet-GRH-v930-Volume-VI.pdf}}
\expandafter\def\csname r@nav:IV-11\endcsname{{32.2}{26}{}{subsection.32.2}{Dirichlet-GRH-v930-Volume-VI.pdf}}
\expandafter\def\csname r@nav:IV-12\endcsname{{32.3}{28}{}{subsection.32.3}{Dirichlet-GRH-v930-Volume-VI.pdf}}
\expandafter\def\csname r@nav:V-00\endcsname{{64.1}{3}{}{subsection.64.1}{Dirichlet-GRH-v930-Volume-XIII.pdf}}
\expandafter\def\csname r@nav:V-05\endcsname{{64.3}{10}{}{subsection.64.3}{Dirichlet-GRH-v930-Volume-XIII.pdf}}
\expandafter\def\csname r@nav:V-06\endcsname{{92.1}{3}{}{subsection.92.1}{Dirichlet-GRH-v930-Volume-XVIII.pdf}}
\expandafter\def\csname r@nav:V-17\endcsname{{66.1}{46}{}{subsection.66.1}{Dirichlet-GRH-v930-Volume-XIII.pdf}}
\expandafter\def\csname r@nav:V-18\endcsname{{66.2}{49}{}{subsection.66.2}{Dirichlet-GRH-v930-Volume-XIII.pdf}}
\expandafter\def\csname r@nav:VI-01\endcsname{{66.9}{64}{}{subsection.66.9}{Dirichlet-GRH-v930-Volume-XIII.pdf}}
\expandafter\def\csname r@nav:bibliography-II\endcsname{{14}{80}{}{section.14}{Dirichlet-GRH-v930-Volume-II.pdf}}
\expandafter\def\csname r@nav:bibliography-III\endcsname{{18}{121}{}{section.18}{Dirichlet-GRH-v930-Volume-III.pdf}}
\expandafter\def\csname r@nav:bibliography-IV\endcsname{{23}{111}{}{section.23}{Dirichlet-GRH-v930-Volume-IV.pdf}}
\expandafter\def\csname r@nav:bibliography-IX\endcsname{{49}{126}{}{section.49}{Dirichlet-GRH-v930-Volume-IX.pdf}}
\expandafter\def\csname r@nav:bibliography-V\endcsname{{29}{106}{}{section.29}{Dirichlet-GRH-v930-Volume-V.pdf}}
\expandafter\def\csname r@nav:bibliography-VI\endcsname{{35}{91}{}{section.35}{Dirichlet-GRH-v930-Volume-VI.pdf}}
\expandafter\def\csname r@nav:bibliography-VII\endcsname{{39}{83}{}{section.39}{Dirichlet-GRH-v930-Volume-VII.pdf}}
\expandafter\def\csname r@nav:bibliography-VIII\endcsname{{43}{85}{}{section.43}{Dirichlet-GRH-v930-Volume-VIII.pdf}}
\expandafter\def\csname r@nav:bibliography-X\endcsname{{53}{84}{}{section.53}{Dirichlet-GRH-v930-Volume-X.pdf}}
\expandafter\def\csname r@nav:bibliography-XI\endcsname{{58}{118}{}{section.58}{Dirichlet-GRH-v930-Volume-XI.pdf}}
\expandafter\def\csname r@nav:bibliography-XII\endcsname{{63}{96}{}{section.63}{Dirichlet-GRH-v930-Volume-XII.pdf}}
\expandafter\def\csname r@nav:bibliography-XIII\endcsname{{69}{117}{}{section.69}{Dirichlet-GRH-v930-Volume-XIII.pdf}}
\expandafter\def\csname r@nav:bibliography-XIV\endcsname{{76}{122}{}{section.76}{Dirichlet-GRH-v930-Volume-XIV.pdf}}
\expandafter\def\csname r@nav:bibliography-XIX\endcsname{{103}{98}{}{section.103}{Dirichlet-GRH-v930-Volume-XIX.pdf}}
\expandafter\def\csname r@nav:bibliography-XV\endcsname{{80}{102}{}{section.80}{Dirichlet-GRH-v930-Volume-XV.pdf}}
\expandafter\def\csname r@nav:bibliography-XVI\endcsname{{85}{93}{}{section.85}{Dirichlet-GRH-v930-Volume-XVI.pdf}}
\expandafter\def\csname r@nav:bibliography-XVII\endcsname{{91}{100}{}{section.91}{Dirichlet-GRH-v930-Volume-XVII.pdf}}
\expandafter\def\csname r@nav:bibliography-XVIII\endcsname{{97}{108}{}{section.97}{Dirichlet-GRH-v930-Volume-XVIII.pdf}}
\expandafter\def\csname r@nav:bibliography-XX\endcsname{{108}{92}{}{section.108}{Dirichlet-GRH-v930-Volume-XX.pdf}}
\expandafter\def\csname r@nav:bibliography-XXI\endcsname{{114}{108}{}{section.114}{Dirichlet-GRH-v930-Volume-XXI.pdf}}
\expandafter\def\csname r@nav:bibliography-XXII\endcsname{{118}{130}{}{section.118}{Dirichlet-GRH-v930-Volume-XXII.pdf}}
\expandafter\def\csname r@nav:bibliography-XXIII\endcsname{{124}{100}{}{section.124}{Dirichlet-GRH-v930-Volume-XXIII.pdf}}
\expandafter\def\csname r@nav:bibliography-XXIV\endcsname{{129}{84}{}{section.129}{Dirichlet-GRH-v930-Volume-XXIV.pdf}}
\expandafter\def\csname r@nav:bibliography-XXV\endcsname{{135}{99}{}{section.135}{Dirichlet-GRH-v930-Volume-XXV.pdf}}
\expandafter\def\csname r@nav:bibliography-XXVI\endcsname{{140}{142}{}{section.140}{Dirichlet-GRH-v930-Volume-XXVI.pdf}}
\expandafter\def\csname r@nav:bibliography-XXVII\endcsname{{147}{109}{}{section.147}{Dirichlet-GRH-v930-Volume-XXVII.pdf}}
\expandafter\def\csname r@nav:carrier-integral\endcsname{{77}{3}{}{section.77}{Dirichlet-GRH-v930-Volume-XV.pdf}}
\expandafter\def\csname r@nav:chapter-coefficient-weyl-lattices\endcsname{{120}{14}{}{section.120}{Dirichlet-GRH-v930-Volume-XXIII.pdf}}
\expandafter\def\csname r@nav:chapter-coincident-characters\endcsname{{55}{18}{}{section.55}{Dirichlet-GRH-v930-Volume-XI.pdf}}
\expandafter\def\csname r@nav:chapter-companions-even-planes\endcsname{{59}{3}{}{section.59}{Dirichlet-GRH-v930-Volume-XII.pdf}}
\expandafter\def\csname r@nav:chapter-descent-strict-orders\endcsname{{41}{30}{}{section.41}{Dirichlet-GRH-v930-Volume-VIII.pdf}}
\expandafter\def\csname r@nav:chapter-dyadic-exceptional-weights\endcsname{{56}{51}{}{section.56}{Dirichlet-GRH-v930-Volume-XI.pdf}}
\expandafter\def\csname r@nav:chapter-dyadic-units-regulators\endcsname{{57}{66}{}{section.57}{Dirichlet-GRH-v930-Volume-XI.pdf}}
\expandafter\def\csname r@nav:chapter-fixed-weight-hecke\endcsname{{44}{3}{}{section.44}{Dirichlet-GRH-v930-Volume-IX.pdf}}
\expandafter\def\csname r@nav:chapter-galois-strict-frobenius\endcsname{{47}{76}{}{section.47}{Dirichlet-GRH-v930-Volume-IX.pdf}}
\expandafter\def\csname r@nav:chapter-integral-trace-quotients\endcsname{{40}{2}{}{section.40}{Dirichlet-GRH-v930-Volume-VIII.pdf}}
\expandafter\def\csname r@nav:chapter-marked-ordinary-traces\endcsname{{113}{84}{}{section.113}{Dirichlet-GRH-v930-Volume-XXI.pdf}}
\expandafter\def\csname r@nav:chapter-marked-period-determinants\endcsname{{143}{50}{}{section.143}{Dirichlet-GRH-v930-Volume-XXVII.pdf}}
\expandafter\def\csname r@nav:chapter-marked-theta-columns\endcsname{{127}{26}{}{section.127}{Dirichlet-GRH-v930-Volume-XXIV.pdf}}
\expandafter\def\csname r@nav:chapter-norms-and-period-ideals\endcsname{{141}{3}{}{section.141}{Dirichlet-GRH-v930-Volume-XXVII.pdf}}
\expandafter\def\csname r@nav:chapter-ordinary-families\endcsname{{37}{32}{}{section.37}{Dirichlet-GRH-v930-Volume-VII.pdf}}
\expandafter\def\csname r@nav:chapter-ordinary-hecke\endcsname{{42}{62}{}{section.42}{Dirichlet-GRH-v930-Volume-VIII.pdf}}
\expandafter\def\csname r@nav:chapter-ordinary-hecke-fibers\endcsname{{62}{64}{}{section.62}{Dirichlet-GRH-v930-Volume-XII.pdf}}
\expandafter\def\csname r@nav:chapter-ordinary-symplectic-equations\endcsname{{48}{104}{}{section.48}{Dirichlet-GRH-v930-Volume-IX.pdf}}
\expandafter\def\csname r@nav:chapter-ordinary-terminal-lattices\endcsname{{115}{3}{}{section.115}{Dirichlet-GRH-v930-Volume-XXII.pdf}}
\expandafter\def\csname r@nav:chapter-projected-period-divisors\endcsname{{145}{80}{}{section.145}{Dirichlet-GRH-v930-Volume-XXVII.pdf}}
\expandafter\def\csname r@nav:chapter-quadratic-induction-symmetric-cubes\endcsname{{61}{48}{}{section.61}{Dirichlet-GRH-v930-Volume-XII.pdf}}
\expandafter\def\csname r@nav:chapter-refined-rational-columns\endcsname{{119}{3}{}{section.119}{Dirichlet-GRH-v930-Volume-XXIII.pdf}}
\expandafter\def\csname r@nav:chapter-signed-terminal-frequencies\endcsname{{146}{92}{}{section.146}{Dirichlet-GRH-v930-Volume-XXVII.pdf}}
\expandafter\def\csname r@nav:chapter-signed-terminal-tests\endcsname{{126}{8}{}{section.126}{Dirichlet-GRH-v930-Volume-XXIV.pdf}}
\expandafter\def\csname r@nav:chapter-signed-transverse-integrals\endcsname{{117}{40}{}{section.117}{Dirichlet-GRH-v930-Volume-XXII.pdf}}
\expandafter\def\csname r@nav:chapter-split-elliptic-jets\endcsname{{68}{87}{}{section.68}{Dirichlet-GRH-v930-Volume-XIII.pdf}}
\expandafter\def\csname r@nav:chapter-symplectic-classes\endcsname{{46}{63}{}{section.46}{Dirichlet-GRH-v930-Volume-IX.pdf}}
\expandafter\def\csname r@nav:chapter-terminal-covectors\endcsname{{128}{55}{}{section.128}{Dirichlet-GRH-v930-Volume-XXIV.pdf}}
\expandafter\def\csname r@nav:chapter-terminal-fourier-heights\endcsname{{144}{61}{}{section.144}{Dirichlet-GRH-v930-Volume-XXVII.pdf}}
\expandafter\def\csname r@nav:chapter-three-character-traces\endcsname{{38}{70}{}{section.38}{Dirichlet-GRH-v930-Volume-VII.pdf}}
\expandafter\def\csname r@nav:chapter-towers-cuspidal-points\endcsname{{60}{24}{}{section.60}{Dirichlet-GRH-v930-Volume-XII.pdf}}
\expandafter\def\csname r@nav:chapter-transverse-determinants\endcsname{{116}{26}{}{section.116}{Dirichlet-GRH-v930-Volume-XXII.pdf}}
\expandafter\def\csname r@nav:chapter-varying-coefficients\endcsname{{54}{3}{}{section.54}{Dirichlet-GRH-v930-Volume-XI.pdf}}
\expandafter\def\csname r@nav:critical-torus-tests\endcsname{{82}{25}{}{section.82}{Dirichlet-GRH-v930-Volume-XVI.pdf}}
\expandafter\def\csname r@nav:first-global-lift\endcsname{{73}{72}{}{section.73}{Dirichlet-GRH-v930-Volume-XIV.pdf}}
\expandafter\def\csname r@nav:first-second-exchange\endcsname{{78}{54}{}{section.78}{Dirichlet-GRH-v930-Volume-XV.pdf}}
\expandafter\def\csname r@nav:frobenius-cups\endcsname{{16}{36}{}{section.16}{Dirichlet-GRH-v930-Volume-III.pdf}}
\expandafter\def\csname r@nav:g-II-0\endcsname{{11}{3}{}{section.11}{Dirichlet-GRH-v930-Volume-II.pdf}}
\expandafter\def\csname r@nav:g-II-1\endcsname{{19}{3}{}{section.19}{Dirichlet-GRH-v930-Volume-IV.pdf}}
\expandafter\def\csname r@nav:g-II-2\endcsname{{20}{22}{}{section.20}{Dirichlet-GRH-v930-Volume-IV.pdf}}
\expandafter\def\csname r@nav:g-II-3\endcsname{{21}{39}{}{section.21}{Dirichlet-GRH-v930-Volume-IV.pdf}}
\expandafter\def\csname r@nav:g-II-4\endcsname{{22}{58}{}{section.22}{Dirichlet-GRH-v930-Volume-IV.pdf}}
\expandafter\def\csname r@nav:g-III-0\endcsname{{24}{3}{}{section.24}{Dirichlet-GRH-v930-Volume-V.pdf}}
\expandafter\def\csname r@nav:g-III-1\endcsname{{25}{24}{}{section.25}{Dirichlet-GRH-v930-Volume-V.pdf}}
\expandafter\def\csname r@nav:g-III-2\endcsname{{26}{46}{}{section.26}{Dirichlet-GRH-v930-Volume-V.pdf}}
\expandafter\def\csname r@nav:g-III-3\endcsname{{27}{74}{}{section.27}{Dirichlet-GRH-v930-Volume-V.pdf}}
\expandafter\def\csname r@nav:g-III-4\endcsname{{28}{84}{}{section.28}{Dirichlet-GRH-v930-Volume-V.pdf}}
\expandafter\def\csname r@nav:g-IV-0\endcsname{{30}{3}{}{section.30}{Dirichlet-GRH-v930-Volume-VI.pdf}}
\expandafter\def\csname r@nav:g-IV-1\endcsname{{31}{12}{}{section.31}{Dirichlet-GRH-v930-Volume-VI.pdf}}
\expandafter\def\csname r@nav:g-IV-2\endcsname{{32}{22}{}{section.32}{Dirichlet-GRH-v930-Volume-VI.pdf}}
\expandafter\def\csname r@nav:g-IV-3\endcsname{{33}{48}{}{section.33}{Dirichlet-GRH-v930-Volume-VI.pdf}}
\expandafter\def\csname r@nav:g-IV-4\endcsname{{34}{64}{}{section.34}{Dirichlet-GRH-v930-Volume-VI.pdf}}
\expandafter\def\csname r@nav:g-IV-5\endcsname{{36}{2}{}{section.36}{Dirichlet-GRH-v930-Volume-VII.pdf}}
\expandafter\def\csname r@nav:g-IX-0\endcsname{{98}{3}{}{section.98}{Dirichlet-GRH-v930-Volume-XIX.pdf}}
\expandafter\def\csname r@nav:g-IX-1\endcsname{{99}{39}{}{section.99}{Dirichlet-GRH-v930-Volume-XIX.pdf}}
\expandafter\def\csname r@nav:g-IX-2\endcsname{{100}{63}{}{section.100}{Dirichlet-GRH-v930-Volume-XIX.pdf}}
\expandafter\def\csname r@nav:g-IX-3\endcsname{{101}{80}{}{section.101}{Dirichlet-GRH-v930-Volume-XIX.pdf}}
\expandafter\def\csname r@nav:g-IX-4\endcsname{{102}{87}{}{section.102}{Dirichlet-GRH-v930-Volume-XIX.pdf}}
\expandafter\def\csname r@nav:g-V-0\endcsname{{64}{3}{}{section.64}{Dirichlet-GRH-v930-Volume-XIII.pdf}}
\expandafter\def\csname r@nav:g-V-1\endcsname{{65}{21}{}{section.65}{Dirichlet-GRH-v930-Volume-XIII.pdf}}
\expandafter\def\csname r@nav:g-V-2\endcsname{{66}{46}{}{section.66}{Dirichlet-GRH-v930-Volume-XIII.pdf}}
\expandafter\def\csname r@nav:g-V-3\endcsname{{67}{75}{}{section.67}{Dirichlet-GRH-v930-Volume-XIII.pdf}}
\expandafter\def\csname r@nav:g-VI-0\endcsname{{70}{3}{}{section.70}{Dirichlet-GRH-v930-Volume-XIV.pdf}}
\expandafter\def\csname r@nav:g-VI-1\endcsname{{71}{13}{}{section.71}{Dirichlet-GRH-v930-Volume-XIV.pdf}}
\expandafter\def\csname r@nav:g-VI-2\endcsname{{72}{44}{}{section.72}{Dirichlet-GRH-v930-Volume-XIV.pdf}}
\expandafter\def\csname r@nav:g-VI-3\endcsname{{74}{100}{}{section.74}{Dirichlet-GRH-v930-Volume-XIV.pdf}}
\expandafter\def\csname r@nav:g-VI-4\endcsname{{75}{112}{}{section.75}{Dirichlet-GRH-v930-Volume-XIV.pdf}}
\expandafter\def\csname r@nav:g-VII-0\endcsname{{86}{3}{}{section.86}{Dirichlet-GRH-v930-Volume-XVII.pdf}}
\expandafter\def\csname r@nav:g-VII-1\endcsname{{87}{32}{}{section.87}{Dirichlet-GRH-v930-Volume-XVII.pdf}}
\expandafter\def\csname r@nav:g-VII-2\endcsname{{88}{52}{}{section.88}{Dirichlet-GRH-v930-Volume-XVII.pdf}}
\expandafter\def\csname r@nav:g-VII-3\endcsname{{89}{73}{}{section.89}{Dirichlet-GRH-v930-Volume-XVII.pdf}}
\expandafter\def\csname r@nav:g-VII-4\endcsname{{90}{83}{}{section.90}{Dirichlet-GRH-v930-Volume-XVII.pdf}}
\expandafter\def\csname r@nav:g-VIII-0\endcsname{{92}{3}{}{section.92}{Dirichlet-GRH-v930-Volume-XVIII.pdf}}
\expandafter\def\csname r@nav:g-VIII-1\endcsname{{93}{21}{}{section.93}{Dirichlet-GRH-v930-Volume-XVIII.pdf}}
\expandafter\def\csname r@nav:g-VIII-2\endcsname{{94}{36}{}{section.94}{Dirichlet-GRH-v930-Volume-XVIII.pdf}}
\expandafter\def\csname r@nav:g-VIII-3\endcsname{{95}{65}{}{section.95}{Dirichlet-GRH-v930-Volume-XVIII.pdf}}
\expandafter\def\csname r@nav:g-VIII-4\endcsname{{83}{52}{}{section.83}{Dirichlet-GRH-v930-Volume-XVI.pdf}}
\expandafter\def\csname r@nav:g-X-0\endcsname{{104}{3}{}{section.104}{Dirichlet-GRH-v930-Volume-XX.pdf}}
\expandafter\def\csname r@nav:g-X-1\endcsname{{105}{32}{}{section.105}{Dirichlet-GRH-v930-Volume-XX.pdf}}
\expandafter\def\csname r@nav:g-X-2\endcsname{{106}{52}{}{section.106}{Dirichlet-GRH-v930-Volume-XX.pdf}}
\expandafter\def\csname r@nav:g-X-3\endcsname{{107}{73}{}{section.107}{Dirichlet-GRH-v930-Volume-XX.pdf}}
\expandafter\def\csname r@nav:g-XI-0\endcsname{{109}{3}{}{section.109}{Dirichlet-GRH-v930-Volume-XXI.pdf}}
\expandafter\def\csname r@nav:g-XI-1\endcsname{{110}{37}{}{section.110}{Dirichlet-GRH-v930-Volume-XXI.pdf}}
\expandafter\def\csname r@nav:g-XI-2\endcsname{{111}{57}{}{section.111}{Dirichlet-GRH-v930-Volume-XXI.pdf}}
\expandafter\def\csname r@nav:g-XI-3\endcsname{{112}{74}{}{section.112}{Dirichlet-GRH-v930-Volume-XXI.pdf}}
\expandafter\def\csname r@nav:g-XII-0\endcsname{{121}{29}{}{section.121}{Dirichlet-GRH-v930-Volume-XXIII.pdf}}
\expandafter\def\csname r@nav:g-XII-1\endcsname{{122}{56}{}{section.122}{Dirichlet-GRH-v930-Volume-XXIII.pdf}}
\expandafter\def\csname r@nav:g-XII-2\endcsname{{123}{76}{}{section.123}{Dirichlet-GRH-v930-Volume-XXIII.pdf}}
\expandafter\def\csname r@nav:g-XII-3\endcsname{{125}{3}{}{section.125}{Dirichlet-GRH-v930-Volume-XXIV.pdf}}
\expandafter\def\csname r@nav:hecke-polynomials\endcsname{{45}{36}{}{section.45}{Dirichlet-GRH-v930-Volume-IX.pdf}}
\expandafter\def\csname r@nav:higher-rank-tests\endcsname{{17}{82}{}{section.17}{Dirichlet-GRH-v930-Volume-III.pdf}}
\expandafter\def\csname r@nav:inert-gram-transforms\endcsname{{138}{67}{}{section.138}{Dirichlet-GRH-v930-Volume-XXVI.pdf}}
\expandafter\def\csname r@nav:local-rational-frequency-sum\endcsname{{139}{93}{}{section.139}{Dirichlet-GRH-v930-Volume-XXVI.pdf}}
\expandafter\def\csname r@nav:marked-open-closed\endcsname{{131}{14}{}{section.131}{Dirichlet-GRH-v930-Volume-XXV.pdf}}
\expandafter\def\csname r@nav:marked-orbital-energy\endcsname{{136}{3}{}{section.136}{Dirichlet-GRH-v930-Volume-XXVI.pdf}}
\expandafter\def\csname r@nav:poincare-spanning\endcsname{{133}{64}{}{section.133}{Dirichlet-GRH-v930-Volume-XXV.pdf}}
\expandafter\def\csname r@nav:selmer-critical\endcsname{{13}{60}{}{section.13}{Dirichlet-GRH-v930-Volume-II.pdf}}
\expandafter\def\csname r@nav:selmer-tests\endcsname{{12}{29}{}{section.12}{Dirichlet-GRH-v930-Volume-II.pdf}}
\expandafter\def\csname r@nav:spanning-and-sphere\endcsname{{134}{79}{}{section.134}{Dirichlet-GRH-v930-Volume-XXV.pdf}}
\expandafter\def\csname r@nav:split-gram-transforms\endcsname{{137}{44}{}{section.137}{Dirichlet-GRH-v930-Volume-XXVI.pdf}}
\expandafter\def\csname r@nav:stationary-cutoff\endcsname{{132}{51}{}{section.132}{Dirichlet-GRH-v930-Volume-XXV.pdf}}
\expandafter\def\csname r@nav:terminal-determinants\endcsname{{84}{77}{}{section.84}{Dirichlet-GRH-v930-Volume-XVI.pdf}}
\expandafter\def\csname r@nav:terminal-whittaker\endcsname{{96}{72}{}{section.96}{Dirichlet-GRH-v930-Volume-XVIII.pdf}}
\expandafter\def\csname r@nav:theta-descent\endcsname{{81}{3}{}{section.81}{Dirichlet-GRH-v930-Volume-XVI.pdf}}
\expandafter\def\csname r@nav:theta-gram\endcsname{{142}{37}{}{section.142}{Dirichlet-GRH-v930-Volume-XXVII.pdf}}
\expandafter\def\csname r@nav:third-exchange\endcsname{{79}{74}{}{section.79}{Dirichlet-GRH-v930-Volume-XV.pdf}}
\expandafter\def\csname r@nav:torus-constituents\endcsname{{50}{3}{}{section.50}{Dirichlet-GRH-v930-Volume-X.pdf}}
\expandafter\def\csname r@nav:torus-modules\endcsname{{52}{59}{}{section.52}{Dirichlet-GRH-v930-Volume-X.pdf}}
\expandafter\def\csname r@nav:torus-primes\endcsname{{51}{39}{}{section.51}{Dirichlet-GRH-v930-Volume-X.pdf}}
\expandafter\def\csname r@nav:trace-extensions\endcsname{{15}{3}{}{section.15}{Dirichlet-GRH-v930-Volume-III.pdf}}
\expandafter\def\csname r@nav:two-mark-flags\endcsname{{130}{3}{}{section.130}{Dirichlet-GRH-v930-Volume-XXV.pdf}}
\expandafter\def\csname r@nav:volume-II\endcsname{{11}{3}{}{section.11}{Dirichlet-GRH-v930-Volume-II.pdf}}
\expandafter\def\csname r@nav:volume-III\endcsname{{15}{3}{}{section.15}{Dirichlet-GRH-v930-Volume-III.pdf}}
\expandafter\def\csname r@nav:volume-IV\endcsname{{19}{3}{}{section.19}{Dirichlet-GRH-v930-Volume-IV.pdf}}
\expandafter\def\csname r@nav:volume-IX\endcsname{{44}{3}{}{section.44}{Dirichlet-GRH-v930-Volume-IX.pdf}}
\expandafter\def\csname r@nav:volume-V\endcsname{{24}{3}{}{section.24}{Dirichlet-GRH-v930-Volume-V.pdf}}
\expandafter\def\csname r@nav:volume-VI\endcsname{{30}{3}{}{section.30}{Dirichlet-GRH-v930-Volume-VI.pdf}}
\expandafter\def\csname r@nav:volume-VII\endcsname{{36}{2}{}{section.36}{Dirichlet-GRH-v930-Volume-VII.pdf}}
\expandafter\def\csname r@nav:volume-VIII\endcsname{{40}{2}{}{section.40}{Dirichlet-GRH-v930-Volume-VIII.pdf}}
\expandafter\def\csname r@nav:volume-X\endcsname{{50}{3}{}{section.50}{Dirichlet-GRH-v930-Volume-X.pdf}}
\expandafter\def\csname r@nav:volume-XI\endcsname{{54}{3}{}{section.54}{Dirichlet-GRH-v930-Volume-XI.pdf}}
\expandafter\def\csname r@nav:volume-XII\endcsname{{59}{3}{}{section.59}{Dirichlet-GRH-v930-Volume-XII.pdf}}
\expandafter\def\csname r@nav:volume-XIII\endcsname{{64}{3}{}{section.64}{Dirichlet-GRH-v930-Volume-XIII.pdf}}
\expandafter\def\csname r@nav:volume-XIV\endcsname{{70}{3}{}{section.70}{Dirichlet-GRH-v930-Volume-XIV.pdf}}
\expandafter\def\csname r@nav:volume-XIX\endcsname{{98}{3}{}{section.98}{Dirichlet-GRH-v930-Volume-XIX.pdf}}
\expandafter\def\csname r@nav:volume-XV\endcsname{{77}{3}{}{section.77}{Dirichlet-GRH-v930-Volume-XV.pdf}}
\expandafter\def\csname r@nav:volume-XVI\endcsname{{81}{3}{}{section.81}{Dirichlet-GRH-v930-Volume-XVI.pdf}}
\expandafter\def\csname r@nav:volume-XVII\endcsname{{86}{3}{}{section.86}{Dirichlet-GRH-v930-Volume-XVII.pdf}}
\expandafter\def\csname r@nav:volume-XVIII\endcsname{{92}{3}{}{section.92}{Dirichlet-GRH-v930-Volume-XVIII.pdf}}
\expandafter\def\csname r@nav:volume-XX\endcsname{{104}{3}{}{section.104}{Dirichlet-GRH-v930-Volume-XX.pdf}}
\expandafter\def\csname r@nav:volume-XXI\endcsname{{109}{3}{}{section.109}{Dirichlet-GRH-v930-Volume-XXI.pdf}}
\expandafter\def\csname r@nav:volume-XXII\endcsname{{115}{3}{}{section.115}{Dirichlet-GRH-v930-Volume-XXII.pdf}}
\expandafter\def\csname r@nav:volume-XXIII\endcsname{{119}{3}{}{section.119}{Dirichlet-GRH-v930-Volume-XXIII.pdf}}
\expandafter\def\csname r@nav:volume-XXIV\endcsname{{125}{3}{}{section.125}{Dirichlet-GRH-v930-Volume-XXIV.pdf}}
\expandafter\def\csname r@nav:volume-XXV\endcsname{{130}{3}{}{section.130}{Dirichlet-GRH-v930-Volume-XXV.pdf}}
\expandafter\def\csname r@nav:volume-XXVI\endcsname{{136}{3}{}{section.136}{Dirichlet-GRH-v930-Volume-XXVI.pdf}}
\expandafter\def\csname r@nav:volume-XXVII\endcsname{{141}{3}{}{section.141}{Dirichlet-GRH-v930-Volume-XXVII.pdf}}
\expandafter\def\csname r@prop:allresidue706\endcsname{{133.19}{77}{}{proposition.133.19}{Dirichlet-GRH-v930-Volume-XXV.pdf}}
\expandafter\def\csname r@prop:armshear887\endcsname{{136.34}{36}{}{proposition.136.34}{Dirichlet-GRH-v930-Volume-XXVI.pdf}}
\expandafter\def\csname r@prop:auxvaluation805\endcsname{{73.6}{78}{}{proposition.73.6}{Dirichlet-GRH-v930-Volume-XIV.pdf}}
\expandafter\def\csname r@prop:blockrows719\endcsname{{45.7}{44}{}{proposition.45.7}{Dirichlet-GRH-v930-Volume-IX.pdf}}
\expandafter\def\csname r@prop:centeredflag341\endcsname{{99.18}{58}{}{proposition.99.18}{Dirichlet-GRH-v930-Volume-XIX.pdf}}
\expandafter\def\csname r@prop:coherentcell391\endcsname{{128.5}{59}{}{proposition.128.5}{Dirichlet-GRH-v930-Volume-XXIV.pdf}}
\expandafter\def\csname r@prop:column110\endcsname{{86.29}{23}{}{proposition.86.29}{Dirichlet-GRH-v930-Volume-XVII.pdf}}
\expandafter\def\csname r@prop:commonaux573\endcsname{{60.7}{28}{}{proposition.60.7}{Dirichlet-GRH-v930-Volume-XII.pdf}}
\expandafter\def\csname r@prop:congruencewave893\endcsname{{137.8}{51}{}{proposition.137.8}{Dirichlet-GRH-v930-Volume-XXVI.pdf}}
\expandafter\def\csname r@prop:descentlayer363\endcsname{{36.28}{32}{}{proposition.36.28}{Dirichlet-GRH-v930-Volume-VII.pdf}}
\expandafter\def\csname r@prop:diag32\endcsname{{121.19}{50}{}{proposition.121.19}{Dirichlet-GRH-v930-Volume-XXIII.pdf}}
\expandafter\def\csname r@prop:dominantexchange344\endcsname{{86.41}{32}{}{proposition.86.41}{Dirichlet-GRH-v930-Volume-XVII.pdf}}
\expandafter\def\csname r@prop:dyadicodd675\endcsname{{11.25}{23}{}{proposition.11.25}{Dirichlet-GRH-v930-Volume-II.pdf}}
\expandafter\def\csname r@prop:effect146\endcsname{{25.16}{37}{}{proposition.25.16}{Dirichlet-GRH-v930-Volume-V.pdf}}
\expandafter\def\csname r@prop:evenplanedeterminants433\endcsname{{46.14}{71}{}{proposition.46.14}{Dirichlet-GRH-v930-Volume-IX.pdf}}
\expandafter\def\csname r@prop:evenreciprocity712\endcsname{{12.3}{31}{}{proposition.12.3}{Dirichlet-GRH-v930-Volume-II.pdf}}
\expandafter\def\csname r@prop:example30\endcsname{{24.15}{14}{}{proposition.24.15}{Dirichlet-GRH-v930-Volume-V.pdf}}
\expandafter\def\csname r@prop:exteriortransport620\endcsname{{117.36}{74}{}{proposition.117.36}{Dirichlet-GRH-v930-Volume-XXII.pdf}}
\expandafter\def\csname r@prop:finite15\endcsname{{11.4}{6}{}{proposition.11.4}{Dirichlet-GRH-v930-Volume-II.pdf}}
\expandafter\def\csname r@prop:finiteordinary558\endcsname{{48.21}{122}{}{proposition.48.21}{Dirichlet-GRH-v930-Volume-IX.pdf}}
\expandafter\def\csname r@prop:flagenergy923\endcsname{{139.26}{123}{}{proposition.139.26}{Dirichlet-GRH-v930-Volume-XXVI.pdf}}
\expandafter\def\csname r@prop:flagsection922\endcsname{{139.24}{120}{}{proposition.139.24}{Dirichlet-GRH-v930-Volume-XXVI.pdf}}
\expandafter\def\csname r@prop:flagtransfer923\endcsname{{139.27}{124}{}{proposition.139.27}{Dirichlet-GRH-v930-Volume-XXVI.pdf}}
\expandafter\def\csname r@prop:four\endcsname{{70.10}{11}{}{proposition.70.10}{Dirichlet-GRH-v930-Volume-XIV.pdf}}
\expandafter\def\csname r@prop:framegeneration446\endcsname{{141.2}{4}{}{proposition.141.2}{Dirichlet-GRH-v930-Volume-XXVII.pdf}}
\expandafter\def\csname r@prop:freeflags380\endcsname{{127.8}{33}{}{proposition.127.8}{Dirichlet-GRH-v930-Volume-XXIV.pdf}}
\expandafter\def\csname r@prop:fulltail621\endcsname{{117.39}{77}{}{proposition.117.39}{Dirichlet-GRH-v930-Volume-XXII.pdf}}
\expandafter\def\csname r@prop:gooddata563\endcsname{{59.8}{9}{}{proposition.59.8}{Dirichlet-GRH-v930-Volume-XII.pdf}}
\expandafter\def\csname r@prop:goodmass868\endcsname{{134.1}{79}{}{proposition.134.1}{Dirichlet-GRH-v930-Volume-XXV.pdf}}
\expandafter\def\csname r@prop:horizontal596\endcsname{{62.6}{72}{}{proposition.62.6}{Dirichlet-GRH-v930-Volume-XII.pdf}}
\expandafter\def\csname r@prop:incidencenorm391\endcsname{{128.4}{59}{}{proposition.128.4}{Dirichlet-GRH-v930-Volume-XXIV.pdf}}
\expandafter\def\csname r@prop:inertloss907\endcsname{{138.22}{88}{}{proposition.138.22}{Dirichlet-GRH-v930-Volume-XXVI.pdf}}
\expandafter\def\csname r@prop:inertoperator900\endcsname{{138.3}{68}{}{proposition.138.3}{Dirichlet-GRH-v930-Volume-XXVI.pdf}}
\expandafter\def\csname r@prop:initialrank392\endcsname{{128.9}{63}{}{proposition.128.9}{Dirichlet-GRH-v930-Volume-XXIV.pdf}}
\expandafter\def\csname r@prop:isotropickernel530\endcsname{{96.17}{84}{}{proposition.96.17}{Dirichlet-GRH-v930-Volume-XVIII.pdf}}
\expandafter\def\csname r@prop:lattices277\endcsname{{90.5}{87}{}{proposition.90.5}{Dirichlet-GRH-v930-Volume-XVII.pdf}}
\expandafter\def\csname r@prop:level241\endcsname{{36.1}{2}{}{proposition.36.1}{Dirichlet-GRH-v930-Volume-VII.pdf}}
\expandafter\def\csname r@prop:level25\endcsname{{121.12}{43}{}{proposition.121.12}{Dirichlet-GRH-v930-Volume-XXIII.pdf}}
\expandafter\def\csname r@prop:levelmass678\endcsname{{130.1}{3}{}{proposition.130.1}{Dirichlet-GRH-v930-Volume-XXV.pdf}}
\expandafter\def\csname r@prop:lineimage522\endcsname{{146.20}{107}{}{proposition.146.20}{Dirichlet-GRH-v930-Volume-XXVII.pdf}}
\expandafter\def\csname r@prop:lowerdimension421\endcsname{{44.27}{24}{}{proposition.44.27}{Dirichlet-GRH-v930-Volume-IX.pdf}}
\expandafter\def\csname r@prop:mainselmer724\endcsname{{12.17}{43}{}{proposition.12.17}{Dirichlet-GRH-v930-Volume-II.pdf}}
\expandafter\def\csname r@prop:markedfourier820\endcsname{{77.30}{33}{}{proposition.77.30}{Dirichlet-GRH-v930-Volume-XV.pdf}}
\expandafter\def\csname r@prop:markedgram625\endcsname{{96.36}{99}{}{proposition.96.36}{Dirichlet-GRH-v930-Volume-XVIII.pdf}}
\expandafter\def\csname r@prop:markedpencil521\endcsname{{146.17}{105}{}{proposition.146.17}{Dirichlet-GRH-v930-Volume-XXVII.pdf}}
\expandafter\def\csname r@prop:markedsampling387\endcsname{{127.20}{48}{}{proposition.127.20}{Dirichlet-GRH-v930-Volume-XXIV.pdf}}
\expandafter\def\csname r@prop:mixedobstruction569\endcsname{{59.22}{21}{}{proposition.59.22}{Dirichlet-GRH-v930-Volume-XII.pdf}}
\expandafter\def\csname r@prop:normalcells388\endcsname{{127.21}{49}{}{proposition.127.21}{Dirichlet-GRH-v930-Volume-XXIV.pdf}}
\expandafter\def\csname r@prop:normalframe349\endcsname{{104.34}{29}{}{proposition.104.34}{Dirichlet-GRH-v930-Volume-XX.pdf}}
\expandafter\def\csname r@prop:odddescent376\endcsname{{38.14}{79}{}{proposition.38.14}{Dirichlet-GRH-v930-Volume-VII.pdf}}
\expandafter\def\csname r@prop:odddivisor384\endcsname{{40.23}{21}{}{proposition.40.23}{Dirichlet-GRH-v930-Volume-VIII.pdf}}
\expandafter\def\csname r@prop:onecompanion560\endcsname{{59.1}{3}{}{proposition.59.1}{Dirichlet-GRH-v930-Volume-XII.pdf}}
\expandafter\def\csname r@prop:orbits273\endcsname{{21.11}{48}{}{proposition.21.11}{Dirichlet-GRH-v930-Volume-IV.pdf}}
\expandafter\def\csname r@prop:orbitsum501\endcsname{{144.13}{71}{}{proposition.144.13}{Dirichlet-GRH-v930-Volume-XXVII.pdf}}
\expandafter\def\csname r@prop:orighecke374\endcsname{{38.7}{74}{}{proposition.38.7}{Dirichlet-GRH-v930-Volume-VII.pdf}}
\expandafter\def\csname r@prop:origideal373\endcsname{{38.2}{71}{}{proposition.38.2}{Dirichlet-GRH-v930-Volume-VII.pdf}}
\expandafter\def\csname r@prop:others77\endcsname{{66.27}{66}{}{proposition.66.27}{Dirichlet-GRH-v930-Volume-XIII.pdf}}
\expandafter\def\csname r@prop:partnerweights531\endcsname{{96.20}{87}{}{proposition.96.20}{Dirichlet-GRH-v930-Volume-XVIII.pdf}}
\expandafter\def\csname r@prop:pellunit506\endcsname{{57.9}{73}{}{proposition.57.9}{Dirichlet-GRH-v930-Volume-XI.pdf}}
\expandafter\def\csname r@prop:pivotflags619\endcsname{{117.33}{71}{}{proposition.117.33}{Dirichlet-GRH-v930-Volume-XXII.pdf}}
\expandafter\def\csname r@prop:positiondrop913\endcsname{{139.9}{100}{}{proposition.139.9}{Dirichlet-GRH-v930-Volume-XXVI.pdf}}
\expandafter\def\csname r@prop:primerange670\endcsname{{141.32}{30}{}{proposition.141.32}{Dirichlet-GRH-v930-Volume-XXVII.pdf}}
\expandafter\def\csname r@prop:quadratic345\endcsname{{36.3}{3}{}{proposition.36.3}{Dirichlet-GRH-v930-Volume-VII.pdf}}
\expandafter\def\csname r@prop:quadraticorders608\endcsname{{22.39}{93}{}{proposition.22.39}{Dirichlet-GRH-v930-Volume-IV.pdf}}
\expandafter\def\csname r@prop:ramifiedorders782\endcsname{{16.44}{80}{}{proposition.16.44}{Dirichlet-GRH-v930-Volume-III.pdf}}
\expandafter\def\csname r@prop:rankonelift599\endcsname{{62.17}{84}{}{proposition.62.17}{Dirichlet-GRH-v930-Volume-XII.pdf}}
\expandafter\def\csname r@prop:ranks47\endcsname{{92.1}{3}{}{proposition.92.1}{Dirichlet-GRH-v930-Volume-XVIII.pdf}}
\expandafter\def\csname r@prop:realizedhyperplanes537\endcsname{{113.17}{97}{}{proposition.113.17}{Dirichlet-GRH-v930-Volume-XXI.pdf}}
\expandafter\def\csname r@prop:realpivotdensity916\endcsname{{139.14}{106}{}{proposition.139.14}{Dirichlet-GRH-v930-Volume-XXVI.pdf}}
\expandafter\def\csname r@prop:repeatedtransport534\endcsname{{113.5}{87}{}{proposition.113.5}{Dirichlet-GRH-v930-Volume-XXI.pdf}}
\expandafter\def\csname r@prop:residue269\endcsname{{21.10}{48}{}{proposition.21.10}{Dirichlet-GRH-v930-Volume-IV.pdf}}
\expandafter\def\csname r@prop:retainedtest535\endcsname{{113.9}{90}{}{proposition.113.9}{Dirichlet-GRH-v930-Volume-XXI.pdf}}
\expandafter\def\csname r@prop:righttransfer721\endcsname{{45.15}{52}{}{proposition.45.15}{Dirichlet-GRH-v930-Volume-IX.pdf}}
\expandafter\def\csname r@prop:root31\endcsname{{30.5}{6}{}{proposition.30.5}{Dirichlet-GRH-v930-Volume-VI.pdf}}
\expandafter\def\csname r@prop:rotationideal611\endcsname{{117.23}{60}{}{proposition.117.23}{Dirichlet-GRH-v930-Volume-XXII.pdf}}
\expandafter\def\csname r@prop:rowrealgram809\endcsname{{73.18}{90}{}{proposition.73.18}{Dirichlet-GRH-v930-Volume-XIV.pdf}}
\expandafter\def\csname r@prop:seedadjoints580\endcsname{{60.25}{45}{}{proposition.60.25}{Dirichlet-GRH-v930-Volume-XII.pdf}}
\expandafter\def\csname r@prop:separatedalgebra573\endcsname{{60.9}{30}{}{proposition.60.9}{Dirichlet-GRH-v930-Volume-XII.pdf}}
\expandafter\def\csname r@prop:sharpminors727\endcsname{{12.22}{47}{}{proposition.12.22}{Dirichlet-GRH-v930-Volume-II.pdf}}
\expandafter\def\csname r@prop:shortpaths352\endcsname{{119.12}{12}{}{proposition.119.12}{Dirichlet-GRH-v930-Volume-XXIII.pdf}}
\expandafter\def\csname r@prop:simultaneousorders779\endcsname{{16.39}{73}{}{proposition.16.39}{Dirichlet-GRH-v930-Volume-III.pdf}}
\expandafter\def\csname r@prop:smallcross590\endcsname{{116.9}{32}{}{proposition.116.9}{Dirichlet-GRH-v930-Volume-XXII.pdf}}
\expandafter\def\csname r@prop:sourcetwist803\endcsname{{52.18}{77}{}{proposition.52.18}{Dirichlet-GRH-v930-Volume-X.pdf}}
\expandafter\def\csname r@prop:splitexception598\endcsname{{62.11}{79}{}{proposition.62.11}{Dirichlet-GRH-v930-Volume-XII.pdf}}
\expandafter\def\csname r@prop:strictdescent404\endcsname{{41.18}{46}{}{proposition.41.18}{Dirichlet-GRH-v930-Volume-VIII.pdf}}
\expandafter\def\csname r@prop:stweight616\endcsname{{32.22}{46}{}{proposition.32.22}{Dirichlet-GRH-v930-Volume-VI.pdf}}
\expandafter\def\csname r@prop:tamecurve365\endcsname{{37.9}{42}{}{proposition.37.9}{Dirichlet-GRH-v930-Volume-VII.pdf}}
\expandafter\def\csname r@prop:tamerigid594\endcsname{{62.1}{64}{}{proposition.62.1}{Dirichlet-GRH-v930-Volume-XII.pdf}}
\expandafter\def\csname r@prop:toproots716\endcsname{{45.2}{37}{}{proposition.45.2}{Dirichlet-GRH-v930-Volume-IX.pdf}}
\expandafter\def\csname r@prop:torsionterm377\endcsname{{40.3}{3}{}{proposition.40.3}{Dirichlet-GRH-v930-Volume-VIII.pdf}}
\expandafter\def\csname r@prop:torusbudget784\endcsname{{52.1}{59}{}{proposition.52.1}{Dirichlet-GRH-v930-Volume-X.pdf}}
\expandafter\def\csname r@prop:tracecrystal375\endcsname{{38.10}{76}{}{proposition.38.10}{Dirichlet-GRH-v930-Volume-VII.pdf}}
\expandafter\def\csname r@prop:twistedfiber919\endcsname{{139.20}{113}{}{proposition.139.20}{Dirichlet-GRH-v930-Volume-XXVI.pdf}}
\expandafter\def\csname r@prop:upperorders382\endcsname{{40.17}{14}{}{proposition.40.17}{Dirichlet-GRH-v930-Volume-VIII.pdf}}
\expandafter\def\csname r@prop:weightchange579\endcsname{{60.22}{41}{}{proposition.60.22}{Dirichlet-GRH-v930-Volume-XII.pdf}}
\expandafter\def\csname r@sec:signedsignal450\endcsname{{141.4}{12}{}{subsection.141.4}{Dirichlet-GRH-v930-Volume-XXVII.pdf}}
\expandafter\def\csname r@thm:abelianperiod744\endcsname{{50.13}{17}{}{proposition.50.13}{Dirichlet-GRH-v930-Volume-X.pdf}}
\expandafter\def\csname r@thm:abslow349\endcsname{{104.32}{28}{}{proposition.104.32}{Dirichlet-GRH-v930-Volume-XX.pdf}}
\expandafter\def\csname r@thm:actualminors804\endcsname{{52.24}{82}{}{proposition.52.24}{Dirichlet-GRH-v930-Volume-X.pdf}}
\expandafter\def\csname r@thm:actualseedcongruences436\endcsname{{55.9}{25}{}{proposition.55.9}{Dirichlet-GRH-v930-Volume-XI.pdf}}
\expandafter\def\csname r@thm:adjacentchart421\endcsname{{44.29}{26}{}{proposition.44.29}{Dirichlet-GRH-v930-Volume-IX.pdf}}
\expandafter\def\csname r@thm:adjacentweights427\endcsname{{54.13}{13}{}{proposition.54.13}{Dirichlet-GRH-v930-Volume-XI.pdf}}
\expandafter\def\csname r@thm:adjcomponent372\endcsname{{37.36}{69}{}{proposition.37.36}{Dirichlet-GRH-v930-Volume-VII.pdf}}
\expandafter\def\csname r@thm:adjoint240\endcsname{{34.16}{78}{}{proposition.34.16}{Dirichlet-GRH-v930-Volume-VI.pdf}}
\expandafter\def\csname r@thm:adjointseed361\endcsname{{36.17}{17}{}{proposition.36.17}{Dirichlet-GRH-v930-Volume-VII.pdf}}
\expandafter\def\csname r@thm:adjtorsion381\endcsname{{40.16}{13}{}{proposition.40.16}{Dirichlet-GRH-v930-Volume-VIII.pdf}}
\expandafter\def\csname r@thm:allcarriers830\endcsname{{78.10}{63}{}{proposition.78.10}{Dirichlet-GRH-v930-Volume-XV.pdf}}
\expandafter\def\csname r@thm:allclasses350\endcsname{{119.8}{8}{}{proposition.119.8}{Dirichlet-GRH-v930-Volume-XXIII.pdf}}
\expandafter\def\csname r@thm:allcollisiontraces445\endcsname{{55.32}{47}{}{proposition.55.32}{Dirichlet-GRH-v930-Volume-XI.pdf}}
\expandafter\def\csname r@thm:allfixermass888\endcsname{{136.36}{39}{}{proposition.136.36}{Dirichlet-GRH-v930-Volume-XXVI.pdf}}
\expandafter\def\csname r@thm:allglobalroots547\endcsname{{48.6}{109}{}{proposition.48.6}{Dirichlet-GRH-v930-Volume-IX.pdf}}
\expandafter\def\csname r@thm:allhalfweights644\endcsname{{57.51}{106}{}{proposition.57.51}{Dirichlet-GRH-v930-Volume-XI.pdf}}
\expandafter\def\csname r@thm:allhorizontal401\endcsname{{41.8}{36}{}{proposition.41.8}{Dirichlet-GRH-v930-Volume-VIII.pdf}}
\expandafter\def\csname r@thm:allinduction586\endcsname{{61.12}{57}{}{proposition.61.12}{Dirichlet-GRH-v930-Volume-XII.pdf}}
\expandafter\def\csname r@thm:allkmass897\endcsname{{137.17}{60}{}{proposition.137.17}{Dirichlet-GRH-v930-Volume-XXVI.pdf}}
\expandafter\def\csname r@thm:alllabels375\endcsname{{38.11}{77}{}{proposition.38.11}{Dirichlet-GRH-v930-Volume-VII.pdf}}
\expandafter\def\csname r@thm:alllevels405\endcsname{{41.23}{51}{}{proposition.41.23}{Dirichlet-GRH-v930-Volume-VIII.pdf}}
\expandafter\def\csname r@thm:allmismatch925\endcsname{{139.29}{127}{}{proposition.139.29}{Dirichlet-GRH-v930-Volume-XXVI.pdf}}
\expandafter\def\csname r@thm:allnorm865\endcsname{{84.10}{85}{}{proposition.84.10}{Dirichlet-GRH-v930-Volume-XVI.pdf}}
\expandafter\def\csname r@thm:allpair689\endcsname{{131.23}{33}{}{proposition.131.23}{Dirichlet-GRH-v930-Volume-XXV.pdf}}
\expandafter\def\csname r@thm:allplacesminor477\endcsname{{143.2}{50}{}{proposition.143.2}{Dirichlet-GRH-v930-Volume-XXVII.pdf}}
\expandafter\def\csname r@thm:allpositions883\endcsname{{136.23}{25}{}{proposition.136.23}{Dirichlet-GRH-v930-Volume-XXVI.pdf}}
\expandafter\def\csname r@thm:allprimeslayers543\endcsname{{57.35}{93}{}{proposition.57.35}{Dirichlet-GRH-v930-Volume-XI.pdf}}
\expandafter\def\csname r@thm:allprimetable862\endcsname{{84.2}{77}{}{proposition.84.2}{Dirichlet-GRH-v930-Volume-XVI.pdf}}
\expandafter\def\csname r@thm:allramifiedlifts761\endcsname{{16.11}{44}{}{proposition.16.11}{Dirichlet-GRH-v930-Volume-III.pdf}}
\expandafter\def\csname r@thm:allraypowers726\endcsname{{12.20}{45}{}{proposition.12.20}{Dirichlet-GRH-v930-Volume-II.pdf}}
\expandafter\def\csname r@thm:allrealwaves924\endcsname{{139.28}{125}{}{proposition.139.28}{Dirichlet-GRH-v930-Volume-XXVI.pdf}}
\expandafter\def\csname r@thm:allscalarfit718\endcsname{{45.5}{41}{}{proposition.45.5}{Dirichlet-GRH-v930-Volume-IX.pdf}}
\expandafter\def\csname r@thm:allspherical798\endcsname{{17.28}{113}{}{proposition.17.28}{Dirichlet-GRH-v930-Volume-III.pdf}}
\expandafter\def\csname r@thm:alltails648\endcsname{{57.59}{113}{}{proposition.57.59}{Dirichlet-GRH-v930-Volume-XI.pdf}}
\expandafter\def\csname r@thm:allterminaltransports641\endcsname{{117.80}{114}{}{proposition.117.80}{Dirichlet-GRH-v930-Volume-XXII.pdf}}
\expandafter\def\csname r@thm:alltorus783\endcsname{{51.23}{58}{}{proposition.51.23}{Dirichlet-GRH-v930-Volume-X.pdf}}
\expandafter\def\csname r@thm:allunmarked909\endcsname{{139.1}{93}{}{proposition.139.1}{Dirichlet-GRH-v930-Volume-XXVI.pdf}}
\expandafter\def\csname r@thm:allzeroweights442\endcsname{{55.24}{40}{}{proposition.55.24}{Dirichlet-GRH-v930-Volume-XI.pdf}}
\expandafter\def\csname r@thm:angularerror610\endcsname{{117.21}{58}{}{proposition.117.21}{Dirichlet-GRH-v930-Volume-XXII.pdf}}
\expandafter\def\csname r@thm:angularlocalization696\endcsname{{132.5}{54}{}{proposition.132.5}{Dirichlet-GRH-v930-Volume-XXV.pdf}}
\expandafter\def\csname r@thm:archbessel684\endcsname{{131.8}{21}{}{proposition.131.8}{Dirichlet-GRH-v930-Volume-XXV.pdf}}
\expandafter\def\csname r@thm:arithmeticclass824\endcsname{{77.41}{48}{}{proposition.77.41}{Dirichlet-GRH-v930-Volume-XV.pdf}}
\expandafter\def\csname r@thm:armopen879\endcsname{{136.14}{16}{}{proposition.136.14}{Dirichlet-GRH-v930-Volume-XXVI.pdf}}
\expandafter\def\csname r@thm:artiniansplit490\endcsname{{47.30}{97}{}{proposition.47.30}{Dirichlet-GRH-v930-Volume-IX.pdf}}
\expandafter\def\csname r@thm:augmented350\endcsname{{119.9}{9}{}{proposition.119.9}{Dirichlet-GRH-v930-Volume-XXIII.pdf}}
\expandafter\def\csname r@thm:augterminal501\endcsname{{144.14}{71}{}{proposition.144.14}{Dirichlet-GRH-v930-Volume-XXVII.pdf}}
\expandafter\def\csname r@thm:automult559\endcsname{{48.25}{125}{}{proposition.48.25}{Dirichlet-GRH-v930-Volume-IX.pdf}}
\expandafter\def\csname r@thm:auxdescent361\endcsname{{36.18}{18}{}{proposition.36.18}{Dirichlet-GRH-v930-Volume-VII.pdf}}
\expandafter\def\csname r@thm:auxdet805\endcsname{{73.5}{78}{}{proposition.73.5}{Dirichlet-GRH-v930-Volume-XIV.pdf}}
\expandafter\def\csname r@thm:auxtriangle713\endcsname{{12.6}{33}{}{proposition.12.6}{Dirichlet-GRH-v930-Volume-II.pdf}}
\expandafter\def\csname r@thm:balancedcount728\endcsname{{12.28}{51}{}{proposition.12.28}{Dirichlet-GRH-v930-Volume-II.pdf}}
\expandafter\def\csname r@thm:bessel29\endcsname{{121.14}{46}{}{proposition.121.14}{Dirichlet-GRH-v930-Volume-XXIII.pdf}}
\expandafter\def\csname r@thm:besselfourier679\endcsname{{130.3}{6}{}{proposition.130.3}{Dirichlet-GRH-v930-Volume-XXV.pdf}}
\expandafter\def\csname r@thm:besselmatrix701\endcsname{{133.2}{65}{}{proposition.133.2}{Dirichlet-GRH-v930-Volume-XXV.pdf}}
\expandafter\def\csname r@thm:besselperiod707\endcsname{{142.2}{37}{}{proposition.142.2}{Dirichlet-GRH-v930-Volume-XXVII.pdf}}
\expandafter\def\csname r@thm:blockwitness575\endcsname{{60.14}{33}{}{proposition.60.14}{Dirichlet-GRH-v930-Volume-XII.pdf}}
\expandafter\def\csname r@thm:bothmarks873\endcsname{{134.11}{96}{}{proposition.134.11}{Dirichlet-GRH-v930-Volume-XXV.pdf}}
\expandafter\def\csname r@thm:bothordinary553\endcsname{{115.6}{7}{}{proposition.115.6}{Dirichlet-GRH-v930-Volume-XXII.pdf}}
\expandafter\def\csname r@thm:bothranks534\endcsname{{113.3}{86}{}{proposition.113.3}{Dirichlet-GRH-v930-Volume-XXI.pdf}}
\expandafter\def\csname r@thm:carrier813\endcsname{{77.6}{7}{}{proposition.77.6}{Dirichlet-GRH-v930-Volume-XV.pdf}}
\expandafter\def\csname r@thm:carrierline831\endcsname{{78.12}{67}{}{proposition.78.12}{Dirichlet-GRH-v930-Volume-XV.pdf}}
\expandafter\def\csname r@thm:carriermultiplicity825\endcsname{{77.44}{50}{}{proposition.77.44}{Dirichlet-GRH-v930-Volume-XV.pdf}}
\expandafter\def\csname r@thm:cells103\endcsname{{20.2}{23}{}{proposition.20.2}{Dirichlet-GRH-v930-Volume-IV.pdf}}
\expandafter\def\csname r@thm:center\endcsname{{70.7}{9}{}{proposition.70.7}{Dirichlet-GRH-v930-Volume-XIV.pdf}}
\expandafter\def\csname r@thm:centralfour636\endcsname{{117.66}{102}{}{proposition.117.66}{Dirichlet-GRH-v930-Volume-XXII.pdf}}
\expandafter\def\csname r@thm:centraltraces635\endcsname{{117.62}{99}{}{proposition.117.62}{Dirichlet-GRH-v930-Volume-XXII.pdf}}
\expandafter\def\csname r@thm:chaintransfer417\endcsname{{44.12}{12}{}{proposition.44.12}{Dirichlet-GRH-v930-Volume-IX.pdf}}
\expandafter\def\csname r@thm:char\endcsname{{24.1}{3}{}{proposition.24.1}{Dirichlet-GRH-v930-Volume-V.pdf}}
\expandafter\def\csname r@thm:charrows719\endcsname{{45.8}{44}{}{proposition.45.8}{Dirichlet-GRH-v930-Volume-IX.pdf}}
\expandafter\def\csname r@thm:circindex267\endcsname{{27.1}{75}{}{proposition.27.1}{Dirichlet-GRH-v930-Volume-V.pdf}}
\expandafter\def\csname r@thm:classden312\endcsname{{122.13}{71}{}{proposition.122.13}{Dirichlet-GRH-v930-Volume-XXIII.pdf}}
\expandafter\def\csname r@thm:classfrob503\endcsname{{56.20}{65}{}{proposition.56.20}{Dirichlet-GRH-v930-Volume-XI.pdf}}
\expandafter\def\csname r@thm:classgauss508\endcsname{{144.20}{76}{}{proposition.144.20}{Dirichlet-GRH-v930-Volume-XXVII.pdf}}
\expandafter\def\csname r@thm:classicaldefect360\endcsname{{36.14}{14}{}{proposition.36.14}{Dirichlet-GRH-v930-Volume-VII.pdf}}
\expandafter\def\csname r@thm:classicalfiber596\endcsname{{62.8}{74}{}{proposition.62.8}{Dirichlet-GRH-v930-Volume-XII.pdf}}
\expandafter\def\csname r@thm:classtrace326\endcsname{{122.17}{75}{}{proposition.122.17}{Dirichlet-GRH-v930-Volume-XXIII.pdf}}
\expandafter\def\csname r@thm:cmcores403\endcsname{{41.15}{43}{}{proposition.41.15}{Dirichlet-GRH-v930-Volume-VIII.pdf}}
\expandafter\def\csname r@thm:codimtwo624\endcsname{{96.31}{95}{}{proposition.96.31}{Dirichlet-GRH-v930-Volume-XVIII.pdf}}
\expandafter\def\csname r@thm:collisionbranch464\endcsname{{56.8}{55}{}{proposition.56.8}{Dirichlet-GRH-v930-Volume-XI.pdf}}
\expandafter\def\csname r@thm:commoncompactgram447\endcsname{{141.6}{8}{}{proposition.141.6}{Dirichlet-GRH-v930-Volume-XXVII.pdf}}
\expandafter\def\csname r@thm:commoncompletion889\endcsname{{136.38}{41}{}{proposition.136.38}{Dirichlet-GRH-v930-Volume-XXVI.pdf}}
\expandafter\def\csname r@thm:commoncritical731\endcsname{{13.4}{62}{}{proposition.13.4}{Dirichlet-GRH-v930-Volume-II.pdf}}
\expandafter\def\csname r@thm:commonfamilies631\endcsname{{117.57}{94}{}{proposition.117.57}{Dirichlet-GRH-v930-Volume-XXII.pdf}}
\expandafter\def\csname r@thm:commonframe365\endcsname{{37.8}{41}{}{proposition.37.8}{Dirichlet-GRH-v930-Volume-VII.pdf}}
\expandafter\def\csname r@thm:commonhecke716\endcsname{{45.3}{38}{}{proposition.45.3}{Dirichlet-GRH-v930-Volume-IX.pdf}}
\expandafter\def\csname r@thm:commonlength267\endcsname{{21.9}{46}{}{proposition.21.9}{Dirichlet-GRH-v930-Volume-IV.pdf}}
\expandafter\def\csname r@thm:commonorbits459\endcsname{{56.5}{53}{}{proposition.56.5}{Dirichlet-GRH-v930-Volume-XI.pdf}}
\expandafter\def\csname r@thm:commonphase441\endcsname{{128.32}{82}{}{proposition.128.32}{Dirichlet-GRH-v930-Volume-XXIV.pdf}}
\expandafter\def\csname r@thm:commonprobes594\endcsname{{62.2}{65}{}{proposition.62.2}{Dirichlet-GRH-v930-Volume-XII.pdf}}
\expandafter\def\csname r@thm:commonsource806\endcsname{{73.8}{81}{}{proposition.73.8}{Dirichlet-GRH-v930-Volume-XIV.pdf}}
\expandafter\def\csname r@thm:commonsplit483\endcsname{{68.15}{96}{}{proposition.68.15}{Dirichlet-GRH-v930-Volume-XIII.pdf}}
\expandafter\def\csname r@thm:commonsptrace433\endcsname{{46.16}{72}{}{proposition.46.16}{Dirichlet-GRH-v930-Volume-IX.pdf}}
\expandafter\def\csname r@thm:compactmass630\endcsname{{117.53}{90}{}{proposition.117.53}{Dirichlet-GRH-v930-Volume-XXII.pdf}}
\expandafter\def\csname r@thm:compare30\endcsname{{24.13}{13}{}{proposition.24.13}{Dirichlet-GRH-v930-Volume-V.pdf}}
\expandafter\def\csname r@thm:compgap787\endcsname{{52.12}{69}{}{proposition.52.12}{Dirichlet-GRH-v930-Volume-X.pdf}}
\expandafter\def\csname r@thm:complementgram812\endcsname{{77.3}{5}{}{proposition.77.3}{Dirichlet-GRH-v930-Volume-XV.pdf}}
\expandafter\def\csname r@thm:completehecke418\endcsname{{44.15}{15}{}{proposition.44.15}{Dirichlet-GRH-v930-Volume-IX.pdf}}
\expandafter\def\csname r@thm:completesystem856\endcsname{{82.8}{36}{}{proposition.82.8}{Dirichlet-GRH-v930-Volume-XVI.pdf}}
\expandafter\def\csname r@thm:complexseparation818\endcsname{{77.25}{27}{}{proposition.77.25}{Dirichlet-GRH-v930-Volume-XV.pdf}}
\expandafter\def\csname r@thm:compressionbudget851\endcsname{{81.19}{23}{}{proposition.81.19}{Dirichlet-GRH-v930-Volume-XVI.pdf}}
\expandafter\def\csname r@thm:conditionedpivot917\endcsname{{139.15}{108}{}{proposition.139.15}{Dirichlet-GRH-v930-Volume-XXVI.pdf}}
\expandafter\def\csname r@thm:conductor232\endcsname{{19.15}{12}{}{proposition.19.15}{Dirichlet-GRH-v930-Volume-IV.pdf}}
\expandafter\def\csname r@thm:conemoments664\endcsname{{117.98}{127}{}{proposition.117.98}{Dirichlet-GRH-v930-Volume-XXII.pdf}}
\expandafter\def\csname r@thm:congruenceorder471\endcsname{{47.11}{84}{}{proposition.47.11}{Dirichlet-GRH-v930-Volume-IX.pdf}}
\expandafter\def\csname r@thm:conjugacyfourier705\endcsname{{133.16}{75}{}{proposition.133.16}{Dirichlet-GRH-v930-Volume-XXV.pdf}}
\expandafter\def\csname r@thm:connectedjets516\endcsname{{146.2}{94}{}{proposition.146.2}{Dirichlet-GRH-v930-Volume-XXVII.pdf}}
\expandafter\def\csname r@thm:cornermatrix414\endcsname{{42.24}{83}{}{proposition.42.24}{Dirichlet-GRH-v930-Volume-VIII.pdf}}
\expandafter\def\csname r@thm:cosetcomparison531\endcsname{{96.21}{88}{}{proposition.96.21}{Dirichlet-GRH-v930-Volume-XVIII.pdf}}
\expandafter\def\csname r@thm:count37\endcsname{{64.2}{4}{}{proposition.64.2}{Dirichlet-GRH-v930-Volume-XIII.pdf}}
\expandafter\def\csname r@thm:count73\endcsname{{66.2}{47}{}{proposition.66.2}{Dirichlet-GRH-v930-Volume-XIII.pdf}}
\expandafter\def\csname r@thm:covectorpaths352\endcsname{{73.1}{73}{}{proposition.73.1}{Dirichlet-GRH-v930-Volume-XIV.pdf}}
\expandafter\def\csname r@thm:critical13\endcsname{{11.1}{3}{}{proposition.11.1}{Dirichlet-GRH-v930-Volume-II.pdf}}
\expandafter\def\csname r@thm:criticalmarks339\endcsname{{22.12}{71}{}{proposition.22.12}{Dirichlet-GRH-v930-Volume-IV.pdf}}
\expandafter\def\csname r@thm:crossfiber919\endcsname{{139.19}{112}{}{proposition.139.19}{Dirichlet-GRH-v930-Volume-XXVI.pdf}}
\expandafter\def\csname r@thm:crossflags923\endcsname{{139.25}{122}{}{proposition.139.25}{Dirichlet-GRH-v930-Volume-XXVI.pdf}}
\expandafter\def\csname r@thm:cubedepth584\endcsname{{61.5}{52}{}{proposition.61.5}{Dirichlet-GRH-v930-Volume-XII.pdf}}
\expandafter\def\csname r@thm:cubeweights584\endcsname{{61.4}{50}{}{proposition.61.4}{Dirichlet-GRH-v930-Volume-XII.pdf}}
\expandafter\def\csname r@thm:cubictower662\endcsname{{26.26}{66}{}{proposition.26.26}{Dirichlet-GRH-v930-Volume-V.pdf}}
\expandafter\def\csname r@thm:cuppath757\endcsname{{16.2}{36}{}{proposition.16.2}{Dirichlet-GRH-v930-Volume-III.pdf}}
\expandafter\def\csname r@thm:cuspidalinitial577\endcsname{{60.18}{37}{}{proposition.60.18}{Dirichlet-GRH-v930-Volume-XII.pdf}}
\expandafter\def\csname r@thm:cut79\endcsname{{24.18}{16}{}{proposition.24.18}{Dirichlet-GRH-v930-Volume-V.pdf}}
\expandafter\def\csname r@thm:cycleorientation674\endcsname{{145.14}{91}{}{proposition.145.14}{Dirichlet-GRH-v930-Volume-XXVII.pdf}}
\expandafter\def\csname r@thm:cycleweyl461\endcsname{{68.2}{88}{}{proposition.68.2}{Dirichlet-GRH-v930-Volume-XIII.pdf}}
\expandafter\def\csname r@thm:cyclicdescent661\endcsname{{22.59}{109}{}{proposition.22.59}{Dirichlet-GRH-v930-Volume-IV.pdf}}
\expandafter\def\csname r@thm:cyclicstrict397\endcsname{{40.30}{28}{}{proposition.40.30}{Dirichlet-GRH-v930-Volume-VIII.pdf}}
\expandafter\def\csname r@thm:cyclicunit495\endcsname{{68.42}{115}{}{proposition.68.42}{Dirichlet-GRH-v930-Volume-XIII.pdf}}
\expandafter\def\csname r@thm:decomploss330\endcsname{{22.8}{67}{}{proposition.22.8}{Dirichlet-GRH-v930-Volume-IV.pdf}}
\expandafter\def\csname r@thm:deepintegrallift767\endcsname{{16.25}{56}{}{proposition.16.25}{Dirichlet-GRH-v930-Volume-III.pdf}}
\expandafter\def\csname r@thm:definiteclass822\endcsname{{77.36}{41}{}{proposition.77.36}{Dirichlet-GRH-v930-Volume-XV.pdf}}
\expandafter\def\csname r@thm:definitelevels678\endcsname{{130.2}{4}{}{proposition.130.2}{Dirichlet-GRH-v930-Volume-XXV.pdf}}
\expandafter\def\csname r@thm:definiteseparation530\endcsname{{96.12}{81}{}{proposition.96.12}{Dirichlet-GRH-v930-Volume-XVIII.pdf}}
\expandafter\def\csname r@thm:descent240\endcsname{{21.16}{53}{}{proposition.21.16}{Dirichlet-GRH-v930-Volume-IV.pdf}}
\expandafter\def\csname r@thm:descentcost665\endcsname{{26.30}{70}{}{proposition.26.30}{Dirichlet-GRH-v930-Volume-V.pdf}}
\expandafter\def\csname r@thm:descenttower846\endcsname{{81.8}{10}{}{proposition.81.8}{Dirichlet-GRH-v930-Volume-XVI.pdf}}
\expandafter\def\csname r@thm:detect241\endcsname{{21.17}{54}{}{proposition.21.17}{Dirichlet-GRH-v930-Volume-IV.pdf}}
\expandafter\def\csname r@thm:detquotient810\endcsname{{73.22}{94}{}{proposition.73.22}{Dirichlet-GRH-v930-Volume-XIV.pdf}}
\expandafter\def\csname r@thm:diagonaltail536\endcsname{{113.12}{92}{}{proposition.113.12}{Dirichlet-GRH-v930-Volume-XXI.pdf}}
\expandafter\def\csname r@thm:discriminant681\endcsname{{130.9}{11}{}{proposition.130.9}{Dirichlet-GRH-v930-Volume-XXV.pdf}}
\expandafter\def\csname r@thm:disjointtransports638\endcsname{{117.70}{106}{}{proposition.117.70}{Dirichlet-GRH-v930-Volume-XXII.pdf}}
\expandafter\def\csname r@thm:displacedreal344\endcsname{{107.18}{90}{}{proposition.107.18}{Dirichlet-GRH-v930-Volume-XX.pdf}}
\expandafter\def\csname r@thm:divconstant513\endcsname{{145.12}{89}{}{proposition.145.12}{Dirichlet-GRH-v930-Volume-XXVII.pdf}}
\expandafter\def\csname r@thm:divisionphase493\endcsname{{68.35}{110}{}{proposition.68.35}{Dirichlet-GRH-v930-Volume-XIII.pdf}}
\expandafter\def\csname r@thm:divtail511\endcsname{{145.7}{85}{}{proposition.145.7}{Dirichlet-GRH-v930-Volume-XXVII.pdf}}
\expandafter\def\csname r@thm:doubleflag546\endcsname{{48.5}{108}{}{proposition.48.5}{Dirichlet-GRH-v930-Volume-IX.pdf}}
\expandafter\def\csname r@thm:doubleorbit532\endcsname{{96.26}{91}{}{proposition.96.26}{Dirichlet-GRH-v930-Volume-XVIII.pdf}}
\expandafter\def\csname r@thm:dualcurve598\endcsname{{62.12}{79}{}{proposition.62.12}{Dirichlet-GRH-v930-Volume-XII.pdf}}
\expandafter\def\csname r@thm:dualrelations712\endcsname{{12.2}{30}{}{proposition.12.2}{Dirichlet-GRH-v930-Volume-II.pdf}}
\expandafter\def\csname r@thm:dyadicall525\endcsname{{57.21}{82}{}{proposition.57.21}{Dirichlet-GRH-v930-Volume-XI.pdf}}
\expandafter\def\csname r@thm:effectivestationary695\endcsname{{132.2}{51}{}{proposition.132.2}{Dirichlet-GRH-v930-Volume-XXV.pdf}}
\expandafter\def\csname r@thm:endtransform894\endcsname{{137.11}{53}{}{proposition.137.11}{Dirichlet-GRH-v930-Volume-XXVI.pdf}}
\expandafter\def\csname r@thm:enormousall559\endcsname{{48.24}{124}{}{proposition.48.24}{Dirichlet-GRH-v930-Volume-IX.pdf}}
\expandafter\def\csname r@thm:entireinterval310\endcsname{{21.14}{52}{}{proposition.21.14}{Dirichlet-GRH-v930-Volume-IV.pdf}}
\expandafter\def\csname r@thm:equallinks869\endcsname{{134.4}{84}{}{proposition.134.4}{Dirichlet-GRH-v930-Volume-XXV.pdf}}
\expandafter\def\csname r@thm:evenplanecompanion560\endcsname{{59.2}{3}{}{proposition.59.2}{Dirichlet-GRH-v930-Volume-XII.pdf}}
\expandafter\def\csname r@thm:evenroot358\endcsname{{36.6}{5}{}{proposition.36.6}{Dirichlet-GRH-v930-Volume-VII.pdf}}
\expandafter\def\csname r@thm:eventower572\endcsname{{60.5}{26}{}{proposition.60.5}{Dirichlet-GRH-v930-Volume-XII.pdf}}
\expandafter\def\csname r@thm:everycyclic663\endcsname{{26.27}{67}{}{proposition.26.27}{Dirichlet-GRH-v930-Volume-V.pdf}}
\expandafter\def\csname r@thm:everydiagonal536\endcsname{{113.10}{91}{}{proposition.113.10}{Dirichlet-GRH-v930-Volume-XXI.pdf}}
\expandafter\def\csname r@thm:exact41\endcsname{{64.6}{8}{}{proposition.64.6}{Dirichlet-GRH-v930-Volume-XIII.pdf}}
\expandafter\def\csname r@thm:exactclassprimes652\endcsname{{57.62}{116}{}{proposition.57.62}{Dirichlet-GRH-v930-Volume-XI.pdf}}
\expandafter\def\csname r@thm:exactcore592\endcsname{{116.15}{37}{}{proposition.116.15}{Dirichlet-GRH-v930-Volume-XXII.pdf}}
\expandafter\def\csname r@thm:exactselmer614\endcsname{{62.18}{85}{}{proposition.62.18}{Dirichlet-GRH-v930-Volume-XII.pdf}}
\expandafter\def\csname r@thm:exactstricttrace470\endcsname{{47.8}{82}{}{proposition.47.8}{Dirichlet-GRH-v930-Volume-IX.pdf}}
\expandafter\def\csname r@thm:exactthree609\endcsname{{22.41}{95}{}{proposition.22.41}{Dirichlet-GRH-v930-Volume-IV.pdf}}
\expandafter\def\csname r@thm:exceptionalperiod766\endcsname{{51.5}{42}{}{proposition.51.5}{Dirichlet-GRH-v930-Volume-X.pdf}}
\expandafter\def\csname r@thm:exteriorcost620\endcsname{{117.38}{76}{}{proposition.117.38}{Dirichlet-GRH-v930-Volume-XXII.pdf}}
\expandafter\def\csname r@thm:factorcomp512\endcsname{{145.9}{87}{}{proposition.145.9}{Dirichlet-GRH-v930-Volume-XXVII.pdf}}
\expandafter\def\csname r@thm:falsedescent849\endcsname{{81.15}{18}{}{proposition.81.15}{Dirichlet-GRH-v930-Volume-XVI.pdf}}
\expandafter\def\csname r@thm:fewtranslates868\endcsname{{134.2}{80}{}{proposition.134.2}{Dirichlet-GRH-v930-Volume-XXV.pdf}}
\expandafter\def\csname r@thm:fiberlength380\endcsname{{127.7}{33}{}{proposition.127.7}{Dirichlet-GRH-v930-Volume-XXIV.pdf}}
\expandafter\def\csname r@thm:filteredperiod755\endcsname{{50.25}{31}{}{proposition.50.25}{Dirichlet-GRH-v930-Volume-X.pdf}}
\expandafter\def\csname r@thm:finite31\endcsname{{30.1}{3}{}{proposition.30.1}{Dirichlet-GRH-v930-Volume-VI.pdf}}
\expandafter\def\csname r@thm:finiteflag546\endcsname{{48.4}{106}{}{proposition.48.4}{Dirichlet-GRH-v930-Volume-IX.pdf}}
\expandafter\def\csname r@thm:finitegram709\endcsname{{142.9}{43}{}{proposition.142.9}{Dirichlet-GRH-v930-Volume-XXVII.pdf}}
\expandafter\def\csname r@thm:finitelevel406\endcsname{{41.28}{56}{}{proposition.41.28}{Dirichlet-GRH-v930-Volume-VIII.pdf}}
\expandafter\def\csname r@thm:finiteunitreduce930\endcsname{{139.38}{140}{}{proposition.139.38}{Dirichlet-GRH-v930-Volume-XXVI.pdf}}
\expandafter\def\csname r@thm:finiteweights515\endcsname{{22.22}{80}{}{proposition.22.22}{Dirichlet-GRH-v930-Volume-IV.pdf}}
\expandafter\def\csname r@thm:first72\endcsname{{32.15}{40}{}{proposition.32.15}{Dirichlet-GRH-v930-Volume-VI.pdf}}
\expandafter\def\csname r@thm:firstflag659\endcsname{{117.91}{122}{}{proposition.117.91}{Dirichlet-GRH-v930-Volume-XXII.pdf}}
\expandafter\def\csname r@thm:firstgram890\endcsname{{137.2}{45}{}{proposition.137.2}{Dirichlet-GRH-v930-Volume-XXVI.pdf}}
\expandafter\def\csname r@thm:firsthol817\endcsname{{77.22}{24}{}{proposition.77.22}{Dirichlet-GRH-v930-Volume-XV.pdf}}
\expandafter\def\csname r@thm:firstlift708\endcsname{{142.5}{40}{}{proposition.142.5}{Dirichlet-GRH-v930-Volume-XXVII.pdf}}
\expandafter\def\csname r@thm:firstmismatch921\endcsname{{139.22}{117}{}{proposition.139.22}{Dirichlet-GRH-v930-Volume-XXVI.pdf}}
\expandafter\def\csname r@thm:firstnormalsum686\endcsname{{131.14}{25}{}{proposition.131.14}{Dirichlet-GRH-v930-Volume-XXV.pdf}}
\expandafter\def\csname r@thm:firstslice395\endcsname{{128.19}{71}{}{proposition.128.19}{Dirichlet-GRH-v930-Volume-XXIV.pdf}}
\expandafter\def\csname r@thm:firstvertex394\endcsname{{128.15}{68}{}{proposition.128.15}{Dirichlet-GRH-v930-Volume-XXIV.pdf}}
\expandafter\def\csname r@thm:fittingcost718\endcsname{{45.4}{40}{}{proposition.45.4}{Dirichlet-GRH-v930-Volume-IX.pdf}}
\expandafter\def\csname r@thm:fixedadic332\endcsname{{83.15}{64}{}{proposition.83.15}{Dirichlet-GRH-v930-Volume-XVI.pdf}}
\expandafter\def\csname r@thm:fixedconductor781\endcsname{{51.20}{55}{}{proposition.51.20}{Dirichlet-GRH-v930-Volume-X.pdf}}
\expandafter\def\csname r@thm:fixedexterior479\endcsname{{143.8}{55}{}{proposition.143.8}{Dirichlet-GRH-v930-Volume-XXVII.pdf}}
\expandafter\def\csname r@thm:fixedsectors493\endcsname{{68.37}{111}{}{proposition.68.37}{Dirichlet-GRH-v930-Volume-XIII.pdf}}
\expandafter\def\csname r@thm:flagcancel659\endcsname{{117.92}{123}{}{proposition.117.92}{Dirichlet-GRH-v930-Volume-XXII.pdf}}
\expandafter\def\csname r@thm:flagdoubling808\endcsname{{73.14}{86}{}{proposition.73.14}{Dirichlet-GRH-v930-Volume-XIV.pdf}}
\expandafter\def\csname r@thm:flagfiber438\endcsname{{55.16}{32}{}{proposition.55.16}{Dirichlet-GRH-v930-Volume-XI.pdf}}
\expandafter\def\csname r@thm:flagfourth658\endcsname{{117.90}{121}{}{proposition.117.90}{Dirichlet-GRH-v930-Volume-XXII.pdf}}
\expandafter\def\csname r@thm:flagisometry544\endcsname{{113.24}{103}{}{proposition.113.24}{Dirichlet-GRH-v930-Volume-XXI.pdf}}
\expandafter\def\csname r@thm:flagmeasure922\endcsname{{139.23}{119}{}{proposition.139.23}{Dirichlet-GRH-v930-Volume-XXVI.pdf}}
\expandafter\def\csname r@thm:flagproject383\endcsname{{127.12}{38}{}{proposition.127.12}{Dirichlet-GRH-v930-Volume-XXIV.pdf}}
\expandafter\def\csname r@thm:flagrecovery526\endcsname{{96.3}{74}{}{proposition.96.3}{Dirichlet-GRH-v930-Volume-XVIII.pdf}}
\expandafter\def\csname r@thm:flagthick385\endcsname{{40.26}{24}{}{proposition.40.26}{Dirichlet-GRH-v930-Volume-VIII.pdf}}
\expandafter\def\csname r@thm:flagvaluation385\endcsname{{40.27}{25}{}{proposition.40.27}{Dirichlet-GRH-v930-Volume-VIII.pdf}}
\expandafter\def\csname r@thm:flagvanish344\endcsname{{86.40}{31}{}{proposition.86.40}{Dirichlet-GRH-v930-Volume-XVII.pdf}}
\expandafter\def\csname r@thm:flattorus853\endcsname{{82.3}{28}{}{proposition.82.3}{Dirichlet-GRH-v930-Volume-XVI.pdf}}
\expandafter\def\csname r@thm:fourblock854\endcsname{{82.6}{32}{}{proposition.82.6}{Dirichlet-GRH-v930-Volume-XVI.pdf}}
\expandafter\def\csname r@thm:fourkernels622\endcsname{{117.43}{82}{}{proposition.117.43}{Dirichlet-GRH-v930-Volume-XXII.pdf}}
\expandafter\def\csname r@thm:frameorientation611\endcsname{{117.22}{59}{}{proposition.117.22}{Dirichlet-GRH-v930-Volume-XXII.pdf}}
\expandafter\def\csname r@thm:framepair509\endcsname{{144.23}{79}{}{proposition.144.23}{Dirichlet-GRH-v930-Volume-XXVII.pdf}}
\expandafter\def\csname r@thm:frameperiods446\endcsname{{141.3}{4}{}{proposition.141.3}{Dirichlet-GRH-v930-Volume-XXVII.pdf}}
\expandafter\def\csname r@thm:freqnode875\endcsname{{136.4}{5}{}{proposition.136.4}{Dirichlet-GRH-v930-Volume-XXVI.pdf}}
\expandafter\def\csname r@thm:frobeniusjet774\endcsname{{16.35}{67}{}{proposition.16.35}{Dirichlet-GRH-v930-Volume-III.pdf}}
\expandafter\def\csname r@thm:full36\endcsname{{31.8}{20}{}{proposition.31.8}{Dirichlet-GRH-v930-Volume-VI.pdf}}
\expandafter\def\csname r@thm:fullfirst898\endcsname{{137.18}{62}{}{proposition.137.18}{Dirichlet-GRH-v930-Volume-XXVI.pdf}}
\expandafter\def\csname r@thm:fullperiodcutoff769\endcsname{{51.7}{44}{}{proposition.51.7}{Dirichlet-GRH-v930-Volume-X.pdf}}
\expandafter\def\csname r@thm:fullrankseparation533\endcsname{{96.28}{93}{}{proposition.96.28}{Dirichlet-GRH-v930-Volume-XVIII.pdf}}
\expandafter\def\csname r@thm:fullrational645\endcsname{{117.83}{117}{}{proposition.117.83}{Dirichlet-GRH-v930-Volume-XXII.pdf}}
\expandafter\def\csname r@thm:fullrealfirst915\endcsname{{139.11}{104}{}{proposition.139.11}{Dirichlet-GRH-v930-Volume-XXVI.pdf}}
\expandafter\def\csname r@thm:fullrealperiod698\endcsname{{132.12}{60}{}{proposition.132.12}{Dirichlet-GRH-v930-Volume-XXV.pdf}}
\expandafter\def\csname r@thm:fullsignal452\endcsname{{141.15}{16}{}{proposition.141.15}{Dirichlet-GRH-v930-Volume-XXVII.pdf}}
\expandafter\def\csname r@thm:gain52\endcsname{{64.14}{15}{}{proposition.64.14}{Dirichlet-GRH-v930-Volume-XIII.pdf}}
\expandafter\def\csname r@thm:gaussianbudget401\endcsname{{41.9}{37}{}{proposition.41.9}{Dirichlet-GRH-v930-Volume-VIII.pdf}}
\expandafter\def\csname r@thm:genuskernel601\endcsname{{22.34}{90}{}{proposition.22.34}{Dirichlet-GRH-v930-Volume-IV.pdf}}
\expandafter\def\csname r@thm:genusranks665\endcsname{{26.31}{70}{}{proposition.26.31}{Dirichlet-GRH-v930-Volume-V.pdf}}
\expandafter\def\csname r@thm:globalamplitudetail700\endcsname{{132.17}{63}{}{proposition.132.17}{Dirichlet-GRH-v930-Volume-XXV.pdf}}
\expandafter\def\csname r@thm:globalequalrank826\endcsname{{77.46}{52}{}{proposition.77.46}{Dirichlet-GRH-v930-Volume-XV.pdf}}
\expandafter\def\csname r@thm:globalgram711\endcsname{{142.15}{48}{}{proposition.142.15}{Dirichlet-GRH-v930-Volume-XXVII.pdf}}
\expandafter\def\csname r@thm:globalhecke423\endcsname{{44.39}{32}{}{proposition.44.39}{Dirichlet-GRH-v930-Volume-IX.pdf}}
\expandafter\def\csname r@thm:globalperiodcontent448\endcsname{{141.9}{10}{}{proposition.141.9}{Dirichlet-GRH-v930-Volume-XXVII.pdf}}
\expandafter\def\csname r@thm:globaltheta43\endcsname{{64.8}{10}{}{proposition.64.8}{Dirichlet-GRH-v930-Volume-XIII.pdf}}
\expandafter\def\csname r@thm:gluedcost811\endcsname{{73.27}{98}{}{proposition.73.27}{Dirichlet-GRH-v930-Volume-XIV.pdf}}
\expandafter\def\csname r@thm:gluedperiodmap844\endcsname{{81.2}{4}{}{proposition.81.2}{Dirichlet-GRH-v930-Volume-XVI.pdf}}
\expandafter\def\csname r@thm:gram25\endcsname{{70.11}{12}{}{proposition.70.11}{Dirichlet-GRH-v930-Volume-XIV.pdf}}
\expandafter\def\csname r@thm:gramrestriction683\endcsname{{131.5}{18}{}{proposition.131.5}{Dirichlet-GRH-v930-Volume-XXV.pdf}}
\expandafter\def\csname r@thm:gramtail891\endcsname{{137.4}{47}{}{proposition.137.4}{Dirichlet-GRH-v930-Volume-XXVI.pdf}}
\expandafter\def\csname r@thm:graph274\endcsname{{94.27}{59}{}{proposition.94.27}{Dirichlet-GRH-v930-Volume-XVIII.pdf}}
\expandafter\def\csname r@thm:graphtransport337\endcsname{{71.21}{38}{}{proposition.71.21}{Dirichlet-GRH-v930-Volume-XIV.pdf}}
\expandafter\def\csname r@thm:grid354\endcsname{{99.22}{62}{}{proposition.99.22}{Dirichlet-GRH-v930-Volume-XIX.pdf}}
\expandafter\def\csname r@thm:growing\endcsname{{19.5}{5}{}{proposition.19.5}{Dirichlet-GRH-v930-Volume-IV.pdf}}
\expandafter\def\csname r@thm:halfclassdigits643\endcsname{{57.47}{103}{}{proposition.57.47}{Dirichlet-GRH-v930-Volume-XI.pdf}}
\expandafter\def\csname r@thm:halfdyadic529\endcsname{{57.26}{86}{}{proposition.57.26}{Dirichlet-GRH-v930-Volume-XI.pdf}}
\expandafter\def\csname r@thm:heckedivisor345\endcsname{{36.4}{4}{}{proposition.36.4}{Dirichlet-GRH-v930-Volume-VII.pdf}}
\expandafter\def\csname r@thm:heckeflag785\endcsname{{52.5}{63}{}{proposition.52.5}{Dirichlet-GRH-v930-Volume-X.pdf}}
\expandafter\def\csname r@thm:heckeseeds595\endcsname{{62.4}{69}{}{proposition.62.4}{Dirichlet-GRH-v930-Volume-XII.pdf}}
\expandafter\def\csname r@thm:hermitiancone900\endcsname{{138.2}{68}{}{proposition.138.2}{Dirichlet-GRH-v930-Volume-XXVI.pdf}}
\expandafter\def\csname r@thm:highallsmall790\endcsname{{17.9}{90}{}{proposition.17.9}{Dirichlet-GRH-v930-Volume-III.pdf}}
\expandafter\def\csname r@thm:highfinite788\endcsname{{17.2}{83}{}{proposition.17.2}{Dirichlet-GRH-v930-Volume-III.pdf}}
\expandafter\def\csname r@thm:highlift788\endcsname{{17.3}{83}{}{proposition.17.3}{Dirichlet-GRH-v930-Volume-III.pdf}}
\expandafter\def\csname r@thm:highmiddle791\endcsname{{17.12}{92}{}{proposition.17.12}{Dirichlet-GRH-v930-Volume-III.pdf}}
\expandafter\def\csname r@thm:highproduct792\endcsname{{17.14}{93}{}{proposition.17.14}{Dirichlet-GRH-v930-Volume-III.pdf}}
\expandafter\def\csname r@thm:highscalars793\endcsname{{17.17}{97}{}{proposition.17.17}{Dirichlet-GRH-v930-Volume-III.pdf}}
\expandafter\def\csname r@thm:highsmall789\endcsname{{17.6}{87}{}{proposition.17.6}{Dirichlet-GRH-v930-Volume-III.pdf}}
\expandafter\def\csname r@thm:horizontalstrict399\endcsname{{41.5}{33}{}{proposition.41.5}{Dirichlet-GRH-v930-Volume-VIII.pdf}}
\expandafter\def\csname r@thm:hyperplanerank603\endcsname{{117.2}{41}{}{proposition.117.2}{Dirichlet-GRH-v930-Volume-XXII.pdf}}
\expandafter\def\csname r@thm:hyperplanerecovery537\endcsname{{113.15}{95}{}{proposition.113.15}{Dirichlet-GRH-v930-Volume-XXI.pdf}}
\expandafter\def\csname r@thm:hyperplanesum610\endcsname{{117.18}{56}{}{proposition.117.18}{Dirichlet-GRH-v930-Volume-XXII.pdf}}
\expandafter\def\csname r@thm:idealordinary554\endcsname{{115.9}{10}{}{proposition.115.9}{Dirichlet-GRH-v930-Volume-XXII.pdf}}
\expandafter\def\csname r@thm:idealprecision608\endcsname{{22.37}{92}{}{proposition.22.37}{Dirichlet-GRH-v930-Volume-IV.pdf}}
\expandafter\def\csname r@thm:identitycoefficient551\endcsname{{113.26}{105}{}{proposition.113.26}{Dirichlet-GRH-v930-Volume-XXI.pdf}}
\expandafter\def\csname r@thm:identityfrequency929\endcsname{{139.35}{138}{}{proposition.139.35}{Dirichlet-GRH-v930-Volume-XXVI.pdf}}
\expandafter\def\csname r@thm:imagedet355\endcsname{{127.3}{28}{}{proposition.127.3}{Dirichlet-GRH-v930-Volume-XXIV.pdf}}
\expandafter\def\csname r@thm:incidenceannihilator413\endcsname{{42.20}{80}{}{proposition.42.20}{Dirichlet-GRH-v930-Volume-VIII.pdf}}
\expandafter\def\csname r@thm:inertfull908\endcsname{{138.24}{91}{}{proposition.138.24}{Dirichlet-GRH-v930-Volume-XXVI.pdf}}
\expandafter\def\csname r@thm:inertgradient860\endcsname{{82.19}{48}{}{proposition.82.19}{Dirichlet-GRH-v930-Volume-XVI.pdf}}
\expandafter\def\csname r@thm:inertiaann654\endcsname{{62.24}{90}{}{proposition.62.24}{Dirichlet-GRH-v930-Volume-XII.pdf}}
\expandafter\def\csname r@thm:inertiatrace406\endcsname{{41.26}{54}{}{proposition.41.26}{Dirichlet-GRH-v930-Volume-VIII.pdf}}
\expandafter\def\csname r@thm:inertjoint907\endcsname{{138.21}{87}{}{proposition.138.21}{Dirichlet-GRH-v930-Volume-XXVI.pdf}}
\expandafter\def\csname r@thm:inertlength902\endcsname{{138.8}{73}{}{proposition.138.8}{Dirichlet-GRH-v930-Volume-XXVI.pdf}}
\expandafter\def\csname r@thm:inertmass462\endcsname{{141.24}{24}{}{proposition.141.24}{Dirichlet-GRH-v930-Volume-XXVII.pdf}}
\expandafter\def\csname r@thm:inertneighbor901\endcsname{{138.5}{71}{}{proposition.138.5}{Dirichlet-GRH-v930-Volume-XXVI.pdf}}
\expandafter\def\csname r@thm:inertprime906\endcsname{{138.19}{85}{}{proposition.138.19}{Dirichlet-GRH-v930-Volume-XXVI.pdf}}
\expandafter\def\csname r@thm:inertquadratic443\endcsname{{55.27}{43}{}{proposition.55.27}{Dirichlet-GRH-v930-Volume-XI.pdf}}
\expandafter\def\csname r@thm:inerttail904\endcsname{{138.13}{80}{}{proposition.138.13}{Dirichlet-GRH-v930-Volume-XXVI.pdf}}
\expandafter\def\csname r@thm:inertterminal903\endcsname{{138.10}{75}{}{proposition.138.10}{Dirichlet-GRH-v930-Volume-XXVI.pdf}}
\expandafter\def\csname r@thm:inertzero905\endcsname{{138.16}{83}{}{proposition.138.16}{Dirichlet-GRH-v930-Volume-XXVI.pdf}}
\expandafter\def\csname r@thm:initialindex396\endcsname{{128.21}{73}{}{proposition.128.21}{Dirichlet-GRH-v930-Volume-XXIV.pdf}}
\expandafter\def\csname r@thm:initialmoment396\endcsname{{128.22}{74}{}{proposition.128.22}{Dirichlet-GRH-v930-Volume-XXIV.pdf}}
\expandafter\def\csname r@thm:integerlayers333\endcsname{{22.9}{69}{}{proposition.22.9}{Dirichlet-GRH-v930-Volume-IV.pdf}}
\expandafter\def\csname r@thm:integerpurity560\endcsname{{59.3}{4}{}{proposition.59.3}{Dirichlet-GRH-v930-Volume-XII.pdf}}
\expandafter\def\csname r@thm:integralblock656\endcsname{{22.57}{107}{}{proposition.22.57}{Dirichlet-GRH-v930-Volume-IV.pdf}}
\expandafter\def\csname r@thm:integralcomparison845\endcsname{{81.5}{7}{}{proposition.81.5}{Dirichlet-GRH-v930-Volume-XVI.pdf}}
\expandafter\def\csname r@thm:integralfamily551\endcsname{{113.27}{106}{}{proposition.113.27}{Dirichlet-GRH-v930-Volume-XXI.pdf}}
\expandafter\def\csname r@thm:integralfreq892\endcsname{{137.7}{49}{}{proposition.137.7}{Dirichlet-GRH-v930-Volume-XXVI.pdf}}
\expandafter\def\csname r@thm:integralleading857\endcsname{{82.11}{39}{}{proposition.82.11}{Dirichlet-GRH-v930-Volume-XVI.pdf}}
\expandafter\def\csname r@thm:integrallifts759\endcsname{{16.5}{39}{}{proposition.16.5}{Dirichlet-GRH-v930-Volume-III.pdf}}
\expandafter\def\csname r@thm:integralmass881\endcsname{{136.18}{19}{}{proposition.136.18}{Dirichlet-GRH-v930-Volume-XXVI.pdf}}
\expandafter\def\csname r@thm:integralrestrictions627\endcsname{{96.39}{102}{}{proposition.96.39}{Dirichlet-GRH-v930-Volume-XVIII.pdf}}
\expandafter\def\csname r@thm:integralroots547\endcsname{{48.7}{110}{}{proposition.48.7}{Dirichlet-GRH-v930-Volume-IX.pdf}}
\expandafter\def\csname r@thm:integratedpeak815\endcsname{{77.15}{17}{}{proposition.77.15}{Dirichlet-GRH-v930-Volume-XV.pdf}}
\expandafter\def\csname r@thm:interior795\endcsname{{17.21}{101}{}{proposition.17.21}{Dirichlet-GRH-v930-Volume-III.pdf}}
\expandafter\def\csname r@thm:intermediateimage392\endcsname{{128.8}{62}{}{proposition.128.8}{Dirichlet-GRH-v930-Volume-XXIV.pdf}}
\expandafter\def\csname r@thm:intersection279\endcsname{{19.25}{21}{}{proposition.19.25}{Dirichlet-GRH-v930-Volume-IV.pdf}}
\expandafter\def\csname r@thm:intrinsicplane567\endcsname{{59.16}{16}{}{proposition.59.16}{Dirichlet-GRH-v930-Volume-XII.pdf}}
\expandafter\def\csname r@thm:irrcurvesupport409\endcsname{{42.7}{68}{}{proposition.42.7}{Dirichlet-GRH-v930-Volume-VIII.pdf}}
\expandafter\def\csname r@thm:iwahori370\endcsname{{37.28}{61}{}{proposition.37.28}{Dirichlet-GRH-v930-Volume-VII.pdf}}
\expandafter\def\csname r@thm:jetlength797\endcsname{{17.25}{108}{}{proposition.17.25}{Dirichlet-GRH-v930-Volume-III.pdf}}
\expandafter\def\csname r@thm:jetnorm487\endcsname{{68.27}{104}{}{proposition.68.27}{Dirichlet-GRH-v930-Volume-XIII.pdf}}
\expandafter\def\csname r@thm:jetscale863\endcsname{{84.5}{80}{}{proposition.84.5}{Dirichlet-GRH-v930-Volume-XVI.pdf}}
\expandafter\def\csname r@thm:jetsmith485\endcsname{{68.19}{99}{}{proposition.68.19}{Dirichlet-GRH-v930-Volume-XIII.pdf}}
\expandafter\def\csname r@thm:jointminima813\endcsname{{77.8}{9}{}{proposition.77.8}{Dirichlet-GRH-v930-Volume-XV.pdf}}
\expandafter\def\csname r@thm:jointprimes484\endcsname{{47.26}{94}{}{proposition.47.26}{Dirichlet-GRH-v930-Volume-IX.pdf}}
\expandafter\def\csname r@thm:jointtamerow725\endcsname{{45.23}{60}{}{proposition.45.23}{Dirichlet-GRH-v930-Volume-IX.pdf}}
\expandafter\def\csname r@thm:labelder494\endcsname{{68.39}{113}{}{proposition.68.39}{Dirichlet-GRH-v930-Volume-XIII.pdf}}
\expandafter\def\csname r@thm:labelledcup765\endcsname{{16.21}{53}{}{proposition.16.21}{Dirichlet-GRH-v930-Volume-III.pdf}}
\expandafter\def\csname r@thm:largeintersection419\endcsname{{44.19}{19}{}{proposition.44.19}{Dirichlet-GRH-v930-Volume-IX.pdf}}
\expandafter\def\csname r@thm:lattice32\endcsname{{121.18}{49}{}{proposition.121.18}{Dirichlet-GRH-v930-Volume-XXIII.pdf}}
\expandafter\def\csname r@thm:latticemean690\endcsname{{131.26}{36}{}{proposition.131.26}{Dirichlet-GRH-v930-Volume-XXV.pdf}}
\expandafter\def\csname r@thm:laws268\endcsname{{27.2}{76}{}{proposition.27.2}{Dirichlet-GRH-v930-Volume-V.pdf}}
\expandafter\def\csname r@thm:lengthmultiplicity550\endcsname{{48.14}{116}{}{proposition.48.14}{Dirichlet-GRH-v930-Volume-IX.pdf}}
\expandafter\def\csname r@thm:levelstrict345\endcsname{{22.17}{75}{}{proposition.22.17}{Dirichlet-GRH-v930-Volume-IV.pdf}}
\expandafter\def\csname r@thm:liftphase604\endcsname{{117.6}{44}{}{proposition.117.6}{Dirichlet-GRH-v930-Volume-XXII.pdf}}
\expandafter\def\csname r@thm:limitkernel407\endcsname{{41.30}{58}{}{proposition.41.30}{Dirichlet-GRH-v930-Volume-VIII.pdf}}
\expandafter\def\csname r@thm:limitselmer407\endcsname{{41.32}{60}{}{proposition.41.32}{Dirichlet-GRH-v930-Volume-VIII.pdf}}
\expandafter\def\csname r@thm:localfinite732\endcsname{{13.5}{64}{}{proposition.13.5}{Dirichlet-GRH-v930-Volume-II.pdf}}
\expandafter\def\csname r@thm:localphase605\endcsname{{117.8}{46}{}{proposition.117.8}{Dirichlet-GRH-v930-Volume-XXII.pdf}}
\expandafter\def\csname r@thm:localrecovery632\endcsname{{117.59}{96}{}{proposition.117.59}{Dirichlet-GRH-v930-Volume-XXII.pdf}}
\expandafter\def\csname r@thm:loginterval319\endcsname{{22.1}{58}{}{proposition.22.1}{Dirichlet-GRH-v930-Volume-IV.pdf}}
\expandafter\def\csname r@thm:logunits319\endcsname{{28.11}{96}{}{proposition.28.11}{Dirichlet-GRH-v930-Volume-V.pdf}}
\expandafter\def\csname r@thm:longcharacter787\endcsname{{52.10}{68}{}{proposition.52.10}{Dirichlet-GRH-v930-Volume-X.pdf}}
\expandafter\def\csname r@thm:longframe867\endcsname{{84.14}{91}{}{proposition.84.14}{Dirichlet-GRH-v930-Volume-XVI.pdf}}
\expandafter\def\csname r@thm:longpoint779\endcsname{{16.40}{74}{}{proposition.16.40}{Dirichlet-GRH-v930-Volume-III.pdf}}
\expandafter\def\csname r@thm:longrootchart422\endcsname{{44.34}{29}{}{proposition.44.34}{Dirichlet-GRH-v930-Volume-IX.pdf}}
\expandafter\def\csname r@thm:lowerfibers384\endcsname{{40.22}{20}{}{proposition.40.22}{Dirichlet-GRH-v930-Volume-VIII.pdf}}
\expandafter\def\csname r@thm:mainorbit776\endcsname{{16.38}{70}{}{proposition.16.38}{Dirichlet-GRH-v930-Volume-III.pdf}}
\expandafter\def\csname r@thm:markedapartments528\endcsname{{96.10}{80}{}{proposition.96.10}{Dirichlet-GRH-v930-Volume-XVIII.pdf}}
\expandafter\def\csname r@thm:markedbottom847\endcsname{{81.10}{12}{}{proposition.81.10}{Dirichlet-GRH-v930-Volume-XVI.pdf}}
\expandafter\def\csname r@thm:markedcurve366\endcsname{{37.13}{46}{}{proposition.37.13}{Dirichlet-GRH-v930-Volume-VII.pdf}}
\expandafter\def\csname r@thm:markeddefinite820\endcsname{{77.32}{34}{}{proposition.77.32}{Dirichlet-GRH-v930-Volume-XV.pdf}}
\expandafter\def\csname r@thm:markedfourier625\endcsname{{96.34}{98}{}{proposition.96.34}{Dirichlet-GRH-v930-Volume-XVIII.pdf}}
\expandafter\def\csname r@thm:markedgauss386\endcsname{{127.15}{42}{}{proposition.127.15}{Dirichlet-GRH-v930-Volume-XXIV.pdf}}
\expandafter\def\csname r@thm:markedglobal339\endcsname{{71.25}{41}{}{proposition.71.25}{Dirichlet-GRH-v930-Volume-XIV.pdf}}
\expandafter\def\csname r@thm:markedimage393\endcsname{{128.11}{65}{}{proposition.128.11}{Dirichlet-GRH-v930-Volume-XXIV.pdf}}
\expandafter\def\csname r@thm:markedindex564\endcsname{{115.19}{18}{}{proposition.115.19}{Dirichlet-GRH-v930-Volume-XXII.pdf}}
\expandafter\def\csname r@thm:markedloss747\endcsname{{15.12}{17}{}{proposition.15.12}{Dirichlet-GRH-v930-Volume-III.pdf}}
\expandafter\def\csname r@thm:markedmatrix715\endcsname{{47.36}{101}{}{proposition.47.36}{Dirichlet-GRH-v930-Volume-IX.pdf}}
\expandafter\def\csname r@thm:markedperiod802\endcsname{{52.17}{75}{}{proposition.52.17}{Dirichlet-GRH-v930-Volume-X.pdf}}
\expandafter\def\csname r@thm:markedpresentation420\endcsname{{44.22}{21}{}{proposition.44.22}{Dirichlet-GRH-v930-Volume-IX.pdf}}
\expandafter\def\csname r@thm:markedranks623\endcsname{{117.46}{85}{}{proposition.117.46}{Dirichlet-GRH-v930-Volume-XXII.pdf}}
\expandafter\def\csname r@thm:markedsecond703\endcsname{{133.10}{71}{}{proposition.133.10}{Dirichlet-GRH-v930-Volume-XXV.pdf}}
\expandafter\def\csname r@thm:markedstrict337\endcsname{{22.11}{70}{}{proposition.22.11}{Dirichlet-GRH-v930-Volume-IV.pdf}}
\expandafter\def\csname r@thm:markedtablegram387\endcsname{{127.18}{45}{}{proposition.127.18}{Dirichlet-GRH-v930-Volume-XXIV.pdf}}
\expandafter\def\csname r@thm:markedtableintegral387\endcsname{{127.19}{47}{}{proposition.127.19}{Dirichlet-GRH-v930-Volume-XXIV.pdf}}
\expandafter\def\csname r@thm:markedtype801\endcsname{{17.33}{119}{}{proposition.17.33}{Dirichlet-GRH-v930-Volume-III.pdf}}
\expandafter\def\csname r@thm:marksections835\endcsname{{79.3}{76}{}{proposition.79.3}{Dirichlet-GRH-v930-Volume-XV.pdf}}
\expandafter\def\csname r@thm:matchedcoordinates631\endcsname{{117.56}{93}{}{proposition.117.56}{Dirichlet-GRH-v930-Volume-XXII.pdf}}
\expandafter\def\csname r@thm:matrixden574\endcsname{{60.11}{31}{}{proposition.60.11}{Dirichlet-GRH-v930-Volume-XII.pdf}}
\expandafter\def\csname r@thm:matrixshift456\endcsname{{120.17}{28}{}{proposition.120.17}{Dirichlet-GRH-v930-Volume-XXIII.pdf}}
\expandafter\def\csname r@thm:meangauss504\endcsname{{144.17}{74}{}{proposition.144.17}{Dirichlet-GRH-v930-Volume-XXVII.pdf}}
\expandafter\def\csname r@thm:middlerange773\endcsname{{51.13}{48}{}{proposition.51.13}{Dirichlet-GRH-v930-Volume-X.pdf}}
\expandafter\def\csname r@thm:middleterminal570\endcsname{{115.22}{20}{}{proposition.115.22}{Dirichlet-GRH-v930-Volume-XXII.pdf}}
\expandafter\def\csname r@thm:minorstrict373\endcsname{{38.3}{72}{}{proposition.38.3}{Dirichlet-GRH-v930-Volume-VII.pdf}}
\expandafter\def\csname r@thm:missedmass671\endcsname{{141.36}{32}{}{proposition.141.36}{Dirichlet-GRH-v930-Volume-XXVII.pdf}}
\expandafter\def\csname r@thm:mixedclasses634\endcsname{{57.40}{97}{}{proposition.57.40}{Dirichlet-GRH-v930-Volume-XI.pdf}}
\expandafter\def\csname r@thm:mixedcup772\endcsname{{16.33}{64}{}{proposition.16.33}{Dirichlet-GRH-v930-Volume-III.pdf}}
\expandafter\def\csname r@thm:mixedflags653\endcsname{{117.86}{119}{}{proposition.117.86}{Dirichlet-GRH-v930-Volume-XXII.pdf}}
\expandafter\def\csname r@thm:mixedfourier334\endcsname{{83.19}{68}{}{proposition.83.19}{Dirichlet-GRH-v930-Volume-XVI.pdf}}
\expandafter\def\csname r@thm:mixedhalf556\endcsname{{115.15}{15}{}{proposition.115.15}{Dirichlet-GRH-v930-Volume-XXII.pdf}}
\expandafter\def\csname r@thm:mixedmass876\endcsname{{136.8}{8}{}{proposition.136.8}{Dirichlet-GRH-v930-Volume-XXVI.pdf}}
\expandafter\def\csname r@thm:mixedopen885\endcsname{{136.28}{31}{}{proposition.136.28}{Dirichlet-GRH-v930-Volume-XXVI.pdf}}
\expandafter\def\csname r@thm:mixedordinary541\endcsname{{113.22}{101}{}{proposition.113.22}{Dirichlet-GRH-v930-Volume-XXI.pdf}}
\expandafter\def\csname r@thm:mixedperiod876\endcsname{{136.7}{7}{}{proposition.136.7}{Dirichlet-GRH-v930-Volume-XXVI.pdf}}
\expandafter\def\csname r@thm:mixedplane578\endcsname{{60.20}{39}{}{proposition.60.20}{Dirichlet-GRH-v930-Volume-XII.pdf}}
\expandafter\def\csname r@thm:mixedtrace359\endcsname{{36.11}{10}{}{proposition.36.11}{Dirichlet-GRH-v930-Volume-VII.pdf}}
\expandafter\def\csname r@thm:mixedtrace569\endcsname{{59.23}{22}{}{proposition.59.23}{Dirichlet-GRH-v930-Volume-XII.pdf}}
\expandafter\def\csname r@thm:modulepatch408\endcsname{{42.3}{64}{}{proposition.42.3}{Dirichlet-GRH-v930-Volume-VIII.pdf}}
\expandafter\def\csname r@thm:mon44\endcsname{{32.6}{28}{}{proposition.32.6}{Dirichlet-GRH-v930-Volume-VI.pdf}}
\expandafter\def\csname r@thm:monodromy476\endcsname{{47.23}{92}{}{proposition.47.23}{Dirichlet-GRH-v930-Volume-IX.pdf}}
\expandafter\def\csname r@thm:movingprimes428\endcsname{{54.18}{17}{}{proposition.54.18}{Dirichlet-GRH-v930-Volume-XI.pdf}}
\expandafter\def\csname r@thm:negativeinterval740\endcsname{{15.3}{4}{}{proposition.15.3}{Dirichlet-GRH-v930-Volume-III.pdf}}
\expandafter\def\csname r@thm:negfit398\endcsname{{41.2}{31}{}{proposition.41.2}{Dirichlet-GRH-v930-Volume-VIII.pdf}}
\expandafter\def\csname r@thm:neggraph794\endcsname{{17.19}{99}{}{proposition.17.19}{Dirichlet-GRH-v930-Volume-III.pdf}}
\expandafter\def\csname r@thm:neighborindex356\endcsname{{90.15}{96}{}{proposition.90.15}{Dirichlet-GRH-v930-Volume-XVII.pdf}}
\expandafter\def\csname r@thm:neighbour93\endcsname{{94.11}{46}{}{proposition.94.11}{Dirichlet-GRH-v930-Volume-XVIII.pdf}}
\expandafter\def\csname r@thm:nocorner40\endcsname{{32.5}{27}{}{proposition.32.5}{Dirichlet-GRH-v930-Volume-VI.pdf}}
\expandafter\def\csname r@thm:noncompacttail697\endcsname{{132.9}{57}{}{proposition.132.9}{Dirichlet-GRH-v930-Volume-XXV.pdf}}
\expandafter\def\csname r@thm:nonspheretail884\endcsname{{136.25}{28}{}{proposition.136.25}{Dirichlet-GRH-v930-Volume-XXVI.pdf}}
\expandafter\def\csname r@thm:nonsplitprefix466\endcsname{{56.11}{58}{}{proposition.56.11}{Dirichlet-GRH-v930-Volume-XI.pdf}}
\expandafter\def\csname r@thm:nonzero74\endcsname{{66.7}{50}{}{proposition.66.7}{Dirichlet-GRH-v930-Volume-XIII.pdf}}
\expandafter\def\csname r@thm:norm28\endcsname{{24.9}{9}{}{proposition.24.9}{Dirichlet-GRH-v930-Volume-V.pdf}}
\expandafter\def\csname r@thm:norm73\endcsname{{66.3}{48}{}{proposition.66.3}{Dirichlet-GRH-v930-Volume-XIII.pdf}}
\expandafter\def\csname r@thm:normalfirst828\endcsname{{78.4}{58}{}{proposition.78.4}{Dirichlet-GRH-v930-Volume-XV.pdf}}
\expandafter\def\csname r@thm:normalframe389\endcsname{{127.25}{54}{}{proposition.127.25}{Dirichlet-GRH-v930-Volume-XXIV.pdf}}
\expandafter\def\csname r@thm:normalizedcut699\endcsname{{132.14}{61}{}{proposition.132.14}{Dirichlet-GRH-v930-Volume-XXV.pdf}}
\expandafter\def\csname r@thm:normalizedimage357\endcsname{{90.18}{98}{}{proposition.90.18}{Dirichlet-GRH-v930-Volume-XVII.pdf}}
\expandafter\def\csname r@thm:normalizedprimes600\endcsname{{22.31}{87}{}{proposition.22.31}{Dirichlet-GRH-v930-Volume-IV.pdf}}
\expandafter\def\csname r@thm:normaljet388\endcsname{{127.22}{51}{}{proposition.127.22}{Dirichlet-GRH-v930-Volume-XXIV.pdf}}
\expandafter\def\csname r@thm:normsum327\endcsname{{22.6}{65}{}{proposition.22.6}{Dirichlet-GRH-v930-Volume-IV.pdf}}
\expandafter\def\csname r@thm:octicrange524\endcsname{{57.18}{80}{}{proposition.57.18}{Dirichlet-GRH-v930-Volume-XI.pdf}}
\expandafter\def\csname r@thm:oddinteger677\endcsname{{11.29}{26}{}{proposition.11.29}{Dirichlet-GRH-v930-Volume-II.pdf}}
\expandafter\def\csname r@thm:oddquarticfield539\endcsname{{22.28}{85}{}{proposition.22.28}{Dirichlet-GRH-v930-Volume-IV.pdf}}
\expandafter\def\csname r@thm:oddrelations411\endcsname{{42.15}{76}{}{proposition.42.15}{Dirichlet-GRH-v930-Volume-VIII.pdf}}
\expandafter\def\csname r@thm:oddrootlocalization430\endcsname{{46.8}{67}{}{proposition.46.8}{Dirichlet-GRH-v930-Volume-IX.pdf}}
\expandafter\def\csname r@thm:oddselmer739\endcsname{{13.20}{77}{}{proposition.13.20}{Dirichlet-GRH-v930-Volume-II.pdf}}
\expandafter\def\csname r@thm:oddstrict675\endcsname{{11.24}{22}{}{proposition.11.24}{Dirichlet-GRH-v930-Volume-II.pdf}}
\expandafter\def\csname r@thm:oddsupport410\endcsname{{42.10}{71}{}{proposition.42.10}{Dirichlet-GRH-v930-Volume-VIII.pdf}}
\expandafter\def\csname r@thm:oldcolumn745\endcsname{{50.16}{20}{}{proposition.50.16}{Dirichlet-GRH-v930-Volume-X.pdf}}
\expandafter\def\csname r@thm:oldtorustraces743\endcsname{{15.8}{12}{}{proposition.15.8}{Dirichlet-GRH-v930-Volume-III.pdf}}
\expandafter\def\csname r@thm:oldweightzero748\endcsname{{15.15}{20}{}{proposition.15.15}{Dirichlet-GRH-v930-Volume-III.pdf}}
\expandafter\def\csname r@thm:openenergy870\endcsname{{134.6}{86}{}{proposition.134.6}{Dirichlet-GRH-v930-Volume-XXV.pdf}}
\expandafter\def\csname r@thm:openindex269\endcsname{{27.3}{77}{}{proposition.27.3}{Dirichlet-GRH-v930-Volume-V.pdf}}
\expandafter\def\csname r@thm:openpair688\endcsname{{131.20}{30}{}{proposition.131.20}{Dirichlet-GRH-v930-Volume-XXV.pdf}}
\expandafter\def\csname r@thm:oppositeaverage660\endcsname{{117.94}{125}{}{proposition.117.94}{Dirichlet-GRH-v930-Volume-XXII.pdf}}
\expandafter\def\csname r@thm:orbitaverage532\endcsname{{96.24}{90}{}{proposition.96.24}{Dirichlet-GRH-v930-Volume-XVIII.pdf}}
\expandafter\def\csname r@thm:orbitfactors871\endcsname{{134.8}{89}{}{proposition.134.8}{Dirichlet-GRH-v930-Volume-XXV.pdf}}
\expandafter\def\csname r@thm:orbitlift764\endcsname{{16.18}{50}{}{proposition.16.18}{Dirichlet-GRH-v930-Volume-III.pdf}}
\expandafter\def\csname r@thm:orbitnoncancellation533\endcsname{{96.30}{94}{}{proposition.96.30}{Dirichlet-GRH-v930-Volume-XVIII.pdf}}
\expandafter\def\csname r@thm:orderedconductor437\endcsname{{55.12}{28}{}{proposition.55.12}{Dirichlet-GRH-v930-Volume-XI.pdf}}
\expandafter\def\csname r@thm:ordinarycross568\endcsname{{59.20}{20}{}{proposition.59.20}{Dirichlet-GRH-v930-Volume-XII.pdf}}
\expandafter\def\csname r@thm:ordinarydigits514\endcsname{{22.19}{78}{}{proposition.22.19}{Dirichlet-GRH-v930-Volume-IV.pdf}}
\expandafter\def\csname r@thm:ordinaryperiod756\endcsname{{50.28}{35}{}{proposition.50.28}{Dirichlet-GRH-v930-Volume-X.pdf}}
\expandafter\def\csname r@thm:ordinarysingular544\endcsname{{113.25}{104}{}{proposition.113.25}{Dirichlet-GRH-v930-Volume-XXI.pdf}}
\expandafter\def\csname r@thm:ordinaryunit552\endcsname{{115.3}{5}{}{proposition.115.3}{Dirichlet-GRH-v930-Volume-XXII.pdf}}
\expandafter\def\csname r@thm:orientationmoment613\endcsname{{117.28}{65}{}{proposition.117.28}{Dirichlet-GRH-v930-Volume-XXII.pdf}}
\expandafter\def\csname r@thm:originalfinite424\endcsname{{44.41}{34}{}{proposition.44.41}{Dirichlet-GRH-v930-Volume-IX.pdf}}
\expandafter\def\csname r@thm:originalorder424\endcsname{{44.44}{35}{}{proposition.44.44}{Dirichlet-GRH-v930-Volume-IX.pdf}}
\expandafter\def\csname r@thm:orthcollapse384\endcsname{{40.24}{22}{}{proposition.40.24}{Dirichlet-GRH-v930-Volume-VIII.pdf}}
\expandafter\def\csname r@thm:outsidedifferent453\endcsname{{141.18}{18}{}{proposition.141.18}{Dirichlet-GRH-v930-Volume-XXVII.pdf}}
\expandafter\def\csname r@thm:paircount691\endcsname{{131.29}{41}{}{proposition.131.29}{Dirichlet-GRH-v930-Volume-XXV.pdf}}
\expandafter\def\csname r@thm:pairedtable861\endcsname{{82.21}{51}{}{proposition.82.21}{Dirichlet-GRH-v930-Volume-XVI.pdf}}
\expandafter\def\csname r@thm:pairedtrace651\endcsname{{62.23}{89}{}{proposition.62.23}{Dirichlet-GRH-v930-Volume-XII.pdf}}
\expandafter\def\csname r@thm:paritydescent358\endcsname{{36.8}{7}{}{proposition.36.8}{Dirichlet-GRH-v930-Volume-VII.pdf}}
\expandafter\def\csname r@thm:partnerlifts531\endcsname{{96.19}{86}{}{proposition.96.19}{Dirichlet-GRH-v930-Volume-XVIII.pdf}}
\expandafter\def\csname r@thm:patchrelations408\endcsname{{42.4}{64}{}{proposition.42.4}{Dirichlet-GRH-v930-Volume-VIII.pdf}}
\expandafter\def\csname r@thm:pathnorm352\endcsname{{119.11}{11}{}{proposition.119.11}{Dirichlet-GRH-v930-Volume-XXIII.pdf}}
\expandafter\def\csname r@thm:periodden312\endcsname{{125.3}{6}{}{proposition.125.3}{Dirichlet-GRH-v930-Volume-XXIV.pdf}}
\expandafter\def\csname r@thm:periodminors400\endcsname{{128.25}{77}{}{proposition.128.25}{Dirichlet-GRH-v930-Volume-XXIV.pdf}}
\expandafter\def\csname r@thm:pintable340\endcsname{{83.25}{72}{}{proposition.83.25}{Dirichlet-GRH-v930-Volume-XVI.pdf}}
\expandafter\def\csname r@thm:pivotintegral619\endcsname{{117.35}{73}{}{proposition.117.35}{Dirichlet-GRH-v930-Volume-XXII.pdf}}
\expandafter\def\csname r@thm:pivotrecursion660\endcsname{{117.95}{125}{}{proposition.117.95}{Dirichlet-GRH-v930-Volume-XXII.pdf}}
\expandafter\def\csname r@thm:planedepth566\endcsname{{59.13}{14}{}{proposition.59.13}{Dirichlet-GRH-v930-Volume-XII.pdf}}
\expandafter\def\csname r@thm:planedesc381\endcsname{{40.14}{11}{}{proposition.40.14}{Dirichlet-GRH-v930-Volume-VIII.pdf}}
\expandafter\def\csname r@thm:planehaar383\endcsname{{127.11}{37}{}{proposition.127.11}{Dirichlet-GRH-v930-Volume-XXIV.pdf}}
\expandafter\def\csname r@thm:planeunits565\endcsname{{59.11}{12}{}{proposition.59.11}{Dirichlet-GRH-v930-Volume-XII.pdf}}
\expandafter\def\csname r@thm:poincaregram704\endcsname{{133.12}{72}{}{proposition.133.12}{Dirichlet-GRH-v930-Volume-XXV.pdf}}
\expandafter\def\csname r@thm:pointlattice353\endcsname{{99.20}{60}{}{proposition.99.20}{Dirichlet-GRH-v930-Volume-XIX.pdf}}
\expandafter\def\csname r@thm:pol49\endcsname{{64.10}{12}{}{proposition.64.10}{Dirichlet-GRH-v930-Volume-XIII.pdf}}
\expandafter\def\csname r@thm:polarweights593\endcsname{{61.18}{61}{}{proposition.61.18}{Dirichlet-GRH-v930-Volume-XII.pdf}}
\expandafter\def\csname r@thm:positivecompare576\endcsname{{115.25}{22}{}{proposition.115.25}{Dirichlet-GRH-v930-Volume-XXII.pdf}}
\expandafter\def\csname r@thm:positivecore696\endcsname{{132.6}{54}{}{proposition.132.6}{Dirichlet-GRH-v930-Volume-XXV.pdf}}
\expandafter\def\csname r@thm:positivejets482\endcsname{{68.11}{94}{}{proposition.68.11}{Dirichlet-GRH-v930-Volume-XIII.pdf}}
\expandafter\def\csname r@thm:potentialperiod746\endcsname{{50.19}{24}{}{proposition.50.19}{Dirichlet-GRH-v930-Volume-X.pdf}}
\expandafter\def\csname r@thm:prefixpack928\endcsname{{139.33}{135}{}{proposition.139.33}{Dirichlet-GRH-v930-Volume-XXVI.pdf}}
\expandafter\def\csname r@thm:present31\endcsname{{30.4}{5}{}{proposition.30.4}{Dirichlet-GRH-v930-Volume-VI.pdf}}
\expandafter\def\csname r@thm:preserveselmer727\endcsname{{12.24}{48}{}{proposition.12.24}{Dirichlet-GRH-v930-Volume-II.pdf}}
\expandafter\def\csname r@thm:primemass650\endcsname{{22.53}{104}{}{proposition.22.53}{Dirichlet-GRH-v930-Volume-IV.pdf}}
\expandafter\def\csname r@thm:primemoment555\endcsname{{115.12}{12}{}{proposition.115.12}{Dirichlet-GRH-v930-Volume-XXII.pdf}}
\expandafter\def\csname r@thm:primeprojections670\endcsname{{141.34}{31}{}{proposition.141.34}{Dirichlet-GRH-v930-Volume-XXVII.pdf}}
\expandafter\def\csname r@thm:primes367\endcsname{{37.17}{51}{}{proposition.37.17}{Dirichlet-GRH-v930-Volume-VII.pdf}}
\expandafter\def\csname r@thm:primitivegram463\endcsname{{141.27}{26}{}{proposition.141.27}{Dirichlet-GRH-v930-Volume-XXVII.pdf}}
\expandafter\def\csname r@thm:productmeasure647\endcsname{{57.56}{111}{}{proposition.57.56}{Dirichlet-GRH-v930-Volume-XI.pdf}}
\expandafter\def\csname r@thm:projectedcups742\endcsname{{50.11}{12}{}{proposition.50.11}{Dirichlet-GRH-v930-Volume-X.pdf}}
\expandafter\def\csname r@thm:projectedinteger668\endcsname{{141.30}{28}{}{proposition.141.30}{Dirichlet-GRH-v930-Volume-XXVII.pdf}}
\expandafter\def\csname r@thm:projectedrank705\endcsname{{133.18}{77}{}{proposition.133.18}{Dirichlet-GRH-v930-Volume-XXV.pdf}}
\expandafter\def\csname r@thm:projectionindex784\endcsname{{52.2}{59}{}{proposition.52.2}{Dirichlet-GRH-v930-Volume-X.pdf}}
\expandafter\def\csname r@thm:projectiveperiod337\endcsname{{83.21}{69}{}{proposition.83.21}{Dirichlet-GRH-v930-Volume-XVI.pdf}}
\expandafter\def\csname r@thm:puncturedisometry518\endcsname{{146.5}{96}{}{proposition.146.5}{Dirichlet-GRH-v930-Volume-XXVII.pdf}}
\expandafter\def\csname r@thm:puncturedrank498\endcsname{{144.8}{67}{}{proposition.144.8}{Dirichlet-GRH-v930-Volume-XXVII.pdf}}
\expandafter\def\csname r@thm:purityblocks549\endcsname{{48.13}{115}{}{proposition.48.13}{Dirichlet-GRH-v930-Volume-IX.pdf}}
\expandafter\def\csname r@thm:quadraticexponent587\endcsname{{61.13}{57}{}{proposition.61.13}{Dirichlet-GRH-v930-Volume-XII.pdf}}
\expandafter\def\csname r@thm:quadraticfamily583\endcsname{{61.1}{48}{}{proposition.61.1}{Dirichlet-GRH-v930-Volume-XII.pdf}}
\expandafter\def\csname r@thm:quantlog646\endcsname{{57.53}{108}{}{proposition.57.53}{Dirichlet-GRH-v930-Volume-XI.pdf}}
\expandafter\def\csname r@thm:ramification549\endcsname{{48.12}{114}{}{proposition.48.12}{Dirichlet-GRH-v930-Volume-IX.pdf}}
\expandafter\def\csname r@thm:ramifiedcoset324\endcsname{{123.13}{94}{}{proposition.123.13}{Dirichlet-GRH-v930-Volume-XXIII.pdf}}
\expandafter\def\csname r@thm:ramifiedcut782\endcsname{{16.43}{79}{}{proposition.16.43}{Dirichlet-GRH-v930-Volume-III.pdf}}
\expandafter\def\csname r@thm:rankcutoff778\endcsname{{51.18}{52}{}{proposition.51.18}{Dirichlet-GRH-v930-Volume-X.pdf}}
\expandafter\def\csname r@thm:rankdichotomy473\endcsname{{47.14}{86}{}{proposition.47.14}{Dirichlet-GRH-v930-Volume-IX.pdf}}
\expandafter\def\csname r@thm:rankfourselmer558\endcsname{{48.22}{123}{}{proposition.48.22}{Dirichlet-GRH-v930-Volume-IX.pdf}}
\expandafter\def\csname r@thm:rankmarks602\endcsname{{116.17}{38}{}{proposition.116.17}{Dirichlet-GRH-v930-Volume-XXII.pdf}}
\expandafter\def\csname r@thm:rankone354\endcsname{{104.36}{31}{}{proposition.104.36}{Dirichlet-GRH-v930-Volume-XX.pdf}}
\expandafter\def\csname r@thm:rankonecost455\endcsname{{120.14}{25}{}{proposition.120.14}{Dirichlet-GRH-v930-Volume-XXIII.pdf}}
\expandafter\def\csname r@thm:rankonepair736\endcsname{{50.3}{4}{}{proposition.50.3}{Dirichlet-GRH-v930-Volume-X.pdf}}
\expandafter\def\csname r@thm:rankonespherical737\endcsname{{50.7}{7}{}{proposition.50.7}{Dirichlet-GRH-v930-Volume-X.pdf}}
\expandafter\def\csname r@thm:rankoperator630\endcsname{{117.54}{91}{}{proposition.117.54}{Dirichlet-GRH-v930-Volume-XXII.pdf}}
\expandafter\def\csname r@thm:rationalflags680\endcsname{{130.6}{9}{}{proposition.130.6}{Dirichlet-GRH-v930-Volume-XXV.pdf}}
\expandafter\def\csname r@thm:rationalgram629\endcsname{{117.50}{88}{}{proposition.117.50}{Dirichlet-GRH-v930-Volume-XXII.pdf}}
\expandafter\def\csname r@thm:rationalinduction585\endcsname{{61.8}{54}{}{proposition.61.8}{Dirichlet-GRH-v930-Volume-XII.pdf}}
\expandafter\def\csname r@thm:rationalperiod342\endcsname{{86.38}{30}{}{proposition.86.38}{Dirichlet-GRH-v930-Volume-XVII.pdf}}
\expandafter\def\csname r@thm:rationalpolynomial538\endcsname{{22.25}{82}{}{proposition.22.25}{Dirichlet-GRH-v930-Volume-IV.pdf}}
\expandafter\def\csname r@thm:raw97\endcsname{{87.2}{33}{}{proposition.87.2}{Dirichlet-GRH-v930-Volume-XVII.pdf}}
\expandafter\def\csname r@thm:rawsaturation672\endcsname{{141.37}{34}{}{proposition.141.37}{Dirichlet-GRH-v930-Volume-XXVII.pdf}}
\expandafter\def\csname r@thm:rayconductor444\endcsname{{55.29}{45}{}{proposition.55.29}{Dirichlet-GRH-v930-Volume-XI.pdf}}
\expandafter\def\csname r@thm:raylength656\endcsname{{22.55}{106}{}{proposition.22.55}{Dirichlet-GRH-v930-Volume-IV.pdf}}
\expandafter\def\csname r@thm:realcutmass916\endcsname{{139.13}{106}{}{proposition.139.13}{Dirichlet-GRH-v930-Volume-XXVI.pdf}}
\expandafter\def\csname r@thm:realgram682\endcsname{{131.1}{14}{}{proposition.131.1}{Dirichlet-GRH-v930-Volume-XXV.pdf}}
\expandafter\def\csname r@thm:realgram710\endcsname{{142.13}{46}{}{proposition.142.13}{Dirichlet-GRH-v930-Volume-XXVII.pdf}}
\expandafter\def\csname r@thm:realnulldescent926\endcsname{{139.30}{130}{}{proposition.139.30}{Dirichlet-GRH-v930-Volume-XXVI.pdf}}
\expandafter\def\csname r@thm:realpair685\endcsname{{131.11}{23}{}{proposition.131.11}{Dirichlet-GRH-v930-Volume-XXV.pdf}}
\expandafter\def\csname r@thm:realpositionball912\endcsname{{139.7}{99}{}{proposition.139.7}{Dirichlet-GRH-v930-Volume-XXVI.pdf}}
\expandafter\def\csname r@thm:realquadrics911\endcsname{{139.3}{95}{}{proposition.139.3}{Dirichlet-GRH-v930-Volume-XXVI.pdf}}
\expandafter\def\csname r@thm:realremotetail914\endcsname{{139.10}{102}{}{proposition.139.10}{Dirichlet-GRH-v930-Volume-XXVI.pdf}}
\expandafter\def\csname r@thm:realtwist404\endcsname{{41.17}{46}{}{proposition.41.17}{Dirichlet-GRH-v930-Volume-VIII.pdf}}
\expandafter\def\csname r@thm:rectangulartheta530\endcsname{{96.16}{83}{}{proposition.96.16}{Dirichlet-GRH-v930-Volume-XVIII.pdf}}
\expandafter\def\csname r@thm:reduciblecollision435\endcsname{{55.3}{20}{}{proposition.55.3}{Dirichlet-GRH-v930-Volume-XI.pdf}}
\expandafter\def\csname r@thm:reduciblefiber376\endcsname{{38.15}{80}{}{proposition.38.15}{Dirichlet-GRH-v930-Volume-VII.pdf}}
\expandafter\def\csname r@thm:reducibleminors378\endcsname{{40.7}{6}{}{proposition.40.7}{Dirichlet-GRH-v930-Volume-VIII.pdf}}
\expandafter\def\csname r@thm:reducibleorder379\endcsname{{40.12}{9}{}{proposition.40.12}{Dirichlet-GRH-v930-Volume-VIII.pdf}}
\expandafter\def\csname r@thm:reduciblepropagation412\endcsname{{42.18}{78}{}{proposition.42.18}{Dirichlet-GRH-v930-Volume-VIII.pdf}}
\expandafter\def\csname r@thm:reduciblerigid379\endcsname{{40.10}{8}{}{proposition.40.10}{Dirichlet-GRH-v930-Volume-VIII.pdf}}
\expandafter\def\csname r@thm:regularcosets597\endcsname{{62.9}{76}{}{proposition.62.9}{Dirichlet-GRH-v930-Volume-XII.pdf}}
\expandafter\def\csname r@thm:regularind579\endcsname{{60.23}{43}{}{proposition.60.23}{Dirichlet-GRH-v930-Volume-XII.pdf}}
\expandafter\def\csname r@thm:relativeprogress325\endcsname{{22.3}{61}{}{proposition.22.3}{Dirichlet-GRH-v930-Volume-IV.pdf}}
\expandafter\def\csname r@thm:relhalf633\endcsname{{57.37}{95}{}{proposition.57.37}{Dirichlet-GRH-v930-Volume-XI.pdf}}
\expandafter\def\csname r@thm:remainingmain771\endcsname{{51.11}{47}{}{proposition.51.11}{Dirichlet-GRH-v930-Volume-X.pdf}}
\expandafter\def\csname r@thm:res100\endcsname{{11.17}{17}{}{proposition.11.17}{Dirichlet-GRH-v930-Volume-II.pdf}}
\expandafter\def\csname r@thm:rescurve364\endcsname{{37.4}{36}{}{proposition.37.4}{Dirichlet-GRH-v930-Volume-VII.pdf}}
\expandafter\def\csname r@thm:residualflag786\endcsname{{52.8}{66}{}{proposition.52.8}{Dirichlet-GRH-v930-Volume-X.pdf}}
\expandafter\def\csname r@thm:residualvaluation673\endcsname{{141.40}{36}{}{proposition.141.40}{Dirichlet-GRH-v930-Volume-XXVII.pdf}}
\expandafter\def\csname r@thm:residueconductor628\endcsname{{96.42}{105}{}{proposition.96.42}{Dirichlet-GRH-v930-Volume-XVIII.pdf}}
\expandafter\def\csname r@thm:residuejet492\endcsname{{68.33}{109}{}{proposition.68.33}{Dirichlet-GRH-v930-Volume-XIII.pdf}}
\expandafter\def\csname r@thm:residueorbits606\endcsname{{117.11}{50}{}{proposition.117.11}{Dirichlet-GRH-v930-Volume-XXII.pdf}}
\expandafter\def\csname r@thm:residuepositive749\endcsname{{15.19}{23}{}{proposition.15.19}{Dirichlet-GRH-v930-Volume-III.pdf}}
\expandafter\def\csname r@thm:restrictedtower848\endcsname{{81.13}{15}{}{proposition.81.13}{Dirichlet-GRH-v930-Volume-XVI.pdf}}
\expandafter\def\csname r@thm:restrictedtrace374\endcsname{{38.6}{74}{}{proposition.38.6}{Dirichlet-GRH-v930-Volume-VII.pdf}}
\expandafter\def\csname r@thm:restrictionrecursion626\endcsname{{96.37}{100}{}{proposition.96.37}{Dirichlet-GRH-v930-Volume-XVIII.pdf}}
\expandafter\def\csname r@thm:ribettrace720\endcsname{{45.11}{48}{}{proposition.45.11}{Dirichlet-GRH-v930-Volume-IX.pdf}}
\expandafter\def\csname r@thm:rigidtrace377\endcsname{{40.2}{2}{}{proposition.40.2}{Dirichlet-GRH-v930-Volume-VIII.pdf}}
\expandafter\def\csname r@thm:ring39\endcsname{{32.1}{23}{}{proposition.32.1}{Dirichlet-GRH-v930-Volume-VI.pdf}}
\expandafter\def\csname r@thm:rotationcharacters612\endcsname{{117.26}{63}{}{proposition.117.26}{Dirichlet-GRH-v930-Volume-XXII.pdf}}
\expandafter\def\csname r@thm:rotationperiod611\endcsname{{117.24}{61}{}{proposition.117.24}{Dirichlet-GRH-v930-Volume-XXII.pdf}}
\expandafter\def\csname r@thm:row277\endcsname{{94.31}{64}{}{proposition.94.31}{Dirichlet-GRH-v930-Volume-XVIII.pdf}}
\expandafter\def\csname r@thm:rowcontent335\endcsname{{126.10}{20}{}{proposition.126.10}{Dirichlet-GRH-v930-Volume-XXIV.pdf}}
\expandafter\def\csname r@thm:rowcuspidal809\endcsname{{73.17}{89}{}{proposition.73.17}{Dirichlet-GRH-v930-Volume-XIV.pdf}}
\expandafter\def\csname r@thm:rowdescent895\endcsname{{137.13}{56}{}{proposition.137.13}{Dirichlet-GRH-v930-Volume-XXVI.pdf}}
\expandafter\def\csname r@thm:rowfactor607\endcsname{{117.14}{52}{}{proposition.117.14}{Dirichlet-GRH-v930-Volume-XXII.pdf}}
\expandafter\def\csname r@thm:rowsumgram637\endcsname{{117.68}{104}{}{proposition.117.68}{Dirichlet-GRH-v930-Volume-XXII.pdf}}
\expandafter\def\csname r@thm:samelink687\endcsname{{131.17}{28}{}{proposition.131.17}{Dirichlet-GRH-v930-Volume-XXV.pdf}}
\expandafter\def\csname r@thm:scalaraugmentation519\endcsname{{146.9}{99}{}{proposition.146.9}{Dirichlet-GRH-v930-Volume-XXVII.pdf}}
\expandafter\def\csname r@thm:scalarconductor655\endcsname{{62.28}{94}{}{proposition.62.28}{Dirichlet-GRH-v930-Volume-XII.pdf}}
\expandafter\def\csname r@thm:scalarcost717\endcsname{{12.12}{39}{}{proposition.12.12}{Dirichlet-GRH-v930-Volume-II.pdf}}
\expandafter\def\csname r@thm:scalarinertiaplanes439\endcsname{{55.20}{36}{}{proposition.55.20}{Dirichlet-GRH-v930-Volume-XI.pdf}}
\expandafter\def\csname r@thm:scalarlevel467\endcsname{{56.14}{60}{}{proposition.56.14}{Dirichlet-GRH-v930-Volume-XI.pdf}}
\expandafter\def\csname r@thm:scalarmeromorphic770\endcsname{{16.29}{61}{}{proposition.16.29}{Dirichlet-GRH-v930-Volume-III.pdf}}
\expandafter\def\csname r@thm:scalarorbits332\endcsname{{119.4}{5}{}{proposition.119.4}{Dirichlet-GRH-v930-Volume-XXIII.pdf}}
\expandafter\def\csname r@thm:scalarseparate636\endcsname{{117.65}{101}{}{proposition.117.65}{Dirichlet-GRH-v930-Volume-XXII.pdf}}
\expandafter\def\csname r@thm:scalartrace859\endcsname{{82.16}{44}{}{proposition.82.16}{Dirichlet-GRH-v930-Volume-XVI.pdf}}
\expandafter\def\csname r@thm:scaleddefinite821\endcsname{{77.34}{38}{}{proposition.77.34}{Dirichlet-GRH-v930-Volume-XV.pdf}}
\expandafter\def\csname r@thm:schurmean692\endcsname{{131.32}{43}{}{proposition.131.32}{Dirichlet-GRH-v930-Volume-XXV.pdf}}
\expandafter\def\csname r@thm:secondallpositions918\endcsname{{139.18}{110}{}{proposition.139.18}{Dirichlet-GRH-v930-Volume-XXVI.pdf}}
\expandafter\def\csname r@thm:secondcomparison829\endcsname{{78.8}{62}{}{proposition.78.8}{Dirichlet-GRH-v930-Volume-XV.pdf}}
\expandafter\def\csname r@thm:secondglobal833\endcsname{{78.16}{72}{}{proposition.78.16}{Dirichlet-GRH-v930-Volume-XV.pdf}}
\expandafter\def\csname r@thm:secondlocal832\endcsname{{78.15}{70}{}{proposition.78.15}{Dirichlet-GRH-v930-Volume-XV.pdf}}
\expandafter\def\csname r@thm:secondpoint775\endcsname{{16.37}{68}{}{proposition.16.37}{Dirichlet-GRH-v930-Volume-III.pdf}}
\expandafter\def\csname r@thm:secondsharp920\endcsname{{139.21}{114}{}{proposition.139.21}{Dirichlet-GRH-v930-Volume-XXVI.pdf}}
\expandafter\def\csname r@thm:section110\endcsname{{86.31}{24}{}{proposition.86.31}{Dirichlet-GRH-v930-Volume-XVII.pdf}}
\expandafter\def\csname r@thm:seedweight580\endcsname{{60.27}{47}{}{proposition.60.27}{Dirichlet-GRH-v930-Volume-XII.pdf}}
\expandafter\def\csname r@thm:selectedopen874\endcsname{{136.1}{3}{}{proposition.136.1}{Dirichlet-GRH-v930-Volume-XXVI.pdf}}
\expandafter\def\csname r@thm:selftwist728\endcsname{{12.25}{50}{}{proposition.12.25}{Dirichlet-GRH-v930-Volume-II.pdf}}
\expandafter\def\csname r@thm:selmerlength368\endcsname{{37.20}{54}{}{proposition.37.20}{Dirichlet-GRH-v930-Volume-VII.pdf}}
\expandafter\def\csname r@thm:selmerlift369\endcsname{{37.25}{58}{}{proposition.37.25}{Dirichlet-GRH-v930-Volume-VII.pdf}}
\expandafter\def\csname r@thm:sexticweights669\endcsname{{26.34}{73}{}{proposition.26.34}{Dirichlet-GRH-v930-Volume-V.pdf}}
\expandafter\def\csname r@thm:shallowwave893\endcsname{{137.9}{51}{}{proposition.137.9}{Dirichlet-GRH-v930-Volume-XXVI.pdf}}
\expandafter\def\csname r@thm:shiftedjets481\endcsname{{68.8}{92}{}{proposition.68.8}{Dirichlet-GRH-v930-Volume-XIII.pdf}}
\expandafter\def\csname r@thm:shortpairs859\endcsname{{82.17}{45}{}{proposition.82.17}{Dirichlet-GRH-v930-Volume-XVI.pdf}}
\expandafter\def\csname r@thm:siegelclasses431\endcsname{{46.12}{70}{}{proposition.46.12}{Dirichlet-GRH-v930-Volume-IX.pdf}}
\expandafter\def\csname r@thm:signal77\endcsname{{66.26}{65}{}{proposition.66.26}{Dirichlet-GRH-v930-Volume-XIII.pdf}}
\expandafter\def\csname r@thm:signedcharacterloss432\endcsname{{128.29}{79}{}{proposition.128.29}{Dirichlet-GRH-v930-Volume-XXIV.pdf}}
\expandafter\def\csname r@thm:signedcompact621\endcsname{{117.41}{80}{}{proposition.117.41}{Dirichlet-GRH-v930-Volume-XXII.pdf}}
\expandafter\def\csname r@thm:signedcomplete488\endcsname{{68.30}{107}{}{proposition.68.30}{Dirichlet-GRH-v930-Volume-XIII.pdf}}
\expandafter\def\csname r@thm:signedcovariance478\endcsname{{143.5}{53}{}{proposition.143.5}{Dirichlet-GRH-v930-Volume-XXVII.pdf}}
\expandafter\def\csname r@thm:signedfibers331\endcsname{{99.12}{53}{}{proposition.99.12}{Dirichlet-GRH-v930-Volume-XIX.pdf}}
\expandafter\def\csname r@thm:signedopen702\endcsname{{133.7}{68}{}{proposition.133.7}{Dirichlet-GRH-v930-Volume-XXV.pdf}}
\expandafter\def\csname r@thm:signedreturn342\endcsname{{83.26}{73}{}{proposition.83.26}{Dirichlet-GRH-v930-Volume-XVI.pdf}}
\expandafter\def\csname r@thm:signedseries823\endcsname{{77.38}{44}{}{proposition.77.38}{Dirichlet-GRH-v930-Volume-XV.pdf}}
\expandafter\def\csname r@thm:signedsignal450\endcsname{{141.12}{13}{}{proposition.141.12}{Dirichlet-GRH-v930-Volume-XXVII.pdf}}
\expandafter\def\csname r@thm:signedtranslation520\endcsname{{146.12}{101}{}{proposition.146.12}{Dirichlet-GRH-v930-Volume-XXVII.pdf}}
\expandafter\def\csname r@thm:simultaneous307\endcsname{{34.22}{85}{}{proposition.34.22}{Dirichlet-GRH-v930-Volume-VI.pdf}}
\expandafter\def\csname r@thm:singleprojection704\endcsname{{133.13}{73}{}{proposition.133.13}{Dirichlet-GRH-v930-Volume-XXV.pdf}}
\expandafter\def\csname r@thm:singularflag878\endcsname{{136.12}{13}{}{proposition.136.12}{Dirichlet-GRH-v930-Volume-XXVI.pdf}}
\expandafter\def\csname r@thm:slthree363\endcsname{{36.26}{30}{}{proposition.36.26}{Dirichlet-GRH-v930-Volume-VII.pdf}}
\expandafter\def\csname r@thm:smallcollision457\endcsname{{55.34}{49}{}{proposition.55.34}{Dirichlet-GRH-v930-Volume-XI.pdf}}
\expandafter\def\csname r@thm:smallcomplement604\endcsname{{117.5}{43}{}{proposition.117.5}{Dirichlet-GRH-v930-Volume-XXII.pdf}}
\expandafter\def\csname r@thm:smallcoords777\endcsname{{51.16}{50}{}{proposition.51.16}{Dirichlet-GRH-v930-Volume-X.pdf}}
\expandafter\def\csname r@thm:smallframes590\endcsname{{116.8}{31}{}{proposition.116.8}{Dirichlet-GRH-v930-Volume-XXII.pdf}}
\expandafter\def\csname r@thm:smallorbits591\endcsname{{116.13}{35}{}{proposition.116.13}{Dirichlet-GRH-v930-Volume-XXII.pdf}}
\expandafter\def\csname r@thm:smallthird834\endcsname{{79.1}{74}{}{proposition.79.1}{Dirichlet-GRH-v930-Volume-XV.pdf}}
\expandafter\def\csname r@thm:smooth69\endcsname{{32.11}{35}{}{proposition.32.11}{Dirichlet-GRH-v930-Volume-VI.pdf}}
\expandafter\def\csname r@thm:sourcecommon563\endcsname{{59.9}{10}{}{proposition.59.9}{Dirichlet-GRH-v930-Volume-XII.pdf}}
\expandafter\def\csname r@thm:sparseapprox814\endcsname{{77.11}{12}{}{proposition.77.11}{Dirichlet-GRH-v930-Volume-XV.pdf}}
\expandafter\def\csname r@thm:sparsecups780\endcsname{{16.41}{76}{}{proposition.16.41}{Dirichlet-GRH-v930-Volume-III.pdf}}
\expandafter\def\csname r@thm:spectralquot499\endcsname{{47.33}{99}{}{proposition.47.33}{Dirichlet-GRH-v930-Volume-IX.pdf}}
\expandafter\def\csname r@thm:spheremean872\endcsname{{134.10}{94}{}{proposition.134.10}{Dirichlet-GRH-v930-Volume-XXV.pdf}}
\expandafter\def\csname r@thm:splitamb506\endcsname{{57.7}{71}{}{proposition.57.7}{Dirichlet-GRH-v930-Volume-XI.pdf}}
\expandafter\def\csname r@thm:splitattachment722\endcsname{{45.16}{54}{}{proposition.45.16}{Dirichlet-GRH-v930-Volume-IX.pdf}}
\expandafter\def\csname r@thm:splitcoherent391\endcsname{{128.6}{60}{}{proposition.128.6}{Dirichlet-GRH-v930-Volume-XXIV.pdf}}
\expandafter\def\csname r@thm:splitintegral573\endcsname{{60.8}{29}{}{proposition.60.8}{Dirichlet-GRH-v930-Volume-XII.pdf}}
\expandafter\def\csname r@thm:splitjets480\endcsname{{68.5}{90}{}{proposition.68.5}{Dirichlet-GRH-v930-Volume-XIII.pdf}}
\expandafter\def\csname r@thm:splitlattice811\endcsname{{73.24}{96}{}{proposition.73.24}{Dirichlet-GRH-v930-Volume-XIV.pdf}}
\expandafter\def\csname r@thm:splitparameters475\endcsname{{47.20}{90}{}{proposition.47.20}{Dirichlet-GRH-v930-Volume-IX.pdf}}
\expandafter\def\csname r@thm:splitprojectors599\endcsname{{62.15}{83}{}{proposition.62.15}{Dirichlet-GRH-v930-Volume-XII.pdf}}
\expandafter\def\csname r@thm:splitweights502\endcsname{{56.17}{63}{}{proposition.56.17}{Dirichlet-GRH-v930-Volume-XI.pdf}}
\expandafter\def\csname r@thm:splitwordfiber438\endcsname{{55.14}{30}{}{proposition.55.14}{Dirichlet-GRH-v930-Volume-XI.pdf}}
\expandafter\def\csname r@thm:squareconductor827\endcsname{{78.2}{56}{}{proposition.78.2}{Dirichlet-GRH-v930-Volume-XV.pdf}}
\expandafter\def\csname r@thm:src109\endcsname{{87.7}{37}{}{proposition.87.7}{Dirichlet-GRH-v930-Volume-XVII.pdf}}
\expandafter\def\csname r@thm:stable102\endcsname{{20.5}{27}{}{proposition.20.5}{Dirichlet-GRH-v930-Volume-IV.pdf}}
\expandafter\def\csname r@thm:stableprime371\endcsname{{37.32}{65}{}{proposition.37.32}{Dirichlet-GRH-v930-Volume-VII.pdf}}
\expandafter\def\csname r@thm:stackedheight819\endcsname{{77.29}{32}{}{proposition.77.29}{Dirichlet-GRH-v930-Volume-XV.pdf}}
\expandafter\def\csname r@thm:standardfirst899\endcsname{{137.21}{65}{}{proposition.137.21}{Dirichlet-GRH-v930-Volume-XXVI.pdf}}
\expandafter\def\csname r@thm:strictbasedescent416\endcsname{{44.7}{8}{}{proposition.44.7}{Dirichlet-GRH-v930-Volume-IX.pdf}}
\expandafter\def\csname r@thm:strictcorner415\endcsname{{44.3}{4}{}{proposition.44.3}{Dirichlet-GRH-v930-Volume-IX.pdf}}
\expandafter\def\csname r@thm:strictcyclic657\endcsname{{26.23}{64}{}{proposition.26.23}{Dirichlet-GRH-v930-Volume-V.pdf}}
\expandafter\def\csname r@thm:strictdigits329\endcsname{{22.7}{66}{}{proposition.22.7}{Dirichlet-GRH-v930-Volume-IV.pdf}}
\expandafter\def\csname r@thm:strictrootring434\endcsname{{46.20}{75}{}{proposition.46.20}{Dirichlet-GRH-v930-Volume-IX.pdf}}
\expandafter\def\csname r@thm:strongtrace474\endcsname{{47.17}{88}{}{proposition.47.17}{Dirichlet-GRH-v930-Volume-IX.pdf}}
\expandafter\def\csname r@thm:ststrict616\endcsname{{32.23}{47}{}{proposition.32.23}{Dirichlet-GRH-v930-Volume-VI.pdf}}
\expandafter\def\csname r@thm:subcriticalperiod763\endcsname{{51.2}{40}{}{proposition.51.2}{Dirichlet-GRH-v930-Volume-X.pdf}}
\expandafter\def\csname r@thm:sunitbase507\endcsname{{57.10}{73}{}{proposition.57.10}{Dirichlet-GRH-v930-Volume-XI.pdf}}
\expandafter\def\csname r@thm:symplecticclasses429\endcsname{{46.4}{65}{}{proposition.46.4}{Dirichlet-GRH-v930-Volume-IX.pdf}}
\expandafter\def\csname r@thm:symplocal655\endcsname{{62.27}{93}{}{proposition.62.27}{Dirichlet-GRH-v930-Volume-XII.pdf}}
\expandafter\def\csname r@thm:sympordinary468\endcsname{{47.2}{77}{}{proposition.47.2}{Dirichlet-GRH-v930-Volume-IX.pdf}}
\expandafter\def\csname r@thm:symprelations469\endcsname{{47.5}{79}{}{proposition.47.5}{Dirichlet-GRH-v930-Volume-IX.pdf}}
\expandafter\def\csname r@thm:tamecups800\endcsname{{17.30}{116}{}{proposition.17.30}{Dirichlet-GRH-v930-Volume-III.pdf}}
\expandafter\def\csname r@thm:tamedivision618\endcsname{{22.48}{100}{}{proposition.22.48}{Dirichlet-GRH-v930-Volume-IV.pdf}}
\expandafter\def\csname r@thm:tameorder753\endcsname{{15.30}{33}{}{proposition.15.30}{Dirichlet-GRH-v930-Volume-III.pdf}}
\expandafter\def\csname r@thm:tameray662\endcsname{{26.24}{65}{}{proposition.26.24}{Dirichlet-GRH-v930-Volume-V.pdf}}
\expandafter\def\csname r@thm:tametower571\endcsname{{60.2}{24}{}{proposition.60.2}{Dirichlet-GRH-v930-Volume-XII.pdf}}
\expandafter\def\csname r@thm:tametransfer721\endcsname{{45.14}{51}{}{proposition.45.14}{Dirichlet-GRH-v930-Volume-IX.pdf}}
\expandafter\def\csname r@thm:target90\endcsname{{94.8}{43}{}{proposition.94.8}{Dirichlet-GRH-v930-Volume-XVIII.pdf}}
\expandafter\def\csname r@thm:tensortrace382\endcsname{{40.19}{16}{}{proposition.40.19}{Dirichlet-GRH-v930-Volume-VIII.pdf}}
\expandafter\def\csname r@thm:termdisc687\endcsname{{131.15}{26}{}{proposition.131.15}{Dirichlet-GRH-v930-Volume-XXV.pdf}}
\expandafter\def\csname r@thm:terminalclass640\endcsname{{117.77}{111}{}{proposition.117.77}{Dirichlet-GRH-v930-Volume-XXII.pdf}}
\expandafter\def\csname r@thm:terminaldegrees454\endcsname{{141.21}{21}{}{proposition.141.21}{Dirichlet-GRH-v930-Volume-XXVII.pdf}}
\expandafter\def\csname r@thm:terminaldirection866\endcsname{{84.12}{88}{}{proposition.84.12}{Dirichlet-GRH-v930-Volume-XVI.pdf}}
\expandafter\def\csname r@thm:terminaldivisor556\endcsname{{115.14}{14}{}{proposition.115.14}{Dirichlet-GRH-v930-Volume-XXII.pdf}}
\expandafter\def\csname r@thm:terminalflag807\endcsname{{73.11}{84}{}{proposition.73.11}{Dirichlet-GRH-v930-Volume-XIV.pdf}}
\expandafter\def\csname r@thm:terminalformula347\endcsname{{127.1}{26}{}{proposition.127.1}{Dirichlet-GRH-v930-Volume-XXIV.pdf}}
\expandafter\def\csname r@thm:terminalfourier877\endcsname{{136.9}{10}{}{proposition.136.9}{Dirichlet-GRH-v930-Volume-XXVI.pdf}}
\expandafter\def\csname r@thm:terminalmarked390\endcsname{{128.2}{56}{}{proposition.128.2}{Dirichlet-GRH-v930-Volume-XXIV.pdf}}
\expandafter\def\csname r@thm:terminalmeet681\endcsname{{130.11}{13}{}{proposition.130.11}{Dirichlet-GRH-v930-Volume-XXV.pdf}}
\expandafter\def\csname r@thm:terminalouter639\endcsname{{117.74}{108}{}{proposition.117.74}{Dirichlet-GRH-v930-Volume-XXII.pdf}}
\expandafter\def\csname r@thm:terminalrange886\endcsname{{136.31}{33}{}{proposition.136.31}{Dirichlet-GRH-v930-Volume-XXVI.pdf}}
\expandafter\def\csname r@thm:terminalreps522\endcsname{{146.19}{106}{}{proposition.146.19}{Dirichlet-GRH-v930-Volume-XXVII.pdf}}
\expandafter\def\csname r@thm:terminalscalar357\endcsname{{96.1}{72}{}{proposition.96.1}{Dirichlet-GRH-v930-Volume-XVIII.pdf}}
\expandafter\def\csname r@thm:terminaltensor927\endcsname{{139.31}{132}{}{proposition.139.31}{Dirichlet-GRH-v930-Volume-XXVI.pdf}}
\expandafter\def\csname r@thm:terminaltypes327\endcsname{{123.15}{96}{}{proposition.123.15}{Dirichlet-GRH-v930-Volume-XXIII.pdf}}
\expandafter\def\csname r@thm:terminalwalk338\endcsname{{71.24}{40}{}{proposition.71.24}{Dirichlet-GRH-v930-Volume-XIV.pdf}}
\expandafter\def\csname r@thm:testedcomparison850\endcsname{{81.18}{22}{}{proposition.81.18}{Dirichlet-GRH-v930-Volume-XVI.pdf}}
\expandafter\def\csname r@thm:thetaalternant486\endcsname{{68.22}{101}{}{proposition.68.22}{Dirichlet-GRH-v930-Volume-XIII.pdf}}
\expandafter\def\csname r@thm:thetacells340\endcsname{{71.27}{42}{}{proposition.71.27}{Dirichlet-GRH-v930-Volume-XIV.pdf}}
\expandafter\def\csname r@thm:thirdadelic842\endcsname{{79.18}{96}{}{proposition.79.18}{Dirichlet-GRH-v930-Volume-XV.pdf}}
\expandafter\def\csname r@thm:thirdglobal843\endcsname{{79.20}{100}{}{proposition.79.20}{Dirichlet-GRH-v930-Volume-XV.pdf}}
\expandafter\def\csname r@thm:thirdholomorphic841\endcsname{{79.15}{93}{}{proposition.79.15}{Dirichlet-GRH-v930-Volume-XV.pdf}}
\expandafter\def\csname r@thm:thirdlargenorm839\endcsname{{79.11}{88}{}{proposition.79.11}{Dirichlet-GRH-v930-Volume-XV.pdf}}
\expandafter\def\csname r@thm:thirdreal838\endcsname{{79.9}{85}{}{proposition.79.9}{Dirichlet-GRH-v930-Volume-XV.pdf}}
\expandafter\def\csname r@thm:thirdsmallnorm840\endcsname{{79.13}{90}{}{proposition.79.13}{Dirichlet-GRH-v930-Volume-XV.pdf}}
\expandafter\def\csname r@thm:thirdsplitprod836\endcsname{{79.6}{79}{}{proposition.79.6}{Dirichlet-GRH-v930-Volume-XV.pdf}}
\expandafter\def\csname r@thm:thirdstandard837\endcsname{{79.8}{83}{}{proposition.79.8}{Dirichlet-GRH-v930-Volume-XV.pdf}}
\expandafter\def\csname r@thm:tiedlength729\endcsname{{12.30}{53}{}{proposition.12.30}{Dirichlet-GRH-v930-Volume-II.pdf}}
\expandafter\def\csname r@thm:topdescent723\endcsname{{45.20}{58}{}{proposition.45.20}{Dirichlet-GRH-v930-Volume-IX.pdf}}
\expandafter\def\csname r@thm:torusnorm852\endcsname{{82.1}{25}{}{proposition.82.1}{Dirichlet-GRH-v930-Volume-XVI.pdf}}
\expandafter\def\csname r@thm:totaldiagonal540\endcsname{{113.19}{99}{}{proposition.113.19}{Dirichlet-GRH-v930-Volume-XXI.pdf}}
\expandafter\def\csname r@thm:totalresidual649\endcsname{{22.51}{103}{}{proposition.22.51}{Dirichlet-GRH-v930-Volume-IV.pdf}}
\expandafter\def\csname r@thm:tower104\endcsname{{19.16}{13}{}{proposition.19.16}{Dirichlet-GRH-v930-Volume-IV.pdf}}
\expandafter\def\csname r@thm:tracefiber420\endcsname{{44.24}{22}{}{proposition.44.24}{Dirichlet-GRH-v930-Volume-IX.pdf}}
\expandafter\def\csname r@thm:tracefitting548\endcsname{{48.10}{113}{}{proposition.48.10}{Dirichlet-GRH-v930-Volume-IX.pdf}}
\expandafter\def\csname r@thm:traceideal362\endcsname{{36.22}{24}{}{proposition.36.22}{Dirichlet-GRH-v930-Volume-VII.pdf}}
\expandafter\def\csname r@thm:traceoverlap602\endcsname{{116.16}{38}{}{proposition.116.16}{Dirichlet-GRH-v930-Volume-XXII.pdf}}
\expandafter\def\csname r@thm:transfer380\endcsname{{127.6}{31}{}{proposition.127.6}{Dirichlet-GRH-v930-Volume-XXIV.pdf}}
\expandafter\def\csname r@thm:transitionrange796\endcsname{{17.23}{105}{}{proposition.17.23}{Dirichlet-GRH-v930-Volume-III.pdf}}
\expandafter\def\csname r@thm:translates706\endcsname{{133.20}{78}{}{proposition.133.20}{Dirichlet-GRH-v930-Volume-XXV.pdf}}
\expandafter\def\csname r@thm:transportproduct589\endcsname{{116.5}{29}{}{proposition.116.5}{Dirichlet-GRH-v930-Volume-XXII.pdf}}
\expandafter\def\csname r@thm:transverseseparate581\endcsname{{115.28}{25}{}{proposition.115.28}{Dirichlet-GRH-v930-Volume-XXII.pdf}}
\expandafter\def\csname r@thm:transversesum582\endcsname{{116.2}{27}{}{proposition.116.2}{Dirichlet-GRH-v930-Volume-XXII.pdf}}
\expandafter\def\csname r@thm:trianglebudget714\endcsname{{12.9}{37}{}{proposition.12.9}{Dirichlet-GRH-v930-Volume-II.pdf}}
\expandafter\def\csname r@thm:trunc15\endcsname{{11.5}{6}{}{proposition.11.5}{Dirichlet-GRH-v930-Volume-II.pdf}}
\expandafter\def\csname r@thm:twelftransfer855\endcsname{{82.7}{34}{}{proposition.82.7}{Dirichlet-GRH-v930-Volume-XVI.pdf}}
\expandafter\def\csname r@thm:twistedselmer730\endcsname{{12.32}{57}{}{proposition.12.32}{Dirichlet-GRH-v930-Volume-II.pdf}}
\expandafter\def\csname r@thm:two17\endcsname{{11.12}{13}{}{proposition.11.12}{Dirichlet-GRH-v930-Volume-II.pdf}}
\expandafter\def\csname r@thm:twocopy858\endcsname{{82.14}{41}{}{proposition.82.14}{Dirichlet-GRH-v930-Volume-XVI.pdf}}
\expandafter\def\csname r@thm:twodepth896\endcsname{{137.15}{58}{}{proposition.137.15}{Dirichlet-GRH-v930-Volume-XXVI.pdf}}
\expandafter\def\csname r@thm:twoenergies592\endcsname{{116.14}{36}{}{proposition.116.14}{Dirichlet-GRH-v930-Volume-XXII.pdf}}
\expandafter\def\csname r@thm:twolayers327\endcsname{{22.4}{62}{}{proposition.22.4}{Dirichlet-GRH-v930-Volume-IV.pdf}}
\expandafter\def\csname r@thm:twomarks557\endcsname{{48.17}{119}{}{proposition.48.17}{Dirichlet-GRH-v930-Volume-IX.pdf}}
\expandafter\def\csname r@thm:twopencils521\endcsname{{146.15}{104}{}{proposition.146.15}{Dirichlet-GRH-v930-Volume-XXVII.pdf}}
\expandafter\def\csname r@thm:twoprimeadjoints436\endcsname{{55.7}{24}{}{proposition.55.7}{Dirichlet-GRH-v930-Volume-XI.pdf}}
\expandafter\def\csname r@thm:typeblocks329\endcsname{{123.16}{98}{}{proposition.123.16}{Dirichlet-GRH-v930-Volume-XXIII.pdf}}
\expandafter\def\csname r@thm:typedindex799\endcsname{{52.13}{71}{}{proposition.52.13}{Dirichlet-GRH-v930-Volume-X.pdf}}
\expandafter\def\csname r@thm:uniformannihilator676\endcsname{{11.28}{24}{}{proposition.11.28}{Dirichlet-GRH-v930-Volume-II.pdf}}
\expandafter\def\csname r@thm:uniformdepth458\endcsname{{56.2}{51}{}{proposition.56.2}{Dirichlet-GRH-v930-Volume-XI.pdf}}
\expandafter\def\csname r@thm:uniformlift760\endcsname{{16.8}{42}{}{proposition.16.8}{Dirichlet-GRH-v930-Volume-III.pdf}}
\expandafter\def\csname r@thm:unitcoordinate751\endcsname{{15.25}{29}{}{proposition.15.25}{Dirichlet-GRH-v930-Volume-III.pdf}}
\expandafter\def\csname r@thm:unitlength505\endcsname{{57.5}{69}{}{proposition.57.5}{Dirichlet-GRH-v930-Volume-XI.pdf}}
\expandafter\def\csname r@thm:unitray472\endcsname{{57.2}{67}{}{proposition.57.2}{Dirichlet-GRH-v930-Volume-XI.pdf}}
\expandafter\def\csname r@thm:unitroots523\endcsname{{57.14}{77}{}{proposition.57.14}{Dirichlet-GRH-v930-Volume-XI.pdf}}
\expandafter\def\csname r@thm:unitweights364\endcsname{{37.1}{33}{}{proposition.37.1}{Dirichlet-GRH-v930-Volume-VII.pdf}}
\expandafter\def\csname r@thm:unmarkedtrace343\endcsname{{126.15}{25}{}{proposition.126.15}{Dirichlet-GRH-v930-Volume-XXIV.pdf}}
\expandafter\def\csname r@thm:unorderedcube588\endcsname{{61.16}{60}{}{proposition.61.16}{Dirichlet-GRH-v930-Volume-XII.pdf}}
\expandafter\def\csname r@thm:unramifiedperiod758\endcsname{{50.31}{37}{}{proposition.50.31}{Dirichlet-GRH-v930-Volume-X.pdf}}
\expandafter\def\csname r@thm:unrestrictedcups762\endcsname{{16.15}{47}{}{proposition.16.15}{Dirichlet-GRH-v930-Volume-III.pdf}}
\expandafter\def\csname r@thm:unrestrictedmass882\endcsname{{136.20}{22}{}{proposition.136.20}{Dirichlet-GRH-v930-Volume-XXVI.pdf}}
\expandafter\def\csname r@thm:variableclasses642\endcsname{{57.43}{100}{}{proposition.57.43}{Dirichlet-GRH-v930-Volume-XI.pdf}}
\expandafter\def\csname r@thm:variabledescent642\endcsname{{57.44}{100}{}{proposition.57.44}{Dirichlet-GRH-v930-Volume-XI.pdf}}
\expandafter\def\csname r@thm:vectorchain353\endcsname{{73.2}{74}{}{proposition.73.2}{Dirichlet-GRH-v930-Volume-XIV.pdf}}
\expandafter\def\csname r@thm:weightchart362\endcsname{{36.21}{23}{}{proposition.36.21}{Dirichlet-GRH-v930-Volume-VII.pdf}}
\expandafter\def\csname r@thm:weylmoments500\endcsname{{144.11}{69}{}{proposition.144.11}{Dirichlet-GRH-v930-Volume-XXVII.pdf}}
\expandafter\def\csname r@thm:weylperiod351\endcsname{{119.10}{10}{}{proposition.119.10}{Dirichlet-GRH-v930-Volume-XXIII.pdf}}
\expandafter\def\csname r@thm:weylrowjet449\endcsname{{120.12}{23}{}{proposition.120.12}{Dirichlet-GRH-v930-Volume-XXIII.pdf}}
\expandafter\def\csname r@thm:weyltransport864\endcsname{{84.8}{83}{}{proposition.84.8}{Dirichlet-GRH-v930-Volume-XVI.pdf}}
\expandafter\def\csname r@thm:whittakerintegral527\endcsname{{96.7}{77}{}{proposition.96.7}{Dirichlet-GRH-v930-Volume-XVIII.pdf}}
\expandafter\def\csname r@thm:whittakerspectrum527\endcsname{{96.5}{75}{}{proposition.96.5}{Dirichlet-GRH-v930-Volume-XVIII.pdf}}
\expandafter\def\csname r@thm:wholeblock661\endcsname{{22.58}{108}{}{proposition.22.58}{Dirichlet-GRH-v930-Volume-IV.pdf}}
\expandafter\def\csname r@thm:wild98\endcsname{{25.1}{24}{}{proposition.25.1}{Dirichlet-GRH-v930-Volume-V.pdf}}
\expandafter\def\csname r@thm:wordgcd510\endcsname{{145.3}{82}{}{proposition.145.3}{Dirichlet-GRH-v930-Volume-XXVII.pdf}}
\expandafter\def\csname r@thm:wordlower348\endcsname{{83.29}{75}{}{proposition.83.29}{Dirichlet-GRH-v930-Volume-XVI.pdf}}
\expandafter\def\csname r@thm:zero106\endcsname{{86.27}{22}{}{proposition.86.27}{Dirichlet-GRH-v930-Volume-XVII.pdf}}
\makeatother
\begin{document}
\setcounter{part}{0}
\fancyhead[L]{\small\color{muted}Dirichlet GRH \textperiodcentered\ Volume I}
\fancyhead[R]{\small\color{muted}Dirichlet zeros and power intervals}
\thispagestyle{plain}\begin{center}
{\LARGE\bfseries The generalized Riemann hypothesis\\[2pt]for Dirichlet \(L\)-functions}\\[6pt]
{\small 27 November 2026 \quad\textperiodcentered\quad Edition 930}\end{center}
\part{Dirichlet zeros and the arithmetic of power intervals}\label{part:I}
\begin{abstract}\noindent The zero criterion in this volume uses a prime sum on power intervals. We obtain its explicit formula, choose the two quadratic twists and specify the two arithmetic estimates on common number fields. We construct rational flag coordinates with integral unmarked root ratios and prescribed auxiliary characters. All their unmarked divisors can be placed above the cohomological cutoff at completely split primes. The exact quadratic-sign decomposition of the centered count at the selected word weights compares its signal part with the signed definite Gram determinant. On the polynomial-conductor intervals an explicit positive contribution from the ramified prime must be retained in the centered identity. Its size and the surviving cohomology orbits are given on the same field. On the common period fields, a single set of weights of density tending to zero contains the support of the truncated local groups at all odd primes up to the cutoff. The estimate holds uniformly for every weight limit on the interval.\end{abstract}
\setcounter{section}{0}

\section{Dirichlet GRH and its proof status}\label{nav:I-00}\label{nav:volume-I}
\label{sec:statement}\label{sec:finish}\label{sec:sum}
\begin{maintheorem}[Dirichlet GRH, proof incomplete]\label{thm:grh16}
For every Dirichlet character \(\chi\) of positive modulus, every zero of \(L(s,\chi)\) in the open critical strip lies on the critical line. Equivalently,
\begin{equation*}
 L(\rho,\chi)=0,\qquad 0<\real\rho<1
                    \quad\Longrightarrow\quad \real\rho=\tfrac12. \tag{G}\label{eq:G}
\end{equation*}
\end{maintheorem}
Theorem~\ref{thm:criterion16} reduces Dirichlet GRH to a one-sided bound for an explicit prime sum on arbitrarily distant power intervals. The unconditional argument below reaches two estimates on the same fields that are still needed to prove that bound.

Theorems~\ref{thm:entireinterval310} and~\ref{thm:relativeprogress325} control exceptional twists and certain full relative orders. Theorems~\ref{thm:twolayers327}, \ref{thm:normsum327}, \ref{thm:strictdigits329} and~\ref{thm:alltrunc329} bound the initial cohomology layers. Theorems~\ref{thm:decomploss330} and~\ref{thm:integerlayers333} bound the local Frobenius contribution and the subgroup killed by the common integer of Theorem~\ref{thm:simultaneous307}. Theorem~\ref{thm:ordinarydigits514} recovers the complementary group at any prescribed precision from ordinary ideal classes in one layer for all fields and weights. Corollary~\ref{cor:allprimes514} bounds the loss at the earlier layers uniformly in that precision. Theorem~\ref{thm:finiteweights515} detects exceptional weights using finite ordinary class groups. Corollary~\ref{cor:smallresidual515} gives simultaneous vanishing of the residual groups at small coefficient primes outside one explicitly bounded set on the interval. Theorem~\ref{thm:fullsparse517} bounds the sum of the full strict lengths in both signs outside finite ordinary class group tests. Corollary~\ref{cor:sameprimes517} gives full strict vanishing at all three primes of the signed periods, for every subfield and every weight below the interval cutoff. Theorem~\ref{thm:flagrecovery526} recovers the finite Whittaker coefficients integrally on each prescribed pair of residue flags at the two auxiliary primes. Corollary~\ref{cor:severalflags526} bounds the cross terms when several output flags contribute. Theorems~\ref{thm:whittakerspectrum527} and~\ref{thm:whittakerintegral527} compare the full finite pairing with its projection onto generic representations for arbitrary signed columns. Corollary~\ref{cor:whittakerterminal527} gives equality of their minor ideals at the same terminal primes as the strict local groups. Theorem~\ref{thm:markedapartments528} bounds the finite pairing from below by explicit alternating sums of the signed amplitudes at the prescribed marking. Corollary~\ref{cor:apartmentideals528} gives integral minor comparisons there and equality on complete flag tables. Theorem~\ref{thm:markedstrict337} allows the three marked primes to be chosen with full vanishing on all descended strict groups of the stated degrees and twists. Theorem~\ref{thm:criticalmarks339} makes the critical sign vanish at those primes on the full fields and throughout the weight cutoff. Corollary~\ref{cor:allsample341} bounds full orders outside explicitly bounded sets of twists in all tame characters, including the decomposition invariants. Corollary~\ref{cor:periodlarge425} bounds the entire complementary sum at primes above \(P(x)\) by a polynomial in \(V\). The local estimate still requires a uniform bound on the residual quotients at primes at most \(P(x)\).

For the periods, Theorems~\ref{thm:ramifiedcoset324}, \ref{thm:classtrace326}, \ref{thm:terminaltypes327} and~\ref{thm:typeblocks329} describe the local sums, integral class tests and finite projections. Theorems~\ref{thm:signedfibers331}, \ref{thm:fixedadic332} and~\ref{thm:scalarorbits332} compute signed local coefficients, finite valuations and a conjugate second moment. Theorems~\ref{thm:mixedfourier334} and~\ref{thm:rowcontent335} give the norm of the split transfer and the integral content of the transferred row away from the marked residue characteristics. Theorem~\ref{thm:graphtransport337} evaluates the terminal move at the first Bessel position. Corollary~\ref{cor:rankobstruction337} shows that this position on graph rows annihilates the refined columns with zero sum on each scalar orbit. Theorem~\ref{thm:markedglobal339} gives two explicit finite translates for which the normalized Gram forms agree through every marked exchange on those columns. Corollary~\ref{cor:markedtests341} computes the coefficient ideal and rank bound for the original and transported tests at all the marked places of the same chain. Theorem~\ref{thm:rationalperiod342} gives the exchange on the original signed auxiliary projector by rational conjugation. Theorem~\ref{thm:flagvanish344} proves that the resulting period vanishes on the original spherical rows if all real components are close to the identity and all new valuations are at good split places. Proposition~\ref{prop:dominantexchange344} gives the exchange for a rational diagonal with integral unmarked root ratios. Section~\ref{sec:missing} still requires a comparison of the designated global Bessel pairing with the signed iterated theta column in the same normalization.

For a primitive Dirichlet character \(\psi\), let \(\Gg_\psi(x)\) denote the real weighted prime sum in Definition~\ref{def:family}. Its terms come from the primitive characters inducing the powers of \(\psi\). For the character of conductor one the sole term is the smoothed zeta prime sum.

For a fixed primitive \(\psi\), consider the assertion
\begin{equation*}
 \begin{gathered}
 \text{for every }h>0\text{ and }Y_0\ge2\text{ there exists }Y\ge Y_0
 \text{ for which}\\[-2pt]
              \Gg_\psi(x)\le x^{1/2+h}\qquad(Y\le x\le Y^{3/2}).
 \end{gathered}                                                \tag{W}\label{eq:window}
\end{equation*}

\begin{theorem}[Reduction to a prime-sum estimate]\leavevmode\label{thm:criterion16}
Assertion \eqref{eq:G} holds for every Dirichlet character if and only if \eqref{eq:window} holds for every primitive character \(\psi\). The equivalence holds with \(3/2\) replaced by any fixed \(v>1\).
\end{theorem}

Write \(\real s\) and \(\imag s\) for the real and imaginary parts of \(s\). A Dirichlet character modulo \(N\) is a character of \((\ZZ/N\ZZ)^\times\), extended by zero to nonunits. Its conductor is the least modulus from which it is induced. The character at that modulus is its \emph{primitive representative}. Write \(\mathbf1\) for the character of conductor one. For \(\real s>1\) the Dirichlet series and its Euler product are
\[
 L(s,\chi)=\sum_{n\ge1}\chi(n)n^{-s}
            =\prod_p(1-\chi(p)p^{-s})^{-1}.
\]
We use the classical continuation and functional equation of primitive Dirichlet functions and the Hadamard factorization theorem for entire functions of order one. The classical analytic inputs and their applications are specified in Sections~\ref{sec:boundary} and~\ref{sec:sharp}. Write \(\gamma\) for Euler's constant.

\begin{definition}[The prime sum and its weights]\label{def:ramp}\label{def:family}
The von Mangoldt function is \(\Lambda(p^j)=\log p\) for \(j\ge1\) and vanishes on all other positive integers. Set \(z_+=\max(z,0)\). For a primitive character and \(x\ge1\), put
\begin{equation}
 R_\chi(x)=\sum_{n\ge1}\Lambda(n)\chi(n)(x/n-1)_+,\qquad
 G_\chi(x)=
 \begin{cases}
 R_{\mathbf1}(x)-x\log x+(1+\gamma)x,&\chi=\mathbf1,\\
 R_\chi(x)+x\,L'(1,\chi)/L(1,\chi),&\chi\ne\mathbf1 .
 \end{cases}                                                    \label{eq:Gchi}
\end{equation}
This sum is finite at each \(x\). Lemma~\ref{lem:boundary} gives \(L(1,\chi)\ne0\) in the second line. Fix a primitive character \(\psi\) of order \(o\). Let \(\mathcal C_\psi\) consist of the primitive representatives of \(1,\psi,\ldots,\psi^{o-1}\). Distinct powers give distinct primitive representatives since each is recovered by inducing its representative to the original modulus. The set contains \(\mathbf1\) and \(\psi\) and is closed under conjugation. Define
\[
 a_{\mathbf1}=1,\qquad a_\chi=(3o+1)^{-1}\ (\chi\ne\mathbf1),
 \qquad \Gg_\psi(x)=\sum_{\chi\in\mathcal C_\psi}a_\chi\,\real G_\chi(x).
\]
\end{definition}

The same \(Y\) works for all real \(x\) in the displayed interval. It may depend on the character and the exponent. For example, for the nonprincipal character \(\chi_{-4}\) modulo four the sum in \eqref{eq:window} is \(G_{\mathbf1}+G_{\chi_{-4}}/7\).

\begin{proof}[Proof of Theorem~\ref{thm:criterion16}]
Assume \eqref{eq:window} for all primitive characters with interval power \(v>1\). If a primitive \(\psi\) has a zero at real part \(\beta>1/2\) in the strip, choose \(1/2<\theta<\beta\). Apply \eqref{eq:window} with \(h=\theta-1/2\) and Theorem~\ref{thm:detect} with \(A=1\). Every zero of every member of \(\mathcal C_\psi\) has real part at most \(\theta\), a contradiction.

For the primitive representative \(\chi^*\) of conductor \(f\), take \(a\in\{0,1\}\) with \(\chi^*(-1)=(-1)^a\). The classical functional equation~\cite{I-I-I-nist,I-I-I-elkies} is
\[
 \Lh(s,\chi^*)=(f/\pi)^{(s+a)/2}\Gamma((s+a)/2)L(s,\chi^*)
              =\omega_{\chi^*}\Lh(1-s,\overline{\chi^*}),
 \qquad |\omega_{\chi^*}|=1.
\]
The factors preceding \(L\) are finite and nonzero in the open strip. A zero to the left of the line would give a zero to its right for the conjugate character. For a character \(\chi\bmod N\) induced by \(\chi^*\), comparison of absolutely convergent Euler products gives, by continuation,
\[
 L(s,\chi)=L(s,\chi^*)\prod_{p\mid N,\ p\nmid f}(1-\chi^*(p)p^{-s}).
\]
Every omitted Euler factor has nonzero value in \(\real s>0\), proving \eqref{eq:G} in general.

Conversely, under \eqref{eq:G}, Theorem~\ref{thm:explicit} gives \(|G_\chi(x)|\le C_\chi\sqrt x\) for every primitive \(\chi\) and \(x\ge2\). Its proof uses only \eqref{eq:G} and classical properties of \(L\)-functions. In the finite sum defining \(\Gg_\psi\), take \(C=\sum_\chi a_\chi C_\chi\). Once \(Y^h\ge C\), the bound in \eqref{eq:window} holds for all \(x\ge Y\). This gives the converse for each \(v\).
\end{proof}
\begin{corollary}[Nonvanishing and the zeta zeros]
Assume \eqref{eq:window} for every primitive character. Then \(L(s,\chi)\ne0\) for \(\real s>1/2\), apart from the principal pole at \(1\), and every zero of \(\zeta\) other than a negative even integer has real part \(1/2\).
\end{corollary}
\begin{proof}
Theorem~\ref{thm:criterion16} settles the open strip and Lemma~\ref{lem:boundary} treats \(\real s\ge1\). For \(\real s\le0\) and \(s\ne0\), use
\(\zeta(s)=2^s\pi^{s-1}\sin(\pi s/2)\Gamma(1-s)\zeta(1-s)\).
All factors except the sine are finite and nonzero there. Its zeros in that region are the negative even integers, and \(\zeta(0)=-1/2\).
\end{proof}
 Theorem~\ref{thm:abscissa} identifies the best interval exponent unconditionally. The upper limits in Section~\ref{sect:sharp} and the prime-ideal formula in Section~\ref{sect:abelian} are consequences under the critical-line hypothesis. The proposed approach to \eqref{eq:window} via a chain of theta lifts of length about \(\sqrt x\) comes from Alp\"oge's \emph{Remark}~\cite{I-I-I-alpoge}. Section~\ref{sec:comparison} states the arithmetic estimate that this construction would have to supply. All deductions below use mathematical arguments or the indicated classical theorems. The locations in the accompanying source identify the proposed constructions.

\subsection{A guide to the twenty-seven volumes}\label{nav:I-01}
\label{seriesguide}\MakeLinkTarget*{series-guide}
Sections and their subdivisions are numbered continuously in reading order across the series.
The first number in a theorem or equation identifies its section. Page numbers start at one in each volume.
The contents and the references in the text link to the corresponding headings and statements.
\par\medskip
\begingroup\small\setlength{\tabcolsep}{4pt}\renewcommand{\arraystretch}{1.12}
\begin{longtable}{@{}l r p{.63\linewidth} r@{}}
\textbf{Vol.}&\textbf{Sections}&\textbf{Subject}&\textbf{Pages}\\[2pt]\hline\endfirsthead
\textbf{Vol.}&\textbf{Sections}&\textbf{Subject}&\textbf{Pages}\\[2pt]\hline\endhead
\hyperref[nav:volume-I]{I} & 1--10 & \hyperref[nav:volume-I]{Dirichlet zeros and the arithmetic of power intervals} & 136 \\
\hyperref[nav:volume-II]{II} & 11--14 & \hyperref[nav:volume-II]{Strict local conditions and Selmer groups of twists} & 80 \\
\hyperref[nav:volume-III]{III} & 15--18 & \hyperref[nav:volume-III]{Trace extensions, Frobenius lifts and local cup products} & 122 \\
\hyperref[nav:volume-IV]{IV} & 19--23 & \hyperref[nav:volume-IV]{Cyclotomic descent and relative strict orders} & 111 \\
\hyperref[nav:volume-V]{V} & 24--29 & \hyperref[nav:volume-V]{Class groups, special values and strict orders} & 106 \\
\hyperref[nav:volume-VI]{VI} & 30--35 & \hyperref[nav:volume-VI]{Marked Galois deformations and trace congruences} & 92 \\
\hyperref[nav:volume-VII]{VII} & 36--39 & \hyperref[nav:volume-VII]{Auxiliary fields and ordinary cuspidal families} & 84 \\
\hyperref[nav:volume-VIII]{VIII} & 40--43 & \hyperref[nav:volume-VIII]{Integral trace quotients and Hecke support} & 85 \\
\hyperref[nav:volume-IX]{IX} & 44--49 & \hyperref[nav:volume-IX]{Fixed-weight Hecke support and symplectic deformations} & 127 \\
\hyperref[nav:volume-X]{X} & 50--53 & \hyperref[nav:volume-X]{Projected torus periods and congruence modules} & 84 \\
\hyperref[nav:volume-XI]{XI} & 54--58 & \hyperref[nav:volume-XI]{Strict cohomology at exceptional weights} & 119 \\
\hyperref[nav:volume-XII]{XII} & 59--63 & \hyperref[nav:volume-XII]{Cuspidal companions and ordinary Hecke fibers} & 96 \\
\hyperref[nav:volume-XIII]{XIII} & 64--69 & \hyperref[nav:volume-XIII]{CM theta derivatives and Frobenius translates} & 117 \\
\hyperref[nav:volume-XIV]{XIV} & 70--76 & \hyperref[nav:volume-XIV]{Unitary branching and first occurrence in the theta tower} & 122 \\
\hyperref[nav:volume-XV]{XV} & 77--80 & \hyperref[nav:volume-XV]{Definite theta kernels and the three global exchanges} & 102 \\
\hyperref[nav:volume-XVI]{XVI} & 81--85 & \hyperref[nav:volume-XVI]{Integral theta projections and terminal determinants} & 93 \\
\hyperref[nav:volume-XVII]{XVII} & 86--91 & \hyperref[nav:volume-XVII]{Bessel models and projected periods} & 100 \\
\hyperref[nav:volume-XVIII]{XVIII} & 92--97 & \hyperref[nav:volume-XVIII]{Terminal characters, measures and Whittaker normalizations} & 108 \\
\hyperref[nav:volume-XIX]{XIX} & 98--103 & \hyperref[nav:volume-XIX]{Local orbital integrals and flag lattices} & 98 \\
\hyperref[nav:volume-XX]{XX} & 104--108 & \hyperref[nav:volume-XX]{Real orbital integrals and matrix coefficients} & 92 \\
\hyperref[nav:volume-XXI]{XXI} & 109--114 & \hyperref[nav:volume-XXI]{Rational flag masses and ordinary traces at the marks} & 108 \\
\hyperref[nav:volume-XXII]{XXII} & 115--118 & \hyperref[nav:volume-XXII]{Signed orbital integrals on the marked rows} & 130 \\
\hyperref[nav:volume-XXIII]{XXIII} & 119--124 & \hyperref[nav:volume-XXIII]{Integral coefficient lattices and tests for periods} & 100 \\
\hyperref[nav:volume-XXIV]{XXIV} & 125--129 & \hyperref[nav:volume-XXIV]{Marked theta columns and terminal covectors} & 84 \\
\hyperref[nav:volume-XXV]{XXV} & 130--135 & \hyperref[nav:volume-XXV]{Rational flags and the second moment of Bessel periods} & 99 \\
\hyperref[nav:volume-XXVI]{XXVI} & 136--140 & \hyperref[nav:volume-XXVI]{Local Bessel correlations and relative Gram frequencies} & 142 \\
\hyperref[nav:volume-XXVII]{XXVII} & 141--147 & \hyperref[nav:volume-XXVII]{Heights, Gram operators and divisors of projected theta periods} & 109 \\
\end{longtable}\endgroup
\par\medskip
For the reduction to the prime sum, begin with Theorem~\ref{thm:criterion16}.
Section~\ref{sec:missing} fixes the fields and intervals for both arithmetic estimates.
Theorem~\ref{thm:twolayers327} controls the number of higher-order factors.
Theorem~\ref{thm:rowcontent335} computes the ideal of tested row values away from the marked residue characteristics. Theorem~\ref{thm:integerlayers333} bounds the classes killed by a common integer. Corollary~\ref{cor:midtangent333} compares full negative cotangent lengths with the remaining strict order.
Theorem~\ref{thm:decomploss330} bounds the total change on passing to ordinary unramified class groups at every odd prime and every positive weight below the cutoff. Theorems~\ref{thm:classtrace326} and~\ref{thm:terminaltypes327} give the integral class tests and the finite type projections on the original terminal tables. Theorem~\ref{thm:signedfibers331} evaluates the local signed coefficients. Theorems~\ref{thm:fixedadic332} and~\ref{thm:scalarorbits332} compute the finite valuations and conjugate second moment for a specified algebraic column.
Theorem~\ref{thm:markedstrict337} gives vanishing at the selected period primes in a stated range of descended groups. Theorems~\ref{thm:projectiveperiod337} and~\ref{thm:markedglobal339} give the kernel on the unshifted graph and compare the two Bessel positions at explicit finite translates. The outstanding steps are a uniform bound for the remaining strict orders and a comparison of the signed theta column with the designated Bessel pairing in the same normalization.

\par\medskip
\clearpage
\subsubsection{Sections in the twenty-seven volumes}\label{nav:series-directory}
\begingroup\small\setlength{\parskip}{1pt}\setlength{\columnsep}{1.4em}\raggedcolumns\begin{multicols}{2}
\noindent\begin{minipage}{\linewidth}\raggedright
\textbf{\hyperref[nav:volume-I]{Volume I}}\enspace\emph{Dirichlet zeros and power intervals}\par\smallskip
\noindent\hangindent=2.2em\hangafter=1\hyperref[nav:I-00]{\makebox[2.2em][l]{1}Dirichlet GRH and its proof status}\par
\noindent\hangindent=2.2em\hangafter=1\hyperref[nav:I-02]{\makebox[2.2em][l]{2}A criterion for zeros on power intervals}\par
\noindent\hangindent=2.2em\hangafter=1\hyperref[nav:I-07]{\makebox[2.2em][l]{3}Explicit formulas and sharp exponents}\par
\noindent\hangindent=2.2em\hangafter=1\hyperref[nav:I-13]{\makebox[2.2em][l]{4}Prime ideals in finite abelian extensions}\par
\noindent\hangindent=2.2em\hangafter=1\hyperref[nav:I-16]{\makebox[2.2em][l]{5}Choosing two auxiliary quadratic twists}\par
\noindent\hangindent=2.2em\hangafter=1\hyperref[nav:I-22]{\makebox[2.2em][l]{6}An arithmetic comparison sufficient for the interval bound}\par
\noindent\hangindent=2.2em\hangafter=1\hyperref[nav:g-I-1]{\makebox[2.2em][l]{7}The field and the auxiliary weights}\par
\noindent\hangindent=2.2em\hangafter=1\hyperref[nav:g-I-2]{\makebox[2.2em][l]{8}Arithmetic averages on the period interval}\par
\noindent\hangindent=2.2em\hangafter=1\hyperref[nav:g-I-3]{\makebox[2.2em][l]{9}The finite data on the common fields}\par
\noindent\hangindent=2.2em\hangafter=1\hyperref[nav:bibliography-I]{\makebox[2.2em][l]{10}References}\par
\end{minipage}\par\medskip
\noindent\begin{minipage}{\linewidth}\raggedright
\textbf{\hyperref[nav:volume-II]{Volume II}}\enspace\emph{Local conditions and Selmer groups}\par\smallskip
\noindent\hangindent=2.2em\hangafter=1\hyperref[nav:g-II-0]{\makebox[2.2em][l]{11}Strict local conditions and conjugation}\par
\noindent\hangindent=2.2em\hangafter=1\hyperref[nav:selmer-tests]{\makebox[2.2em][l]{12}Selmer groups and marked local conditions}\par
\noindent\hangindent=2.2em\hangafter=1\hyperref[nav:selmer-critical]{\makebox[2.2em][l]{13}Critical values and Selmer lengths on a common interval}\par
\noindent\hangindent=2.2em\hangafter=1\hyperref[nav:bibliography-II]{\makebox[2.2em][l]{14}References}\par
\end{minipage}\par\medskip
\noindent\begin{minipage}{\linewidth}\raggedright
\textbf{\hyperref[nav:volume-III]{Volume III}}\enspace\emph{Frobenius lifts and local cup products}\par\smallskip
\noindent\hangindent=2.2em\hangafter=1\hyperref[nav:trace-extensions]{\makebox[2.2em][l]{15}Scalar trace extensions at the original level}\par
\noindent\hangindent=2.2em\hangafter=1\hyperref[nav:frobenius-cups]{\makebox[2.2em][l]{16}Cup products along Frobenius orbits}\par
\noindent\hangindent=2.2em\hangafter=1\hyperref[nav:higher-rank-tests]{\makebox[2.2em][l]{17}Higher-rank local tests at spherical and ramified primes}\par
\noindent\hangindent=2.2em\hangafter=1\hyperref[nav:bibliography-III]{\makebox[2.2em][l]{18}References}\par
\end{minipage}\par\medskip
\noindent\begin{minipage}{\linewidth}\raggedright
\textbf{\hyperref[nav:volume-IV]{Volume IV}}\enspace\emph{Cyclotomic descent and norms}\par\smallskip
\noindent\hangindent=2.2em\hangafter=1\hyperref[nav:g-II-1]{\makebox[2.2em][l]{19}Descent in the real cyclotomic tower}\par
\noindent\hangindent=2.2em\hangafter=1\hyperref[nav:g-II-2]{\makebox[2.2em][l]{20}Uniform ranges of coefficient primes}\par
\noindent\hangindent=2.2em\hangafter=1\hyperref[nav:g-II-3]{\makebox[2.2em][l]{21}Residue data and exceptional twists}\par
\noindent\hangindent=2.2em\hangafter=1\hyperref[nav:g-II-4]{\makebox[2.2em][l]{22}Relative orders through the norm map}\par
\noindent\hangindent=2.2em\hangafter=1\hyperref[nav:bibliography-IV]{\makebox[2.2em][l]{23}References}\par
\end{minipage}\par\medskip
\noindent\begin{minipage}{\linewidth}\raggedright
\textbf{\hyperref[nav:volume-V]{Volume V}}\enspace\emph{Class groups and special values}\par\smallskip
\noindent\hangindent=2.2em\hangafter=1\hyperref[nav:g-III-0]{\makebox[2.2em][l]{24}Kummer classes and real class groups}\par
\noindent\hangindent=2.2em\hangafter=1\hyperref[nav:g-III-1]{\makebox[2.2em][l]{25}Special values and integral annihilators}\par
\noindent\hangindent=2.2em\hangafter=1\hyperref[nav:g-III-2]{\makebox[2.2em][l]{26}Norms and cyclic character fields}\par
\noindent\hangindent=2.2em\hangafter=1\hyperref[nav:g-III-3]{\makebox[2.2em][l]{27}Cyclotomic indices at finite primes}\par
\noindent\hangindent=2.2em\hangafter=1\hyperref[nav:g-III-4]{\makebox[2.2em][l]{28}Positive twists and explicit strict orders}\par
\noindent\hangindent=2.2em\hangafter=1\hyperref[nav:bibliography-V]{\makebox[2.2em][l]{29}References}\par
\end{minipage}\par\medskip
\noindent\begin{minipage}{\linewidth}\raggedright
\textbf{\hyperref[nav:volume-VI]{Volume VI}}\enspace\emph{Marked deformations and traces}\par\smallskip
\noindent\hangindent=2.2em\hangafter=1\hyperref[nav:g-IV-0]{\makebox[2.2em][l]{30}Extension classes and trace families}\par
\noindent\hangindent=2.2em\hangafter=1\hyperref[nav:g-IV-1]{\makebox[2.2em][l]{31}Orbits in the marked deformation space}\par
\noindent\hangindent=2.2em\hangafter=1\hyperref[nav:g-IV-2]{\makebox[2.2em][l]{32}The ordinary locus and its tangent spaces}\par
\noindent\hangindent=2.2em\hangafter=1\hyperref[nav:g-IV-3]{\makebox[2.2em][l]{33}Rigidity and orthogonal monodromy}\par
\noindent\hangindent=2.2em\hangafter=1\hyperref[nav:g-IV-4]{\makebox[2.2em][l]{34}Trace congruences on the common interval}\par
\noindent\hangindent=2.2em\hangafter=1\hyperref[nav:bibliography-VI]{\makebox[2.2em][l]{35}References}\par
\end{minipage}\par\medskip
\noindent\begin{minipage}{\linewidth}\raggedright
\textbf{\hyperref[nav:volume-VII]{Volume VII}}\enspace\emph{Auxiliary fields and cuspidal families}\par\smallskip
\noindent\hangindent=2.2em\hangafter=1\hyperref[nav:g-IV-5]{\makebox[2.2em][l]{36}Auxiliary fields and cuspidal congruences}\par
\noindent\hangindent=2.2em\hangafter=1\hyperref[nav:chapter-ordinary-families]{\makebox[2.2em][l]{37}Ordinary curves and cuspidal components}\par
\noindent\hangindent=2.2em\hangafter=1\hyperref[nav:chapter-three-character-traces]{\makebox[2.2em][l]{38}Recovering classes from three-character traces}\par
\noindent\hangindent=2.2em\hangafter=1\hyperref[nav:bibliography-VII]{\makebox[2.2em][l]{39}References}\par
\end{minipage}\par\medskip
\noindent\begin{minipage}{\linewidth}\raggedright
\textbf{\hyperref[nav:volume-VIII]{Volume VIII}}\enspace\emph{Integral trace quotients}\par\smallskip
\noindent\hangindent=2.2em\hangafter=1\hyperref[nav:chapter-integral-trace-quotients]{\makebox[2.2em][l]{40}Integral trace quotients and root modules}\par
\noindent\hangindent=2.2em\hangafter=1\hyperref[nav:chapter-descent-strict-orders]{\makebox[2.2em][l]{41}Strict orders and descent to auxiliary fields}\par
\noindent\hangindent=2.2em\hangafter=1\hyperref[nav:chapter-ordinary-hecke]{\makebox[2.2em][l]{42}Hecke support and ordinary extension classes}\par
\noindent\hangindent=2.2em\hangafter=1\hyperref[nav:bibliography-VIII]{\makebox[2.2em][l]{43}References}\par
\end{minipage}\par\medskip
\noindent\begin{minipage}{\linewidth}\raggedright
\textbf{\hyperref[nav:volume-IX]{Volume IX}}\enspace\emph{Hecke support and deformations}\par\smallskip
\noindent\hangindent=2.2em\hangafter=1\hyperref[nav:chapter-fixed-weight-hecke]{\makebox[2.2em][l]{44}Hecke support at the prescribed weight}\par
\noindent\hangindent=2.2em\hangafter=1\hyperref[nav:hecke-polynomials]{\makebox[2.2em][l]{45}Auxiliary Hecke polynomials and trace determinants}\par
\noindent\hangindent=2.2em\hangafter=1\hyperref[nav:chapter-symplectic-classes]{\makebox[2.2em][l]{46}Symplectic deformations at strict classes}\par
\noindent\hangindent=2.2em\hangafter=1\hyperref[nav:chapter-galois-strict-frobenius]{\makebox[2.2em][l]{47}Galois representations and Frobenius traces}\par
\noindent\hangindent=2.2em\hangafter=1\hyperref[nav:chapter-ordinary-symplectic-equations]{\makebox[2.2em][l]{48}Ordinary symplectic equations and trace multiplicities}\par
\noindent\hangindent=2.2em\hangafter=1\hyperref[nav:bibliography-IX]{\makebox[2.2em][l]{49}References}\par
\end{minipage}\par\medskip
\noindent\begin{minipage}{\linewidth}\raggedright
\textbf{\hyperref[nav:volume-X]{Volume X}}\enspace\emph{Torus periods and congruences}\par\smallskip
\noindent\hangindent=2.2em\hangafter=1\hyperref[nav:torus-constituents]{\makebox[2.2em][l]{50}Local constituents of the projected torus pairing}\par
\noindent\hangindent=2.2em\hangafter=1\hyperref[nav:torus-primes]{\makebox[2.2em][l]{51}Cup bounds and torus periods at the coefficient primes}\par
\noindent\hangindent=2.2em\hangafter=1\hyperref[nav:torus-modules]{\makebox[2.2em][l]{52}Congruence modules and integral projections}\par
\noindent\hangindent=2.2em\hangafter=1\hyperref[nav:bibliography-X]{\makebox[2.2em][l]{53}References}\par
\end{minipage}\par\medskip
\noindent\begin{minipage}{\linewidth}\raggedright
\textbf{\hyperref[nav:volume-XI]{Volume XI}}\enspace\emph{Strict cohomology at exceptional weights}\par\smallskip
\noindent\hangindent=2.2em\hangafter=1\hyperref[nav:chapter-varying-coefficients]{\makebox[2.2em][l]{54}Strict orders at varying coefficient primes}\par
\noindent\hangindent=2.2em\hangafter=1\hyperref[nav:chapter-coincident-characters]{\makebox[2.2em][l]{55}Colliding ordinary characters}\par
\noindent\hangindent=2.2em\hangafter=1\hyperref[nav:chapter-dyadic-exceptional-weights]{\makebox[2.2em][l]{56}Exceptional weights in the cyclotomic tower}\par
\noindent\hangindent=2.2em\hangafter=1\hyperref[nav:chapter-dyadic-units-regulators]{\makebox[2.2em][l]{57}Units, regulators and relative class numbers}\par
\noindent\hangindent=2.2em\hangafter=1\hyperref[nav:bibliography-XI]{\makebox[2.2em][l]{58}References}\par
\end{minipage}\par\medskip
\noindent\begin{minipage}{\linewidth}\raggedright
\textbf{\hyperref[nav:volume-XII]{Volume XII}}\enspace\emph{Cuspidal companions and Hecke fibers}\par\smallskip
\noindent\hangindent=2.2em\hangafter=1\hyperref[nav:chapter-companions-even-planes]{\makebox[2.2em][l]{59}Cuspidal companions and the two even planes}\par
\noindent\hangindent=2.2em\hangafter=1\hyperref[nav:chapter-towers-cuspidal-points]{\makebox[2.2em][l]{60}Auxiliary towers and cuspidal points}\par
\noindent\hangindent=2.2em\hangafter=1\hyperref[nav:chapter-quadratic-induction-symmetric-cubes]{\makebox[2.2em][l]{61}Quadratic induction and symmetric cubes}\par
\noindent\hangindent=2.2em\hangafter=1\hyperref[nav:chapter-ordinary-hecke-fibers]{\makebox[2.2em][l]{62}Hecke fibers and local Selmer conditions}\par
\noindent\hangindent=2.2em\hangafter=1\hyperref[nav:bibliography-XII]{\makebox[2.2em][l]{63}References}\par
\end{minipage}\par\medskip
\noindent\begin{minipage}{\linewidth}\raggedright
\textbf{\hyperref[nav:volume-XIII]{Volume XIII}}\enspace\emph{CM theta derivatives}\par\smallskip
\noindent\hangindent=2.2em\hangafter=1\hyperref[nav:g-V-0]{\makebox[2.2em][l]{64}Complex multiplication and integral sections}\par
\noindent\hangindent=2.2em\hangafter=1\hyperref[nav:g-V-1]{\makebox[2.2em][l]{65}Elliptic derivatives on definite classes}\par
\noindent\hangindent=2.2em\hangafter=1\hyperref[nav:g-V-2]{\makebox[2.2em][l]{66}Frobenius translates and higher derivatives}\par
\noindent\hangindent=2.2em\hangafter=1\hyperref[nav:g-V-3]{\makebox[2.2em][l]{67}Signed theta periods and CM derivatives}\par
\noindent\hangindent=2.2em\hangafter=1\hyperref[nav:chapter-split-elliptic-jets]{\makebox[2.2em][l]{68}Marked elliptic jets and their determinants}\par
\noindent\hangindent=2.2em\hangafter=1\hyperref[nav:bibliography-XIII]{\makebox[2.2em][l]{69}References}\par
\end{minipage}\par\medskip
\noindent\begin{minipage}{\linewidth}\raggedright
\textbf{\hyperref[nav:volume-XIV]{Volume XIV}}\enspace\emph{Branching and first occurrence}\par\smallskip
\noindent\hangindent=2.2em\hangafter=1\hyperref[nav:g-VI-0]{\makebox[2.2em][l]{70}Unitary branching and the global seesaw}\par
\noindent\hangindent=2.2em\hangafter=1\hyperref[nav:g-VI-1]{\makebox[2.2em][l]{71}The first core and the prescribed row}\par
\noindent\hangindent=2.2em\hangafter=1\hyperref[nav:g-VI-2]{\makebox[2.2em][l]{72}Global theta lifting through the chain}\par
\noindent\hangindent=2.2em\hangafter=1\hyperref[nav:first-global-lift]{\makebox[2.2em][l]{73}Terminal Hecke operators and the first global lift}\par
\noindent\hangindent=2.2em\hangafter=1\hyperref[nav:g-VI-3]{\makebox[2.2em][l]{74}Horizontal lifts and first occurrence}\par
\noindent\hangindent=2.2em\hangafter=1\hyperref[nav:g-VI-4]{\makebox[2.2em][l]{75}Norms on the even core}\par
\noindent\hangindent=2.2em\hangafter=1\hyperref[nav:bibliography-XIV]{\makebox[2.2em][l]{76}References}\par
\end{minipage}\par\medskip
\noindent\begin{minipage}{\linewidth}\raggedright
\textbf{\hyperref[nav:volume-XV]{Volume XV}}\enspace\emph{Definite kernels and theta exchanges}\par\smallskip
\noindent\hangindent=2.2em\hangafter=1\hyperref[nav:carrier-integral]{\makebox[2.2em][l]{77}Definite periods and the carrier integral}\par
\noindent\hangindent=2.2em\hangafter=1\hyperref[nav:first-second-exchange]{\makebox[2.2em][l]{78}The first and second theta exchanges}\par
\noindent\hangindent=2.2em\hangafter=1\hyperref[nav:third-exchange]{\makebox[2.2em][l]{79}The third theta exchange}\par
\noindent\hangindent=2.2em\hangafter=1\hyperref[nav:bibliography-XV]{\makebox[2.2em][l]{80}References}\par
\end{minipage}\par\medskip
\noindent\begin{minipage}{\linewidth}\raggedright
\textbf{\hyperref[nav:volume-XVI]{Volume XVI}}\enspace\emph{Theta projections and determinants}\par\smallskip
\noindent\hangindent=2.2em\hangafter=1\hyperref[nav:theta-descent]{\makebox[2.2em][l]{81}Integral projection and descent in the theta tower}\par
\noindent\hangindent=2.2em\hangafter=1\hyperref[nav:critical-torus-tests]{\makebox[2.2em][l]{82}Critical torus periods and their finite tests}\par
\noindent\hangindent=2.2em\hangafter=1\hyperref[nav:g-VIII-4]{\makebox[2.2em][l]{83}Integral positions in the terminal tests}\par
\noindent\hangindent=2.2em\hangafter=1\hyperref[nav:terminal-determinants]{\makebox[2.2em][l]{84}Local divisors and archimedean scales of the projected determinant}\par
\noindent\hangindent=2.2em\hangafter=1\hyperref[nav:bibliography-XVI]{\makebox[2.2em][l]{85}References}\par
\end{minipage}\par\medskip
\noindent\begin{minipage}{\linewidth}\raggedright
\textbf{\hyperref[nav:volume-XVII]{Volume XVII}}\enspace\emph{Projected Bessel periods}\par\smallskip
\noindent\hangindent=2.2em\hangafter=1\hyperref[nav:g-VII-0]{\makebox[2.2em][l]{86}The global Bessel integral}\par
\noindent\hangindent=2.2em\hangafter=1\hyperref[nav:g-VII-1]{\makebox[2.2em][l]{87}Discrete series and real Bessel periods}\par
\noindent\hangindent=2.2em\hangafter=1\hyperref[nav:g-VII-2]{\makebox[2.2em][l]{88}Unfolding and comparison along the row}\par
\noindent\hangindent=2.2em\hangafter=1\hyperref[nav:g-VII-3]{\makebox[2.2em][l]{89}The prescribed initial and terminal periods}\par
\noindent\hangindent=2.2em\hangafter=1\hyperref[nav:g-VII-4]{\makebox[2.2em][l]{90}Restriction, levels and Fourier coefficients}\par
\noindent\hangindent=2.2em\hangafter=1\hyperref[nav:bibliography-XVII]{\makebox[2.2em][l]{91}References}\par
\end{minipage}\par\medskip
\noindent\begin{minipage}{\linewidth}\raggedright
\textbf{\hyperref[nav:volume-XVIII]{Volume XVIII}}\enspace\emph{Terminal characters and measures}\par\smallskip
\noindent\hangindent=2.2em\hangafter=1\hyperref[nav:g-VIII-0]{\makebox[2.2em][l]{92}Weights, terminal lines and finite levels}\par
\noindent\hangindent=2.2em\hangafter=1\hyperref[nav:g-VIII-1]{\makebox[2.2em][l]{93}The common terminal prime and its characters}\par
\noindent\hangindent=2.2em\hangafter=1\hyperref[nav:g-VIII-2]{\makebox[2.2em][l]{94}Measures, formal degrees and integral weights}\par
\noindent\hangindent=2.2em\hangafter=1\hyperref[nav:g-VIII-3]{\makebox[2.2em][l]{95}Compact coefficients and projection norms}\par
\noindent\hangindent=2.2em\hangafter=1\hyperref[nav:terminal-whittaker]{\makebox[2.2em][l]{96}Whittaker normalizations and terminal restriction spaces}\par
\noindent\hangindent=2.2em\hangafter=1\hyperref[nav:bibliography-XVIII]{\makebox[2.2em][l]{97}References}\par
\end{minipage}\par\medskip
\noindent\begin{minipage}{\linewidth}\raggedright
\textbf{\hyperref[nav:volume-XIX]{Volume XIX}}\enspace\emph{Local orbital integrals}\par\smallskip
\noindent\hangindent=2.2em\hangafter=1\hyperref[nav:g-IX-0]{\makebox[2.2em][l]{98}Local measures on Hermitian frames}\par
\noindent\hangindent=2.2em\hangafter=1\hyperref[nav:g-IX-1]{\makebox[2.2em][l]{99}Integral lattices at the marked primes}\par
\noindent\hangindent=2.2em\hangafter=1\hyperref[nav:g-IX-2]{\makebox[2.2em][l]{100}Local factors for a reversed block}\par
\noindent\hangindent=2.2em\hangafter=1\hyperref[nav:g-IX-3]{\makebox[2.2em][l]{101}Split and inert Hecke operators}\par
\noindent\hangindent=2.2em\hangafter=1\hyperref[nav:g-IX-4]{\makebox[2.2em][l]{102}Unramified and transverse factors}\par
\noindent\hangindent=2.2em\hangafter=1\hyperref[nav:bibliography-XIX]{\makebox[2.2em][l]{103}References}\par
\end{minipage}\par\medskip
\noindent\begin{minipage}{\linewidth}\raggedright
\textbf{\hyperref[nav:volume-XX]{Volume XX}}\enspace\emph{Real orbital integrals}\par\smallskip
\noindent\hangindent=2.2em\hangafter=1\hyperref[nav:g-X-0]{\makebox[2.2em][l]{104}Real theta norms and projected coefficients}\par
\noindent\hangindent=2.2em\hangafter=1\hyperref[nav:g-X-1]{\makebox[2.2em][l]{105}Bounds for the first-plane integral}\par
\noindent\hangindent=2.2em\hangafter=1\hyperref[nav:g-X-2]{\makebox[2.2em][l]{106}Matrix coefficients and Hermite moments}\par
\noindent\hangindent=2.2em\hangafter=1\hyperref[nav:g-X-3]{\makebox[2.2em][l]{107}Transverse pairs and higher-rank blocks}\par
\noindent\hangindent=2.2em\hangafter=1\hyperref[nav:bibliography-XX]{\makebox[2.2em][l]{108}References}\par
\end{minipage}\par\medskip
\noindent\begin{minipage}{\linewidth}\raggedright
\textbf{\hyperref[nav:volume-XXI]{Volume XXI}}\enspace\emph{Rational flag masses and traces}\par\smallskip
\noindent\hangindent=2.2em\hangafter=1\hyperref[nav:g-XI-0]{\makebox[2.2em][l]{109}Rational flags and the mass formula}\par
\noindent\hangindent=2.2em\hangafter=1\hyperref[nav:g-XI-1]{\makebox[2.2em][l]{110}Terminal frequencies and stabilizers}\par
\noindent\hangindent=2.2em\hangafter=1\hyperref[nav:g-XI-2]{\makebox[2.2em][l]{111}Rational orbits on the first plane}\par
\noindent\hangindent=2.2em\hangafter=1\hyperref[nav:g-XI-3]{\makebox[2.2em][l]{112}The mean terminal trace}\par
\noindent\hangindent=2.2em\hangafter=1\hyperref[nav:chapter-marked-ordinary-traces]{\makebox[2.2em][l]{113}Ordinary traces at the two marked primes}\par
\noindent\hangindent=2.2em\hangafter=1\hyperref[nav:bibliography-XXI]{\makebox[2.2em][l]{114}References}\par
\end{minipage}\par\medskip
\noindent\begin{minipage}{\linewidth}\raggedright
\textbf{\hyperref[nav:volume-XXII]{Volume XXII}}\enspace\emph{Signed orbital integrals}\par\smallskip
\noindent\hangindent=2.2em\hangafter=1\hyperref[nav:chapter-ordinary-terminal-lattices]{\makebox[2.2em][l]{115}Terminal frequencies and common flag lattices}\par
\noindent\hangindent=2.2em\hangafter=1\hyperref[nav:chapter-transverse-determinants]{\makebox[2.2em][l]{116}Gram determinants and height truncation}\par
\noindent\hangindent=2.2em\hangafter=1\hyperref[nav:chapter-signed-transverse-integrals]{\makebox[2.2em][l]{117}Integral lifts and the signed transverse period}\par
\noindent\hangindent=2.2em\hangafter=1\hyperref[nav:bibliography-XXII]{\makebox[2.2em][l]{118}References}\par
\end{minipage}\par\medskip
\noindent\begin{minipage}{\linewidth}\raggedright
\textbf{\hyperref[nav:volume-XXIII]{Volume XXIII}}\enspace\emph{Coefficient lattices and period tests}\par\smallskip
\noindent\hangindent=2.2em\hangafter=1\hyperref[nav:chapter-refined-rational-columns]{\makebox[2.2em][l]{119}Refined columns and integral class pairings}\par
\noindent\hangindent=2.2em\hangafter=1\hyperref[nav:chapter-coefficient-weyl-lattices]{\makebox[2.2em][l]{120}Branching and elementary divisors of coefficient lattices}\par
\noindent\hangindent=2.2em\hangafter=1\hyperref[nav:g-XII-0]{\makebox[2.2em][l]{121}Arithmetic functionals on the prescribed rows}\par
\noindent\hangindent=2.2em\hangafter=1\hyperref[nav:g-XII-1]{\makebox[2.2em][l]{122}Integral tests, both coefficient lattices and their minors}\par
\noindent\hangindent=2.2em\hangafter=1\hyperref[nav:g-XII-2]{\makebox[2.2em][l]{123}The signed finite tests along the column}\par
\noindent\hangindent=2.2em\hangafter=1\hyperref[nav:bibliography-XXIII]{\makebox[2.2em][l]{124}References}\par
\end{minipage}\par\medskip
\noindent\begin{minipage}{\linewidth}\raggedright
\textbf{\hyperref[nav:volume-XXIV]{Volume XXIV}}\enspace\emph{Marked theta columns and covectors}\par\smallskip
\noindent\hangindent=2.2em\hangafter=1\hyperref[nav:g-XII-3]{\makebox[2.2em][l]{125}Denominators and norms of the endpoint pairing}\par
\noindent\hangindent=2.2em\hangafter=1\hyperref[nav:chapter-signed-terminal-tests]{\makebox[2.2em][l]{126}Signed tests and terminal restrictions}\par
\noindent\hangindent=2.2em\hangafter=1\hyperref[nav:chapter-marked-theta-columns]{\makebox[2.2em][l]{127}Arithmetic theta columns on marked classes}\par
\noindent\hangindent=2.2em\hangafter=1\hyperref[nav:chapter-terminal-covectors]{\makebox[2.2em][l]{128}Terminal covectors and initial class tables}\par
\noindent\hangindent=2.2em\hangafter=1\hyperref[nav:bibliography-XXIV]{\makebox[2.2em][l]{129}References}\par
\end{minipage}\par\medskip
\noindent\begin{minipage}{\linewidth}\raggedright
\textbf{\hyperref[nav:volume-XXV]{Volume XXV}}\enspace\emph{Rational flags and the second moment}\par\smallskip
\noindent\hangindent=2.2em\hangafter=1\hyperref[nav:two-mark-flags]{\makebox[2.2em][l]{130}Rational flags at the two marked primes}\par
\noindent\hangindent=2.2em\hangafter=1\hyperref[nav:marked-open-closed]{\makebox[2.2em][l]{131}Open and closed orbital integrals}\par
\noindent\hangindent=2.2em\hangafter=1\hyperref[nav:stationary-cutoff]{\makebox[2.2em][l]{132}Stationary real coefficients and truncation}\par
\noindent\hangindent=2.2em\hangafter=1\hyperref[nav:poincare-spanning]{\makebox[2.2em][l]{133}Poincar\'e series and their marked second moment}\par
\noindent\hangindent=2.2em\hangafter=1\hyperref[nav:spanning-and-sphere]{\makebox[2.2em][l]{134}Spanning by translates and the terminal norm sphere}\par
\noindent\hangindent=2.2em\hangafter=1\hyperref[nav:bibliography-XXV]{\makebox[2.2em][l]{135}References}\par
\end{minipage}\par\medskip
\noindent\begin{minipage}{\linewidth}\raggedright
\textbf{\hyperref[nav:volume-XXVI]{Volume XXVI}}\enspace\emph{Local Bessel correlations}\par\smallskip
\noindent\hangindent=2.2em\hangafter=1\hyperref[nav:marked-orbital-energy]{\makebox[2.2em][l]{136}Fixed vectors and local position integrals}\par
\noindent\hangindent=2.2em\hangafter=1\hyperref[nav:split-gram-transforms]{\makebox[2.2em][l]{137}The first relative Gram transform at split places}\par
\noindent\hangindent=2.2em\hangafter=1\hyperref[nav:inert-gram-transforms]{\makebox[2.2em][l]{138}The first relative Gram transform at inert places}\par
\noindent\hangindent=2.2em\hangafter=1\hyperref[nav:local-rational-frequency-sum]{\makebox[2.2em][l]{139}Local transforms in the rational frequency sum}\par
\noindent\hangindent=2.2em\hangafter=1\hyperref[nav:bibliography-XXVI]{\makebox[2.2em][l]{140}References}\par
\end{minipage}\par\medskip
\noindent\begin{minipage}{\linewidth}\raggedright
\textbf{\hyperref[nav:volume-XXVII]{Volume XXVII}}\enspace\emph{Heights of projected theta periods}\par\smallskip
\noindent\hangindent=2.2em\hangafter=1\hyperref[nav:chapter-norms-and-period-ideals]{\makebox[2.2em][l]{141}Complex norms and finite period ideals}\par
\noindent\hangindent=2.2em\hangafter=1\hyperref[nav:theta-gram]{\makebox[2.2em][l]{142}Gram operators in the first theta lifts}\par
\noindent\hangindent=2.2em\hangafter=1\hyperref[nav:chapter-marked-period-determinants]{\makebox[2.2em][l]{143}Marked determinants and signed coefficient matrices}\par
\noindent\hangindent=2.2em\hangafter=1\hyperref[nav:chapter-terminal-fourier-heights]{\makebox[2.2em][l]{144}Terminal Fourier components in the period height}\par
\noindent\hangindent=2.2em\hangafter=1\hyperref[nav:chapter-projected-period-divisors]{\makebox[2.2em][l]{145}Common divisors of the projected columns}\par
\noindent\hangindent=2.2em\hangafter=1\hyperref[nav:chapter-signed-terminal-frequencies]{\makebox[2.2em][l]{146}Signed terminal frequencies and their complex norms}\par
\noindent\hangindent=2.2em\hangafter=1\hyperref[nav:bibliography-XXVII]{\makebox[2.2em][l]{147}References}\par
\end{minipage}\par\medskip
\end{multicols}\endgroup\clearpage
\begingroup\small\setlength{\parskip}{0pt}\tableofcontents\endgroup
\clearpage

\section{A criterion for zeros on power intervals}\label{nav:I-02}
\label{sect:analytic}
\begin{theorem}[A one-sided criterion]\label{thm:detect}
Fix a primitive character \(\psi\), \(v>1\), \(0<\theta<1\) and \(A>0\). Suppose there are arbitrarily large \(Y\) such that
\[
              \Gg_\psi(x)\le A x^\theta\qquad(Y\le x\le Y^v).
\]
Then every zero \(\rho\) in the open critical strip of every \(L(s,\chi)\), \(\chi\in\mathcal C_\psi\), satisfies \(\real\rho\le\theta\).
\end{theorem}
The interval can vary with \(Y\). In logarithmic coordinates it has the form \([U,vU]\). We shall choose the order of a derivative so that the integral giving that derivative is concentrated inside this interval.

\subsection{Nonvanishing on the boundary of the strip}\label{nav:I-03}
\label{sec:boundary}
\begin{lemma}[Boundary nonvanishing]\label{lem:boundary}
For every Dirichlet character \(\chi\bmod N\), \(L(s,\chi)\ne0\) on \(\real s\ge1\), with the principal pole at \(s=1\) excluded.
\end{lemma}
\begin{proof}
Absolute convergence of the Euler product settles \(\real s>1\). The partial sums of a nonprincipal character are bounded by periodicity, so partial summation makes \(L(s,\chi)\) holomorphic on \(\real s>0\). For a principal character, the finite Euler factors and the continuation of \(\zeta\) show that the only singularity there is its simple pole at \(1\).

For \(\sigma>1\) and \(t\in\RR\), take logarithms of the Euler products in
\[
       \zeta(\sigma)^3|L(\sigma+it,\chi)|^4
                             |L(\sigma+2it,\chi^2)|\ \ge\ 1 .
\]
To see the inequality, put \(z=\chi(p)^j p^{-ijt}\) at a prime power \(p^j\). Its contribution to the logarithm is
\(p^{-j\sigma}(3+4\real z+\real(z^2))/j\).
If \(|z|=1\), the last parenthesis equals \(2(1+\real z)^2\). If \(p\mid N\), it is \(3\). Every term is nonnegative. Suppose \(L(1+it,\chi)=0\). If \(t\ne0\) or \(\chi^2\ne\chi_0\), the last \(L\)-factor in the display is bounded as \(\sigma\downarrow1\). The fourth power then cancels the third-order pole with a factor \(O(\sigma-1)\) to spare. This contradicts the lower bound.

The remaining case is \(t=0\) and \(\chi\) real and nonprincipal. If \(L(1,\chi)=0\), the function \(D(s)=\zeta(s)L(s,\chi)\) is holomorphic on \(\real s>0\). Its Dirichlet coefficients \(a_n=\sum_{d\mid n}\chi(d)\) are multiplicative and nonnegative. Indeed \(a_{p^j}=1+\chi(p)+\cdots+\chi(p)^j\), and \(\chi(p)\in\{-1,0,1\}\). For even \(j\) this sum is at least one, so \(a_{m^2}\ge1\). Taylor expansion about \(2\) would converge at \(1/2\). Differentiate the absolutely convergent Dirichlet series at \(2\), and interchange the nonnegative terms of its Taylor series. This gives
\[
 D(1/2)
  =\sum_{k\ge0}\sum_{n\ge1}\frac{a_n}{n^2}
                         \frac{(\tfrac32\log n)^k}{k!}
  =\sum_{n\ge1}\frac{a_n}{\sqrt n}
  \ \ge\ \sum_{m\ge1}\frac1m .
\]
The value on the left is finite. This contradiction proves the lemma.
\end{proof}

\subsection{The Mellin transform}\label{nav:I-04}
\label{sec:mellin}
Put \(g(u)=\Gg_\psi(e^u)\) for \(u\ge0\) and define
\[
 \Phi_\psi(s)=\int_1^\infty\Gg_\psi(x)x^{-s-1}\,dx
             =\int_0^\infty g(u)e^{-su}\,du\qquad(\real s>1).
\]
Indeed \(|R_\chi(x)|\le x\sum_{n\le x}(\log n)/n\ll x(1+\log^2x)\). In particular \(g\) is real and \(|g(u)|\le C_\psi e^{2u}\).

\begin{lemma}[Residues of the Mellin transform]\label{lem:mellin}
The function \(\Phi_\psi\) extends meromorphically to \(\real s>0\) and is holomorphic at \(1\). Its poles in that half-plane are exactly the zeros of the family. They are simple, even when an \(L\)-function has a multiple zero. The residue at a real pole is positive.
\end{lemma}
\begin{proof}
For \(\real s>1\),
\(\int_n^\infty(x/n-1)x^{-s-1}\,dx=n^{-s}/[s(s-1)]\).
Absolute convergence and the logarithmic derivative of the Euler product show that the Mellin transform of \(R_\chi\) is \(- (L'/L)(s,\chi)/[s(s-1)]\). Denote the transform of \(G_\chi\) by \(H_\chi\). For \(\chi\ne\mathbf1\) add \((L'/L)(1,\chi)/(s-1)\) to this expression. For \(\chi=\mathbf1\) add \(-1/(s-1)^2+(1+\gamma)/(s-1)\). The first correction cancels the simple pole at \(1\). For the second, substitute \(\zeta(s)=(s-1)^{-1}+\gamma+O(s-1)\) to see both singular coefficients cancel.

We have \(\overline{G_\chi(x)}=G_{\bar\chi}(x)\). Thus
\[
 \Phi_\psi(s)=\tfrac12\sum_{\chi\in\mathcal C_\psi}
                       a_\chi\big(H_\chi(s)+H_{\bar\chi}(s)\big).
\]
The family is closed under conjugation. For any one of its zeros \(\rho\) with \(\real\rho>0\), let \(m_\chi(\rho)\) be its multiplicity for \(L(s,\chi)\), with value zero if the function does not vanish there. Then
\begin{equation}
 \Res_{s=\rho}\Phi_\psi(s)=-\frac{1}{\rho(\rho-1)}
     \,\frac12\sum_{\chi\in\mathcal C_\psi}
          a_\chi\big(m_\chi(\rho)+m_{\bar\chi}(\rho)\big). \label{eq:res}
\end{equation}
The weights are strictly positive and at least one multiplicity is positive. Hence the pole survives, including when several characters vanish at the same point. The formula above exhausts the possible poles in \(\real s>0\). Lemma~\ref{lem:boundary} puts every real pole in \((0,1)\), where \(-1/[\rho(\rho-1)]>0\).
\end{proof}

\subsection{High derivatives on a vertical line}\label{nav:I-05}
\label{sec:landau}
For \(\sigma>2\) and \(k\in\ZZ_{\ge0}\), set
\[
 A_k(s)=\frac{(-1)^k}{k!}\Phi_\psi^{(k)}(s),\qquad
 K_k(u)=\frac{u^k e^{-\sigma u}}{k!}.
\]
\begin{lemma}[A derivative estimate]\label{lem:kernel}
Under the interval hypothesis of Theorem~\ref{thm:detect}, there exist a real \(\sigma>2\), a constant \(C_*>0\), and a sequence of integers \(k\to\infty\) such that, for every \(\tau\in\RR\),
\begin{equation}
 |A_k(\sigma+i\tau)|+\real A_k(\sigma)\le C_*(\sigma-\theta)^{-k}.
 \label{eq:der}
\end{equation}
\end{lemma}
\begin{proof}
Differentiation under the Laplace integral gives \(A_k(\sigma+i\tau)=\int_0^\infty g(u)K_k(u)e^{-i\tau u}\,du\). Write \(g_+=\max(g,0)\). The triangle inequality, \(K_k\ge0\), and \(|g|+g=2g_+\) yield
\begin{equation}
 |A_k(\sigma+i\tau)|+\real A_k(\sigma)
               \le 2\int_0^\infty g_+(u)K_k(u)\,du. \label{eq:plus}
\end{equation}
Write \(B=2\). On a selected interval \(U=\log Y\le u\le vU\), we have \(g_+(u)\le A e^{\theta u}\). Its contribution is at most \(A(\sigma-\theta)^{-k-1}\). Off the interval use \(g_+(u)\le C_\psi e^{Bu}\).

Choose \(c=(1+v)/2\) and \(k=\lceil c\sigma U\rceil\). For \(\sigma\) large the numbers \(q_-=B+(c-1)\sigma\) and \(q_+=(1-c/v)\sigma-B\) are positive. Bound the left tail by multiplying its integrand by \(e^{q_-(U-u)}\), and the right tail by \(e^{q_+(u-vU)}\). On the respective tails these factors are at least one. Extending each majorant to \(0\le u<\infty\) and using \(\int_0^\infty u^k e^{-au}\,du/k!=a^{-k-1}\) gives the two bounds
\[
 C_\psi e^{q_-U}(c\sigma)^{-k-1},\qquad
 C_\psi e^{-q_+vU}(c\sigma/v)^{-k-1}.
\]
After multiplication by \(\sigma^k\), their logarithms have the form \(b_-U+O(1)\) and \(b_+U+O(1)\), where \(b_-\) and \(b_+\) are, respectively,
\[
 B+\sigma(c-1-c\log c),\qquad vB+\sigma(c-v+c\log(v/c)).
\]
Both coefficients of \(\sigma\) in these expressions are negative. The first inequality follows from \(\log c>1-1/c\) for \(c>1\), and the second from \(\log(v/c)<v/c-1\). Fix \(\sigma\) large enough that both displayed quantities are negative. Each tail is then \(O(\sigma^{-k})\) uniformly as \(U\to\infty\). Since \(0<\sigma-\theta<\sigma\), \eqref{eq:plus} gives \eqref{eq:der} along a sequence of the selected intervals. The constant is independent of \(\tau\).
\end{proof}

\subsection{Isolating a pole}\label{nav:I-06}
\label{sec:near}
\begin{proof}[Proof of Theorem~\ref{thm:detect}]
Let \(Z\) be the pole set from Lemma~\ref{lem:mellin}. Suppose \(\rho_0\in Z\) has \(\real\rho_0>\theta\), and let \(\sigma,k\) be as in Lemma~\ref{lem:kernel}. Set \(D_*=\sigma-\theta\) and
\(D=\min\{D_*,\operatorname{dist}(\sigma,Z)\}\).
The disk \(B(\sigma,D)\) is pole-free. Whenever a closed disk of radius \(r\) centered at \(z\) contains just one simple pole \(\rho\) in its interior and none on its boundary, Cauchy's estimate after subtracting its polar part gives
\begin{equation}
 A_k(z)=\frac{\Res_{s=\rho}\Phi_\psi(s)}{(z-\rho)^{k+1}}+O(r^{-k}).
 \label{eq:cauchy}
\end{equation}
For a pole-free closed disk only the error term occurs. All disks used below are contained in \(\real s>0\).

Suppose first that \(\sigma-\real\rho_1<D\) for some pole \(\rho_1\). Choose a radius strictly between these two numbers. At every height \(\tau\) close enough to \(\imag\rho_1\), the disk of that radius about \(z=\sigma+i\tau\) contains \(\rho_1\). The union of the closed disks over a short closed interval of heights is compact, so it contains finitely many poles. Two with different imaginary parts are equidistant from \(z\) at at most one height. Two distinct poles of the same height are never equidistant, since both are to the left of \(\real s=\sigma\). Avoiding the finitely many ties gives a unique nearest pole \(\rho\) to \(z\), at a distance \(d<D\). Choose \(d<r<D\) small enough that \(\rho\) is the only pole in the closed disk about \(z\). The closed disk of radius \(r\) about \(\sigma\) is pole-free. Using \eqref{eq:cauchy} at both centers in \eqref{eq:der} gives
\[
 |\Res_{s=\rho}\Phi_\psi(s)|\,d^{-k-1}-O(r^{-k})
                       \le C_*D_*^{-k}.
\]
Multiply by \(d^k\) and take the limit along the sequence. The residue would vanish, contradicting \eqref{eq:res}.

It remains to consider the case \(\sigma-\real\rho\ge D\) for every \(\rho\in Z\). The putative pole \(\rho_0\) gives \(D\le\sigma-\real\rho_0<D_*\). Thus \(D=\operatorname{dist}(\sigma,Z)\), and the distance is attained, by local finiteness in any closed disk of radius between \(D\) and \(D_*\). Euclidean distance is at least horizontal distance. Equality forces the nearest pole to be the unique real point \(\beta=\sigma-D\). Choose \(D<r<D_*\) so that it is the only pole in the closed disk of radius \(r\). Its residue \(c_\beta\) is positive by Lemma~\ref{lem:mellin}. Equations~\eqref{eq:cauchy} and \eqref{eq:der} at \(\tau=0\) give
\(c_\beta D^{-k-1}-O(r^{-k})\le C_*D_*^{-k}\).
Multiplying by \(D^k\) and taking the same limit gives \(c_\beta/D\le0\). This last contradiction proves the theorem.
\end{proof}

\section{Explicit formulas and sharp exponents}\label{nav:I-07}
\label{sect:sharp}

\subsection{The explicit formula on both sides of the critical line}\label{nav:I-08}
\label{sec:sharp}
We obtain an explicit formula which applies to any fixed primitive character. Its coefficients are absolutely summable. This permits passage to upper limits and exhibits the role of the rightmost zeros.

\begin{lemma}[Local zero count and logarithmic derivative]\label{lem:logder}
Fix a primitive character \(\chi\), including the character of conductor one. Count nontrivial zeros \(\rho=\beta+i\gamma_\rho\) with multiplicity. For \(|t|\ge2\),
\[
 \#\{\rho:|\gamma_\rho-t|\le1\}\ll_\chi\log(2+|t|).
\]
For \(s=\sigma+it\) with \(1/4\le\sigma\le2\), \(|t|\ge2\) and \(L(s,\chi)\ne0\), one has
\[
 \frac{L'}{L}(s,\chi)=
       \sum_{|\gamma_\rho-t|\le1}\frac1{s-\rho}
                           +O_\chi(\log(2+|t|)).
\]
Both bounds hold prior to imposing any critical-line conclusion.
\end{lemma}
\begin{proof}
Use the classical entire completion \(F_\chi(s)=\Lh(s,\chi)\) for a nonprincipal primitive character and \(F_{\mathbf1}(s)=\tfrac12s(s-1)\Lh(s,\mathbf1)\). These are entire functions of order one. By the functional equation and Lemma~\ref{lem:boundary}, their zeros are precisely the nontrivial zeros with \(0<\real\rho<1\), and \(F_\chi(0)F_\chi(1)\ne0\). Hadamard factorization gives
\[
 \frac{F_\chi'}{F_\chi}(s)
   =b_\chi+\sum_\rho\left(\frac1{s-\rho}+\frac1\rho\right),
             \qquad \sum_\rho|\rho|^{-2}<\infty .
\]
Here \(b_\chi\) is a constant and the series converges locally uniformly off the zeros. The stated continuation, order and functional equation follow from the Mellin integral of the theta series and Poisson summation~\cite{I-I-I-nist,I-I-I-elkies}.

At \(s_0=2+it\), the Euler product bounds \(L'/L\) uniformly, and Stirling's formula bounds the logarithmic derivative of the gamma factor by \(O_\chi(\log(2+|t|))\). The extra logarithmic derivatives of \(s\) and \(s-1\) for \(\zeta\) are bounded. Since \(\sum_\rho |\real(1/\rho)|\le\sum_\rho|\rho|^{-2}\), taking real parts gives
\[
     \sum_\rho \frac{2-\beta}{(2-\beta)^2+(t-\gamma_\rho)^2}
                         \ll_\chi\log(2+|t|).
\]
For \(|t-\gamma_\rho|\le1\) each summand is at least \(1/5\). This proves the count. Enlarging the constant on a compact set of heights makes the count valid for all real \(t\).

Subtract the Hadamard logarithmic derivatives at \(s\) and \(s_0\). For the distant zeros, each difference is bounded by \(2|t-\gamma_\rho|^{-2}\). Group the ordinates by \(j<|t-\gamma_\rho|\le j+1\), \(j\ge1\). The first estimate gives a total
\[
       \ll_\chi \sum_{j\ge1}\frac{\log(3+|t|+j)}{j^2}
                   \ll_\chi\log(2+|t|).
\]
For the nearby zeros the terms \(1/(s_0-\rho)\) are bounded individually by one. Removing the gamma and, for \(\zeta\), the rational factors again costs \(O_\chi(\log(2+|t|))\), uniformly on the stated strip. This proves the second estimate.
\end{proof}

\begin{proposition}[A formula without a critical-line hypothesis]\label{prop:allzeros}
For a fixed primitive \(\chi\), let \(\lambda=\ord_{s=0} L(s,\chi)\) and \(d_0=\lim_{s\to0}\{L'(s,\chi)/L(s,\chi)-\lambda/s\}\). For \(x\ge2\), with every nontrivial zero repeated according to multiplicity,
\begin{equation}
 G_\chi(x)=\sum_\rho\frac{x^\rho}{\rho(1-\rho)}
             +\lambda\log x+d_0+\lambda+O_\chi(x^{-1/2}).
                                                               \label{eq:uncond12}
\end{equation}
In particular \(G_\chi(x)=o_\chi(x)\).
\end{proposition}
\begin{proof}
The inverse Mellin kernel for \((v-1)_+\) is \(1/[s(s-1)]\) on \(\real s=2\). This follows by integrating the ramp as in Lemma~\ref{lem:mellin}, or by closing the inverse contour and taking its residues at zero and one. At \(v=1\) continuity gives zero. Absolute convergence of the prime series and integrability of the kernel now give
\[
 R_\chi(x)=\frac{1}{2\pi i}\int_{\real s=2}
                  -\frac{L'}{L}(s,\chi)\frac{x^s}{s(s-1)}\,ds .
\]
Move the line to \(\real s=-1/2\). On this line \(L'/L\ll_\chi\log(2+|\imag s|)\). Indeed, logarithmically differentiating the functional equation expresses it through \(L'/L\) on \(\real s=3/2\) and the two gamma logarithmic derivatives, whose arguments avoid the gamma poles. Lemma~\ref{lem:logder} allows us to choose heights in each of \([n,n+1]\) and \([-n-1,-n]\) at distance at least \(c_\chi/\log n\) from all zero ordinates by deleting short intervals around the \(O_\chi(\log n)\) zeros within distance one of these intervals. At these heights \(L'/L=O_\chi(\log^2n)\) uniformly for \(-1/2\le\real s\le2\). The subtraction in the proof of Lemma~\ref{lem:logder} works on this wider horizontal segment, where the distant-zero differences are \(O(|t-\imag\rho|^{-2})\) and the gamma terms satisfy Stirling's bound. Consequently the two horizontal integrals tend to zero for each fixed \(x\). The new vertical integral is \(O_\chi(x^{-1/2})\) for all \(x\ge2\).

The residue at one is \(-xD_\chi\) for nonprincipal \(\chi\), with \(D_\chi=L'(1,\chi)/L(1,\chi)\). For the principal character the Laurent expansion at the simple pole of \(\zeta\) gives \(x\log x-(1+\gamma)x\). Each is precisely the term removed from \(R_\chi\) in \eqref{eq:Gchi}. At zero the factor preceding \(x^s\) is \(\lambda/s^2+(d_0+\lambda)/s+O(1)\), giving the two elementary terms in \eqref{eq:uncond12}. The functional equation and boundary nonvanishing exclude other zeros with \(-1/2\le\real s\le0\). Every zero in the open strip contributes \(x^\rho/[\rho(1-\rho)]\) per unit of multiplicity. The local zero count implies \(\sum_\rho|\rho(1-\rho)|^{-1}<\infty\). This justifies the limit along the chosen contours. After division by \(x\), each zero term tends to zero since \(\real\rho<1\), and its absolute value is bounded by this same summable sequence. Dominated convergence proves the last assertion.
\end{proof}

\begin{theorem}[Explicit formula for a primitive character]\label{thm:explicit}
Assume all nontrivial zeros of the fixed primitive \(\chi\) lie on the critical line. The choice \(\chi=\mathbf1\) is allowed. List its nontrivial zeros with multiplicity as \(\rho=1/2+it_\rho\). For every real \(x\ge2\),
\begin{equation}
 G_\chi(x)=\sqrt{x}\sum_{\rho}
       \frac{e^{it_\rho\log x}}{1/4+t_\rho^2}+O_\chi(x^{1/4}).
 \label{eq:sharp}
\end{equation}
After division by \(\sqrt{x}\), the series converges absolutely and uniformly. In particular, \(|G_\chi(x)|\ll_\chi\sqrt{x}\) for each fixed primitive character. 

\end{theorem}
\begin{proof}
Put \(\rho=1/2+it_\rho\) in \eqref{eq:uncond12}. Then \(\rho(1-\rho)=1/4+t_\rho^2\), and the elementary terms are \(O_\chi(1+\log x)\), which is \(O_\chi(x^{1/4})\). The local zero count makes the displayed series absolutely convergent after division by \(\sqrt x\), uniformly for \(x\ge2\). Formula~\eqref{eq:uncond12} also gives its secondary term \(\lambda\log x+d_0+\lambda\) with remainder \(O_\chi(x^{-1/2})\).
\end{proof}

\subsection{The limiting upper value}\label{nav:I-09}
\label{sec:limsup}
\begin{proposition}[The upper limit for a primitive character]\label{prop:lim}
Assume the critical-line condition for the fixed primitive character under consideration. The real part of the smoothed sum has the exact upper limit
\begin{equation}
 \limsup_{x\to\infty}\frac{\real G_\chi(x)}{\sqrt x}
       =\sum_\rho\frac1{1/4+t_\rho^2}.                     \label{eq:lim}
\end{equation}
\end{proposition}
\begin{proof}
The upper bound follows from \eqref{eq:sharp} by taking absolute values term by term. For the reverse bound, fix \(\delta>0\). Keep finitely many zeros so that the discarded weights sum to less than \(\delta\). For the retained ordinates \(t_1,\ldots,t_b\), consider the sequence on a compact torus
\[
       z_n=(e^{i2^nt_1},\ldots,e^{i2^nt_b}).
\]
A subsequence \(z_{n_j}\) converges, so the coordinatewise quotients of consecutive members tend to the identity. Put \(u_j=2^{n_{j+1}}-2^{n_j}\). Then \(u_j\to\infty\) and \(e^{iu_jt_k}\to1\) for all retained ordinates. Substitute \(x_j=e^{u_j}\) in \eqref{eq:sharp}. Each retained term, after normalization, tends to its positive weight. The real part of the tail is at least \(-\delta\), and \(O_\chi(x_j^{-1/4})\) tends to zero. The limsup is therefore at least the total weight minus \(2\delta\). Letting \(\delta\downarrow0\) proves equality.

\end{proof}

\lead{The constant for the trivial character}Assume RH for \(\zeta\). The sum on page~3 of the \emph{Remark} is \(G_{\mathbf1}(x)\). In this case the right side of \eqref{eq:lim} can be evaluated. Write
\[
 \xi(s)=\tfrac12s(s-1)\pi^{-s/2}\Gamma(s/2)\zeta(s),\qquad
 \xi(s)=\xi(1-s).
\]
Take the logarithmic derivative of the Hadamard product and subtract its value at \(0\) from its value at \(1\). The exponential prefactor and the regularizing constant in each summand cancel. Thus
\[
 \sum_\rho\frac{1}{\rho(1-\rho)}
    =\frac{\xi'}{\xi}(1)-\frac{\xi'}{\xi}(0)=2\frac{\xi'}{\xi}(1).
\]
Using \((s-1)\zeta(s)=1+\gamma(s-1)+O((s-1)^2)\) and \(\Gamma'(1/2)/\Gamma(1/2)=-\gamma-2\log2\) gives \(2\xi'(1)/\xi(1)=2+\gamma-\log(4\pi)\). On the critical line, \(\rho(1-\rho)=1/4+t_\rho^2\). We obtain
\[
 \limsup_{x\to\infty}
 \frac{\sum_{n\ge1}\Lambda(n)(x/n-1)_+-x\log x+(1+\gamma)x}{\sqrt x}
       =2+\gamma-\log(4\pi)\simeq0.0461914 .
\]
The constant is strictly positive. It is the sum of the positive weights of all nontrivial zeros in the explicit formula. Hence \(G_{\mathbf1}(x)\) cannot be \(O(x^\alpha)\) for any \(\alpha<1/2\).\enspace\href{https://people.math.harvard.edu/~elkies/M229.15/prod.pdf}{Product formula for the completed zeta function}

\subsection{Evaluation of the constant for a primitive character}\label{nav:I-10}
\label{sec:closedconstant}
Let \(\chi\) be nonprincipal and primitive of conductor \(q\), and let \(a\in\{0,1\}\) be determined by \(\chi(-1)=(-1)^a\). 
\begin{proposition}[The constant in terms of the conductor]\label{prop:const}
Assuming the critical-line condition for \(\chi\), the sum \(S_\chi\) on the right of \eqref{eq:lim} is
\begin{equation}
 S_\chi=\log(q/\pi)+\frac{\Gamma'}{\Gamma}\!\left(\frac{1+a}{2}\right)
                     +2\real\frac{L'}{L}(1,\chi).
 \label{eq:constchi}
\end{equation}
This expression also evaluates \(\limsup_{x\to\infty}|G_\chi(x)|/\sqrt{x}\). Indeed the limsup of the real part already attains the termwise upper bound on the absolute value.

\end{proposition}
\begin{proof}
The completed function
\[
 \Lambda_\chi(s)=(q/\pi)^{(s+a)/2}\Gamma((s+a)/2)L(s,\chi)
\]
is entire of order one. Its zeros are the nontrivial zeros of \(L(s,\chi)\). Its values at 0 and 1 are nonzero, by nonvanishing at 1 and the functional equation. Subtract the logarithmic derivatives of its Hadamard product at these two points. The constant and the genus-one regularizing terms cancel, giving
\[
 \frac{\Lambda_\chi'}{\Lambda_\chi}(1)-
 \frac{\Lambda_\chi'}{\Lambda_\chi}(0)
 =\sum_\rho\left(\frac1{1-\rho}+\frac1{\rho}\right)
 =\sum_\rho\frac1{\rho(1-\rho)}.
\]
The last sum converges absolutely by the local zero count. The functional equation gives
\(\Lambda_\chi'/\Lambda_\chi(0)=-\overline{\Lambda_\chi'/\Lambda_\chi(1)}\).
Taking the logarithmic derivative of the displayed completion at 1 gives \eqref{eq:constchi}. On the critical line \(\rho(1-\rho)=1/4+(\imag\rho)^2\), so this sum is exactly \(S_\chi\).
\end{proof}

Assume the critical-line condition for every member of \(\mathcal C_\psi\). For this family the calculation gives the sharp formula
\begin{equation}
       \limsup_{x\to\infty}\frac{\Gg_\psi(x)}{\sqrt{x}}
            =\sum_{\chi\in\mathcal C_\psi}a_\chi S_\chi>0,
 \qquad S_{\mathbf1}=2+\gamma-\log(4\pi).       \label{eq:weightedlim}
\end{equation}
The upper bound follows from the individual formulas. For the lower bound keep finitely many zeros in each of the finitely many series so the total discarded weight is less than \(\delta\). Apply the compact-torus argument of Section~\ref{sec:limsup} to the union of their ordinates. Along the resulting common sequence \(x_j\) all retained phases tend to 1. The total loss is at most \(2\delta\) in the limit. The same sequence approaches the upper limit for every member of this finite family.

\subsection{Simultaneous recurrence near the upper value}\label{nav:I-11}
\label{sec:peaks}
The exact upper limits of Sections~\ref{sec:limsup} and \ref{sec:closedconstant} can be approached simultaneously and with bounded gaps in the logarithm of \(x\).

\begin{corollary}[Simultaneous near-attainment]\label{cor:peaks}
Let \(\mathcal F\) be a fixed finite nonempty set of primitive Dirichlet characters, all of whose nontrivial zeros lie on the critical line. Write \(S_\chi\) for the constant in \eqref{eq:lim}. For every \(\delta>0\) there exist \(C_\delta>1\) and \(Y_\delta\) such that every \([Y,C_\delta Y]\), \(Y\ge Y_\delta\), contains a single \(x\) satisfying
\[
                   \frac{\real G_\chi(x)}{\sqrt{x}}>S_\chi-\delta
                       \qquad(\chi\in\mathcal F).
\]
In particular, for each fixed \(v>1\),
\[
 \lim_{Y\to\infty}\ \max_{Y\le x\le Y^v}\frac{\Gg_\psi(x)}{\sqrt x}
              =\sum_{\chi\in\mathcal C_\psi}a_\chi S_\chi .
\]
\end{corollary}
\begin{proof}
Use the absolutely convergent series \eqref{eq:sharp}. For each \(\chi\) discard a tail of total weight below \(\delta/8\). Collect all retained ordinates in a finite list \(t_1,\ldots,t_b\). Let \(z(u)=(e^{it_1u},\ldots,e^{it_bu})\) and \(K=\overline{z(\RR)}\). This is a compact subgroup of the torus. There is a symmetric open neighborhood \(V\subset K\) of the identity on which the retained real sums differ from their values at the identity by less than \(\delta/4\), simultaneously for the finite family.

The sets \(z(u)V,\ u\in\RR\), cover \(K\) by density. Choose a finite subcover with parameters \(u_1,\ldots,u_m\), and put \(a=\min u_j,\ b'=\max u_j\). For any \(U\), choose \(j\) such that \(z(U+b')\in z(u_j)V\). Then
\[
                    U\le U+b'-u_j\le U+b'-a,\qquad z(U+b'-u_j)\in V .
\]
Take \(U=\log Y\), \(x=\exp(U+b'-u_j)\), \(C_\delta=\exp(1+b'-a)\). At this \(x\) each tail costs at most twice its total weight in a comparison with \(S_\chi\). By taking \(Y_\delta\) large, all the normalized \(O_\chi(x^{-1/4})\) remainders are less than \(\delta/4\). Adding the three errors proves the first assertion. For the second, use \(\mathcal F=\mathcal C_\psi\) and \(C_\delta Y\le Y^v\) for all sufficiently large \(Y\). The termwise upper bound in \eqref{eq:sharp} is uniform for \(x\ge Y\). Let \(\delta\downarrow0\) in the lower bound and combine it with that upper bound.
\end{proof}
The compact-group argument also covers coincident ordinates and arbitrary rational relations among them.

\subsection{The least admissible interval exponent}\label{nav:I-12}
\label{sec:abscissa}
The interval criterion also determines the endpoint exponent without assuming any location for the zeros. Let \(Z_\psi\) be the union of the nontrivial zeros of the family and \(b=\sup_{\rho\in Z_\psi}\real\rho\). This set contains a zero of \(\zeta\). Indeed, if the entire completion \(\xi\) had no zeros, its order and Hadamard factorization would give \(\xi(s)=e^{A+Bs}\). The symmetry \(\xi(s)=\xi(1-s)\) forces \(B=0\), contrary to its growth on the positive real axis by Stirling's formula. It is symmetric about the critical line by the functional equation, so \(1/2\le b\le1\). For fixed \(v>1\), call \(a\in\RR\) admissible if for some \(C>0\) there are arbitrarily large \(Y\) with \(\Gg_\psi(x)\le Cx^a\) throughout \([Y,Y^v]\).

\begin{theorem}[The optimal exponent and the boundary zeros]\label{thm:abscissa}
The least admissible exponent is \(b\), and \(|\Gg_\psi(x)|\ll_\psi x^b\) for all \(x\ge2\). If \(a<b\) and \(C>0\), every sufficiently distant \([Y,Y^v]\) contains an \(x\) with \(\Gg_\psi(x)>Cx^a\). At the endpoint, \(\Gg_\psi(x)=o(x^b)\) holds if and only if the family has no zero of real part \(b\).
\end{theorem}
\begin{proof}
In \eqref{eq:uncond12} every term \(x^\rho\) has absolute value at most \(x^b\) and the coefficients are absolutely summable. The elementary terms are \(O_\chi(1+\log x)\). Summing over the fixed family proves the global upper bound. If \(a<b\) were admissible, take \(\max(0,a)<\theta<b\). Theorem~\ref{thm:detect} would put every real part at most \(\theta\), contradicting the supremum. This also proves the assertion about every sufficiently distant interval by negating the quantifiers in admissibility.

Divide \eqref{eq:uncond12} by \(x^b\). Dominated convergence makes the terms with \(\real\rho<b\) tend to zero in sum, since \(|x^{\rho-b}|\le1\). Thus the normalized family sum equals, up to \(o(1)\),
\[
 \real\sum_{\chi\in\mathcal C_\psi} a_\chi
          \sum_{\substack{L(\rho,\chi)=0\\\real\rho=b\\0<\real\rho<1}}
               \frac{e^{i\imag\rho\log x}}{\rho(1-\rho)} .
\]
An empty sum gives the asserted decay. If it is nonempty then \(0<b<1\). Every coefficient has positive real part because
\(\real(\rho(1-\rho))=b(1-b)+(\imag\rho)^2>0\). The sum \(c\) of their weighted real parts is strictly positive. Retain finitely many terms so the absolute tail is less than \(c/4\), and use the compact-torus argument in Section~\ref{sec:limsup} to make their phases tend simultaneously to one along \(x\to\infty\). The resulting limsup is at least \(c/2>0\). This proves the endpoint assertion.
\end{proof}
For characters satisfying \eqref{eq:window}, Theorem~\ref{thm:detect} and symmetry give \(b=1/2\). On that line the residues are positive real numbers, and \eqref{eq:weightedlim} evaluates the exact upper limit.

\section{Prime ideals in finite abelian extensions}\label{nav:I-13}
\label{sect:abelian}

\subsection{Factorization into Dirichlet functions}\label{nav:I-14}
\label{sec:abel}
\begin{corollary}[Abelian Dedekind zeta functions]\label{cor:abelian}
Let \(K/\QQ\) be a finite abelian extension. The Dirichlet result implies that every zero of \(\zeta_K(s)\) in \(0<\real s<1\) has real part \(1/2\).

\end{corollary}
\begin{proof}
By the Kronecker--Weber theorem, \(K\) is contained in some \(\QQ(\zeta_M)\). Set \(G=\Gal(K/\QQ)\). Pulling each \(\xi\in\widehat G\) back through the cyclotomic quotient gives a Dirichlet character. Write \(\chi_\xi^*\) for its primitive representative. We will prove the factorization by comparing the Euler factors at each rational prime.
\begin{equation}
                 \zeta_K(s)=\prod_{\xi\in\widehat G}L(s,\chi_\xi^*).
                 \label{eq:dedekind}
\end{equation}

Fix \(p\). Let \(I\subset G\) be the inertia group, \(e=|I|\), \(F\in G/I\) the Frobenius element, and \(f=\ord(F)\). There are \(|G|/(ef)\) prime ideals above \(p\), each with residue field of size \(p^f\). With \(z=p^{-s}\), the Euler factor on the left is \((1-z^f)^{-|G|/(ef)}\). On the right, if \(\xi|_I\ne1\), the character is ramified at \(p\), its primitive conductor is divisible by \(p\), and its Euler factor is \(1\). The remaining characters run through \(\widehat{G/I}\). Their values at \(F\) run through all \(f\)-th roots of unity, each \(|G|/(ef)\) times, since every character of a subgroup of a finite abelian group extends to the group. Hence the local product on the right is
\[
       \left(\prod_{\omega^f=1}(1-\omega z)\right)^{-|G|/(ef)}
                         =(1-z^f)^{-|G|/(ef)} .
\]
This calculation covers ramified primes as well. It proves \eqref{eq:dedekind} on \(\real s>1\), hence everywhere by meromorphic continuation. All factors are holomorphic in the open strip. The multiplicity of a zero of their product is the sum of its multiplicities in the factors. Under \eqref{eq:G}, this proves the asserted location of the zeros. Conversely, by taking cyclotomic fields we obtain every primitive Dirichlet character in some factorization. Thus the assertion for finite abelian extensions also implies Dirichlet GRH.

\end{proof}

For \(K=\QQ(i)\) the factorization is \(\zeta_K(s)=\zeta(s)L(s,\chi_{-4})\). The same two \(L\)-functions appeared in the positively weighted sum of Section~\ref{sec:sum}. The factorization into one-dimensional characters in \eqref{eq:dedekind} relies on the abelian hypothesis. The \emph{Remark} describes an extension to general number fields under the name \emph{Extended RH} as an expectation. The prime-ideal formula below uses the abelian factorization.

\subsection{The sharp prime-ideal error}\label{nav:I-15}
\label{sec:primeideal}
The factorization in Section~\ref{sec:abel} gives an exact upper limit for a prime-ideal counting error. Let \(K/\QQ\) be finite abelian, of signature \((r_1,r_2)\) and absolute discriminant \(D_K\). Its Euler--Kronecker constant is
\[
 \gamma_K=\lim_{s\to1}\left(\frac{\zeta_K'}{\zeta_K}(s)+\frac1{s-1}\right).
\]
For a nonzero prime ideal \(\mathfrak p\) of \(\mathcal O_K\), let \(N\mathfrak p=|\mathcal O_K/\mathfrak p|\), and put
\[
 G_K(x)=\sum_{\mathfrak p}\sum_{m\ge1}\log N\mathfrak p
             \left(\frac{x}{(N\mathfrak p)^m}-1\right)_+
                   -x\log x+(1+\gamma_K)x .
\]
\begin{corollary}[The exact prime-ideal constant]\label{cor:idealpeaks}
Assume the critical-line condition for every primitive character in \eqref{eq:dedekind}. For every fixed \(\nu>1\), the following limits have the same value \(S_K>0\), given by
\[
 \limsup_{x\to\infty}\frac{G_K(x)}{\sqrt x}
       =\lim_{Y\to\infty}\max_{Y\le x\le Y^\nu}\frac{G_K(x)}{\sqrt x}
       =\sum_{\zeta_K(\rho)=0,\ 0<\real\rho<1}\frac1{1/4+(\imag\rho)^2},
\]
where zeros are counted with multiplicity. In field invariants this value is
\[
 S_K=2+2\gamma_K+\log D_K-r_1\big(\gamma+\log(4\pi)\big)
                               -2r_2\big(\gamma+\log(2\pi)\big).
\]
\end{corollary}
\begin{proof}
Write \(\mathcal F=\{\chi_\xi^*:\xi\in\widehat{\Gal(K/\QQ)}\}\). Logarithmically differentiate \eqref{eq:dedekind} on \(\real s>1\). Comparing the absolutely convergent prime-power coefficients gives the uncorrected prime-ideal sum as \(\sum_{\chi\in\mathcal F}R_\chi(x)\). Comparing the constant terms at the pole gives \(\gamma_K=\gamma+\sum_{\chi\ne\mathbf1}L'(1,\chi)/L(1,\chi)\), with the sum over \(\mathcal F\). Thus \(G_K(x)=\sum_{\chi\in\mathcal F}G_\chi(x)\). This sum is real since \(\mathcal F\) is closed under conjugation. Corollary~\ref{cor:peaks} gives simultaneous lower bounds for the real parts on every sufficiently distant \([Y,C_\delta Y]\). The termwise upper bound in \eqref{eq:sharp} is uniform for \(x\ge Y\). Since \(C_\delta Y\le Y^\nu\) eventually, both limits equal \(\sum_{\chi\in\mathcal F} S_\chi\). The factorization adds zero multiplicities in the strip, and the zeta factor gives strict positivity.

For the second formula let \(n=r_1+2r_2\). There are \(r_2\) odd and \(r_1+r_2\) even characters. Indeed if complex conjugation acts nontrivially then \(K\) is totally imaginary and exactly half its characters are odd. The conductor--discriminant formula gives \(\prod_{\chi\in\mathcal F}\operatorname{cond}\chi=D_K\)~\cite[Ch.~V, Thm.~3.27]{I-I-I-milne}. In the sum of \eqref{eq:constchi}, give the principal character its value \(2+\gamma-\log(4\pi)\) from Section~\ref{sec:limsup}. Use \((\Gamma'/\Gamma)(1)=-\gamma\), \((\Gamma'/\Gamma)(1/2)=-\gamma-2\log2\) and the identity for \(\gamma_K\) above. Collecting the terms gives \(2+2\gamma_K+\log D_K-n(\gamma+\log\pi)-2(r_1+r_2)\log2\), which is the stated value.
\end{proof}
For \(K=\QQ(i)\), writing \(\beta(s)=L(s,\chi_{-4})\), this specializes to \(S_K=2-2\log\pi+2\beta'(1)/\beta(1)\). This is the sharp upper limit for the smoothed prime-ideal error of the Gaussian integers.

\section{Choosing two auxiliary quadratic twists}\label{nav:I-16}
\label{sec:twists}

The character average in Section~\ref{sec:comparison} contains the original powers of \(\psi\) and additional characters. We prove a selection theorem for the additional characters of the form \(\xi(\cdot)(\cdot/q)^i(\cdot/d)^j\) with \(i+j>0\). Throughout this section \(F\) is a fixed positive integer, \(K\ge1\), \(0<\eps<1/52\), and \(Y\) tends to infinity. Constants and thresholds may depend on these three fixed parameters. Put
\[
 z=Y^{\eps/5},\quad z'=z(\log Y)^2,\quad
 U_1=e^z,\quad U_2=e^{z'},\quad Q_1=8F U_1,\quad Q_2=8F U_1U_2 .
\]

\begin{theorem}[Simultaneous auxiliary twists]\label{thm:twists}
For all sufficiently large \(Y\) there are primes \(\sqrt{U_1}<q\le U_1\), \(\sqrt{U_2}<d\le U_2\) with
\[
 q\equiv3\pmod8,\qquad d\equiv7\pmod{16},\qquad d\equiv-1\pmod p,
\]
where the last congruence holds for \(p=q\) and for every odd prime \(p\le z\).
For every primitive \(\xi\) of conductor dividing \(8F\) and \(i,j\in\{0,1\}\), \(i+j>0\), the primitive character
\(\alpha=\xi(\cdot)(\cdot/q)^i(\cdot/d)^j\) satisfies, simultaneously for \(Y\le y\le Y^K\),
\begin{equation}
 |G_\alpha(y)|\le y^{1/2+\eps/2},\qquad |D_\alpha|\le 3\log y,
 \quad D_\alpha=L'(1,\alpha)/L(1,\alpha).                         \label{eq:tw11}
\end{equation}
Here an exponent zero omits that factor, including at multiples of its conductor. These functions are zero-free in \(\real s\ge49/50,\ |\imag s|\le Q_1\) for \(j=0\), and in the same strip up to height \(Q_2\) for \(j=1\).
\end{theorem}

\subsection{Moments of a finite prime sum}\label{nav:I-17}
Put \(Z=Y^{400K}\). For coefficients \(|c_n|\le1\), a prime \(\ell\), and \(y\in[Y,Y^K]\), define
\[
 S_{\ell,c}(y)=\sum_{n\le Z}\Lambda(n)c_n\min(1,y/n)(n/\ell).
\]
The coefficients may depend on \(Y\) and on choices already fixed when \(\ell\) is being selected.

\begin{lemma}[Counting large twisted sums]\label{lem:high}
Let \(U_1\le U\le U_2\) and let \(M\) be a multiple of \(8\). Assume each prime divisor of \(M\) at most \(Z\) is at most \(z\). For each unit residue \(a\bmod M\), every \(\tau>\eps/10\), and fixed \(c\) as above, the number of primes \(\sqrt U<\ell\le U,\ \ell\equiv a\bmod M\) such that
\[
       |S_{\ell,c}(y)|>y^{1/2+\tau}/2\quad
                         \text{for some }Y\le y\le Y^K
\]
is \(O_{K,\eps,\tau}(U^{1-\eta})\) with \(\eta=\eta(K,\eps,\tau)>0\). The constant and \(\eta\) are independent of \(M,a,c,U,Y\).
\end{lemma}
\begin{proof}
For every \(t\ge2\) the total mass \(\sum_{p^j\le t,\ j\ge2}\log p\) is \(O(\sqrt t\log^2(2t))\). Indeed \(j\le\log_2t\) and each \(j\) contributes at most \(t^{1/j}\log t\). For \(p^j>y\) split into intervals \((2^k y,2^{k+1}y]\) and use the factor \(y/p^j\le2^{-k}\). The resulting geometric sum shows that all \(j\ge2\) contribute \(O(\sqrt y\log^2(2Z))\) to \(S\). Primes dividing \(M\), including \(2\), contribute at most \(z\log(2z)\) to the \(j=1\) terms. Both bounds together are \(o(y^{1/2+\tau})\), uniformly on the interval.

It remains to bound the sum over odd \(p\le Z\) coprime to \(M\). Fix \(y\). By quadratic reciprocity the sum at a prime \(\ell\equiv a\bmod M\) equals \(B(\ell)\), where for every positive integer \(m\equiv a\bmod M\) we put
\[
 B(m)=\sum_{p\le Z,\ p\nmid M} b_p(m/p),\qquad
 b_p=(-1)^{((p-1)/2)((a-1)/2)}c_p\log p\,\min(1,y/p).
\]
Use an odd representative for \(a\). Since \(4\mid M\), the sign is well-defined, and \(\sqrt U>Z\) for large \(Y\). Write \(V=\sum_p|b_p|^2\). The trivial count of integers up to \(y\) and the bound \(\sum_{n>y}n^{-2}\ll1/y\) give \(V\ll y\log^2(2Z)\). Also \(\sum_p|b_p|\le Z\log(2Z)\).

For each integer \(r\ge1\) we claim
\begin{equation}
 \sum_{\substack{m\le U\\m\equiv a\bmod M}}|B(m)|^{2r}
 \le (U/M+1)(2rV)^r+Z^{6r}                                     \label{eq:mom11}
\end{equation}
for large \(Z\). Expand \(B(m)^r\overline{B(m)}^r\). If the product of the \(2r\) primes in a tuple is a square, equal primes can be paired. Summing the absolute weights for a fixed pairing gives \(V^r\). There are at most \((2r)^r\) pairings, and the absolute product of symbols is at most one. For a nonsquare tuple, take the modulus \(D\) equal to the product of its distinct primes. In \(m=a+Mk\), every full period of \(D\) successive \(k\)'s gives all residues modulo \(D\). The sum on such a period vanishes by the Chinese remainder theorem, because some prime has odd multiplicity. The last incomplete period contributes at most \(D\le Z^{2r}\). Multiplication by \((\sum|b_p|)^{2r}\le Z^{4r}\) proves the claim. The bound also holds when the progression has fewer than \(D\) terms.

Take \(r=\lfloor\log U/(20\log Z)\rfloor\), which tends to infinity. Then \(Z^{6r}\le U^{3/10}\). For \(\tau>\eps/10\) choose \(0<\lambda<2\tau-\eps/5\). Since \(r\ll_{K}Y^{\eps/5}\log Y\), for all sufficiently large \(Y\)
\[
       \frac{128rV}{y^{1+2\tau}}\le Y^{-\lambda}.
\]
By \eqref{eq:mom11}, the number of indices with \(|B(m)|>y^{1/2+\tau}/8\) is at most \((U+1)Y^{-\lambda r}+U^{3/10}\). Since \(r\ge \log U/(40\log Z)\) and \(\log Z=400K\log Y\), the first term is \(O(U^{1-\lambda/(16000K)})\). Decrease the saving exponent if needed. The discarded prime powers and the primes dividing \(M\) cost less than \(y^{1/2+\tau}/8\) for large \(Y\). At a fixed \(y\) the same exceptional count therefore suffices for \(|S_{\ell,c}(y)|>y^{1/2+\tau}/4\).

Test all integers in \([Y,Y^K]\) and its two endpoints. There are \(O(Y^K+1)=U^{o(1)}\) tests. On each intervening segment \(S_{\ell,c}\) is affine, so its norm is bounded by the greater endpoint norm. The ratio of the larger endpoint to any point in the segment is at most \(1+1/Y\), whose power \(1/2+\tau\) is eventually less than two. A union bound, reducing the saving exponent once more, proves the result.
\end{proof}

\subsection{Prime counts and zero-free rectangles}\label{nav:I-18}
We use two classical prime and zero distribution theorems here. Their uniformity matters because the second progression will depend on the first prime.

\begin{lemma}[Size of the progression and the exceptional set]\label{lem:dens11}
If \(M\ge2\), \((a,M)=1\) and \(\log M=o(\log U)\), then the primes \(\sqrt U<\ell\le U\) congruent to \(a\bmod M\) are \(U^{1-o(1)}\) in number, uniformly when the ratio \(\log M/\log U\) tends uniformly to zero. Among all primitive characters of conductor at most \(Q\), at most \(O(Q^{1/3})\) have a zero in \(\real s\ge49/50,\ |\imag s|\le Q\), for \(Q\ge3\).
\end{lemma}
\begin{proof}
For the first bound use Thorner and Zaman's prime number theorem in progressions~\cite[Cor.~1.4 and (1.2), (1.7)]{I-I-I-tz}. For \(U\ge M^{12}\) it gives a main term \(\lambda U/\varphi(M)\) for \(\sum_{p\le U,\ p\equiv a\bmod M}\log p\). The relative error tends to zero when \(\log M=o(\log U)\). The factor is \(1\) in the absence of an exceptional zero. In the other case their lower bounds give \(\lambda\gg\min(1,(1-\beta)\log U)\) and \(1-\beta\gg M^{-1}\), where \(\beta\) denotes the exceptional zero. In either case \(\lambda/\varphi(M)\gg M^{-2}\). Thus the weighted count is at least a constant times \(U/M^2=U^{1-o(1)}\). Primes at most \(\sqrt U\) contribute at most \(\sqrt U\log U\). Each remaining logarithm is at most \(\log U\). This proves the required lower bound, and \(U\) gives an upper bound.

For the second bound let \(N_\chi(\sigma,T)\) count zeros with real part at least \(\sigma\) and ordinate of absolute value at most \(T\). The classical density estimate of Montgomery, at fixed \(1/2<\sigma<1\), gives, for every \(\nu>0\),
\[
  \sum_{f\le Q}\ \sum_{\chi\bmod f}^{*}N_\chi(\sigma,T)
       \ll_{\sigma,\nu}(Q^2T)^{3(1-\sigma)/(2-\sigma)+\nu}
       \qquad(Q,T\ge 2).
\]
The star restricts the sum to primitive characters. One can absorb the logarithmic factor in the published estimate into \(\nu\)~\cite[Thm.~1]{I-I-I-mont}. At \(\sigma=49/50,\ T=Q\), the exponent of \(Q\) before this loss is \(3/17\). Choose \(\nu\) with \(3/17+3\nu<1/3\). Each exceptional character contributes at least one zero.
\end{proof}

\subsection{A smoothed formula for the logarithmic derivative at one}\label{nav:I-19}
Set \(w(v)=2(1-v)_+-(1-2v)_+\) for \(v>0\). It equals one up to \(1/2\), lies in \([0,1]\), and vanishes from \(1\) onwards.

\begin{lemma}[Approximation at one]\label{lem:tail11}
Let \(\alpha\) be nonprincipal and primitive of conductor at most \(Q\) with no zero in the rectangle of Lemma~\ref{lem:dens11}. For \(X\ge e^2,\ Q\ge8\),
\[
 \left|D_\alpha+\sum_{n\le X}\frac{\Lambda(n)\alpha(n)}n w(n/X)\right|
 \ll X^{-1/100}\log(2Q)+\frac{\log(2X)+\log(2Q)}{Q}.
\]
The implied constant is absolute.
\end{lemma}
\begin{proof}
We first give the conductor dependence in the local logarithmic derivative estimate. Let \(F_\alpha\) denote the entire completion used in Lemma~\ref{lem:logder}, and \(f\) its conductor. For its Hadamard constant \(b_\alpha\) one has
\(\real b_\alpha+\sum_\rho\real(1/\rho)=0\).
Indeed the functional equation gives \(F_\alpha(s)=\omega_\alpha\overline{F_\alpha(1-\bar s)}\), so its logarithmic derivative has real part zero at any nonzero value on the critical line. At such a point \(s\), the real parts of \(1/(s-\rho)\) cancel in pairs under \(\rho\mapsto1-\bar\rho\), and their sum converges absolutely. Taking real parts of the Hadamard formula proves the identity. Consequently at \(s_0=2+it\),
\[
 \sum_\rho\real\frac1{s_0-\rho}
   =\real\frac{F_\alpha'}{F_\alpha}(s_0)
                 \ll \log(2f(2+|t|)).
\]
Here the Euler product bounds \(L'/L\) by an absolute constant and the completion contributes \(\frac12\log(f/\pi)\) and a digamma term. The latter is \(O(\log(2+|t|))\) by Stirling's formula at large height and continuity at bounded height. Repeating the subtraction at \(s_0\) in Lemma~\ref{lem:logder} gives the same local formula with error \(O(\log(2f(2+|t|)))\), uniformly for \(99/100\le\real s\le3/2\). The bound holds at small heights too, since the gamma arguments in this strip stay away from poles. Each zero with \(|\imag\rho-t|\le1\) is at horizontal distance at least \(1/100\) for \(|t|\le Q/2\) by the rectangle hypothesis. Thus \(L'/L\ll\log(2Q)\) on this whole smaller rectangle.

The Mellin transform of \(w\) is \(W(s)=(2-2^{-s})/[s(s+1)]\) for \(\real s>0\), by integrating each linear ramp. It has residue \(1\) at zero and decays quadratically on vertical lines away from its poles. With \(c=1/\log X\) and by absolute convergence, Mellin inversion gives the finite sum in the lemma as
\[
 \frac1{2\pi i}\int_{\real s=c}
              -\frac{L'}{L}(1+s,\alpha)X^s W(s)\,ds ,
\]
where the line is traversed upwards. On this line \(|L'/L|\le-\zeta'(1+c)/\zeta(1+c)\ll \log(2X)\), by the pole at one. Truncate at heights \(\pm Q/2\) with cost \(O(\log(2X)/Q)\), then move the segment to \(\real s=-1/100\). The only residue is \(-D_\alpha\) at zero. Since \(X^c=e\), the horizontal sides cost \(O(\log(2Q)/Q^2)\) and the new vertical side costs \(O(X^{-1/100}\log(2Q))\). This proves the claim.
\end{proof}

\subsection{The successive choices}\label{nav:I-20}
\begin{proof}[Proof of Theorem~\ref{thm:twists}]
At the first stage take \(U=U_1\), \(M=8,\ a=3,\ Q=Q_1\) and let the base characters \(\kappa\) be all primitive characters of conductor dividing \(8F\). At the second stage \(q\) is already fixed. Put
\[
 U=U_2,\qquad M=16q\prod_{3\le p\le z}^{*}p,\qquad Q=Q_2 ,
\]
where the star indicates primes. The Chinese remainder theorem specifies one unit class modulo \(M\) by the congruences for \(d\). Since \(q>\sqrt{U_1}>Z\), its factor is distinct from the small primes. Here \(\log M\ll z\log(2z)+\log q=o(\log U_2)\), uniformly in \(q\). Take each previous base character and also its product with \((\cdot/q)\) as a base \(\kappa\). The number of tests at each stage is bounded solely in terms of \(F\).

Let \(\ell\) range over prime candidates at the current stage. It eventually exceeds every base conductor, since at stage two \(8Fq<\sqrt{U_2}\), and at stage one \(8F<\sqrt{U_1}\). The product \(\alpha=\kappa(\cdot)(\cdot/\ell)\) is primitive of conductor \(\operatorname{cond}(\kappa)\ell\). Indeed the coprime local factors form a product on the unit group by the Chinese remainder theorem, and each local conductor is still minimal on its own factor. Thus, for fixed \(\kappa\), different candidate primes produce different primitive characters. Their conductors are at most \(Q\). We have \(Q\le U^{3/2}\) at either stage for large \(Y\), so Lemma~\ref{lem:dens11} excludes at most \(O_F(U^{1/2})\) candidates for the zero conditions.

For each base character apply Lemma~\ref{lem:high} to \(c_n=\kappa(n)w(n/Z)\) with \(\tau=\eps/2\). The required condition on the prime divisors of \(M\) holds at both stages. A further \(O_F(U^{1-\eta})\) primes fail some finite-sum test. Lemma~\ref{lem:dens11} supplies \(U^{1-o(1)}\) candidates in all. For fixed \(F,K,\eps\) the union of both exceptional sets has proportion tending to zero, uniformly in the first-stage choice at the second stage. Choose any remaining prime at each stage.

It remains to pass from the finite tests to \eqref{eq:tw11}. For all \(Y\le y\le Y^K\) one has \(y<Z/2\) eventually, so \(w(n/Z)=1\) whenever \(n\le y\). Splitting the defining sum of \(G_\alpha\) at \(y\) gives the exact identity
\[
 G_\alpha(y)=-S_{\ell,c}(y)
    +y\left(D_\alpha+\sum_{n\le Z}\Lambda(n)\alpha(n)w(n/Z)/n\right).
\]
At both stages \(\log(2Q)\ll_F z(\log Y)^2\) and \(Q\ge e^z\). With \(X=Z\), Lemma~\ref{lem:tail11} makes the last term \(o(1)\), uniformly for \(y\le Y^K\), since \(Y^K Z^{-1/100}=Y^{-3K}\). The finite-sum test already gives \(|S_{\ell,c}(y)|\le y^{1/2+\eps/2}/2\), proving the first bound. With \(X=Y\), the lemma again gives \(o(1)\) since \(\eps/5<1/100\). Finally, putting \(J=\lfloor Y\rfloor\), the identity \(\sum_{n\le Y}\Lambda(n)\lfloor Y/n\rfloor=\log(J!)\) follows by summing \(\sum_{d\mid m}\Lambda(d)=\log m\). Since \(\lfloor Y/n\rfloor\ge Y/(2n)\) for \(n\le Y\),
\[
           \sum_{n\le Y}\Lambda(n)/n\le 2\log(J!)/Y\le2\log Y .
\]
Use \(0\le w\le1\) and \(|\alpha(n)|\le1\) to get \(|D_\alpha|\le2\log Y+o(1)\le3\log y\). This proves both bounds for the same primes and all the required characters.
\end{proof}

\subsection{One pair of primes on an exponential range}\label{nav:I-21}
Put \(z=Y^{\varepsilon/5}\), as in Theorem~\ref{thm:twists}.
We shall choose the pair before the subinterval for the cohomology
estimate is known. This requires bounds for the character sums on
a longer range. The modulus of the second progression still contains
the first prime.

\begin{lemma}[A uniform lower bound for the number of candidates]
\label{lem:many232}
For each \(0<\delta<1\) there are \(a_\delta>0\) and
\(U_\delta\) such that, for \(U\ge U_\delta\), \(2\le M\le
U^{a_\delta}\), and every unit class modulo \(M\), at least
\(U^{1-\delta}\) primes in \((\sqrt U,U]\) lie in that class.
\end{lemma}
\begin{proof}
If this failed, for every positive integer \(j\) there would
be \(U_j\ge e^j\), \(2\le M_j\le U_j^{1/j}\), and a unit class
with fewer than \(U_j^{1-\delta}\) primes. Along this sequence
\(\log M_j/\log U_j\to0\), so Lemma~\ref{lem:dens11} gives
\(U_j^{1-o(1)}\) primes in these very classes. This contradicts
the fixed positive \(\delta\). The proof of that lemma includes
the possible exceptional real zero in the progression estimate.
\end{proof}

\begin{theorem}[Fixed auxiliary twists on a longer interval]
\label{thm:long232}
Fix \(F\ge1\) and \(0<\varepsilon<1/52\). There are constants
\(c=c(\varepsilon)>0\) and \(C=C(\varepsilon)>4\) with the
following property. Set \(U_1=e^z\), and choose either
\[
            U_2=e^{Cz}
       \quad\hbox{or}\quad U_2=\exp\!\big(z(\log Y)^2\big).
\]
For every sufficiently large \(Y\) there are primes
\(\sqrt{U_1}<q\le U_1\), \(\sqrt{U_2}<d_0\le U_2\) with
\(q\equiv3\bmod8\), \(d_0\equiv7\bmod16\), and
\(d_0\equiv-1\) modulo \(q\) and every odd prime at most \(z\).
For the same pair, all the auxiliary characters of
Theorem~\ref{thm:twists} satisfy \eqref{eq:tw11} for
\emph{every} real \(Y\le y\le e^{cz}\). They satisfy the
zero-free rectangle of that theorem with \(Q_1=8F U_1\) or
\(Q_2=8F U_1U_2\), according as the factor of conductor \(d_0\)
is absent or present.
\end{theorem}
\begin{proof}
Write \(\lambda=\varepsilon/3\), \(\rho=\lambda/32000\), and fix
\(0<c<\min(1/32000,\rho/8)\). At each stage let \(U\) be its
search bound. Cover \([Y,e^{cz}]\) by consecutive bands with left
endpoints \(H=Y,Y^2,Y^4,\ldots\) and right endpoints
\(\min(H^2,e^{cz})\), stopping at the last nonempty band.
For each band use the cutoff \(Z=H^{800}\). Its logarithm is at
most \(800cz\), and \(Z<\sqrt{U_1}\).

Fix one unit progression and one sequence of coefficients
\(|c_n|\le1\) throughout a band. Assume that all primes dividing
the modulus \(M\) and at most \(Z\) are at most \(z\), and
\(8\mid M\). Define \(B(m)\) and \(V=\sum_p|b_p|^2\) as in the proof
of Lemma~\ref{lem:high} at a fixed \(y\) in this band. Removing
the prime powers of exponent at least two costs
\(O(\sqrt y\log^2(2Z))\) and the primes dividing \(M\) cost at most
\(z\log(2z)\). Both are \(o(y^{1/2+\varepsilon/2})\), uniformly,
since \(H\ge Y\). Quadratic reciprocity identifies the remaining sum
at a candidate prime with \(B(\ell)\), and
\(V\ll y\log^2(2Z)\). Equation~\eqref{eq:mom11} holds for every
integer moment order. Use \(l=\lfloor\log U/(20\log Z)\rfloor\).
Our choice of \(c\) gives \(l\ge\log U/(40\log Z)\) for both
stages. With \(C\) fixed, for all sufficiently large \(Y\) one has
\[
 l\log^2(2Z)\ll (\log U)\log Z
       \ll_C z^2(\log Y)^2=o(Y^{\varepsilon/2}).
\]
Consequently \(128lV/y^{1+\varepsilon}\le H^{-\lambda}\)
uniformly. Markov's inequality at \(y^{1/2+\varepsilon/2}/8\)
bounds the number of exceptional indices by
\[
  (U+1)H^{-\lambda l}+Z^{6l}
     \ll U^{1-\rho}+U^{3/10}.
\]
Here \(\log Z=800\log H\), which makes the saving independent
of the band and of the fixed constant \(C\).

Test each integer in each band and the two endpoints, using
\(c_n=\kappa(n)w(n/Z)\) for a fixed base character \(\kappa\)
at the stage under consideration. There are at most
\(\exp(2cz)\) tests for large \(Y\), since the bands have disjoint
interiors and their number is \(O(1+\log z)\). The finite sum is
affine between successive tests within a band. Its norm is at
most the larger endpoint norm, and the ratio of adjacent arguments
is at most \(1+1/Y\). Adding back the discarded terms shows that,
apart from \(O(U^{1-\rho/2})\) candidates, the finite sum in each
band is at most \(y^{1/2+\varepsilon/2}/2\) in absolute value.
The constants in the power saving hold before \(C\) is chosen.
The lower threshold for \(Y\) may depend on \(C\).

At the first stage take \(M=8\). Having chosen \(q\), take
\(M=16q\prod_{3\le p\le z}p\), the product being over primes.
The Chinese remainder theorem supplies the required unit class.
The elementary Chebyshev bound \(\sum_{p\le z}\log p\ll z\)
gives \(\log M\le C_0z\) with an absolute \(C_0\). For instance,
the primes in \((n,2n]\) divide \(\binom{2n}{n}\le4^n\), and
summing the resulting bounds over dyadic intervals gives the
claim. Also \(q>Z\) in every band, as required in the moment
argument. Apply Lemma~\ref{lem:many232} with \(\delta=\rho/4\).
Choose \(C>\max(4,C_0/a_\delta)\). Either value of \(U_2\)
then admits at least \(U_2^{1-\delta}\) candidates, uniformly in
the first prime. The same lower bound with \(U_1\) holds at
the first stage.

The base characters are exactly those in the proof of
Theorem~\ref{thm:twists}, and there are \(O_F(1)\) of them at
each stage. The candidate prime exceeds every base conductor
and gives distinct primitive twists for each fixed base.
At both stages \(Q_i\le U_i^{3/2}\) for sufficiently large \(Y\).
Lemma~\ref{lem:dens11} excludes \(O_F(U_i^{1/2})\) candidates
for zeros in the prescribed rectangles. Together with the moment
exceptions this still leaves a prime at each stage.

Finally apply Lemma~\ref{lem:tail11} with \(X=Z\) on each band.
In the identity for the prime sum in the proof of
Theorem~\ref{thm:twists}, the error in replacing \(G_\alpha(y)\)
by the negative of the finite sum is bounded by
\[
 O\!\left(H^{-6}\log(2Q_i)
       +e^{cz}\frac{\log(2Z)+\log(2Q_i)}{Q_i}\right)=o(1).
\]
Indeed \(y\le H^2\), \(Q_i\ge e^z\), and
\(\log(2Q_i)\ll_{F,C} z(\log Y)^2\). The weight is one for
\(n\le y\) since \(y<Z/2\). Applying that same lemma with
\(X=Y\) gives \(|D_\alpha|\le2\log Y+o(1)\) by the elementary
von Mangoldt sum used in Theorem~\ref{thm:twists}, as
\(\varepsilon/5<1/100\). Both parts of \eqref{eq:tw11} follow.
\end{proof}

\section{An arithmetic comparison sufficient for the interval bound}\label{nav:I-22}
\label{sec:comparison}

We now express the arithmetic step as a comparison between an actual count and its main term. This gives a precise place to test the proposed construction, already for the case \(\psi=\mathbf1\).

\subsection{Fourier expansion at a prime}\label{nav:I-23}
Let \(E/\QQ\) be finite abelian with Galois group \(B=(\ZZ/2\ZZ)^2\times A\), written additively. Set \(n=|A|\), let \(\Phi\) be the subgroup of elements with first coordinate zero, and put \(f=|\Phi|=2n\). For \(\xi\in\widehat B\) let \(\xi^*\) be the associated primitive Dirichlet character. This identification is furnished by embedding \(E\) in a cyclotomic field. At \(p\), write \(I_p\) for inertia and \(F_p\) for any lift of Frobenius in \(B/I_p\). For \(j\ge1\) we have
\[
 \xi^*(p^j)=|I_p|^{-1}\sum_{z\in I_p}\xi(jF_p+z).
\]
Indeed, if the restriction to \(I_p\) is nontrivial the average vanishes and \(p\) divides the primitive conductor. Otherwise the character is unramified and equals its value on Frobenius. These assertions also follow by projecting \((\ZZ/M\ZZ)^\times\) to its prime-power factors in a cyclotomic field containing \(E\). Orthogonality on \(B\) now gives, for every \(g\in B\),
\begin{equation}
 \frac1f\sum_{\xi\in\widehat B}\overline{\xi(g)}\,\xi^*(p^j)
   =\frac{2}{|I_p|}\mathbf1_{g-jF_p\in I_p}.                       \label{eq:frob}
\end{equation}

Give each \(\alpha\in\Phi\) a weight \(w_\alpha>0\). For \(t\in A\) put \(p_t(g)=(1,1,t)+g\), a bijection \(B\setminus\Phi\to\Phi\). For \(a>0\) define
\[
 C(t,a,w)=\sum_{g\notin\Phi}\sum_{\substack{p\ {\rm prime}\\ j\ge1}}
       \frac{2\log p}{|I_p|}\mathbf1_{g-jF_p\in I_p}
                         (a/p^j-w_{p_t(g)})_+ .
\]
The cutoff makes the sum finite. Let \(\xi_c\) denote the sign of the first binary coordinate and write \(D_\chi=L'(1,\chi)/L(1,\chi)\) for nonprincipal primitive characters.

\begin{proposition}[The main term of the count]\label{prop:count}
Define \(M(a,w)=a\{\log a-1-\gamma+D_{\xi_c^*}-f^{-1}\sum_{\alpha\in\Phi}\log w_\alpha\}\). Then
\begin{equation}
 C(t,a,w)-M(a,w)=\frac1f\sum_{g\notin\Phi}\sum_{\xi\in\widehat B}
   \overline{\xi(g)}\,w_{p_t(g)}\,G_{\xi^*}(a/w_{p_t(g)}).         \label{eq:count}
\end{equation}
In this formula extend the finite sum and the correction in \eqref{eq:Gchi} to \(x>0\) by the same formulas.
\end{proposition}
\begin{proof}
Insert \eqref{eq:frob} and use \((a/p^j-w)_+=w((a/w)/p^j-1)_+\). This gives the right side with \(R\) in place of \(G\). For the trivial character the difference, averaged over \(g\), is \(a\{\log a-1-\gamma-f^{-1}\sum\log w_\alpha\}\). For a nontrivial \(\xi\) the difference per summand before multiplication by the character is \(-aD_{\xi^*}\), independent of \(w\). The normalized sum of \(\overline{\xi(g)}\) over the nonidentity coset of \(\Phi\) vanishes unless \(\xi=\xi_c\), when it equals \(-1\). This contributes \(aD_{\xi_c^*}\). Conjugate terms pair, so both sides of \eqref{eq:count} are real.
\end{proof}

\subsection{Which characters survive the average}\label{nav:I-24}
Fix the original \(\psi\) of order \(o\). Let \(A_0\subset A\) have cyclic quotient \(H\) of order \(o\). Let \(T\subset A\) and \(T_2\subset A_0\) be nonempty. Write \(h=|T|,\ h_2=|T_2|,\ \delta=h_2/n\). Choose \(a_t>0\) for each \(t\in T\), and positive weights \(w_{uj,\alpha}\) for \(u\in A_0\), \(1\le j\le m\), \(\alpha\in\Phi\). At scale \(x\ge3\), put \(b=\exp(x^\eps)\), \(Q=hmb\), and set
\[
 \Delta_u=\sum_{t\in T}\sum_{j=1}^m
      \{C(t,a_t,w_{uj})-M(a_t,w_{uj})\}.
\]
Here \(w_{uj}\) denotes the whole list of its weights. The distinguished weights occur at \(L_u=\{(0,1,z+u):z\in T_2\}\). Define \(P_{ut}\) as the right side of \eqref{eq:count} restricted to \(p_t(g)\in L_u\), with all distinguished weights replaced by \(b\). Write \(y_t=a_t/b\) and \(e(t)=(1,0,t)\).

\begin{lemma}[Averaging a translate]\label{lem:av10}
For every \(t\) the exact average is
\[
 |A_0|^{-1}\sum_{u\in A_0}P_{ut}
    =\frac{b\delta}{2}\sum_{\xi|_{A_0}=1}\xi(e(t))G_{\xi^*}(y_t).
\]
\end{lemma}
\begin{proof}
For \(p_t(g)=(0,1,z+u)\) one has \(g=(1,0,z+u-t)\). Consequently \(\overline{\xi(g)}=\xi(e(t))\xi(0,0,-z)\xi(0,0,-u)\). Averaging the last factor over \(u\) kills a nontrivial restriction. On the surviving characters the middle factor is also one, since \(T_2\subset A_0\). The sum over \(z\) has \(h_2\) terms, and \(h_2/f=\delta/2\).
\end{proof}

We specify the hypotheses needed at this single \(x\). These are conditions on data which can be checked before applying the next proposition.
\begin{enumerate}
\item\label{cond:sizes} There are \(a_E\ge0\) and \(k>0\) with
\[
 \begin{gathered}
 a_E\le2xb,\quad h\ge\tfrac12n x^{-\eps/8},\quad \delta\ge x^{-\eps},
 \qquad |y_t-x|\le1,\\
 x^{1/2-\eps}\le m\le x^{1/2+\eps/2},\qquad
 xm\le k\le n^2x^{1/2-\eps}.
 \end{gathered}
\]
\item\label{cond:chars} Among the \(4o\) characters trivial on \(A_0\), a subset is in bijection with \(\mathcal C_\psi\) by \(\xi\mapsto\xi^*\). On this subset \(\xi(1,0,0)=1\), and the restrictions to \(H\) run through all its characters, with the trivial restriction paired with \(\mathbf1\). Each other surviving character satisfies \(|G_{\xi^*}(y_t)|\le 2x^{1/2+\eps}\) for every \(t\). The proportions \(\nu(\alpha)=|T\cap(\alpha+A_0)|/h\) satisfy
\[
 \left|\nu(\alpha)-\frac{3+\mathbf1_{\alpha=0}}{3o+1}\right|
                 \le x^{-1+\eps}\qquad (\alpha\in H).
\]
\item\label{cond:weights} For each \(u,j,t\) the difference between \(C(t,a_t,w_{uj})-M(a_t,w_{uj})\) and \(P_{ut}\) has absolute value at most \(2b/x\).
\item\label{cond:gap} For every \(u\in A_0\) the prime-power count satisfies
\begin{align}
  \Delta_u &\le x^{5\eps}\mathcal B,                                      \label{eq:gap}\\
  \mathcal B&=h a_E(1+m^2/k)
        +m\left(2nmbx^\eps+\frac{h_2 k b x^{2\eps}}{n^2}\right).           \notag
\end{align}
\end{enumerate}
For clarity, condition~\ref{cond:weights} can be tested on individual weights. If \(b\le w\le b+\eta\), the derivative of \(wR_\chi(a_t/w)\), away from its finitely many breakpoints, has absolute value at most \(\sum_{n\le a_t/w}\log n\). The nonprincipal correction \(a_tD_\chi\) is constant in \(w\). The principal correction adds a derivative \(a_t/w\). Integration shows that each distinguished term changes by at most \(H_x\eta\), with \(H_x=3(x+1)\log(2x+4)\). It suffices that \(|B|h_2H_x\eta/f\le b/x\) and that the sum of \(|w_{uj,\alpha}G_{\xi^*}(a_t/w_{uj,\alpha})|/f\) over the other weights and all \(\xi\) be at most \(b/x\).

\begin{proposition}[Arithmetic reduction]\label{prop:balance}
Fix \(\psi\) and \(0<\eps<1/52\). If conditions~\ref{cond:sizes}--\ref{cond:gap} hold at \(x\), then for all sufficiently large such \(x\),
\[
                         \Gg_\psi(x)\le x^{1/2+9\eps}.
\]
The threshold depends only on \(\psi,\eps\). In particular, if for every sufficiently small fixed \(\eps>0\) these data exist at every point of arbitrarily distant intervals \([Y,Y^{3/2}]\), then \(\psi\) satisfies \eqref{eq:window}.
\end{proposition}
\begin{proof}
For retained characters, replacing \(y_t\) by \(x\) costs \(O_\psi(1+\log^2x)\) per character. Indeed the ramp functions have slopes bounded in total by \(\sum_{n\le x+1}\log(n)/n\), and the explicit corrections have derivatives \(O_\psi(1+\log x)\). Also \(|G_\chi(x)|\ll_\psi x(1+\log^2x)\). Replacing \(\nu\) by its target therefore costs \(O_\psi(x^\eps(1+\log^2x))\) in the normalized average over \(t\). By orthogonality on \(H\), the target's Fourier coefficients are \(1\) and \(1/(3o+1)\) at the trivial and nontrivial characters respectively. The retained terms now sum to \(\Gg_\psi(x)\) after taking real parts. The discarded terms have total absolute value at most \(8o x^{1/2+\eps}\).

Sum Lemma~\ref{lem:av10} over \(j,t\) and use conditions~\ref{cond:weights} and~\ref{cond:gap}. After division by \(Q\), the resulting inequality is
\[
 \frac{\delta}{2}\big(\Gg_\psi(x)-O_\psi(x^{1/2+\eps})\big)
                       \le x^{5\eps}\mathcal B/Q+2/x .
\]
The four terms of \(\mathcal B/Q\), in order, are at most
\[
  2x^{1/2+\eps},\qquad 2,\qquad
  4x^{1/2+13\eps/8},\qquad 2x^{1/2+9\eps/8}.
\]
For the second use \(a_E\le2xb\) and \(k\ge xm\). For the last use \(h_2\le n\), \(n/h\le2x^{\eps/8}\) and \(k/n^2\le x^{1/2-\eps}\). Using \(\delta^{-1}\le x^\eps\) gives \(O_\psi(x^{1/2+(61/8)\eps})\). Since \(61/8<9\), the constant is absorbed for fixed \(\eps>0\). Finally choose \(\eps<h'/9\) for the desired loss \(h'>0\) in \eqref{eq:window} and take \(Y\) beyond the pointwise threshold.
\end{proof}

\subsection{Constructing the field and the character weights}\label{nav:I-25}
\label{sec:field12}
The next construction realizes the character identifications and the first three conditions in Proposition~\ref{prop:balance}. Its last inequality is a separate requirement on the resulting count.

\begin{proposition}[A field carrying the required characters]\label{prop:field12}
Fix \(\psi\) of conductor \(F\), order \(o\), and \(0<\eps<1/52\). For every sufficiently large \(Y\) there are fixed primes \(q,d\) such that at every \(x\in[2Y,(2Y)^{3/2}]\) an abelian field, translates, and positive weights can be chosen to satisfy conditions~\ref{cond:sizes}--\ref{cond:weights}. For these choices \(\delta=1/o\) and, uniformly in \(u\in A_0\),
\begin{equation}
 \left|\frac{2\Delta_u}{Q\delta}-\Gg_\psi(x)\right|
             \le 3o\,x^{1/2+\eps/2}+\frac{4}{\delta x}.              \label{eq:actual12}
\end{equation}
\end{proposition}
\begin{proof}
Choose \(q,d\) by Theorem~\ref{thm:twists} with \(K=2\). Let \(K_\psi\) be the subfield of \(\QQ(\zeta_F)\) fixed by \(\ker\psi\). Its Galois group is cyclic of order \(o\). Choose the least integer \(l\ge0\) with \(2^l\ge x^{2+8\eps}\) and form the real cyclotomic layer \(C_l=\QQ(\zeta_{2^{l+2}})^+\) of degree \(2^l\). Write \(E_0=K_\psi C_l\), \(A=\Gal(E_0/\QQ)\), \(A_0=\Gal(E_0/K_\psi)\). Then \(A/A_0\simeq\Gal(K_\psi/\QQ)\) and \(2^l\le n=|A|\le o2^l\).

Put \(E=E_0(\sqrt{-q},\sqrt{-d})\). Both primes are \(3\bmod4\), so their quadratic discriminants are \(-q,-d\). The field \(E_0\) is unramified at \(q,d\), since these primes do not divide \(2F\). Each of the three quadratic subfields of \(\QQ(\sqrt{-q},\sqrt{-d})\) ramifies at one or both of \(q,d\). Any nontrivial intersection with \(E_0\) would contain one of these subfields. Hence the compositum has group \(B=(\ZZ/2\ZZ)^2\times A\), with first and second factors cut out by the quadratic characters of conductors \(q,d\). Quadratic reciprocity identifies these characters with \((\cdot/q),(\cdot/d)\), also at \(2\) by its supplementary law.

A character of \(B\) trivial on \(A_0\) is a product of the two quadratic characters, each to exponent zero or one, and a character of \(A/A_0\). The last factor runs through precisely the primitive representatives of powers of \(\psi\). Taking both binary exponents zero gives all \(o\) retained characters and \(\xi(1,0,0)=1\) for them. Every other character has at least one quadratic factor and hence obeys \eqref{eq:tw11}. Here \(Y\le x\le Y^2\) once \(Y\) is large.

Let \(s=|A_0|\), \(v=\lfloor s/4\rfloor\). In the identity coset of \(A_0\) choose \(4v\) elements, and in each other coset choose \(3v\), to make \(T\). These counts fit in their cosets. They realize exactly the proportions in condition~\ref{cond:chars} and give \(h=(3o+1)v\ge n/2\) for large \(s\). Take \(T_2=A_0\), so \(\delta=1/o\). Put \(m=\lfloor\sqrt{x}\rfloor,\ k=xm,\ a_E=xb,\ a_t=xb\) for all \(t\). Since \(n\ge x^{2+8\eps}\), the upper bound \(k\le n^2x^{1/2-\eps}\) follows. The remaining size bounds follow from \(m/\sqrt{x}\to1\), \(h\ge n/2\) and fixed \(o\).

Give the distinguished positions weight exactly \(b\) and all other positions the same positive weight \(w\). At each fixed \(x\) the field and the list of characters are finite. Proposition~\ref{prop:allzeros} gives \(wG_{\xi^*}(a_t/w)=a_t\,o_{\xi^*}(1)\) as \(w\downarrow0\). Thus one can choose \(w\) so that the sum of the absolute values of all the omitted terms in \eqref{eq:count}, including the factor \(1/f\), is at most \(b/x\) per \(u,j,t\). This is enough for condition~\ref{cond:weights}. This qualitative choice uses no uniform estimate as the conductor or \(x\) grows.

Finally \(L_u\) is independent of \(u\) because \(T_2=A_0\). The idealized sum for each \(u\) therefore equals its average in Lemma~\ref{lem:av10}. Here \(y_t=x\) exactly. The Fourier coefficients of \(T\) are the exact prescribed ones on the retained characters. All the other \(3o\) coefficients have modulus at most one after division by \(h\). Applying \eqref{eq:tw11}, and allowing the full error of condition~\ref{cond:weights}, gives \eqref{eq:actual12}.
\end{proof}
The weights in this proposition are freely chosen real numbers. In the theta-lift construction the weights and translates have geometric constraints. Formula~\eqref{eq:actual12} explains what remains even with the freedom used here. Controlling the twists leaves the entire retained sum in the normalized count.

\subsection{The local restriction kernels in the cyclotomic tower}\label{nav:I-26}
\label{sec:coho12}
We make explicit one part of the divisibility problem for the fields above. For a finite number field \(E\), let \(G=G_{E,S_\ell}\) be the Galois group of its maximal extension unramified at finite places outside \(S_\ell=\{v:v\mid\ell\}\). Write \(\kappa\) for the \(\ell\)-adic cyclotomic character. Give \(M_r=\QQ_\ell/\ZZ_\ell\) the discrete topology and let \(G\) act by \(\kappa^{-r}\), for an integer \(r\ge1\). The strict local restriction kernel is
\[
 \mathcal S_{\ell,r}(E)=
 \ker\left(H^1(G,M_r)\longrightarrow
                          \prod_{v\mid\ell}H^1(E_v,M_r)\right).
\]
All cohomology and homomorphisms here are continuous.

Set \(E_\infty=E(\mu_{\ell^\infty})\), \(\Gamma=\Gal(E_\infty/E)\) and \(H=\ker(G\to\Gamma)\). For each \(v\mid\ell\) choose a decomposition subgroup \(D_v\) in \(G\). Let \(X'_\ell(E)\) be the maximal abelian pro-\(\ell\) quotient of \(H\) on which all subgroups \(H\cap gD_vg^{-1}\), \(g\in G,\ v\mid\ell\), become trivial. Equivalently, first pass to the maximal abelian pro-\(\ell\) quotient, then divide out the closed subgroup generated by those images. This compact module carries an action of \(\Gamma\) by conjugation. Write
\[
 X'_r=X'_\ell(E)\big/
   \overline{\langle \sigma b-\kappa(\sigma)^{-r}b:
                          \sigma\in\Gamma,\ b\in X'_\ell(E)\rangle}.
\]
Here the brackets denote the \(\ZZ_\ell\)-submodule generated by the displayed elements. For a compact abelian pro-\(\ell\) group \(B_0\), write \(B_0^\vee=\operatorname{Hom}_{\rm cts}(B_0,\QQ_\ell/\ZZ_\ell)\).

\begin{proposition}[Specialization of strict classes]\label{prop:spec12}
If \(\ell\) is odd, restriction gives an isomorphism \(\mathcal S_{\ell,r}(E)\simeq(X'_r)^\vee\). For \(\ell=2\) it gives a map into \((X'_r)^\vee\) with kernel of order at most two. In particular, whenever \(X'_r\) is finite,
\[
 \len_{\ZZ_\ell}\mathcal S_{\ell,r}(E)
        \le\len_{\ZZ_\ell} X'_r+\mathbf1_{\ell=2}.
\]
\end{proposition}
\begin{proof}
On \(H\) the module \(M_r\) is trivial, so a cocycle restricts to a homomorphism and a coboundary restricts to zero. The cocycle identity makes the homomorphism \(\Gamma\)-equivariant. A locally trivial class vanishes on \(H\cap gD_vg^{-1}\) after restriction. Its image is a finite \(\ell\)-group by continuity and compactness. Thus it factors through \(X'_\ell(E)\), and equivariance is exactly the displayed relation for \(X'_r\). Using \(D_v\) in place of the absolute local Galois group gives the same condition, since inflation in degree one for their surjection is injective.

For odd \(\ell\) the group \(\Gamma\) is an open subgroup of \(\ZZ_\ell^\times\). Its intersection \(\Gamma_1\) with \(1+\ell\ZZ_\ell\) is procyclic and isomorphic to \(\ZZ_\ell\), and the index is prime to \(\ell\). These statements follow because \(E\) is finite and \(1+\ell\ZZ_\ell\simeq\ZZ_\ell\) by the logarithm. If \(\gamma_1\) is a topological generator, continuous cohomology of \(\Gamma_1\) in degree one is \(M_r/(\gamma_1-1)M_r\) and it vanishes in degree two. Indeed, with \(T=\gamma_1-1\) the two-term resolution of the trivial module is multiplication by \(T\) on \(\ZZ_\ell[[T]]\), with cokernel \(\ZZ_\ell\). On \(M_r\) this map is multiplication by the nonzero \(\ell\)-adic number \(\kappa(\gamma_1)^{-r}-1\) in \(\ZZ_\ell\), hence is surjective. Averaging over the finite quotient of order prime to \(\ell\) makes its higher cohomology vanish. The group extension sequence for \(\Gamma_1\subset\Gamma\) thus gives \(H^1(\Gamma,M_r)=H^2(\Gamma,M_r)=0\).

Inflation and restriction now identify \(H^1(G,M_r)\) with \(\operatorname{Hom}(H,M_r)^\Gamma\), since the kernel is \(H^1(\Gamma,M_r)\) and the obstruction to surjectivity lies in \(H^2(\Gamma,M_r)\). A homomorphism corresponding to \((X'_r)^\vee\) vanishes on \(H\cap D_v\). The class it defines on \(D_v\) must therefore inflate from the image \(\Gamma_v\) of \(D_v\). This is again open in \(\ZZ_\ell^\times\), since the local field \(E_v/\QQ_\ell\) is finite. The same proof gives \(H^1(\Gamma_v,M_r)=0\). Consequently the global class is locally trivial at each \(v\), proving the isomorphism.

At \(\ell=2\) use \(\Gamma_1=\Gamma\cap(1+4\ZZ_2)\). It is isomorphic to \(\ZZ_2\), and the quotient \(C=\Gamma/\Gamma_1\) has order at most two. Again \(H^1(\Gamma_1,M_r)=0\). Inflation identifies \(H^1(\Gamma,M_r)\) with \(H^1(C,M_r^{\Gamma_1})\). Its coefficient group is the kernel of multiplication by \(\kappa(\gamma_1)^{-r}-1\) on \(\QQ_2/\ZZ_2\), hence finite cyclic. Degree-one cohomology for \(C\) is a subquotient of this cyclic group, by evaluation of a cocycle on a generator. It is killed by \(|C|\), so has at most two elements. The kernel of global restriction has at most that size as well. Duality for finite abelian \(\ell\)-groups preserves their order, which proves the length bound.
\end{proof}

\begin{lemma}[A determinant bound and exceptional twists]\label{lem:det12}
Let \(\Gamma_1=\langle\gamma_1\rangle\) as above, topologically. Suppose \(X'_\ell(E)\), as a module over \(\ZZ_\ell[[T]]\) with \(T=\gamma_1-1\), is a quotient of the cokernel of a square matrix \(P(T)\) over this ring. Put \(u=\kappa(\gamma_1)\), \(\theta_r=u^{-r}-1\). If \(\det P(\theta_r)\ne0\), then
\[
  \len_{\ZZ_\ell}\mathcal S_{\ell,r}(E)
       \le v_\ell\big(\det P(\theta_r)\big)+\mathbf1_{\ell=2}.
\]
Moreover let \(p(T)\in\ZZ_\ell[T]\) be monic of degree \(d\ge1\) and \(e=v_\ell(u-1)\). For integer \(j\ge0\), all \(r\ge1\) with \(v_\ell(p(\theta_r))>d(e+j)\) lie in at most \(d\) residue classes modulo \(\ell^{j+1}\). Here \(v_\ell(0)=+\infty\).
\end{lemma}
\begin{proof}
Specialization \(T\mapsto\theta_r\) is defined because \(\theta_r\in\ell\ZZ_\ell\). Its kernel is \((T-\theta_r)\), as is seen by translating the variable in the power series ring. Passing to the quotient is right exact. The cokernel of \(P(\theta_r)\) therefore surjects onto \(X'_\ell(E)/(T-\theta_r)X'_\ell(E)\), which in turn surjects onto \(X'_r\). The images involved are compact and hence closed. Over the discrete valuation ring \(\ZZ_\ell\), Smith normal form gives \(\len\operatorname{coker}P(\theta_r)=v_\ell(\det P(\theta_r))\). Apply Proposition~\ref{prop:spec12}.

For positive \(m\), \(v_\ell(u^m-1)=e+v_\ell(m)\). For \(m\) prime to \(\ell\) this follows from the first term of the binomial expansion. Raising a principal unit to the \(\ell\)-th power increases this valuation by exactly one, again by the binomial expansion, using \(e\ge2\) at \(\ell=2\). Thus \(v_\ell(\theta_r-\theta_s)=e+v_\ell(r-s)\) for distinct \(r,s\). Factor \(p(T)\) over \(\overline{\QQ}_\ell\), with the extended valuation. If the displayed threshold is exceeded, some root \(\alpha\) satisfies \(v_\ell(\theta_r-\alpha)>e+j\). Two indices attached to this root have \(v_\ell(\theta_r-\theta_s)>e+j\) by the ultrametric inequality, so \(r\equiv s\bmod\ell^{j+1}\). There are at most \(d\) roots.
\end{proof}

\begin{proposition}[A common integer for Frobenius trace comparisons]\label{prop:frob12}
Fix \(A=p^k\ge7\), with \(p\) prime and \(k\ge1\). Let \(Z_0\) be a finite multiset of roots of unity and \(B_0\) a finite multiset of algebraic numbers. Suppose both are stable under Galois conjugation, \(A b\) is an algebraic integer for \(b\in B_0\), and \(|b|\le3\) in every complex embedding. Set \(t_z=A^{-1}z+1+A z^{-1}\). There is a nonzero integer \(N_p\) with
\[
 \log|N_p|\le |Z_0|\,|B_0|\log(2A^2).
\]
For every \(\ell\ne p\), any embedding of the data into \(\overline{\QQ}_\ell\), and any submultiset \(D\subset Z_0\times B_0\),
\[
                \sum_{(z,b)\in D}v_\ell(b-t_z)\le v_\ell(N_p),
 \qquad v_\ell(\ell)=1 .
\]
\end{proposition}
\begin{proof}
Take \(N_p=\prod_{z,b}A(b-t_z)\). Each factor \(Ab-z-A-A^2z^{-1}\) is integral. The factor with \(A^2\) has modulus \(A^2\), and the other three terms have total modulus at most \(4A+1<A^2\). Thus every factor is nonzero of modulus at most \(2A^2\). Galois invariance of the product makes it a rational algebraic integer. This proves the asserted archimedean estimate, including the empty product. At \(\ell\ne p\), \(A\) is a unit, so all the summand valuations are nonnegative and the sum over the whole multiset is \(v_\ell(N_p)\). Discarding terms proves the inequality.
\end{proof}

In the proposed cohomological comparison, \(b\) is a Frobenius trace of a three-dimensional representation and \(t_z\) is the trace of a sum of three characters. The role of a comparison with two distinct primes \(p_1,p_2\) can be stated precisely. If nonnegative local lengths \(a_\ell\) have been bounded by \(2v_\ell(N_{p_i})\) for each \(\ell\ne p_i\), then
\[
             \sum_{\ell} a_\ell\log\ell
                    \le 2\log|N_{p_1}|+2\log|N_{p_2}|.
\]
Use the second integer at \(\ell=p_1\) and the first at all other primes. For either integer, the sum of all its valuations weighted by \(\log\ell\) equals the logarithm of its absolute value. The missing arithmetic input in this use of the proposition is the bound of actual cohomology lengths by trace valuations for a Galois-stable list of controlled size. Knowing only a determinant degree cannot yield such a bound. For example, at any fixed index \(r_0\) the degree-one polynomial \(T-\theta_{r_0}-\ell^L\) has specialized valuation \(L\) there for every positive integer \(L\).

For the tower of fields, a required estimate in the arithmetic construction takes the following form at a scale \(x\), with \(t=\lfloor x^{401}\rfloor\),
\begin{equation*}
 \sup_{S\ \mathrm{finite}}\sum_{\ell\in S}(\log\ell)
           \sum_{1\le r<t}\len_{\ZZ_\ell}\mathcal S_{\ell,r}(E)
                \ \le\ \exp(x^\eps/2),                                  \tag{L}\label{eq:length12}
\end{equation*}
where \(S\) varies over sets of rational primes. The specialization proposition gives an exact group to bound at odd primes and an upper bound at two. The determinant lemma controls the ordinary twists if suitable presentations, with quantitative bounds on their determinants, are available. Counting exceptional residue classes puts no upper bound on the valuation \emph{inside} such a class. A proof of \eqref{eq:length12} must control these terms, the dependence on \(E\), and the sum over \(\ell\) on the same intervals as the period comparison.

\subsection{Local divisibility and the rank of a period map}\label{nav:I-27}
\label{sec:geom11}
Here are two deductions from the kinds of local and archimedean estimates needed for \eqref{eq:gap}. They apply to arbitrary lattices and linear maps satisfying the stated hypotheses.

\begin{lemma}[Divisibility of a pairing]\label{lem:lattice11}
Let \(N,P\) be finite projective lattices over the integers of a number field \(J\), and let \(B:N_J\times P_J\to J\) be bilinear, where \(N_J=J\otimes N\) and likewise for \(P_J\). For finite places \(w\) suppose \(v_w(B(N_w,P_w))\ge a_w\), with \(a_w\) integers of finite support and \(v_w(0)=+\infty\). Put \(I=\prod_w\mathfrak p_w^{-a_w}\). Then \(B\) gives an integral map \(IN\to P^\vee\). For a fixed Euclidean metric at infinity and \(d=\operatorname{rank} N>0\), set \(\mu(N)=-2\log\operatorname{covol}(N)/([J:\QQ]d)\). One has
\[
 \mu(IN)-\mu(N)=\frac{2}{[J:\QQ]}\sum_w a_w\log N\mathfrak p_w .
\]
\end{lemma}
\begin{proof}
At \(w\), the first lattice is multiplied by \(\pi_w^{-a_w}\), where \(\pi_w\) is a uniformizer. Every pairing is then integral at \(w\). Intersection of the local integer rings inside \(J\) gives the asserted global integral map. The determinant line is multiplied by \(I^d\). For an integral ideal this changes covolume by its index to the power \(d\), as follows by counting local quotients over the residue fields. Taking ratios gives the same rule for a fractional ideal. In this case the factor is \(\prod_w(N\mathfrak p_w)^{-da_w}\), proving the formula.
\end{proof}

\begin{proposition}[Loss of rank under restriction]\leavevmode\label{prop:rank11}
Let \(A_i:V_i\to\CC^D\) be linear maps between Hermitian spaces, \(D>0\), \(R=\sum_i\dim V_i>0\), and \(Q_0=\sum_i A_iA_i^*\). Write \(\|M\|_{\rm HS}^2=\operatorname{tr}(M^*M)\). Suppose \(0\le e\le1/4,\ \eta\ge0\) and
\[
 \sum_i\|A_i^*A_i-I\|_{\rm HS}^2\le eR,\qquad
 D\operatorname{tr}(Q_0^2)\le(1+\eta)(\operatorname{tr}Q_0)^2.
\]
If \(U_i\subset V_i\) have total codimension \(\epsilon_0R\), then \(W=\sum_i A_i U_i\) satisfies
\[
              \operatorname{codim}(W)/D\le\eta+4\epsilon_0+4\sqrt e .
\]
\end{proposition}
\begin{proof}
Put \(a=\sqrt e\), \(T_0=\operatorname{tr}Q_0\). The first inequality and Cauchy--Schwarz give \(T_0\ge(1-a)R>0\). Let \(P,P_i\) be the orthogonal projections to \(W^\perp,U_i^\perp\) respectively. Since \(PA_i=PA_iP_i\),
\[
 \operatorname{tr}(PQ_0)=\sum_i\|PA_iP_i\|_{\rm HS}^2
 \le \epsilon_0R+\sum_i \|P_i\|_{\rm HS}\|A_i^*A_i-I\|_{\rm HS}
 \le R(\epsilon_0+a\sqrt{\epsilon_0}).
\]
Let \(d_0=\operatorname{rank}(P)/D,\ p_0=\operatorname{tr}(PQ_0)/T_0\).
The second hypothesis gives \(\operatorname{tr}((DQ_0/T_0-I)^2)\le D\eta\), whence \(|p_0-d_0|\le\sqrt{d_0\eta}\) by Cauchy--Schwarz. As \(2\sqrt{d_0\eta}\le d_0+\eta\), it follows that \(d_0\le2p_0+\eta\). Now \(p_0\le2(\epsilon_0+a)\), proving the claim.
\end{proof}
\begin{corollary}[From the ranks to mean slopes]\label{cor:slopes12}
In Proposition~\ref{prop:rank11}, give the source spaces finite decreasing real filtrations and list their jumps, with dimension multiplicities, as shifted slopes. Suppose their overall mean is at least \(s\) and each is at most \(s+v\) with \(v\ge0\). Let a quotient of \(\CC^D\) of dimension \(d\ge(1-\delta_0)D>0\) have a filtration with all slopes at least \(s-H_0\), \(H_0\ge0\). Assume that for each threshold the projected maps carry its source filtration spaces into the corresponding target filtration space. For the mean \(\mu\) of the target slopes one has
\[
 s-\mu\le
 \frac{4v\log(1+H_0/v)+(\eta+4\sqrt e)H_0}{1-\delta_0},
 \qquad 0\le\delta_0<1 .
\]
At \(v=0\) the first term of the numerator is defined to be zero.
\end{corollary}
\begin{proof}
A fraction \(p_\tau\) of source slopes below \(s-\tau\), \(\tau>0\), forces the mean to be at most \(p_\tau(s-\tau)+(1-p_\tau)(s+v)\). Hence \(p_\tau\le v/(v+\tau)\). Use the source spaces at threshold \(s-\tau\) as the \(U_i\) in Proposition~\ref{prop:rank11}. Linear projection to a quotient does not increase the absolute codimension of a span. Therefore the fraction of target slopes below \(s-\tau\) is at most \((\eta+4\sqrt e+4v/(v+\tau))/(1-\delta_0)\). For each target slope \(a\), the integral of \(\mathbf1_{a<s-\tau}\) over \(0\le\tau\le H_0\) is \((s-a)_+\). Average this identity and integrate the displayed bound. When \(v=0\) the bound for \(p_\tau\) is zero for each \(\tau>0\).
\end{proof}
In particular, a local divisibility computation giving a lower bound \(\mu+\Gamma_0-D_0\le s\) for this \emph{same} collection of maps bounds its gain \(\Gamma_0\) by \(D_0\) plus the displayed error. Applying this to integral periods requires both the lower bound and the comparison of filtrations and Gram matrices for the chosen maps.

At the archimedean places the raw quantities in one such comparison are block Gram matrices \(K_{ij}=A_i^*A_j\), before restriction, with \(r_i=\dim V_i\). For a comparison dimension \(D_*>0\), the proposed estimates (after fixing the size normalization of the blocks) include
\[
 \|K_{ij}\|_{\rm HS}^2\le(1+\eta_0)r_ir_j/D_*\ (i\ne j),\qquad
 \|K_{ii}-I\|_{\rm HS}^2\le e_0r_i,\qquad D\le(1+\eta_1)D_* ,
\]
with \(\eta_0,\eta_1,\sqrt{e_0}\le\exp(-x^{\eps/8})\). These are quantitative estimates for actual theta periods. For instance, the reason for controlling different blocks in the first inequality is the exact identity \(\operatorname{tr}((\sum_i A_iA_i^*)^2)=\sum_{i,j}\|K_{ij}\|_{\rm HS}^2\), obtained by cyclic invariance of the trace. It is this sum that governs the rank estimate above.

\subsection{The estimates on a common interval}\label{nav:I-28}
\label{sec:missing}
\noindent The following paragraphs give the common field, interval and comparison inequalities. The companion volumes develop the two estimates by subject. The \hyperref[seriesguide]{volume guide} gives their locations.
The preceding calculation explains the choice of chain length in the \emph{Remark}. Two terms of the normalized error have sizes governed by \(x/m\) and \(m\). A choice of \(m\) near \(\sqrt{x}\) balances them. This calculation starts after the bound \eqref{eq:gap} for the actual counts has been obtained.

In the construction, \(E\) varies with \(x\). One needs control of the weighted lengths of local restriction kernels in the cohomology of \(E\), and integral period maps for a chain of theta lifts over that same field. Local divisibility is to give a lower bound for the gain in the period lattice. An upper bound involving its covolume, its ranks and the archimedean norms of the period maps must give \eqref{eq:gap} for every translate. Both bounds must apply to data obeying conditions~\ref{cond:sizes}--\ref{cond:weights} at each point of one whole interval. Two independently unbounded sets of admissible scales can be disjoint. Even at \(o=1\), when the only retained term is \(G_{\mathbf1}\), the count estimate is needed.

\lead{The residual local groups}Fix a finite set \(\Pi\) of odd primes containing all odd divisors of \(o=\operatorname{ord}(\psi)\). Use the fields and the single pair \(q,d\) on the arbitrarily distant intervals \([V,V^{3/2}]\) of Corollary~\ref{cor:tower27}. Write \(\eta=\eps/128,\ t=\lfloor x^{401}\rfloor,\ C_j=\QQ(\zeta_{2^{j+2}})^+\), and set
\[
 L_x=K_\psi C_{\lfloor\eta\log_2 x\rfloor}(\sqrt{-q},\sqrt{-d}),\quad
 B(x)=K_\psi C_l,\quad
 P=\exp(x^\eps/12),\qquad b_0=\lfloor\exp(x^\eps/8)\rfloor.
\]
Here \(l\) is the least integer with \(2^l\ge x^{2+8\eps}\). Both subfields lie in \(E_x\) for all sufficiently large \(x\). For an odd prime \(\ell\), put \(S=\mathcal S_{\ell,r}(E_x)\) and \(M=eS\), where \(e=(1+(-1)^rc)/2\) is the complementary-sign idempotent. Define an operator \(J_{\ell,r}\) on \(M\) as follows. For \(r\le x^\eta\) use the average of \(\Gal(E_x/L_x)\). For \(r>x^\eta\) and \(\ell\in\Pi\) use the average of \(\Gal(E_x/B(x))\). For the other pairs use the zero map. Both averages are integral idempotents at odd primes since the relative degrees are powers of two. Let \(U_{\ell,r}=(1-J_{\ell,r})M\) and put
\[
 R_{\ell,r}=
 \begin{cases}
 U_{\ell,r}/U_{\ell,r}[\ell^{b_0}],& \ell\le P,\\
 U_{\ell,r},&\ell>P .
 \end{cases}
\]

\begin{proposition}[The part still needed in the local-length estimate]\label{prop:res23}
Denote the left side of \eqref{eq:length12} by \(\mathcal L(x)\) and define
\(\mathcal R(x)=\sum_{\ell>2}(\log\ell)\sum_{1\le r<t}\len_{\ZZ_\ell}R_{\ell,r}\), where \(\ell\) runs through primes. On these intervals,
\begin{equation}
 0\le\mathcal L(x)-\mathcal R(x)
 \le \tfrac12\exp(x^\eps/2)+2\exp(x^{2\eps/3})+\exp(x^\eps/3).
 \label{eq:res23}
\end{equation}
In particular \(\mathcal R(x)\le\tfrac14\exp(x^\eps/2)\) throughout arbitrarily distant such intervals would suffice for \eqref{eq:length12}.
\end{proposition}
\begin{proof}
All the groups are finite and for each \(x\) the sum has finite support, by Proposition~\ref{prop:finite15}. Its supremum over finite sets of primes is its full sum. The prime two and the sign \((-1)^{r+1}\) at all odd primes use at most \(\tfrac14\exp(x^\eps/2)\) each by Theorems~\ref{thm:two17} and~\ref{thm:critical13}. At each odd prime the other sign is \(M=J_{\ell,r}M\oplus U_{\ell,r}\), since the field averages commute with \(c\). The weighted total of the first summands is at most \(2\exp(x^{2\eps/3})\) by Theorem~\ref{thm:growing} and Corollary~\ref{cor:tower27}. These apply to the two disjoint twist ranges in the definition of \(J_{\ell,r}\). For \(\ell\le P\), the kernel of \(U_{\ell,r}\to R_{\ell,r}\) is contained in \(S[\ell^{b_0}]\); Theorem~\ref{thm:trunc15}, bounds the sum of their lengths with weights by \(\exp(x^\eps/3)\). Additivity of length proves \eqref{eq:res23}. Both further terms there are \(o(\exp(x^\eps/2))\) as \(x\) tends to infinity, giving the sufficient bound for the remainder.
\end{proof}

Theorem~\ref{thm:norm28}, gives an exponent test at every odd prime for components with nontrivial residual character. Its Corollary~\ref{cor:interval} bounds the congruence ideals of characteristic-zero Hecke traces. Proposition~\ref{prop:root31}, exhibits actual extension classes whose marked deformations have identical traces. An upper bound from an integral deformation model needs information about its relations in addition to its characteristic-zero trace values.

\begin{lemma}[Characteristic-zero points and an integral length]\label{lem:vertical23}
For every prime \(\ell\) and every \(b\ge1\) there is a finite local \(\ZZ_\ell\)-algebra with a single characteristic-zero geometric point and whose augmentation ideal modulo its square has length \(b\).
\end{lemma}
\begin{proof}
Use \(A_b=\ZZ_\ell[X]/(X^2,\ell^b X)\) and the augmentation \(X\mapsto0\). Its nilpotent ideal \(I=(X)\) has quotient \(\ZZ_\ell\), so \(A_b\) is local with maximal ideal \((\ell,X)\). As a module it is \(\ZZ_\ell\oplus(\ZZ_\ell/\ell^b)X\). Every characteristic-zero point over \(\ZZ_\ell\) sends the nilpotent \(X\) to zero, and the rational generic fiber is \(\QQ_\ell\). Yet \(I^2=0\) and \(I/I^2\) has length \(b\).
\end{proof}
This illustrates why the model of the deformation ring, including its integral relations, enters the length inequality. The example by itself makes no assertion about the deformation rings of the present fields.

At every odd prime, Theorem~\ref{thm:levelstrict345} compares the full complementary length over the period field with scalar groups having the local conditions at the conductor. The total discrepancy is polynomial in the interval scale. Corollary~\ref{cor:matchedlength617} gives the exact elementary divisors added by relaxing the strict condition at any number of primes in a common split Frobenius class of prescribed cyclotomic depths. A single prime in that class preserves every indicated complementary group at every precision. Theorem~\ref{thm:tamedivision618} chooses another common set using residual dual classes. It makes every original strict class divisible by a prescribed power in the relaxed group. At cyclotomic depth one, Corollary~\ref{cor:tamelayers618} recovers each successive layer of the length sum from the two sets of primes. Corollary~\ref{cor:periodblock346} proves full vanishing for the stated residual characters by a two-prime decomposition-group test on the original fields. Theorem~\ref{thm:heckedivisor345} constructs the common divisor of the Frobenius tests over one auxiliary quadratic field. Corollary~\ref{cor:heckebudget345} states the full-length inequality still needed at the large primes. On the same period fields, Theorem~\ref{thm:evenroot358} gives the full integral ring of lower triangular lifts preserving the marked symmetric form. Its specializations count the complementary strict classes at every depth. Theorem~\ref{thm:paritydescent358} determines which irreducible orthogonal points can be automorphic over a CM extension. Theorem~\ref{thm:mixedtrace359} realizes the dual of the full large-prime complementary group as a quotient of the cotangent module of the integral trace algebra. Theorem~\ref{thm:classicaldefect360} bounds, by \(\log J_*\) on the entire interval, the contribution after quotienting by the image of the trace relations that vanish at the listed classical systems. At each large prime Theorem~\ref{thm:adjointseed361} constructs cuspidal ordinary adjoints over a single auxiliary CM field for all the period data on the interval. Theorem~\ref{thm:auxdescent361} retains the entire integral trace quotient there and descends automorphy for any strongly irreducible point for which it holds on that field. At the cuspidal adjoints, Theorem~\ref{thm:weightchart362} identifies the full characteristic-zero ordinary deformation ring with its weight ring. Theorem~\ref{thm:traceideal362} computes the ideal of inverse-trace differences there, including at every infinitesimal order in symmetric weights. Theorem~\ref{thm:slthree363} constructs cuspidal representations with full special linear monodromy at the prescribed weights on one auxiliary field for the entire interval. Corollary~\ref{cor:skewsize363} bounds a nonzero inverse-trace difference in valuation. Proposition~\ref{prop:descentlayer363} computes the exact local obstruction to descending these points to the original field with crystalline type. In residue characteristic, Theorem~\ref{thm:unitweights364} gives the equations imposed by global units on sums of three characters. Theorem~\ref{thm:rescurve364} uses them to construct a nonabelian trace curve with distinct characters on the kernel of cyclotomic inertia at every place above the coefficient prime. Theorem~\ref{thm:commonframe365} establishes the residual cohomology hypotheses for the entire interval on one sufficiently deep auxiliary field. Proposition~\ref{prop:tamecurve365} places the trace curve on its specified Galois quotient. Theorem~\ref{thm:markedcurve366} realizes that curve in the marked ordinary deformation ring with a perfect symmetric pairing, and verifies its local and generic point conditions. Theorem~\ref{thm:primes367} constructs simultaneous split Frobenius tests for the dual cohomology on these curves at every parameter precision, with a single loss independent of the precision. Theorem~\ref{thm:selmerlength368} computes the leading term of the first-order Selmer length there, uniformly through the auxiliary prime levels and all powers of the curve parameter. Theorem~\ref{thm:selmerlift369} bounds the obstruction to lifting ordinary cocycles on the same levels and gives the second numerical input for patching. Theorem~\ref{thm:iwahori370} constructs integral comparison maps between hyperspecial and Iwahori levels. Corollary~\ref{cor:levelcurve370} computes their loss at the Frobenius polynomials of the same curves. Theorem~\ref{thm:stableprime371} places every selected curve on the closure of a family with absolutely irreducible generic trace on the deep auxiliary field. Corollary~\ref{cor:curvesupp371} proves that the actual ordinary modules on that field have positive rank upon specialization to the curve at every selected auxiliary level. Theorem~\ref{thm:adjcomponent372} identifies the unique full ordinary determinant component through the cuspidal adjoint with the corresponding integral Hecke component. Corollary~\ref{cor:classfiber372} proves classicality of its characteristic-zero points at the prescribed weights. Proposition~\ref{prop:origideal373} identifies the integral augmentation ideal on the original quadratic field with the four-by-four trace minors. Theorem~\ref{thm:minorstrict373} recovers the full complementary strict group from these minors even after restriction to the auxiliary field. Theorem~\ref{thm:restrictedtrace374} proves equality of the complete fixed-weight trace rings on the original and auxiliary fields. Proposition~\ref{prop:orighecke374} constructs a finite quotient on the original field, and Corollary~\ref{cor:irreduciblelabel374} gives the original cuspidal labels for its irreducible marked points. Proposition~\ref{prop:tracecrystal375} proves crystallinity for global trace specializations, including the first twist. Theorem~\ref{thm:alllabels375} puts all non-self-dual characteristic-zero points of the finite quotient on the original cuspidal list, giving the remaining length in Corollary~\ref{cor:componentdefect375}. Theorem~\ref{thm:reduciblefiber376} shows that reducible characteristic-zero points contribute zero inverse-trace difference throughout their infinitesimal fibers. Its proof uses the odd descent and Selmer vanishing of Proposition~\ref{prop:odddescent376}. Theorem~\ref{thm:rigidtrace377} proves reducedness at every irreducible point of the finite quotient. Proposition~\ref{prop:torsionterm377} computes the exact multiplicity and torsion terms in its inverse-trace ideal, and Corollary~\ref{cor:twointegral377} identifies the two integral contributions still to be bounded. Theorem~\ref{thm:reducibleminors378} constructs a finite ring on the original quadratic field using the determinant of the trace pairing. Its torsion maps onto the full complementary group by Corollary~\ref{cor:torsionimage378}. Theorem~\ref{thm:reduciblerigid379} identifies the rational local rings at all reducible points. Theorem~\ref{thm:reducibleorder379} bounds the cotangent contribution of their torsion-free order by one explicit integer throughout the interval. Theorem~\ref{thm:planedesc381} extends the plane trace to an odd two-dimensional determinant over the full integral quotient. Theorem~\ref{thm:adjtorsion381} realizes the full dual strict group as an integral direct summand in the other factor of a fiber product on the original field. Theorem~\ref{thm:tensortrace382} determines that factor by an integral tensor product. Proposition~\ref{prop:upperorders382} computes the upper elementary divisors using the preceding twist, with an ideal class group formula at the first twist. Theorem~\ref{thm:lowerfibers384} counts the fibers above each strict class in the full lower triangular deformation problem. Proposition~\ref{prop:odddivisor384} bounds their total logarithmic contribution at every coefficient depth on the same interval. Outside the prime divisors of one explicit integer, Theorem~\ref{thm:orthcollapse384} identifies the entire even orthogonal deformation ring with the lower triangular ring. Theorems~\ref{thm:flagthick385} and~\ref{thm:flagvaluation385} compare the integral trace ring of families with an invariant direct summand to the ring cut out by the largest trace minors. The latter is a nilpotent extension of the former of index at most eight. Corollary~\ref{cor:fourcoeff385} realizes the complementary elementary divisors on finite coefficient algebras of minimal rank. Theorem~\ref{thm:cyclicstrict397} shows how a nonzero cyclic complementary group gives either a nonconstant irreducible trace in residual characteristic or a finite flat intersection whose Frobenius determinant bounds its full length. Theorem~\ref{thm:negfit398} recovers the full order at the split trace for any number of generators away from the critical determinant. Corollary~\ref{cor:negwedge398} identifies the higher exterior powers there with those of the strict group dual. Theorem~\ref{thm:horizontalstrict399} bounds the entire strict length by the Frobenius determinant of a finite trace algebra, with geometric multiplicities, in the absence of a residual trace curve. Corollary~\ref{cor:frobeniusstrict399} realizes the minimum determinant valuation at split primes. Theorem~\ref{thm:allhorizontal401} gives the multiplicity bound at every large prime. Theorem~\ref{thm:gaussianbudget401} bounds by a single integer on the interval the full lengths accounted for by the listed classical traces, including components with several strict generators. Corollary~\ref{cor:gausssum402} includes in that estimate every finite component whose irreducible points become cuspidal over solvable CM extensions. The auxiliary extensions can vary with the point. Theorem~\ref{thm:cmcores403} also permits a varying intermediate CM field of bounded degree and controlled discriminant. Its proof counts all such fields and tests every trace with a common Frobenius power. Theorem~\ref{thm:realtwist404} removes all finite ramification of the conductor characters on a common totally real field. Proposition~\ref{prop:strictdescent404} recovers the strict groups on the original period fields by taking invariants. Theorem~\ref{thm:alllevels405} constructs the ordinary determinant on the three integral modules used in the change of level. Corollary~\ref{cor:annram405} gives an annihilator for the extra ramification on diamond coinvariants. Theorem~\ref{thm:inertiatrace406} pairs Frobenius roots with their inertia characters on the integral diagonal cover. Theorem~\ref{thm:finitelevel406} gives the resulting ordered finite-level comparison on each of the auxiliary sets. Theorem~\ref{thm:limitkernel407} identifies the kernel created by the varying primes in a common limiting group. Theorem~\ref{thm:limitselmer407} establishes the ordinary tangent bound, the simultaneous lifting relation and integral descent for the off-diagonal cohomology on that group. Theorems~\ref{thm:modulepatch408} and~\ref{thm:patchrelations408} construct a finite patched module at the discriminant of these curves and identify its fiber with the ordered hyperspecial module. Theorem~\ref{thm:irrcurvesupport409} identifies the local ordinary determinant ring with the local Hecke algebra at each selected curve whose generic representation is absolutely irreducible. Every component through such a curve is in Hecke support. At a curve with a character and a plane, Theorem~\ref{thm:oddsupport410} identifies the irreducible and the nonsemisimple fibers of the ordinary Galois algebra. The minimum number of generators of its off-diagonal corner is the ordinary extension rank on the original field by Corollary~\ref{cor:selmergenerators410}. The lift of mixed ordinary cocycles in Theorem~\ref{thm:limitselmer407} gives the square matrix of relations in Theorem~\ref{thm:oddrelations411}. Theorem~\ref{thm:reduciblepropagation412} then places all irreducible and all reducible nonsemisimple primes below each selected two-constituent curve in ordinary Hecke support. Theorem~\ref{thm:incidenceannihilator413} computes the exact annihilator of the off-diagonal module in the completed patched ring. After setting the auxiliary diamonds to one, Theorem~\ref{thm:cornermatrix414} computes the integral extension module at reducible specializations of the completed ordinary ring by a square matrix and its minors. The fixed arithmetic weight is disjoint from that localization. Theorem~\ref{thm:strictcorner415} and Corollary~\ref{cor:strictminors415} recover the actual large-prime lengths on each period field from quotients of the unlocalized ordinary algebra at the split weight. Uniform upper bounds for their minors remain necessary. Theorem~\ref{thm:strictbasedescent416} gives exact descent of strict cohomology at the auxiliary prime, at every coefficient precision. Corollary~\ref{cor:intervalbase416} applies this to all the period fields on the interval with one common auxiliary extension. Theorem~\ref{thm:chaintransfer417} propagates the ordinary Hecke support along finite chains of separated curves with a line and an irreducible plane from one supported irreducible prime at the first curve. Theorem~\ref{thm:completehecke418} identifies the ordinary ring with the Hecke algebra and proves it is a complete intersection at every supported irreducible curve with the stated inertia separation. Corollary~\ref{cor:mixedchain418} allows both types of curves in one chain. Theorem~\ref{thm:largeintersection419} gives an explicit dimension bound on a component intersection that guarantees transfer of support. Corollary~\ref{cor:componentgraph419} applies it starting from the full component through the cuspidal adjoint. Theorems~\ref{thm:markedpresentation420} and~\ref{thm:tracefiber420} give a presentation of the marked ring and a dimension bound for the map to traces on the one- and two-constituent loci. Proposition~\ref{prop:lowerdimension421} bounds the dimension of the triangular locus by means of strict cohomology on the same field. Theorem~\ref{thm:adjacentchart421} transfers support throughout a marked chart with a nonzero adjacent lower entry. Theorem~\ref{thm:longrootchart422} proves transfer to every irreducible determinant in a chart with only a long lower entry at the residual point. Its proof treats intersections with reducible components by a first-order deformation of the same Galois algebra. Theorem~\ref{thm:globalhecke423} proves Hecke support at every geometrically irreducible determinant over the common layer. Corollary~\ref{cor:irreducibleweight423} proves finiteness of their reduced closure after specializing to the original arithmetic weight. Theorem~\ref{thm:originalfinite424} proves that the entire original fixed-weight trace ring is finite. Theorem~\ref{thm:originalorder424} bounds the determinant on its torsion-free augmentation ideal by one integer throughout the period interval. Corollary~\ref{cor:periodlarge425} bounds the entire large-prime complementary length on each original period field by a polynomial in the interval scale. Corollary~\ref{cor:remainingresidues426} also eliminates the high-power tails outside the stated set of coefficient primes in all but at most eight residue classes of twists modulo \(\ell-1\). Theorem~\ref{thm:adjacentweights427} controls every twist with three distinct residual inertia characters at the indicated primes. Corollary~\ref{cor:smallorder427} leaves only the two colliding residues there. Theorem~\ref{thm:movingprimes428} also controls both moving ramified coefficient primes, including all the characters ramified at those primes. Corollary~\ref{cor:unboundedprimes428} locates the remaining small-prime tails at a fixed finite set of primes and at the two collision residues. Theorem~\ref{thm:symplecticclasses429} realizes every scalar component in that range with a nontrivial residual character and twist at least two by a symplectic matrix entry, with an explicit bound on the precision needed for the local splittings. This construction alone gives no bound on the lengths of these components. Theorem~\ref{thm:oddrootlocalization430} computes the global localization map for each of the six odd symplectic roots at every coefficient depth. The remaining scalar entry is even. Theorem~\ref{thm:siegelclasses431} computes the two ordinary extension spaces preserving either paired plane also at a nonsplit strict class. Corollary~\ref{cor:siegeltrace431} shows exactly which characteristic polynomials persist in both spaces. Proposition~\ref{prop:evenplanedeterminants433} excludes any two-dimensional constituent of scalar complex conjugation at a point satisfying the stated Frobenius spread bound. Theorem~\ref{thm:commonsptrace433} gives one integer controlling the trace congruences for any specified integral cuspidal spectrum on the interval. Theorem~\ref{thm:strictrootring434} identifies the exact finite-depth strict group in the normalized lower root using the scheme-theoretic images of local conjugation. Its dual is a quotient of the cotangent module in Corollary~\ref{cor:saturatedcotangent434}. At two coincident inertia weights, Theorem~\ref{thm:reduciblecollision435} proves finiteness of the reducible trace locus whenever Frobenius separates the three characters. Corollary~\ref{cor:undistinguished435} puts the exceptional characters in explicit dyadic subfields of the same period fields. Theorem~\ref{thm:twoprimeadjoints436} uses a second auxiliary prime to construct a non-CM ordinary form at the original weight even for these exceptional characters. Theorem~\ref{thm:actualseedcongruences436} bounds its trace congruences on the common interval. Theorem~\ref{thm:orderedconductor437} bounds the integral loss in recovering all ordered Frobenius characters from good traces at the original weights. Corollary~\ref{cor:seedcomponent437} identifies the ordinary Galois component through the constructed adjoint under the stated tame conditions. Theorems~\ref{thm:splitwordfiber438} and~\ref{thm:flagfiber438} compute the complete residual fiber of the split trace before and after imposing the local characteristic polynomials of a flag. Positive-dimensional fibers of the word equations occur on the original intervals. Corollary~\ref{cor:flagcotangent438} bounds the contribution of ordered character labels to the cotangent module of every such flag quotient at the original weight. Theorem~\ref{thm:scalarinertiaplanes439} proves integral finiteness of the reducible symmetric trace quotient at both collision residues for every globally distinct residual triple on the auxiliary quadratic field, even when its full local triple is scalar. Corollary~\ref{cor:flagfinitered439} gives the same finiteness with the local characters of an actual flag adjoined. Theorem~\ref{thm:allzeroweights442} removes the residual-order condition at every twist divisible by \(\ell-1\). The field used for the descent is the compositum of the real quadratic subfields of the largest period field on the interval. Theorem~\ref{thm:inertquadratic443} treats the other collision residue at primes inert in the auxiliary imaginary quadratic field when the reduced conductor character is nontrivial and the outer character is quadratic. The ramified plane has two fixed local characters, allowing a global deformation with unchanged inertia at the coefficient prime. Theorem~\ref{thm:rayconductor444} computes the ramification of the induced residual representation at the coefficient and discriminant primes. Corollary~\ref{cor:rayinterval444} bounds its conductor and field degree by fixed powers of the original interval scale. Theorem~\ref{thm:allcollisiontraces445} proves finiteness of the entire reducible trace quotient at every weight collision for primes at least five prime to the prescribed level and quadratic discriminant. At split primes the local line gives an ordinary deformation and at inert primes the entire local trace is constant. On the polynomial-conductor intervals, Theorem~\ref{thm:markedstrict337} makes all the stated descended strict groups vanish at the same three primes used in the period data. Theorem~\ref{thm:criticalmarks339} covers the entire critical sign on the full period fields at these three coefficient primes. Corollary~\ref{cor:allsample341} bounds the sum of the full complementary orders in logarithmic scale outside class norm exceptions in every tame character through its prime cutoff. By Corollary~\ref{cor:periodlarge425} the total residual contribution above \(P(x)\) is polynomial in \(V\). In Proposition~\ref{prop:res23} the remaining odd-prime contribution comes from the quotients \(U_{\ell,r}/U_{\ell,r}[\ell^{b_0}]\) for \(\ell\le P(x)\). Theorem~\ref{thm:rationalpolynomial538} bounds the logarithmic orders of the rational even-twist groups by a polynomial in the weight. Corollary~\ref{cor:rationalperiod538} removes the invariant part of these residual quotients on every original field. Theorem~\ref{thm:oddquarticfield539} removes every odd character of order two or four inherited from the fixed dyadic tower. Corollary~\ref{cor:fixedimaginary539} accounts for every complementary weight on each fixed imaginary quadratic subfield of that tower. Corollary~\ref{cor:complexallweights542} removes all weights on a real or complex subfield of degree of order \((\log V)^2\) at the coefficient primes \(V\le p\le P(x)\) congruent to three or five modulo eight. Corollary~\ref{cor:octicspectrum543} treats the odd weights on the fixed subfield of exponent eight. Theorem~\ref{thm:allprimeslayers543} treats all weights on a single growing subfield down to the smallest primes in both progressions. Corollary~\ref{cor:collisionloss545} bounds the total extra length admitted by the polynomial ordinary equations on the strict symplectic root at the colliding weights, uniformly in the finite precision. Theorems~\ref{thm:finiteflag546} and~\ref{thm:doubleflag546} recover an ordinary flag from finite matrix equations and determine the sharp uniqueness precision when the residual character is locally trivial. Theorems~\ref{thm:allglobalroots547} and~\ref{thm:integralroots547} recover the strict groups at every depth from the global symplectic equations for all locally nontrivial even residual characters at the two colliding residues. Theorem~\ref{thm:tracefitting548} identifies the full strict length with the valuation of a Fitting ideal at the split trace, for every residual rank in that calculation. Theorem~\ref{thm:ramification549} gives the least ramification index needed to realize each integral trace by a global ordinary representation; Theorem~\ref{thm:purityblocks549} computes the Frobenius determinant obstruction on its two-plane locus. Theorem~\ref{thm:lengthmultiplicity550} bounds the entire strict length by Frobenius valuations with generic multiplicities when the trace algebra is finite. Corollary~\ref{cor:pointalternative550} gives a representation in characteristic zero or a nonconstant representation in characteristic~\(\ell\) whenever this strict group is nonzero. Theorem~\ref{thm:twomarks557} recovers the closed trace algebra from the source's two marks even with repeated local residual eigenvalues. Corollary~\ref{cor:finmark557} identifies its finiteness clause for each covered datum with the one in the strict-length inequality. Theorem~\ref{thm:rankfourselmer558} embeds the tangent at an irreducible characteristic-zero trace into the finite symplectic adjoint Selmer group. Proposition~\ref{prop:finiteordinary558} verifies the coefficient-prime condition on first-order lifts, including weight differences of one. Theorem~\ref{thm:enormousall559} verifies the cyclotomic image condition for every finite extension at the covered characteristic-zero points. Theorem~\ref{thm:automult559} gives multiplicity one whenever a rank-four point has a polarized cuspidal companion over a totally real or CM field. A cuspidal match over the fixed field \(\QQ(i)\) already gives multiplicity one and integral Frobenius eigenvalues of absolute value \(p^{(r+1)/2}\) by Proposition~\ref{prop:onecompanion560} and Theorem~\ref{thm:integerpurity560}. Theorem~\ref{thm:evenplanecompanion560} excludes a cuspidal match for either even plane by the sign of its pairing. Corollary~\ref{cor:uniformstrict561} gives the polynomial bound for the sum of strict lengths on each period field for triples with finite trace algebra and cuspidal matches at all nonsplit points. The levels, Frobenius prime, multiplicities and common integer are supplied by the argument. Proposition~\ref{prop:gooddata563} proves the odd strict vanishings for the exact character conductors prescribed in both source cases. Theorem~\ref{thm:sourcecommon563} gives the common bound for those conductors under the three stated global arithmetic hypotheses. Theorem~\ref{thm:planeunits565} excludes the even plane by either of two explicit residue sums and proves full strict vanishing when the negative-root sum is nonzero. Corollary~\ref{cor:planeinterval565} removes the plane hypothesis on the corresponding triples. Theorem~\ref{thm:planedepth566} computes the trace congruence ideal of an even plane and bounds its depth by a critical Dirichlet value. In Corollary~\ref{cor:planeweight566} a polynomial bound on the total multiplicity of such points suffices in place of excluding them. When the two strict residual dimensions are zero and one, Theorem~\ref{thm:intrinsicplane567} bounds the sum of all congruence depths inside the plane locus by a single critical value. Corollary~\ref{cor:transverseplane567} accounts for the two transverse coordinates separately. At any integral two-plane point, Theorem~\ref{thm:ordinarycross568} computes the ordinary extensions in each off-diagonal direction at every finite coefficient depth. Their ranks are two and one, and each one-sided extension has the original traces. Proposition~\ref{prop:mixedobstruction569} removes the scalar part of the obstruction to mixing the two extensions. Theorem~\ref{thm:mixedtrace569} identifies the kernel of the remaining traceless equation with the transverse trace tangent space. Theorem~\ref{thm:tametower571} constructs the totally real and CM fields required for the odd-order case with two auxiliary quadratic fields common to the interval. Corollary~\ref{cor:towerdescent571} descends the strict groups to each original period field. Theorem~\ref{thm:eventower572} supplies the two local witnesses and the nontrivial unramified quadratic character in the even-order tower. Corollary~\ref{cor:bothaux572} gives the field construction and descent on the same interval in both cases. Proposition~\ref{prop:commonaux573} constructs one field of polynomial Galois exponent for the entire interval. Theorem~\ref{thm:splitintegral573} proves integral finiteness over the traces at primes split in that field, and Proposition~\ref{prop:separatedalgebra573} gives equality over each separated auxiliary base. Theorem~\ref{thm:matrixden574} recovers the four diagonal projectors over this common splitting field with an explicit integral denominator. Corollary~\ref{cor:splitloss574} bounds the loss in trace valuations and preserves the complete generic trace algebra. Theorem~\ref{thm:blockwitness575} constructs an ordinary nonsplit extension at every two-plane point in either characteristic. Corollary~\ref{cor:commonwitness575} detects a failure of finiteness at a rational prime split in the common field. Theorem~\ref{thm:cuspidalinitial577} constructs a cuspidal symmetric-cube initial point with the required residual lattice, multiplier and ordered flags for every odd-order triple, using one level prime throughout the interval. Theorem~\ref{thm:mixedplane578} constructs an even-order Hilbert cuspidal plane with opposite ordinary quotients at the split coefficient prime. Its quadratic induction is cuspidal and remains irreducible on the auxiliary cyclotomic kernel by Corollary~\ref{cor:parallelind578}. Proposition~\ref{prop:weightchange579} changes one Hilbert weight at the same tame level. Theorem~\ref{thm:regularind579} and Corollary~\ref{cor:bothinitial579} give regular cuspidal initial points with the prescribed residual matrices and flags in both character cases. At both initial points Proposition~\ref{prop:seedadjoints580} gives vanishing of the two finite adjoint Selmer groups. Theorem~\ref{thm:seedweight580} identifies the unframed ordinary deformation ring in characteristic zero with its formal weight space. In the even-order case Theorem~\ref{thm:quadraticfamily583} identifies the entire ordinary family at the initial point with the quadratic inductions of the ordinary plane over the specified field. Corollary~\ref{cor:quadraticcomponent583} gives the resulting trace identities on every component through that point. Theorem~\ref{thm:cubeweights584} gives the equations for the symmetric cubes at the odd-order point. Theorem~\ref{thm:cubedepth584} computes their ordered weight obstruction on the original period fields, and Corollary~\ref{cor:cubecutoff584} compares its exact depth with the torsion cutoff. Theorem~\ref{thm:rationalinduction585} excludes the even-order quadratic trace identities at all characteristic-zero rational fixed-weight points. Corollary~\ref{cor:integralinduction585} gives a torsion intersection and proves finite length on the two-plane locus and at all even orders greater than two. Theorem~\ref{thm:allinduction586} proves finite length of this intersection for every even order by proving that the local ordinary line of the remaining inducing plane is globally invariant. Theorem~\ref{thm:quadraticexponent587} computes the exact power of the coefficient prime annihilating the full integral intersection. Corollary~\ref{cor:quadraticcutoff587} places its cohomological images below the cutoff on the original period fields. Theorem~\ref{thm:unorderedcube588} bounds the precision of symmetric-cube characteristic polynomials on the period fields using a single inertia element, without an ordering of the eigenvalues. Corollary~\ref{cor:unorderedcutoff588} gives the corresponding uniform length estimate. Theorem~\ref{thm:polarweights593} determines the full polarized ordinary deformation rings at both cuspidal starting points and the equations for symplectic descent. Corollary~\ref{cor:polarsplit593} gives coordinates from Frobenius traces at rational primes splitting in all the period fields. Proposition~\ref{prop:tamerigid594} computes the ramified local cohomology at the cuspidal points to every precision. Theorem~\ref{thm:commonprobes594} obtains their Frobenius coordinates from one progression on the interval whose primes preserve the strict cohomology on every original field for the indicated weights. Theorem~\ref{thm:heckeseeds595} identifies the ordinary Hecke algebra and the polarized deformation ring after completion at either cuspidal point. Corollary~\ref{cor:seedcomponents595} gives the Hecke relations on every marked component through that point. Proposition~\ref{prop:horizontal596} gives an integral Hecke component reaching the prescribed weight. Theorem~\ref{thm:classicalfiber596} realizes every characteristic-zero eigensystem in that Hecke fiber by a classical automorphic system over the same auxiliary CM field after restoring the multiplier. Theorem~\ref{thm:regularcosets597} proves density of simple Frobenius spectra on the cyclotomic kernels of the irreducible ordinary curves with separated weights in characteristic~\(\ell\). Corollary~\ref{cor:commoncurves597} fixes their discriminant valuations at common split primes for every cyclotomic depth. Theorem~\ref{thm:dualcurve598} bounds the finite-coefficient dual classes passing the Frobenius tests uniformly in both precisions. Corollary~\ref{cor:ordinarycurve598} bounds the ordinary Selmer lengths on the irreducible curves. Proposition~\ref{prop:splitexception598} identifies the quadratic case requiring extra primes inert in the inducing field. Theorem~\ref{thm:splitprojectors599} detects the full dual cohomology, including this quadratic case, at primes split in every period field by testing one matrix coefficient at each assigned prime. Theorem~\ref{thm:exactselmer614} prescribes the exact cyclotomic depths at these same primes. Corollary~\ref{cor:selmercongruence614} computes the integral cost of isolating the unramified tuple at commuting Frobenius specializations. Corollary~\ref{cor:rankonedual599} gives dual Selmer vanishing with these local conditions. Proposition~\ref{prop:rankonelift599} computes their tame lifting obstruction. For the original fields, Theorem~\ref{thm:normalizedprimes600} recovers every strict class and its exact order from normalized Frobenius values. Corollary~\ref{cor:originalinterval600} bounds the number of common test primes polynomially on the interval for good coefficient primes of polynomial size. The orders of the coordinate images in that corollary remain to be bounded. Theorem~\ref{thm:genuskernel601} identifies the kernel of restriction to the absolute commutator subgroup at every coefficient prime and bounds its full weighted contribution by \((t(x)-1)\log g_\psi\). By Corollary~\ref{cor:allprimesstrict601} a common list of split primes detects the quotient also at ramified primes and at two. Theorem~\ref{thm:idealprecision608} computes the strict subgroup at arbitrary precision from ideal classes and local units on one cyclotomic layer for all period fields. Proposition~\ref{prop:quadraticorders608} bounds the full order of the quadratic summand of discriminant \(-1208\) at three by nine at every odd weight. Theorem~\ref{thm:exactthree609} determines this order to be three at every positive odd weight. Corollary~\ref{cor:quadraticrange609} bounds the total contribution at three for the indicated quadratic characters with real discriminant below \(100000\).

\lead{The period estimate}Sections~\ref{nav:I-37} and~\ref{nav:I-39} give an integral denominator for weight projections on definite lattice theta coefficients and a seesaw formula for a unitary lift. For an iterated chain Theorem~\ref{thm:center} computes the output central character. Corollary~\ref{cor:orth} gives an explicit large family of mutually orthogonal images on the same fields \(E_x\). The position of the final projection is decisive. If \(A_i\) has image in one central character and \(P_\nu\) denotes the orthogonal projection to another, then \(P_\nu A_i=0\). In particular, on a source of dimension \(r_i>0\),
\[
          \|(P_\nu A_i)^*(P_\nu A_i)-I\|_{\rm HS}^2=r_i .
\]
This follows from orthogonality of distinct characters and \(\operatorname{tr}(I)=r_i\). Thus a family that separates into many central characters before projection does not on that account satisfy the near-isometry estimate for nonzero sources in a prescribed common Bessel component. For the first Bessel position on the prescribed terminal graph rows, Corollary~\ref{cor:rankobstruction337} gives this exact loss on the separately zero-sum columns counted in Corollary~\ref{cor:rank111}. The rank on the full refined table is at most its multiplicity divided by \((T-1)^f\). Theorem~\ref{thm:terminalwalk338} compares all the terminal exchanges on these columns after a specified translation, retaining the original graph at the symmetric endpoint. Theorem~\ref{thm:markedglobal339} includes the two auxiliary primes on a specified common translation there. Theorem~\ref{thm:thetacells340} transports every theta test to a common integral subdivision. Theorem~\ref{thm:pintable340} bounds the rank on pairing the transported fine first test space back with its original space. Corollary~\ref{cor:markedtests341} includes both auxiliary primes at the same time and gives the coefficient ideal on the marked tensor of tests. Theorem~\ref{thm:signedreturn342} rules out returning a matched translate to its original signed projector. Theorem~\ref{thm:rationalperiod342} instead performs all exchanges on the original auxiliary rows for a specified rational conjugate of the Bessel character. Corollary~\ref{cor:avoidrange343} chooses its unmarked support above the cohomological cutoff. The common denominator for the additional level projection is given in Theorem~\ref{thm:unmarkedtrace343}. Theorem~\ref{thm:flagvanish344} proves vanishing on the original spherical rows for the choice close to the identity at infinity with good split support. Theorem~\ref{thm:dominantflag344} and Proposition~\ref{prop:dominantexchange344} retain the exchange using integral unmarked root ratios and an explicit archimedean displacement. Theorem~\ref{thm:displacedreal344} bounds the actual analytic integral at that displacement. Theorem~\ref{thm:terminalformula347} evaluates the last restriction in the arithmetic class sum. Corollary~\ref{cor:normaldegree347} gives its exact normal-line factor on these fields. Theorem~\ref{thm:wordlower348} bounds the full finite and archimedean variation of this last normalization across the selected words by one integer of polynomial logarithm. Theorem~\ref{thm:abslow349} bounds the absolute scalar of the last restriction in the highest-vector metrics, and Proposition~\ref{prop:normalframe349} identifies its integral coefficient maps even above the terminal prime. At the base tuple and matching marked levels, Theorems~\ref{thm:allclasses350} and~\ref{thm:augmented350} determine the last restriction on the entire arithmetic class table, including its norm on columns with separate scalar sums zero. Theorem~\ref{thm:weylperiod351} bounds the integral loss of this last step at all primes on the actual projected test lattice, with one integer for all words and translates at the given field. Theorems~\ref{thm:covectorpaths352} and~\ref{thm:pathnorm352} express the finite signed theta column on the scalar zero-sum input in terms of the reductions of the composite matrix paths. Proposition~\ref{prop:shortpaths352} separates distinct rational directions below a stated height at the chosen terminal prime. Theorem~\ref{thm:vectorchain353} carries the arithmetic recurrence to vector markings by finite Fourier inversion. Theorem~\ref{thm:pointlattice353} and Corollary~\ref{cor:isolatedcost353} compute the actual coefficient ideals for an individual vector value and for a prescribed isolated difference. Theorem~\ref{thm:grid354} gives simultaneous integral interpolation on prescribed residue grids. Corollary~\ref{cor:fouriermean354} computes the contribution of the full initial table to the mean determinant at the terminal prime. Theorem~\ref{thm:rankone354} expresses the Gram operator on the actual signed arithmetic column as a positive sum over matrix frequencies of rank one. Theorem~\ref{thm:imagedet355} and Corollary~\ref{cor:jointimage355} bound the contribution of the integral input at the terminal primes to the determinant of the actual joint image, even when some periods vanish. The estimate is relative to the same signed maps on the unrefined vector lattices. Theorem~\ref{thm:transfer380} retains the row stabilizers when the last restriction forgets auxiliary frames. Theorem~\ref{thm:fiberlength380} computes the resulting multiplicities, and Proposition~\ref{prop:freeflags380} estimates them from below using the residue fields of those same frames. Corollary~\ref{cor:rowhaar380} gives the character factor in a Haar integral for the specified class sum. Theorem~\ref{thm:planehaar383} computes the signed terminal column itself at the extra ramified places and gives its Haar pairing in the original row space. Theorem~\ref{thm:flagproject383} and Corollary~\ref{cor:flagcost383} separate its flag translates and bound the integral cost in terms of the same conductor. Theorem~\ref{thm:markedgauss386} evaluates the full local Fourier integral of the last two edges at each marked place, including the base tuple. Corollary~\ref{cor:coupledglobal386} identifies its frequency support from the preceding matrix and retains its phase and denominator in the signed column on the same fields. Theorem~\ref{thm:markedtablegram387} computes the inner products of the marked Fourier tables, including paths with the same supporting hyperplane. Theorem~\ref{thm:markedtableintegral387} gives a sharp integral lower bound for a nonzero group of signed paths averaged over its full marked frequency table and a bound at each nonzero sampled coefficient. Proposition~\ref{prop:markedsampling387} proves that the terminal matrix index samples every marked residue covector equally often before the theta derivatives are applied. Proposition~\ref{prop:normalcells388} realizes the terminal residue cells as a complete family of normal theta sections. Theorem~\ref{thm:normaljet388} identifies the precise normal derivatives in the projected operators, and Corollary~\ref{cor:normaltranslate388} gives their signed sum in integral combinations of torsion translates. Theorem~\ref{thm:normalframe389} bounds the full matrix of normal derivatives on the common interval. Corollary~\ref{cor:weightedkernel389} transfers this estimate to the projected terminal operators and identifies the component of a signed coefficient vector that contributes to their sum. Theorem~\ref{thm:terminalmarked390} uses the larger prime at the endpoint to retain the signed sum at both auxiliary primes. Corollary~\ref{cor:signedterminal390} bounds its mean square over the nonzero intermediate covectors at the endpoint. Propositions~\ref{prop:incidencenorm391} and~\ref{prop:coherentcell391} impose the preceding covector equations for a fixed path. Theorem~\ref{thm:splitcoherent391} estimates the coefficients after summing every terminal covector with the same initial pullback when the image at the terminal places comes from a common orthogonal summand. Theorem~\ref{thm:intermediateimage392} proves a near-isometry after the terminal covectors with a common intermediate pullback are summed, outside an explicitly bounded set of matrix residues. Proposition~\ref{prop:initialrank392} gives the rank obstruction at the initial vertex even on the Fourier support of a varying marked path. Theorem~\ref{thm:markedimage393} determines how the actual signed tables with a fixed penultimate image vary inside that support. Corollary~\ref{cor:lastblock393} sums the redundant terminal coordinates and gives a single prescribed test for the last varying determinant block. Theorem~\ref{thm:firstvertex394} estimates the signed contraction at the first intermediate vertex for marked tables supported on the specified affine subspaces. Corollary~\ref{cor:markedpaths394} applies the estimate to every individual marked test and to common sums over later determinant blocks. Corollary~\ref{cor:firstaverage395} measures the loss at the initial covectors when the first matrix varies uniformly. Theorem~\ref{thm:firstslice395} expresses the actual weighted first-edge sum in terms of translates of its theta distribution at rank-one matrix frequencies. Theorems~\ref{thm:initialindex396} and~\ref{thm:initialmoment396} compute the integral index of initial character tests and their values on the full signed column to every prescribed precision at the terminal prime. Theorem~\ref{thm:periodminors400} gives the terminal valuations for minors of signed character tests through finite moments of the actual composite covectors. Corollary~\ref{cor:fullperiodmoments400} computes the determinantal ideals on the full integral initial table with separate scalar sums zero. Theorem~\ref{thm:signedcharacterloss432} also compares the actual projected minors on that table with their values on pure scalar characters at every finite place. The total loss per minor dimension is at most \(f\log(T-1)\), and the archimedean Gram operator is the sum of the character Gram operators with the induced norms. Theorem~\ref{thm:commonphase441} computes the row normalization in the two analytic evaluations of the prescribed periodized columns. A single phase works for every word and disappears from every mixed Gram entry by Corollary~\ref{cor:mixedphase441}. Theorem~\ref{thm:frameperiods446} computes the period ideals on the image of the integral tests under the finite projector from actual lowered class coefficients at all three marked primes. Rational frames on the original Hermitian spaces realize the least coefficient valuation in each exterior degree. Proposition~\ref{prop:framegeneration446} also gives a kernel description for the saturated projected tests. Theorem~\ref{thm:commoncompactgram447} expresses each definite Gram determinant as a positive sum of the squared absolute values of actual projected period minors. The same rational weights work at all complex embeddings and are units at the three marked primes. Corollary~\ref{cor:threeminor447} gives an archimedean bound for one minor attaining the three selected local minima simultaneously. Theorem~\ref{thm:globalperiodcontent448} bounds the full ideal norms of the actual definite periods by their complex Gram determinants at all embeddings. Corollary~\ref{cor:saturatedcontent448} gives this comparison on the image test lattice with zero additive cost and on the saturated projected lattice with cost at most the logarithm of the marked group order per exterior degree. Theorem~\ref{thm:weylrowjet449} proves divisibility of each definite theta operator on integral Weyl lattices at the good primes interchanging the two halves of the CM type. The loss at each odd Frobenius translate is at most the length of the first row of its partner partition. Theorem~\ref{thm:signedsignal450} carries the actual word count through the signed class maps and both projected test lattices at the good signal primes. Corollary~\ref{cor:gramsignal450} compares it with the Gram determinants of those very columns, retaining the scalar and outside-prime costs. Theorem~\ref{thm:centerwords451} evaluates the complementary quadratic sign at the actual word weights. Corollary~\ref{cor:centerheight451} bounds the centered count by the definite Gram determinant minus its explicit signal main term, together with costs computed from the finite input data. Lemma~\ref{lem:signalsupport452} and Theorem~\ref{thm:fullsignal452} include every active signal prime for the prescribed finite tests. Corollary~\ref{cor:fullcenter452} removes the omitted-prime sum and computes the signal denominators for elementary columns. Theorem~\ref{thm:outsidedifferent453} computes the outside loss away from the terminal primes in terms of the different and the original split test functions. Corollary~\ref{cor:outsidecenter453} sums this estimate on the same two period ideals used for the centered count. Theorem~\ref{thm:terminaldegrees454} computes a bound at the terminal prime after removing the initial and normal-coordinate scalars. Corollary~\ref{cor:allcost454} gives the centered inequality with every finite term read directly from the fields, weights, levels and input tests. Theorem~\ref{thm:rankonecost455} uses the rank of the actual terminal frequencies to replace the total partition size by its first row. Corollary~\ref{cor:terminalrow455} gives this improvement for the original signed columns on the two projected test lattices. Theorem~\ref{thm:matrixshift456} treats arbitrary split matrix shifts using their elementary divisors. Corollary~\ref{cor:splitcost456} gives the corresponding costs of the original row tests from their finite Fourier transforms. Theorem~\ref{thm:smallcollision457} bounds the class number of the cyclic field of a colliding character. Corollary~\ref{cor:fixedcollision457} eliminates all quadratic characters above a polynomial conductor threshold at the half residue, and treats explicit congruence classes for characters of odd prime order. Theorem~\ref{thm:uniformdepth458} bounds the exponent at colliding residues with constants independent of the character order. Corollary~\ref{cor:dyadicdepth458} applies it to faithful characters of the growing real dyadic layers in the original intervals. Theorem~\ref{thm:commonorbits459} counts all faithful Frobenius orbits against one class group even when their twists differ. Corollary~\ref{cor:dyadicorbit459} leaves at most one cyclic scalar group up to conjugation at a half-residue twist on the indicated dyadic layers. Theorem~\ref{thm:polycenter460} retains the ramified contribution on the polynomial-conductor family and determines its mean over translates. Corollary~\ref{cor:shortcenter460} applies the orbit bound on precisely its dyadic layer. Theorem~\ref{thm:cycleweyl461} bounds the integral Weyl operator jointly in its rank-one directions on every even Frobenius cycle. Corollary~\ref{cor:inertcorrection461} carries the resulting extra divisibility at all unramified inert primes to the centered comparison. Theorem~\ref{thm:inertmass462} gives the lower bound \((1/4-o(1))hma\log X'\) for the extra projected-period divisor. Corollary~\ref{cor:dyadicchoice462} improves its coefficient by choosing the permitted high cosets within the same fields. Theorem~\ref{thm:primitivegram463} gives the centered bound after dividing the actual columns by explicit rational integers. Corollary~\ref{cor:heightfloor463} determines a necessary leading lower bound for the original Gram term together with its finite costs. Theorem~\ref{thm:allplacesminor477} compares one marked determinant with the archimedean Gram determinants of the same columns. Corollary~\ref{cor:markedheight477} gives a comparison with error independent of the coefficient-field degree using a finite family of such determinants. Theorem~\ref{thm:signedcovariance478} computes the translation averages of the full signed columns from their finite coefficients. Corollary~\ref{cor:signedspan478} determines exactly which lists can be made independent by allowing separate torsion translations. Theorem~\ref{thm:fixedexterior479} bounds the Gram determinant at the prescribed tests in terms of finite coefficient vectors and the reference vacuum norms. Corollary~\ref{cor:fixedheight479} inserts this bound into the centered comparison on the same interval. Theorem~\ref{thm:splitjets480} computes every reduced elliptic theta derivative from orders below the split prime. Corollary~\ref{cor:splitrank480} applies this to minors at the prescribed tests, with the coupling and signs in the original column. Theorem~\ref{thm:shiftedjets481} determines the rank of the reduced elliptic evaluations at a list of shifts. Corollary~\ref{cor:productjets481} gives independent control of their low derivatives on an entire product edge when additional input translates are allowed. Theorem~\ref{thm:positivejets482} and Corollary~\ref{cor:fixedproductjets482} recover the coefficient span of a fixed signed elliptic product using positive derivative orders at one prescribed tuple of shifts. Corollary~\ref{cor:fixedhankel482} gives the rank of its one-coordinate derivative matrices in terms of a finite set of nonzero roots. Theorem~\ref{thm:commonsplit483} finds common split primes with long Frobenius orbits. Corollary~\ref{cor:projectsymbol483} proves that a nonzero bounded-degree symbol on any fixed nonzero torsion array survives among its positive mixed derivatives, even with arbitrary correlations among the shifts. Theorem~\ref{thm:jetsmith485} computes the exact integral cost of prescribing jets at the original torsion shifts using a common input level. Corollary~\ref{cor:projectedminors485} gives the determinantal ideals after a specified signed linear projection on those jets. Theorem~\ref{thm:thetaalternant486} and Corollary~\ref{cor:confluent486} evaluate the archimedean determinant of that same elliptic table with its normalization. Corollary~\ref{cor:fixedminor486} evaluates minors with prescribed signed rows and input columns by derivatives of a single theta product. Theorem~\ref{thm:jetnorm487} computes the absolute rational norm of the complete marked elliptic interpolation determinant from the CM annihilator ideals. Corollary~\ref{cor:localjetnorm487} computes its mean content at each rational prime, including primes dividing the level or the shift orders. Theorem~\ref{thm:signedcomplete488} bounds signed sums of complete elliptic blocks with a fixed multiplicity partition, uniformly in the shifts. Corollary~\ref{cor:arbitraryorders488} treats specified derivative orders on the same complete input basis. Corollary~\ref{cor:commoncapacity489} bounds the archimedean contribution of the original projected columns using their signed coefficient matrices and the exact CM area contribution of the total degree. Corollary~\ref{cor:orthometric491} gives this contribution in any chosen test metric. The optimum recovers the signed coefficient mixed determinant within one unit in its logarithm per column. Theorem~\ref{thm:residuejet492} determines the valuation of each prescribed derivative of an ordinary translate whose full primary order survives reduction. Corollary~\ref{cor:residuecolumn492} computes the leading residue of the signed columns at the same test. Theorem~\ref{thm:divisionphase493} gives each translated germ on one explicit division field, including its phase. Theorem~\ref{thm:fixedsectors493} then separates the character components of the original period determinant under field automorphisms. Theorem~\ref{thm:labelder494} identifies a character by the product of the two split coordinates of a shift. Corollary~\ref{cor:indepder494} proves independence for a family of prescribed derivative orders and identifies the trivial character on a pure terminal frequency. Theorem~\ref{thm:cyclicunit495} sums a complete cyclic terminal Fourier projection. At a place where its frequency subgroup is connected, every derivative for a nonzero character is a unit in the original differential. Corollary~\ref{cor:centeredterm496} removes half the terminal Fourier index from the finite-place bound for the prescribed signed periods, retaining the first-row derivative cost and the original fields and intervals. Corollary~\ref{cor:terminalheight497} cancels the remaining explicit terminal Fourier index against the squared norms of the same complete projections at the complex embeddings. Its coefficient matrices retain the signed covector sums. Theorem~\ref{thm:puncturedrank498} computes the residue rank of a terminal derivative table on its nonzero covectors. Corollary~\ref{cor:homogeneousrank498} gives the ranks at every homogeneous order, before the intermediate covectors are summed. Theorem~\ref{thm:spectralquot499} cuts the exceptional trace algebras to finite modules at common split primes while preserving any prescribed finite integral spectra. Corollary~\ref{cor:strictvertical499} bounds the strict lengths by a torsion exponent and a determinant on the resulting characteristic-zero factors. Theorem~\ref{thm:weylmoments500} computes the reduction after the full sum over a shared covector at one connected block. Its moments use exponents bounded by the first row of the prescribed Weyl shape. Corollary~\ref{cor:orbitkernel500} determines the loss on a scalar orbit for the unshifted moment map. Proposition~\ref{prop:orbitsum501} sums the scalar covectors through the whole signed chain. Theorem~\ref{thm:augterminal501} proves a valuation gain at every edge for a separately zero-sum input. Corollary~\ref{cor:augheight501} carries it to the upper bound for the divided projected columns. Lemma~\ref{lem:primitiveclass502} computes a cyclotomic specialization from an ordinary ideal class group, including characters trivial at the coefficient prime. Theorem~\ref{thm:splitweights502} and Corollary~\ref{cor:fullcollision502} bound the full lengths away from a sparse set on the common interval. Theorem~\ref{thm:classfrob503} realizes the visible part of the exceptional dyadic ideal classes at a common prime testing the elliptic periods. Corollary~\ref{cor:classmass503} bounds the total loss in their weighted lengths by the logarithm of the degree of the test field over its maximal abelian subfield. Theorem~\ref{thm:meangauss504} gives a uniform saving per terminal block for every pure nontrivial scalar character, averaged over coefficient places. Corollary~\ref{cor:packetsave504} gives a saving for a prescribed sum in terms of its scalar support. Theorem~\ref{thm:unitlength505} computes the unit norm index and capitulation kernel between cyclotomic layers at a surviving dyadic orbit. Corollary~\ref{cor:unitcap505} gives a test with ordinary units modulo the coefficient prime that bounds every weight of the orbit when the tested norm map vanishes. Theorem~\ref{thm:splitamb506} computes the invariant class group lengths from units of the residual field. Corollary~\ref{cor:zeromass506} bounds all nonsplit zero-residue exceptions by one class number product, and Proposition~\ref{prop:pellunit506} proves full vanishing at seven in the real quadratic dyadic layer. Theorem~\ref{thm:sunitbase507} computes the constant specialization from units supported at the coefficient prime. Corollary~\ref{cor:intervalroots523} computes the number of generators of the complementary group at every positive multiple of the residue period, outside one integer of polynomial logarithmic size. Theorem~\ref{thm:unitroots523} also gives the complete scalar length when the two unit logarithms determine it within the available precision. Theorem~\ref{thm:octicrange524} proves full strict vanishing at positive multiples of the residue period in the first three real dyadic layers for coefficient primes from five through ten million. For primes outside the two indicated classes modulo sixty-four it applies to every dyadic layer. Theorem~\ref{thm:dyadicall525} treats every dyadic layer at all odd primes through ten million and isolates two possibly nonzero characters at the specified weights. Corollary~\ref{cor:actualdyadic525} removes their restrictions from the residual sum on the common period intervals. Theorem~\ref{thm:halfdyadic529} gives vanishing at the half-residue weights throughout a growing range of dyadic layers by comparing relative regulators with a joint mean of Dirichlet values. Corollary~\ref{cor:periodhalf529} applies it on the same fields and intervals as the signed periods. Corollary~\ref{cor:allweightlayer529} removes all weights on one common dyadic subfield of growing degree for the stated range of coefficient primes. At primes three modulo eight, Corollary~\ref{cor:fullreal529} removes every weight on the full real dyadic part in the same range. Corollary~\ref{cor:quadraticrange507} proves vanishing at every positive zero-residue weight for the real quadratic dyadic character and all coefficient primes from five through ten million. Theorem~\ref{thm:classgauss508} evaluates the mean terminal Gauss valuation on the actual CM half in terms of the auxiliary imaginary quadratic class number. Corollary~\ref{cor:cmcomp508} gives the resulting uniform saving on the common period intervals. Theorem~\ref{thm:framepair509} computes the signed Fischer pairing of class paths in common product coordinates. Corollary~\ref{cor:frameheight509} bounds the prescribed period determinant in terms of explicit factorials and the norms of its initial forms, finite transports and endpoint tests. Theorem~\ref{thm:wordgcd510} constructs a common divisor for all the signed columns at each translate and evaluates its logarithm from the Frobenius cycles. Corollary~\ref{cor:commonminors510} controls its effect on every projected minor and Gram determinant on the same intervals. Theorem~\ref{thm:divtail511} evaluates the large-prime part of the common divisor from Frobenius cycles, explicit real integrals and logarithmic derivatives at one. The remaining prime evaluations run only up to a fixed power of the interval variable. Theorem~\ref{thm:factorcomp512} compares the compact factorials of every word with the incoming real hook products on the same interval. Lemma~\ref{lem:frameend512} identifies the normalized terminal restriction in the signed coefficient matrix. Theorem~\ref{thm:divconstant513} computes the leading constant of the common integral divisor from the order of the original Dirichlet character. Corollary~\ref{cor:signedfactor513} inserts it in the determinant of the prescribed divided columns. Theorem~\ref{thm:connectedjets516} determines the contribution of the connected terminal Fourier coordinates modulo every power of the marked prime. Corollary~\ref{cor:signedprecision516} gives the truncated elementary divisors after the original signed contraction. Theorem~\ref{thm:puncturedisometry518} gives a two-sided complex norm estimate at the nonzero terminal covectors for each fixed translated product. Corollary~\ref{cor:terminalgram518} bounds its Gram determinant and its pairing with the projected incoming table. Theorem~\ref{thm:scalaraugmentation519} preserves this rank after subtracting the mean on each actual nonzero scalar orbit. Corollary~\ref{cor:multipleaugmentations519} gives the same comparison after the separate projections at all terminal blocks of one product term. Theorem~\ref{thm:signedtranslation520} controls signed sums of distinct torsion translates at the prescribed point. Corollary~\ref{cor:flagshifts520} gives a rank and Gram determinant comparison within each coset of torsion supported at the two auxiliary primes with bounded powers. Corollary~\ref{cor:actualsource520} computes the first contraction on the full refined initial class in its automorphic norm. Theorem~\ref{thm:twopencils521} evaluates the inner product of two prescribed matrix pullbacks. Proposition~\ref{prop:markedpencil521} counts the terminal markings at which their actual input lines can coincide. Theorem~\ref{thm:terminalreps522} bounds the height of representatives for every tuple at the prescribed marking. Proposition~\ref{prop:lineimage522} computes the integral image and the complex norm from the signed paths on those same lines. Theorem~\ref{thm:artiniansplit490} gives finite quotients bounding all exceptional strict lengths at at most three common split primes on the period interval, also for positive-dimensional original trace algebras. Corollary~\ref{cor:jetslength490} gives an explicit length bound from the degrees of suitable initial forms of their relations. Theorem~\ref{thm:collisionbranch464} describes the possible large lengths of the faithful dyadic scalar groups by one cyclotomic parameter. Corollary~\ref{cor:onedeepweight464} leaves at most one weight per indicated prime and dyadic level on the entire common interval. Theorem~\ref{thm:simultaneousweights465} concentrates the residual cohomology simultaneously at all primes up to the cutoff on a sparse set of weights. Corollary~\ref{cor:supportmass465} gives the weighted sum still needed on that set. Theorem~\ref{thm:nonsplitprefix466} improves the exceptional-weight count for the characters nontrivial at the coefficient prime. Corollary~\ref{cor:fullnonsplit466} bounds their full orders away from those weights on the period fields. Theorem~\ref{thm:scalarlevel467} detects a length beyond any prescribed depth using one finite class group and its Galois action. Corollary~\ref{cor:leveltwentyfive467} isolates the only possible length above twenty-five for the faithful dyadic half-residue weights in the stated prime window. Theorem~\ref{thm:sympordinary468} constructs the integral local symplectic ring in six coordinates. Corollary~\ref{cor:sympcomponent468} identifies its flag condition on the component through the split representation. Theorem~\ref{thm:symprelations469} counts the global symplectic generators and relations for the surviving dyadic orbit. Corollary~\ref{cor:firsttorus469} gives its presentation after taking the first torus invariants. Theorem~\ref{thm:exactstricttrace470} identifies the full strict length with an integral congruence quotient. Corollary~\ref{cor:frobideal470} generates its augmentation ideal by at most three good-prime trace differences. Theorem~\ref{thm:congruenceorder471} describes the torsion and integral congruences when the trace algebra is finite. Corollary~\ref{cor:strictdet471} then bounds the exceptional length by a determinant of actual trace differences. Theorem~\ref{thm:unitray472} detects the opposite ordinary direction using global units modulo the coefficient prime. Corollary~\ref{cor:eliminateupper472} removes one trace parameter whenever this unit projection is nonzero. Theorem~\ref{thm:rankdichotomy473} realizes the nonzero trace coordinates by an irreducible rank-four representation or a pair of even planes. Corollary~\ref{cor:finitecurve473} shows exactly how a positive-dimensional special fiber gives such a representation in characteristic equal to the coefficient prime. Theorem~\ref{thm:strongtrace474} proves that irreducibility persists after restriction to each period field. Corollary~\ref{cor:twoplane474} identifies the two planes in a reducible trace as a single even plane and its prescribed twist. Theorem~\ref{thm:splitparameters475} chooses at most three parameter primes that split in all the period fields of the interval. Corollary~\ref{cor:splitdet475} separates all the finite generic spectra at one such split prime. Theorem~\ref{thm:monodromy476} determines the connected algebraic monodromy at the trace witnesses. Corollary~\ref{cor:cyclomonodromy476} verifies the distinct-eigenvalue projector test on the cyclotomic kernel over each period field. Theorem~\ref{thm:jointprimes484} chooses the trace parameters with the split-prime theta estimates on the same fields. Corollary~\ref{cor:strictmindet484} bounds the strict length of a finite exceptional trace algebra by its least determinant valuation and computes that valuation branch by branch. Theorem~\ref{thm:neighborindex356} computes the index of the additional Hecke modification at the terminal prime and bounds its effect on the projected image on the same fields. Theorems~\ref{thm:terminalscalar357} and~\ref{thm:normalizedimage357} compute the scalar normalization at a common terminal prime and its exact contribution to the determinant after signed projection. A lower bound for their signed Gram sum and a comparison with the analytic incoming row are still required. Proposition~\ref{prop:four} controls the content of a finite Fourier transform when the determinant fibers of the relevant period are rational on the test lattices.

One can also put all the determinant twists into one output central character from the outset. For the ranks \(n_i=i+1\) and chain length \(m\), Theorem~\ref{thm:gram25}, writes the two second moments required in Proposition~\ref{prop:rank11} as correlations of determinant fibers after the prescribed output map. There are \(|C(P_0)|^{m-2}\) twists in a given central character, where \(P_0\) is a fixed unramified auxiliary prime in Volume~I. Its Proposition~\ref{prop:level25} shows that if both neighboring kernels at each intermediate group are invariant under the integral determinant units at \(P_0\), the proportion of nonzero maps on the common fields is at most \(\exp(-c_{P_0}x^{5/2+8\eps})\). Zero maps alone then violate the first error bound in Proposition~\ref{prop:rank11} for that full family. The computation identifies a restriction on the choice of levels. The common ramified vector in Theorem~\ref{thm:bessel29}, gives vanishing cross terms after restriction. Its individual norms are treated below.

To finish the geometric argument one must identify the finite-place divisibility of the specified Bessel periods with the weighted count \(\Delta_u\), and bound their norms, ranks and covolumes after projection in that same target. The required result is \eqref{eq:gap} simultaneously for all translates and all \(x\) in a selected interval. The following observation measures the strength of that estimate for the explicit counts.

\begin{proposition}[The strength of the remaining count estimate]\leavevmode\label{prop:strength26}
Fix a primitive \(\psi\). Use the explicit data of Proposition~\ref{prop:field12} on the common intervals of Volume~IV, with \(\Pi\) fixed as above. The following two assertions are equivalent.
\begin{enumerate}
\item Every zero in the open strip of every \(L(s,\chi)\) with \(\chi\in\mathcal C_\psi\) has real part \(1/2\).
\item For every sufficiently small \(\eps>0\), there are arbitrarily distant such intervals on which \eqref{eq:gap} holds for all \(u\) at every \(x\).
\end{enumerate}
At conductor one the first assertion is the Riemann hypothesis for \(\zeta(s)\).
\end{proposition}
\begin{proof}
Under the critical-line assertion, Theorem~\ref{thm:explicit} gives \(G_\chi(x)=O_\chi(\sqrt{x})\) for each of the finitely many retained characters. Summing their fixed weights gives the same bound for \(\Gg_\psi\). By \eqref{eq:actual12} the explicit count then satisfies \(\Delta_u\ll_{\psi,\eps} Q x^{1/2+\eps/2}\), uniformly in \(u\) wherever the data are chosen. Their parameters have \(a_E=xb\), \(m=\lfloor\sqrt{x}\rfloor\), \(Q=hmb\), and \(\mathcal B\ge h a_E\). Thus \(\mathcal B/Q\ge x/m\ge\sqrt{x}\). Since \(5\eps>\eps/2\), inequality \eqref{eq:gap} holds everywhere on every sufficiently distant common interval.

Conversely the first three conditions in Proposition~\ref{prop:balance} hold for the explicit data. The assumed fourth condition gives \(\Gg_\psi(x)\le x^{1/2+9\eps}\) throughout arbitrarily distant power intervals. Apply Theorem~\ref{thm:detect} with exponent \(1/2+9\eps<1\), and then let \(\eps\) decrease to zero. Every open-strip zero of every member of the family has real part at most \(1/2\). The primitive functional equation reflects any zero of smaller real part to one of larger real part for the conjugate character, which is also in this family. This proves the converse. At conductor one the retained family consists only of the trivial character.
\end{proof}

\lead{Which global lengths remain}Fix \(x\) on the interval and \(E=E_x\). For each odd prime \(\ell\) set \(h_\ell=b_0\) if \(\ell\le P\), and \(h_\ell=0\) otherwise. Use the product decomposition \(\operatorname{Gal}(E/\QQ)=P_\ell\times H_\ell\) into its \(\ell\)-Sylow and prime-to-\(\ell\) factors. For a character \(\alpha\) of \(H_\ell\), take its extension trivial on \(P_\ell\), and restrict \(\vartheta=\kappa_\ell^{-r}\alpha^{-1}\) to the absolute Galois group of \(E^{H_\ell}\). In the following sums use only the complementary sign \(\alpha(c)=(-1)^r\). The unramified coefficient ring \(\mathcal O\) may depend on \(E,\ell\) and contains all values of the characters of \(H_\ell\).

We retain triples \((\ell,r,\alpha)\) with \(1\le r<t\) subject to three conditions. First, the field average \(J_{\ell,r}\) above acts by zero on this component. A nonzero field average acts by one exactly when \(\alpha\) is trivial on the subgroup being averaged, since that subgroup has order a power of two and lies in \(H_\ell\). Second, \(\bar\vartheta\ne1\) on the absolute Galois group of \(E^{H_\ell}\). Third, \(v_\ell(r+1)<e_{\ell,r}(E)\), where \(e_{\ell,r}(E)\) is the exponent index of the finite Kummer group in Theorem~\ref{thm:compare30}. It can be read from ordinary ideal classes with their local generator conditions. The example of Proposition~\ref{prop:example30}, shows why positive lengths genuinely occur in these families.

For a retained triple let \(B\subset\mathcal O^n\) be the integral relation module of Theorem~\ref{thm:present31}. For \(0\le j\le n\) let \(D_j\) be the minimum valuation of the \(j\)-rowed minors of a relation matrix, with \(D_0=0\). These are finite because \(B\) has rank \(n\). Put
\begin{equation}
 \mathcal R_*(x)=\sum_{\mathrm{retained}\ (\ell,r,\alpha)}(\log\ell)
       \max_{0\le j\le n}\{D_n-D_j-(n-j)h_\ell\}.              \label{eq:minors33}
\end{equation}

\begin{proposition}[The remaining lengths in terms of global relations]\label{prop:minors33}
On all sufficiently distant common intervals chosen above,
\begin{equation}
 0\le\mathcal L(x)-\mathcal R_*(x)
 \le\tfrac12\exp(x^\eps/2)+3\exp(x^{2\eps/3})+2\exp(x^\eps/3).
                                                               \label{eq:minerr33}
\end{equation}
In particular \(\mathcal R_*(x)\le\tfrac14\exp(x^\eps/2)\) at every \(x\) on arbitrarily distant such intervals would suffice for \eqref{eq:length12}. For \(\ell>P\) the finite local conditions already recover every finite layer of the strict group at its original exponent. At \(\ell\le P\) one can recover the layer at \(\ell^k\) using finite local conditions at \(\ell^{k+s_*}\) with \(s_*\ll_{\psi,\eps}\log(2x)\).
\end{proposition}
\begin{proof}
Extending scalars from \(\ZZ_\ell\) to the unramified \(\mathcal O\) does not change the normalized module length, by filtering into factors of length one. The idempotents of \(H_\ell\) decompose the resulting module as a direct sum; positive-degree cohomology of this prime-to-\(\ell\) group vanishes by averaging. Lemma~\ref{lem:des28}, identifies the summands with the indicated strict groups over \(E^{H_\ell}\), also at primes dividing the order of the original Dirichlet character. The groups have finite total prime support for each \(x\) by Proposition~\ref{prop:finite15}.

By Theorem~\ref{thm:present31}, the elementary divisors of the retained group are those of \(\mathcal O^n/B\). Write their valuation exponents \(0\le b_1\le\cdots\le b_n\), including zeros. Smith normal form gives \(D_j=\sum_{i\le j}b_i\): integral invertible row and column operations preserve the ideals of minors, and the statement follows for the diagonal matrix. Removal of the subgroup killed by \(\ell^{h_\ell}\) leaves length \(\sum_i(b_i-h_\ell)_+\). This equals the maximum in \eqref{eq:minors33}, by including in the tail exactly the exponents greater than \(h_\ell\). Thus the summands are actual lengths and nonnegative. A determinant in Theorem~\ref{thm:norm28}, bounding the full exponent by \(\ell^{h_\ell}\) makes its summand zero.

Proposition~\ref{prop:res23} bounds the removal of the dyadic groups, the critical sign, the field averages and the indicated finite truncations. Among what remains, the globally trivial residual characters contribute at most \(\exp(x^{2\eps/3})\) by Corollary~\ref{cor:res28}. Its hypothesis uses our choice of \(\Pi\). All pairs removed by the third condition contribute in total at most \(\exp(x^{\eps/3})\) to the full complementary length, by Corollary~\ref{cor:family30}. Removing already-counted subgroups or truncated summands only decreases these bounds. Addition proves \eqref{eq:minerr33}. The two latter exponential orders are \(o(\exp(x^\eps/2))\), which gives the stated sufficient condition.

For each retained component let \(s_*\) be the largest local loss of Theorem~\ref{thm:finite31}. Its exact finite-level identification in the proof gives the last two assertions. Indeed \(s_*\le 1+\log_\ell(t[E:\QQ])\ll_{\psi,\eps}\log(2x)\). For all sufficiently large \(x\), both \(t\) and the cyclotomic modulus \(N\) are below \(P\), since \(\log(2N)\ll_{\psi,\eps}x^{\eps/4}\). For \(\ell>P\) any nonzero local loss would require \(\ell\mid N\) or \(\ell-1\mid r\), by Corollary~\ref{cor:sum31}. Both are impossible.
\end{proof}

\lead{The individual global periods}Use an unramified prime \(P_0\) and the standard definite chain of Section~\ref{nav:I-39} on \(E_x\). In Theorem~\ref{thm:lattice32}, one fixed orbit of finite Bessel level classes gives a common integral test space of rank \(n_0\), with one line for each of the \(n_0\) twist characters \(\gamma\). Let \(h_B\) be the number of such orbits. The covolume in the scalar coordinates and its bound are proved there. Let \(p_\gamma(y)\) be the integral mean of the actual restricted theta lift on the orbit representative \(y\) against the common archimedean test, in Haar probability. Then the projected column has squared norm \(|p_\gamma(y)|^2/h_B\) on a unit source, and all distinct columns are orthogonal. Proposition~\ref{prop:diag32}, reduces both archimedean moments for this model to a bound for the mean squared deviation of these numbers from one. Orthogonality by itself allows every one of them to vanish. Theorems~\ref{thm:definiteseparation530} and~\ref{thm:rectangulartheta530} recover coefficients of definite frames with bounded denominators, including fixed positive Gram coefficients when their residue image subspaces are retained. Proposition~\ref{prop:isotropickernel530} exhibits complete cancellation on singular frame orbits with radical of dimension at least two at the same terminal congruences. Theorem~\ref{thm:partnerlifts531} separates every integral lift of one rational partner flag. Proposition~\ref{prop:partnerweights531} identifies their signed operators, and Theorem~\ref{thm:cosetcomparison531} bounds interference between distinct residue cosets. Theorem~\ref{thm:orbitaverage532} computes the signed sum at each residue in a two-sided arithmetic orbit as an explicit projection. Theorem~\ref{thm:doubleorbit532} recovers these operators integrally on the common fields. Theorems~\ref{thm:fullrankseparation533} and~\ref{thm:orbitnoncancellation533} separate distinct full-rank rational transports below an explicit denominator bound and recover the individual operator on a single two-sided orbit. Theorem~\ref{thm:bothranks534} includes both common positive ranks for equal terminal norms. Proposition~\ref{prop:repeatedtransport534} gives infinitely many transverse partners with the same transport and height one. Corollary~\ref{cor:actualgram535} proves separation under a cross Gram determinant bound for both original terminal frequencies. Proposition~\ref{prop:retainedtest535} shows why the retained columns require a sharper denominator estimate. Theorem~\ref{thm:everydiagonal536} constructs infinitely many terms at one residue on every retained diagonal row. Theorem~\ref{thm:diagonaltail536} identifies their signed branch sums and bounds the ordinary operator tail for each fixed row. Theorem~\ref{thm:hyperplanerecovery537} gives integral recovery and asymptotically equal finite norms for integral transports with fixed marked hyperplanes. Proposition~\ref{prop:realizedhyperplanes537} realizes each choice of marked hyperplanes on retained diagonal rows. Theorem~\ref{thm:totaldiagonal540} bounds the sum over all ordinary partners outside the rational decorated stabilizer on a diagonal row, using congruences at both original marked primes. Corollary~\ref{cor:projecteddiagonal540} gives the diagonal error after the prescribed finite projection, with its dependence on the row and real section. Theorem~\ref{thm:mixedordinary541} treats distinct original rows and both terminal norms, using explicit counting constants. Corollary~\ref{cor:mixedmatrix541} bounds their simultaneous projected operator error by row and column sums. Theorems~\ref{thm:flagisometry544} and~\ref{thm:ordinarysingular544} compute the singular values of the flag-preserving terms on distinct decorated rows and bound the mean square error of the full projected ordinary period there. Theorems~\ref{thm:identitycoefficient551} and~\ref{thm:integralfamily551} recover the identity coefficient on these rows, with an explicit bound for every other marked index and a simultaneous integral coordinate choice at the same two primes. Theorem~\ref{thm:ordinaryunit552} bounds the full ordinary diagonal error on every norm-one row by an explicit real-section norm times a uniform negative power of the original marked level. Theorem~\ref{thm:bothordinary553} proves the corresponding arithmetic bound for all pairs of retained primitive-frame rows at both terminal norms. It includes the full signed coefficient at the identity marking. Theorem~\ref{thm:idealordinary554} extends the bound to the specified outer powers of the inverse different. Their ideal norms cancel in both height separation and the exterior lattice count. Theorem~\ref{thm:primemoment555} proves that the normalized absolute upper bound grows on the prime-terminal rows, using concentration of their actual lowest-type sections. The signed period comparison still requires a sharper argument. For mixed terminal frequencies Theorems~\ref{thm:terminaldivisor556} and~\ref{thm:mixedhalf556} give a terminal-prime divisor for every nonzero intersection and a marked-period estimate in unit-frequency norms when the intersection contains at least half the outer directions. Corollary~\ref{cor:denomperiod562} gives a single bound for the large-denominator terms at every proper nonzero intersection, including at equal prime terminal norms. Lemma~\ref{lem:scalarden562} bounds the number of remaining scalar values on those strata. Theorem~\ref{thm:markedindex564} supplies a divisor of the common flag index at both marked primes. Corollary~\ref{cor:markedperiod564} incorporates it in the signed period estimate, including terms with small scalar denominator. Theorem~\ref{thm:middleterminal570} extracts a further terminal divisor from the common middle lattice of intersecting flags. Corollary~\ref{cor:middleperiod570} bounds the mixed period on the full range of prefix dimensions with \(2t\le s+1\). Theorem~\ref{thm:positivecompare576} separates distinct rational isometries by the indices of the common positive lattice, including at marked primes dividing those indices. Corollary~\ref{cor:indexperiod576} gives integral recovery on small-index collections and a uniform bound on the large-index terms. For transverse hyperplanes Lemma~\ref{lem:cofactormark581} proves a determinant congruence even at a mark dividing the pairing discriminant. Theorem~\ref{thm:transverseseparate581} and Corollary~\ref{cor:transversetrace581} give a separation bound and a signed orbit trace with the actual local denominator ideals. Theorem~\ref{thm:transversesum582} bounds the absolute sum of transverse terms above a cutoff given by the frame determinants. Corollaries~\ref{cor:terminalprefix582} and~\ref{cor:cyclicuniform582} include terms detected at the terminal prime and the union of the indicated repeated-transport families. Theorem~\ref{thm:transportproduct589} separates distinct transports using the product of their discriminant ideals. Corollary~\ref{cor:smalltransverse589} bounds their overlap in the remaining low-height transverse sum on the original rows. Theorem~\ref{thm:smallframes590} and Proposition~\ref{prop:smallcross590} count the rational classes and integral cross matrices at low frame height in both terminal frequencies. Corollary~\ref{cor:smallgroups590} gives the resulting comparison for the signed groups. Theorem~\ref{thm:smallorbits591} bounds the number of individual integral double cosets below the height cutoff, including terms with the same marked index, for both terminal norms. Theorem~\ref{thm:twoenergies592} bounds the signed ordinary orbital integrals by normalized squared section moments. Theorem~\ref{thm:exactcore592} computes the projected diagonal with both real heights restricted and bounds its mass at prime terminal frequency. Theorem~\ref{thm:traceoverlap602} separates distinct transports in the proper-intersection period up to positive lattice index bound \(Q/(2k\sqrt u)\) on the \(f\)-th root scale. Theorem~\ref{thm:rankmarks602} compares their rational and marked ranks. Corollary~\ref{cor:properperiod602} bounds the complementary signed contribution by moments of the original sections. Theorem~\ref{thm:hyperplanerank603} gives a unique rational transport at each full transverse index below discriminant height \(Q/2\), including at singular marks. Corollaries~\ref{cor:rankpacking603} and~\ref{cor:oneindexperiod603} bound the multiplicity above that height and its contribution to the signed hyperplane period. Theorem~\ref{thm:smallcomplement604} describes all transverse rational classes at fixed projected hyperplanes and transport in the special unitary group of their orthogonal complement. Theorem~\ref{thm:liftphase604} sums the phases of the integral lifts of each such class. Corollary~\ref{cor:signedslice604} gives the exact squared norm of the resulting signed operator. Theorem~\ref{thm:localphase605} factors the signed coefficient of every supported transverse rational class at a fixed index into local algebraic integers after division by the covolume of a single actual lift and the common row volume. Corollary~\ref{cor:phasevaluation605} tests exact local cancellation and bounds its normalized terminal valuation. Theorem~\ref{thm:residueorbits606} computes the cancellation test on finite residue classes using the original row matrices. Corollary~\ref{cor:orbitideal606} bounds all conjugates of the normalized arithmetic coefficient and the number of its integral lifts by an ideal norm computed from the rational class. Theorem~\ref{thm:rowfactor607} evaluates the class coefficient with the covolumes of the two original rows. Its complex bound is independent of the precision used for the local residue sums. Corollary~\ref{cor:rowperiod607} gives the signed operator formula with this normalization. Theorem~\ref{thm:hyperplanesum610} sums different projected hyperplanes at a fixed marking and transport. Corollary~\ref{cor:heightoneplanes610} computes their cross terms exactly at normal height one. Theorem~\ref{thm:angularerror610} bounds those terms at larger heights in each compact representation. Theorem~\ref{thm:frameorientation611} classifies the remaining rational classes above a fixed cross matrix by norm-one rotations. Proposition~\ref{prop:rotationideal611} bounds their integral support by a divisor count, and Theorem~\ref{thm:rotationperiod611} computes their full contribution using one real tensor. Theorem~\ref{thm:rotationcharacters612} evaluates the sum over the frame determinants by finite local character sums with the actual marking. Corollary~\ref{cor:rotationroots612} gives its exact value at normalized Gram determinant one and preserves the terminal valuation for identical columns. Theorem~\ref{thm:orientationmoment613} sums the supported rational frame determinants at all Gram heights for either terminal frequency, with a positive moment of the normalized Gram determinant. On the cell where every leading minor is nonzero, Corollary~\ref{cor:openprofiles615} bounds the sum of squared real orbital profiles over supported frame determinants by explicit mixed derivative energies. The estimate stays uniform down to zero terminal Gram value. Proposition~\ref{prop:pivotflags619} supplies integral flag minors at every prefix also on the other positive-Gram blocks. Theorem~\ref{thm:pivotintegral619} bounds the original transverse orbital operators on all these blocks by two explicit exterior-power moments of the prescribed rows. Proposition~\ref{prop:exteriortransport620} proves finiteness in the indicated range and computes the full terminal-frequency factor in both moments. Theorem~\ref{thm:exteriorcost620} shows that their absolute majorant grows at prime frequency even after extracting the known lower Gram cutoff.

\begin{proposition}[Finite divisibility and the nonzero period values]\leavevmode\label{prop:prod33}
Fix such data, and suppose their normalized scalar periods \(p_\gamma\) are elements of a common number field \(J\). Use the same scalar identification at each embedding \(\sigma:J\hookrightarrow\CC\). Let \(G_0=\{\gamma:p_\gamma\ne0\}\). Suppose for some \(0\le e\le1/4\) one has
\[
 \frac{1}{[J:\QQ] n_0}\sum_{\sigma,\gamma}
       \left(\frac{|\sigma p_\gamma|^2}{h_B}-1\right)^2\le e .
\]
For \(\gamma\in G_0\) let the integers \(a_{\gamma,w}\) have finite support in the finite places of \(J\) and satisfy \(a_{\gamma,w}\le v_w(p_\gamma)\) at every such place; the valuations assign one to a uniformizer. Then \(|G_0|\ge(1-e)n_0\) and
\begin{equation}
 \frac{\sum_{\gamma\in G_0}\sum_{w\nmid\infty} a_{\gamma,w}
                        \log(\operatorname{Norm} w)}
      {[J:\QQ]|G_0|}
 \le\tfrac12\log h_B+
             \tfrac12\log\!\left(\frac{1+\sqrt e}{1-e}\right).
                                                               \label{eq:scal33}
\end{equation}
\end{proposition}
\begin{proof}
A zero algebraic period contributes one to the squared deviation at every embedding, so the stated lower bound for \(|G_0|\) follows. By Cauchy--Schwarz the average of \(|\sigma p_\gamma|^2/h_B\) over all \(\gamma,\sigma\) is at most \(1+\sqrt e\). For a nonzero algebraic number, the sum of its valuation times the log residue norm over finite places is the log absolute norm to \(\QQ\), hence the sum of the logs of the absolute values across all embeddings. Apply this identity for each \(\gamma\in G_0\), using the valuation bounds. Concavity of the logarithm bounds the resulting mean by one half of the log of the mean of \(|\sigma p_\gamma|^2\) over the nonzero periods. That latter mean is at most \(h_B(1+\sqrt e)n_0/|G_0|\). This proves \eqref{eq:scal33}.
\end{proof}

To obtain \eqref{eq:gap} from the periods, their finite valuations and archimedean sizes must satisfy the comparison in the same normalization. The preceding proposition proves the passage from estimates on nonzero period values to the corresponding finite divisibility. We next compare the finite-place count with a divisibility statement that can already be proved on the same fields.

\lead{A local comparison at the distinguished deformation}Theorem~\ref{thm:full36} proves the integral cotangent estimate for presentations with arbitrary numbers of coordinates of each of the specified weights. More precisely, for an augmented integral graded algebra with the smooth chart and point hypotheses stated there, the length of the weight-minus-one cotangent quotient is at most twice the sum of the normalized valuations of a separating trace difference on its mixed scaling orbits. The proof includes components supported at the residue characteristic by taking the difference between the cokernel and kernel lengths of multiplication by that trace. For the large primes, Sections~\ref{nav:IV-09} and~\ref{nav:IV-11} construct a marked arithmetic presentation and analyze its points. The chart, rigidity and remaining arithmetic bounds are discussed at the end of this section.

\lead{The count carried by integral CM derivatives}There is now a finite-place bound over the actual fields \(E_x\). Their common imaginary quadratic subfield on each interval is \(K=\QQ(\sqrt{-q})\). Section~\ref{nav:V-00} constructs CM abelian schemes of the type induced from \(K\) and bounds an integral derivative of order \(a\) at their identity. For a nonsplit prime, let \(e_p,f_p\) be the ramification index and residue degree in \(K\). In the normalization \(v_p(p)=1\) the lower bound is
\[
     v_p(a!)-e_p^{-1}\left\lfloor\frac{a-1}{p^{f_p}-1}\right\rfloor .
\]
For weights at least \(2(1+\lceil\log_2(a+1)\rceil)\), every integral derivative satisfies \(v_p(D^a f(0))\ge C_p(t,a,w)/(2\log p)\) at every prime, including the ramified prime. Here \(C_p\) is the summand with rational prime \(p\) in the count \(C\) of Section~\ref{nav:I-23}. The extension field supporting the CM schemes may be fixed throughout the interval. The next result computes the size of the division that a passage from such jets to a normalized period would have to account for.

\begin{theorem}[The integer quotient for the weighted count]\leavevmode\label{thm:divide38}
Use the fields and translates of Proposition~\ref{prop:field12} on the intervals of Section~\ref{sec:missing}. Write
\[
 b=\exp(x^\eps),\quad a=\lfloor bx\rfloor,\quad
 B'=\lceil b\rceil,\quad W=\lceil\sqrt b\rceil,\quad
 \omega=h_2/f=1/(2o).
\]
Give the distinguished positions weight \(B'\) and the other positions weight \(W\) at every stage. Denote this list by \(w\). It is independent of \(u\) since \(T_2=A_0\) and \(A_0+u=A_0\). Write \(w_*\) for the list identically equal to \(W\). For each \(t\in T\), define positive integers
\[
 D_t=\prod_p p^{\left\lfloor C_p(t,a,w)/(2\log p)\right\rfloor},
 \qquad D_*=\prod_p p^{\left\lfloor C_p(t,a,w_*)/(2\log p)\right\rfloor},
 \qquad R_x=\prod_{t\in T}(D_*/D_t).
\]
The second product is independent of \(t\), and every \(D_t\) divides \(D_*\). Every integral CM jet of this order in Corollary \ref{cor:ideal37}, is divisible by \(D_*\) in the ring of integers of its coefficient field. On all sufficiently distant common intervals, simultaneously for every \(u\),
\begin{equation}
  \left|\Delta_u-
       \left(hm a\omega\log(B'/W)-2m\log R_x\right)\right|
                        \le hm\,b/x .                          \label{eq:divide38}
\end{equation}
\end{theorem}
\begin{proof}
We have \(B'\ge W\ge2(1+\lceil\log_2(a+1)\rceil)\) for large \(x\). Also \(a/W\ge\exp(x^{\eps/3})/3\) for large \(x\). Proposition \ref{prop:weights14}, therefore applies to this choice, in particular to the comparison with the retained prime sums. Every term has weight at least \(W\), so primes \(p>a/W\) contribute zero. The products are finite. Monotonicity in each weight gives \(D_t\mid D_*\). With \(w_*\), the local inequality of Theorem \ref{thm:count37}, gives divisibility by \(p^{\lfloor C_p/(2\log p)\rfloor}\) at every prime over \(p\). At the other places this rational integer is a unit. Taking the product gives the assertion for jets.

Set \(Y=a/W\). The constant-weight count \(C_*\) is independent of \(t\) and obeys
\[
        C_*=W\{R_{\mathbf1}(Y)-R_{\xi_c^*}(Y)\},\qquad
        C_*-M(a,w_*)=W\{G_{\mathbf1}(Y)-G_{\xi_c^*}(Y)\}.
\]
In fact the sum over the entire nontrivial coset of \(B/\Phi\), including the inertia average, gives \(1-\xi_c^*(p^j)\) at every prime power. Substitution of the corrections in \eqref{eq:Gchi} gives the second identity. By the zero-free estimate of Lemma \ref{lem:far14}, its absolute value is at most \(2a\exp(-c x^{\eps/8})\) for some \(c>0\). The precise main-term difference is
\[
            M(a,w_*)-M(a,w)=a(h_2/f)\log(B'/W)
\]
because precisely \(h_2\) of the \(f\) positions change weight.

For a fixed \(t\), the error in replacing \(C_*-C(t,a,w)\) by \(2\log(D_*/D_t)\) is at most \(2\sum_{p\le Y}\log p\le 2Y\log(2Y)\). At each prime the two fractional parts differ by less than one. Finally \(2a\exp(-c x^{\eps/8})+2Y\log(2Y)\le b/x\) for large \(x\), since \(a\le bx\) and \(W\ge\exp(x^\eps/2)\). Compute \(C(t,a,w)-M(a,w)\) by adding and subtracting the constant-weight count and its main term. Summing over the \(h\) choices of \(t\) and the \(m\) identical stages proves \eqref{eq:divide38}.
\end{proof}

\begin{corollary}[The excess divisibility of an unnormalized jet]\label{cor:excess38}
On this family every nonzero integral CM jet \(z\) of order \(a\), over a number field \(J\) as in Volume~XIII, satisfies
\[
 \frac{2}{[J:\QQ]}\log|N_{J/\QQ}z|
          \ge M(a,w)+a\omega\log(B'/W)-b/x .
\]
The excess term \(a\omega\log(B'/W)\) is asymptotic to \(a x^\eps/(4o)\). It exceeds the allowed average error per summand \(x^{5\eps}\mathcal B/(hm)\) by a factor tending to infinity. Thus an upper estimate for the unnormalized heights at the weighted main term with that error is impossible on this family.
\end{corollary}
\begin{proof}
The integer \(D_*\) divides \(z\) in \(\mathcal O_J\). For a nonzero algebraic integer \(z/D_*\) its norm has absolute value at least one. The rounding estimate for \(2\log D_*\) and the zero-free estimate for \(C_*\) in the preceding proof give \(2\log D_*\ge M(a,w_*)-b/x\). This proves the inequality. Here \(\log(B'/W)=x^\eps/2+o(1)\) and \(\omega=1/(2o)\). In \eqref{eq:gap} substitute \(m=\lfloor\sqrt x\rfloor,\ k=xm,\ a_E=xb,\ h\ge n/2,\ h_2\le n\), and \(n\ge x^{2+8\eps}\). The three resulting terms give \(x^{5\eps}\mathcal B/(hm)=O_\psi(bx^{1/2+6\eps})=o(a)\) because \(\eps<1/52\). This proves the claimed separation in size.
\end{proof}

The integer quotient has a direct interpretation in terms of ordinary prime sums. For \(t\in T\) let \(S_t=p_t^{-1}(L_u)\), independent of \(u\), and put
\[
 \Psi_t(s)=\sum_{p^j\le s}\frac{2\log p}{|I_p|}
                           |(jF_p+I_p)\cap S_t| .
\]
For a distinguished position the difference of its two ramps equals
\(\int_W^{B'}\mathbf1_{y<a/p^j}\,dy\). The nondistinguished positions contribute zero to the difference. Interchanging finite sums and integration and writing \(s=a/y\) therefore gives the exact formula
\[
     C_*-C(t,a,w)=a\int_{a/B'}^{a/W}\frac{\Psi_t(s)}{s^2}\,ds .
\]
The expected linear function in this identity is \(\omega s\), whose integral contributes \(a\omega\log(B'/W)\). Dividing a jet by the integer \(D_*/D_t\) preserves integrality and leaves divisibility by \(D_t\). An independent lower bound for the product of the factors removed, or an independent height estimate for the actual periods after their arithmetic normalization, is still needed to prove \eqref{eq:gap}. The lower bound required for \(2m\log R_x\) follows explicitly from \eqref{eq:divide38} up to \(hm b/x\). Neither the separation in Corollary~\ref{cor:excess38} nor the formal derivative estimate gives that lower bound. The cohomology bound on the same family remains an additional requirement of the proposed construction.

\begin{theorem}[Least scalar-jet norm on the common intervals]\label{thm:height42}
Keep the fields, weights and notation of Theorem~\ref{thm:divide38}. On each sufficiently distant common interval choose the field \(J\) and CM abelian schemes of Theorem \ref{thm:exact41}. At each \(x\) let
\[
  \mathcal H_K(a)=\min_{z\ne0}\frac{2}{[J:\QQ]}\log|N_{J/\QQ}z|
\]
over all integral scalar derivatives \(z=(D^a f)(0)\) on that scheme. Write \(\xi_K\) for the primitive quadratic character of \(K=\QQ(\sqrt{-q})\), \(c_K(n)=(1-\xi_K(n))\Lambda(n)\), and \(\Psi_K(y)=\sum_{n\le y}c_K(n)\). Put \(Y=a/W\). Then
\begin{align}
 \mathcal H_K(a)&=\sum_{n\le a}c_K(n)\lfloor a/n\rfloor,             \label{eq:height42}\\
 \mathcal H_K(a)-C_*&=\sum_{k=1}^W\Psi_K(a/k)
                 -\sum_{n\le Y}c_K(n)\{a/n\}.                       \label{eq:harm42}
\end{align}
Here \(\{z\}=z-\lfloor z\rfloor\). On these intervals, uniformly in \(x\) and every translate \(t\),
\begin{equation}
 \mathcal H_K(a)-C(t,a,w)
     =a(\log W+\gamma)+2\log(D_*/D_t)+O_{\psi,\eps}(b/x).
                                                                    \label{eq:minper42}
\end{equation}
The minimum is attained by a globally defined jet. The vector may be fixed for all orders used on each field of the interval.
\end{theorem}
\begin{proof}
Theorem \ref{thm:exact41}, computes the image ideal as \(H_K(a)\mathcal O_J\) for one such vector, and proves universal divisibility for all vectors. A nonzero algebraic integer has absolute norm at least one. Thus \(H_K(a)\) itself attains the minimum, which equals \(2\log|H_K(a)|\) in either complex embedding of \(K\). At inert \(p\) its normalized valuation is
\[
  \sum_{j\ge1}\lfloor a/p^j\rfloor-
  \sum_{j\ge1}\lfloor a/p^{2j}\rfloor
       =\sum_{j\ge1,\ j\ {\rm odd}}\lfloor a/p^j\rfloor .
\]
At the ramified prime it is \(\tfrac12\sum_{j\ge1}\lfloor a/q^j\rfloor\), and at a split prime it is zero. These are precisely the coefficients of the prime powers in \eqref{eq:height42}, including the factor two in the definition of the minimum.

By Theorem~\ref{thm:divide38}, \(C_*=\sum_n c_K(n)(a/n-W)_+\). For \(n\le Y\) the difference between the floor and the ramp equals \(W-\{a/n\}\). For \(Y<n\le a\) the floor counts the integers \(1\le k<W\) such that \(n\le a/k\). Summing first in \(n\), then in \(k\), cancels \((W-1)\Psi_K(Y)\) against all but one copy of the term \(W\Psi_K(Y)\). Since \(Y=a/W\), this proves \eqref{eq:harm42} with both endpoints included.

We give the uniform estimate for the unsmoothed sum here because the conductor varies with the interval. Write \(R(y)=R_{\mathbf1}(y)-R_{\xi_K}(y)\). By \eqref{eq:Gchi} and Lemma \ref{lem:far14}, for some \(c>0\) and a real number \(A_K\) independent of \(y\),
\[
 R(y)= y\log y + A_K y + O(y\delta),\qquad
     \delta=\exp(-c x^{\eps/8}),\qquad
     y\ge \tfrac13\exp(x^{\eps/3}).
\]
The constant in the error is uniform on the common family. The coefficients \(c_K(n)\) are nonnegative, so \(R(y)=\sum_n c_K(n)(y/n-1)_+\) is convex and its right derivative satisfies
\(yR'_+(y)-R(y)=\Psi_K(y)\), also at integers. For \(Y\le y\le a\), both \(y(1\pm\sqrt\delta)\) are in the range of the estimate for large \(x\), since \(Y\asymp x\exp(x^\eps/2)\). The two adjacent secant slopes at these three arguments bound \(R'_+(y)\). The secants for the main term are \(\log y+1+A_K+O(\sqrt\delta)\) by Taylor's formula, and the error terms in the secants are \(O(\delta/\sqrt\delta)\). Hence
\[
                   \Psi_K(y)=y+O(y\sqrt\delta)
                \qquad (Y\le y\le a)
\]
uniformly. The constant \(A_K\) cancels. This use of convexity only requires the quadratic character of the common subfield.

The last sum in \eqref{eq:harm42} lies between zero and \(\Psi_K(Y)=O(Y)\). Write \(H_W=\sum_{k\le W}1/k\). The first sum is \(aH_W+O(a H_W\sqrt\delta)\). Since \(H_W=\log W+\gamma+O(1/W)\), we have
\[
           \mathcal H_K(a)-C_* =a(\log W+\gamma)
                          +O(Y+aH_W\sqrt\delta).
\]
Both errors are \(o(b/x)\) for fixed \(\eps>0\). The difference between \(C_*-C(t,a,w)\) and \(2\log(D_*/D_t)\) is at most \(2Y\log(2Y)=o(b/x)\) by the prime-by-prime rounding in Theorem~\ref{thm:divide38}. Adding proves \eqref{eq:minper42}.
\end{proof}

\begin{corollary}[What persists after dividing the scalar derivatives]\label{cor:divide42}
On these same CM schemes one has the following two statements.
\begin{enumerate}\setlength{\itemsep}{1pt}
\item With \(D\) as in the theorem, integral scalar jets divided by their full common factor \(H_K(a)\) have every value in \(\mathcal O_J\) as a normalized derivative, including one. Thus local integrality of this quotient alone imposes no proper ideal of divisibility depending on \(t\).
\item Using the smaller divisor \(D_t\), the least normalized log norm \(2[J:\QQ]^{-1}\log|N(z/D_t)|\) among nonzero jets is
\[
       a(\log W+\gamma)+2\log(D_*/D_t)+O_{\psi,\eps}(b/x).
\]
\end{enumerate}
\end{corollary}
\begin{proof}
The first assertion is the image ideal calculation in Volume~XIII. For the second the minimum is \(\mathcal H_K(a)-2\log D_t\). Replace \(\mathcal H_K(a)\) by \(C_*+a(\log W+\gamma)+O(b/x)\) as in the proof above and use \(0\le C_*-2\log D_*\le2Y\log(2Y)\).
\end{proof}

\lead{The remaining arithmetic inequalities}
For the large primes \(\ell>\exp(x^\eps/12)\), Theorems \ref{thm:ring39} and \ref{thm:nocorner40} realize the strict classes of each complementary character as first-order deformations with a marked inertia element and fixed local matrices. These results exclude the two-dimensional mixed corner on restriction to \(G_{\QQ(i)}\). Theorem~\ref{thm:ststrict616} realizes the strict classes in weight one at integral lifts with regular tame monodromy at two primes split in every interval field. Proposition~\ref{prop:stweight616} gives the exact depth to which an unramified lower triangular lift in each higher weight satisfies the same local equation. An integral mixed lift that survives must be absolutely irreducible there, cannot be induced from a character and has three different absolute values for crystalline Frobenius at \(\ell\). To use the cotangent estimate, one must still establish the presentation and point hypotheses of Theorem \ref{thm:full36} for this arithmetic model, as well as the required bound on the surviving traces. The stronger fixed local matrices used in the construction retain the strict first-order classes, but can add relations to the presentation. The actual weighted cohomology bound in Proposition~\ref{prop:minors33} also includes the other retained primes.

For the theta comparison, the arithmetic norm in Theorem~\ref{thm:height42} is now attained on the actual fields, while the distinguished Frobenius classes enter only through \(D_*/D_t\). The prescribed Bessel periods in Proposition~\ref{prop:prod33} have additional constraints from their kernels and their common projection. A bound for their finite valuations and archimedean size in that normalization, or an independent lower estimate for \(2m\log R_x\) sufficient in \eqref{eq:divide38}, is still required on the common intervals. In particular \eqref{eq:height42} holds even when the twist-dependent lower estimate is unknown. The unconditional implication from these constructions to \eqref{eq:G} remains incomplete.

\begin{proposition}[The exact Frobenius trace ideal of a mixed lift]\label{prop:trace45}
For the same large primes, fix an integral mixed lift of the marked model with coefficient discrete valuation ring \(R\). Put \(H=G_{\QQ(i)}\) and let \(U:H\to\mathrm{SL}_3(R)\) be its matrix in the marked basis. Let \(d_1=\Psi,d_2=1,d_3=\Psi^{-1}\) on \(H\); write \(t_0(h)=\sum_i d_i(h)\). For distinct indices let \(I_{ij}\subset R\) be the ideal generated by the entries \(U(h)_{ij}\) for all \(h\in H\). Then
\begin{equation}
 \big(\operatorname{tr}U(h)-t_0(h):h\in H\big)
       = I_{12}I_{21}+I_{13}I_{31}+I_{23}I_{32}
       = I_{12}I_{21}+I_{13}I_{31}.                         \label{eq:cycles45}
\end{equation}
Write \(v(I)=\min_{0\ne z\in I}v_\ell(z)\) and \(v_\ell(\ell)=1\). Outside any prescribed finite set, the minimum valuation of the trace difference at Frobenius elements of degree-one primes of \(\QQ(i)\) is exactly
\[
               \min\{v(I_{12})+v(I_{21}),\ v(I_{13})+v(I_{31})\}.
\]
All four ideals here are nonzero. Each minimum is attained, with infinitely many possible underlying rational primes.
\end{proposition}
\begin{proof}
By Sections \ref{nav:IV-09}--\ref{nav:IV-12}, the residual representation is diagonal with entries \(\bar d_i\), the chosen local decomposition group at \(\ell\) has the exact prescribed diagonal, and the restriction to \(H\) is absolutely irreducible. Write \(g_*\in H\) for the marked inertia element. The three entries \(k_i=d_i(g_*)\) are pairwise distinct modulo the maximal ideal. Indeed \(\bar\Psi(g_*)\) has order \((\ell-1)/\gcd(\ell-1,r)>2\), and the three entries are \(\bar\Psi(g_*),1,\bar\Psi(g_*)^{-1}\).

Denote by \(J_0\) the left side of \eqref{eq:cycles45} and by \(J_d\) the ideal of the diagonal differences \(U(h)_{ii}-d_i(h)\). For \(s=0,1,2\), the trace difference at \(g_*^s h\) is \(\sum_i k_i^s(U(h)_{ii}-d_i(h))\). The Vandermonde determinant is a unit in \(R\). Thus all diagonal differences are \(R\)-linear combinations of trace differences, while their sum gives the trace difference. Hence \(J_0=J_d\).

Let \(J_c=\sum_{i<j}I_{ij}I_{ji}\). We show first that \(J_c\subset J_d\). For a fixed \(j\), let \(P_j(X)=\prod_{i\ne j}(X-k_i)/(k_j-k_i)\in R[X]\). Evaluating at \(U(g_*)\) gives the projector onto coordinate \(j\). For \(i\ne j\) and \(h,h'\in H\), the \(ii\) entry of \(U(h)P_j(U(g_*))U(h')\) is \(U(h)_{ij}U(h')_{ji}\). Expand \(P_j\) in powers of its variable. Modulo \(J_d\) each entry \(U(hg_*^s h')_{ii}\) is \(d_i(hh')k_i^s\), so the resulting entry is \(d_i(hh')P_j(k_i)=0\). This gives the inclusion for every pair of factors.

Conversely, modulo \(J_c\), the diagonal entries of a product are the products of its diagonal entries. Thus \(u_i(h)=U(h)_{ii}\bmod J_c\) is a continuous character to \((R/J_c)^\times\), since these entries reduce to units. Each \(I_{ij}\) is nonzero: the span of the orbit of coordinate vector \(e_j\) is the whole representation by irreducibility, and so some orbit vector has nonzero \(i\)-th entry. All off-diagonal entries reduce to zero, so \(J_c\) is a nonzero proper ideal. The characters \(u_i/d_i\) are consequently valued in the principal units of a finite local quotient of \(R\), an \(\ell\)-group. On inertia away from \(\ell\) the marked representation and the reference characters factor through a finite group of order prime to \(\ell\), by their defining inertia kernels and the large-prime cutoff. Thus each ratio is trivial there. It is trivial on the chosen decomposition group above \(\ell\) by the fixed local matrices. At a conjugate place the same is true, including a place reached by \(c\): the antidiagonal inverse transpose sends the diagonal \(\operatorname{diag}(s,1,s^{-1})\) to itself, and the reference characters come from \(G_\QQ\). Thus each ratio is everywhere unramified. The ordinary ideal class group of the imaginary quadratic field \(\QQ(i)\) is trivial, so global class field theory kills these characters. Therefore \(J_d\subset J_c\), proving the first equality.

The marked symmetry \(U(ch^{-1}c^{-1})=J_3 U(h)^{\mathsf T}J_3\) gives \(I_{12}=I_{23}\) and \(I_{21}=I_{32}\), which proves the other equality. A sum of nonzero ideals in a discrete valuation ring has valuation equal to the least of their valuations. The ideal generated by all the traces thus has the stated valuation, and one of the generating traces attains it by discreteness.

Take such an \(h\) and reduce the matrices and reference characters modulo a power of the uniformizer strictly higher than the attained integral valuation. Their kernels contain an open normal subgroup of \(H\); intersect with its conjugate under \(c\) to obtain an open subgroup normal in \(G_\QQ\), still contained in \(H\). In the resulting finite Galois extension of \(\QQ\), the Chebotarev density theorem realizes the conjugacy class of \(h\) at infinitely many unramified rational primes outside the excluded set. Since \(h\in H\), all such primes split in \(\QQ(i)\). Choose a prime of the extension whose Frobenius is the specified representative of the conjugacy class. At the prime below it in \(\QQ(i)\), the evaluated trace is the specified one modulo the chosen power, hence has the same difference valuation. One may fix either consistent Frobenius convention in this argument.
\end{proof}

\begin{corollary}[The trace on the even alternative]\label{cor:even45}
In the second alternative of Theorem \ref{thm:mon44}, write \(B,C\) for the ideals generated, respectively, by the upper- and lower-off-diagonal entries of its even two-dimensional representation \(\rho(h)\), \(h\in H\). The ideal in \eqref{eq:cycles45} is exactly \(BC\). For any finite collection of mixed points in the same marked model, one integral linear combination of split Frobenius trace differences attains the pointwise minimum of Proposition~\ref{prop:trace45} at every point, after finite unramified extension of the base coefficient field if needed.
\end{corollary}
\begin{proof}
Write a matrix of \(\rho\) as \(\left(\begin{smallmatrix}a&b\\c&d\end{smallmatrix}\right)\) with unit determinant and unit diagonal entries, since its reduction is \(\operatorname{diag}(\bar\Psi,1)\). In the trace-zero basis \(E_{12},\operatorname{diag}(1,-1),2E_{21}\), conjugation of the middle basis vector has first coordinate \(-2ab/\det\rho\), and conjugation of \(E_{12}\) has middle coordinate \(-ac/\det\rho\). The two extreme entries are \(-2b^2/\det\rho\) and \(-c^2/(2\det\rho)\). Consequently \(I_{12}=B,I_{21}=C,I_{13}=B^2,I_{31}=C^2\), using discreteness to attain the least valuation among the squares. Formula \eqref{eq:cycles45} now gives \(BC+B^2 C^2=BC\).

For each point select one Frobenius that attains its minimum, excluding the finitely many bad primes of all the points. Use the union of these finitely many elements. At a fixed point, divide their trace differences by a generator of its trace ideal in a common coefficient extension and reduce. At least one coefficient in this row is nonzero. A combination fails to attain the minimum exactly when its coefficient vector lies on the hyperplane defined by this row. The product of the finitely many nonzero linear forms is a nonzero polynomial. Some vector over the algebraic closure of the residue field avoids all its zeros, and all of its entries lie in a finite residue extension. They lift to the integers of a finite unramified extension of the base coefficient field, enlarged enough to realize the chosen residues in the common extension. Their linear combination has the asserted valuation simultaneously. Trace differences are zero at the diagonal basepoint and are invariant under the scaling action, so this also supplies a separating function on the stated finite set.
\end{proof}

Equations \eqref{eq:cycles45} and \eqref{eq:divide38} concern the same family of fields and intervals. The first computes the individual trace valuations entering the graded cotangent comparison if its arithmetic presentation hypotheses hold. It gives no uniform upper bound for the off-diagonal ideals' valuations or for their sum over all points, primes and twists. On the other side Theorem \ref{thm:globaltheta43} attains the norm in Theorem~\ref{thm:height42} on a genuine global algebraic theta section of level \(3a\) relative to the chosen theta bundle, in one integral frame. Corollary~\ref{cor:parity43} controls a projection to inversion parity as well. The normalized Bessel projection in Proposition~\ref{prop:prod33} and its distinguished prime counts remain a separate estimate. These are the outstanding arithmetic steps in this route to Main Theorem~\ref{thm:grh16}.

\subsubsection{The CM type in the supplied construction}\label{nav:I-29}
There are two imaginary quadratic subfields in Proposition~\ref{prop:field12}.
Write \(K_q=\QQ(\sqrt{-q})\), \(K_d=\QQ(\sqrt{-d})\), and let
\(\xi_q,\xi_d\) denote their primitive quadratic characters.
The CM embeddings specified in the source send \(\sqrt{-d}\) to
\(i\sqrt d\). The first sign in its Galois group records the action on
\(\sqrt{-q}\), and the second sign records the action on \(\sqrt{-d}\).
These conventions are explicit in
\srcloc{Formalism/LocalRingCardinality/DataHost}{43} and
\srcloc{Formalism/SharpRestriction/RamificationFrobenius/TotientGaloisClasses/ConjEMemHostFIff}{66}.
The count in Theorem~\ref{thm:height42} uses \(K_q\). The distinction
can already be measured at the prime two.

\begin{proposition}[The two quadratic types at the prime two]\label{prop:dyad46}
Use \(a=\lceil xb\rceil,\ B'=\lceil b\rceil,\ W=\lceil\sqrt b\rceil\).
For the count indexed by the complement of \(\Gal(E/K_q)\), with weights
\(B',W\), the term at two satisfies, uniformly for every translate,
\[
 \frac{C_{2,q}}{2\log 2}=\frac{2a}{3}
                       +O\bigl(B'\log(2+a/B')\bigr).
\]
For the count indexed by the complement of \(\Gal(E/K_d)\), its term
at two is zero. On an abelian scheme of the induced \(K_d\)-type over
a common coefficient field, a global integral theta section can have
odd-order or even-order derivative of the prescribed order \(a\) which
is a unit at every place above two, in the frame specified below.
\end{proposition}
\begin{proof}
The discriminants are \(-q\equiv5\bmod8\) and \(-d\equiv1\bmod8\).
Thus two is inert in \(K_q\) and split in \(K_d\). Although \(E\)
itself is ramified at two, its inertia maps to the identity in both
quadratic quotients. For \(K_q\) precisely the odd powers of Frobenius
contribute, with total inertia weight one at each power after division
by \(2\log2\). Since every weight is between zero and \(B'\), the
normalized sum lies between \(\sum_{j\ {\rm odd}}(a/2^j-B')_+\)
and \(\sum_{j\ {\rm odd}}a/2^j=2a/3\). There are at most
\(1+\log_2(a/B')\) indices below the cutoff for large \(x\), and the
geometric tail past it is at most \(2B'\). This proves the estimate.
For \(K_d\) every Frobenius coset is contained in \(\Gal(E/K_d)\).

Here is the scalar ideal for the second type. For \(Q\) a prime power put
\(F_Q(a)=\sum_{j\ge1}\lfloor a/Q^j\rfloor\). With splitting and
inertia taken in \(K_d\), define
\[
 H_d(a)=
 \frac{a!}{\displaystyle
   \prod_{p\ {\rm split}}p^{F_p(a)}
   \prod_{p\ {\rm inert}}p^{F_{p^2}(a)}
   (\sqrt{-d})^{F_d(a)}} .
\]
Take an elliptic curve \(C_d\) of maximal CM by \(\mathcal O_{K_d}\)
and form \(\mathcal A_E=\mathcal O_E\otimes_{\mathcal O_{K_d}}C_d\).
The argument of Sections \ref{nav:V-00}--\ref{nav:V-05} applies with this field. We give the hypotheses that
change. Its only ramified prime is \(d\), of index two and residue
cardinality \(d\). At a split prime its formal elliptic factor has height
one, including at two. At an inert prime the scalar uniformizer is \(p\)
with residue cardinality \(p^2\). The local factorial calculation of
Lemmas~\ref{lem:factorial41} and~\ref{lem:unit41} consequently gives
\(v_p(a!)-F_{p^f}(a)/e\) for the attainable valuation at each place,
where \(e,f\) are the local index and residue degree. The number of
elliptic factors is \(2n>1+\log_2 a\) for large \(x\), so the separate
directions used for the nonzero base-\(p^f\) digits exist. There are
finitely many fields and integer orders on the bounded interval. The
residue extensions and Chinese remainder construction of Theorem~\ref{thm:exact41}
therefore give one coefficient field \(J\) and one tangent vector per
scheme for all these orders. The global image ideal is \(H_d(a)\mathcal O_J\).
For any relatively ample symmetric rigidified bundle \(P\), use a
section \(s_0\) of \(P^3\) with value one at the identity. Degree-\(a\)
products in that linear system realize every order-\(a\) jet in the
frame \(s_0^a\), by Lemma~\ref{lem:very43} and Theorem~\ref{thm:globaltheta43}. Some
section of \(P^{3a}\) thus gives \(H_d(a)\) itself. At two the denominator
contains \(2^{F_2(a)}\), with \(F_2(a)=v_2(a!)\). The asserted unit
property follows.
\end{proof}

\begin{theorem}[Changing the quadratic type in the centered count]\label{thm:orientation46}
Keep the fields, the sets \(T,T_2=A_0\), the weights and the integer
order of Proposition~\ref{prop:dyad46} on the common intervals of
Section~\ref{sec:missing}. Work first in the product coordinates
\((i,j,z)\in(\ZZ/2)^2\times A\) of Proposition~\ref{prop:field12}, in
the order \(q,d,A\). For \(e=q,d\), use the subgroup fixing \(K_e\)
for the count \(C_e\) and use \(D_{\xi_e}\) in its main term \(M_e\).
The distinguished Frobenius classes in the two counts at translate \(t\)
are, respectively,
\[
  \{(1,0,z-t):z\in A_0\},\qquad \{(0,1,z-t):z\in A_0\}.
\]
For \(\eta\in\widehat{A/A_0}\), write \(\alpha_{ij\eta}\) for the
primitive representative of \(\xi_q^i\xi_d^j\eta\). Then at each
individual \(t\), with \(o=|A/A_0|\),
\begin{align}
 &(C_q-M_q)-(C_d-M_d) \notag\\
 &\quad=\frac{b}{2o}\sum_{\eta\in\widehat{A/A_0}}\eta(t)
       \sum_{\substack{i,j\in\{0,1\}\\i+j=1}}
       \bigl((-1)^i-(-1)^j\bigr)G_{\alpha_{ij\eta}}(x)
       +O_{\psi,\eps}(b/x).                                      \label{eq:change46}
\end{align}
Its absolute value is at most \(2b x^{1/2+\eps/2}+O_{\psi,\eps}(b/x)\).
These assertions also hold when the original character \(\psi\) is odd.
\end{theorem}
\begin{proof}
Let \(c_A\in A\) be the restriction of complex conjugation to \(E_0\);
it has order dividing two. The change of coordinates
\[
 (i,j,z)\longmapsto(i,j,z+i c_A)
\]
is its own inverse. It sends complex conjugation to \((1,1,0)\) and
preserves both quadratic signs, the subgroup \(A\) and its projection
to the real cyclotomic layer. In these coordinates, addition of
\((1,1,t)\) sends the second distinguished set in the theorem to
\(\{(1,0,z):z\in A_0\}\) in the CM subgroup \(j=0\). Its preimage has
first sign zero and so has exactly the same coordinates before the
change. The first distinguished set is taken in the original coordinates
as in Proposition~\ref{prop:field12}. This accounts for conjugation
without assuming \(E_0\) totally real.

Apply character orthogonality as in \eqref{eq:count}. The sum over the
whole complementary coset selects just the trivial character and the
quadratic sign with coefficient \(-1\), giving the stated main term for
each orientation. For a nondistinguished weight one has
\(a/W\gg\exp(x^\eps/3)\). The bound of Lemma \ref{lem:far14}, times
the number \(O_\psi(n)\) of terms after division by \(2n\), is
\(O_\psi(na \exp(-c x^{\eps/8}))=o(b/x)\).
For the distinguished classes one can replace \(B'\) by \(b\) with
total error \(O_{\psi,\eps}(b/x)\), using the integral-weight
estimate of Proposition \ref{prop:weights14}. Its proof uses only
\(|B|,|A_0|\), the uniform zero-free bound and the slopes of the ramps,
which are the same for the two choices of subgroup.

Summing over \(z\in A_0\) annihilates characters nontrivial on \(A_0\).
The remaining ones are exactly \(\alpha_{ij\eta}\), and the factor
from this sum divided by \(2n\) is \(1/(2o)\).
Their conjugate values at the two distinguished sets supply
\((-1)^i\eta(t)\) and \((-1)^j\eta(t)\). The difference vanishes for
\(i=j\). The resulting argument of \(G\) is \(a/b\).
To replace it by \(x\), note that the slope of the prime sum on
\([x,x+1]\) is bounded by \(\sum_{v\le x+1}\Lambda(v)/v\ll\log(2x)\).
For the surviving nonprincipal characters their correction has slope
\(D_{\alpha_{ij\eta}}=O(\log(2x))\) by Theorem~\ref{thm:twists}.
The change is \(O_\psi(\log(2x))\) after multiplication by \(b\),
as \(0\le a/b-x\le b^{-1}\). This is \(o(b/x)\). Finally the two
coefficients per \(\eta\) each have absolute value two, and
Theorem~\ref{thm:twists} bounds both functions at \(x\) itself.
\end{proof}

\begin{corollary}[The count and scalar norm for the prescribed type]\label{cor:prescribed46}
In the \(d\)-orientation put
\(D_{t,d}=\prod_p p^{\lfloor C_{p,d}(t,a,w)/(2\log p)\rfloor}\)
and define \(D_{*,d}\) with the constant weight \(W\). For the minimum
normalized logarithmic norm of a nonzero scalar derivative on the
\(K_d\)-type schemes one has, uniformly in \(t\),
\begin{align}
 \frac{2}{[J:\QQ]}\log|N_{J/\QQ}H_d(a)|
    &=\sum_{v\le a}(1-\xi_d(v))\Lambda(v)\lfloor a/v\rfloor, \notag\\
 \frac{2}{[J:\QQ]}\log|N_{J/\QQ}H_d(a)|-C_d(t,a,w)
    &=a(\log W+\gamma)+2\log(D_{*,d}/D_{t,d})
                         +O_{\psi,\eps}(b/x).                     \label{eq:newheight46}
\end{align}
The replacement of the first centered count by the second changes
\(\Delta_u\) by \(O_\psi(hm b x^{1/2+\eps/2})\) for all \(u\).
When \(m\le\sqrt x,\ a_E=xb\), this is \(o(x^{5\eps}\mathcal B)\).
\end{corollary}
\begin{proof}
The ideal in Proposition~\ref{prop:dyad46} realizes the minimum by the
integrality of the absolute norm. At inert primes its normalized
valuation is the sum of \(\lfloor a/p^j\rfloor\) for odd \(j\).
At \(d\) it is half the sum over all \(j\), and at split primes it is
zero. This gives the first identity. The constant-weight count is
\(\sum_v(1-\xi_d(v))\Lambda(v)(a/v-W)_+\) by inertia averaging.
Subtract it from the first identity and apply the floor-sum identity
\eqref{eq:harm42} with \(K_d\). The secant-slope proof following that
identity applies because the coefficients are nonnegative and Lemma \ref{lem:far14}, applies to the same \(\xi_d\) on these intervals. Thus the
difference is \(a(\log W+\gamma)+O_{\psi,\eps}(b/x)\). Rounding the
exponents for primes at most \(a/W\) contributes at most
\(2(a/W)\log(2a/W)=o(b/x)\) to the difference of the two counts.
This gives \eqref{eq:newheight46}. Sum \eqref{eq:change46} over \(h\)
translates and \(m\) stages. Since \(\mathcal B\ge hxb\), the ratio of
the resulting bound to \(x^{5\eps}\mathcal B\) is
\(O_\psi(x^{-9\eps/2})\).
\end{proof}

The dyadic calculation rules out transferring the \(q\)-count
prime by prime to all integral sections of the \(d\)-type. The
centered count permits the change on the full interval by
Theorem~\ref{thm:orientation46}. Its terms with \(i=j=0\), including the
original Dirichlet characters, persist in each orientation. Estimating
their contribution still requires the normalized projected periods or
another proof of \eqref{eq:gap}. The strict cohomology groups depend on
the field \(E_x\) itself, which is the same in both counts.

\subsubsection{Sparse distinguished embeddings on every translate}\label{nav:I-30}
The source imposes a restriction on its distinguished set that fails
for an entire cyclic subgroup. In integer coordinates, three distinct
members \(s_1,s_2,s_3\) must avoid
\(j(s_2-s_1)=i(s_3-s_1)\) for \(0<|i|,|j|\le L,\ i\ne j\).
The integer representatives must also satisfy \(4Ls<n_0\), where
\(n_0\) is the cyclic order. See
\srcloc{Formalism/FrobeniusOrbits/HostThickExists}{13} and
\srcloc{Formalism/LocalRingCardinality/DataHost}{106}.
The length in
\srcloc{Development/AutomorphicArithmetic/CyclotomicBounds/ClassNumberBounds/ExplicitConstants/UniformIffAboveThreshold}{60}
grows logarithmically with \(x\) for fixed parameters. We construct
such a set while retaining the original prime sum at every translate.

\begin{lemma}[A sparse set balanced modulo four]\label{lem:sphere48}
Let \(s\) be a power of two tending to infinity. Let
\(2\le L\le C_0(1+\log s)^{C_1}\) be an integer, with \(C_0,C_1\)
fixed. There is a set \(S\) of nonnegative integers \(<s/(8L)\)
with the same number of elements in each residue class modulo four and
\[
 |S|\ \ge\ s\exp\bigl(-C\sqrt{\log s\,\log(2L+2)}\bigr)
\]
for all sufficiently large \(s\). If
\(\alpha+\beta+\gamma=0\) and the three coefficients are nonzero
integers of absolute value at most \(2L\), every solution of
\(\alpha u+\beta v+\gamma w=0\) in \(S\) has \(u=v=w\).
In \(\ZZ/s\ZZ\) a translate of \(S\) meets each segment of \(2L\)
successive points of an orbit of order \(>2L\) at most twice. It meets
each entire orbit of order \(\le8L\) at most once, counting distinct
points.
\end{lemma}
\begin{proof}
Choose four low digits
\[
 b_0=0,\quad b_1=1,\quad b_2=64L+2,\quad b_3=(64L)^2+3 .
\]
They represent the four residues modulo four. If all three chosen
digits are distinct, the largest exceeds \(4L\) times the next
largest, so its term cannot be canceled by the other two. If two
digits coincide, the zero sum of the coefficients forces all three to
coincide. Let \(D=32L(b_3+1)\), and let \(X=\lfloor s/(16LD)\rfloor\).
We first construct a set \(S'\subset[0,X)\).

Put \(p=\lfloor\sqrt{\log X/\log(2L+2)}\rfloor\),
\(M=\lfloor X^{1/p}\rfloor\), and \(H=\lfloor M/(8L)\rfloor\).
For sufficiently large \(s\) we have \(p,H\ge2\). Encode vectors
\((z_0,\ldots,z_{p-1})\) with \(0\le z_i<H\) by
\(\sum_i z_i M^i<X\). Among the at most \(1+pH^2\) values of
\(\sum_i z_i^2\), some sphere contains at least \(H^p/(1+pH^2)\)
vectors. Use these for \(S'\), following Behrend's sphere
construction.~\cite{I-I-I-behrend48}
An integer relation with the stated coefficients gives the same
relation in each digit. Indeed reducing modulo \(M\) forces the
lowest-digit combination to vanish since its absolute value is
\(<6LH<M\); division by \(M\) repeats the argument. One of the three
vectors is then a strict convex combination of the other two unless
they coincide. On a sphere this is impossible, by strict convexity
of the squared Euclidean norm.

For large \(s\), \(H\ge X^{1/p}/(32L)\). The logarithm of the loss
from \(X\) to \(H^p/(1+pH^2)\) is at most
\(p\log(32L)+2\log M+\log(2p)\), which is
\(O(\sqrt{\log X\,\log(2L+2)})\) by the choice of \(p\).
Put \(S=\{Dv+b_i:v\in S',\ 0\le i<4\}\). All these integers are
distinct and \(<s/(8L)\) for large \(s\). A relation reduced modulo
\(D\) has a low-digit combination of magnitude at most \(6Lb_3<D\),
so it vanishes over the integers. The low digits coincide, and then
the high digits coincide. The residues are exactly balanced.
The factor \(s/X=O(L^4)\) is absorbed in the displayed loss.

For the orbit assertion, suppose three distinct points with indices
\(i_0,i_1,i_2\) occur in a segment of \(2L\) consecutive indices,
and subtract the translate from their values to obtain representatives
\(s_0,s_1,s_2\in S\). Eliminating the common step gives
\[
 (i_2-i_0)(s_1-s_0)\equiv(i_1-i_0)(s_2-s_0)\pmod{s}.
\]
The difference has absolute value \(<s\), so the congruence is an
integer equality. Its three coefficients are nonzero differences
of indices, each smaller than \(2L\), a contradiction. For an orbit
of order \(b_0'\le8L\), distinct elements have circular separation
at least \(s/b_0'\ge s/(8L)\), larger than the diameter of \(S\).
\end{proof}

\begin{theorem}[The exact prime sum with a sparse distinguished set]\label{thm:sparse48}
Fix the fields, \(a,b,B',W\) and the \(K_d\)-orientation of
Theorem~\ref{thm:orientation46}. Take \(m,k\) as in
Proposition \ref{prop:ranks47},
and put \(a_E=xb\). For each \(x\) and any integer \(L\) in the range
of Lemma~\ref{lem:sphere48}, with \(s=|A_0|\), one can choose
\(T_2\subset A_0,\ T\subset A\) as follows. The integer representatives
of \(T_2\) satisfy that lemma, \(h_2=|T_2|\), \(\delta=h_2/n\ge x^{-\eps}\),
and the proportions of \(T\) in \(A/A_0\) are exactly those of
condition~\ref{cond:chars}. Set \(y=a/B'\), \(Y'=a/W\),
\(H=A/A_0\), \(w_{\mathbf1}=1\) and \(w_\eta=1/(3o+1)\) for
\(\eta\in\widehat H\) nontrivial. With \(\alpha_{ij\eta}\) as in
Theorem~\ref{thm:orientation46}, at every \(u\in A_0\) one has
\begin{align}
 \frac{\Delta_u}{hm}
 &=W\bigl(G_{\mathbf1}(Y')-G_{\xi_d}(Y')\bigr) \notag\\
 &\quad+\frac{\delta}{2}\sum_{\eta\in\widehat H}w_\eta
       \sum_{i,j=0}^1(-1)^j
       \bigl(B'G_{\alpha_{ij\eta}}(y)-W G_{\alpha_{ij\eta}}(Y')\bigr).
                                                               \label{eq:sparse48}
\end{align}
In particular \(\Delta_u\) is the same for all \(u\), although the
distinguished sets themselves vary with \(u\).
\end{theorem}
\begin{proof}
Restriction embeds \(A_0=\Gal(K_\psi C_l/K_\psi)\) in the cyclic
group \(\Gal(C_l/\QQ)\), whose order is \(2^l\). Hence \(A_0\) is
cyclic of order \(s=n/o\asymp_\psi x^{2+8\eps}\). Fix a generator
and identify it with \(1\in\ZZ/s\ZZ\). Its restriction to the
cyclotomic layer has the same order. Use the integers of \(S\) for
\(T_2\) in these coordinates. Since \(L\) is at most polylogarithmic,
\(\log(n/h_2)=O_{\psi,C_0,C_1}(\sqrt{\log s\,\log(2L+2)})=o(\log x)\).
This gives the claimed density for fixed \(\eps\).

Include the entire identity coset of \(A_0\) in \(T\). In each
nonidentity coset choose a representative and include exactly the
three nonzero residue classes modulo \(4A_0\) relative to it.
There are four classes of equal size since \(4\mid s\). Thus
\(h=(3o+1)s/4\ge n/2\), and the normalized character sum
\(h^{-1}\sum_{t\in T}\eta(t)\) equals \(w_\eta\) for characters of
\(H\), by orthogonality on \(H\).

For a character \(\theta\) of \(A\) nontrivial on \(A_0\) the product
\[
 \left(\sum_{t\in T}\theta(t)\right)
 \left(\sum_{z\in T_2}\overline{\theta(z)}\right)=0.
\]
If \(\theta\) is nontrivial on \(4A_0\), the first sum vanishes by
translation invariance of \(T\). In the remaining case
\(\theta|_{A_0}\) induces a nontrivial character of \(A_0/4A_0\);
the second sum vanishes by the four equal multiplicities. Translation
by \(u\) simply multiplies the second sum by \(\overline{\theta(u)}\).

With weight \(W\) everywhere, the centered count is the first term
of \eqref{eq:sparse48} at each \(t\). Correcting the weight to \(B'\)
on the distinguished set changes each Fourier summand by
\(B'G(y)-W G(Y')\), by \eqref{eq:count}. In the original product
coordinates the distinguished Frobenius elements are exactly
\((0,1,z+u-t)\), \(z\in T_2\), using the coordinate change in
Theorem~\ref{thm:orientation46}. Their character sum gives the product
above with the further factor \((-1)^j\overline{\theta(u)}\).
Only \(\theta\) trivial on \(A_0\) survives. Then its normalized
first factor is \(w_\eta\), and its second is \(h_2\). Division by
\(2n\) and summation over the \(m\) identical stages proves the
identity with the primitive representatives at every prime, including
the ramified primes.
\end{proof}

\begin{corollary}[The remaining estimate with these embeddings]\label{cor:window48}
For sufficiently large \(x\) on the selected intervals, uniformly in
the sets just constructed,
\[
 \left|\frac{2\Delta_u}{hm b\delta}-\Gg_\psi(x)\right|
       \le 3o\,x^{1/2+\eps/2}+o_{\psi,\eps,C_0,C_1}(1).
                                                                    \label{eq:rem48}
\]
For fixed \(\psi\), the critical-line assertion for the primitive
powers of \(\psi\) is equivalent to \eqref{eq:gap} on arbitrarily
distant such intervals for every sufficiently small \(\eps>0\).
\end{corollary}
\begin{proof}
At \(Y'=a/W\) every \(WG_{\alpha}(Y')\) in the identity is
\(O(a\exp(-c x^{\eps/8}))\) by Lemma \ref{lem:far14}. Division by
\(b\delta\) still gives \(o(1)\). Replacing
\((B'/b)G_\alpha(y)\) by \(G_\alpha(x)\) costs \(o(1)\) for each
surviving character. In fact \(|y-x|=O(x/b)\); the slope bound in
Theorem~\ref{thm:orientation46} applies to the auxiliary twists, and
the retained characters have fixed conductor with slopes
\(O_\psi(\log(2x))\) including the principal correction. Their values
in this range are \(O_\psi(x\log(2x))\), also for the auxiliary twists
by the bound on their corrections. The terms \(i=j=0\) now give
\(\Gg_\psi(x)\) by conjugate pairing. The other \(3o\) terms have
weights of absolute value at most one and obey Theorem~\ref{thm:twists}.

Under the critical-line assertion the explicit formula gives
\(\Gg_\psi(x)=O_\psi(\sqrt x)\). Since \(\delta\le1\),
\(\mathcal B/(hmb)\ge x/m\ge\sqrt x\), and \(5\eps>\eps/2\),
\eqref{eq:rem48} implies \eqref{eq:gap} for all sufficiently large \(x\).
Conversely \(k=2n,\ h\ge n/2,\ m\sim\sqrt{x}\) give
\(\mathcal B/(hmb)=O_\psi(x^{1/2+\eps})\) in \eqref{eq:gap}.
Using \(\delta\ge x^{-\eps}\), the assumed inequality and
\eqref{eq:rem48} imply \(\Gg_\psi(x)\ll_{\psi,\eps}x^{1/2+7\eps}\)
on the whole interval. The interval theorem and functional equation
used in Proposition~\ref{prop:strength26}, with \(\eps\) arbitrarily
small, give the critical-line assertion.
\end{proof}

The sparse sets satisfy the short-orbit and three-point restrictions
throughout the cohomology intervals. Equation~\eqref{eq:sparse48}
computes their weighted counts for every translate. An upper bound
for those counts from the specified projected periods, and the strict
local cohomology bound on these fields, are still required.

\subsubsection{The numerical field degrees, ranks and densities}\label{nav:I-31}
We now impose the growth parameters of
\srcloc{Development/AbelianVarietyTheory/Definitions/GlobalTwisting/PairRowIffQdRows}{28}
and the cardinalities and auxiliary-prime window at lines 78--88
there. The absolute field degree in these conditions is \(4n\). The
constant factor in \(n=[K_\psi C_l:\QQ]\) matters for its prescribed
upper bound of \(2x^{2+8\eps}\).

\begin{proposition}[Choosing the field degree and transferring strict classes]\label{prop:desc50}
On the same selected intervals put \(h_0=2+8\eps\) and take the least
index \(l'\) with \([K_\psi C_{l'}:\QQ]\ge x^{h_0}\). Let
\(\widetilde E=K_\psi C_{l'}(\sqrt{-q},\sqrt{-d})\) and
\(\widetilde n=[\widetilde E:\QQ]/4\). For large \(x\),
\[
 x^{h_0}\le\widetilde n<2x^{h_0},\qquad
 \widetilde E\subset E_x,\qquad [E_x:\widetilde E]=2^e\le 2o .
\]
At every odd prime \(\ell\) and every \(r\ge1\), restriction identifies
\(\mathcal S_{\ell,r}(\widetilde E)\) with the
\(\Gal(E_x/\widetilde E)\)-invariants of \(\mathcal S_{\ell,r}(E_x)\).
At two its kernel has \(\ZZ_2\)-length at most \(e\). Consequently the
full weighted sum of strict lengths up to \(r<t=\lfloor x^{401}\rfloor\)
over \(\widetilde E\) is at most that over \(E_x\) plus
\(t\log(2o)\).
\end{proposition}
\begin{proof}
The increasing intersections \(K_\psi\cap C_j\) stabilize, since the
cyclic field \(K_\psi\) has finitely many subfields. The degrees of
\(K_\psi C_j\) consequently double at each sufficiently large step.
Minimality gives the two bounds on \(\widetilde n\). The index used
for \(E_x\) already meets this lower bound, since \(2^l\ge x^{h_0}\).
The ratio of the two degrees is at most \(2o\) since
\([K_\psi C_l:\QQ]\le o2^l<2o x^{h_0}\).
The independent quadratic factors are disjoint from this tower
by ramification at \(q,d\). The relative group is cyclic of the
displayed order because it embeds in the cyclic cyclotomic layer.

Write \(M=\QQ_\ell/\ZZ_\ell(-r)\). A strict class is an absolute
Galois \(H^1\)-class unramified outside \(\ell\) at finite places
and zero locally above \(\ell\). Its image by restriction is strict.
At odd \(\ell\) the cohomology in positive degrees of the finite
relative group on any \(\ell\)-primary module vanishes, by averaging.
Inflation--restriction therefore identifies the full \(H^1\) downstairs
with the invariants upstairs. If the upstairs class is locally zero
above \(\ell\), its descended restriction is locally zero by injectivity
of restriction for each local extension of two-power degree. At any
other prime the same argument on inertia proves that the downstairs
class is unramified. This establishes the assertion for strict classes.
At two the kernel is contained in \(H^1(C_{2^e},M^{G_{E_x}})\).
For a cyclic group with generator \(\sigma\) its first cohomology is
\(\ker(1+\cdots+\sigma^{2^e-1})/(\sigma-1)M^{G_{E_x}}\).
It is killed by \(2^e\) and is a subquotient of a subgroup of the
quasicyclic group \(\QQ_2/\ZZ_2\). Such a subquotient of bounded
exponent is cyclic, as every finitely generated subgroup is cyclic
and the orders of its elements are bounded. Its length is at most
\(e\). Summation gives the result, including its uniform additive
cost at two.
\end{proof}

From now on write \(E=\widetilde E,\ n=\widetilde n,\ f=2n\),
\(A=\Gal(K_\psi C_{l'}/\QQ)\), and \(A_0=\Gal(K_\psi C_{l'}/K_\psi)\).
All analytic large-argument bounds of Lemma \ref{lem:far14}, pass to
these fields because their characters are characters of \(E_x\) and
the primes \(q,d\) are the same. The sparse embeddings, cohomology
groups and polarized abelian schemes below all use this one field.

\begin{theorem}[Simultaneous sparse cardinalities at the specified ranks]\label{thm:rangenew50}
Set \(\delta_1=x^{-\eps/8},\ b=\exp(x^\eps),\ L_0=x^\eps\), and choose
\[
 \begin{gathered}
 m=2\lceil\sqrt{x\delta_1}/2\rceil,\qquad
 k=2\lceil x^{9/2+10\eps}/2\rceil,\qquad r_0=k+m-1,\\
 n_W=8\left\lceil\frac{10 r_0}{8}
                  \left(f+\frac{f^2(k-1)^3L_0^2}{m}\right)\right\rceil,
 \qquad
 t_{\rm r}=8fr_0+4+\lceil x^{377}\rceil n_W-r_0,\\
 L=10\lceil405\log x\rceil+2\lceil30/\eps\rceil+22 .
 \end{gathered}
\]
The rank and signature construction and residue-character conclusions
of Section \ref{nav:V-06}
hold for these ranks, with \(t_{\rm r}<x^{401}\) and
\(N_m=k+m(t_{\rm r}+1)<x^{402}\). The sets in
Theorem~\ref{thm:sparse48} can be chosen to satisfy simultaneously
\[
 \tfrac12\delta_1 n\le h\le2\delta_1 n,\qquad
 h_2\ge n x^{-\eps},\qquad
 \nu(\alpha)=\frac{3+\mathbf1_{\alpha=0}}{3o+1}.
\]
The short-orbit and three-point conclusions of Lemma~\ref{lem:sphere48}
hold at this \(L\). The exact identity \eqref{eq:sparse48} holds for
every \(u\) at these new densities.
\end{theorem}
\begin{proof}
The rounded \(k,m\) are even, \(r_0\) is odd and \(8\mid n_W\), so
\(t_{\rm r}\) is odd. All the lower rank inequalities of Section \ref{nav:V-06}, follow from \(2n\le k,\ m^2\le k\) and
\(t_{\rm r}>3k\) for large \(x\). Also
\(xm\le k\le n^2x^{1/2-\eps}\) and \(q,d\) each exceed
\(4t_{\rm r}r_0+1\) eventually, using the fixed initial power window
as in that section. For the upper bounds, the leading power of \(x\)
in \(r_0 f^2(k-1)^3L_0^2/m\) is
\(43/2+(58+1/16)\eps<23\). Thus \(n_W=O(x^{23})\), and the
formula gives \(t_{\rm r}=O(x^{400})\). The proofs of the rank
embeddings, residue characters, Haar norms and Fourier denominators
in Volume~XIII use the dimension formulas for the towers for arbitrary
such ranks. The logarithms of the finite costs there are still
polynomial in \(x\), since \(f,m,N_m\) are. The same observation shows
that \(G=f r_0(r_0+t_{\rm r})\) in
Corollary \ref{cor:jet49}
is polynomial in \(x\).

Let \(s=|A_0|=n/o\), a power of two. Choose the least power of two
\(D\ge2/\delta_1\). Then \(2/\delta_1\le D<4/\delta_1\) and \(D\mid s\)
for large \(x\). Choose \(T\) to consist of four residue classes modulo
\(DA_0\) in the identity coset and three in each other coset. Its size
is \(h=(3o+1)s/D\), which satisfies the two bounds and the asserted
proportions. We need a large three-point-free set exactly balanced
in all \(D\) residues.

For \(0\le v<D\) put \(c_v=v+8LD\,v^2\). These represent the \(D\)
residues. Suppose \(\alpha c_u+\beta c_v+\gamma c_w=0\) with the
coefficients of Lemma~\ref{lem:sphere48}. Reduction modulo \(8LD\)
gives \(\alpha u+\beta v+\gamma w=0\) over the integers because
the absolute value is \(<6LD\). The original relation then gives the
same identity in \(u^2,v^2,w^2\). If \(u,v,w\) are distinct, the
Vandermonde matrix with rows \(1,z,z^2\) has nonzero determinant.
This contradicts the nonzero coefficient triple. With two equal
entries the relations force the third equal also.

Let \(M\) be the least multiple of \(D\) exceeding \(8L(1+\max c_v)\);
thus \(M=O(L^2D^3)\). Take \(X=\lfloor s/(16LM)\rfloor\), and construct
the high digits \(S'\subset[0,X)\) by the sphere proof of
Lemma~\ref{lem:sphere48}. Its hypotheses hold because
\(\log X\asymp_\psi\log x\) and \(L=O_\eps(\log x)\).
Set \(T_2=\{Mz+c_v:z\in S',\,0\le v<D\}\) in the generator coordinates
of \(A_0\). Distinct pairs give distinct integers \(<s/(8L)\).
Reduction modulo \(M\) first forces the low-digit relation over
the integers, since \(6L\max c_v<M\). The preceding argument makes
the low digits coincide, and the sphere argument then makes the
high digits coincide. The orbit conclusions follow verbatim from
the elimination of the common step in Lemma~\ref{lem:sphere48}.
Each residue modulo \(D\) occurs \(|S'|\) times, and
\[
 \frac{h_2}{n}\gg_\psi
    \frac{\exp(-C\sqrt{\log s\,\log(2L+2)})}{L^3D^2}
    \ \ge\ x^{-\eps/4-o(1)}\ \ge\ x^{-\eps}
\]
for large \(x\). A character of \(A\) nontrivial on \(DA_0\) sums to
zero on \(T\). One trivial there but nontrivial on \(A_0\) sums to
zero on \(T_2\). These are precisely the two cancellations needed
in the proof of \eqref{eq:sparse48}, which now applies at every \(u\).
\end{proof}

\begin{proposition}[The two auxiliary primes in their prescribed interval]\label{prop:ray50}
Let \(N=\operatorname{lcm}(F,2^{l'+2})qd\) and \(U=(kb)^8\).
For every sufficiently large \(x\) on the selected intervals there
are distinct primes \(g_1,g_2\in[U,2U]\), both \(1\bmod N\), with
\(g_i=\pi_i\bar\pi_i\) in \(\mathcal O_K\), \(K=\QQ(\sqrt{-d})\).
In fact the interval contains \(U^{1-o(1)}\) such primes. One may
also choose two distinct smaller primes \(1\bmod 15N\) of size
\(O(N^{C})\) for an absolute \(C\). All four split completely in \(E\).
\end{proposition}
\begin{proof}
Write \(H_K\) for the Hilbert class field of \(K\), \(C_N=\QQ(\zeta_N)\)
and \(L=H_K C_N\), a Galois extension of \(\QQ\). The class group of
\(K\) has odd order. To see this, conjugation acts as inversion. If the
class of an ideal \(I\) has order dividing two, then \(\bar I/I=(a)\)
with \(a\bar a=1\), because this is a positive rational unit. Write
\(a=z/\bar z\), taking \(z=1+a\) unless \(a=-1\), in which case take
\(z=\sqrt{-d}\). The fractional ideal \(zI\) is conjugation invariant.
At a split prime its two exponents agree, and an inert prime itself is
rational. The only ramified prime is generated by \(\sqrt{-d}\).
Consequently every such invariant ideal is principal, proving the claim.
Artin reciprocity identifies \(\Gal(H_K/K)\) with the class group and
identifies conjugation on it with inversion. Since the squaring map on
this odd group is onto, commutators with complex conjugation generate
the whole group. Thus the largest subextension of \(H_K/\QQ\) abelian
over \(\QQ\) is \(K\). As \(K\subset C_N\), it follows that
\(H_K\cap C_N=K\) and
\(\Gal(L/K)=\Gal(H_K/K)\times\Gal(C_N/K)\).

Apply Thorner and Zaman's Chebotarev theorem to \(L/\QQ\), identity
conjugacy class, with abelian subgroup \(\Gal(L/K)\).
Its maximal relative character conductor norm is at most \(N^2\).
Indeed the first factor is unramified. A character on the second
factor extends to the abelian group \(\Gal(C_N/\QQ)\); on \(K\) it
is obtained by the local norm from a Dirichlet character of conductor
dividing \(N\). The norm sends \(1+N\mathcal O_{K,v}\) to
\(1+N\ZZ_p\), as follows by expanding the local norm or at a split
place by multiplying the two coordinates.

The only possible exceptional zero in the relative theorem is simple,
real, and attached to a real character of \(\Gal(L/K)\). Such a
character kills the odd class group, so its zero would also be a zero
of \(\zeta_{C_N}\). The other Hecke factors of \(\zeta_L\) are
holomorphic there. Factoring \(\zeta_{C_N}\) over \(\QQ\), a zero of
a nonreal Dirichlet character at a real argument is repeated by its
conjugate and therefore cannot give this simple zero. Every real
primitive character appearing is a product of powers of \(\xi_q,\xi_d\)
and a real character with conductor dividing \(8F\). This follows from
the square subgroup of the two-adic units and the fact that \(q,d\)
are distinct odd primes prime to \(2F\). The finitely many base
characters have a fixed zero-free interval next to one. The selected
auxiliary twists have no zero in \([49/50,1]\) by
Theorem~\ref{thm:twists}. Hence the exceptional term, if present, has
its zero bounded away from one and is negligible here.

The theorem gives, uniformly for \(Z=U,2U\),
\[
     \pi_{\rm split}(Z)=\frac{\operatorname{Li}(Z)}{[L:\QQ]}(1+o(1)).
\]
In fact its two relative errors are bounded by constants times
\(\exp(-c\log Z/\log(4dN^2))\) and
\(\exp(-c\sqrt{\log Z})\), since the base degree is two.
Here \(\log(dN^2)=O_{\psi,\eps}(x^{\eps/5}\log^2(2x)+\log(2x))
=o(\log U)\). Subtract the estimates at the two endpoints.
Minkowski's bound represents each ideal class of \(K\) by an integral
ideal of norm \(O(\sqrt d)\). There are \(O(d)\) such ideals since
\(\zeta_K(2)\le\zeta(2)^2\). Thus \([L:\QQ]\ll dN=U^{o(1)}\), and
the difference is \(U^{1-o(1)}\). A prime counted splits in \(C_N\),
so it is \(1\bmod N\), and its prime ideals over \(K\) are principal
by splitting in \(H_K\). A generator of either has norm \(g\).
The two conjugate ideals are coprime. For the two smaller primes apply
Linnik's least-prime bound in the progression \(1\bmod 15N\), and
then in the progression \(1\bmod 15Np_1\), where \(p_1\) is the
first prime. Both are polynomial in \(N\), and are below \(U\) for
large \(x\). Finally \(E\subset C_N\).~\cite[Thm.~1.4 and
Thm.~3.1]{I-I-I-tz50}\cite[Thm.~1.1]{I-I-I-tzleast50}
\end{proof}

\begin{corollary}[Stability under rounding the orders and weights]\label{cor:actual50}
At the fields, ranks and densities just chosen, give each \(t\in T\)
an integer order \(a_t\) with \(|a_t/b-x|\le1\). At each stage and
translate suppose the distinguished weights lie in \([b,b+\sqrt b]\)
and all other weights \(w\) satisfy
\[
           1\le w,\qquad a_t/w\ge\tfrac13\exp(x^{\eps/3}).
\]
Then \(\left|2\Delta_u/(hm b\delta)-\Gg_\psi(x)\right|
\ll_{\psi,\eps}x^{1/2+\eps/2}\) for every \(u\).
The upper estimate \eqref{eq:gap} on arbitrarily distant such intervals
for every sufficiently small \(\eps>0\) would imply Dirichlet GRH for
the retained characters.
\end{corollary}
\begin{proof}
First replace the weights at fixed \(a_t\) by the two values \(B',W\).
The large-argument bound on the nondistinguished terms and the slope
bound on the distinguished terms in Proposition \ref{prop:weights14}, cost
\(O(b/x)\) per stage and translate for either choice of weights. The
same proof applies at any \(a_t\) in the stated range and bounds the
difference via the common idealized sum. Now hold the weights \(B',W\)
fixed and change \(a_t\) to \(a=\lceil bx\rceil\). In the prime-power
formula before Fourier expansion the slope in \(a_t\) is bounded by
\(2\sum_{v\le (a+2b)/W}\Lambda(v)/v\): at each prime power the
normalized multiplicity of a Frobenius coset is at most two. This is
\(O(\log(2a))\) by the factorial estimate in the proof of
Theorem~\ref{thm:twists}. The slope of \(M_d\) has the same bound
because \(1\le w\le B'\), \(D_{\xi_d}=O(\log(2x))\), and the real
order stays within \(2b\) of \(a\). Thus the change costs
\(O(b\log(2a))\) per stage and translate, with no factor \(n\).
After normalizing, the error is
\(O(\delta^{-1}\log(2a))=O_{\psi,\eps}(x^{2\eps})\).
Equation~\eqref{eq:rem48} applies to the constant orders and weights
by the exact cancellation in Theorem~\ref{thm:rangenew50}. This
proves the first assertion.

For these parameters \(h\ge\delta_1 n/2,\ h_2\le n\),
\(m\asymp x^{1/2-\eps/16}\), and \(k\asymp x^{9/2+10\eps}\).
Substitution in \eqref{eq:gap} gives
\(\mathcal B/(hmb)=O_\psi(x^{1/2+17\eps/16})\).
Indeed its three contributions after division are
\((x/m)(1+m^2/k)\), \(2nm x^\eps/h\), and
\(h_2 k x^{2\eps}/(h n^2)\).
Using \(\delta^{-1}\le x^\eps\), the assumed upper estimate gives
\(\Gg_\psi(x)\le x^{1/2+8\eps}\) for large \(x\) on the interval.
The interval theorem and the primitive functional equation give the
critical-line assertion as \(\eps\) decreases to zero.
\end{proof}

At \(g_1,g_2\) the polarized construction now has full Fourier
denominator with logarithm at most
\(f r_0(r_0+t_{\rm r})\log(g_1g_2)=o(b/x)\), since the ranks
and \(\log U\) are polynomial in \(x\). The primes may vary with
\(x\). There are only finitely many fields and possible primes on
each bounded interval, so one can take a common coefficient field for
the geometric calculations of Volume~XIII. Proposition~\ref{prop:desc50}
transfers the known cohomology control to precisely these fields and
quantifies its dyadic cost. Neither the remaining odd-primary strict
lengths nor the comparison of the specified normalized row and
column periods is estimated by these arguments. They are still needed
for the proposed proof of Main Theorem~\ref{thm:grh16}.

\subsubsection{The exact dyadic kernel in the change of field degree}\label{nav:I-32}
Proposition~\ref{prop:desc50} loses at most \(e\) dyadic length per
twist when the two-power extension has degree \(2^e\). We compute the
kernel exactly. The loss can occur for a fixed conductor at every
sufficiently large scale.

\begin{theorem}[Dyadic strict descent on the prescribed fields]\label{thm:sharp51}
Keep the fields \(E'=K_\psi C_{l'}(\sqrt{-q},\sqrt{-d})\subseteq
E=K_\psi C_l(\sqrt{-q},\sqrt{-d})\) of
Proposition~\ref{prop:desc50}. At a dyadic place write \(E_w\) for the
completion of \(E\). For all sufficiently large scales and all integers
\(r\ge1\), the kernel of the restriction
\[
        \mathcal S_{2,r}(E')\longrightarrow\mathcal S_{2,r}(E)
\]
is cyclic of order two precisely when
\[
 l>l',\qquad r\ {\rm odd},\qquad i\notin E,\qquad i\in E_w,
                  \qquad i^2=-1.                              \label{eq:roots51}
\]
It is zero in all other cases. The condition on \(E_w\) has the same
answer at all dyadic places. When the levels differ and the two
conditions on \(i\) hold, the sum of weighted kernel lengths for
\(1\le r<t\) is exactly \(\lfloor t/2\rfloor\log 2\). Otherwise it is
zero. The odd-prime restriction remains the invariant-group
identification of Proposition~\ref{prop:desc50}.
\end{theorem}
\begin{proof}
We first describe the cyclotomic images, using a superscript \(+\) on
a two-power cyclotomic field for the fixed field of inversion on roots
of unity. At the global place this is its maximal real subfield.
In the local argument the same notation denotes the corresponding
totally ramified subfield of the two-power extension of \(\QQ_2\).
Let \(T_\infty\) denote the union of these full cyclotomic fields,
globally or locally as appropriate. At every sufficiently large
level \(j\), the intersection of the corresponding field
\(E_j=K_\psi C_j(\sqrt{-q},\sqrt{-d})\) with the global \(T_\infty\)
is either \(C_j\) or \(\QQ(\zeta_{2^{j+2}})\). Its dyadic completion
has intersection with the local \(T_\infty\) equal to the analogous
one of the two local fields.

For the global assertion, a character of \(E_j\) on a pure two-power
conductor has trivial exponents on both auxiliary quadratic characters:
each ramifies at its own prime \(q\) or \(d\), whereas \(K_\psi C_j\)
is unramified there. Hence adjoining those two square roots adds
nothing to the two-power cyclotomic intersection. For a fixed abelian
base field \(B\) and \(C_j\subset T_\infty\), the intersection
\((BC_j)\cap T_\infty=(B\cap T_\infty)C_j\). Indeed the Galois
action on \(T_\infty\) fixing \(B\) has image
\(\Gal(T_\infty/(B\cap T_\infty))\); imposing also that it fix
\(C_j\) intersects this subgroup with \(\Gal(T_\infty/C_j)\).
Take \(B=K_\psi\). Its fixed finite intersection with \(T_\infty\)
lies in the full level \(2^{j+2}\) for every sufficiently large \(j\).
The latter field has degree two over \(C_j\); the larger possibility
occurs exactly when \(i\in E_j\).

Locally take \(B\) to be the compositum of a dyadic completion of
\(K_\psi\) and the unramified quadratic extension of \(\QQ_2\).
Indeed \(K_q\) gives this unramified quadratic field and \(K_d\)
splits at two throughout the interval. The local base is thus fixed
even as \(q,d\) vary. The same Galois argument proves the local
intersection assertion. At both levels \(j=l',l\) the choice between
the two alternatives is constant once \(l'\) is large: it is determined
by whether the fixed base intersection is contained in the fixed
field of inversion in \(T_\infty\). Locally the relative degree
from \(l'\) to \(l\) is \(2^{l-l'}\). In fact the intersections of
the fixed local base with the increasing \(C_j\) stabilize, so taking
the compositum preserves the full relative degree. This equals the
global relative degree. Thus at every dyadic place the decomposition
group in \(E/E'\) is the entire group \(G\), and \(G\) is cyclic.
All assertions are uniform in the choice of the two auxiliary primes.

Write \(M=\QQ_2/\ZZ_2(-r)\), and regard \(M^{G_E}\subseteq M^{G_{E_w}}\)
as subgroups of the same torsion group. The global and local
cyclotomic images at level \(j\) are either \(1+2^{j+2}\ZZ_2\),
when the field contains \(i\), or
\(\{\pm1\}(1+2^{j+2}\ZZ_2)\), otherwise. Under the first image
the invariants are \(M[2^{j+2+v_2(r)}]\). This follows by taking
a generator \(1+2^{j+2}\) and applying the binomial expansion, which
gives \(v_2((1+2^{j+2})^r-1)=j+2+v_2(r)\). Inverting the generator's
value preserves that valuation. Under the second image one imposes
in addition the action of \((-1)^r\). It makes no change for even
\(r\) and gives exactly \(M[2]\) for odd \(r\).

By inflation--restriction the absolute global restriction kernel is
\(H^1(G,M^{G_E})\). This group maps to the absolute local restriction
kernel \(H^1(G,M^{G_{E_w}})\) by the inclusion of the coefficients.
A class in the former group is automatically unramified outside two:
the extension \(E/E'\) is unramified there and the inflated cocycle
is zero on inertia. Its membership in the strict kernel is therefore
exactly the vanishing of its image at every dyadic place. If the two
groups of coefficients coincide, this map is the identity under the
common group identification, and the strict kernel is zero. This
includes all even \(r\) and all cases where the two cyclotomic images
agree. If \(l=l'\), restriction itself is the identity.

It remains to consider \(l>l'\), \(r\) odd, \(i\notin E\) but
\(i\in E_w\). The global coefficients are \(M[2]\) with trivial
\(G\)-action. Thus \(H^1(G,M[2])\) has order two. Locally the
coefficients \(A=M[2^{l+2}]\) form a cyclic group on which a generator
of \(G\) acts by a unit \(z\) with \(v_2(z-1)=l'+2\).
Indeed the relative action on the full two-power cyclotomic layers
is generated by \(1+2^{l'+2}\) modulo \(2^{l+2}\), and raising to
the odd power \(-r\) preserves this first valuation.
Let \(e=l-l'>0\). In the cyclic cohomology formula the norm operator
acts by \(1+z+\cdots+z^{2^e-1}\), of valuation \(e\) by the same
binomial identity. Its kernel on \(A\) has order \(2^e\), equal to
the order of \((z-1)A\), which it contains. Hence \(H^1(G,A)=0\).
The global class vanishes locally at every dyadic place. The local
test is independent of \(w\) since \(E/\QQ\) is abelian. This proves
the formula and summing over the odd \(r<t\) gives the weighted total.
\end{proof}

\begin{corollary}[A conductor attaining the dyadic loss]\label{cor:example51}
Let \(\psi\) be the primitive real quadratic character of conductor
\(28\) cutting out \(\QQ(\sqrt7)\). On every sufficiently distant
common interval the strict restriction kernel in
Theorem~\ref{thm:sharp51} has order two for every odd \(r\ge1\).
\end{corollary}
\begin{proof}
The quadratic field \(\QQ(\sqrt7)\) ramifies at seven and meets the
global two-power cyclotomic extension trivially. Therefore
\([K_\psi C_j:\QQ]=2^{j+1}\), so \(l'=l-1\) for all large \(x\) by
the two minimality definitions. The global intersection in the proof
of the theorem is \(C_l\), which excludes \(i\). But \(-7\equiv1\bmod8\)
is a square in \(\QQ_2\), so \(\QQ_2(\sqrt7)=\QQ_2(i)\). Thus every
condition in the theorem holds for odd \(r\).
\end{proof}

This determines the loss in descending strict classes even when the
local and global two-power roots of unity differ. The remaining
large odd-prime strict groups and the normalized global periods are
the two estimates in Section~\ref{sec:missing} still needed for
Main Theorem~\ref{thm:grh16}.

\subsubsection{A simultaneous divisor for projected theta sections}\label{nav:I-33}
Fix one of the intervals \(I=[V,V^{3/2}]\) of
Section \ref{nav:II-15}.
Use throughout the fields \(E=\widetilde E_x\) of
Proposition~\ref{prop:desc50}, the sets of Theorem~\ref{thm:rangenew50}
and the two auxiliary primes of Proposition~\ref{prop:ray50}. Put
\(K=K_d\). On each field take the trace-polarized abelian scheme
\(\mathcal A=M\otimes_{\mathcal O_K}C_d\) with \(\mathfrak c=\mathcal O_K\),
and its rigidified symmetric ample bundle \(P\), as in
Theorem \ref{thm:pol49}.
All coefficient fields below may be enlarged by a single finite
extension on \(I\). We first identify the primes at which several
CM directions can contribute to the same divisibility estimate.

\begin{lemma}[Primes with the specified quadratic signs]\label{lem:primes53}
Let \(F\) and \(o\) be the conductor and order of \(\psi\), respectively,
and \(N=\operatorname{lcm}(2^{l'+2},F)qd\). Put
\[
 \mathcal P_x=\{5x/4<p\le7x/4:\ p\text{ prime},\ p\nmid N,\ 
           \psi(p)=1,\ \xi_q(p)=1,\ \xi_d(p)=-1\}.
\]
Uniformly for \(x\in I\) as the left endpoints tend to infinity,
\[
        \sum_{p\in\mathcal P_x}\log p=\frac{x}{8o}+o(x),
        \qquad |\mathcal P_x|\sim\frac{x}{8o\log x}.
\]
For every such prime, the permutation \(\sigma_p\) of
\(\Sigma=\operatorname{Hom}_K(E,J)\) given by
\(\tau\mapsto\tau\circ\Frob_p^2\) has no fixed point.
\end{lemma}
\begin{proof}
Write \(c(n)=\Lambda(n)\) when \(n\) is prime to \(N\) and satisfies
the same three character conditions as \(p\), and put \(c(n)=0\)
otherwise. Write \(S(y)=\sum_n c(n)(y/n-1)_+\).
On the units modulo \(N\), the indicator of the three conditions is
\[
     \frac{1}{4o}\left(\sum_{j=0}^{o-1}\psi(n)^j\right)
                         (1+\xi_q(n))(1-\xi_d(n)).
\]
The interval construction chooses \(q,d\) with
Theorem~\ref{thm:twists} on \([Z,Z^{K_0+1}]\), where \(Z\le V\)
and \(V^2\le Z^{K_0}\) for a fixed \(K_0\). Thus
\([x,2x]\subset[Z,Z^{K_0+1}]\) throughout \(I\), and \(q,d>2x\)
for large \(Z\). When each character in the expansion is replaced
by its primitive representative, the correction to the ramp sum
on this range comes from powers of primes dividing \(2F\).
For one such prime \(v\) and character \(\chi\), its omitted
contribution has the form \(c_{v,\chi}y+O_F(\log(2y))\).
Indeed the terms involving \(y/v^j\) extend to a convergent
geometric series with tail \(O_F(1)\) after multiplication by \(y\),
and the remaining bounded summands number \(O_F(\log(2y))\).

The only principal character in the expansion is obtained by taking
the principal term in all three factors. Every nontrivial quadratic
twist satisfies \eqref{eq:tw11}. The remaining primitive characters
are fixed and satisfy \(G_\chi(y)=o_\psi(y)\) by
Proposition~\ref{prop:allzeros}, including the character of conductor one.
It follows from \eqref{eq:Gchi} that on \([x,2x]\)
\[
       S(y)=\frac{y\log y}{4o}+C_Iy+O_\psi(y\eta_x),
              \qquad \sup_{x\in I}\eta_x\longrightarrow0 .
\]
The real constant \(C_I\) includes the logarithmic derivatives and
the contributions at primes dividing \(2F\). Its size will cancel.
The function \(S\) is convex and satisfies
\(yS'_+(y)-S(y)=\sum_{n\le y}c(n)\), also at integers. Bound
\(S'_+(y)\) by the secants to \((1-\delta)y\) and \((1+\delta)y\)
at \(y=5x/4,7x/4\), with
\(\delta=\sqrt{\eta_x}+1/\sqrt{\log x}\). The endpoints lie in
\([x,2x]\) for sufficiently large \(V\). For the displayed main
term each secant differs from its derivative by \(O_\psi(\delta)\).
The two errors in forming each secant contribute \(O_\psi(\eta_x/\delta)\).
After subtracting \(S(y)\) this gives
\(\sum_{n\le y}c(n)=y/(4o)+o(x)\) at both points. The contribution
from powers \(p^j\le2x\) with \(j\ge2\) is
\(O(\sqrt{x}\log^2(2x))\), by counting at most \(\sqrt{2x}\) possible
bases and \(O(\log x)\) exponents. Subtraction proves the first
asymptotic. Each prime counted has \(\log p=\log x+O(1)\), giving
the second.

Since \(n=[K_\psi C_{l'}:\QQ]\le o\,2^{l'}\) and
\(n\ge x^{2+8\eps}\), for large \(x\) one has
\(1<p^2<2^{l'+2}-1\). Thus \(p^2\) represents neither sign of
one modulo \(2^{l'+2}\). Its Frobenius acts nontrivially on the
real cyclotomic layer \(C_{l'}\). Since \(p\) is inert in \(K\),
\(\Frob_p^2\in\Gal(E/K)\). This nonidentity element acts freely by
translation on the embeddings over \(K\).
\end{proof}

\begin{theorem}[One mixed derivative for all the primes and translates]\label{thm:sim53}
Set \(a=\lceil xb\rceil\) and
\[
      r_x=\left\lfloor\frac{h_2}{8|\mathcal P_x|}\right\rfloor,
      \qquad d_x=\prod_{p\in\mathcal P_x}p^{r_x}.
\]
At every sufficiently large \(x\in I\) there are nonnegative orders
\(\boldsymbol a=(a_\tau)_{\tau\in\Sigma}\) with sum \(a\).
Put \(\mathcal H=\prod_\tau H_d(a_\tau)\), with \(H_d(0)=1\) and
\(H_d\) as in Proposition~\ref{prop:dyad46}. One can take integral
invariant tangent vectors \(D_\tau\) in the corresponding CM
eigenlines, primitive at every place above \(\mathcal P_x\), such
that for every integral scalar jet \(F_0\) through order \(a\) and
every \(u\in A_0\),
\begin{equation}
    \mathcal D_u F_0(0)\ \in\ \mathcal H\,d_x\,\mathcal O_J,
       \qquad \mathcal D_u=\prod_{\tau\in\Sigma}D_{\tau+u}^{a_\tau}.
                                                               \label{eq:whole53}
\end{equation}
The vectors can be fixed for each field on the whole interval.
The additional divisor satisfies
\(\log d_x=(1/8+o(1))h_2\log x\) uniformly there.
\end{theorem}
\begin{proof}
In the coordinates of Theorem~\ref{thm:orientation46} identify the
embeddings with the coset \(j=0\) extending the chosen embedding of
\(K\). The distinguished positions before translation are
\(S_0=\{(1,0,z):z\in T_2\}\), of size \(h_2\).
For each \(p\in\mathcal P_x\) choose \(r_x\) pairs of positions
\((\nu,\sigma_p^{-1}\nu)\), with \(\nu\in S_0\), so that all the
positions in all the pairs are distinct. To do this successively,
if \(j\) pairs have already been chosen, their used positions form a
set \(U\) of size \(2j\). The forbidden centers in
\(U\cup\sigma_p U\) number at most \(4j<h_2\). Any remaining
center gives two different positions by Lemma~\ref{lem:primes53}.
The total number of pairs is at most \(h_2/8\), so the construction
reaches the last prime.

Assign order \(p-1\) to each center for \(p\) and \(p^2\) to its
partner. The total assigned is \(O(h_2x^2)\), which is
\(O_\psi(x^{4+8\eps})<a\) by Proposition~\ref{prop:desc50}.
Put the entire remainder at one unused position and set all
remaining orders equal to zero. Each center contributes \(p-1\)
low digits at that position in base \(p^2\). The low digit of
the \(p^2\)-term moves from its partner to the same center under
\(\sigma_p\). Thus the integer \(L_\nu\) in
Theorem \ref{thm:gain52}
is at least \(p\) at each of these \(r_x\) centers. All other orders
have nonnegative digits. Translating the whole list by \(u\)
commutes with \(\sigma_p\) because \(E/\QQ\) is abelian. The same
bound holds at the \(r_x\) distinct centers of \(S_0+u\).

To choose the vectors, intersect each generic tangent eigenline
with the integral tangent module at the identity. It is a projective
rank-one module over the Dedekind ring \(\mathcal O_J\). At a prime
in \(\mathcal P_x\) the CM algebra is unramified and splits over the
coefficient ring into its integral idempotent factors. The intersection
there is precisely the primitive eigenline of the local calculation.
By the Chinese remainder theorem for this module, one global element
has nonzero reduction at every selected place. The union of the
selected places for all \(x\in I\) is finite. We can make the
choice once for each eigenline and each of the finitely many fields.

At any finite place the formal group in the scalar CM coordinates
is a power of the formal group of \(C_d\). Indeed \(M\) is projective
over \(\mathcal O_K\) and free after localization at the scalar prime.
The scalar factorial bound of Proposition~\ref{prop:dyad46} shows
that \(D_\tau^c/H_d(c)\) preserves integral formal functions for
every \(c\). To pass from evaluation at the identity to this
operator assertion, apply the coefficient calculation to the
function \(F_0(X+_{\mathcal A}Y)\) in \(Y\) with integral power
series coefficients in \(X\). One may choose a formal lift of the
jet locally since only orders through \(a\) enter the final value.
Composing the commuting operators gives divisibility by
\(\mathcal H\) everywhere. At every place above \(p\in\mathcal P_x\),
the valuation of \(\mathcal H\), normalized by \(v_p(p)=1\), is
\(\sum_\tau(v_p(a_\tau!)-F_{p^2}(a_\tau))\). Here
\(F_Q(c)=\sum_{i\ge1}\lfloor c/Q^i\rfloor\). The additional
\(r_x\) from Theorem~\ref{thm:gain52} applies at every such place, proving
\eqref{eq:whole53}. Finally \(h_2\ge nx^{-\eps}\) and
\(|\mathcal P_x|\asymp_\psi x/\log x\). Rounding \(r_x\) therefore
has relative error tending to zero. Summing \(\log p=\log x+O(1)\)
gives the asserted logarithm.
\end{proof}

\begin{theorem}[Integral jets in every translation character]\label{thm:project53}
Write \(g_i=\pi_i\bar\pi_i\) for the two simultaneous primes of
Proposition~\ref{prop:ray50}. Set \(G=g_1g_2\),
\(\varphi=[\pi_1\pi_2]\) on \(\mathcal A\), and \(B=\ker\varphi\).
Take a section \(s_0\) of \(P^3\) with value one at zero in the
rigidification. After one finite extension on \(I\), all characters
\(\lambda\) of \(B(J)\) are defined. Write \(e_\lambda\) for their
translation projectors on sections of \(\varphi^*P^{3a}\).
At every finite place \(w\nmid G\), with integer ring \(R\) in the
completion, the scalar jets through order \(a\) of
\[
        \frac{e_\lambda s}{\varphi^*s_0^a},
       \qquad s\in H^0(\mathcal A,\varphi^*P^{3a})\otimes_{\mathcal O_J}R,
\]
fill the entire integral jet module at the origin, for every
\(\lambda\). At every \(x\in I\) sufficiently large, there is also
a single global integral section \(s\) for which
\begin{equation}
   0\ne
      \frac{1}{\mathcal H}\mathcal D_u
        \left(\frac{G^{2n}e_\lambda s}{\varphi^*s_0^a}\right)(0)
                 \ \in\ d_x\mathcal O_J
             \qquad(\lambda\in\widehat{B(J)},\ u\in A_0).
                                                            \label{eq:project53}
\end{equation}
\end{theorem}
\begin{proof}
The two kernels of \([\pi_i]\) have coprime orders. Each is killed
by \(g_i\) and has order \(g_i^{2n}\) by its index on the lattice
\(M\). Their product is the kernel of \(\varphi\), by the degree
formula and the trivial intersection. After finite extension all
these torsion points and their character values are rational.
Every point extends integrally by properness. Translation acts on
the pullback bundle since \(\varphi\circ t_z=\varphi\) for \(z\in B(J)\).
Consequently
\[
           G^{2n}e_\lambda=\sum_{z\in B(J)}\lambda(z)^{-1}t_z^*
\]
preserves global integral sections. This is the simultaneous
projector for the two individual translation characters after
clearing the displayed denominator.

Over \(R\) at \(w\nmid G\) the isogeny is finite \'etale, with
constant kernel, and is a torsor under that kernel. In fact the
kernel lies in the prime-to-residue-characteristic torsion \(\mathcal A[G]\),
and the generic splitting of
the finite \'etale kernel extends over the normal base.
In \(\varphi_*\mathcal O_{\mathcal A_R}\) the summand of character
\(\lambda\) is an invertible sheaf \(L_\lambda\) on the target.
This follows by trivializing the torsor in the \'etale topology,
where its functions are the regular representation and the
\(\lambda\)-functions are multiples of a single unit character.
The same trivializations give
\(L_\lambda\otimes L_\mu\simeq L_{\lambda\mu}\) by multiplication.
Thus \(L_\lambda^G\) is trivial. By the projection formula the
sections of character \(\lambda\) identify on the target with the
sections of \(P^{3a}\otimes L_\lambda\).

We show that this last bundle supplies all integral jets.
On the geometric generic fiber \(L_\lambda\) lies in
\(\operatorname{Pic}^0\). Indeed its associated homomorphism to
the dual abelian variety is killed by \(G\), and a homomorphism
from a connected abelian variety to a finite group in characteristic
zero is zero. The kernel of the map from line bundles to these
homomorphisms is \(\operatorname{Pic}^0\). The homomorphism
associated with \(P^{3a}\) is an isogeny. Hence there is a geometric
point \(v\) such that \(P^{3a}\otimes L_\lambda\simeq t_v^*P^{3a}\)
on the generic fiber. These properties of the dual and the
polarization are recalled in
\href{https://www.jmilne.org/math/xnotes/AVs.pdf}{Milne,
\emph{Abelian Varieties}, Sections 9--10 and 13}.
There are finitely many orders, fields, pairs of primes and characters
on the bounded interval. Enlarge \(J\) once so these points and
isomorphisms are defined. Each \(v\) extends over \(R\) by properness.
The quotient of the two bundles is then represented by a vertical
divisor. The abelian scheme over \(R\) is regular with one integral
special fiber, the divisor of a uniformizer, so this quotient is
trivial over \(R\) as well.

Thus \(P^{3a}\otimes L_\lambda\) is the \(a\)-th power of the
relatively very ample bundle \(t_v^*P^3\) over \(R\). For clarity,
choose a section nonzero at the origin modulo the maximal ideal.
In the associated affine chart the ratios of the other sections,
after subtracting their values at zero, generate the ideal of the
origin in its local ring because the linear system gives a closed
immersion. Products of at most \(a\) of these functions therefore
span the truncated local ring over \(R\), by induction on the
powers of that ideal. Homogenizing the products with the chosen
section shows that sections of the \(a\)-th power give every jet
through order \(a\). Since \(\varphi\) is \'etale and carries zero
to zero with the same residue field, it identifies the completed
local neighborhoods of the two origins. In a local trivialization
of \(L_\lambda\), passing to the \(\lambda\)-part on the source
multiplies by an invertible function. Changing to the stated frame
also multiplies the entire jet module by a unit. Sections over \(R\)
are obtained by flat base change for \(H^0\), and \(e_\lambda\)
projects onto that integral summand because \(G\) is a unit there.
This proves the surjectivity claimed in the theorem.

In characteristic zero each functional obtained by applying
\(\mathcal D_u\) to this character component is nonzero. In tangent
coordinates dual to the eigenvectors, evaluate on the homogeneous
monomial with the corresponding exponents of total degree \(a\).
All terms involving fewer than \(a\) ordinary derivatives vanish at
zero and the value is \(\prod_\tau a_\tau!\ne0\).
Surjectivity on jets realizes this test. The kernels of all these
functionals form a finite union of proper hyperplanes over the
infinite field \(J\). A vector outside their union exists because
the product of their defining nonzero linear polynomials cannot
vanish everywhere on a vector space over an infinite field.
Choose such a section and clear one common denominator to make it
integral. Every \(G^{2n}e_\lambda s\) is integral. Taking its
scalar jet in the indicated integral frame and applying
Theorem~\ref{thm:sim53} proves \eqref{eq:project53}.
\end{proof}

The two translation characters, the CM lattice and the prime interval
in \eqref{eq:project53} have now been imposed together. The choice of
orders in Theorem~\ref{thm:sim53} used the positions of the Frobenius
pairs. For the periods of the iterated lifts in the \emph{Remark},
the harmonic coefficients and the sum over the global period classes
prescribe a particular section and its differentiations. Establishing
their integral comparison and an upper bound for their archimedean
norms remains necessary for \eqref{eq:gap}. The uniform odd-primary
strict cohomology estimate on these same fields also remains open. For the ordinary signed operators, Proposition~\ref{prop:fulltail621} bounds the discarded real tails on all retained pairs. Theorem~\ref{thm:signedcompact621} gives a compact integral for the full squared norm with the original marked columns and a uniform error. Theorem~\ref{thm:fourkernels622} and Corollary~\ref{cor:gramerror622} extend this approximation to the off-diagonal Gram squares and diagonal rank statistic of the ordinary projected blocks, using a fixed power of the terminal prime for each transported real row cut. Theorem~\ref{thm:markedranks623} bounds the large-rank terms of both signed moments at the diagonal identity normalization. Corollary~\ref{cor:integralranks623} retains this saving on one simultaneous integral choice of marked coordinates for the ordinary blocks. For a fixed restriction on a codimension-two subspace at each mark, Corollary~\ref{cor:codimmarked624} compares the projected table with its mixed differences in both norm and minor ideals. Theorem~\ref{thm:codimtwo624} identifies an entire determinant fiber on which the unnormalized projection cancels a constant coefficient. Theorem~\ref{thm:markedfourier625} computes the finite Gram operator on the entire codimension-two support. Corollary~\ref{cor:markednearisometry625} and Proposition~\ref{prop:markedgram625} give a relative error tending to zero on the detected subspace and bound the change of each mixed Gram block there. Theorem~\ref{thm:restrictionrecursion626} reduces the marked Gram operator for any fixed restriction to smaller general linear groups. Corollary~\ref{cor:mixedrestrictions626} makes the reduction simultaneously for signed inputs on different supports and bounds the dimensions required by their common restriction. Theorem~\ref{thm:integralrestrictions627} gives a reduction of the unnormalized finite matrices over the integers away from the auxiliary primes. Corollary~\ref{cor:integralmixed627} preserves all terminal determinantal ideals of the original signed input columns on a common support. Theorem~\ref{thm:residueconductor628} bounds the integral loss at the auxiliary primes themselves. Corollary~\ref{cor:globalreducedheight628} preserves the full mixed Gram matrices in a specified metric and bounds the all-place discrepancy of the usual exterior heights by a polynomial per rank. Theorem~\ref{thm:rationalgram629} removes the terms of large rational difference rank from the signed second and mixed fourth ordinary moments, at the diagonal identity normalization. Corollary~\ref{cor:rationalfixed629} identifies a common integral lattice in the fixed space of the composite transport. Theorem~\ref{thm:compactmass630} bounds orbital operator norms at the scalar identity normalization. Theorem~\ref{thm:rankoperator630} reduces the twofold and fourfold ordinary operators at this normalization to rational composites of difference rank at most a constant times \(d^3\). Theorem~\ref{thm:matchedcoordinates631} recovers each marked coefficient integrally at fixed source lattice, target lattice and isometry. Theorem~\ref{thm:commonfamilies631} gives a common list of integral coordinates in which collisions on the retained block require rational difference rank at most a constant times \(d^2\). Theorem~\ref{thm:localrecovery632} retrieves a signed coefficient in the presence of other classes when sufficiently few classes cover its nearby marked indices. Corollary~\ref{cor:fullinput632} gives a lower bound for the full ordinary period on each incoming vector and a necessary condition for exact cancellation in the integral projection. Theorem~\ref{thm:relhalf633} bounds the relative class group contribution at the half residue using the conductor of the squared character. Corollary~\ref{cor:mixedhalf633} proves half-residue vanishing on the growing dyadic layers together with the full fixed subfield of two-power degree at primes three or five modulo eight. Theorem~\ref{thm:mixedclasses634} excludes large primes in those progressions from the class groups of dyadic layers over a fixed real cyclic field. Corollary~\ref{cor:intervalzero634} descends the full zero-residue contribution on the indicated subfields to a finite list of cyclic decomposition fields. Theorems~\ref{thm:variableclasses642} and~\ref{thm:variabledescent642} give class number control and strict descent for varying dyadic Frobenius orders. Corollaries~\ref{cor:exceptionalweights642} and~\ref{cor:realhighprime642} bound the exceptional pairs on the common intervals and treat the entire real dyadic tower in the indicated high-prime range. Theorem~\ref{thm:halfclassdigits643} identifies the combined torsion of the odd-power half-residue characters with a single relative class group at each available precision. Corollary~\ref{cor:mixedhalfexponent643} removes the full mixed-field contribution at the specified weights for varying Frobenius orders. Theorem~\ref{thm:allhalfweights644} bounds the total weight divisibility of the exceptionally long half-residue components across all characters and all weights by one relative class number. Corollary~\ref{cor:periodfieldweights644} applies the bound on the full period fields throughout the common interval. Theorem~\ref{thm:quantlog646} bounds the split unit indices of several faithful characters using a simultaneous estimate for their local logarithms. Corollary~\ref{cor:splitlayers646} bounds invariant ideal classes uniformly in every cyclotomic layer for the stated degrees and characters. Theorem~\ref{thm:productmeasure647} measures the sum of the half-residue lengths by one polynomial in the action on an ordinary class group, with error independent of the precision. Corollary~\ref{cor:halftail647} includes the exceptional lengths and gives a finite test for their total on the original fields. Theorem~\ref{thm:alltails648} treats all surviving weights at a good prime by one polynomial on an ordinary class group, including the locally split zero residue. Corollary~\ref{cor:goodtest648} bounds the width of the resulting two-sided test by a polynomial uniformly on the interval. Theorem~\ref{thm:totalresidual649} measures the entire residual local sum, including all odd primes dividing the conductor or degree, by finite ordinary class group quotients with polynomial uncertainty at the proposed budget. Theorem~\ref{thm:primemass650} bounds the entire contribution of the prime ideals above any specified finite set of rational primes. Corollary~\ref{cor:detectingsplit650} gives a necessary cyclotomic splitting condition for a prime class to survive the residual cutoff. Theorem~\ref{thm:pairedtrace651} constructs an integral polynomial in the symplectic inertia traces which separates the unramified systems and bounds the loss on their congruence ideals and cotangent spaces. Theorem~\ref{thm:exactclassprimes652} realizes simultaneous prime ideal generators of finite strict class quotients at exact cyclotomic depths. When the residual characters are locally nontrivial, it proves that the specified common depth is the least possible. Theorem~\ref{thm:mixedflags653} determines the mixed Gram block at prescribed input flags, retaining all refinements at the other marked places. Corollary~\ref{cor:flagmoments653} computes its trace moments from the signed scalar matrix throughout the stated height range. Theorem~\ref{thm:inertiaann654} constructs the sharp integral trace annihilator for finite symplectic inertia on two planes. Corollary~\ref{cor:verticaltrace654} controls the loss on passing to unramified traces even in the presence of nilpotents and coefficient torsion. At the ramified primes of the character conductor, Theorem~\ref{thm:symplocal655} bounds the inertia restriction on the specified symplectic component. Theorem~\ref{thm:scalarconductor655} computes the exact cost of the corresponding scalar local condition, and Corollary~\ref{cor:sumconductor655} sums the resulting bounds over coefficient primes and weights. Theorem~\ref{thm:raylength656} computes strict torsion from a ray class group with explicit local equations. Corollary~\ref{cor:raystable656} detects the whole length when two successive precisions agree. Theorem~\ref{thm:integralblock656} applies this test on the original period fields, including primes dividing their degrees. Theorem~\ref{thm:strictcyclic657} computes complete strict groups at all positive weights divisible by the coefficient prime minus one for a class of cyclic character fields. The cubic character of conductor 3367 gives a group of order nine throughout the even weights. Lemma~\ref{lem:flagsign658} gives the signed weight of two specified input flags in the fourth moment. Theorem~\ref{thm:flagfourth658} bounds the contribution of pairs far from opposition at every orbital height on the original marked fields. Theorem~\ref{thm:firstflag659} finds the least relative flag deficit supporting a unipotent marking paired with its inverse and computes its signed contribution. Theorem~\ref{thm:flagcancel659} evaluates exactly how much the remaining positions cancel in that coefficient. Theorem~\ref{thm:oppositeaverage660} evaluates a transvection against every opposite input flag average, with an arbitrary first relative marking. Theorem~\ref{thm:pivotrecursion660} gives all other relative positions by a recurrence in two pivot indices. Theorems~\ref{thm:wholeblock661} and~\ref{thm:cyclicdescent661} compute torsion and detect the vanishing of an entire integral Sylow block by ordinary class groups. Corollary~\ref{cor:dyadiccyclic661} verifies the local hypotheses throughout the stated dyadic layers of the original fields. Theorem~\ref{thm:tameray662} identifies the extra tame inertia allowed in cyclic descent, and Corollary~\ref{cor:unitrows662} computes its residual contribution from global units. Theorem~\ref{thm:cubictower662} determines the strict group over every real dyadic layer of the cubic field of conductor \(3367\) at every positive even weight. Theorem~\ref{thm:everycyclic663} computes finite torsion at every intermediate cyclic field through local norms. Corollaries~\ref{cor:splitray663} and~\ref{cor:allmoving663} give ray class tests for all moving quadratic components at the fixed odd primes dividing the character order. Theorem~\ref{thm:conemoments664} localizes the untruncated signed moments at all real places with an error at the identity normalization. Corollary~\ref{cor:crossangles664} bounds the normalized real cross pairings for the small determinants at prime terminal frequency. Theorems~\ref{thm:descentcost665} and~\ref{thm:genusranks665} bound the change in length under cyclic descent and the elementary divisors forced by tame ramification. Corollary~\ref{cor:movingcubic665} proves that an actual moving component at conductor \(3367\) is nonzero at every odd weight on the common intervals. Theorem~\ref{thm:nodecount666} bounds positive integrals of the retained prime sum using only counts on a sparse set of scales. Corollary~\ref{cor:sparsenodes666} gives the corresponding interpolation of the signed mean between neighboring fields. Theorem~\ref{thm:threeaverages667} detects a zero to the right of the proposed line using three signed averages along the common intervals at a single rational frequency. Theorem~\ref{thm:projectedinteger668} computes a finite integer upper comparison for each chosen translate from its projected integral columns. Corollary~\ref{cor:intervalintegers668} combines these comparisons across scales with the ramified and higher-prime-power terms retained. Theorem~\ref{thm:sexticweights669} computes the complementary sign of the strict group at every weight over infinitely many sextic character fields and all their real dyadic layers. Corollary~\ref{cor:sexticmoving669} bounds the ranks of moving components on those same intervals from below. Proposition~\ref{prop:primerange670} evaluates the terms outside the first-power range between the square root and the scale to the required accuracy. Theorem~\ref{thm:primeprojections670} uses the actual projected entries in that range and allows the translate to vary with the prime. Lemma~\ref{lem:excessmass671} computes the total prime mass in this range on the common intervals. Theorem~\ref{thm:missedmass671} proves that the integral upper comparison requires a power saving in the weighted mass on which the projected entries fail to attain the specified gain. Theorem~\ref{thm:rawsaturation672} shows that the undivided columns reach the upper cutoff at every such prime. Corollary~\ref{cor:rawfailure672} proves that the proposed upper comparison for these columns therefore fails on every sufficiently distant interval. Theorem~\ref{thm:residualvaluation673} removes the common divisor and evaluates the excess of the new comparison in terms of the residual valuations of the projected columns. Theorem~\ref{thm:cycleorientation674} identifies the maximizing orientation on every cycle containing a distinguished embedding in this prime range. Corollary~\ref{cor:fivepowers674} evaluates the normalized logarithm of the removed divisor there using the first five odd powers of Frobenius, with absolute error tending to zero. Theorem~\ref{thm:oddstrict675} identifies the finite odd cocycles of the second archive with one strict eigensubgroup at odd primes. Proposition~\ref{prop:dyadicodd675} controls the discrepancy at two, and Corollary~\ref{cor:finiteoddtotal675} bounds the entire weighted sum of these finite cocycles on the period fields. Theorem~\ref{thm:uniformannihilator676} gives an explicit integer annihilating those odd cocycles for every coefficient prime and every CM subfield of a given cyclotomic field, by a finite-level form of Stickelberger's theorem. Theorem~\ref{thm:oddinteger677} bounds the full lengths by the valuations of one integer, uniformly over all CM subfields of the cyclotomic field. Corollary~\ref{cor:oddwindow677} supplies this integer on the original period intervals and the odd-cocycle input of the shorter source argument. Theorem~\ref{thm:dualrelations712} constructs independent dual Selmer classes at every coefficient precision. Proposition~\ref{prop:evenreciprocity712} gives their relations on tame inertia, and Corollary~\ref{cor:fieldrelations712} applies them on both families of period fields. Theorem~\ref{thm:auxtriangle713} selects distinct split primes on those fields and gives a triangular bound on the tame values of even cocycles. Corollary~\ref{cor:scalarfrob713} preserves a unit determinant condition on an auxiliary representation at the same primes. Theorem~\ref{thm:trianglebudget714} bounds the total valuation loss of these primes on an entire period interval. Corollary~\ref{cor:actualmarks714} bounds it by the two prescribed level group orders in the shorter construction. Theorem~\ref{thm:markedmatrix715} computes the integral matrix algebra of the marked tame types and produces the scalar inertia used in Corollary~\ref{cor:scalarfrob713}. Corollary~\ref{cor:markeddet715} gives the unit determinant at the auxiliary primes for each representation of these types. Proposition~\ref{prop:toproots716} proves which factors of the spherical Hecke polynomial annihilate the Iwahori invariants. Theorem~\ref{thm:commonhecke716} constructs a common integral polynomial for the compatible weak lifts at the primes used by the cohomology tests. Theorem~\ref{thm:scalarcost717} bounds the scalar groups after the paired inertia and Frobenius tests, including the cost of integral representatives at two. Corollary~\ref{cor:scalarsum717} sums these bounds over the weights on the period fields. Theorem~\ref{thm:fittingcost718} turns the paired relations into a Fitting determinant with a bounded coefficient. Theorem~\ref{thm:allscalarfit718} produces one such coefficient for the full set of normalized scalar cocycles, uniformly in the precision. Proposition~\ref{prop:blockrows719} realizes the main summand's cocycle bound by a matrix of group-algebra relations. Theorem~\ref{thm:charrows719} realizes the scalar determinants for all weights by matrices of those same relations. Theorem~\ref{thm:ribettrace720} gives an integral trace multiplier and bounds its squarefree coefficient at a chosen lift. Corollary~\ref{cor:augdet720} gives one augmentation determinant that detects every system without a trivial composition factor. Theorem~\ref{thm:tametransfer721} identifies the tame trace after transfer at each split auxiliary prime for locally compatible systems. Proposition~\ref{prop:righttransfer721} controls right translation of the same operator by the base stabilizer orders. Theorem~\ref{thm:splitattachment722} constructs the Galois representations for the definite unitary groups on the period fields and proves local compatibility on every type entering the auxiliary transfer. Corollary~\ref{cor:normalizedtransfer722} verifies the conjugate relation and the spherical support after the theta normalization. Theorem~\ref{thm:topdescent723} descends the squarefree trace coefficient to the integral lattice at base level. Corollary~\ref{cor:baseseparation723} makes the separation of base systems independent of the subsequent choice of auxiliary primes. Corollary~\ref{cor:scalarfl724} proves that all scalar crystalline colengths vanish at the spherical unramified coefficient primes in a common Fontaine--Laffaille range on the period fields. Proposition~\ref{prop:mainselmer724} imposes the crystalline local condition at these primes on the main summand. Theorem~\ref{thm:jointtamerow725} constructs simultaneous inertia rows from Frobenius at every old place. Corollary~\ref{cor:levelcost725} bounds their combined contribution at the two level primes by twice the logarithm of the level group order. Theorem~\ref{thm:allraypowers726} gives one anticyclotomic character for simultaneous polynomial tests on all its prescribed powers. Corollary~\ref{cor:raytamagawa726} bounds the local Selmer cost at its new primes on the period field. Proposition~\ref{prop:sharpminors727} expresses each finite Selmer length by integral polynomial minors. Theorem~\ref{thm:preserveselmer727} constructs a twist on a single pair whose powers preserve the original upper bound up to an arbitrarily small normalized error. Theorem~\ref{thm:selftwist728} computes the global congruence term of a marked summand from its Frobenius traces. Theorem~\ref{thm:balancedcount728} bounds the primary Selmer length by its dual and the original ramification terms on the period field. Theorem~\ref{thm:tiedlength729} chooses split prime pairs and local equations for which the dual Selmer group injects into the computed global congruence term, simultaneously for all the indicated main lattices. Theorem~\ref{thm:twistedselmer730} transfers these bounds to genuine anticyclotomic twists and all their prescribed powers. Corollary~\ref{cor:maincost730} bounds the higher-rank trace cocycles at the spherical coefficient primes on the shorter period fields. Theorem~\ref{thm:commoncritical731} bounds the critical Hecke-value ratios for the very same twists, including the endpoint at one, when the coefficient prime is also prime to the old conductor and the field discriminant. Theorem~\ref{thm:localfinite732} and Corollary~\ref{cor:fixedinterval732} give a finite local cohomology loss for arbitrary fixed weight intervals, uniform in the auxiliary levels, coefficient extensions and precision. Its growth as the fields and weights vary still needs a bound. Corollary~\ref{cor:allcritical733} removes the restrictions on the coefficient prime arising from the CM discriminant and the old Hecke conductor for characters ramified at both marked levels. Corollary~\ref{cor:quotientchainheight734} bounds the height of the product of the critical Hecke ratios for the full chain by any prescribed positive number on the shorter fields. Corollary~\ref{cor:labelledtrace735} removes the local trace defects at each spherical completion whose filtration intervals have the required widths separately at its embeddings. Theorem~\ref{thm:rankonepair736} turns the trace bound into a denominator estimate for the inner product of two projected torus tests. Its additional constants depend only on the original levels. Theorem~\ref{thm:rankonespherical737} sums these projected inner-product denominators over the spherical coefficient primes in the labelled lifting range, uniformly in the tests and the auxiliary moduli. Corollary~\ref{cor:alltwistcost738} bounds the new twist-place row terms and the integral index of both inertia summands in every block rank for these same twists. Theorem~\ref{thm:oddselmer739} bounds the marked odd-rank Selmer lengths with the supremum taken over a whole shorter period interval inside the sum over the coefficient primes. Theorem~\ref{thm:negativeinterval740} gives a summable bound at every negative-jump completion of the designated rank-one main characters on this interval and gives zero there at the two level primes. Corollary~\ref{cor:torustests741} bounds the torsion error at all coefficient places and counts the integral cup products still to be bounded for the rank-one main and scalar trace classes. Theorem~\ref{thm:projectedcups742} compares the denominator of the projected torus pairing with those free cup products at every prime, with a summable error on the same period interval. Theorem~\ref{thm:oldtorustraces743} proves the integral local condition for every main trace at both original level primes and for the indicated scalar weights there. Theorem~\ref{thm:abelianperiod744} bounds the full prime sum for abelian local trace families and expresses the remaining free cup products through integral tuples supported on members with noncommuting local matrices. Theorem~\ref{thm:oldcolumn745} identifies the entire remaining joint cup image at either original level prime by a single column of tame blocks of a fixed wild inertia element. Theorem~\ref{thm:potentialperiod746} proves the projected prime sum for families potentially abelian at each coefficient place, with no common abelianizing extension required. Theorem~\ref{thm:markedloss747} removes the entire old-prime scalar contribution when the marked character occurs on a Galois subline, a quotient line, or a line with extremal jump. Corollary~\ref{cor:oldlogperiod748} removes the entire weight-zero scalar term at both original level primes for all marked families by the finite-precision cyclotomic logarithm calculation of Theorem~\ref{thm:oldweightzero748}. Theorem~\ref{thm:residuepositive749} computes the first residue of each remaining positive scalar generator on the rational uniformizer tower, and Corollary~\ref{cor:finitepositive749} gives finite coefficient equations there for bounding its integral contribution. Corollary~\ref{cor:rationaldepth750} bounds the positive scalar loss for every old-prime family by finite rational-tower depths independent of the tame extension, the far jump and the trace precision. Theorem~\ref{thm:unitcoordinate751} evaluates each positive scalar class at all precisions on one norm-compatible sequence of explicit cyclotomic units. Corollary~\ref{cor:oldperiod752} proves the integral projected period estimate at both original coefficient primes for the full marked torus family on the same shorter interval. Theorem~\ref{thm:tameorder753} gives simultaneous necessary inequalities on two tame characters for a nonfinite scalar class, including at primitive labels of large inertia order. Corollary~\ref{cor:twoperiod754} removes wide two-dimensional spherical blocks from the positive scalar trace equations on the period interval and isolates the two-dimensional weight pattern still needed for the main line. Theorem~\ref{thm:filteredperiod755} uses the residual constituents of each rational filtration to omit local cup coordinates, including the ramified weight-zero coordinate, from the same projected period estimate. Theorem~\ref{thm:ordinaryperiod756} proves the full projected period estimate for the marked torus data with potentially ordered crystalline flags at the spherical coefficient places. Theorem~\ref{thm:cuppath757} proves the Frobenius cochain calculation for arbitrary local matrices, and Corollary~\ref{cor:pathperiod757} applies it to the integral free cup tests of the same period data. Theorem~\ref{thm:unramifiedperiod758} bounds the weight-zero period loss for potentially triangular local matrices in arbitrary flag order and treats every target at their places with nonnegative jumps. Theorem~\ref{thm:integrallifts759} bounds the power of the Kisin parameter required to lift a quotient generator at each torsion precision for arbitrary crystalline trace matrices. Corollary~\ref{cor:periodlifts759} gives integral lifts for every scalar dual test on the period fields. Theorem~\ref{thm:uniformlift760} turns the finite-precision lifts into a single integral lift of a specified power of the coefficient prime when the absolute ramification index is at most that prime minus two. Theorem~\ref{thm:allramifiedlifts761} gives a lift independent of precision at every ramification index, and Corollary~\ref{cor:threeextensions761} applies it simultaneously to the trace, the test and their torsion pairing. Theorem~\ref{thm:unrestrictedcups762} constructs a two-step Frobenius path and bounds the order of each local cup in terms of the actual Hodge jumps and ramification, for arbitrary rank and torsion precision. Corollary~\ref{cor:sphericalcupbound762} bounds the resulting contribution to the joint integral period image. Theorem~\ref{thm:subcriticalperiod763} sums the new cup bound on the original period interval for all local targets within its stated height bound. Corollary~\ref{cor:shortscalar763} includes every scalar test supported on the short weights, at every spherical place and in arbitrary rank. Theorem~\ref{thm:orbitlift764} tracks the upper jumps separately around each unramified Frobenius orbit. Corollary~\ref{cor:orbitconstants764} bounds the total length lost on the lattice of constants by the sum of the labelled heights. Theorem~\ref{thm:labelledcup765} compares the cup cochains at successive embeddings with different integral lift exponents, yielding an annihilator at the original coefficient precision. Theorem~\ref{thm:exceptionalperiod766} proves the full marked torus period estimate, for arbitrary local ranks and all the given weights, on the large coefficient primes having a nonzero multiple in the specified neighborhoods of the weight combinations. Theorem~\ref{thm:deepintegrallift767} raises the order of the Galois displacements at the Eisenstein prime. The cost is a sum of valuations over specified ranges of integers near the jumps, and Corollary~\ref{cor:secondlatticedepth767} preserves divisibility at the adjacent Frobenius prime. Corollary~\ref{cor:actualperiodmultiples768} bounds the exponent of every original local cup at a spherical prime above the first auxiliary scale by the valuations of integers in five intervals. If these intervals contain no nonzero multiple of the prime the cup vanishes at the original precision. Theorem~\ref{thm:fullperiodcutoff769} bounds the full marked torus period on every coefficient prime above the square root of the weight times the displayed polynomial and logarithmic factors, on the original fields and interval. Theorem~\ref{thm:scalarmeromorphic770} allows a pole in the graph point during the cup comparison. Corollary~\ref{cor:allscalartests770} thereby bounds every scalar-target cup at every unramified spherical prime of residue characteristic at least five. Theorem~\ref{thm:remainingmain771} bounds the marked period after retaining only specified free local coordinates. Between the polynomial threshold and the first auxiliary scale, only main targets at places with both signs and both base weights are retained. Theorem~\ref{thm:mixedcup772} bounds the original main cups at mixed unramified places by a second Taylor expansion. Restriction to the degree-prime extension fixing a root of the uniformizer and corestriction back to the original field cost one further power of the prime. Theorem~\ref{thm:middlerange773} confines the remaining mixed main coordinates above the polynomial threshold to a shorter interval about the ratio of the two base weights. Corollary~\ref{cor:coarsemixed773} bounds their contribution on that interval. Theorem~\ref{thm:frobeniusjet774} constructs an integral vector representing a prescribed horizontal Taylor coefficient at the uniformizer after Frobenius. Its cost decreases with the number of Frobenius steps and is independent of the weights. Theorem~\ref{thm:secondpoint775} constructs a common integral point in two Frobenius lattices across an intervening label. Its displacement is divisible by prescribed powers at both primes. Theorem~\ref{thm:mainorbit776} bounds every original main cup at a spherical prime above the polynomial cutoff, with weight terms of order the larger weight divided by the square of the prime and the smaller weight divided by the prime. Theorem~\ref{thm:smallcoords777} controls the full projected integral torus period at the source normalization for every coefficient prime above the polynomial threshold, on the original fields and interval. Theorem~\ref{thm:rankcutoff778} extends the full integral torus-period estimate down to the product of the source rank and the string bound. The primewise interval supremum costs at most an inverse square root of the initial node. Proposition~\ref{prop:simultaneousorders779} imposes several local Taylor orders on the same integral point. Theorem~\ref{thm:longpoint779} supplies a point spanning an arbitrary string of bounded intermediate jumps. Theorem~\ref{thm:sparsecups780} bounds the main and scalar cups along an unramified orbit in which three larger weights cannot occur with two short successive gaps. Theorem~\ref{thm:fixedconductor781} proves the integral torus-period bound at the source normalization outside the prime divisors of the fixed conductor, for the sparse choices on the original fields and interval. Theorem~\ref{thm:ramifiedcut782} quantifies the quotient lift across bounded residue-field labels at any ramification index. Proposition~\ref{prop:ramifiedorders782} gives the integral simultaneous Taylor orders there, including at two. Theorem~\ref{thm:alltorus783} bounds the full projected integral marked torus period at all coefficient primes at the source normalization, with a single choice of the designated data and simultaneous twists on the period interval. Proposition~\ref{prop:torusbudget784} puts the torus tests within the determinant budget of the chain. Theorem~\ref{thm:projectionindex784} identifies the minimum loss for each higher-rank integral projection with the length of its tested congruence module. Theorem~\ref{thm:heckeflag785} bounds this length by the valuations of split Hecke values on the tested complement, with separate multiplicities for the competing systems. Corollary~\ref{cor:markbudget785} sums the cost of the marking on the period interval. Theorem~\ref{thm:residualflag786} confines the remaining congruence length to shared residual representations and bounds it there by the full cyclotomic polynomial remainder and by pairwise characteristic-coefficient differences. Theorem~\ref{thm:longcharacter787} identifies the global character block on the higher-rank lifts and its full cyclotomic Galois summand. Theorem~\ref{thm:compgap787} excludes every block of this length from the tested complement and bounds its split Hecke values at the complex place. Theorem~\ref{thm:highfinite788} computes the corrected finite local lengths of the higher-rank summand at every prime. Theorem~\ref{thm:highlift788} eliminates its integral trace-class loss at spherical split completions in the labelled Fontaine--Laffaille range. Theorem~\ref{thm:highsmall789} bounds the sum of the higher-rank local trace defects up to the product of the ranks at spherical places away from the fixed conductor, including inert places and weights outside the Fontaine--Laffaille range. Theorem~\ref{thm:highallsmall790} includes the primes of the fixed conductor, allowing their ramification indices to grow, and gives a single interval bound for the local quotients through the product of the ranks. Theorem~\ref{thm:highmiddle791} controls the higher-rank local quotients up to \(x_{1,x}/x^{12}\) and at every prime in the two classes permitting three successive large weights. Theorem~\ref{thm:highproduct792} restricts the defects at \(p\ge 8b_x/r_x^8\) to prime divisors of an explicit product of short shifts of \(b_x\) and bounds their exponent. Theorem~\ref{thm:highscalars793} bounds the full spherical local sum at all coefficient primes for the scalar summands of the higher-rank graph. Theorem~\ref{thm:neggraph794} bounds the full remaining local sum for the main summand at completions with negative far graph weights, using the multiplicity of the obstructing Hodge weights. Theorem~\ref{thm:interior795} handles every spherical completion and every target in the range \(8x_{1,x}/x\le p\le b_x/x^{45}\), including changes of orientation on a single orbit. Theorem~\ref{thm:transitionrange796} gives a single local estimate through \(p=b_x/x^{45}\) at all spherical places on the original interval. Its proof treats the transition where the possible pole order crosses the smaller base weight. Theorem~\ref{thm:jetlength797} bounds the length of a Taylor condition at the first Frobenius prime by the multiplicity of the affected weight. Corollary~\ref{cor:negativejets797} applies this bound simultaneously to the integral dual tests. Theorem~\ref{thm:allspherical798} bounds the spherical local cohomology quotients for the main and scalar targets together, at every coefficient prime on the same fields and period interval. Theorem~\ref{thm:typedindex799} compares the global tested theta index with the sum of its indices for individual marked residue types, including at the two level primes. Corollary~\ref{cor:typedbudget799} places the error in the original determinant budget. Theorem~\ref{thm:tamecups800} controls the integral cup at a split level prime for a common tame inertia spectrum. Corollary~\ref{cor:tamebudget800} sums the resulting losses for fixed spectra on the same fields and intervals. Theorem~\ref{thm:markedtype801} identifies the common smooth inertia characters for every constituent in a fixed marked component at each original level prime. Corollary~\ref{cor:markedlocal801} applies the local cohomology bound to these components on the period interval. Theorem~\ref{thm:markedperiod802} bounds the higher-rank integral projection by the divisible Selmer lengths of its true summands and the coefficient-place defects on the same marked fields and interval. Proposition~\ref{prop:sourcetwist803} identifies the summands of the actual higher-rank lifts before choosing their common twist. Corollary~\ref{cor:actualdenom803} applies the resulting simultaneous Selmer estimate to the determinant and torus denominators on the shorter interval. Theorem~\ref{thm:actualminors804} bounds every minor of the projected theta-period pairing at the complex and finite places, with one projection index in the finite bound for the whole minor. Theorem~\ref{thm:auxdet805} identifies the common rational level factor in the Gram determinants of all auxiliary single lifts. Proposition~\ref{prop:auxvaluation805} computes its coefficient-prime valuation, including at two. Theorem~\ref{thm:commonsource806} identifies the finite Hecke modules of the marked row sources. Corollary~\ref{cor:scalarpower806} computes the resulting exponent in the product of the auxiliary determinants. Theorem~\ref{thm:terminalflag807} counts the terminal flag and evaluates its normalized level index. Corollary~\ref{cor:remaininggram807} cancels the explicit conductor factor and isolates the Gram determinant still requiring an analytic comparison. Theorem~\ref{thm:flagdoubling808} evaluates all split row point factors together with the terminal flag index. After normalization by the split doubling denominators their ratio to the auxiliary factors is a coefficient-prime unit at every new twist place. Theorem~\ref{thm:rowcuspidal809} proves cuspidality at every step of the prescribed bilinear row on these same fields. Proposition~\ref{prop:rowrealgram809} proves strict positivity of its specified archimedean factors, including the steps outside the range of absolute convergence of the Siegel Eisenstein series. Theorem~\ref{thm:detquotient810} computes the change of the row Gram determinant relative to the auxiliary determinants under a tame twist. Corollary~\ref{cor:detheight810} bounds its positive coefficient-prime valuation on the interval. Theorem~\ref{thm:splitlattice811} and Corollary~\ref{cor:gluedperiod811} give the product tests for an individual auxiliary period. Theorem~\ref{thm:gluedcost811} bounds the denominator of their algebraic coefficients at the old marks. Theorem~\ref{thm:complementgram812} computes the Gram operator of the complementary factor. Corollary~\ref{cor:complementreproduce812} proves that its period evaluations span the dual of the specified archimedean projection at the actual level. Theorem~\ref{thm:carrier813} constructs the nonzero cuspidal form used in the first factors. Theorem~\ref{thm:jointminima813} gives a simultaneous lower bound for the shortest vectors at its peaks and at those of every good form on the same field. Theorem~\ref{thm:sparseapprox814} realizes a real oscillator coefficient at these peaks by a linear combination of tests at the prescribed finite level, with an explicit uniform error. Corollary~\ref{cor:paddedminor814} applies it to the determinant polynomial of the first factor. Theorem~\ref{thm:integratedpeak815} computes the main term after integration against the carrier on its real group and proves that an elementary residue test survives that integration at every good-form peak. Corollary~\ref{cor:adelicfirst816} expresses this integral as a combination of the prescribed adelic first factors and proves that some first factor is nonzero at each such peak. Theorem~\ref{thm:firsthol817} determines the holomorphic weight of those adelic first factors by transferring noncompact root derivatives directly to the carrier. Theorem~\ref{thm:complexseparation818} proves that the products of these first factors with the complementary class tests separate every good form over the complex numbers on the original intervals. Theorem~\ref{thm:stackedheight819} gives common denominators and archimedean bounds for the actual determinant coefficient polynomials on the definite side of the comparison. Proposition~\ref{prop:markedfourier820} computes the finite Fourier coefficients at every twist, and Theorem~\ref{thm:markeddefinite820} bounds all conjugates of the three normalized positive definite class kernels. Theorem~\ref{thm:scaleddefinite821} gives one normalization of these evaluated kernels that is integral at integral arguments away from each new conductor and has the same conjugate size bound. Theorem~\ref{thm:definiteclass822} proves integrality in the moving target Weyl lattice and bounds all conjugates of the positive definite double class sum for every prescribed twist. Theorem~\ref{thm:signedseries823} transfers the integral and archimedean estimates to the definite series with its prescribed negative source, using a single scalar across all twists. Theorem~\ref{thm:arithmeticclass824} inserts the bounded vectors into the theta class integral on the prescribed group with the exact weights at the refined level. Theorem~\ref{thm:carriermultiplicity825} computes the old-level multiplicity of the separating carrier period for the matched residue and archimedean data. Theorem~\ref{thm:globalequalrank826} proves that the first equal-rank class integral is constant on the entire adelic source for every prescribed twist. Theorem~\ref{thm:squareconductor827} computes the absolute value of the continued relative split doubling product at equal rank, including the exact power of the tame conductor. Theorem~\ref{thm:normalfirst828} and Corollary~\ref{cor:signedfirst828} prove nonvanishing and determine the twist dependence of the normalized first global exchange, including its negative-source form. Theorem~\ref{thm:secondcomparison829} proves that the full-level cuspidal form in the second exchange satisfies the same split Hecke equations as the carrier, also at all marks. Theorem~\ref{thm:allcarriers830} identifies the spherical parameters at every finite place for the two sides of the second exchange and for the adjoint image of the carrier. Theorem~\ref{thm:carrierline831} proves that each of these spherical and archimedean systems occurs on a single automorphic line. Theorem~\ref{thm:secondlocal832} evaluates the local doubling integrals for the adjoint carrier, with a positive integrable value at each real place. Theorem~\ref{thm:secondglobal833} proves the normalized carrier identity for every prescribed twist and evaluates its two conductor and Hecke-value factors on the given field. Theorem~\ref{thm:smallthird834} identifies the contracted smaller-rank integral on the third pair with the separating carrier period for every twist, with its exact scalar. Theorem~\ref{thm:marksections835} identifies the two marked frame sections in the third exchange on the whole split group, including nonintegral translates. Theorem~\ref{thm:thirdsplitprod836} evaluates the relative split frame multiplier at the third critical block on the same fields and twist sets. Theorem~\ref{thm:thirdstandard837} realizes both actual global integrals in the third comparison by the same finite sum of translates of one pair of standard residue tests. Theorem~\ref{thm:thirdreal838} identifies the full real compact types for the given polynomials and computes their relative Gaussian norm. Corollary~\ref{cor:thirdseed838} reduces the third comparison to a single explicit pair of global theta functions per twist. Theorem~\ref{thm:thirdlargenorm839} evaluates the squared norm of the larger standard global theta integral and proves its nonvanishing independently of the smaller integral. Theorem~\ref{thm:thirdsmallnorm840} evaluates the norm on the indefinite smaller side. Corollary~\ref{cor:thirdnormratio840} gives the full third conductor and Hecke-value multiplier as a ratio of norms of the two actual global functions. Theorem~\ref{thm:thirdholomorphic841} proves the same holomorphy equations on the two global standard functions by integration by parts on the smaller group. Corollary~\ref{cor:thirdrealmatrix841} identifies their normalized real matrix coefficients. Theorem~\ref{thm:thirdadelic842} places both standard global periods on the same specified vector in abstractly equivalent discrete automorphic representations, including their ramified finite components. Theorem~\ref{thm:thirdglobal843} identifies the two global theta functions and gives the common scalar relating the definite class sum to every carrier test on the original field and interval. Theorem~\ref{thm:gluedperiodmap844} realizes the product periods on the original definite target by a map preserving the Weyl lattices away from the new conductor and identifies its adjoint on good lifts with one injective map for all twists. Theorem~\ref{thm:integralcomparison845} constructs a positive comparison operator integral away from the twist conductor, bounds every archimedean conjugate at the trivial twist, and computes the exact quadratic form it induces on the good source. Theorem~\ref{thm:descenttower846} excludes every first nonzero descent strictly between the original rank and the rank needed for the complementary periods, using the actual split Hecke root on the same field. Theorem~\ref{thm:markedbottom847} proves that vanishing of the block descent passes to arbitrary Schwartz functions at all its finite marked places, on the actual selected types. Theorem~\ref{thm:restrictedtower848} propagates bottom vanishing up the required tower with the same unmarked vacuum and the same definite coefficient, by a pointwise Hecke identity and an operator norm bound. Theorem~\ref{thm:falsedescent849} proves zero descent on every required member of the tower for the actual false-lift space at one split place, including all real coefficient vectors and all marked Schwartz data. Theorem~\ref{thm:testedcomparison850} places the integral comparison of every true lift in the tested space by unfolding a Fourier coefficient over the complementary definite group at the given level. Theorem~\ref{thm:compressionbudget851} bounds the actual compressed determinants with the local cohomology estimate on the same interval. Corollary~\ref{cor:terminalcomparison851} gives their common multiplier for the terminal Gram ratio after removing the explicit critical Hecke factors. Theorem~\ref{thm:torusnorm852} computes the conductor power in the norm of the actual rank-one lift. Corollary~\ref{cor:torusvaluation852} expresses the valuation of an inverse critical ratio using two tests of that same lift. Theorem~\ref{thm:flattorus853} gives a single normalization for the flat torus coefficients, integral away from each new conductor. Corollary~\ref{cor:flatmean853} bounds the common algebraic multiplier for two auxiliary characters. Theorem~\ref{thm:fourblock854} constructs integral tested forms for each derivative monomial and gives their exact pairing with the critical torus kernel, including the three class factors. Theorem~\ref{thm:twelftransfer855} bounds a critical inverse ratio by the valuations of a single untwisted twelve-row derivative table and its conjugate, with the projected-period loss on the same family. Theorem~\ref{thm:completesystem856} identifies the actual residue sections of the twelve-row test with a basis of a complete linear system and proves surjectivity onto all Taylor coefficients through the prescribed order. Corollary~\ref{cor:nonzerotable856} proves full row rank of the algebraic derivative table on the same fields and intervals. Theorem~\ref{thm:integralleading857} bounds the least valuation in every row of the actual derivative table by a common scalar, factorials and explicit CM lattice ideals. Corollary~\ref{cor:largeprimetable857} absorbs the table contribution to the projected period at primes larger than the weights and away from the relative different, including the original marked primes. Theorem~\ref{thm:twocopy858} compares the twelve-row table with integral scalar functions on two copies of the original formal abelian group. Corollary~\ref{cor:splitperiod858} proves the projected-period estimate at every prime that splits in the imaginary quadratic field, in all derivative weights and with the same twist and local cohomological bound. Theorem~\ref{thm:scalartrace859} bounds the derivative table at every nonsplit prime by the derivative rate of one integral elliptic formal group. Theorem~\ref{thm:shortpairs859} gives the directional local prime sum at the ramified prime and isolates the remaining terms for small inert primes as disjoint pairs of larger-weight embeddings. Theorem~\ref{thm:inertgradient860} computes the reduced logarithmic derivatives of an integral formal frame of the actual one-entry theta line at an inert prime. At each label the surviving terms have exactly the odd Frobenius exponents, and every coefficient is nonzero. Theorem~\ref{thm:pairedtable861} realizes one saving in the valuation of the actual twelve-row table for every repetition of a prescribed Frobenius pair. It applies to the small inert-prime pairs of Theorem~\ref{thm:shortpairs859} and to the whole Frobenius permutation at exponent one. Theorem~\ref{thm:allprimetable862} proves the directional upper bound for the actual Taylor table at all finite primes with one multiplier of controlled size. Corollary~\ref{cor:allprimeperiod862} combines it with the local cohomology estimate on the identical fields, weights, intervals and twists, using a CM normalization of absolute square \(M^h\). Theorem~\ref{thm:jetscale863} incorporates the relative different and the elliptic area into the normalization at every embedding while preserving the finite estimate. Corollary~\ref{cor:identitynorm863} bounds the actual untwisted norm ratio at the distinguished embedding. Theorem~\ref{thm:weyltransport864} identifies the conjugates of arbitrary linear combinations of finite Weyl translates, with a single semilinear map on the marked theta model. Theorem~\ref{thm:allnorm865} compares the actual algebraic class norms at every complex embedding. It identifies the two conjugate theta sums in the Hermitian pairing and cancels their common scalar against a flat lift of the same rank. Theorem~\ref{thm:terminaldirection866} gives a signed bound for the terminal determinant ratio itself on the shorter interval. Its directional term is the same explicit prime-power sum at every main critical pair. Theorem~\ref{thm:longframe867} converts the signed terminal determinant estimate into a covolume lower bound for full or completed frames. Lemma~\ref{lem:terminalintegral867} proves integrality of the terminal class pairing at the actual marked level without a stabilizer denominator. Theorem~\ref{thm:fewtranslates868} constructs independent inputs for the terminal determinant with a logarithmic bound on the number of translates. Proposition~\ref{prop:goodmass868} gives the positive source dimension and an exact finite class formula for it. Theorem~\ref{thm:equallinks869} compares the full residue moment on the same nodes to signed integrals grouped by equal full links, with an error exponentially small in \(kf\log b\) relative to the mass normalization. Theorem~\ref{thm:openenergy870} reduces the sharp spanning estimate to equal-link integrals with nonzero successive minors at one marked place, and proves that fixing frame vectors at a quarter of the first marked places already gives negligible mass. Theorem~\ref{thm:orbitfactors871} gives a finite formula for the signed equal-link sums on the open residues. It identifies the middle isometry and separates the factors at the finite and infinite places before the square is taken. Theorem~\ref{thm:spheremean872} evaluates the actual sum over definite classes in the open-link energy by an integral on a single norm sphere. The factor outside the selected finite places is \(1+O(f2^{1-k})\), uniformly in the represented norm. Theorem~\ref{thm:bothmarks873} imposes unit pivots in both components at both original marks outside fewer than \(f/4\) places. Corollary~\ref{cor:integratedmark873} evaluates the entire integral at each unused mark and gives the common pivot congruence there. Theorem~\ref{thm:selectedopen874} evaluates the integral for a fixed frame vector at a split mark with unit pivots, and Corollary~\ref{cor:singularselected874} bounds the energy by the part where every selected mark has a singular minor. Theorem~\ref{thm:freqnode875} sums the arithmetic frequency weight at the two scales of the shorter construction, and Corollary~\ref{cor:pivotoperator875} bounds its matrix on distinct last pivots. Theorem~\ref{thm:mixedperiod876} replaces one cut in the signed link energy by the full discrete-series vector, and Theorem~\ref{thm:mixedmass876} evaluates the resulting position integral with a uniform tail estimate. Theorem~\ref{thm:terminalfourier877} computes the constant Fourier coefficient of the actual unramified terminal density and bounds all other coefficients, including those at unused split marks. Theorem~\ref{thm:singularflag878} gives depth bounds on the actual flag determinants at a selected mark, and Corollary~\ref{cor:terminalfixer878} reduces the singular energy to fixers of the terminal vector. Theorem~\ref{thm:armopen879} evaluates the fixed-vector loss whenever the full arm pairing is invertible, including singular internal flags, and Corollary~\ref{cor:armdefect879} separates the remaining rank defects. Theorem~\ref{thm:integralmass881} bounds the integral part of the local position mass at a terminal fixer on every arm stratum. Theorem~\ref{thm:unrestrictedmass882} gives the identity contribution and a power-saving error on all integral positions at an unrestricted mark. Theorem~\ref{thm:allpositions883} gives the cancellation of the two middle root coordinates at arbitrary depth, including positions outside the integral group. Theorem~\ref{thm:nonspheretail884} bounds the deep middle-coordinate contribution on the nonintegral shells. Theorem~\ref{thm:mixedopen885} evaluates the position correlation for each original open residue and every pair of heads, and Corollary~\ref{cor:unusedposition885} applies it at all the unselected marks. Theorem~\ref{thm:terminalrange886} controls all terminal-stabilizer position correlations, even in absolute value, throughout a growing range of denominator exponents. Proposition~\ref{prop:armshear887} gives a measure-preserving change of unipotent variables inside each principal double coset and a cancellation test for every interior arm row. Theorem~\ref{thm:allfixermass888} bounds the full local position mass at a terminal fixer over the entire marked group. Corollary~\ref{cor:allmarks888} gives the product with the original open tests on the same fields. Theorem~\ref{thm:commoncompletion889} gives a positive majorant for the actual common-pivot terms at the unselected places, with integral dual frequencies and an explicit cost in the local zero coefficients. Theorem~\ref{thm:firstgram890} computes the principal-position term of the first relative Gram transform and proves the finite frequency weight for that term at a selected mark. Theorem~\ref{thm:gramtail891} proves absolute convergence and a conductor-dependent bound on the distant positions of the first relative Gram transform. Theorem~\ref{thm:integralfreq892} bounds the first relative Gram transform on every integral position at both types of split mark and every nonintegral frequency. Proposition~\ref{prop:congruencewave893} gives exact vanishing at unequal conductors on a range of positions. Theorem~\ref{thm:shallowwave893} includes the terminal saving for small positive heights relative to conductor. Theorem~\ref{thm:endtransform894} proves the first-coordinate frequency and terminal mass bounds together at every matrix height when both coefficient factors satisfy the stated end-ideal test. Theorem~\ref{thm:rowdescent895} obtains a saving at the full depth of the first end by integrating successive rows against the original unipotent character. Theorem~\ref{thm:twodepth896} controls both truncated end depths on every shell, including the group fixing the terminal vector. Theorem~\ref{thm:allkmass897} isolates the principal term of the unrestricted local mass on the entire split group and bounds the remaining mass. Theorem~\ref{thm:fullfirst898} proves the first relative Fourier bound on the entire split group with all original coefficients. Corollary~\ref{cor:firstsample898} sums the resulting marked factors on the shorter interval. Theorem~\ref{thm:standardfirst899} proves the full first-coordinate estimate at standard unmarked split places. Corollary~\ref{cor:standardeuler899} bounds their product on the same fields. Theorem~\ref{thm:hermitiancone900} evaluates the principal first transform at a standard inert place in every residue characteristic. Proposition~\ref{prop:inertoperator900} controls the associated operator on the full compact orbit. Theorem~\ref{thm:inertneighbor901} calculates the zero-frequency mass of the first inert lattice shell and bounds its full contribution at every conductor. Theorem~\ref{thm:inertlength902} bounds the unweighted inert shell energy by the residual length after accounting for both middle blocks of the unipotent. Corollary~\ref{cor:inertsum902} sums the part where this length grows with height. Theorem~\ref{thm:inertterminal903} proves cancellation from the terminal row on each refined inert orbit. Corollary~\ref{cor:inertterminal903} gives a summable bound even when the two middle lengths in Lemma~\ref{lem:inertmiddle902} exhaust the lattice length. Theorem~\ref{thm:inerttail904} bounds the entire distant part of the first inert transform at every conductor, including its zero-frequency mass. The proof chooses an interior simple root by the successive lengths of a lattice flag. Theorem~\ref{thm:inertzero905} estimates the complete zero-frequency mass at a standard inert place. Corollary~\ref{cor:inertproduct905} combines it with the split calculation in an absolutely convergent product on the same fields. Theorem~\ref{thm:inertprime906} proves the full first transform at a standard inert place for a pole of order one. Corollary~\ref{cor:primeproduct906} gives the product at good split and inert places for frequencies with squarefree inert denominator. Theorem~\ref{thm:inertjoint907} imposes all simple-root character tests at once on each original refined inert orbit. Proposition~\ref{prop:inertloss907} bounds their combined loss by one index on a fixed summand of the middle lattice. Theorem~\ref{thm:inertfull908} bounds the original first relative Gram transform at every conductor at standard unmarked inert places. Corollary~\ref{cor:standardall908} gives the standard split and inert Euler product for arbitrary rational frequency on the period fields. Theorem~\ref{thm:allunmarked909} verifies the local hypotheses at every unmarked finite place of the shorter period fields. Corollary~\ref{cor:firstglobal909} sums the resulting finite first-coordinate factors, with the specified real frequency weight, over all rational denominators allowed by the unselected marks. Theorem~\ref{thm:realquadrics911} bounds the plane wave on every real frame quadric. Corollary~\ref{cor:realgramposition911} gives the required frequency power for the first Gram wave at a fixed position, with the conditional correlation measured in its orbit density. Theorem~\ref{thm:realpositionball912} integrates the first real Gram wave over a ball of positions. Its compact-type factor and Haar constant are computed in the normalizations of the period. Lemma~\ref{lem:commondirection913} bounds negative moments of a common dual pairing. Proposition~\ref{prop:positiondrop913} relates this pairing to both the next real position and the unipotent integration variables. Theorem~\ref{thm:realremotetail914} controls the first real correlation over all positions outside a norm ball, uniformly in the frequency. Its proof uses the first small common dual pairing and the square integral on the remaining flag. Theorem~\ref{thm:fullrealfirst915} proves the first real Gram-frequency estimate on the whole position group. Corollary~\ref{cor:adelicfirst915} inserts these real integrals into the rational first-frequency sum on the same period fields as the finite factors. Theorem~\ref{thm:realcutmass916} compares the zero-frequency mass of the two cut tests to the discrete-series formal degree. Proposition~\ref{prop:realpivotdensity916} recovers the continuous first-pivot density from the integrated real transform. Theorem~\ref{thm:conditionedpivot917} gives the fiber measure at a prescribed nonzero first real pivot and identifies the second Gram entry on that fiber with a residual Hermitian pairing. Corollary~\ref{cor:conditionedwave917} bounds its Fourier operator at a fixed position. Theorem~\ref{thm:secondallpositions918} integrates the second real Gram wave over all positions at each prescribed nonzero first pivot. The bound retains the cost of that pivot explicitly. Theorem~\ref{thm:crossfiber919} bounds the two branches on a prescribed first Gram fiber by crossed common-dual weights. Proposition~\ref{prop:twistedfiber919} controls the weighted position measure, including different first coordinates at the two ends. Theorem~\ref{thm:secondsharp920} bounds the entire second real Gram transform for arbitrary prescribed first entries at both ends. Its losses at small first entries have exponent independent of the large rank \(k\). Theorem~\ref{thm:firstmismatch921} proves decay in the difference of the reciprocal first pivots, including the full position integral when the next frequency is zero. Theorem~\ref{thm:flagmeasure922} computes the measure for a pair of ordered real null flags at any list of nonzero prescribed pivots. Proposition~\ref{prop:flagsection922} bounds a common isometry in terms of the prefix determinants. Theorem~\ref{thm:crossflags923} bounds both conditional moments for prescribed real flags. Propositions~\ref{prop:flagenergy923} and~\ref{prop:flagtransfer923} control the residual position with arbitrary endpoint scales. Theorem~\ref{thm:allrealwaves924} controls every nonzero real Gram wave, after any number of prescribed null pivots, by the required product of prefix-determinant weights on the shorter interval. Theorem~\ref{thm:allmismatch925} proves decay in the difference of the reciprocal last pivots for arbitrary prescribed real heads, with the same prefix weights and on the same interval. Theorem~\ref{thm:realnulldescent926} identifies adjoining a common real pivot with inverse Fourier transformation at every nonzero pivot, with absolute convergence through the position integral. Theorem~\ref{thm:terminaltensor927} identifies the terminal density with the Gram matrix of the normalized orbital tensors in the original Haar measures. Corollary~\ref{cor:realsystem927} assembles the real-place estimates and Fourier identities for every row. Theorem~\ref{thm:prefixpack928} sums the actual scaled real and finite determinant weights over all rational prefixes on the period fields. Its bound is independent of the exceptional set of finite primes. Theorem~\ref{thm:identityfrequency929} sums the real transforms and the finite frequency weights over all rational prefixes in units of the real identity mass. Corollary~\ref{cor:actualidentity929} obtains the corresponding bound for the actual finite integrals at the first row. Theorem~\ref{thm:finiteunitreduce930} reduces the original two-head form with unit prefixes at any standard unmarked finite place, over all positions, to the first form for the residual lattice. Lemma~\ref{lem:finiteunitdensity930} computes its measure factor. Proposition~\ref{prop:levelmass678} computes the Plancherel mass of the permitted residue types at the two level primes in that construction. Theorem~\ref{thm:definitelevels678} bounds both the discarded dimension on the definite automorphic space and the integral index of its finite level projection. Theorem~\ref{thm:besselfourier679} computes the Fourier transform of the conjugation average of the Bessel character at a split place. Corollary~\ref{cor:markedbessel679} uses it to bound the mass discarded by the two projections on the finite Bessel character of the actual period. Theorem~\ref{thm:rationalflags680} classifies the intersections of agreeing rational frames. Corollary~\ref{cor:parabolicslice680} removes the off-Bessel terms preserving the isotropic space on either marked slice. Theorems~\ref{thm:discriminant681} and~\ref{thm:terminalmeet681} give the split-prime divisors of the joint Gram determinant and of the two terminal columns after saturation. Theorem~\ref{thm:realgram682} computes the Haar density over a pair of real frames. Corollary~\ref{cor:realmarked682} evaluates it at the rational terms on the marked slices. Theorem~\ref{thm:gramrestriction683} bounds the value of a localized Bessel convolution at a prescribed Gram matrix in terms of its quotient norms. Theorem~\ref{thm:archbessel684} proves the formal-degree bound for the actual translated discrete-series average. Theorem~\ref{thm:realpair685} bounds the cut kernel of one rational open-cell term over all real places, including weights coupling those places. Theorem~\ref{thm:firstnormalsum686} sums the squared open-cell main bound over distinct rational frame products and first positive normals in a fixed class position. Theorems~\ref{thm:termdisc687} and~\ref{thm:samelink687} compute the integral lift discriminant and prove the marked-prime gain for two terms with a common positive link. Theorem~\ref{thm:openpair688} gives the level gain for any pair of open terms from the codimension of their common full link. Theorem~\ref{thm:allpair689} proves the arithmetic discriminant estimate for every pair of surviving rational cell ranks on the same class lattices and levels. Theorem~\ref{thm:latticemean690} bounds the full weighted sum of lattices of each corank, averaged over the actual marked definite classes on this interval. Theorem~\ref{thm:paircount691} bounds the weighted count of all surviving near pairs in agreeing blocks, including the multiplicity of rational frames over each lattice. Theorem~\ref{thm:schurmean692} bounds the translation mean of the projected frame weight over the marked positive column lattice at every nonzero cell rank. Corollary~\ref{cor:closedmoment693} bounds the absolute cut contributions of all positive closed ranks in the squared period sum on agreeing blocks. Corollary~\ref{cor:normalrange694} estimates the cut second moment from open terms in an explicit range of normal lengths, using the compact angular integral of Lemma~\ref{lem:terminalrotation694}. Theorem~\ref{thm:effectivestationary695} constructs a stationary position of length at most \(b^{5n}\) at every real place, and Corollary~\ref{cor:exactphase695} proves its exact trace identity. Theorem~\ref{thm:angularlocalization696} localizes the integrand modulo the common compact subgroup. Theorem~\ref{thm:positivecore696} bounds its positive integral on a specified neighborhood at all real places. Theorem~\ref{thm:noncompacttail697} bounds the full absolute integral outside a shrinking tube about the maximal compact subgroup, uniformly at each real place. Theorem~\ref{thm:fullrealperiod698} gives a positive quantitative lower bound for the entire real Bessel integral at every place of the common field. Theorem~\ref{thm:normalizedcut699} normalizes the nonzero real test vectors, proves their decay, and estimates the error of the cut used in the global sum. Theorem~\ref{thm:globalamplitudetail700} proves absolute convergence over rational translates and bounds the difference between the full and cut sums on every residue row. Theorem~\ref{thm:besselmatrix701} evaluates the Bessel terms and Corollary~\ref{cor:parabolicmass701} computes the entire parabolic contribution on either marked slice. Theorem~\ref{thm:signedopen702} sums the signed contributions of the remaining open orbits on each marked slice with an explicit archimedean factor. Theorem~\ref{thm:markedsecond703} bounds the full signed second moment on both marked slices, and Corollary~\ref{cor:identityblock703} bounds the deviation of the identity block from the identity operator. Theorem~\ref{thm:poincaregram704} realizes the matrices as Petersson products. Theorem~\ref{thm:singleprojection704} compares the actual projected series at one mark with the corresponding finite Bessel operator. Theorem~\ref{thm:conjugacyfourier705} computes a finite Fourier bound by counting cyclic vectors. Theorem~\ref{thm:projectedrank705} bounds the rank lost after both prescribed cuspidal projections. Proposition~\ref{prop:allresidue706} bounds the second moment over all residues. Theorem~\ref{thm:translates706} and Corollary~\ref{cor:avoidspace706} give dimensions in the cuspidal target obtained by finite translates. Theorem~\ref{thm:besselperiod707} identifies each Petersson pairing with its global Bessel period, including its Haar and formal-degree factors. Corollary~\ref{cor:periodrank707} gives the dimension of the projected period functionals. Theorem~\ref{thm:firstlift708} evaluates the first theta lift of the projected Poincar\'e series by an absolutely convergent Bessel integral. Corollary~\ref{cor:firstprojection708} gives its finite projection as one integral combination of translated first-lift values after scaling. Theorem~\ref{thm:finitegram709} computes the finite factors of the first and auxiliary theta lifts and bounds the normalized convolution in operator norm. Corollary~\ref{cor:finitedet709} controls their total logarithmic determinant error. Lemma~\ref{lem:bilinearsource710} identifies the conjugate source required by the bilinear lift. Theorem~\ref{thm:realgram710} gives a positive lower bound for its real Gram factor with the specified harmonic test. Theorem~\ref{thm:globalgram711} computes the Petersson Gram operator of the first and auxiliary lifts. Corollary~\ref{cor:firstliftrank711} proves the first-lift rank and compares the actual Petersson determinants across tame twists. Theorem~\ref{thm:centraltraces635} separates the signed full-dimensional contributions at distinct projective isometries of polynomial lattice index in the second trace and bounds the noncentral fourth-trace terms. At the full retained index, Corollary~\ref{cor:retainedtraces635} bounds the terms outside a small rank neighborhood of rational scalars. Theorem~\ref{thm:scalarseparate636} separates distinct scalar rotations in the original integral coordinates. Theorem~\ref{thm:centralfour636} bounds the fourth-trace contribution of nonidentity eigenvalues of large multiplicity and reduces the central sum at polynomial matched index to exact products equal to the identity. Theorem~\ref{thm:rowsumgram637} computes the ordinary Gram norms and diagonal defects of the full-dimensional portion at polynomial matched index from signed correlations at rational products, with relative error uniform in the number of target rows. Theorem~\ref{thm:disjointtransports638} gives disjoint integral coordinates for distinct rational transports below the stated matched index. Corollaries~\ref{cor:exactdiagonal638} and~\ref{cor:alignedrows638} compute the full-dimensional diagonal Gram block and the mixed block for source rows with a small marked identification. Theorem~\ref{thm:terminalouter639} bounds the sum of distinct isotropic spaces with deep terminal pairings. Corollary~\ref{cor:terminalidentity639} evaluates the repeated-row identity coefficient for equal terminal reductions with an error uniform on the common intervals. Theorem~\ref{thm:terminalclass640} recovers the rational unipotent double class in each fixed positive hyperplane and transport from terminal reduction at precision linear in the flag length. Corollary~\ref{cor:terminalamplitudes640} gives the aggregate error for the original signed coefficients. Theorem~\ref{thm:allterminaltransports641} recovers the positive transport in both common dimensions and the two hyperplanes when present. Corollary~\ref{cor:allterminalcoeff641} groups the signed coefficients for a whole ordinary block with uniform total error. Theorem~\ref{thm:fullrational645} recovers the full-dimensional rational unipotent class from a bordered frame minor and the positive transport. Corollary~\ref{cor:allrational645} recovers the signed operators class by class at both common dimensions with one error for the original block.

\subsection{How the source states the final implication}\label{nav:I-34}
The theorem on line 10 of \nolinkurl{Main.lean} selects the
critical-strip assertion from the three conclusions of
\srcloc{Supplements/Kernels/MainTheorems}{15}.
The analytic deduction takes a zero in the strip to a zero
of the inducing primitive character. The missing Euler factors
cannot vanish there, since \(\lvert\chi(p)p^{-s}\rvert<1\).
The functional equation for the primitive character takes
a zero with \(\real s<1/2\) to a zero at \(1-s\) for the
inverse character. Its real part exceeds \(1/2\).
Consequently the asserted absence of zeros in the right half
of the strip excludes both possibilities off the line.
The source gives this implication in
\srcloc{Development/AutomorphicArithmetic/CharacterBoxes/ClaimStrip}{59}.

The source supplies the right-half-plane assertion through
\srcloc{Development/AutomorphicArithmetic/CharacterBoxes/MainConclusion}{10}.
One of the two branches in line 8 there supplies the
cohomological window. In the other branch, the theta
construction uses the simultaneous bounds for a common
family on the window. Its five clauses are written in
\srcloc{Formalism/SharpRestriction/GlobalTheoremInputs}{158}.
The family occurs as the same argument in each clause.
Specializing the cohomological estimate to one family and
the theta comparison to another would leave this inference
without the common family it requires.

\subsection{Locations in the accompanying source}\label{nav:I-35}
The paths below are relative to \nolinkurl{lean/src/} in the supplied archive.
\begin{itemize}\setlength{\itemsep}{1pt}
\item Local restriction kernel. \srcloc{Formalism/GaloisCohomology/Statement}{75}.
\item From a cohomological interval to an auxiliary field. \srcloc{Development/AutomorphicArithmetic/CharacterBoxes/SupplyChi}{8}.
\item The five simultaneous requirements for the period comparison. \srcloc{Formalism/SharpRestriction/GlobalTheoremInputs}{158}.
\item The comparison with its numerical exponents. \srcloc{Development/AutomorphicArithmetic/ThetaCorrespondence/ThetaPairLifts/ThirdPlaceSeries/TopShape}{52}.
\end{itemize}
\clearpage


\subsection{Interpolating the arithmetic count between selected fields}
\label{nav:count-interpolation-666}

Fix a primitive character \(\psi\) of order \(o\). The subtraction of
\(x\log x\) in Definition~\ref{def:family} determines how large the
prime sum can be between two scales where arithmetic estimates are
available. We first give the interpolation error with its sign.

\begin{lemma}[The interpolation error for the family prime sum]
\leavevmode\label{lem:convexcount666}
Let \(\mu\) be a probability measure supported on \([a,b]\), where
\(1\le a\le b<\infty\), and let \(x=\int y\,d\mu(y)\). Put
\[
 c_\psi(n)=1+\frac{1}{3o+1}
             \sum_{\substack{\chi\in\mathcal C_\psi\\\chi\ne\mathbf1}}
                    \real\chi(n),\qquad
 H_n=\int (y-n)_+\,d\mu(y)-(x-n)_+ .
\]
Then \(H_n\ge0\), \(c_\psi(n)>2/3\), and
\[
 \Gg_\psi(x)-\int\Gg_\psi(y)\,d\mu(y)
  =\int y\log(y/x)\,d\mu(y)
    -\sum_{a<n<b}\frac{c_\psi(n)\Lambda(n)}{n}H_n
  \ \le\ \frac{1}{2a}\int(y-x)^2\,d\mu(y).
\]
For the probability on \(a<b\) with atoms only at its two endpoints,
the last upper bound equals \((x-a)(b-x)/(2a)\).
\end{lemma}
\begin{proof}
There are \(o-1\) terms besides the principal one, and each real
part is at least \(-1\), including at primes dividing a conductor.
Thus \(c_\psi(n)\ge 1-(o-1)/(3o+1)>2/3\).
The definitions give a real constant \(L_\psi\) such that
\[
 \Gg_\psi(y)+y\log y-L_\psi y
   =\sum_{n\le b}\frac{c_\psi(n)\Lambda(n)}{n}(y-n)_+
       \qquad (a\le y\le b).
\]
Integrate this identity and subtract its value at the mean.
The constant times \(y\) cancels. Each ramp is convex and is
affine on the support when \(n\le a\) or \(n\ge b\). This proves
the identity and the claimed sign. On \([a,b]\) the second
derivative of \(y\log y\) is \(1/y\le1/a\). Taylor's inequality
about \(x\), after the linear term integrates to zero, gives
the upper bound. The endpoint masses are \((b-x)/(b-a)\) and
\((x-a)/(b-a)\), whose variance is \((x-a)(b-x)\).
\end{proof}

For a node \(t\) at which conditions~\ref{cond:sizes}--\ref{cond:weights}
hold, use the data at \(t\) to define the real number
\[
 D(t)=\frac{2}{Q\delta\,|A_0|}\sum_{u\in A_0}\Delta_u .
\]
Here and below the symbols in a node value refer to that node's
field, weights and translates. A node can therefore be used even
if a neighboring node has a different dyadic layer.

\begin{theorem}[A positive quadrature bound on the common intervals]
\leavevmode\label{thm:nodecount666}
Fix \(0<\eps<1/52\) and a common interval \([Y,Y^{3/2}]\).
Use the same two auxiliary primes throughout. Take nodes
\(Y=t_0<\cdots<t_N=Y^{3/2}\) with
\[
 t_{j+1}-t_j\le C t_j^{3/4+2\eps}.
\]
Assume conditions~\ref{cond:sizes}--\ref{cond:weights} at the
nodes. Write \(\lambda_j(x)\) for the continuous, piecewise affine
function on this partition taking value one at \(t_j\) and zero
at the other nodes. For any finite positive measure \(\nu\) on the
interval put \(w_j=\int\lambda_j\,d\nu\). Uniformly over \(\nu\)
and the nodes, for sufficiently large \(Y\),
\[
  \int\Gg_\psi(x)\,d\nu(x)
   \le \sum_{j=0}^N w_jD(t_j)
          +O_{\psi,\eps,C}\left(\int x^{1/2+4\eps}\,d\nu(x)\right).
\]
One can choose such nodes with \(O_{\eps}(X^{1/4-2\eps}+1)\)
in each dyadic shell \([X,2X]\) inside the interval, including
every transition of the cyclotomic layer of
Proposition~\ref{prop:field12}.
\end{theorem}
\begin{proof}
At a node \(t\), average the defining count over translates,
as in Lemma~\ref{lem:av10}. Replacing the count by the idealized
sum costs at most \(2Q/t\) before normalization. After division
by \(Q\delta/2\), the retained terms give \(\Gg_\psi(t)\) up to
\(O_\psi((1+t^\eps)(1+\log^2 t))\) from replacing the arguments
and the character proportions. The other terms are bounded in
absolute value by \(8o\,t^{1/2+\eps}\) by
condition~\ref{cond:chars}. Since \(\delta\ge t^{-\eps}\), this
proves \(|D(t)-\Gg_\psi(t)|\ll_{\psi,\eps}t^{1/2+\eps}\).

For \(t_j\le x\le t_{j+1}\), apply Lemma~\ref{lem:convexcount666}
with the endpoint masses. Its upper error is at most
\((t_{j+1}-t_j)^2/(8t_j)\ll_C x^{1/2+4\eps}\).
Since \(3/4+2\eps<1\), both endpoints are between \(x/2\)
and \(2x\) for \(Y\) large. The endpoint masses are precisely the
two linear interpolation weights. Substitution of the estimate
at the nodes and integration against \(\nu\) give the assertion.

For the count of nodes split the interval at the powers of two.
In the part of each shell \([X,2X]\), \(X\) a power of two,
divide into equal subintervals of length at most
\(X^{3/4+2\eps}\), using the ceiling of the length divided by
this quantity. Include the layer transition points
\(x^{2+8\eps}=2^l\). There are \(O(1)\) of these per shell because
consecutive transition points have ratio \(2^{1/(2+8\eps)}\).
Adding nodes can only reduce the gaps. Any interval \([X,2X]\)
meets a bounded number of these shells.
\end{proof}

\begin{corollary}[The count comparison at selected scales]
\leavevmode\label{cor:sparsenodes666}
Suppose that for a fixed \(\psi\) and arbitrarily small
\(\eps>0\) there are arbitrarily distant common intervals with
such nodes, on each of which the signed mean satisfies
\(D(t_j)\le A_{\psi,\eps} t_j^{1/2+9\eps}\) for every \(j\).
Then every open-strip zero of the retained family lies on the
critical line. This applies in particular if
condition~\ref{cond:gap} holds at just those nodes with the
specified data.
\end{corollary}
\begin{proof}
Take \(\nu\) to be the point mass at an arbitrary \(x\) in the
interval. Only its two neighboring node values enter. The theorem
gives \(\Gg_\psi(x)\ll_{\psi,\eps,C,A_{\psi,\eps}}x^{1/2+9\eps}\)
on the whole interval, for a mesh constant uniform in the chosen
intervals. Theorem~\ref{thm:detect}, followed by
\(\eps\downarrow0\), bounds all the real parts by \(1/2\).
The functional equation and closure under conjugation give the
other inequality. Under condition~\ref{cond:gap} the bound on
\(D(t_j)\) follows directly from the four estimates for
\(\mathcal B/Q\) in Proposition~\ref{prop:balance} and
\(\delta^{-1}\le t_j^\eps\).
\end{proof}

This removes the cost of interpolating a cohomology group or a
period lattice between dyadic layers. The upper estimate on the
signed node values is still required.


\subsection{Three averages of the signed count on a common interval}
\label{nav:three-averages-667}

We retain the arithmetic mean \(D(t_j)\) and the interpolation
functions \(\lambda_j\) of Theorem~\ref{thm:nodecount666}. There
is an alternative to an upper estimate at each node. A zero to
the right of a proposed line can be detected by three signed
averages at a single rational frequency.

Set \(\sigma=256\). For \(U=\log Y\), put \(k=\lceil5\sigma U/4\rceil\)
and \(K_k(u)=u^k e^{-\sigma u}/k!\). For a rational \(\tau\) and
\(l=0,1,2\) use the following nonnegative weights on the nodes,
which depend only on the mesh, the interval and the frequency.
\[
 w_{lj}=\int_U^{3U/2} K_k(u)
   \bigl(1+\cos(\tau u+2\pi l/3)\bigr)\lambda_j(e^u)\,du .
\]
In particular the same tuple of fields and translates is used in
the three averages. The node counts themselves can have either sign.

\begin{theorem}[Detection by three signed averages]
\leavevmode\label{thm:threeaverages667}
Fix \(\psi\), \(0<\eps<1/52\), and \(1/2+4\eps\le\theta<1\).
Suppose that for every rational \(\tau\) there are arbitrarily
distant common intervals, with nodes and data as in
Theorem~\ref{thm:nodecount666}, on which
\[
 \max_{0\le l\le2}\sum_{j=0}^N w_{lj}D(t_j)
     \ \le\ B_\tau(\sigma-\theta)^{-k-1}.
\]
The constant \(B_\tau\) and the mesh constant may depend on
\(\psi,\eps,\theta,\tau\) and are uniform over the selected
intervals at that frequency. Then every open-strip zero of the
retained family has real part at most \(\theta\).
Conversely, if these zeros have real part at most \(\theta\),
the displayed bound holds on every sufficiently distant such
interval, with a constant independent of \(\tau\) and of its nodes.
\end{theorem}
\begin{proof}
Put \(g(u)=\Gg_\psi(e^u)\), and define \(A_k(s)\) from its Laplace
transform as in Section~\ref{sec:landau}. Write \(D_*=\sigma-\theta\).
Use Theorem~\ref{thm:nodecount666} with the image under \(u\mapsto e^u\)
of \(K_k(u)(1+\cos(\tau u+2\pi l/3))\,du\) restricted to the
interval. The error is at most a constant times
\[
 2\int_0^\infty K_k(u)e^{\theta u}\,du=2D_*^{-k-1}.
\]
This is an upper estimate for a signed integral since the measure
used for interpolation is positive.

We can extend all three integrals to \(0\le u<\infty\) with
an absolute error \(O_\psi(\sigma^{-k})\).
Indeed \(|g(u)|\le C_\psi e^{2u}\). In the tail estimates of
Lemma~\ref{lem:kernel} take \(B=2,\ c=5/4,\ v=3/2\). The
two exponential rates after multiplication by \(\sigma^k\) are
\[
   2+\sigma(1/4-(5/4)\log(5/4)),\qquad
   3+\sigma(-1/4+(5/4)\log(6/5)).
\]
Both are negative at \(\sigma=256\). For example the alternating
series for \(\log(1+t)\) gives
\(\log(5/4)>1/4-(1/4)^2/2\) and
\(\log(6/5)<1/5-(1/5)^2/2+(1/5)^3/3\).
The bound on the cosine factor is two at every frequency.

The extended integrals equal
\(A_k(\sigma)+\real(e^{-2\pi i l/3}A_k(\sigma+i\tau))\) by
differentiation under the integral, since \(g\)
is real. Of three equally spaced points on the unit circle at
least one has angular distance at most \(\pi/3\) from any given
direction. Maximizing therefore gives
\[
 A_k(\sigma)+\tfrac12|A_k(\sigma+i\tau)|
       \ll_{\psi,\eps,\theta,\tau}D_*^{-k}.
\]
For each fixed \(\tau\) this holds along an unbounded sequence of
\(k\). We explain why such separate sequences suffice.

Let \(Z\) be the set of poles of the transform in \(\real s>0\).
They are simple and their residues are nonzero by
Lemma~\ref{lem:mellin}. Suppose a pole satisfies \(\real\rho>\theta\).
Set \(D=\min(D_*,\operatorname{dist}(\sigma,Z))\).
If some \(\rho_1\) has \(\sigma-\real\rho_1<D\), take \(R\)
strictly between these numbers. A disk of radius \(R\) about
\(\sigma+i\tau\) contains \(\rho_1\) for all \(\tau\) in an open
interval about \(\imag\rho_1\). Shrink to a bounded interval.
Only finitely many poles meet the resulting closed disks, by
compactness in \(\real s>0\). Each pair of distinct poles has
an equidistance at at most one height, since both are to the
left of \(\real s=\sigma\) and poles of the same height have
different horizontal distances. We can thus choose a rational
\(\tau\) for which there is a unique nearest pole \(\rho\), at
distance \(d<R\). Choose \(d<r<R\) so there is no further pole
in its closed disk. At the real center the closed disk of radius
\(r\) is pole-free. Subtracting the simple polar part and using
Cauchy's inequality as in \eqref{eq:cauchy} yields
\[
 \tfrac12|\Res_\rho\Phi_\psi|d^{-k-1}-O(r^{-k})
     \ll_{\psi,\eps,\theta,\tau}D_*^{-k}.
\]
Multiplication by \(d^k\) along the sequence for this rational
frequency contradicts the nonzero residue.

In the remaining case the horizontal distance of every pole is
at least \(D\). The proposed pole gives \(D<D_*\), so a pole
attains the distance \(D\) from \(\sigma\). It must have zero
imaginary part, and is the unique nearest pole. Its residue is
positive by Lemma~\ref{lem:mellin}. Choose \(D<r<D_*\) giving
no other pole in that closed disk. At \(\tau=0\) Cauchy's
estimate gives \(A_k(\sigma)=(\Res_{\sigma-D}\Phi_\psi)D^{-k-1}
+O(r^{-k})\), which contradicts the same upper bound along the
sequence for zero frequency.

For the converse, the absolutely convergent explicit formula
\eqref{eq:uncond12} gives \(|\Gg_\psi(t)|\ll_{\psi,\theta}t^\theta\).
The node error in Theorem~\ref{thm:nodecount666} gives the same
bound on \(|D(t)|\). If adjacent nodes have ratio at most two,
their \(\theta\)-powers are at most \(2^\theta x^\theta\) at each
intermediate \(x\). This ratio condition holds on sufficiently
distant intervals for each fixed mesh constant. Now sum with the
nonnegative weights, use \(\sum_j\lambda_j=1\), and integrate
as in the first display of the proof.
\end{proof}

In particular one may take \(\theta=1/2+4\eps\) and let \(\eps\)
decrease to zero. The functional equation then yields the critical
line for every character whose family satisfies the three bounds.
The computations establish the weights and the detection theorem.
The upper bound on those three averages of the actual arithmetic
count is the unresolved estimate in this formulation.

\section{The field and the auxiliary weights}\label{nav:g-I-1}

\subsection{Integral weights and theta coefficients on the same interval}\label{nav:II-04}
\label{sec:period14}

Fix the field \(E\) of Proposition~\ref{prop:field12} at each \(x\) of its interval. In the next two results the auxiliary primes \(q,d\) stay fixed throughout that interval. First we choose integral weights for the prime sum. We then give an integral period comparison for theta coefficients of positive definite Hermitian lattices over \(E\).

\subsubsection{Choosing the small weights uniformly}\label{nav:II-05}
\begin{lemma}\label{lem:far14}
There are constants \(c>0\) and \(Y_0\), depending only on \(F,\eps\), such that, for \(Y\ge Y_0\), every primitive character \(\alpha\) belonging to \(\Gal(E/\QQ)\) satisfies
\[
 |G_\alpha(y)|\le y\exp(-c x^{\eps/8})
 \qquad\left(y\ge \tfrac13\exp(x^{\eps/3})\right).
\]
\end{lemma}
\begin{proof}
Write \(N=2^{l+2}Fqd\), so every conductor in question divides \(N\). We first remove the possible exceptional real zero from the classical zero-free region. Every real character of this field is the product of the two auxiliary quadratic characters to powers zero or one and a real character of \(E_0=K_\psi C_l\). The latter has odd conductor part dividing \(F\). Its local character at two has order at most two, hence factors modulo eight, since every unit congruent to one modulo eight is a square in \(\ZZ_2^\times\). Its primitive conductor therefore divides \(8F\). If an auxiliary factor occurs, Theorem~\ref{thm:twists} excludes every real zero in \([49/50,1]\). Otherwise there are finitely many possible real primitive characters, independently of \(Y\), and Lemma~\ref{lem:boundary} gives a fixed zero-free neighborhood to the left of one for each. The principal character is covered by the classical zeta zero-free region. Thus the usual de la Vall\'ee Poussin theorem, whose sole possible exception for a primitive character is a real zero of a real character, gives
\[
    1-\real\rho\ge
        \frac{c_F}{\log(2N(3+|\imag\rho|))}
\]
for every nontrivial zero under consideration. The same lower bound for \(\real\rho\) follows by the functional equation applied to the conjugate character.~\cite{I-I-I-elkies,I-I-I-tz}

Here are the uniformities needed in \eqref{eq:uncond12}. In a band \(|\imag\rho-u|\le1\) the number of zeros is \(O(\log(2N(3+|u|)))\), with an absolute constant. This follows by taking the real part of the completed logarithmic derivative at \(2+iu\), as in Section~\ref{sec:sharp}; the Euler product there is bounded absolutely and the gamma term and conductor give the displayed logarithm. Zeta has two additional rational terms, both bounded on that line. For \(|\imag\rho|\le1\) each reciprocal \(|\rho(1-\rho)|^{-1}\) is \(O_F(\log^2(2N))\) by the two zero-free bounds. Summing the bands at larger ordinates gives
\[
  \sum_\rho\frac{1}{|\rho(1-\rho)|}\ll_F\log^3(2N),\qquad
  \sum_{|\imag\rho|>T}\frac{1}{|\rho(1-\rho)|}
       \ll\frac{\log(2N(3+T))}{T}\quad(T\ge2).
\]
For the constant \(d_0\) in \eqref{eq:uncond12}, logarithmically differentiate the functional equation at zero and remove its simple gamma pole when present. For a nonprincipal character this bounds \(|d_0|\) by \(O(1+\log(2N)+|L'(1,\bar\alpha)/L(1,\bar\alpha)|)\). The local zero expansion obtained by subtracting the completed logarithmic derivative at \(2\) from the one at \(1\) bounds the last term by \(O_F(\log^2(2N))\). Indeed the nearby zeros have distance at least \(c_F/\log(6N)\) from \(1\), and the other bands contribute \(O(\log(2N))\) by summation with weights \(j^{-2}\). Zeta has a fixed \(d_0\). The coefficient \(\lambda\) in the same formula is either zero or one. Finally, on \(\real s=-1/2\) the functional equation and absolute convergence at \(1-s\) give \(|L'/L(s,\alpha)|\ll\log(2N(3+|\imag s|))\). Integration against \(1/|s(s-1)|\) makes the remainder \(O(y^{-1/2}\log(2N))\).

Set \(L=\log y\), \(T=\exp(\sqrt L)\). For \(|\imag\rho|\le T\) we have
\[
 y^{\real\rho-1}
 \le \exp\!\left(-c'_F\min\!\left(\sqrt L,\frac{L}{\log(2N)}\right)\right).
\]
By construction \(\log(2N)\ll_{F,\eps} x^{\eps/5}(\log x)^2+\log(2x)
\ll_{F,\eps} x^{5\eps/24}\). Also \(L\ge x^{\eps/3}-\log3\). The minimum is bounded below by a constant times \(x^{\eps/8}\), since \(1/3-5/24=1/8\). Insert this estimate, the mass and tail bounds, and the preceding remainder estimates into \eqref{eq:uncond12} divided by \(y\). Powers of \(\log(2N)\) are absorbed by decreasing \(c\); the tail is \(O((\log(2N)+\sqrt L)e^{-\sqrt L})\). The same bound with a smaller \(c\) absorbs \((L+1+\log^2(2N))e^{-L}\). This proves the lemma uniformly for all larger \(y\).
\end{proof}

\begin{proposition}[Integral prime-sum weights]\label{prop:weights14}
In Proposition~\ref{prop:field12}, take the distinguished weights to be \(\lceil b\rceil\), all others to be \(1\), and \(a_t=\lfloor bx\rfloor\). These choices satisfy conditions~\ref{cond:sizes}--\ref{cond:weights} for every \(x\) on the same interval. More generally, the weight estimate in condition~\ref{cond:weights} holds if \(|a_t/b-x|\le1\), every distinguished weight is in \([b,b+\sqrt b]\), and each other positive weight \(w\) satisfies \(a_t/w\ge\exp(x^{\eps/3})/3\).
\end{proposition}
\begin{proof}
For a nondistinguished weight, Lemma~\ref{lem:far14} gives
\[
 |w G_{\xi^*}(a_t/w)|\le b(x+1)e^{-c x^{\eps/8}}.
\]
There are at most \(f|B|\) terms before division by \(f\) in \eqref{eq:count}, and \(|B|\ll_F x^3\). Their combined absolute value is at most \(b/x\) for large \(Y\). For the distinguished weights the derivative estimate after condition~\ref{cond:gap} bounds their total change by \(|B|h_2 H_x\sqrt b/f\). Since \(h_2\le f\), \(H_x=3(x+1)\log(2x+4)\), and \(b=\exp(x^\eps)\), this too is at most \(b/x\). This proves the more general assertion. For the stated integer choices \(|a_t/b-x|\le1/b\), and \(a_t\ge\exp(x^{\eps/3})/3\). All the other size conditions are those already proved in Proposition~\ref{prop:field12}. The points \(a_t/b\) still belong to \([Y,Y^2]\) for large \(Y\), so Theorem~\ref{thm:twists} supplies the required character inequalities.
\end{proof}

\subsubsection{An integral theta period for chains of Hermitian lattices}\label{nav:II-06}
Write \(k=E^+\) for the fixed field of complex conjugation. For \(0\le j\le m=\lfloor\sqrt{x}\rfloor\), put \(V_j=E^{j+1}\) with the standard positive definite Hermitian form \(H_j\), and fix the lattice \(\mathcal O_E^{j+1}\). Let \(\mathcal C_j\) be its genus, taken up to global isometry. Concretely this is the set of isometry classes of full \(\mathcal O_E\)-lattices locally isometric to the specified lattice at every finite place of \(k\). It is the double coset set for \(U(V_j)(k)\) acting on the finite adelic group modulo the compact open stabilizer of the lattice. The set is finite by the finiteness theorem for algebraic modular forms on a group compact at infinity, which applies since \(H_j\) is definite at every real place of \(k\). The automorphism group of a representative \(A\) is finite as well. It preserves the positive definite real form \(\operatorname{Tr}_{k/\QQ}H_j(v,v)\) on its underlying \(\ZZ\)-lattice. Write \(w(A)=|\operatorname{Aut}(A)|\).~\cite{I-I-II-gross14}

For \(A\in\mathcal C_{j+1}\) and \(B\in\mathcal C_j\), define the theta coefficient
\[
 K_j(A,B)=\#\{s:B\to A\text{ \(\mathcal O_E\)-linear}\mid s^*H_A=H_B\}.
\]
This is the coefficient at \(H_B\) of the Hermitian theta series summing the pullback forms \(s^*H_A\) over \(\operatorname{Hom}_{\mathcal O_E}(B,A)\). The coefficient is finite. For each vector in a fixed \(\ZZ\)-basis of \(B\), its image belongs to a positive definite lattice and has prescribed trace norm, leaving only finitely many choices. Every map counted here is injective because \(H_B\) is positive definite.

For functions on the finite sets use the bilinear period
\[
     [u,v]_j=\sum_{B\in\mathcal C_j}\frac{u(B)v(B)}{w(B)} .
\]
These weights also give integration for which the finite adelic lattice stabilizer has volume one and the compact group at infinity has volume one. Indeed, on the orbit of \(B\) the stabilizer from rational points has cardinality \(w(B)\). Define
\[
 (T_j u)(A)=\sum_B\frac{K_j(A,B)}{w(B)}u(B),
 \qquad (T_j^\dagger v)(B)=\sum_A\frac{K_j(A,B)}{w(A)}v(A).
\]

\begin{theorem}[An endpoint denominator for a chain]\label{thm:end14}
Let \(J\) be any number field and take \(\mathcal O_J\)-valued functions \(f\) on \(\mathcal C_0\) and \(h\) on \(\mathcal C_m\). Each \(T_j\) carries integral functions to integral functions. With \(s_m=(m+1)[E:\QQ]\) and \(D_m=|\operatorname{GL}_{s_m}(\mathbf F_3)|\), the common period
\begin{equation}
 \begin{split}
   P(f,h)&=[T_{m-1}\cdots T_0 f,h]_m\\
         &=[f,T_0^\dagger\cdots T_{m-1}^\dagger h]_0
                                  \ \in\ D_m^{-1}\mathcal O_J
 \end{split}\label{eq:end14}
\end{equation}
satisfies \(\log D_m\ll_F x^7\) at every point of the interval in Proposition~\ref{prop:field12}.
\end{theorem}
\begin{proof}
The group \(\operatorname{Aut}(B)\) acts on the maps counted by \(K_j(A,B)\) by precomposition. If \(s\circ u=s\), injectivity of \(s\) gives \(u=1\). All orbits therefore have \(w(B)\) elements. In fact each orbit specifies exactly one sublattice of \(A\) isometric to \(B\), since two isometries onto the same image differ by an automorphism of their domain. Thus \(K_j(A,B)/w(B)\) is an integer. The matrices of all forward maps have integral entries, so the entire forward composite preserves integral functions.

Expand \([T_ju,v]_{j+1}\). The term indexed by \(A,B\) is \(K_j(A,B)u(B)v(A)/(w(A)w(B))\). Summing first over \(A\) gives \([u,T_j^\dagger v]_j\). All sums are finite. Applying this identity successively yields the two expressions in \eqref{eq:end14} with the same normalization.

An automorphism of an endpoint lattice acts faithfully on its underlying free \(\ZZ\)-module of rank \(s_m\). Reduction modulo three is injective on every finite subgroup of \(\operatorname{GL}_{s_m}(\ZZ)\). To see this, suppose that its kernel contained an element of prime order \(p\). Write it as \(I+3^a C\) with \(a\ge1\) and some entry of \(C\) prime to three. If \(p\ne3\), the binomial expansion to power \(p\) has a nonzero term \(p3^a C\) modulo \(3^{a+1}\). For \(p=3\), the first term after the identity is \(3^{a+1}C\) modulo \(3^{a+2}\). Both contradict the element having order \(p\). Every nontrivial finite kernel would contain such an element. Consequently \(w(A)\mid D_m\) for every endpoint class \(A\). Multiplying the forward period by \(D_m\) clears all its denominators simultaneously. Counting all matrices over \(\mathbf F_3\) gives \(\log D_m\le s_m^2\log3\). On the common family \([E:\QQ]\ll_F x^3\) and \(m+1\le2\sqrt x\), giving the asserted uniform estimate.
\end{proof}

\subsubsection{Integral normalizations in the Bessel comparison}\label{nav:II-07}
The comparison in the \emph{Remark} also uses periods with nontrivial archimedean weights and Hecke characters. The following calculation specifies what an identity for such periods would have to control. Let \(M,L\) be full \(\mathcal O_J\)-lattices in finite-dimensional \(J\)-spaces and let \(P\ne0\) be a \(J\)-bilinear form. Define its content \(I(P;M,L)\) to be the \(\mathcal O_J\)-ideal generated by the values \(P(M,L)\), allowing a fractional ideal. This ideal is finitely generated because the lattices are. For \eqref{eq:end14} the two formulas have equal content and \(D_m I\subseteq\mathcal O_J\).

\begin{proposition}\label{prop:content14}
Suppose forms \(P,P'\) are related over \(J\) by \(P(f,h)=aP'(Uf,Vh)\) with \(a\in J^\times\) and \(J\)-linear isomorphisms \(U,V\). Then
\[
 I(P;M,L)=a\,I(P';UM,VL).
\]
At a finite place \(v\) write \(\pi\) for a uniformizer. Suppose endpoint lattices \(M',L'\) obey
\[
 \pi^i M'_v\subseteq UM_v\subseteq\pi^{-i'}M'_v,\qquad
 \pi^j L'_v\subseteq VL_v\subseteq\pi^{-j'}L'_v
\]
for nonnegative integers \(i,i',j,j'\). Put \(b_v=v(I(P;M,L))-v(a)-v(I(P';M',L'))\). Then \(-i'-j'\le b_v\le i+j\).
\end{proposition}
\begin{proof}
The first equality holds on the generating values in both directions. Localize it at \(v\). Bilinearity and the inclusions imply that the image content lies between \(\pi^{i+j}I(P';M',L')_v\) and \(\pi^{-i'-j'}I(P';M',L')_v\), before multiplication by \(a\). Taking valuations gives the inequalities.
\end{proof}

At infinity fix Hermitian norms and let \(\|P\|_v\) denote the operator norm. If \(U,V\) are isometries at every such place, then \(\|P\|_v=|a|_v\|P'\|_v\). The product formula
\[
 \sum_{v\mid\infty}[J_v:\RR]\log|a|_v
       =\sum_{v\nmid\infty}v(a)\log|\mathcal O_J/\mathfrak p_v|
\]
shows exactly how the scalar cancels in the difference between the finite content and the logarithmic archimedean norms. The remaining finite discrepancies are the \(b_v\). Even in dimension one an integral theta coefficient may be multiplied by an arbitrary integer by scaling its integral test function. Its finite content and its archimedean norm both acquire that factor.

The coefficient maps in Theorem~\ref{thm:end14} count Hermitian lattice embeddings. The periods required in \eqref{eq:gap} involve specific automorphic representations, projections onto specified weights, and a common Bessel lattice for all translates \(u\). To apply Proposition~\ref{prop:content14} one must identify the projected periods over a coefficient field, bound the lattice indices after projection, and estimate the archimedean norms on that same family. The endpoint denominator in \eqref{eq:end14} accounts for stabilizers in the definite lattice sum. It gives no bound for the indices or norms of those additional projections.

\subsection{The local ring and the nonzero auxiliary period on one family}\label{nav:II-23}
Keep \(E=K_\psi C_l(\sqrt{-q},\sqrt{-d})\) on the common
power interval. The index \(l\) may be the shortened one of
Proposition~\ref{prop:desc50}. The real subfield has degree
\([C_l:\QQ]\gg_\psi x^{2+8\eps}\) in either choice.
Let \(K=\QQ(\sqrt{-d})\), \(\Phi=\Gal(E/K)\),
\(b=\exp(x^\eps)\), \(a=\lceil xb\rceil\),
\(W=\lceil\sqrt b\rceil\), and \(B'=\lceil b\rceil\).
Use the orientation and translation by complex conjugation of
Theorem~\ref{thm:orientation46}. In the following statement
\(j\) denotes a cyclotomic twist and \(r\le s\) are matrix
ranks at one definite theta step.

\begin{theorem}[Simultaneous conclusions for the strict groups and the auxiliary norm]\label{thm:joint75}
At every \(x\) on a sufficiently distant common interval the
following statements hold for this same \(E\).
\begin{enumerate}
\item For \(\ell>\exp(x^\eps/12)\) and
\(1\le j<\lfloor x^{401}\rfloor\), put
\(\Psi=\kappa_\ell^j\alpha\) for any character summand of the
complementary sign \(\alpha(c)=(-1)^j\), with coefficient
discrete valuation ring \(O\), fraction field \(J_\ell\), and
\(D_\ell=J_\ell/O\). Write \(\eta\) for the quadratic character
of \(\QQ(i)\). The local component through the diagonal
representation has an integral formal power series ring with
\(1,2,3,4,2\) variables in weights \(-2,-1,0,1,2\).
The quotient \(C_{-1}\) of the augmentation cotangent module
of the marked global ring representing deformations on its two
lower simple-root lines satisfies
\[
     \operatorname{Hom}_O(C_{-1},D_\ell)\simeq
      \mathcal S(\QQ,D_\ell(\Psi^{-1}))\oplus
      \mathcal S(\QQ,D_\ell(\eta\Psi^{-1})).
\]
Here \(\mathcal S\) means locally zero above \(\ell\) and
unramified at the other finite primes. The first summand
is the corresponding strict character summand for \(E\)
with the convention of Section \ref{nav:IV-09}.
\item Take polynomially bounded ranks \(r\le s\), with \(rs\)
even, and the fixed class representatives and positive split
normalization in
Sections \ref{nav:V-17}--\ref{nav:V-18}.
For a finite family of threshold lists \(W\le w_\tau\le B'\)
and translates \(t\), there is one torsion input at this step
giving nonzero algebraic numbers \(Y_t\) from the corresponding
class sums. Enlarge the coefficient field to a common finite
Galois extension \(J/\QQ\). Then
\begin{equation}
 C_d(t,a,w)\le
       \frac{\log|N_{J/\QQ}Y_t|/[J:\QQ]+B_-^{(S)}(Y_t)}r
              + O_{\psi,\eps}(b/x).                     \label{eq:joint75}
\end{equation}
The allowed threshold depends on the fixed polynomial bound for
the ranks. The error is uniform in the listed translates. The
set \(S\) consists of the good inert primes up to \(a/W\) in
Section \ref{nav:V-17}, and \(B_-^{(S)}\) sums the negative
valuations outside it with the normalized local degrees and
logarithms of the rational primes.
\end{enumerate}
\end{theorem}
\begin{proof}
For the cohomology statement, the three-factor inertia equation
at the two distinct outer eigenvalues detects local coboundaries
in both lower simple-root directions over every finite
coefficient module. The conjugated diagonal family belongs to
the saturated component, so the converse holds there as well.
The marking at complex conjugation and the two further entry
conditions give a unique representative of each pair of strict
classes. Passing to the direct limit over the finite coefficients
gives the displayed integral dual. For \(j\ge2\) the residual
upper-triangular obstruction vanishes by local Tate duality.
For \(j=1\) the two adjacent cocycles split by Kummer theory
into a unit term and a valuation term. Saturating by the two
Frobenius differences kills the valuation coordinates exactly.
This supplies the same twelve coordinates also in the exceptional
first twist. The proofs, including the finite coefficient
passage, are in
Corollary \ref{cor:dual68} and Theorems \ref{thm:smooth69} and \ref{thm:first72}.

For the period statement, place degree \(a\) at \(t\), with
one further nonzero minor degree for each nontrivial cyclic
subgroup of the 2-primary part of \(\Phi\). For a Frobenius
cycle of even order its low half contains exactly one of those
embeddings in the cycle through \(t\). Its degree is at most
the prescribed threshold at every odd power of that Frobenius.
Corollary~\ref{cor:minor67} therefore gives the weighted
divisibility of the actual derivatives. The primes with no
2-primary Frobenius component after composing with conjugation
are \(\pm1\) modulo the conductor of the real cyclotomic layer.
Their total cost and the accumulated rounding error are
\(O_{\psi,\eps}(b/x)\), as proved in Theorems
\ref{thm:count73}--\ref{thm:norm73}. For the positive split normalization each class sum
has strictly positive inner product with the initial vacuum.
The mean squared derivative over translations equals its
positive squared norm times the scaled Fischer norm of the
minor. A torsion grid supported at one fixed split prime is
dense in the common period torus. The intersection of the
finitely many open dense nonvanishing sets therefore contains
one such torsion point. Theorem \ref{thm:nonzero74}, verifies that
the corresponding input gives these values at all the fixed
representatives. Applying the norm formula now proves
\eqref{eq:joint75}. All field selections in this argument precede
the local and theta constructions.
\end{proof}

There is a quantitative obstruction to reading the required
centered estimate directly from this auxiliary norm. The
obstruction applies even after the nonvanishing just proved.

\begin{proposition}[The excess for a character with nontrivial odd part]\label{prop:excess75}
Suppose the order \(o\) of \(\psi\) has a nontrivial odd divisor.
Let \(w\) equal \(B'\) on a subset of a single coset of
\(A_0=\Gal(K_\psi C_l/K_\psi)\) in \(\Phi\), and equal
\(W\) elsewhere. Write \(\omega\) for the proportion of
positions carrying \(B'\). For the actual numbers of
Theorem~\ref{thm:joint75} put
\[
 H_t=\bigl([J:\QQ]^{-1}\log|N_{J/\QQ}Y_t|
                           +B_-^{(S)}(Y_t)\bigr)/r .
\]
Then uniformly throughout the interval
\begin{equation}
       H_t-M_d(a,w)\ \ge\ a\omega\log(B'/W)
                                      -O_{\psi,\eps}(b/x).
                                                        \label{eq:excess75}
\end{equation}
For the sparse densities of Theorem~\ref{thm:rangenew50},
the positive main term here is \(\gg bx\).
It exceeds the permitted error \(x^{5\eps}\mathcal B/(hm)\)
per translate and step for those ranks.
\end{proposition}
\begin{proof}
The odd-order subgroup \(P'\) of \(\Phi\) has nontrivial
image in \(\Phi/A_0\), since \(A_0\) is a 2-group and
\(o\) has an odd divisor. Thus \(t+g+P'\) cannot be
contained in a single coset of \(A_0\). For every nontrivial
cyclic 2-subgroup the minimum defining its minor degree in
equation \ref{eq:degrees73}, is therefore \(W\).
The entire degree list, its scalars \(Y_t\), and the clearing
integers are exactly the same as those for the constant
threshold list \(w_*\equiv W\), using the very same input
and representatives. Apply \eqref{eq:joint75} to \(w_*\).

For this constant list the inertia average summed over all
classes outside \(\Phi\) gives \(1-\xi_d(p^k)\) at each
prime power. Consequently, with \(X=a/W\),
\[
       C_d(t,a,w_*)=W\{R_{\mathbf1}(X)-R_{\xi_d}(X)\}.
\]
The corrections in \eqref{eq:Gchi} and the selected-field estimate
of Lemma \ref{lem:far14}
give \(C_d(t,a,w_*)=M_d(a,w_*)+O(ae^{-c x^{\eps/8}})\).
Here \(X\asymp x e^{x^\eps/2}\) lies in its large-argument
range, \(\xi_d\) is one of its selected twists, and the error
is \(o(b/x)\). The main-term formula of Proposition
\ref{prop:count}, used in the \(K_d\)-orientation, gives
\(M_d(a,w_*)-M_d(a,w)=a\omega\log(B'/W)\). This proves
\eqref{eq:excess75}.

At the specified sparse densities \(\omega=h_2/(2n)\ge
x^{-\eps}/2\), and \(\log(B'/W)\sim x^\eps/2\).
For comparison, dividing the formula in \eqref{eq:gap} by
\(hm\) gives three terms after multiplication by \(x^{5\eps}\).
Use \(a_E=xb\), \(m\asymp x^{1/2-\eps/16}\),
\(h\asymp n x^{-\eps/8}\), \(h_2\le n\),
\(n\asymp x^{2+8\eps}\), and \(k\asymp x^{9/2+10\eps}\).
The three terms are respectively
\[
 O_{\psi,\eps}\!\left(bx^{1/2+(5+1/16)\eps}\right),\quad
 O_{\psi,\eps}\!\left(bx^{1/2+(6+1/16)\eps}\right),\quad
 O_{\psi,\eps}\!\left(bx^{1/2+(1+1/8)\eps}\right).
\]
Each is \(o(bx)\) since \(0<\eps<1/52\).
This establishes the claimed comparison.
\end{proof}

The proposition concerns the positive single-step auxiliary
norm with the particular minor degrees of Volume~XIII. It shows
exactly where that construction loses the odd component of the
distinguished weights. The originally prescribed Bessel period
retains further row, level and character data. Its quantitative
comparison and the uniform sum of the strict lengths in
\eqref{eq:minors33} are still required for Main Theorem~\ref{thm:grh16}.

\subsection{Removing finite cohomology layers and a common theta divisor}\label{nav:II-25}
Keep the fields and ranks of Theorem~\ref{thm:rangenew50},
\(a=\lceil xb\rceil\), \(b=\exp(x^\eps)\), and put
\(P=b^{1/12}\). The prime-to-\(\ell\) characters and the
relations in \eqref{eq:minors33} are those attached to these
fields. We first remove a range of actual lengths from that sum.

\begin{proposition}[The residual sum after the real class-group bound]
\label{prop:res80}
For \(\ell\le P\) retain \(h_\ell\) from \eqref{eq:minors33}.
For \(\ell>P\) replace its value zero by
\[
  h'_\ell=\max\bigl(\{0\}\cup
       \{j\ge1:\ell^{j+1}\le x^{-12}\exp(x^\eps/2)\}\bigr).
\]
Let \(\mathcal R_{**}(x)\) be \(\mathcal R_*(x)\) with these
new thresholds. Then, uniformly on the common intervals,
\[
     0\le \mathcal R_*(x)-\mathcal R_{**}(x)
           \ll_{\psi,\eps} x^{-8}\exp(x^\eps/2).
\]
In particular \(\mathcal R_{**}(x)\le\frac14\exp(x^\eps/2)\)
on such intervals would suffice for \eqref{eq:length12}.
\end{proposition}
\begin{proof}
At a retained \(\ell>P\) decompose the finite relation
quotient into cyclic modules with exponents \(d_i\), as in
Proposition~\ref{prop:minors33}. Replacing the threshold zero
by \(h'_\ell\) decreases its length contribution by precisely
\(\sum_i\min(d_i,h'_\ell)\). This is the length of the subgroup
killed by \(\ell^{h'_\ell}\). The sum over the retained
characters is bounded by the corresponding sum over all
complementary characters since the omitted lengths are
nonnegative. Theorem \ref{thm:cut79} bounds this full sum with its factors \(\log\ell\)
by the displayed error. That theorem uses ordinary ideal class
groups of \((E_x\QQ(\mu_{\ell^{h'_\ell}}))^+\) and distinct
prime-to-\(\ell\) eigenspaces for the different twists. The
terms at \(\ell\le P\) agree. Now use \eqref{eq:minerr33}.
Its first term uses half of \(\exp(x^\eps/2)\) and its remaining
terms, together with the new error, are \(o(\exp(x^\eps/2))\).
This proves the sufficiency assertion.
\end{proof}

Next consider primes which conjugate both quadratic fields.
Under the degree list of Proposition~\ref{prop:center78} they
can give a factor of a period norm which is independent of the
translate \(u\). We determine an explicit such factor without
knowing the centered prime estimate.

\begin{lemma}[Primes moving the high set off itself]\label{lem:move80}
Let \(\mathcal P_x\) consist of the unramified primes
\(5x/4<p\le 7x/4\) inert in both \(K_q,K_d\) for which
\(\pi_p T\cap T=\varnothing\), where \(\pi_p=\sigma_pc_E\)
acts on the CM embeddings. Then
\[
                  \sum_{p\in\mathcal P_x}\log p=x/8+o_{\psi,\eps}(x)
\]
uniformly as the left endpoints of the intervals tend to infinity.
Every \(\pi_p\) in this set has even order.
\end{lemma}
\begin{proof}
On integers coprime to \(qd\), the indicator of the two negative
quadratic signs is \((1-\xi_q)(1-\xi_d)/4\). Consider the ramp
sum \(\sum_n c(n)(y/n-1)_+\) with \(c(n)=\Lambda(n)\) times
this indicator. In \(x\le y\le2x\) both \(q,d\) exceed \(2x\).
Expanding the indicator, using \eqref{eq:Gchi} and applying
Theorem~\ref{thm:twists} to the three nonprincipal characters
gives
\[
                \tfrac14 y\log y+C_I y+o_{\psi,\eps}(y).
\]
For the principal remainder use \(G_{\mathbf1}(y)=o(y)\)
from Proposition~\ref{prop:allzeros}. The error is uniform on
this range. The same hypotheses on \(q,d\) used in
Lemma~\ref{lem:primes53} apply. The ramp sum \(S(y)\) is convex
and \(yS'_+(y)-S(y)=\sum_{n\le y}c(n)\). At \(y=5x/4,7x/4\),
bound the derivative by the two secants with displacements
\(\pm\eta y\). If the relative remainder is at most \(e_x=o(1)\),
take \(\eta=\sqrt{e_x}+1/\sqrt{\log x}\). The resulting error
in \(yS'_+(y)-S(y)\) is \(O(y(\eta+e_x/\eta))=o(y)\).
The linear term cancels and the derivative of the logarithmic
term gives \(y/4\) in this expression. On subtraction, and
discarding higher prime powers of total weight
\(O(\sqrt x\log^2(2x))\), this gives \(x/8+o(x)\) before
the condition involving \(T\). The other possible ramified
primes divide the fixed integer \(2\operatorname{cond}(\psi)\).

In the construction of Theorem~\ref{thm:rangenew50}, \(T\)
is the union of \(3o+1\) cosets of \(D A_0\) in \(A\), where
\(o=\operatorname{ord}\psi\) and \(D\) is a power of two
with \(2x^{\eps/8}\le D<4x^{\eps/8}\). Here the groups use
additive notation. Write \(G_l=\Gal(C_l/\QQ)\), a cyclic
2-group of order \(2^l\) divisible by \(D\). Restriction of
\(D A_0\) lies in \(D G_l\). The image of \(T-T\) in
\(G_l/DG_l\) has at most \((3o+1)^2\) elements. Since
conjugation fixes the real field \(C_l\), the image of
\(\pi_p\) there is the image of \(p\). The quotient of order
\(D\) is the Galois group of the maximal real subfield of the
cyclotomic field of conductor \(4D\). Thus \(\pi_pT\cap
T\ne\varnothing\) puts \(p\) in at most \(2(3o+1)^2\)
residue classes modulo \(4D\). Those classes contain
\(O_\psi(1+x/D)\) integers in this interval, of total
logarithmic weight \(o(x)\). This proves the stated asymptotic.
Finally \(1<p<2^{l+2}-1\) for large \(x\), so \(p\) is
neither sign of one modulo the conductor of \(C_l\).
Its nontrivial image in the 2-group \(G_l\) forces even order.
\end{proof}

\begin{theorem}[A common rational divisor of the simultaneous norms]
\leavevmode\label{thm:quot80}
For Proposition~\ref{prop:center78} choose at each definite
step its one basis integral at the primes in \(S_j\cup\mathcal P_x\).
This is allowed by Lemma \ref{lem:basis63}. Write \(s_j\) for
the larger rank. At this step put
\[
    D_j^{\rm per}=\prod_{p\in\mathcal P_x}
        p^{\,2r_j\lceil h/2\rceil d_p(a,s_j)},\qquad
    d_p(a,s)=\max(0,\lceil (a-s(p-1))/p\rceil).
\]
The same integer can be used for every translate \(u\).
For all \(p\mid D_j^{\rm per}\), every coefficient place over
\(p\) gives a nonnegative valuation to
\(\widetilde Y_{u,j}=Y_{u,j}/D_j^{\rm per}\). One has
\begin{align}
 \frac{1}{r_j}\log D_j^{\rm per}
    &\ge (1/14-o_{\psi,\eps}(1))ha,                         \label{eq:div80}\\
 \mathcal H_{S_j}(\widetilde Y_{u,j})
    &=\mathcal H_{S_j}(Y_{u,j})-\log D_j^{\rm per}.           \label{eq:htdiv80}
\end{align}
The centered estimate \eqref{eq:center78} holds with every
\(Y_{u,j}\) replaced by \(\widetilde Y_{u,j}\).
For the unmodified norms of this common-basis construction one has
\[
       \sum_{j=1}^m \mathcal H_{S_j}(Y_{u,j})/r_j-hm M_{\rm sp}
             \ \ge (1/14-o_{\psi,\eps}(1))hma .
\]
\end{theorem}
\begin{proof}
At \(p\in\mathcal P_x\) the permutation \(\pi_p\) fixes the
\(K_q\)-sign of an embedding. Thus \(\pi_p t\) belongs to
\(\Phi_0\setminus T\) for \(t\in T\) and has minor degree
zero. In every cycle choose the half containing at least as
many members of \(T\) as its partner. These choices cover
at least \(\lceil h/2\rceil\) members in total. Use just the
first power in
Corollary \ref{cor:minor77}.
It gives valuation at least
\(r_j\lceil h/2\rceil d_p(a,s_j)\) for each cleared coordinate.
Squaring and summing doubles the exponent at every place over
\(p\). The input is supported at a fixed split prime and the
representatives move lattices only at split primes, so the
local hypotheses hold at all of \(\mathcal P_x\). This proves
the integrality after division, including possible zeros of
individual coordinates. The scalar norms themselves are nonzero
by the torsion existence theorem in Section \ref{nav:V-18}.

The inequality \(d_p(a,s_j)\ge a/p-s_j\), Lemma~\ref{lem:move80},
and \(p\le7x/4\) give
\[
 \sum_{p\in\mathcal P_x}d_p(a,s_j)\log p
       \ge \frac{4a}{7x}(x/8+o(x))-O(s_j x)
       = (1/14-o(1))a.
\]
All ranks have a fixed polynomial bound and \(s_jx=o(a)\)
uniformly. Multiplication by \(2\lceil h/2\rceil\ge h\)
gives \eqref{eq:div80}. The primes in \(S_j\) split in
\(K_q\), hence are disjoint from \(\mathcal P_x\). At the
dividing primes no negative valuation has been introduced.
At every other finite place the valuation is unchanged.
Division by the positive rational integer subtracts exactly
its logarithm from the normalized logarithm of the field norm.
This proves \eqref{eq:htdiv80}. All the divisibilities used for
the two-sign count persist on \(S_j\). The product formula
applied to the quotient therefore gives \eqref{eq:center78}
with the claimed replacements.

Finally \eqref{eq:sparse48} gives \(\Delta_u=o_{\psi,\eps}(hma)\)
unconditionally. Indeed \(Y'=a/W\asymp x\sqrt b\) is in the
large-argument estimate of Lemma \ref{lem:far14}. At \(a/B'\sim x\)
the finitely many primitive powers of the fixed character satisfy
\(G_\rho(y)=o_\psi(y)\) by Proposition~\ref{prop:allzeros},
and the remaining selected twists satisfy Theorem~\ref{thm:twists}.
The coefficients in \eqref{eq:sparse48} are bounded in terms
of \(\psi\). Insert this lower bound on \(\Delta_u\) in the
estimate for the quotient norms and use \eqref{eq:htdiv80}
and \eqref{eq:div80}. This yields the final inequality.
\end{proof}

Here \(x^{5\eps}\mathcal B=o(hma)\) by the ranks in
Theorem~\ref{thm:rangenew50} and the computation in
Proposition~\ref{prop:excess75}. Thus in this construction a
term larger than the permitted error has to be removed before
seeking an upper estimate for the periods. The explicitly divided
norms still need an upper comparison at the precision of
\eqref{eq:gap}. For the cohomology sum the remaining requirement
is the bound on \(\mathcal R_{**}\) of Proposition~\ref{prop:res80}.

\subsection{Principal generators at the two auxiliary primes and in the terminal window}\label{nav:II-26}
Write \(K=K_d=\QQ(\sqrt{-d})\), so \(|D_K|=d\), and keep the
integer \(N=\operatorname{lcm}(F,2^{l'+2})qd\) and the scale
\(U=(kb)^8\) of Proposition~\ref{prop:ray50}. If an algebraic
integer \(\pi\in K\) generates a degree-one prime ideal, the
condition \(\pi\equiv1\bmod M\mathcal O_K\) specifies its residue in
the quadratic field. This is the condition on the generator in
\srcloc{Formalism/SharpRestriction/GlobalSchurBounds/ShortIntervalPrimes}{83}.

\begin{lemma}[Principal ray primes in a relative interval]
\label{lem:ray88}
Let \(K\) be any imaginary quadratic field of absolute
discriminant \(d\) and \(M\ge1\) an integer. Put \(B=4dM^2\).
There are absolute \(c,C>0\) with the following property. For
\(Y\ge \max(16,B^C)\), put
\[
     \rho(Y,B)=\exp(-c\log Y/\log B)+\exp(-c\sqrt{\log Y}).
\]
If \(0<\vartheta\le1\) and \(\vartheta\ge C\rho(Y,B)\),
the number of rational primes in \((Y,(1+\vartheta)Y]\)
of the form \(p=\pi\bar\pi\), \(\pi\in1+M\mathcal O_K\), is at least
\[
              \frac{\vartheta Y}{B^C\log Y}.
\]
\end{lemma}
\begin{proof}
Let \(L\) be the ray class field modulo \(M\mathcal O_K\). Conjugation
preserves the ray subgroup, so \(L/\QQ\) is Galois. Its subgroup
\(\Gal(L/K)\) is abelian, and the conductor norm of each of its
characters is at most \(M^2\). Also \([L:\QQ]\ll dM^2\). Indeed
the ray class exact sequence bounds \([L:K]\) by
\(h_K|(\mathcal O_K/M\mathcal O_K)^\times|\), and the second factor is at
most \(M^2\). Minkowski's bound represents each class by an
integral ideal of norm \(O(\sqrt d)\). The sum of the inverse
squares of all ideal norms is \(\zeta_K(2)\le\zeta(2)^2\), so
there are \(O(d)\) such ideals.

Apply the unconditional Theorem 1.4 of Thorner and Zaman to
\(L/\QQ\), the identity conjugacy class, and this abelian
subgroup.~\cite{I-I-I-tz50} At \(z=Y,(1+\vartheta)Y\) its formula
for the number of completely split unramified rational primes
is
\[
   [L:\QQ]^{-1}A(z)(1+O(\rho(Y,B))).
\]
Here \(A(z)=\operatorname{Li}(z)\) in the absence of an
exceptional zero. In the presence of the possible simple real
zero \(\beta\) it is
\(A(z)=\operatorname{Li}(z)-\operatorname{Li}(z^\beta)\).
The character value on the identity is one. Its defining region gives \(\beta>1/2\).

We record an elementary bound \(1-\beta\gg B^{-C_0}\) for an
absolute \(C_0\) in this quadratic situation. The zeta function
of \(L\) factors over \(K\) into the Hecke functions of the ray
characters, with the trivial factor \(\zeta_K\). A zero of the
trivial factor in this real interval comes from the quadratic
Dirichlet character of \(K\), since the alternating series for
the rational zeta function excludes a real zero between zero
and one. For a nontrivial factor the exceptional character
\(\eta\) is real. Conjugation of its ideals preserves their
norms, so the conjugate character has the same Hecke function.
Simplicity of the zero forces that character to equal \(\eta\).
All the factors are holomorphic at \(\beta\). Thus the quadratic
extension \(K'/K\) cut out by \(\eta\) is normal over \(\QQ\).
Complex conjugation has order two on \(K'\) and acts
nontrivially on \(K\), so \(\Gal(K'/\QQ)\) is the Klein four
group. The conductor--discriminant formula gives
\(|D_{K'}|=d^2\mathrm N\mathfrak f_\eta\le d^2M^2\).
Factoring \(\zeta_{K'}\) over \(\QQ\) shows that \(\beta\)
is a zero of a primitive quadratic Dirichlet character of
conductor \(D\le d^2M^2\). The conductor bound follows also by
the discriminant formula in the tower over its quadratic
subfield. The same bound includes the trivial-factor case.

The quadratic class number formula gives
\(L(1,\chi)\gg D^{-1/2}\). For a real quadratic field the
fundamental unit exceeds an absolute constant greater than one,
by its integral quadratic equation and norm \(\pm1\). For an
imaginary quadratic field the number of roots of unity is at
most six. The sum of a nontrivial character over a complete
period is zero, so all its partial sums have absolute value at
most \(D\). Partial summation gives
\[
 L(\sigma,\chi)=\sigma\int_1^\infty
         \left(\sum_{n\le t}\chi(n)\right)t^{-\sigma-1}\,dt .
\]
Differentiation under the integral shows
\(|L'(\sigma,\chi)|\le6D\) for \(1/2\le\sigma\le1\).
The mean value theorem proves the asserted lower bound for
\(1-\beta\).

To subtract the two counting formulas set \(s(Y)=1\) when
there is no exceptional zero and \(s(Y)=1-Y^{\beta-1}\)
otherwise. Uniformly for \(Y\le t\le2Y\),
\(A'(t)\gg s(Y)/\log Y\) and \(A(t)\ll Ys(Y)/\log Y\).
In the exceptional case differentiation gives
\((1-t^{\beta-1})/\log t\). For the upper bound integrate
\(1/\log v\) from \(t^\beta\) to \(t\) and use
\(\log(t^\beta)\ge\frac12\log Y\) and
\(1-t^{\beta-1}\le2(1-Y^{\beta-1})\). This last inequality
follows from \(\log t\le2\log Y\). Subtraction and
\(\vartheta\ge C\rho(Y,B)\) now give a lower bound
\(\gg\vartheta Ys(Y)/([L:\QQ]\log Y)\).
The bound on \(1-\beta\) gives \(s(Y)\gg B^{-C_0}\) with an
increase of the constant if needed. Each prime counted splits
in \(K\). A prime above it is trivial in the ray class group,
so it has a generator \(\pi\equiv1\bmod M\mathcal O_K\). The ideal
is integral and of norm \(p\), giving
\(\pi\in\mathcal O_K\) and \(\pi\bar\pi=p\). Increasing \(C\) absorbs
the absolute constants.
\end{proof}

For a rational prime \(g\) and a positive integer \(r\), put
\[
 \mathcal B_r(g)=\{g\}\ \cup\
  \{\ell\text{ prime}\mid
       \ell^{v_\ell(g^i-1)}>g^{i/4}
              \text{ for some }1\le i\le r\}.
\]

\begin{theorem}[Compatible split levels on the common fields]
\label{thm:levels88}
At every sufficiently large \(x\) on the common intervals, for
any \(r\ge1\) bounded by a fixed power of \(x\), the two primes
of Proposition~\ref{prop:ray50} can be chosen in \((U,2U]\)
with \(\pi_i\in1+N\mathcal O_K\) and
\(\mathcal B_r(g_1)\cap\mathcal B_r(g_2)=\varnothing\).
Fix such a choice and set \(M=N(g_1g_2)^2\). For any fixed
\(c_0>0\) and every real \(Y\) satisfying
\(\log Y\ge c_0 m x^\eps\), the interval
\[
                  Y<T\le Y\bigl(1+1/(fk)\bigr)
\]
contains \(Y^{1-o_{\psi,\eps,c_0}(1)}\) rational primes
\(T=\pi\bar\pi\) with \(\pi\in1+M\mathcal O_K\). The error in the
exponent tends to zero as \(x\) increases, uniformly in \(Y\)
and the two auxiliary primes. Every such \(T\) is \(1\bmod M\)
and splits completely in \(E\).
\end{theorem}
\begin{proof}
For \(M=N,\ Y=U,\ \vartheta=1\), the estimates in
Proposition~\ref{prop:ray50} give
\(\log(4dN^2)\ll_{\psi,\eps}x^{\eps/5}\log^2(2x)+\log(2x)\)
and \(\log U\asymp_\eps x^\eps\).
Lemma~\ref{lem:ray88} gives \(U^{1-o(1)}\) choices.
For each \(i\) at most three prime powers exactly dividing
\(g_1^i-1\) can exceed \(g_1^{i/4}\), since the product of
four would exceed \(g_1^i\). Thus
\(|\mathcal B_r(g_1)|\le3r+1\).
Fix \(\ell\) in this set. A candidate \(g\ne\ell\) for which
\(\ell\in\mathcal B_r(g)\) satisfies
\(g^i\equiv1\bmod\ell^e\) for some \(i\le r\), where \(e\)
is the smallest positive integer with \(\ell^e>U^{1/4}\).
There are at most \(2i\) such residue classes for each \(i\).
This follows from cyclicity of the unit group for odd prime
powers and its decomposition into at most one cyclic factor and
a factor of order at most two for powers of two. Each class
contains at most \(1+U^{3/4}\) integers in \((U,2U]\). Including
\(g=\ell\), the number of all excluded candidates is therefore
\(O(r^3(1+U^{3/4}))\). This is smaller than \(U^{1-o(1)}\)
since \(r\) is polynomially bounded in \(x\). A remaining
candidate supplies the second prime.

For the terminal modulus, \(\log(4dM^2)\ll_{\psi,\eps}x^\eps\)
uniformly in the two choices. By Theorem~\ref{thm:rangenew50},
\(m\asymp x^{1/2-\eps/16}\), and \(f,k\) are polynomially bounded.
The exponent \(\log Y/\log(4dM^2)\) is \(\gg m\) and
\(\sqrt{\log Y}\gg\sqrt{m x^\eps}\). Each exceeds every fixed
multiple of \(\log x\) eventually. Thus
\(\rho(Y,4dM^2)=o(1/(fk))\) throughout the specified range.
Apply Lemma~\ref{lem:ray88} with \(\vartheta=1/(fk)\). The
logarithms of \(fk\), of the conductor factor \((4dM^2)^C\),
and of \(\log Y\) are all \(o(\log Y)\) uniformly in \(Y\).
This proves the count and allows exclusion of any list of size
bounded by a fixed power of \(x\). Taking norms of
\(\pi\equiv1\bmod M\mathcal O_K\) gives \(T\equiv1\bmod M\), since
\(M\mathcal O_K\cap\ZZ=M\ZZ\). Finally \(E\subset\QQ(\zeta_N)\).
\end{proof}

\begin{proposition}[The effect of a prime-power division on the height]
\label{prop:height88}
Let \(z\ne0\) lie in a number field \(J\), let \(S\) be a set of
rational primes, and let \(T\notin S\) be a rational prime.
Normalize \(v_w(T)=1\) for \(w\mid T\) and set
\(q_w=[J_w:\QQ_T]/[J:\QQ]\). For each integer \(a_0\ge0\),
\[
 \mathcal H_S(z)-\mathcal H_S(z/T^{a_0})
    =\log T\sum_{w\mid T}q_w\min\{a_0,\max(0,v_w(z))\}.
\]
For \(a_0>0\) the decrease equals \(a_0\log T\) exactly when
\(v_w(z)\ge a_0\) at every \(w\mid T\).
\end{proposition}
\begin{proof}
By the product formula the height of Proposition
~\ref{prop:center78} is the sum of the normalized signed
valuations above \(S\) and the positive parts of valuations
elsewhere, in each case with local-degree and \(\log p\) weights.
Only the places above \(T\) change. At such a place, writing
\(t=v_w(z)\), the difference of the positive parts is
\(\max(0,t)-\max(0,t-a_0)=\min(a_0,\max(0,t))\).
All the \(q_w\) are positive and their sum is one. This also
gives the equality condition.
\end{proof}

The source places the terminal level near
\(T_*=(D_*/(d_{u,*}\Gamma_P^2))^{1/(fk)}\) where \(D_*,d_{u,*}\)
are its dimension parameters and \(\Gamma_P>1\) its loss factor.
See \srcloc{Formalism/LocalRingCardinality/DataKernels}{157}
and \srcloc{Formalism/SharpRestriction/GlobalSchurBounds/ShortIntervalPrimes}{329}.
Theorem~\ref{thm:levels88} supplies the primes for every target
in its stated range, independently of how that target is obtained.
Corollary~\ref{cor:real91} below verifies this lower bound for the specified unitary row groups and discrete-series parameters. The negative term involving \(\log T\) in
\srcloc{Formalism/LocalRingCardinality/DataDIVDefs}{83} also
requires a period estimate. Proposition~\ref{prop:height88}
specifies the valuation gain that would be needed to realize it
by a rational division of the auxiliary scalar.

\section{Arithmetic averages on the period interval}\label{nav:g-I-2}

\subsection{Exact descent of strict classes on the period interval}\label{nav:I-50}
\begin{theorem}[Exact preservation at every primary exponent]
\label{thm:exact241}
Fix the entire interval and the nonzero large-prime summands
of Theorem~\ref{thm:detect241}. Choose one of its primes \(s\).
For each coefficient prime \(\ell\) let \(L_\ell\) be the
degree-\(\ell\) subfield of \(\QQ(\zeta_s)\), and let \(L\)
be their compositum. For a datum \((E_x,\ell,r,\alpha)\), put
\[
 F_x=E_x\QQ(i)L,\qquad \Delta_x=\operatorname{Gal}(F_x/\QQ),
 \qquad A=(J/O)(\kappa_\ell^{-r}\alpha^{-1}),
\]
where \(J=\operatorname{Frac}(O)\). Use the strict group
\(S_D\) of Theorem~\ref{thm:descent240} over a field \(D\).
For \(M=A[\ell^n]\), \(n\ge1\), and for \(M=A\), there are
equalities and restriction isomorphisms
\[
    S_\QQ^{\{s\}}(M)=S_\QQ(M)
        \ \xrightarrow{\ \sim\ }\ S_{F_x}(M)^{\Delta_x}.
\]
Here the superscript \(\{s\}\) allows arbitrary ramification
at \(s\) and keeps all other strict conditions. On each
\(F_x\), for every \(m\ge0\), there is an ordinary crystalline
cuspidal adjoint with cyclotomic exponents \(u,0,-u\) for
\(u=r+m(\ell-1)\), with the original residual characters,
and unramified away from \(\ell\). All the original
splitting conditions of Corollary~\ref{cor:common240} hold in
this one \(L\). In particular the invariant groups recover the
original strict character summands, including every primary
exponent, throughout the given support.
\end{theorem}
\begin{proof}
First take finite coefficients. Via the trace pairing of the
unramified ring \(O\) over \(\ZZ_\ell\), the Cartier dual of
\(M=A[\ell^n]\) is \(D_n=T/\ell^nT\), with \(T\) as in
Lemma~\ref{lem:dual241}. Let \(d_n\) be the reduction of its
integral class. Both modules are unramified at \(s\). The
unramified local cohomology of \(D_n\) is
\[
 H^1_{\rm ur}(\QQ_s,D_n)
   =D_n/(\operatorname{Fr}_s-1)D_n .
\]
It is a cyclic \(O\)-module. Frobenius acts trivially on the
residual character at \(s\), and the reduction of \(d_n\)
evaluates there to a nonzero element of the residue field by
Theorem~\ref{thm:detect241}. Hence the single class \(d_n|_s\)
generates this module over \(O\) for every \(n\).

At \(s\) the unramified subgroups for \(M\) and \(D_n\)
annihilate one another under the perfect local Tate pairing.
Indeed their cup products inflate from the residue Galois
group, of \(\ell\)-cohomological dimension one. For either
unramified finite module, its unramified subgroup has the
order of its Frobenius invariants, by the displayed formula.
Local duality and the local Euler characteristic at \(s\ne\ell\)
give \(|H^1(\QQ_s,M)|=|H^0(\QQ_s,M)|\,|H^0(\QQ_s,D_n)|\).
These orders give exact annihilators.

For \(c\in S_\QQ^{\{s\}}(M)\) take the global cup product of
\(c\) with any \(O\)-multiple of \(d_n\). At \(\ell\) the
restriction of \(c\) is zero. At each other finite place
apart from \(s\), both classes are unramified, so the local
product vanishes by inflation from the residue group, also
when the coefficient action is ramified. At infinity the
positive-degree cohomology vanishes since \(\ell\) is odd.
Global reciprocity says the local invariants sum to zero.
Thus \(c|_s\) is orthogonal to all of \(H^1_{\rm ur}(\QQ_s,D_n)\).
It is unramified at \(s\) by the preceding paragraph. This
proves the first equality for every \(n\).
\href{https://www.jmilne.org/math/Books/adt.html}{Milne,
\emph{Arithmetic Duality Theorems}, Part I, Sections 2 and 4}.

For \(A\), a continuous cocycle has finite image and takes
values in \(A[\ell^n]\) for some \(n\). Strictness of its
class lifts to these finite coefficients at each required
place. Indeed \(A^{G_{\QQ_\ell}}=0\) by the nontrivial
character on residual cyclotomic inertia. At each \(p\ne\ell\),
the group \(A^{I_p}\) is divisible since the averaging projector
for the finite prime-to-\(\ell\) inertia action splits it off
from \(A\). The multiplication-by-\(\ell^n\) sequence thus
gives injective maps on the local \(H^1\) at \(\ell\) and
on \(H^1(I_p,-)\) elsewhere. The finite-coefficient equality
now applies.

For each \(\ell\) choose a character of order \(\ell\)
cutting out \(L_\ell\). The one-prime form in
Proposition~\ref{prop:level241} has this auxiliary character.
In the proof of Theorem~\ref{thm:adjoint240}, the same
primitive conductor argument gives inertia characters
\(\chi\xi,1,(\chi\xi)^{-1}\) and zero monodromy away from
\(\ell\). The exponents of the ordinary crystalline flag
are \(u,0,-u\). The residual cyclotomic character to exponent
\(u\) equals the one to exponent \(r\). Its inertia order
therefore still excludes a quadratic self-twist.
The adjoint stays cuspidal
by solvable base change. Restriction to \(F_x\) kills
\(\eta,\alpha,\xi\), giving the claimed unramifiedness.
Each \(L_\ell\) is totally real because \(\ell\) is odd.
The power conditions of Theorem~\ref{thm:detect241} give
complete splitting in \(L\) of every prime on the specified
list, including the coefficient primes.

Finally \(E_x\QQ(i)\) is unramified at \(\ell\) and has
degree prime to \(\ell\). The inertia above \(\ell\) still
contains an element acting on \(M\) through a scalar whose
difference from one is a unit. Hence \(M^{G_{F_x}}=0\),
and inflation--restriction identifies the global \(H^1\)
over \(\QQ\) with the invariants over \(F_x\). At \(\ell\)
the local degree is prime to \(\ell\). Away from \(\ell,s\)
the inertia index is prime to \(\ell\) because \(L\) is
unramified there. Restriction at these places is injective
by corestriction. Thus the antecedent of an invariant strict
class belongs to \(S_\QQ^{\{s\}}(M)=S_\QQ(M)\).
Restriction of a strict class is strict, proving surjectivity
and the isomorphism. Prime-to-degree character descent in
Theorem~\ref{thm:char} identifies the group for \(A\) with
the corresponding original summand over \(E_x\).
Unramified enlargement of the coefficient DVR follows by
finite free scalar extension, which commutes with cochains,
restriction kernels and invariants.
\end{proof}

In particular the contribution of a supported summand to
\eqref{eq:region198} can be computed on \(F_x\) after taking
invariants and multiplying by the prescribed power of \(\ell\),
with exactly the same logarithmic order. An upper bound for the
sum of those orders, the global theta-period upper bound, and
the lower bound for the designated integral pairing on \(E_x\)
are still required. The auxiliary prime depends on the support
of the strict groups. The construction gives no size bound
for it in terms of the interval parameter.

\subsection{Three primes throughout the dyadic tower}\label{nav:I-59}
\begin{theorem}[Vanishing at \(37,59,67\) in every twist]\leavevmode
\label{thm:towerzero267}
Let \(p\in\{37,59,67\}\), let \(D\) be any subfield of
\(\QQ(\zeta_{2^j})\) for any \(j\ge1\), and let \(r\ge1\).
The \((-1)^r\)-eigenspace of complex conjugation in
\(\mathcal S_{p,r}(D)\) is zero. In particular \(K_{4j}(\ZZ)\{p\}=0\)
for each \(j\ge1\) and each of the three primes.

On each original period field \(E_x\), the complementary
strict classes invariant over any such \(D\subset E_x\) vanish
whenever \(p\nmid[E_x:D]\). For a character \(\psi\) of two-power
conductor this applies to the entire invariant summand over
\(K_\psi C_J\), for every twist on every common interval.
\end{theorem}
\begin{proof}
First take the full cyclotomic field and split the integral
group ring over an unramified extension of \(\ZZ_p\).
Character descent in Theorem~\ref{thm:char} reduces the assertion
to \(S(A)=0\) for each primitive character \(\alpha\) of
two-power conductor \(F\) with \(\alpha(-1)=(-1)^r\).
The strict group uses \(A=(J/O)(\alpha^{-1}\kappa_p^{-r})\).
Averaging is integral since the field degree is a power of two.

We first remove all high conductors independently of \(r\).
The odd character \(\omega^{r+1}\alpha\), with \(\omega\) the
Teichm\"uller character, has odd conductor part \(f_o\mid p\).
Its odd-conductor factor has order dividing \(p-1\). In
Proposition~\ref{prop:effect146} take \(Q=37\) for \(p=37\)
and \(Q=p^2\) for \(p=59,67\). Then \(e=v_2(Q-1)\) is
\(2,3,3\), respectively. The maximal \(p\)-subfield is \(\QQ\).
All conductors in the following table are powers of two. Above
the displayed \(F_0\) one has \(F\ge p4^e\) and
\(F/2^e>p+1\). Hence \(F\ge f_o4^e\), and the only possible
extra odd prime \(p\) is neither sign of one modulo \(F/2^e\).
All the hypotheses of that proposition hold, so the strict
group for each such character vanishes.

For \(F\le F_0\) use Proposition~\ref{prop:moment267} at
precision one. The following are the auxiliary primes \(s\)
and generators \(g\) used. A single pair suffices in the last
row.
\[
\begin{array}{c|r|rr|rr}
p&F_0&s_1&g_1&s_2&g_2\\ \hline
37&512&37889&3&151553&3\\
59&2048&120833&3&724993&5\\
67&4096&1097729&3& & 
\end{array}
\]
Each \(s\equiv1\bmod p F_0\). Primality can be checked by trial
division through \(\lfloor\sqrt{s}\rfloor\). The primitive-root
check is \(g^{(s-1)/q}\not\equiv1\bmod s\) for the prime
divisors \(q\) of \(s-1\). In order down the displayed pairs
these sets are \(\{2,37\},\{2,37\},\{2,59\},\{2,3,59\},\{2,67\}\).

Here are details of the remaining finite calculation. Write
\(u\) for the representative of \(r\bmod(p-1)\) in
\([0,p-2]\). At \(F=1\) retain even \(u\), and at \(F=4\)
retain odd \(u\). At \(F=2^j\ge8\) retain all \(u\).
For such an \(F\) put \(d=F/4\) and write each odd residue
as \((-1)^b5^v\bmod F\) with \(0\le v<d\).
Indeed \(v_2(5^{2^a}-1)=a+2\), so \(5\) generates the
subgroup congruent to one modulo four. A character with the
given parity is primitive precisely when its value \(Y\) at
five has order \(d\), since \(5^{d/2}=1+F/2\bmod F\).
Thus all admissible \(Y\)'s are the roots of
\(H_F(X)=X^{d/2}+1\) over \(\overline{\mathbb F}_p\).
For \(F=1,4\) use a dummy variable with \(H_F(X)=X\),
\(v=0\), and \(b=0\) at \(F=1\) or \((-1)^b=a\bmod4\)
at \(F=4\).

For a listed pair \((s,g)\) set \(w=g^{(s-1)/(Fp)}\) in
\(\mathbb F_s\). The polynomial
\[
 P_{s,F,u}(X)=
     \sum_{\substack{1\le a<Fp\\(a,Fp)=1}}
     (-1)^{bu} a^u\,\operatorname{ind}_g(1-w^a)\,X^v
                  \quad\text{in }\mathbb F_p[X]
\]
uses \((b,v)\) for \(a\bmod F\). Its value at the residual
character is exactly \(D_{s,1}(\alpha,r)\). The table below
lists every nonconstant monic gcd \((H_F,P_{s_1,F,u})\).
Each has the single root \(Y\) displayed, and the last column is
the value at that root for the second pair.
\[
\begin{array}{r|r|r|r|r}
p&F&u&Y&P_{s_2,F,u}(Y)\\ \hline
37&1&20&0&4\\
37&8&6&36&18\\
37&16&1&6&13\\
37&16&5&6&28\\
37&16&25&31&2\\
37&16&32&6&7\\
59&8&6&58&57
\end{array}
\]
For reproducibility the complete calculation of the gcd column
is given by the following loops. All polynomial operations and
accumulations are modulo \(p\). The powering used to find the
index is modulo \(s\).
\begin{quote}\small
For each \(F=1,4,8,\ldots,F_0\), initialize one zero
polynomial \(P_u\) for each allowed \(u\). Form
\(L[(w^F)^j]=j\) for \(0\le j<p\). For each
\(1\le a<Fp\) prime to \(Fp\), read \(b,v\) from the list of
signed powers of five, with the two small cases specified above.
Put \(z=L[(1-w^a)^{(s-1)/p}]\). For each allowed \(u\), add
\(z((-1)^ba)^u X^v\) to \(P_u\). For each resulting polynomial
start with \(B=H_F,\ C=P_u\bmod H_F\). Replace the pair by
\(C,B\bmod C\) while \(C\ne0\). Divide \(B\) by its leading
coefficient. Record the resulting polynomial when its
degree is positive.
\end{quote}
In the displayed cases this gives \(X-Y\). At every root of
\(H_F\) at least one of the available moments is therefore
nonzero. The character values have order prime to \(p\), so
their reductions preserve their orders. The last assertion
of Proposition~\ref{prop:moment267} proves \(S(A)=0\) for
every positive twist of this parity. In particular this
includes the rational twists \(r\equiv1048\bmod1332\) at
\(37\) for which the earlier Bernoulli test had an exception.

Restriction to the full dyadic field has prime-to-\(p\)
degree for every subfield \(D\) in the theorem, so it is
injective on the strict group, by corestriction.
The assertion for \(\ZZ\) follows from
Corollary~\ref{cor:transfer100}, or directly from
Theorem~\ref{thm:wild98} since global and local invariants
over \(\QQ\) agree. For \(D\subset E_x\) the usual
prime-to-degree descent gives exactly the invariant group
over \(E_x\). Indeed averaging makes restriction on global
cohomology an isomorphism onto the invariants and detects
each local condition by restriction. Finally \(K_\psi C_J\)
is dyadic when \(\psi\) has two-power conductor, and adjoining
the two auxiliary square roots has relative degree dividing four.
\end{proof}

\subsection{The full cohomology estimate from pairs of residue primes}\label{nav:I-60}
Fix an original interval as in Theorem~\ref{thm:conductor232}.
There is an a priori integer \(B=\lceil\exp(V^{A_0})\rceil\),
with \(A_0\) depending only on \(\psi,\varepsilon\), that may
be used in the following construction. List \emph{every} odd
prime \(p\le\exp B\) and take \(h_p=B\). Choose \(K\) as in
Proposition~\ref{prop:residue269} with \(N_*\) and these
precisions. The cyclotomic choice alone contains all the
original period fields. Let \(s_1,\ldots,s_a\) be sampled
independently among the primes completely split in \(K\).
Probabilities and averages of functions of this tuple mean
limits over the \(a\)-fold product of the split primes up to
a common bound tending to infinity. Repeated sampling of the
same rational prime has limiting density zero when \(a\ge2\).

For \(x\) on the interval let \(I_x\) consist of all triples
\((p,r,[\alpha])\) in the preceding list with \(r<t(x)\),
where \(\alpha\) runs over the character orbits of \(E_x\)
of sign \((-1)^r\) and \(r\ge1\). Set
\[
 U_{x,a}=\sum_{j\in I_x}b_j\log p_j\,
                       \min_{1\le i\le a}W_j(s_i).
\]
Each summand for \(a=2\) is the logarithm of the norm of
the ideal \((D_j(s_1),D_j(s_2),p_j^B)\) in the local integer
ring, with the norm understood as its additive index.

\begin{theorem}[A common residue calculation for the full strict sum]\leavevmode
\label{thm:average269}
For \(C_x=C(E_x,t(x);N_*)\) from Theorem~\ref{thm:lattices269},
\[
 U_{x,2}\ge C_x,\qquad
 \mathbb E\,\sup_{x\in[V,V^{3/2}]}(U_{x,2}-C_x)
             \ll n_*t_* .
\]
The implied constant is absolute, independently of the prime
cutoff and all precisions used above. With the full cohomology
sum \(\mathcal L(x)\) of \eqref{eq:length12},
\[
 \mathbb E\,\sup_{V\le x\le V^{3/2}}
  \frac{|U_{x,2}-\mathcal L(x)|}{\exp(x^\varepsilon/2)}
                 \ \longrightarrow\ 0 .
\]
For \(a\ge2\), the relative density that
\(U_{x,a}=C_x\) simultaneously for every \(x\) is at least
\(1-4n_*t_*3^{-a}\). In particular \(O_{\psi,\varepsilon,A}
(\log(2V))\) samples suffice for density at least \(1-V^{-A}\)
for any fixed \(A>0\). The prime list and its precision can
be fixed before any of the strict groups are known.
\end{theorem}
\begin{proof}
Theorem~\ref{thm:cells103} bounds the sum of the logarithmic
strict orders on \(E_*\) at twist \(r\) by
\(\exp(C[n_*(r+2)]^5\log(2N_*n_*(r+2)))\).
The lattice comparison in Theorem~\ref{thm:lattices269}
was proved for any twist cutoff on this field. Apply it with
\(t_*\). The preceding bound summed through \(t_*\), together
with its polynomial comparison error, is less than \(B\)
if \(A_0\) is sufficiently large. Indeed \(n_*,t_*,N_*\)
are polynomially bounded in \(V\). Every nonzero index in
the resulting sum contributes \(b_jd_j\log p_j\). Therefore
none requires \(p>\exp B\), and \(e_\alpha B\ge d_j\) at
every listed prime, since \(p\ge3\). The characters on the
smaller fields inflate to this list with the identical open
set, primitive conductor and lattice. Hence these choices
cover the entire interval.

For each triple the excess of the minimum over \(d_j\) is
nonnegative. By Proposition~\ref{prop:residue269} the chance
that it is at least \(v>0\) equals \(p_j^{-a b_j v}\)
when \(v\le e_\alpha B-d_j\), and zero otherwise. This uses
independence between the sampled primes only. All events
in question come from finite Galois quotients, so the
individual densities give these product limits. Summing over
positive \(v\), the expected excess is at most
\(1/(p_j^{a b_j}-1)\). For any one \(p\) the sum of the
\(b_j\) over the largest field and all its twists is at
most \(n_*t_*\). Indeed the sum of \(e_\alpha b_\alpha\)
over all character orbits in the regular representation is
\(n_*\). Dominate the supremum by the sum of the nonnegative
excesses on this entire list. For \(a=2\) the resulting bound
is at most
\[
                 n_*t_*\sum_{m\ge3}\frac{\log m}{m^2-1}
                 \ll n_*t_* .
\]
This summable bound allows the prime list to include even
those primes whose every index is zero.

The deterministic comparison error with \({\cal A}(E_x,t(x))\)
is polynomial in \(V\) by Theorem~\ref{thm:lattices269}.
After division by \(\exp(x^\varepsilon/2)\) this polynomial
and the expected full excess tend uniformly to zero.
Theorem~\ref{thm:res100} controls the remaining difference
between \({\cal A}\) and \(\mathcal L\) on the same intervals.
This proves the second assertion. Finally the union bound
for any individual excess to be positive is at most
\(n_*t_*\sum_{m\ge3}m^{-a}\). Integral comparison bounds this
sum by \(3^{-a}(1+3/(a-1))\le4\cdot3^{-a}\), proving the
density statement.
\end{proof}

Thus a bound for the stated mean, with a fixed margin below
\(\exp(x^\varepsilon/2)\) uniformly on the interval, would
supply the cohomology input in \eqref{eq:length12}. The
error in passing to the mean is already below this scale.
An upper estimate for the mean itself remains open.

\subsection{An exact mean over primes split in the common conductor}\label{nav:I-61}
Use the original interval \([V,V^{3/2}]\), its largest modulus
\(N_*\), degree \(n_*\) and twist cutoff \(t_*\).
For primes \(s_1,s_2\equiv1\bmod N_*\) put
\(g_k=\gcd(s_1^k-1,s_2^k-1)\). For \(M\ge1\) define
\[
  A_{x,M}=\sum_{1\le r<t(x)}
   \sum_{\substack{p\mid g_{r+1},\,p\le M\\p\ {\rm odd}}}
    \ \sum_{[\alpha]} b_\alpha\log p
    \sum_{v=1}^{e_\alpha\min(M,v_p(g_{r+1}))} H_v .
\]
Take one \(p\)-adic orbit of characters of \(E_x\) of sign
\((-1)^r\) in each inner character sum. Use its primitive
conductor and the weights in Proposition~\ref{prop:law273}.
Write \(A^+_{x,M}\) for the sum with \(P_v\) in place of \(H_v\).
For a fixed cutoff the mean is the limit over all ordered pairs
of the allowed primes up to \(Y\), divided by the square of their
number, as \(Y\) tends to infinity.

\begin{theorem}[A two-prime formula with support on ordinary integer divisors]\leavevmode
\label{thm:mean273}
For each fixed \(M\) the exact mean is
\[
       \mathbb E A_{x,M}=
       \sum_{1\le r<t(x)}\ \sum_{\substack{p\le M\\p\ {\rm odd}}}
       \ \sum_{[\alpha]} b_\alpha\log p\,\min(d_{p,r,\alpha}(N_*),
                                             e_\alpha M).
\]
The mean of \(A^+_{x,M}\) exceeds or equals this expression and
differs from it by at most \(C n_*t_*\) for an absolute \(C\).
As \(M\) increases the exact mean eventually equals
\(C_x=C(E_x,t(x);N_*)\) simultaneously on the full interval.
With \(M\) past that threshold, the supremum over the interval of
\(\big|\mathbb E A^+_{x,M}-\mathcal L(x)\big|/\exp(x^\varepsilon/2)\)
tends to zero as \(V\) tends to infinity. Each term for a given
twist uses only the two fields of degrees \((r+1)N_*\) over the
sampled residue fields or their subfields.
\end{theorem}
\begin{proof}
For fixed \(M\) one can sum over all odd \(p\le M\) and all
\(v\le e_\alpha M\), putting ineligible terms equal to zero.
Indeed the eligibility requirement at both primes is
\(p^{\lceil v/e_\alpha\rceil}\mid g_{r+1}\). This is exactly
the divisor and exponent cutoff in the definition. All
summands are functions on finite Galois data. For this limit
one may omit the finitely many sampled primes dividing the
moduli of these data, with density zero. Independence between
the two prime coordinates holds in the limiting product even
when the events for different characters on one prime are
dependent. Linearity of the finite sum and
Proposition~\ref{prop:law273} now give the exact formula.

For the positive sum the excess at \((p,r,[\alpha])\),
before its weight, is at most \(\sum_{j\ge1}p^{-2b_\alpha j}\).
For any one \(p,r\) the sum of \(b_\alpha\) is at most \(n_*\).
Thus the whole error is bounded by
\(n_*t_*\sum_{j\ge3}(\log j)/(j^2-1)\), as claimed.
Theorem~\ref{thm:average269} gives an a priori \(B\) for the
entire interval such that \(p\le\exp B\) contains every nonzero
index and \(d_{p,r,\alpha}(N_*)\le e_\alpha B\).
For example \(M=\lceil\exp B\rceil\) covers both bounds.
This proves stabilization of the means. Their deterministic
comparison to \(\mathcal L\) follows from
Theorems~\ref{thm:lattices269} and~\ref{thm:res100}, since
\(n_*t_*\) and the character-lattice comparison error grow
at most polynomially in \(V\).
Proposition~\ref{prop:orbits273} proves the finite-field
degree assertion. In particular \(N_*\) in the sampled prime
progression is independent of the precision and prime cutoff.
\end{proof}

The limits here take \(Y\) to infinity for fixed \(M\).
The weights \(m_h^2\) can be large when their eligibility
conditions are rare. The uniform upper estimate needed in
\eqref{eq:length12} would require control of the magnitude
of this mean on the common intervals.

\subsection{The centered estimate on the same fields}\label{nav:II-24}
We now record the part of the desired inequality that follows
for the simultaneous high orders of Theorem~\ref{thm:rangenew50}.
In the \(K_d\)-orientation write \(\Phi_1\) for the half of the
embeddings which conjugate \(K_q\). Give the translated \(T_2\)
in this half the weight \(B'\), and all other positions weight
\(W\). Put \(\omega=h_2/(2n)\), and again take \(a=\lceil xb\rceil\).
Write \(D_{\xi_{qd}}\) for the logarithmic derivative at one of
the primitive product of the two quadratic characters.

\begin{proposition}[The surviving arithmetic norm and its required upper bound]
\label{prop:center78}
For the simultaneous minor list of Section \ref{nav:VI-01}, let
\(Y_{u,j}\) denote its nonzero fixed-representative norm at step
\(j\), \(r_j\) the smaller rank, and \(S_j\) the selected primes
inert in \(K_d\) and split in \(K_q\). The choices of torsion
input at distinct steps are independent. Define the real number
\[
 M_{\rm sp}=\frac a2\{\log(a/W)-1-\gamma+
               D_{\xi_d}-D_{\xi_q}+D_{\xi_{qd}}\}
                       -a\omega\log(B'/W).
\]
Uniformly in \(u\) on the common intervals,
\begin{equation}
 \Delta_u\ \le\ \sum_{j=1}^m
            \frac{\mathcal H_{S_j}(Y_{u,j})}{r_j}
                 -hm M_{\rm sp}+O_{\psi,\eps}(hmb/x).
 \label{eq:center78}
\end{equation}
Here \(\mathcal H_S\) is the normalized logarithm of the field
norm with the negative valuations outside \(S\) added back, as
in Section \ref{nav:VI-01}. On these very fields the large-prime
strict groups have the integral cotangent description of
Theorem~\ref{thm:joint75}.
\end{proposition}
\begin{proof}
The mean of the logarithms of the thresholds on \(\Phi\) is
\(\log W+\omega\log(B'/W)\). Insert this value into Proposition
\ref{prop:count}, with \(\xi_d\) as its quadratic sign. From
the result subtract the explicit \(M_{\rm both}\) of Proposition \ref{prop:others77}. The difference is exactly \(M_{\rm sp}\)
displayed above. Theorem \ref{thm:signal77}, bounds the sum over
all \(t\) of the other sign contribution by the one norm at step
\(j\), up to \(O_{\psi,\eps}(hb/x)\). Proposition~\ref{prop:others77} also
estimates the error for the omitted sign by \(O_{\psi,\eps}(b/x)\)
per \(t\). Subtracting the two main terms and summing over the
\(m\) steps proves \eqref{eq:center78}. The large-prime assertion
uses the same \(x\) and \(E_x\) in Theorem~\ref{thm:joint75}.
\end{proof}

For instance an upper bound by \(hm M_{\rm sp}+
\tfrac12 x^{5\eps}\mathcal B\) for the sum of norms on the right
of \eqref{eq:center78} would suffice for \eqref{eq:gap} on a
sufficiently distant interval. Indeed \(hmb/x=o(x^{5\eps}\mathcal B)\)
by the positive term \(h a_E\) for the chosen ranks and \(a_E=xb\).
The low-order factorial cost in Proposition \ref{prop:others77}, is
already \(o(b\sqrt x)\) per high position and step. The upper
estimate on the full arithmetic norms at this precision, or on
the normalized iterated Bessel periods giving \eqref{eq:gap},
and the uniform estimate for \eqref{eq:minors33} remain open here.

\section{The finite data on the common fields}\label{nav:g-I-3}

\subsection{Prime representatives with prescribed archimedean sizes}\label{nav:III-34}
We need a generator of a ray-principal ideal whose absolute values
are almost equal at infinity. The degree and the modulus both vary.
We give the estimates with absolute constants. We use the analytic
continuation, functional equation and order-one factorization of
primitive Hecke \(L\)-functions, including the pole of the Dedekind
zeta function. These are the usual Hecke \(L\)-function theorems.
For the relation with counting prime elements in archimedean regions
see \href{https://arxiv.org/abs/2209.11816}{W. Kai, Sections 2--3}.

Let \(K\) be totally real or CM, of degree \(n\). Use one embedding
from each complex pair, or all real embeddings in the totally real
case. Denote their number by \(f\), so \(n=ef\) with \(e=2\) or \(1\).
In \(\RR^f\) give the hyperplane \(H\) of coordinate sum zero its
Euclidean metric. Write \(\lambda(a)\) for the orthogonal projection
of \((\log|\sigma a|)_\sigma\) onto \(H\). Let \(\mathfrak m\) be a
nonzero integral ideal and use the narrow ray group in the totally
real case. Put \(d_K=|\operatorname{disc}K|\).

\begin{lemma}[Cancellation for a nonzero archimedean frequency]
\label{lem:angular234}
Let \(B\) be sufficiently large, \(n+\log(3d_K N\mathfrak m)\le B\),
\(X=\exp(B^{20})\), and \(e^{-B}\le \delta\le1/10\). Fix a smooth
nonnegative function \(g\) supported in \((0,1)\), equal to one
on \([1/3,2/3]\), and set \(w(y)=g(\log y/\delta)\). Consider a
unitary ideal character \(\chi\) of conductor dividing \(\mathfrak m\)
whose value on ray-principal ideals is
\(\exp(2\pi\mathrm i\langle v,\lambda(a)\rangle)\), where
\(v\in H\), \(0<\|v\|\le e^{2B}\). Extend it by zero at ideals
not prime to \(\mathfrak m\). Then, for an absolute \(c>0\),
\[
 \left|\sum_{\mathfrak a}\Lambda_K(\mathfrak a)\chi(\mathfrak a)
                         w(N\mathfrak a/X)\right|
          \ll_g X\exp(-cB^8).
\]
Here \(\Lambda_K(\mathfrak p^j)=\log N\mathfrak p\) for \(j\ge1\),
and it is zero on all other ideals.
\end{lemma}
\begin{proof}
We detail the uniformity. Write \(t\) for the imaginary part of \(s\)
and put
\[
  L(t)=\log(3d_K N\mathfrak m)+
             n\log(10+|t|+20\|v\|).
\]
For the primitive character inducing \(\chi\) the gamma factors are
real or complex gamma factors with fixed parity shifts and imaginary
shifts of size at most \(2\pi\|v\|\). This follows by writing its
infinite character in the individual absolute-value coordinates.
The parameters of its square satisfy the same bound with \(4\pi\).
Neither character is a power of the absolute norm. Indeed their
archimedean exponents have sum zero and are nonzero. Equality on
ray elements detects these exponents by weak approximation at the
infinite places with the finite ray conditions imposed.

For \(1<\sigma\le2\), logarithmic differentiation of the completed
product gives, for the primitive \(L\)-function,
\[
 -\operatorname{Re}\frac{L'}{L}(\sigma+\mathrm it,\chi)
 \le C L(t)-\sum_\rho
       \frac{\sigma-\operatorname{Re}\rho}{|\sigma+\mathrm it-\rho|^2}.
\]
The sum runs over the zeros of the completed function. In its
Hadamard product the real constant cancels by the symmetry
\(s\mapsto1-\bar s\), as is seen on the line of symmetry away
from zeros. Stirling's formula bounds the gamma logarithmic
derivatives by \(CL(t)\) in this strip. For \(\zeta_K(\sigma)\)
the pole adds \(1/(\sigma-1)\) to the upper bound. Omitting Euler
factors at the modulus costs \(O(\log(2N\mathfrak m))\), by summing
\(\log N\mathfrak p/(N\mathfrak p^\sigma-1)\).

The inequality \(3+4\cos u+\cos(2u)\ge0\), applied term by term
to the Euler logarithmic derivatives of the trivial character,
\(\chi\) and \(\chi^2\) at respective heights \(0,t,2t\), now
gives, at a zero \(\beta+\mathrm it\) of the primitive function,
\[
         \frac4{\sigma-\beta}\le\frac3{\sigma-1}+C_1 L(t).
\]
Take \(\sigma=1+a/L(t)\), with \(a>0\) a sufficiently small
absolute constant. Rearranging gives \(\beta\le1-c_1/L(t)\).
Also at most \(C_2L(t)\) zeros have heights in \([t-1,t+1]\).
To see this use the completed logarithmic derivative at \(2+\mathrm
it\), where the Euler logarithmic derivative is \(O(n)\) by absolute
convergence, and each of these zeros contributes at least \(1/5\).

Set \(T_0=\exp(B^8)\) and \(L_*=\max_{|t|\le T_0+2}L(t)\ll B^9\).
Throughout \(|t|\le T_0,\ 1-c_1/(2L_*)\le\sigma\le2\) we have
\(|L'/L(\sigma+\mathrm it,\chi)|\ll L_*^2\).
Indeed subtract the logarithmic derivative at \(2+\mathrm it\)
from the product at the point in question. Each zero within one
unit of height contributes \(O(L_*)\) by the zero-free region.
The distant terms are bounded in total by a constant times
\(\sum_{j\ge1}(L(t)+n\log(2+j))/j^2\), using the count in unit
intervals on both sides. The gamma functions have positive real
arguments before the imaginary shifts and contribute \(O(L_*)\).
The same omitted local factors are bounded here by \(O(\log(2N\mathfrak m))\),
since the real part is at least \(1/2\).

Mellin inversion uses the transform
\(\widehat w(s)=\delta\int_0^1g(u)e^{s\delta u}\,du\).
Four integrations by parts give
\(|\widehat w(\sigma+\mathrm it)|\ll_g \delta(1+\delta |t|)^{-4}\)
for \(1/2\le\sigma\le2\). Start at \(1+1/\log X\), where absolute
convergence bounds the logarithmic derivative by \(O(n\log X)\)
using the rational zeta function. Truncate at \(\pm T_0\) and
move to \(1-c_1/(2L_*)\). The integral on this side is
\(O_g(X L_*^2 e^{-c_1\log X/(2L_*)})\), since the transform has
bounded integral on each vertical line. The two horizontal sides
and truncated tails together are
\(O_g(X(n\log X+L_*^2)\delta^{-3}T_0^{-3})\). These estimates
give the assertion.
\end{proof}

\begin{theorem}[A ray prime with prescribed logarithms of its generator]
\label{thm:sector234}
Let \(A\) be any fractional ideal prime to \(\mathfrak m\), let
\(y\in H\), and let \(0<\delta\le1/10\). Suppose the log vectors
of the ray units contain \(f-1\) linearly independent vectors of
norm at most \(e^B\). Take \(B\) above an absolute constant with
\[
     n+\log(3d_K N\mathfrak m)+\log(1/\delta)\le B,
       \qquad X=\exp(B^{20}).
\]
There is a prime ideal \(P\) of absolute residue degree one and an
\(s\in K^\times\), positive at all real places, such that
\[
  X<NP<e^\delta X,\qquad P=(s)A,\qquad
  s\equiv1\bmod\mathfrak m,\qquad \|\lambda(s)-y\|<\delta .
\]
In particular each \(\log|\sigma s|\) differs from
\(n^{-1}\log(X/NA)+y_\sigma\) by at most \(2\delta\).
\end{theorem}
\begin{proof}
Let \(h\) be the ray class number and \(\Lambda\) the lattice of
logs of the ray units in \(H\). One has \(\log h\ll B^2\).
Minkowski's bound represents each ordinary ideal class by an
integral ideal of norm at most \(2^n\sqrt{d_K}\). The number of
ideals of norm at most \(M\) is at most \(M^2\zeta(2)^n\) by the
Euler product. Passing to narrow ray classes adds at most the
factor \(2^n N\mathfrak m\).

We first count primes in the finite ray class of \(A\). Uniformly
for \(e^{\delta/3}X\le t\le e^{2\delta/3}X\), the abelian case of
\href{https://arxiv.org/pdf/1803.02823}{Thorner and Zaman, Theorem~1.4},
counts primes through \(t\) by
\(h^{-1}(\operatorname{Li}(t)-\theta\operatorname{Li}(t^\beta))\)
with relative error \(O(e^{-c B^9})\). Here the exceptional term
can be absent, and otherwise \(\theta=1\) or \(-1\). Their base
field is \(K\), and their largest finite conductor norm is at most
\(N\mathfrak m\). Thus the logarithm entering their range is
\(O(B\log(2B))\). The two exponential errors of their theorem
give the asserted weaker rate.

If an exceptional zero occurs, \(1-\beta\ge e^{-C B^2}\).
Indeed its real character cuts out \(K\) itself or a quadratic
extension \(K'/K\). Its zero is a zero of \(\zeta_{K'}\), with
\([K':\QQ]\le2n\) and \(d_{K'}\le d_K^2 N\mathfrak m\) by the
conductor--discriminant formula. Apply Stark's effective bound
as stated in \href{https://arxiv.org/pdf/1803.02823}{Thorner
and Zaman, (1.3)--(1.4)}, to this field. A zero outside their
exceptional interval already satisfies the displayed weaker bound.
For \(t\) in the indicated interval the derivative of
\(\operatorname{Li}(t)-\theta\operatorname{Li}(t^\beta)\)
is \((1-\theta t^{\beta-1})/\log t\). Its numerator is at least
\(e^{-C' B^2}\), also for the unfavorable sign, by
\(1-e^{-z}\ge\min(z,1)/2\) for \(z\ge0\).
The difference of the main terms at the endpoints therefore
dominates the two errors, since \(\delta\ge e^{-B}\).
The sum of \((\log NP)w(NP/X)\) over primes in the specified
class is at least \(X e^{-C''B^2}\). These counts may exclude
ramified primes by definition. Omitting any additional primes
over \(\mathfrak m\) costs \(O(\log(2N\mathfrak m))\).

Put \(r=f-1\). If \(r=0\) there is no archimedean condition apart
from the norm. In general use the quotient \(H/\Lambda\). Its
covolume is at most \(e^{Br}\), by the given independent vectors.
Every nonzero lattice vector has length at least an absolute
constant. In the totally real case a positive unit with all logs
of modulus less than \((\log 2)/3\) has every conjugate within
one of \(1\), so the integral norm of its difference from \(1\)
forces equality. In the CM case apply the same norm argument to
the totally positive unit \(a\bar a\), whose logs are twice
those of \(a\). A unit has absolute norm one, so its ordinary
and centered log vectors agree.

For an ideal \(\mathfrak a\) in the class of \(A\) let its point
in \(H/\Lambda\) be \(\lambda(s)\), where \(\mathfrak a=(s)A\)
and \(s\) satisfies the ray conditions. With \(\epsilon_0=\delta
B^{-5}\), use the periodic Gaussian
\[
 \Phi(z)=\sum_{u\in\Lambda}
              \exp\bigl(-\pi\|z-y+u\|^2/\epsilon_0^2\bigr).
\]
It is bounded by \(2\) for large \(B\). For example balls of
radius half the minimum separation are disjoint, so a ball of
radius \(R\) contains at most \((1+CR)^r\) lattice points.
At most one is closer than half the separation. Summing the
Gaussian over the successive remaining shells gives \(o(1)\),
also when its width is \(\sqrt2\,\epsilon_0\). Consequently
\(\Phi(z)\le2e^{-\pi B^{10}/2}\) if the distance on the quotient
from \(z\) to \(y\) is at least \(\delta\).

By Poisson summation the coefficient at \(v\in\Lambda^*\) has
modulus \(c_0e^{-\pi\epsilon_0^2\|v\|^2}\), where
\(c_0=\epsilon_0^r/\operatorname{covol}(\Lambda)\ge e^{-C B^2}\).
In dimension zero use \(c_0=\Phi=1\). The sum of the absolute
coefficients is at most \(2\), by evaluating the unshifted Gaussian
at zero. The part with \(\|v\|> B^4/\epsilon_0\) has sum at most
\(2^{1+r/2}e^{-\pi B^8/2}\), by halving the Fourier exponent
and applying Poisson summation at width \(\epsilon_0/\sqrt2\).

On the subgroup of ray-principal ideals, each \(v\) gives the
unitary character \((a)\mapsto e^{2\pi\mathrm i\langle
v,\lambda(a)\rangle}\). Extend it to the ideals prime to the
modulus. This is possible by adjoining generators of the finite
quotient and taking roots in the circle. The resulting character
is a Hecke character with the indicated archimedean parameters
and finite conductor dividing \(\mathfrak m\), by the ray
condition. Multiplying by the \(h\) characters of the finite
ray group and averaging restricts its sum to the class of \(A\).
The factor from dividing by \(A\) has modulus one. For every
retained nonzero \(v\) the preceding lemma thus bounds the
weighted ideal sum by \(O_g(Xe^{-cB^8})\), since
\(B^4/\epsilon_0\le e^{2B}\). For the omitted frequencies use
the bound \(O(nX\log X)\) for the total von Mangoldt weight
up to \(2X\). It follows for example by summing over powers of
rational primes and using that the sum of residue degrees above
each rational prime is at most \(n\).

Multiply the finite-class lower bound by \(c_0\). This is a lower
bound for the zero-frequency contribution to the nonnegative
weighted sum including prime powers. It dominates all nonzero
Fourier terms. Terms at prime powers \(\mathfrak p^j\) with
\(j[\OO_K/\mathfrak p:\mathbb F_p]\ge2\) contribute in total
\(O(n\sqrt{2X}\log^2(2X))\), since the underlying \(p\) is at
most \(\sqrt{2X}\). The bound outside the ball of radius \(\delta\)
removes at most \(O(nX\log X\,e^{-\pi B^{10}/2})\). A prime of
degree one inside the ball remains. Multiplying a ray generator
by a ray unit realizes a representative within \(\delta\) of
\(y\) in \(H\). Finally the mean of its logarithms is
\(n^{-1}\log(NP/NA)\), proving the coordinate assertion.
\end{proof}

\subsection{The level and class count of the integral column}\label{nav:II-57}
Retain the notation of Theorem~\ref{thm:row277}. Define \(J_W\)
on the original definite column by its prescribed covector
level \(K_{\pi,v}(\ell)\) above \(T\), its principal subgroup
of depth \(a_v+2j_v\) at \(v\in S_{\rm m}\), and its full
lattice stabilizer elsewhere. Put \(q_v=Nv\), and let \(N_x\)
be the common cyclotomic conductor of
Theorem~\ref{thm:conductor232}. Write \(\mathfrak m_l\) for
the mass of the standard positive Hermitian genus in rank \(l\).
For \(b\ge1\) put
\(G_l(q,b)=q^{b l^2}\prod_{i=1}^l(1-q^{-i})\).

\begin{corollary}[A common column level with a computable class count]\leavevmode
\label{cor:mass277}
Let \(h_W\) count the definite finite classes at \(J_W\).
The large class count fixing \(j\) and the small class count
in the integral pairing are
\begin{align*}
 h_H&=\mathfrak m_N\,G_N(T,1)^f
                  \prod_{v\in S_{\rm m}}G_N(q_v,a_v),\\
 h_W&=\mathfrak m_k\,(T^k-1)^f
                  \prod_{v\in S_{\rm m}}G_k(q_v,a_v+2j_v).
\end{align*}
At this \(J_W\), for each prescribed word and character there
is an integral column \(Q_\lambda\) with the bound
\[
  |\langle Q_\lambda,F_0 A_\lambda\rangle_{[H_0]}|
  \ \ge\ \frac{|Z|}{h_W}
       \sqrt{\frac{b_V b_L}{d_\mu I J'}}\,
       \frac{\sqrt{|Z|}(1-c_T)^{f/2}}{2\sqrt2}
       \sqrt{\frac{m_\nu}{d_\rho d_\mu}}\,M_B .
\]
Here \(I,J'\) are the sizes of the two integral generating
lists and \(b_V,b_L\) their lower frame bounds as in
Corollary~\ref{cor:column276}. All words on the fixed grid
use the same finite level. For a finite list, a single
integer linear combination of the rows retains at least
half of each individual lower bound at this level.
\end{corollary}
\begin{proof}
Both lattices are locally in the standard genus. At a split
place use one lattice basis with its dual. The orthonormal
basis proof in Proposition~\ref{prop:lattices277} applies at
every nonsplit place. At all real places both forms are positive
definite. The classification of Hermitian spaces used in
Theorem~\ref{thm:graph274} thus identifies each rational space
with the standard space of its rank. A finite adelic translate
of a lattice has the same genus and mass.

Refining its full finite stabilizer \(K\) to \(J_0\subset K\)
multiplies the mass by \([K:J_0]\). Indeed the refinements of
one class with automorphism group \(\Gamma\) are its orbits
on \(K/J_0\), after choosing a representative. The sum of the
reciprocal stabilizer orders on these orbits is
\([K:J_0]/|\Gamma|\). For each of the two levels in question
the stabilizers after refinement are trivial, by the marked
odd principal subgroup. Mass is therefore the number of
classes. At a split place the congruence index in rank \(l\)
is \(|\mathrm{GL}_l(\mathcal O_v/\pi_v^b)|=G_l(q_v,b)\).
Reduction modulo \(\pi_v\) counts successive independent
columns, and every invertible residue matrix has exactly
\(q_v^{(b-1)l^2}\) lifts. At a terminal place the small
covector stabilizer has index \(T^k-1\) by transitivity on
the nonzero residue covectors. This proves both formulas.

These expressions bound the cost using field and grid data
alone. For example the mass formula and Bernoulli estimate in
Lemma~\ref{lem:mass209} give
\(\log\mathfrak m_l\ll f l^2\log(2lN_x)\) as an upper
bound. The primitive factors at \(j=1\) also satisfy the
stated Bernoulli upper bound directly from \(B_1(y)=y-1/2\).
With \(D_{\rm m}=\sum_{v\in S_{\rm m}}a_v\log q_v\) the
preceding formulas give
\begin{align*}
 \log h_H&\ll fN^2\log(2NN_x)+fN^2\log T+N^2D_{\rm m},\\
 j\log Q&\ll \log Q+\log(2h_H d_\rho^3),\\
 \log h_W&\ll fk^2\log(2kN_x)+fk\log T
                         + k^2(D_{\rm m}+2j\log Q).
\end{align*}
For the middle inequality, \(\log Q/(-\log c_Q)\) is bounded
for \(Q\ge3\), by its formula and its limit one at infinity.
For several detecting types use the maximum \(d_\rho\).

Theorem~\ref{thm:row277} gives right \(J_W\)-invariance of
the projected coefficient, using centrality of \(Z\) on the
column and equivariance of the Fourier integral. The \(|Z|\)
translates of the identity class remain distinct by the
definite-lattice and torsion-free argument of
Corollary~\ref{cor:column276}. Each has mass \(1/h_W\).
Applying the integral matrix coefficient construction of that
corollary on the identity class and summing its translates
against \(\bar\lambda\) gives the first factor
\((|Z|/h_W)\sqrt{b_Vb_L/(d_\mu I J')}\) times \(\|p\|_2\).
Now insert Theorem~\ref{thm:row277}. Its row levels are also
common across the prescribed words and characters. The
integer coefficient argument of Corollary~\ref{cor:column276}
therefore applies on these levels.
\end{proof}

This bounds the loss due to the finite class chosen in extracting
the restricted coefficient. The arithmetic normalization of
\(M_B\), survival in the prescribed global cuspidal row at its
original marked levels, the upper estimate for the full theta
period and the uniform strict cohomology estimate still require
proof on the common intervals.

\subsection{The integral correction on all the original fields}\label{nav:II-59}
Use \(E=E_x\), its degree \(n\), a conductor multiple \(N\)
from Theorem~\ref{thm:conductor232}, and \(t=\lfloor x^{401}\rfloor\).
For odd \(p\) let \(A=(\QQ_p/\ZZ_p)(-r)\),
\(G=\Gal(E/\QQ)=P\times H\), with \(P\) the \(p\)-Sylow group.
It is cyclic of order \(p^a\). Indeed restriction embeds it in
the cyclic group for \(K_\psi\), because the relative degree
over \(K_\psi\) is a power of two. Put \(\eta=(-1)^r\).
For a finite group with involution \(c\), superscript \(\eta\)
below means its \(\eta\)-eigenspace for complex conjugation.

Let \([\chi]\) run over \(p\)-adic Galois orbits of characters
of \(H\) of sign \(\eta\). Put \(b_\chi=[\ZZ_p[\chi]:\ZZ_p]\).
Use Theorem~\ref{thm:intersection279} with this \(P,\chi,N\),
writing \(H_i(\chi)\) and \(I_{0,i,\chi},I_{1,i,\chi}\) for
its character lengths and residue subspaces. Define
\[
 \begin{split}
 A_{p,r}&=(\log p)\sum_{[\chi]} b_\chi\sum_{i=0}^a H_i(\chi),\\
 B_{p,r}&=\log\left|\left(\bigoplus_{w\mid Np}A^{G_{E_w}}\right)^\eta\right|
                  -\log\left|(A^{G_E})^\eta\right|.
 \end{split}
\]
Here \(G_{E_w}\) and \(G_E\) are the absolute local and global
Galois groups. All these invariant groups are finite, since the
local cyclotomic image is open at \(p\), and at \(w\nmid p\)
Frobenius acts by \((\operatorname{Norm}w)^{-r}\ne1\).
At \(i=0\) the finite character is \(\chi\). At \(i>0\) it is
\(\chi\) times the chosen character of order \(p^i\).
Thus each \(H_i(\chi)\) is given by the index and explicit
Euler factors of Theorem~\ref{thm:openindex269}.

\begin{theorem}[A finite residue correction to the sum of character lengths]\leavevmode
\label{thm:common279}
For every odd \(p\) and \(1\le r<t\),
\[
 \log|\mathcal S_{p,r}(E)^\eta|
       = A_{p,r}-B_{p,r}-D_{p,r},
 \qquad
 D_{p,r}=(\log p)\sum_{[\chi]} b_\chi
       \sum_{i=1}^a
       \bigl(p^{i-1}-\dim(I_{0,i,\chi}\cap I_{1,i,\chi})\bigr).
\]
The dimension in each term is over the residue field of
\(\ZZ_p[\chi]\). The total correction everywhere on the given
interval satisfies
\[
  0\ \le\ \sum_{p>2}\sum_{r=1}^{t-1}D_{p,r}
   \ \le\ n(t-1)
     \sum_{\substack{p\mid\operatorname{ord}\psi\\p\ {\rm odd}}}
                           \frac{\log p}{p-1}.
\]
\end{theorem}
\begin{proof}
On \(U=\operatorname{Spec}\ZZ[1/(Np)]\) the field \(E\)
defines a finite étale cover. Shapiro's lemma represents its
integral degree-two cohomology at twist \(r+1\) by the
permutation lattice \(\ZZ_p[G](r+1)\) over \(\QQ\).
The \(H\)-group algebra is the product of the rings
\(\ZZ_p[\chi]\) over the indicated orbits before restricting
the sign. Indeed after an unramified splitting extension the
usual character idempotents have the invertible denominator
\(|H|\). Frobenius permutes their factors, and descending an orbit gives
its own unramified integer ring. Complex conjugation lies in
\(H\), so taking sign \(\eta\) selects the stated factors.
Choosing the inverse labeling of the regular action permutes
the same set of characters. For each orbit the remaining group
ring is precisely \(R_a\) with its character twist. Its
degree-two length is \(M_a\) of
Theorem~\ref{thm:intersection279}. The logarithm of the order
of this factor is \(b_\chi M_a\log p\). The theorem's identity
therefore gives \(A_{p,r}-D_{p,r}\) for the logarithm of the
order of the integral degree-two group in this sign.

At any \(w\nmid p\) the twist \(A\) is unramified and its
unramified degree-one subgroup is
\(A/((\operatorname{Norm}w)^{-r}-1)A=0\), by divisibility.
Opening the primes over \(N\) and imposing zero localization
there thus gives exactly the original strict group. The
Poitou--Tate sequence identifies its dual with the kernel of
the map from the integral degree-two group to the sum of the
local degree-two groups. The cokernel of that map is
\((A^{G_E})^\vee\), and the local groups are
\((A^{G_{E_w}})^\vee\) by local duality. The real-place
modified groups vanish at odd \(p\). All maps commute with
the field action, and the two sign projectors are integral.
Their finite duals have the same orders in each sign.
Consequently the logarithm to subtract is \(B_{p,r}\).

The sum of \(b_\chi\) over the retained orbits is at most
\(|H|=n/p^a\), since it counts their individual characters.
The correction for each orbit is between zero and
\((p^a-1)/(p-1)\) in normalized length. Summing this bound
over the twists gives the displayed estimate, since a nontrivial
\(P\) forces \(p\mid\operatorname{ord}\psi\). On the common
intervals \(n\ll_\psi x^{2+8\varepsilon}\), so this contribution
is polynomial throughout. The character indices within
\(A_{p,r}\) still require the full upper estimate.
\end{proof}

\subsection{Reduction of the cyclotomic class at a split prime}\label{nav:IV-40}
Keep \(f,\alpha,\ell,r,O\) from Theorem~\ref{thm:circindex267}.
For \(h\ge1\) and \(z\in O/\ell^hO\) set
\[
 w_h(z)=\max\{0\le j\le h:z\in\ell^j(O/\ell^hO)\}.
\]
Thus \(w_h(0)=h\). For a prime \(s\equiv1\bmod f\ell^h\),
take a generator \(g\) of \(\mathbb F_s^\times\) and put
\(w=g^{(s-1)/(f\ell^h)}\). Write \(\operatorname{ind}_g y\)
for the exponent of \(y\) with respect to \(g\), reduced
modulo \(\ell^h\). Define
\[
 D_{s,h}(\alpha,r)=
 \sum_{\substack{1\le a\le f\ell^h\\(a,f\ell)=1}}
       \alpha(a)\,a^r\,\operatorname{ind}_g(1-w^a)
       \quad\text{in }O/\ell^hO .
\]
The arguments of the indices are nonzero since \(w\) is
primitive of order \(f\ell^h>1\). One can compute the index
using only the subgroup of order \(\ell^h\). It is the unique
\(j\bmod\ell^h\) for which
\((1-w^a)^{(s-1)/\ell^h}=(w^f)^j\).

\begin{proposition}[A finite-field bound for the whole primary order]\leavevmode
\label{prop:moment267}
At any such prime the Frobenius value of \(c_{\alpha,r+1}\)
at precision \(\ell^h\), in compatible bases, is
\(D_{s,h}(\alpha,r)\). Put \(d=v_\ell(a_{\alpha,r})\),
where \(a_{\alpha,r}\) is the coordinate of the cyclotomic
class in a basis of \(H^1(U,T)\). Then
\[
  w_h(D_{s,h}(\alpha,r))\ge\min(h,d).
\]
Equality holds when the basis class has nonzero residual
Frobenius value at \(s\). In particular, any nonzero moment
gives \(\operatorname{length}_O S(A)\le w_h(D_{s,h}(\alpha,r))\).
At \(h=1\), a nonzero moment proves vanishing of the full
group for every positive twist in the same residue class
modulo \(\ell-1\) and of the prescribed parity.
\end{proposition}
\begin{proof}
At \(s\) the Frobenius action on \(T/\ell^hT\) is trivial,
as is inertia. The value on arithmetic Frobenius of an
unramified class therefore gives an element of this module
independent of its cocycle representative. Identify
\(\zeta_{f\ell^h}\) with \(w\) in one residue field and use
\(w^f\) as basis for the local roots of order \(\ell^h\).
If the globally chosen primitive root reduces to a different
one, relabel the embeddings over \(s\) and change the local
basis. This can multiply the coordinate by a unit and preserves
\(w_h\).

The restriction of a corestriction at a completely split
prime is the sum over all embeddings into \(\QQ_s\).
This is the double-coset formula for restriction and
corestriction, with each local corestriction the identity.
The embeddings send \(\zeta_{f\ell^h}\) to \(w^a\), for the
unit classes \(a\) in the sum. For the Kummer class of a
residue unit \(y\), choose \(z\) in the algebraic closure
with \(z^{\ell^h}=y\). Then
\(\operatorname{Frob}_s(z)/z=y^{(s-1)/\ell^h}\).
Thus its coordinate in the basis \(w^f\) is
\(\operatorname{ind}_g y\). The additional \(r\) tensor
factors of \(\zeta_{\ell^h}\) contribute \(a^r\), and the
action on \(t_\alpha\) contributes \(\alpha(a)\).
Adding gives the stated finite sum. This computation is also
the residue method for higher circular units at odd weights
over real abelian fields in
\href{https://arxiv.org/pdf/0912.0819\#page=19}{Beliaeva and
Belliard, Sections 6--7}.

Evaluation is \(O\)-linear. Multiplication of the basis class
by \(a_{\alpha,r}\) therefore gives the asserted inequality,
and gives equality if the basis evaluation is a unit.
If the moment is nonzero, \(d<h\) and
\(d\le w_h(D_{s,h})\). Use Theorem~\ref{thm:circindex267} to
bound the strict order. Finally at \(h=1\) every \(a\) in
the sum is prime to \(\ell\), so that \(a^r\bmod\ell\) only
depends on \(r\bmod(\ell-1)\). The same sum gives the
assertion at every twist indicated.
\end{proof}

\subsection{An average for the strict orders on the full interval}\label{nav:IV-41}
Use the list of large-prime triples in Theorem~\ref{thm:laws268}.
Choose \(h_\ell\) at least the largest length at that prime.
At \(x\in[V,V^{3/2}]\) let \(I_x\) consist of the character
orbits of \(E_x\), for \(1\le r<t(x)\) and
\(\ell>\exp(x^\varepsilon/12)\), at the coefficient primes
that occur in the interval list. Include the zero components
at these primes if desired. Enlarge the interval list by those
orbits before applying the theorem. A prime entirely absent
from the support contributes zero to the strict order.
Define the two real numbers
\[
 L_x=\sum_{j\in I_x}b_jd_j\log\ell_j,\qquad
 U_x(s)=\sum_{j\in I_x}b_jW_j(s)\log\ell_j .
\]
The character decomposition over the unramified integer rings
identifies \(L_x\) with the sum of the full logarithmic orders
of the large-prime complementary groups on \(E_x\) when all
supported orbits are included. Restriction to a larger period
field identifies the original character component with its
copy upstairs, by prime-to-degree descent. Repetitions of that
component are counted just once at any given \(x\).

\begin{theorem}[One average with polynomial additive error]\leavevmode
\label{thm:mean268}
For a constant \(C=C(\psi,\varepsilon)\), for all sufficiently
distant original intervals,
\[
 0\le {\mathbb E}_{s\ {\rm split\ in}\ K}
       \sup_{V\le x\le V^{3/2}}\bigl(U_x(s)-L_x\bigr)
       \le V^C .
\]
The same estimate holds simultaneously for every sublist of
the given triples, with the supremum taken over those sublists
as well. In particular the mean of \(U_x\) recovers \(L_x\)
with additive error \(o(\exp(x^\varepsilon/2))\), uniformly
on this interval. For every fixed \(A>0\), a set of relative
density at least \(1-V^{-A}\) among primes split in \(K\)
has simultaneous error at most \(V^{C+A}\).
\end{theorem}
\begin{proof}
The summands in the difference are nonnegative. For the
supremum over all sublists it suffices to add the excesses
over the entire list once. These are finitely many
finite-valued variables, so their mean is the ordinary
limit of prime averages by Theorem~\ref{thm:laws268}. For
one summand, summing the probabilities of the positive
integer levels gives the exact unweighted excess
\[
 {\mathbb E}(W_j-d_j)
    =\sum_{a=1}^{h_j-d_j}\ell_j^{-b_j a}
    =\frac{1-\ell_j^{-b_j(h_j-d_j)}}{\ell_j^{b_j}-1}.
\]
In particular the error estimate is independent of all the
precisions and lengths. At a given prime the sum of the
\(b_j\) is at most \(n_*t_*\). This follows by decomposing
the regular representation of the abelian Galois group of
the largest original period field into residue-field
orbits at each twist.

The support bound in the proof of
Theorem~\ref{thm:commonlength267} gives at most
\(Z=\exp(V^{C_0})\) different coefficient primes, with \(C_0\)
fixed independently of the interval and of the precisions.
For the primes \(\ell\le Z^2\), bound
\(\sum \log\ell/(\ell-1)\) by summing over all integers
between \(2\) and \(Z^2\). This costs
\(O((1+\log Z)^2)\) by the harmonic sum. Above \(Z^2\),
the function \(\log y/(y-1)\) is decreasing, so the
remaining sum is at most \(Z\log(Z^2)/(Z^2-1)\).
Multiplying by \(n_*t_*\), polynomial in \(V\), gives the
asserted bound after increasing \(C\). A polynomial in \(V\)
is \(o(\exp(x^\varepsilon/2))\) everywhere on the stated
interval. Markov's inequality for the full excess proves the
relative-density assertion. The finite set of possible
sublists on the entire interval justifies each supremum.
\end{proof}

This places all of the large-prime contribution under one
average of explicit finite-field quantities on the given
period fields. A bound for that average itself on the scale
of \eqref{eq:region198} remains necessary. The splitting field
and the set of coefficient primes can both grow with the interval.

\subsection{Integral character comparison at every odd prime}\label{nav:IV-42}
Use an original interval of Theorem~\ref{thm:conductor232}.
Let \(E=E_x\), \(G=\operatorname{Gal}(E/\mathbb Q)\),
\(n=|G|\), and let \(N=N_*\) be the modulus of the largest
field on the interval. Write \(c_\infty\in G\) for complex
conjugation and \(t=t(x)\). In the following sum let \(\alpha\)
range once over the \(p\)-adic Galois orbits of characters of
\(G\) satisfying \(\alpha(c_\infty)=(-1)^r\). Use their
primitive conductors. With the residue degree and the index of
Theorem~\ref{thm:openindex269}, put
\[
 C(E,t;N)=\sum_{\substack{p\ {\rm odd}\\1\le r<t}}
                 \ \sum_{[\alpha]} b_\alpha
                      d_{p,r,\alpha}(N)\log p .
\]

\begin{theorem}[All odd primary orders on the original fields]\leavevmode
\label{thm:lattices269}
This sum has finite support. Its difference from the sum of
the logarithmic orders of the complementary odd strict
groups on \(E\) is bounded in absolute value by
\([\,2nt\log(2N)\,]^{C_\psi}\) for a constant \(C_\psi\).
The same bound holds for its difference from
\({\cal A}(E,t)\) in \eqref{eq:A100}. These comparisons use
the full primary orders, including primes dividing \(|G|\).
\end{theorem}
\begin{proof}
Fix \(p,r\), put \(k=r+1\), and write
\(e_r=(1+(-1)^r c_\infty)/2\). On \(U_N\) compare the two
Galois lattices
\[
 L=e_r\mathbb Z_p[G](k)\ \lhook\joinrel\longrightarrow\
 L'=\bigoplus_{[\alpha]}O_\alpha(\alpha\kappa_p^k).
\]
Use the regular action and evaluation at the characters.
Fourier inversion over \(\overline{\mathbb Q}_p\) proves
injectivity. Multiplication by \(n\) clears the only
denominator of its inverse. Indeed the inverse coefficients
are sums of values times inverse roots of unity divided by
\(n\). The sums are invariant under coefficient Galois action
and are \(p\)-integral. Thus the map is onto when \(p\nmid n\).
Otherwise \(D=L'/L\) has order at most \(p^{an}\), where
\(a=v_p(n)\). Since \([E:K_\psi]\) is a power of two, these
odd primes divide the fixed order of \(\psi\), and \(a\) is
bounded in terms of \(\psi\).

We bound the error in degree two at those fixed primes.
The action on \(D\) becomes trivial on
\(F=E(\mu_{p^a})\), of degree \(O_\psi(n)\) and conductor
dividing \(Np^a\). Inflation--restriction bounds the first
global cohomology with coefficient \(D\) by the product of
\(|D|^{[F:\mathbb Q]}\) and
\(|\operatorname{Hom}(G_{F,S},D)|\), where \(S\) comprises
the primes over \(Np\) and infinity. In fact every cocycle
of the finite quotient is one of the indicated functions,
and the restricted cocycle is a homomorphism. Since
\(\mu_p\subset F\), the Kummer sequence bounds
\(\dim_{\mathbb F_p}H^1(G_{F,S},\mathbb F_p)\) by
\([F:\mathbb Q]+|S|+1+\log|\operatorname{Cl}(F)|/\log p\).
The terms come from \(S\)-units and the \(p\)-torsion of the
\(S\)-class group. The conductor--discriminant formula gives
\(\log|\operatorname{disc}F|\ll_\psi n\log(2N)\).
Minkowski's ideal bound and the estimate
\(\zeta_F(2)\le\zeta(2)^{[F:\mathbb Q]}\) thus give
\(\log|\operatorname{Cl}(F)|\ll_\psi n\log(2N)\).
Also \(|S|\ll_\psi n\log(2N)\). Filtration by groups of order
\(p\) bounds the logarithm of the homomorphism count by this
dimension times \(\log|D|\). Hence
\(\log|H^1(U_N,D)|\ll_\psi n^2\log(2N)\).
The finite-module global Euler characteristic formula over
\(\mathbb Q\) gives the same bound for
\(\log|H^2(U_N,D)|\), increasing the constant if necessary.
The long exact sequence then bounds the absolute difference
between the logarithmic degree-two orders for \(L,L'\) by
the sum of these two bounds.

For \(L\), Shapiro identifies the open degree-two group
with \(e_rH^2(\mathcal O_E[1/Np],\mathbb Z_p(k))\).
Pass back to \(\mathcal O_E[1/p]\). At each removed place
\(w\nmid p\), purity gives zero in the preceding residue
term \(H^0(\kappa(w),\mathbb Z_p(r))\). The next term has
order the \(p\)-part of \((Nw)^r-1\), by arithmetic
Frobenius on \(\mathbb Z_p(r)\). Thus the change in the
logarithmic order in either sign, summed over odd \(p\),
is at most \(nr\log N\), since for each \(q\mid N\) one has
\(\sum_{w\mid q}\log Nw\le n\log q\).
Degree two on \(\mathcal O_E[1/p]\) is the \(p\)-part of
\(K_{2r}(\mathcal O_E)\). Its complementary sign gives the
order appearing in \({\cal A}(E,t)\), by
Theorem~\ref{thm:res100}.

For \(L'\), apply Theorem~\ref{thm:openindex269} to each
summand. At a fixed \(q\mid N\), form the product of
\(1-\alpha(q)q^r\) over all geometric characters in the
specified sign that are unramified at \(q\). It is a nonzero
rational integer by Galois invariance and \(r>0\).
At each \(p\ne q\) its \(p\)-valuation is the sum of the
character valuations weighted by \(b_\alpha\), by the local
norm formula even in ramified coefficient fields. Its
absolute value is at most \((1+q^r)^n\).
The sum of all these local-factor costs over \(q\mid N\)
and all odd \(p\) is therefore \(O(nr\log(2N))\).

Finiteness of the even \(K\)-groups, finiteness of the Euler
products at each \(r\), and the finite list of primes dividing
\(n\) now also prove finite support for the indices. Summing
the preceding bounds for \(1\le r<t\) proves the comparison
with \({\cal A}\).
To recover the strict group, Theorem~\ref{thm:wild98} bounds
the change by the local invariants above \(p\). If such an
invariant has order \(p^j>1\) over \(E_w\), the cyclotomic
image in \((\mathbb Z/p^j)^\times\) has index at most \(n\)
and exponent dividing \(r\). As that group is cyclic,
\(\varphi(p^j)\le nr\). The invariant group is cyclic of
order at most \(2nr\), and it can occur only at \(p\le nr+1\).
There are at most \(n\) places over each such prime.
The resulting total is \(O(n^2t^2\log(2nt))\).
Taking sign summands preserves this upper bound. All errors
have the asserted polynomial size.
\end{proof}

\subsection{The distribution when only one twist is split}\label{nav:IV-43}
Take \(N_*,p,k,\alpha,O,\pi,e,v,h\) as in
Proposition~\ref{prop:orbits273}. Write \(q=|O/\pi|\),
\(d=d_{p,r,\alpha}(N_*)\), and \(\nu=\min(h,v_p(N_*))\).
Put
\[
       t_h=\frac{\varphi(p^h)}{\varphi(p^\nu)},\qquad
       m_h=\frac{t_h}{\gcd(t_h,k)},\qquad\varphi(1)=1 .
\]
All densities below are relative natural densities among rational
primes \(s\equiv1\bmod N_*\), omitting ramified primes in each
fixed finite calculation.

\begin{proposition}[A conditional Frobenius law and its finite difference]\leavevmode
\label{prop:law273}
The condition \(p^h\mid s^k-1\) has density \(m_h^{-1}\).
Given this condition, for \(0\le j\le v\) one has
\[
         \Pr\bigl(D_v(s)=0\bmod\pi^j\bigr)
                              = q^{-\max(0,j-d)} .
\]
For an independent pair \(s_1,s_2\) write \(I_v=1\) if
both satisfy the condition and zero otherwise. On \(I_v=1\)
let \(J_j\) equal one if both orbit sums vanish modulo
\(\pi^j\), and zero otherwise. Take \(J_0=1\) there. Set
the following variables to zero off the eligible event:
\[
     P_v=m_h^2 I_vJ_v,\qquad
     H_v=m_h^2 I_v\frac{J_v-q^{-2}J_{v-1}}{1-q^{-2}} .
\]
Then \(\mathbb E P_v=q^{-2\max(0,v-d)}\) and
\(\mathbb E H_v=\mathbf1_{\{v\le d\}}\).
\end{proposition}
\begin{proof}
Modulo \(p^h\) the allowed residue classes are the units equal
to one modulo \(p^\nu\), by the Chinese remainder theorem.
They form a cyclic group of order \(t_h\) since \(p\) is odd.
The \(k\)-th power homomorphism has image of order \(m_h\).
Equidistribution of prime residues modulo
\(\operatorname{lcm}(N_*,p^h)\) gives the first density.

For the second calculation fix any eligible residue class
also modulo \(n=fp^h\) compatible with \(s=1\bmod N_*\).
Put \(K=\mathbb Q(\mu_{\operatorname{lcm}(N_*,n,p^h)})\)
and first consider a basis cocycle \(u\) of the free
\(H^1\) of Theorem~\ref{thm:openindex269}.
Its values on \(G_K\) map onto the additive group \(O/\pi^v\).
Here is the argument. The residual basis class is nonzero
by the multiplication sequence. Restriction to \(G_K\) is
nonzero because the abelian cyclotomic quotient has zero
first cohomology on the residual line. Indeed complex
conjugation is central there and acts by \(-1\) on that
line. Evaluating a cocycle on a product in both orders
expresses it as a coboundary, since two is invertible.
The action on the coefficient is trivial on \(G_K\) at
the indicated precision, so the image is additive and
contains a unit. Conjugation by a global element preserves
it and acts on it by \(\alpha\kappa_p^k\). It is already
stable under the scalar integers modulo \(p^h\), so it
is stable under all values of \(\alpha\) and hence under
\(O=\mathbb Z_p[\alpha]\). An ideal containing a unit
is the full ring.

On any fixed lift of the eligible cyclotomic Frobenius,
multiplication by elements of \(G_K\) now makes the basis
cocycle uniform on \(O/\pi^v\) in the finite quotient
recording the coefficient, the cocycle and \(K\).
Conjugation multiplies its value by a unit on this
fiber, since the coefficient of the eligible Frobenius
is one modulo \(\pi^v\). All divisibility tests are
therefore conjugacy invariant. Chebotarev's density
theorem gives the finite-group proportions as prime
densities. The actual class is a scalar of valuation
\(d\) times the basis class. A possible coboundary
vanishes on eligible Frobenius. Reduction to level \(j\)
gives the displayed probability, independently of the
chosen compatible residue. Averaging the residues proves
the conditional law.

For two sampled primes the conditional probabilities multiply.
The event \(I_v=1\) has density \(m_h^{-2}\), proving the
formula for \(P_v\). For \(H_v\) the numerator of its
expectation before division by \(1-q^{-2}\) is
\[
        q^{-2\max(0,v-d)}
            -q^{-2}q^{-2\max(0,v-1-d)} .
\]
It equals \(1-q^{-2}\) for \(v\le d\) and zero otherwise.
Both indicators use the same eligibility condition at
level \(v\), including when \(j=v-1\).
\end{proof}

\subsection{The period field and the compact types of its integral}\label{nav:VI-25}
We distinguish the period field
\(\widetilde E=K_\psi C_{l'}(\sqrt{-q},\sqrt{-d_0})\) of
Proposition~\ref{prop:desc50} from the upper field \(E_x\) in
Theorem~\ref{thm:tower104}. We write \(s_0=\lfloor u\eps\log_2 x\rfloor\)
for the subfield index of that theorem and reserve \(s=m/2-1\)
for the Witt flag length. All the choices below are made on its
original selected intervals, with their primes \(q,d_0\).

\begin{proposition}[The bounded summand on the period field]\label{prop:short107}
For every sufficiently large \(x\) in these intervals, every odd
prime \(p\), and every \(r\ge1\), restriction from \(L_{u,x}\)
identifies the two invariant strict groups
\[
 \mathcal S_{p,r}(\widetilde E)^{\Gal(\widetilde E/L_{u,x})}
 \ \simeq\ \mathcal S_{p,r}(L_{u,x})\ \simeq\
 \mathcal S_{p,r}(E_x)^{\Gal(E_x/L_{u,x})}.
\]
In particular the bounds \eqref{eq:pub104} and \eqref{eq:vor104}
hold with \(\widetilde E\) as the upper field. The residual
decomposition \eqref{eq:rest104} holds there with its own strict
groups and averaging projectors.
\end{proposition}
\begin{proof}
Put \(o=[K_\psi:\QQ]\). Proposition~\ref{prop:desc50} gives
\(x^{2+8\eps}\le[K_\psi C_{l'}:\QQ]\le o2^{l'}\).
For \(0<u<1\), this implies \(s_0\le l'\le l\) for all
sufficiently large \(x\). Both fields in the display are
therefore extensions of precisely the same \(L_{u,x}\) of
two-power degree. In positive group cohomological degrees the
two finite Galois groups have zero cohomology on \(p\)-primary
modules, by averaging. Inflation--restriction thus identifies
the unrestricted degree-one group downstairs with the invariants
upstairs. At \(p\), local restriction is injective since the
local extension degree is prime to \(p\). Away from \(p\) the
coefficient module is unramified and an inertia homomorphism
that becomes zero factors through a group of two-power order.
Thus the local conditions descend in both extensions. This gives
the identifications, through the full primary exponent.

At odd primes the averaging operator on the shortened field and
the two sign projectors commute, since the field is abelian.
The critical-sign lengths on the shortened field are bounded by
those on the upper field, by injectivity. At two,
Proposition~\ref{prop:desc50} bounds the extra weighted length
by \(t\log(2o)\), which is polynomial in \(x\) and hence
\(o(\exp(x^\eps/2))\). The proof of Corollary~\ref{cor:rest104}
now gives the assertion for \(\widetilde E\). To omit individually
certified components one applies
Theorem~\ref{thm:stable102} on the field in question.
\end{proof}

\begin{lemma}[An integral embedding of a compact type]\label{lem:shuffle107}
Let \(A_1\ge\cdots\ge A_P\) be integers and
\(I=\{i_1<\cdots<i_k\}\subset\{1,\ldots,P\}\). Put
\(\nu=(A_{i_1},\ldots,A_{i_k})\). In the coordinate tensor
lattices defined below there is a primitive \(U(k)\)-equivariant
isometric embedding \(J_I:V_\nu\longrightarrow V_A\).
The induced invariant Hermitian pairing is integral on these
lattices and takes value one on a pair of primitive vectors.
\end{lemma}
\begin{proof}
For \(W=\ZZ^P\) write \(d_j=A_j-A_{j+1}\) for \(1\le j<P\)
and use the ambient tensor
\[
 {\cal T}_A(W)=
  \bigotimes_{j=1}^{P-1}(\bigwedge\nolimits^j W)^{\otimes d_j}
               \otimes(\det W)^{A_P}.
\]
For a negative power of a determinant use the dual line. Give
exterior coordinate monomials norm one and make them orthogonal.
Let \(w_A\) be the unit coordinate tensor obtained from the
first \(j\) basis vectors in every degree \(j\), and the chosen
generator of the determinant line. It has torus weight \(A\)
and is killed by all upper triangular raising operators. By
complete reducibility it generates an irreducible complex
representation of highest weight \(A\). Indeed, the projections of
this vector to irreducibles can be nonzero only on their
highest weight lines of weight \(A\), and together they generate
a single diagonal copy. Define the lattice by intersecting this
representation with the integral ambient tensor. It spans the
representation, since it is generated over \(\CC\) by products
of matrix Lie operators with integral coefficients applied to
\(w_A\). Use the identical construction for \(\nu\).

Write \(W=X\oplus Z\) for the coordinate splitting into ranks
\(k,P-k\). Order a basis of \(W\) by putting the ordered basis
of \(X\) in the slots \(I\) and that of \(Z\) in the others.
Put \(b_j=|I\cap[1,j]|\). In each row factor of exterior degree
\(j\), send an exterior tensor of degree \(b_j\) in \(X\) to its
wedge with the first \(j-b_j\) basis vectors of \(Z\), with the
sign for this ordering. This sends distinct basis monomials to
distinct signed basis monomials, isometrically and integrally.
For \(1\le t<k\), the multiplicity of degree \(t\) in the
domain of these maps is
\[
     \sum_{j=i_t}^{i_{t+1}-1} d_j
          = A_{i_t}-A_{i_{t+1}}.
\]
Degree zero contributes unit lines. The total determinant
power in \(X\) is
\(A_P+\sum_{j=i_k}^{P-1}d_j=A_{i_k}\).
The fixed determinant of \(Z\) accounts for the remaining
line. Grouping these factors identifies the domain with
\({\cal T}_\nu(X)\) by a signed permutation of coordinate
tensors, including the canonical identifications of determinant
lines and their duals. All identifications commute with \(U(k)\)
after extending scalars to \(\CC\).

The tensor embedding sends \(w_\nu\) to \(w_A\) written in the
permuted basis, up to a fixed sign. This is a vector of \(V_A\),
since the permutation matrix lies in \(U(P)\). Its \(U(k)\)-span
gives the required copy of \(V_\nu\). The embedding identifies
the integral ambient domain with a coordinate direct summand.
Consequently an element of the complex image that lies in the
integral row lattice has integral preimage, proving primitivity.
Pairing a row vector with the embedded column by the ambient
Hermitian form gives integral coefficients. Pairing \(J_Iw_\nu\)
with itself gives one, and Cauchy--Schwarz gives operator norm one.
\end{proof}

\begin{theorem}[One compact row for every translated column]\label{thm:pair107}
Keep the fields, intervals, column parameters, and finite levels
of Theorem~\ref{thm:target90}. There is one common regular row
parameter and one chamber, independent of the translate, for
which its lowest compact type contains both the high and low
column types at every real place. All prescribed products of
these types have integral compact pairings of norm one in a
single row coefficient lattice. Each pairing takes the value
one there. For a common admissible \(\Gamma\) one terminal prime
works simultaneously for all translates at the formal-degree
scale of this row.
\end{theorem}
\begin{proof}
Write \(P=k+s,\ Q=s+1,\ r=P+Q\). The old column parameters are
those in the proof of
Theorem \ref{thm:zero106}.
Let \(\nu^{\,0}\) be the compact highest weight for a low place,
and \(\nu^{\,a}=\nu^{\,0}+(a,0,\ldots,0)\). Keep every shared
entry \(v_i=2(k-1-i)\), \(1\le i<k\). At every real place take
the same ordered row list
\[
 v_1+a,\quad v_1+1,\quad v_1,v_2,\ldots,v_{k-1},\quad
      -1,\quad -(t_{\rm r}+4),\ldots,-(m-3)(t_{\rm r}+4).
\]
It has \(r=k-1+m\) distinct integer entries for \(a\ge2\).
It retains the original interlacing
\(u_i>v_i>u_{i+1}\) with each column. Give the second slot
negative signature and all the other slots through \(-1\)
positive signature. Of the remaining \(2s-1\) slots choose
the first \(s-1\) positive and the rest negative. This yields
the required signature \((P,Q)\). For the positive entries
\(j_1=1\) and \(j_i=i+1\) for \(2\le i\le k+1\).
The compact highest weight formula of equation \ref{eq:mu106}, now gives
\[
          (A_1,\ldots,A_{k+1})=
  (\nu^{\,a}_1,\nu^{\,0}_1,\nu^{\,0}_2,\ldots,\nu^{\,0}_k).
\]
Indeed a slot occupied by \(v_i\) gives
\(A_{i+1}=v_i-k/2+i\), and the slot \(-1\) gives
\(A_{k+1}=k/2-1=\nu^{\,0}_k\).

Apply Lemma~\ref{lem:shuffle107} to \(I_0=\{2,\ldots,k+1\}\)
and \(I_a=\{1,3,\ldots,k+1\}\). Tensor these two embeddings
separately with a fixed integral highest vector of norm one
in the \(U(Q)\)-factor, on which \(H_\tau\) acts trivially.
The scalar circle of \(U(k)\) acts on the images by characters
whose exponents differ by \(a\). Invariance of the inner product
therefore makes the two images orthogonal. Tensoring the appropriate embeddings over the real
places preserves both the norm and the value one on the chosen
integral vectors. Distinct sets of high places give orthogonal
images by the same circle calculation at a place where they
differ. The lattice used throughout is the tensor product of
these split coordinate coefficient lattices.

The same row at every place fixes the product of its formal
degrees. Every translate has \(h\) high places and \(f-h\) low
places, so the column dimension product is always
\(W_k(u^{\,a})^h W_k(u^{\,0})^{f-h}\), where \(u^{\,a},u^{\,0}\)
denote the two column parameters. Hence \(D_*^{\rm row}\),
\(d_{u,*}\), and the target \(T_*\) of Corollary~\ref{cor:real91}
are independent of the translate for a common \(\Gamma\).
That corollary, using the interlacing checked above and the
original column spans, gives a prime in
\((T_*,(1+1/(fk))T_*]\) satisfying all its finite congruences.
In Lemma~\ref{lem:ray88} one may double the modulus as before
Theorem~\ref{thm:neighbour93}, adding only \(\log 4\) to
\(\log(4dM^2)\). Choose the prime with this stronger congruence.
The row parameter has size at most \(a+2k+m(t_{\rm r}+4)\),
which is at most \(\exp(Cx^\eps)\) for large \(x\), by the
definitions of these ranks and orders. The parameter bound
in Lemma~\ref{lem:size97} and the analytic period estimate of
Theorem~\ref{thm:raw97} therefore apply at this same prime.
\end{proof}

The prime window here is recalculated using the replacement row.
The compact pairing of Theorem~\ref{thm:pair107} is explicitly
defined on its tensor lattice. To use it in \eqref{eq:gap} one
still needs a global automorphic realization with a nonzero
Bessel integral, and a comparison of that integral on the
arithmetic lattices with its analytic normalization. Its
height is precisely the issue in Corollary~\ref{cor:alg97}.
On this very field the complementary strict groups in
Proposition~\ref{prop:short107} still require a uniform bound.

\subsection{One terminal prime for a family of coded columns}\label{nav:VI-26}
Fix the field \(\widetilde E\), the ranks and the auxiliary primes
of Proposition~\ref{prop:short107} and Theorem~\ref{thm:levels88},
at one \(x\) on their common intervals. Write
\[
 b=\exp(x^\eps),\quad X'=\exp(x^{\eps/3}),\quad
 b'=xb/X',\quad w''=b'/(xn),\quad a=\lceil xb\rceil,\quad
 b_*=\lceil b'\rceil .
\]
Write \(r=k+m-1\) and \(t=t_{\rm r}\) below. Fix the \(h\) high
real places indexed by the order set of Theorem~\ref{thm:rangenew50}.
For each translate \(\ell\in A_0\) take its \(h_2\) distinguished
places \(L_\ell\) in the other half of the real-place labels, as
in Lemma~\ref{lem:av10}. A word consists of
\(n_0=k-2\) bits \(c_2,\ldots,c_{k-1}\), used at every place.
The Hamming distance is the number of bits at which two words
differ. We will select words before selecting the terminal prime.

We use the lists prescribed in
\srcloc{Development/AbelianVarietyTheory/Definitions/AutomorphicMembership/EvenCeil}{103}
and lines 171--191 there. Take any positive integers \(\Delta,H_0\)
bounded by fixed powers of \(x\), with
\(H_0\ge k+m(t+3)+r^{19}f^5x^\eps\). Such an \(H_0\) is
polynomially bounded by Theorem~\ref{thm:rangenew50}. At every place set
\begin{equation}
 v_i=4\Delta(k-1-i)\ (1\le i<k),\qquad
 u_i=v_{i-1}-\tfrac12-\Delta c_i\ (2\le i<k).
                                                               \label{eq:letters111}
\end{equation}
Set the pair \((u_1,u_k)\) equal to
\((a+(k-1)/2,-b_*-(k-1)/2)\) at a high place and to
\((v_1+4k\Delta+1/2,-1/2)\) elsewhere. Consecutive \(v_i\)
differ by \(4\Delta\) and \(a\) dominates every fixed power of
\(x\). These facts prove strict interlacing and give
\(u_1-v_1\ge4k\Delta\) at every place. Since \(b'=o(xb)\),
we also have \(u_1-u_k\le3a\).

At a high place put \(B_\tau^0=b_*+\lfloor b'/x\rfloor\), on
\(L_\ell\) put \(B_\tau^0=\lceil b\rceil\), and at the remaining
places put \(B_\tau^0=\lceil w''\rceil\). If their sum is odd
add one at the distinguished place indexed by the least element
of the untranslated distinguished set followed by translation by
\(\ell\). Otherwise make no adjustment. Call the result \(B_\tau\).
Then set
\[
 b_{j,\tau}=B_\tau+2(j-1)H_0,\qquad
 e_{j,\tau}=B_\tau+(2j-1)H_0\qquad(1\le j\le m).
\]
The count uses \(e_{j,\tau}\) itself by
\srcloc{Formalism/LocalToGlobal/Divisibility/Numbers}{160}.
The shifted representation weight is \(e_{j,\tau}-k/2-(2j-1)(t+1)/2\).
The successive gaps from \(b_j\) to \(e_j\) and \(e_j\) to
\(b_{j+1}\) are both \(H_0\). In particular they exceed
\((t+3)/2\). Both \(b'/x\) and \(w''\) grow faster than every
fixed power of \(x\), as follows by taking logarithms. Thus
\(u_k>-b_1\) also at high places, \(4k\Delta\le b_1/x\),
and \(H_0\le b_j/x\). Both the count weight and the shifted
weight lie above \(b_j\) and below \(b_{j+1}\) for \(j<m\),
and their last entries are at most \(2b\) eventually. These are
the weight inequalities of
\srcloc{Formalism/SharpRestriction/GlobalSchurBounds/StandardLattice}{496}.
The parity adjustment is independent of the word. At a fixed
translate the same signed row \eqref{eq:src109} therefore serves
the entire family. Every parameter has absolute value at most
\(\exp(Cx^\eps)\).

\begin{lemma}[Variation of the column dimensions]\label{lem:spread111}
For two arbitrary words \(c,c'\) in \eqref{eq:letters111}, at
each real place one has
\[
       \left|\log W_k(u_\tau(c))-\log W_k(u_\tau(c'))\right|\le Ck
\]
with an absolute constant. The same estimate holds for their
logarithmic Vandermonde products.
\end{lemma}
\begin{proof}
The denominator of the Weyl dimension formula is independent
of the word. Renumber the interior bits from 1 to \(n_0\) in
decreasing parameter order. Relative to the all-zero word, an
interior pair \(i<j\) changes its logarithmic factor by
\[
       \log\left(1+\frac{c_j-c_i}{4(j-i)}\right).
\]
The argument added to one has absolute value at most \(1/4\).
By Taylor's theorem replacing each logarithm by its linear
term costs at most \(C n_0\sum_{d\ge1}d^{-2}\) in total.
Four times the sum of the linear terms is
\[
        \sum_{i=1}^{n_0}c_i(H_{i-1}-H_{n_0-i}),\qquad
        H_j=\sum_{l=1}^j l^{-1},\quad H_0=0 .
\]
The coefficients are increasing and change sign once. The
maximum over binary words is attained by putting ones on the
positive coefficients, and the minimum by putting ones on
the negative coefficients. For either step function, at distance
\(d\) at most \(d\) pairs cross the step. The absolute contribution
of these pairs to four times the linear sum is at most one.
Summing over the distances bounds both extrema by \(n_0-1\).

In a pair with the top endpoint the gap is at least \(4k\Delta\)
and changes by at most \(\Delta\). The sum of the changes of
these logarithms is \(O(1)\). For the bottom endpoint the
smaller possible gap is at least \((4(k-i)-1)\Delta\) at interior
index \(i\) in the original numbering, since \(u_k\le-1/2\).
Its logarithm changes
by at most \(1/(4(k-i)-1)\), using \(\log(1+y)\le y\).
These terms sum to \(O(\log k)\). The pair of endpoints
is fixed. Adding the three estimates and comparing two words
through the zero word proves the lemma.
\end{proof}

\begin{theorem}[A common terminal prime for exponentially many words]
\label{thm:common111}
Put \(d_c=\mathfrak m_k I_\pi\prod_\tau W_k(u_\tau(c))\) and
\(D=\mathcal V(G)I_\sigma\prod_\tau W_r(\lambda_\tau)\).
Assign real numbers \(1\le\Gamma_c\le\exp(C f m^5 x^\eps)\)
to the words in advance, with fixed \(C\). Define
\[
     Y_c=\left(D/(d_c\Gamma_c^2)\right)^{1/(fk)},\qquad
     \delta_x=\exp(-x^{1/4}).
\]
For sufficiently large \(x\) there is a set \(\mathcal A\)
of at least \(\exp(k/5)\) words, all of the same weight
\(\lfloor n_0/2\rfloor\), with pairwise Hamming distance
at least \(n_0/8\). Writing \(Y=\max_{c\in\mathcal A}Y_c\), there
are \(Y^{1-o(1)}\) rational primes \(T=\varpi\bar\varpi\), with
\(\varpi\equiv1\bmod 2N(g_1g_2)^2\mathcal O_{K_d}\), such that
simultaneously
\[
            0<fk\log(T/Y_c)\le 3\delta_x/8
                         \qquad(c\in\mathcal A).
\]
In particular \(Y_c<T\le(1+1/(fk))Y_c\) for every selected
word. Their column representations have the same central
character at each real place.
\end{theorem}
\begin{proof}
Let \(p_c=\log(d_c\Gamma_c^2)\).
Lemma~\ref{lem:spread111} bounds the spread of \(\log d_c\)
by \(C_1fk\). From \(f\ll_\psi x^{2+8\eps}\),
\(m\asymp x^{1/2-\eps/16}\), \(k\asymp x^{9/2+10\eps}\) we
obtain
\[
            f m^5 x^\eps/k\ll_{\psi,\eps} x^{-21\eps/16}.
\]
Consequently the full range of the \(p_c\) is contained in
an interval of length \(O(fk)\). The layer of weight
\(\lfloor n_0/2\rfloor\) contains at least \(2^{n_0}/(n_0+1)\)
words, since its binomial coefficient is the largest. Subdivide
the interval for \(p_c\) into bins of width \(\delta_x/8\).
At most \(O(1+fk/\delta_x)\) bins meet this range.
One contains at least the size of the layer divided by this
number of bins.

From this bin choose words successively, discarding at each
choice all words at distance at most \(\lfloor n_0/8\rfloor\)
from the chosen one. Each choice discards at most
\(\exp(n_0 H(1/8))\) words, where
\(H(p)=-p\log p-(1-p)\log(1-p)\). Indeed, for \(p=1/8\)
and \(z=p/(1-p)\) the binomial expansion of \((1+z)^{n_0}\)
gives
\(\sum_{j\le n_0p}\binom{n_0}{j}\le z^{-n_0p}(1+z)^{n_0}\).
The logarithm of the number of bins is \(O(\log x+x^{1/4})=o(k)\).
Thus the logarithm of the selected cardinality is at least
\((\log 2-H(1/8)-o(1))k\), which exceeds \(k/5\).
Equal weights fix \(\sum_i u_i\) at every place. Subtracting
the same compact half sum of positive roots gives equal sums of
highest weights, hence equal scalar central characters.

The shared bin gives \(0\le\log(Y/Y_c)\le\delta_x/(8fk)\).
Corollary~\ref{cor:real91} applies to every list by the verified
span, interlacing and range of \(\Gamma_c\). In particular
\(\log Y\ge(16-o(1))m x^\eps\). Let \(M=2N(g_1g_2)^2\) and
\(\mathfrak B=4dM^2\). The proof of Theorem~\ref{thm:levels88}
gives \(\log\mathfrak B\ll_{\psi,\eps}x^\eps\), also for this
doubled modulus. Hence \(\log Y/\log\mathfrak B\gg m\) and
\(\sqrt{\log Y}\gg x^{1/4+15\eps/32}\). Each grows faster than
\(x^{1/4}+\log(fk)\) for fixed \(0<\eps<1/52\).
Thus the error \(\rho(Y,\mathfrak B)\) of Lemma~\ref{lem:ray88}
is \(o(\delta_x/(fk))\). Apply that lemma with relative width
\(\delta_x/(4fk)\). Its lower bound is \(Y^{1-o(1)}\) since
\(\log\mathfrak B+\log(fk)+x^{1/4}+\log\log Y=o(\log Y)\).
The trivial count of all integers in the interval supplies the
corresponding upper exponent. For any resulting prime,
\(0<\log(T/Y)\le\delta_x/(4fk)\). Adding the bound
on \(\log(Y/Y_c)\) gives the assertion. Finally
\(\log(1+1/(fk))\ge1/(2fk)\), which yields the stated
inclusion in each original window.
\end{proof}

\begin{corollary}[The actual refined column multiplicities]\label{cor:rank111}
At this same \(T\), let \(R_c\) be the multiplicity of the
column type after refinement by a nonzero residue covector at
each place above \(T\). Let \(R_c^\circ\) be the multiplicity
on the subspace with zero sum on each local scalar orbit, in
each coordinate separately. Uniformly over \(\mathcal A\),
\[
    \frac{\Gamma_c^2 R_c}{D}=1+O(\delta_x),
       \qquad \frac{\Gamma_c^2 R_c^\circ}{D}=1+O(\delta_x).
\]
For any fixed positive \(C_0,C_H\), the size of \(\mathcal A\)
exceeds \(C_0x^{C_H}\max_c\Gamma_c^2\) eventually.
\end{corollary}
\begin{proof}
By Proposition \ref{prop:column110}, \(d_c\) is the multiplicity at the
unrefined terminal level, including the finite levels at
\(g_1,g_2\). Its classes have trivial stabilizers. Since \(T\)
splits completely, every class has independently \(T^k-1\)
refinements at each of the \(f\) places. Each scalar orbit
contains \(T-1\) nonzero covectors, and restricting functions
on one such orbit to have zero sum leaves dimension \(T-2\).
Tensoring these local choices yields
\[
 R_c=d_c(T^k-1)^f,\qquad
 R_c^\circ=R_c(1-1/(T-1))^f.
\]
It follows from the definition of \(Y_c\) that
\(\Gamma_c^2R_c/D=(T/Y_c)^{fk}(1-T^{-k})^f\).
The first factor equals \(1+O(\delta_x)\) by
Theorem~\ref{thm:common111}. For \(0\le y\le1\),
\(0\le 1-(1-y)^f\le fy\). Apply this to both remaining
factors, using \(f/(T-1)=o(\delta_x)\) by the lower bound
on \(\log Y\). Finally \(\log(C_0 x^{C_H}\max_c\Gamma_c^2)=o(k)\)
by the rank calculation in the preceding proof. This proves
the size assertion requested numerically in
\srcloc{Formalism/SharpRestriction/GlobalSchurBounds/Chosen}{56}.
\end{proof}

\begin{corollary}[The same prime and the same count on every translate]
\label{cor:count111}
Use the preceding lists for every \(\ell\in A_0\), with the same
\(\Gamma_c\) for a given word. One set of words and one terminal
prime as in Theorem~\ref{thm:common111} give both rank estimates of
Corollary~\ref{cor:rank111} for every translate. For the weights
and orders of these very lists the centered prime-power count
\(\Delta_\ell\) of Section~\ref{sec:comparison} satisfies
\[
 \left|\frac{2\Delta_\ell}{hm b\delta}-\Gg_\psi(x)\right|
       \ll_{\psi,\eps}x^{1/2+\eps/2},\qquad \delta=h_2/n .
\]
\end{corollary}
\begin{proof}
Translating the distinguished set permutes its positions, preserves
its size and moves the marked position with it. Its positions
are disjoint from the high places. Thus the multiset of the
\(B_\tau\) stays fixed, including the possible parity adjustment.
The row product in \(D\) depends only on this multiset. The
column product in \(d_c\) depends on \(c\) and on the number
of high places. Hence every \(Y_c\), the bins of the proof of
Theorem~\ref{thm:common111}, and the common prime interval can
be taken identical for all translates.

The order at a high place is \(u_1-(k-1)/2=a\), precisely the
definition in \srcloc{Formalism/LocalToGlobal/Divisibility/Numbers}{247}.
At a distinguished place the count weight \(e_j\) lies between
\(b\) and \(b+2+2mH_0\le b+\sqrt b\) eventually. At every other
place it lies between one and \(2b'\) eventually. Indeed at a
high place it is at most \(b'+1+b'/x+2mH_0\), and elsewhere in
the complement it is at most \(w''+1+2mH_0\). Hence \(a/e_j\ge
X'/2\) there. Also \(|a/b-x|\le1/b\). Corollary~\ref{cor:actual50}
now gives exactly the displayed estimate on the same field and
interval, for each translate.
\end{proof}

Theorem~\ref{thm:src109} supplies the selected columns' individual
real Bessel models. On the unrefined terminal quotient the
smooth functions of Theorem \ref{thm:section110} realize both frequencies simultaneously.
Proposition~\ref{prop:short107} bounds the indicated invariant
cohomology summand on the same field. A rank comparison does
not evaluate the Bessel coefficient on the integral theta image.
That evaluation and the complementary cohomology estimate are
still needed for \eqref{eq:window}.

\subsection{One cyclotomic degree bound for all the complementary twists}\label{nav:degreebudget329}
Fix \(x\) on one of the intervals of Section~\ref{sec:missing}, and
keep its prescribed pair of auxiliary primes. Put \(n_*=[E_*:\QQ]\),
\(t=\lfloor x^{401}\rfloor\), and \(Y=x^{-3}\exp(x^\varepsilon/4)\).
For every odd prime \(p\le Y\) take the largest integer \(j_p\ge0\)
with \(p^{j_p+1}\le Y\). At this prime use the tower and the full
complementary dual \(A_{x,r}\) of Theorem~\ref{thm:strictdigits329}.
The prime is included as a subscript below when it varies.

\begin{theorem}[A uniform truncated sum including the decomposition invariants]\leavevmode
\label{thm:alltrunc329}
For \(0<\varepsilon<1/52\), uniformly at all sufficiently large
\(x\) on these intervals,
\begin{equation}\label{eq:alltrunc329}
 \sum_{\substack{2<p\le Y\\p\ {\rm prime}}}\;
 \sum_{1\le r<t}\log\bigl|
 A_{p,x,r}/p^{s_{p,r}+j_p}A_{p,x,r}\bigr|
 \ll_{\psi,\varepsilon} x^{-2}\exp(x^\varepsilon/2).
\end{equation}
This includes all characters of the tame cyclotomic action.
If at each prime one replaces the logarithmic class group at level
\(j_p\) by its image in the ordinary class group with \(p\) inverted,
the sum decreases by at most
\begin{equation}\label{eq:localdegree329}
 O\bigl(n_*\log(2n_*Yt)\,[Y+tn_*\log(2Y)]\bigr).
\end{equation}
For each prime and weight for which the norm-level test of
Theorem~\ref{thm:strictdigits329} holds at \(j_p\), the corresponding
summand in \eqref{eq:alltrunc329} is the full logarithmic order.
\end{theorem}

\begin{proof}
Fix \(p\), put \(j=j_p\), \(F=F_j\) and \(T=\widetilde C(F)\).
Use the integer \(u\) of Theorem~\ref{thm:strictdigits329}. A
single precision bound for all the indicated weights is
\[
 H=u+1+j+\lfloor\log(t)/\log(p)\rfloor,
 \qquad H\log p\le\log(n_*Yt).
\]
Indeed \(p^u\le n_*\), \(p^{1+j}\le Y\), and \(p^{v_p(r)}\le t\).
Let \(\Delta_0\) be the torsion subgroup of
\(\Gal(E_*(\mu_{p^\infty})/E_*)\), and set
\(d_p=|\Delta_0|=[E_*(\mu_p):E_*]\ge(p-1)/n_*\).
The torsion subgroup maps isomorphically to the reduction modulo
\(p\) of the cyclotomic image. This follows from
\(\ZZ_p^\times=\mu_{p-1}\times(1+p\ZZ_p)\). A closed subgroup
contains the torsion part of each element by taking powers
\(p^m\) with \(p^m\equiv1\bmod (p-1)\) and passing to the limit.

For \(\tau=\kappa_p^{-r}|_{\Delta_0}\), denote the usual character
idempotent by \(e_\tau\). It is integral since \(p\nmid d_p\) and
the character values belong to \(\ZZ_p^\times\).
The quotient in \eqref{eq:alltrunc329} is \(Q_{x,r}(T)\) by
\eqref{eq:strictdigits329}. It is a quotient of \(e_\tau(T/p^HT)\).
Each \(\tau\) occurs at most \(1+t/d_p\) times. Orthogonality and
the sum-to-one identity for the idempotents give
\[
 \sum_{1\le r<t}\log |Q_{x,r}(T)|
 \le (1+t/d_p)\log|T/p^HT|.
\]

Let \(C'\) denote the image of \(T\) in
\(\Cl(\mathcal O_F[1/p])\{p\}\) and let \(K\) be its kernel.
The local degree sequence of Theorem~\ref{thm:logunits319}
shows that \(K\) is generated over \(\ZZ_p\) by at most \(g-1\)
elements, where \(g\) is the number of primes of \(F\) over \(p\).
Indeed the local degree group at each such prime is a subgroup
of \(\ZZ_p\). The kernel of the sum of these \(g\) groups is
free of rank \(g-1\), and surjects onto \(K\). Reduction modulo
\(p^H\) is right exact, so
\[
 \log|T/p^HT|\le\log |\Cl(F)|+(g-1)H\log p.
\]
More precisely \(Q_{x,r}(T)\to Q_{x,r}(C')\) is onto, and its
kernel is the image of \(K\). To see this, lift the character
relations from \(C'\) to \(T\). The quotient of \(T\) by \(K\)
and these relations is exactly \(Q_{x,r}(C')\). The image of the
compact group \(K\) is closed. Only \(e_\tau(K/p^HK)\) can
contribute to that image. Summing over \(r\) as above bounds the
total change at \(p\) by \((1+t/d_p)(g-1)H\log p\).

There is a single prime above \(p\) in the real part of
\(\QQ(\mu_{p^{u+1+j}})\), since the cyclotomic field is totally
ramified at \(p\). Its real part is contained in \(F\) with
relative degree at most \(n_*\), by the compositum description
of \(F\) and the equality of the two CM degrees with their
real degrees over this cyclotomic base. Hence \(g\le n_*\).
Summing with \(1/d_p\le n_* /(p-1)\) and
\(\sum_{2<p\le Y}1/(p-1)\le1+\log Y\) proves
\eqref{eq:localdegree329}.

Finally \([F:\QQ]=p^j[F_0:\QQ]\le n_*p^{j+1}\le n_*Y\). If
\(N_*\) is the cyclotomic modulus of Theorem~\ref{thm:conductor232},
then \(F\) has modulus at most \(N_*p^{u+1+j}\le N_*n_*Y\).
The conductor--discriminant formula gives
\(\log|\operatorname{disc}(F)|\le [F:\QQ]\log(N_*n_*Y)\).
Minkowski's ideal bound and the estimate
\(\zeta_F(2)\le\zeta(2)^{[F:\QQ]}\), as in
Theorem~\ref{thm:normsum327}, now give
\[
 \log|\Cl(F)|\ll n_*Y\log(2N_*n_*Y).
\]
After summing over primes, this is multiplied by at most
\(Y+tn_*(1+\log Y)\). In the present family,
\(n_*\ll_\psi x^{3+12\varepsilon}\) and
\(\log(2N_*n_*Yt)\ll_{\psi,\varepsilon}x^\varepsilon\), by the
field and conductor bounds in Theorem~\ref{thm:normsum327}.
The term quadratic in \(Y\) is therefore
\(O_{\psi,\varepsilon}(x^{-3+13\varepsilon}\exp(x^\varepsilon/2))\).
The other terms, including \eqref{eq:localdegree329}, are bounded
by a fixed power of \(x\) times \(\exp(x^\varepsilon/4)\).
Both bounds imply \eqref{eq:alltrunc329}.
\end{proof}

At the primes in this range, subtracting the left side of
\eqref{eq:alltrunc329} from the sum of the full orders in logarithmic
scale gives
\[
 \sum_{\substack{2<p\le Y\\p\ {\rm prime}}}
 \sum_{1\le r<t}\log|p^{s_{p,r}+j_p}A_{p,x,r}|.
\]
A uniform bound for this excess, the contribution of the other primes
and the comparison of signed periods on the same fields remain to be
established in Section~\ref{sec:missing}.


\subsection{A rational flag translation with prescribed auxiliary residues}\label{sec:rationalflag342}
Write $E=E_x$, $F=E^+$, $n=[E:\QQ]$ and $\Delta_E=|\mathrm{disc}(E)|$.
Work with the rational hyperbolic pairs $e_i,f_i$ in the row space,
the marked line $z$ and the column space $A$. A rational flag
diagonal fixes $A,z$ and acts by $c_i\in E^\times$ on $e_i$ and by
$\bar c_i^{-1}$ on $f_i$. This is an element of the unitary group
$G(F)$ and centralizes the column group.

\begin{lemma}[An integral representative close to all real and complex targets]\leavevmode\label{lem:archlift342}
For a positive integer $Q$, each residue class in $\mathcal O_E/Q\mathcal O_E$
has a representative within distance $nQ2^{n+1}\sqrt{\Delta_E}$ of
any prescribed point of $E\otimes_\QQ\mathbb R$. Distance is the
maximum of the absolute values at the embeddings, with the usual
complex modulus.
\end{lemma}
\begin{proof}
Use the product of the unit intervals and unit complex disks in
Minkowski space. Every nonzero algebraic integer has maximum absolute
value at least one by its nonzero integral norm. The first and hence
all successive minima of this body on $\mathcal O_E$ are at least
one. Minkowski's second theorem bounds their product by $2^n$ times
the lattice covolume divided by the volume of the body. The covolume
is $2^{-r_2}\sqrt{\Delta_E}$ and the volume is $2^{r_1}\pi^{r_2}$,
where $r_1,r_2$ count real embeddings and conjugate pairs.
Consequently there are $n$ real-linearly independent integers of
maximum absolute value at most $2^{n+1}\sqrt{\Delta_E}$. Take any
representative of the given residue class. Express its displacement
from the desired point in the real basis formed by $Q$ times these
integers and round each coefficient to a nearest integer. Subtracting
the resulting integral linear combination preserves the residue
class and gives the claimed distance.
\end{proof}

\begin{theorem}[One rational diagonal at both auxiliary primes]\leavevmode\label{thm:rationalflag342}
Fix the signed auxiliary levels and terminal prime of
Theorems~\ref{thm:levels88} and~\ref{thm:common111}. Choose an
integer $b\ge1$ and $0<\delta<1/2$. There is a rational flag
diagonal $d_*$ with these properties.
\begin{enumerate}
\item At every auxiliary half its string multipliers, up to first
principal units, are those of Proposition~\ref{prop:centeredflag341}
after dividing them all by the unit on its middle line $z$. Thus
it gives precisely the same character on every integral string root
after conjugation.
\item At every half over $T$ the diagonal is the identity modulo
$T^b$ in the split flag coordinates. At all archimedean embeddings,
$|\sigma(c_i)-1|\le\delta$ and
$|\sigma(\bar c_i^{-1})-1|\le2\delta$.
\item With $Q=g_1g_2T^b$ one has, for each $i$,
\[
 \sum_{\mathfrak p}|v_{\mathfrak p}(c_i)|\log N\mathfrak p
 \ll n\bigl(s\log(g_1g_2)+\log Q+n+\log\Delta_E+\log(1/\delta)\bigr).
\]
The sum is over all finite primes of $E$, with integer-valued
valuations. For $b=1$ and $\delta=\exp(-x^B)$ with fixed $B>0$,
this bound is polynomial in $x$ on the same intervals. It holds
simultaneously for all the string lines.
\end{enumerate}
The prescribed finite and archimedean conditions also force a
displacement away from the marked primes. With $a_i=s+1-i$, every
such choice, for each $1\le i\le s$, satisfies
\[
 \sum_{\mathfrak p\nmid g_1g_2T}v_{\mathfrak p}(c_i)\log N\mathfrak p
 \ \ge\ n\log\bigl((g_1g_2)^{a_i}(1-\delta)\bigr)>0 .
\]
\end{theorem}
\begin{proof}
Put $a_i=s+1-i$ for $1\le i\le s$ and $D_i=(g_1g_2)^{a_i}$.
Normalize the local centered diagonal at each auxiliary place by
dividing its entire string by its middle eigenvalue. That eigenvalue
is a unit and the simple ratios are unchanged. At the chosen split
half prescribe $c_i$ by the normalized eigenvalue on $e_i$. At the
other half prescribe it by the conjugate of the inverse eigenvalue
on $f_i$ in the chosen half. Both prescriptions have valuation
$-a_i$. They use different primes of $E$, so they impose independent
conditions. Only the unit leading residues of these prescribed
values need to be matched. At the terminal halves prescribe
$c_i\equiv1\bmod T^b$.

Write $R=nQ2^{n+1}\sqrt{\Delta_E}$ and choose the least nonnegative
integer $k$ with $2^k\ge R/\delta$. The marked primes are odd.
Multiplication of the local prescriptions by $2^kD_i$ gives unit
residues at both auxiliary primes and gives the residue $2^kD_i$
at $T^b$. Since all marked primes split completely, the Chinese
remainder theorem specifies one class in $\mathcal O_E/Q\mathcal O_E$.
Apply Lemma~\ref{lem:archlift342} with the archimedean target
$2^kD_i$ in every embedding, and denote the resulting integer by
$u_i$. Then $c_i=u_i/(2^kD_i)$ has all the required finite residues
and $|\sigma(c_i)-1|\le R/(2^kD_i)\le\delta$. In particular $u_i$
is nonzero. The bound for the inverse conjugate follows by division
by an absolute value at least $1-\delta$.

The sum of absolute finite valuations, weighted by log norms, is
bounded by $\log|N_{E/\QQ}(u_i)|+n\log(2^kD_i)$. This follows by
the triangle inequality for the exponent of each ideal prime.
The archimedean bounds give $|N(u_i)|\le((1+\delta)2^kD_i)^n$.
Using the minimal choice of $k$ proves the estimate. The field
degree, $s$, $\log T$ and $\log(g_1g_2)$ are polynomially bounded
by the results cited in the statement. On these abelian fields
$\log\Delta_E\le n\log N_*$ by the conductor--discriminant formula.
Substitution gives the polynomial bound. The tuple of resulting
$c_i$ defines a single rational unitary element. For any admissible
choice the sum of $v_{\mathfrak p}(c_i)\log N\mathfrak p$ over the
auxiliary primes equals $-na_i\log(g_1g_2)$, since both rational
primes split completely. The terms over $T$ vanish. The sum over
all finite primes is $\log|N_{E/\QQ}(c_i)|$, which is at least
$n\log(1-\delta)$ by the archimedean condition. Subtracting the
marked terms proves the lower bound. It is positive because the
auxiliary primes are odd, $a_i\ge1$ and $\delta<1/2$.
\end{proof}


\subsection{Prescribing the splitting of the compensating primes}\label{sec:splitflag343}
Retain $E=E_x$, $F=E^+$ and $n=[E:\QQ]$ from
Theorem~\ref{thm:rationalflag342}. The field $E$ is Galois over
$\QQ$. For $c\in E^\times$ write
\[
                  h_f(c)=\sum_{\mathfrak p}|v_{\mathfrak p}(c)|
                                      \log N\mathfrak p .
\]
Here $v_{\mathfrak p}$ assigns one to a uniformizer. We shall realize
the prescribed flag with every additional nonzero valuation over a
rational prime splitting completely in $E$.

\begin{lemma}[Primes dividing a linear factor in a Galois field]\leavevmode\label{lem:splitfactor343}
Let $\theta$ be an integral primitive element of a finite Galois
extension $E/\QQ$, and let $f$ be its minimal polynomial. If a
rational prime $\ell\nmid\operatorname{disc}(f)$ divides a nonzero
integer $N_{E/\QQ}(a-\theta)$ with $a\in\ZZ$, then $\ell$ splits
completely in $E$.
\end{lemma}
\begin{proof}
There is a prime $\mathfrak l$ over $\ell$ at which $\theta$ has
residue $a$. The prime $\ell$ is unramified because the field
discriminant divides $\operatorname{disc}(f)$. All the conjugates of
$\theta$ have distinct residues at $\mathfrak l$, since the product
of their squared differences is a unit there. Arithmetic Frobenius
at $\mathfrak l$ carries the residue $a$ to $a^\ell=a$ in the
residue field. Its value on $\theta$ must therefore be $\theta$.
As $\theta$ generates $E$, Frobenius is the identity. In a Galois
extension this gives residue degree one at every prime over $\ell$.
\end{proof}

\begin{theorem}[A rational flag with all new valuations at split primes]\leavevmode\label{thm:splitflag343}
Fix the two auxiliary primes, the terminal prime $T$, the signed
characters and the precision $0<\delta<1/2$ of
Theorem~\ref{thm:rationalflag342}. Let $M$ be any positive integer
coprime to $g_1g_2$ and divisible by $T^b$ for some $b\ge1$.
On the same fields and intervals, for all sufficiently large $x$,
one can choose a single rational flag diagonal with the prescribed
leading residues at both auxiliary primes and
$|\sigma(c_i)-1|\le\delta$ at every embedding. At every
$\mathfrak p\mid M$ it satisfies $c_i\equiv1\bmod M\mathcal O_{E,\mathfrak p}$.
Every rational prime supporting a finite valuation of any $c_i$
splits completely in $E$. With
\[
 L=2+\log\Delta_E+n\log(2n)+\log(g_1g_2)+\log M+\log(1/\delta),
\]
one has, simultaneously in $1\le i\le s$,
\[
                   h_f(c_i)\ll n^2(n^2+s+2)L
\]
with an absolute implied constant. In particular one may avoid any
prescribed finite set of unmarked rational primes. When $\log M$ and
$\log(1/\delta)$ are bounded by fixed powers of $x$, the displayed
height bound is polynomial in $x$.
\end{theorem}
\begin{proof}
Theorem~\ref{thm:levels88} gives $\log g_j\asymp_\varepsilon
x^\varepsilon$, while $n$ is polynomially bounded, so $g_1,g_2>n$
for large $x$. At each $g=g_1,g_2$ prescribe distinct elements of
$\mathbb F_g$ at the $n$ primes of $E$ over $g$. At every rational
prime dividing $M$ or at most $n$, prescribe zero for $\theta$ modulo
the rational prime. These data specify a class in
$\mathcal O_E/Q_0\mathcal O_E$ by the Chinese remainder theorem,
where $Q_0$ is the product of the distinct rational primes involved.
We have $\log Q_0\le\log(g_1g_2)+\log M+n\log(2n)$.
Lemma~\ref{lem:archlift342}, with target zero, supplies an integral
representative $\theta$ with $\log A\ll L$, where
$A=\max(2,\max_\sigma|\sigma\theta|)$. Its monic minimal polynomial
$f$ has at least $n$ distinct roots modulo $g_1$, one at each
prescribed residue, and therefore has degree $n$. In particular
$\theta$ is primitive. The differences between conjugates give
\[
                     \log|\operatorname{disc}(f)|
                                \le n(n-1)\log(2A).
\]
At either auxiliary prime $f$ has distinct roots, so its
discriminant is prime to $g_1g_2$.

Fix $i$ and put $a_i=s+1-i$. For each $\mathfrak p$ above either
auxiliary prime $g$, let $w_{\mathfrak p}\in\mathbb F_g^\times$ be
the required residue of $g^{a_i}c_i$ in the construction of
Theorem~\ref{thm:rationalflag342}. We construct one factor
$(u-\theta)/(v-\theta)$ for this prime. In
$E_{\mathfrak p}=\QQ_g$ write $\theta_{\mathfrak p}$ for the image of
$\theta$. Prescribe the rational integers $u,v$ by
\[
 u\equiv\theta_{\mathfrak p}+g\widetilde w_{\mathfrak p}\pmod{g^2},
 \qquad
 v\equiv\theta_{\mathfrak p}+g^{a_i+1}\pmod{g^{a_i+2}},
\]
using any integral lift of $w_{\mathfrak p}$. The ratio has valuation
$-a_i$ and the specified leading residue at $\mathfrak p$. At every
other prime above the same $g$, the distinct residues of $\theta$
make both differences units and the ratio is one modulo $g$.
At the other auxiliary rational prime choose the same residue for
$u,v$ outside the $n$ roots of $f$. This makes every corresponding
ratio one modulo that prime.

Impose $u\equiv v\equiv1\bmod M$. For all primes dividing $M$, the
two differences are units since $\theta$ has residue zero. This gives
the stated congruence for the ratio, to the full precision of $M$.
For each remaining prime $\ell\mid\operatorname{disc}(f)$, choose
the same residue for $u,v$ outside the roots of $f$ modulo $\ell$.
If $\ell>n$ this is possible by degree. If $\ell\le n$, the choice
of $\theta$ gives $f(X)\equiv X^n\bmod\ell$, so residue one works.
All these congruences have rational integer solutions by the
Chinese remainder theorem. The moduli for both integers can be
replaced by one common multiple $J$ with
\[
   \log J\le (s+2)\log(g_1g_2)+\log M+\log|\operatorname{disc}(f)|.
\]
For example take the least common multiple of $M$,
$(g_1g_2)^{s+2}$ and the product of the primes dividing the
discriminant. At primes already covered by $M$ the specified
residue one applies.

Choose representatives $u_0,v_0\in[0,J)$. Add to both the least
multiple $H$ of $J$ which is at least $2A+8nJ/\delta$. In every
embedding the differences from $\theta$ are nonzero and
\[
 \left|\frac{u-\sigma\theta}{v-\sigma\theta}-1\right|
       \le \frac{J}{H-A}\le\frac{\delta}{8n}.
\]
No prime dividing $\operatorname{disc}(f)$ divides either norm.
Every prime that does divide either norm therefore splits completely
by Lemma~\ref{lem:splitfactor343}. The sum of the two logarithms of
the absolute norms is at most $2n\log(H+J+A)$. It bounds the
absolute finite valuation sum of the ratio because the numerator and
denominator are algebraic integers. The definitions of $H,J,A$ give
$\log(H+J+A)\ll(n^2+s+2)L$.

Now multiply these $2n$ factors to make $c_i$. At each auxiliary
prime of $E$ exactly one factor prescribes its leading residue and
all the others are one in its residue field. Multiplication preserves
the congruences at $M$. At each embedding the distance of the product
from one is at most $(1+\delta/(8n))^{2n}-1\le e^{\delta/4}-1\le\delta$.
Adding the height bounds proves the estimate. Take the multipliers
$c_i$ on the isotropic lines and $\bar c_i^{-1}$ on their duals as
in Theorem~\ref{thm:rationalflag342}. The original bounds on $n,s$,
$\log(g_1g_2)$ and $\log\Delta_E$ give the polynomial assertion.
\end{proof}


\subsection{An integral choice of all the unmarked root ratios}\label{sec:dominantflag344}
Put $g=g_1g_2$, set $c_{s+1}=1$, and write $\rho_j=c_j/c_{j+1}$
for $1\le j\le s$. Thus the positive half of the flag reads
$c_1,\ldots,c_s,1$. Across the other half the successive ratios are
the conjugates of these ratios in reverse order. At each auxiliary
prime of $E$ the matching condition prescribes valuation $-1$ and
a unit leading residue for every $\rho_j$. At unmarked places
the valuations of these ratios determine whether the Bessel
character can be trivial on all the original integral root groups.

\begin{theorem}[Simultaneous integral root ratios off the two auxiliary primes]\leavevmode\label{thm:dominantflag344}
Use the fields and auxiliary residue data of
Theorem~\ref{thm:splitflag343}. Let $s\ge1$, $0<\delta<1/2$, and
let $M$ be a positive multiple of $T^b$, coprime to $g$. One can
choose a rational unitary flag diagonal with the prescribed leading
residues at both auxiliary primes such that every $\rho_j$ is
integral away from $g$ and congruent to one modulo $M$. All primes
with nonzero valuations split completely in $E$. There is a
single real number $W>1$ for which
\[
      \left|\frac{\sigma(c_i)}{W^{s+1-i}}-1\right|\le\delta
                     \qquad(1\le i\le s,\ \sigma:E\hookrightarrow\CC).
\]
In terms of $n=[E:\QQ]$, put
\[
    L'=2+\log\Delta_E+sn\log(2sn)+\log(gM)+\log(1/\delta).
\]
Then $\log(2gW)\ll(1+sn^2)L'$ and
\[
           h_f(c_i)\le (s+1-i)n\log(2g^2W).
\]
In particular the set of new valuations can be required to lie
above the same cohomological cutoff $Y$ as in
Corollary~\ref{cor:avoidrange343}. With its choices $b=1$ and
$\delta=\exp(-x^B)$ and no further level congruences, both $\log W$
and $\sum_i h_f(c_i)$ are at most $\exp((1/4+o(1))x^\varepsilon)$.
\end{theorem}
\begin{proof}
For each $j$ and each $\mathfrak p\mid g$ let $w_{j,\mathfrak p}$
be the specified residue of $g\rho_j$. These data come by taking
ratios of the desired diagonal entries on the positive lines at
each prime of $E$. Products of the successive ratios recover
exactly the prescribed leading residues of all $c_i$ since the
middle line has multiplier one.

First construct a primitive algebraic integer $\omega$ with
$\log\max(2,\max_\sigma|\sigma\omega|)\ll L'$. At the $n$ primes
over $g_1$ prescribe distinct residues, which is possible because
$g_1>n$. Apply Lemma~\ref{lem:archlift342} modulo $g_1$ with
archimedean target zero. Its height bound gives the assertion,
and the degree is $n$ because the minimal polynomial has $n$
different roots modulo $g_1$.

For each $j$ prescribe $\theta_j\equiv-w_{j,\mathfrak p}$ at each
$\mathfrak p\mid g$, $\theta_j\equiv-g\bmod M$ and
$\theta_j\equiv0\bmod\ell$ for every rational prime
$\ell\le sn$ not dividing $gM$. All these congruences specify
one class modulo the integer
\[
                     Q=gM\prod_{\ell\le sn,\ \ell\nmid gM}\ell ,
\]
where the product is over primes. By
Lemma~\ref{lem:archlift342}, choose an integral representative
$\theta_j^0$ with logarithmic maximum norm $O(L')$, using
$\log Q\le\log(gM)+sn\log(2sn)$. Among the integers
$k=0,\ldots,n(n-1)/2$ choose one for which
$\theta_j=\theta_j^0+kQ\omega$ has distinct conjugates. Each
pair of embeddings forbids at most one $k$ because their values
on $\omega$ differ. There are more available integers than
pairs. The congruences are preserved and
$A=\max(2,\max_{j,\sigma}|\sigma\theta_j|)$ still satisfies
$\log A\ll L'$. Write $f_j$ for the monic minimal polynomial of
$\theta_j$. It has degree $n$ and
$\log|\operatorname{disc}(f_j)|\le n(n-1)\log(2A)$.

We choose one rational integer $a$ for all the numerators. Impose
$a\equiv0\bmod gM$. For primes there, all $\theta_j$ are units
at every place above the prime, so no norm $N(a-\theta_j)$ can be
divisible by it. At every other prime $\ell$ dividing a
$\operatorname{disc}(f_j)$, exclude the roots of all the $f_j$
modulo $\ell$. If $\ell>sn$, at most $sn$ residues are excluded.
If $\ell\le sn$, all the $\theta_j$ are divisible by $\ell$ and
their polynomials are $X^n\bmod\ell$, so residue one is available.
The Chinese remainder theorem gives a residue class for $a$ modulo
an integer $J$ with
\[
                \log J\le\log(gM)+sn(n-1)\log(2A).
\]
Let $\eta=\delta/(4s)$ and take the first positive $a$ in this
class exceeding $\max(2g,A/\eta)$. It is at most this maximum
plus $J$. In particular $W=a/g>1$ satisfies the stated logarithmic
bound.

Set $\rho_j=(a-\theta_j)/g$. At $\mathfrak p\mid g$ its valuation
is $-1$ and $g\rho_j$ reduces to $w_{j,\mathfrak p}$. At primes
over $M$ the numerator is $g\bmod M$, giving the full congruence.
At all other primes the denominator is a unit. A prime dividing
any numerator norm divides none of the polynomial discriminants
by the choice of $a$. It splits completely by
Lemma~\ref{lem:splitfactor343}. At an embedding,
$|\sigma(\rho_j)/W-1|\le A/a\le\eta$. Now put
$c_i=\prod_{j=i}^s\rho_j$. The error relative to $W^{s+1-i}$ is
at most $(1+\eta)^s-1\le e^{\delta/4}-1\le\delta$.
These coordinates and their inverse conjugates define the desired
unitary diagonal.

The negative valuations of each $\rho_j$ contribute exactly
$n\log g$ to its absolute finite valuation sum. Its signed sum
is $\log|N(\rho_j)|\le n\log(2W)$ by the archimedean estimate.
Adding twice the negative contribution gives
$h_f(\rho_j)\le n\log(2g^2W)$. Subadditivity proves the height
bound for $c_i$. To avoid primes through $Y$, take $M$ to be $T$
times the product of all primes up to $Y$ outside $gT$. The same
estimate $\log M\le\log T+Y\log(2Y)$ as in
Corollary~\ref{cor:avoidrange343}, together with the polynomial
bounds on $n,s$, gives both final bounds.
\end{proof}



\subsection{The centered prime sum for the actual period weights}
\label{sec:centeredwordperiods451}

Keep the \(K_d\)-oriented field, translate \(u\), and ordinary
word weights of Corollary~\ref{cor:count111}. Write
\(h=|T_{\rm hi}|\), \(\delta=h_2/n\), \(a=\lceil xb\rceil\),
and \(X'=\exp(x^{\varepsilon/3})\). Denote by \(\xi_q,\xi_d\)
the primitive quadratic characters of \(K_q,K_d\), and by
\(\xi_{qd}\) the primitive character of their product. We use
\(D_\chi=L'(1,\chi)/L(1,\chi)\). The centered count
\(\Delta_u\) is the sum of \(C_d(t,a,e_i)-M_d(a,e_i)\)
over all \(t\in T_{\rm hi}\) and all \(m\) edges.
Define the explicit real number
\[
 A_{\rm sig}(x)=\frac{hma}{2}
 \big\{\log x+(1-\delta)(\log n+x^{\varepsilon/3})
     -1-\gamma-D_{\xi_q}+D_{\xi_d}+D_{\xi_{qd}}\big\}.
\]

\begin{theorem}[The complementary quadratic sign for the word weights]\leavevmode
\label{thm:centerwords451}
At every \(x\) on the common intervals and uniformly over
all translates, one has, for every fixed \(A>0\) as \(x\to\infty\),
\[
 \Delta_u=
  2\!\!\sum_{\substack{p\ge5\ {\rm unramified\ in}\ E_x\\
             p\ {\rm split\ in}\ K_q,\ {\rm inert\ in}\ K_d}}
                     U_p\log p
       -A_{\rm sig}(x)+o_{\psi,\varepsilon,A}(b/x^A).
\]
Here \(U_p\) is the sum for this translate in
Section~\ref{sec:signedsignal450}. The prime sum has no term
with \(p>2nxX'\).
\end{theorem}
\begin{proof}
At a fixed edge write \(w_\alpha=e_{i,\alpha}\).
Let \(C_{--}\) be the part of the count from group elements
nontrivial on both quadratic fields. This definition uses the
inertia coset average also at ramified primes. For a fixed high
\(t\), the map \(g\mapsto t+c_E+g\) carries this coset bijectively
onto \(\Phi_0\). The sign assertion holds for either parity of the
original \(\psi\), since complex conjugation reverses both quadratic
signs. Apply character orthogonality \eqref{eq:frob} before summing
over prime powers. As in the proof of Proposition~\ref{prop:count},
the part of \(C_{--}\) with a given \(g,\xi\) is
\[
 (2n)^{-1}\overline{\xi(g)}\, w_\alpha
                       R_{\xi^*}(a/w_\alpha),
               \qquad \alpha=t+c_E+g.
\]
Here \(R\) is the uncentered ramp sum in \eqref{eq:Gchi}.
For \(\alpha\in\Phi_0\) we have \(w_\alpha\le2b'\) eventually,
as \(b'=xb/X'\) dominates the rank shifts and the largest
weight there is at most \(\lceil b'\rceil+\lfloor b'/x\rfloor
+1+2mH_0\). Thus \(a/w_\alpha\ge X'/2\). Lemma~\ref{lem:far14}
applies to every character of this field, including those
of a shortened cyclotomic layer by inflation to the larger field.
After replacing each \(R\) by its main term, the error at this
edge and \(t\) is \(O_{\psi,\varepsilon}(na\exp(-c x^{\varepsilon/8}))\),
since there are \(n\) group elements in the coset and \(4n\)
characters.

For the principal character the main term at \(\alpha\) before
division by \(2n\) is \(a(\log a-\log w_\alpha-1-\gamma)\).
For a nonprincipal character it is \(-aD_{\xi^*}\) times the
character coefficient. A nontrivial character on the subgroup
fixing both quadratic fields sums to zero on this coset. The
three remaining nonprincipal characters have respective values
\(-1,-1,1\). Consequently the main term of \(C_{--}\) is
\[
 \frac a2\left\{\log a-1-\gamma-\frac1n
   \sum_{\alpha\in\Phi_0}\log w_\alpha+
   D_{\xi_q}+D_{\xi_d}-D_{\xi_{qd}}\right\}.
\]
Subtract this from \(M_d(a,w)\) in Proposition~\ref{prop:count}.
The result is
\[
 \frac a2\left\{\log a-1-\gamma-\frac1n
   \sum_{\alpha\in\Phi_1}\log w_\alpha-
   D_{\xi_q}+D_{\xi_d}+D_{\xi_{qd}}\right\}.
\]
On \(\Phi_1\) the distinguished \(h_2\) weights differ from \(b\)
by at most \(2+2mH_0\); the other weights differ from
\(w''=b/(nX')\) by at most this amount. As they are all at
least their respective reference values, the mean of their logs
differs from \(\delta\log b+(1-\delta)\log w''\) by at most
\((2+2mH_0)/w''\), using \(\log(1+v)\le v\).
Also \(\log a=\log x+\log b+O(1/(xb))\). Multiplication by
\(hma\) makes these rounding errors \(O(x^B X')\) for a fixed
\(B\); the error from the character sums is
\(O(hmna\exp(-c x^{\varepsilon/8}))\). Both estimates are
\(o(b/x^A)\). This gives \(A_{\rm sig}(x)\) for the total
complementary main term.

The remaining coset consists of elements split in \(K_q\) and
inert in \(K_d\). At an unramified prime an even Frobenius power
has trivial quadratic signs, while an odd power retains the
signs of that prime. Its exact contribution is therefore the
displayed \(2U_p\log p\) by \eqref{eq:cp73}. The primes \(2,3\)
and the fixed conductor primes split in \(K_d\) in the chosen
family, as in the proof of Theorem~\ref{thm:signal77}, so have
no contribution in this coset. The other ramified primes can
only be the two auxiliary quadratic conductor primes. For each
of them, the inertia average is at most one after removing the
factor two and the sum of \(a/p^j\) is \(a/(p-1)\).
Their total contribution is \(o(b/x^A)\) since those primes
grow faster than every power of \(x\) and \(hma/b\) is polynomial
in \(x\). Finally \(e_{i,\alpha}\ge b/(nX')\) on \(\Phi_1\),
so every active prime power is at most \(2nxX'\).
\end{proof}

To describe the finite denominator term explicitly, keep the
\(q\) independent columns and the set of good signal primes
\(\mathcal S\) from Corollary~\ref{cor:gramsignal450}.
For each \(D'=S_0,eS\), choose finite generators of the
integral test module \(D'\) over \(\mathcal O_J\). They exist
because the class coefficient lattices are finitely generated
and \(\mathcal O_J\) is Noetherian. They can be obtained on the
saturated space by intersecting the integral lattice with the
kernel of \(1-e\), and on the image space by applying \(e\)
to generators of \(S\).

\begin{proposition}[An input bound for the outside-prime loss]\leavevmode
\label{prop:outsidewordcost451}
For each column \(i\) and generator of \(D'\), expand its signed
period by Theorem~\ref{thm:covectorpaths352}, using independent
symbols for the undivided jets of every edge and class coset.
Write \(\beta\) for a nonzero scalar coefficient of a term, and
\(J_1,\ldots,J_m\) for its jet symbols, one per edge. Set
\[
 H^{D'}_{i,w}=
 \max\left(\{0\}\cup
 \left\{-v_w(\beta)+\sum_{j=1}^m\Gamma_w(J_j):
             \text{terms for these generators}\right\}\right),
\]
where \(\Gamma_w(J_j)\) is the nonnegative finite bound of
Theorem~\ref{thm:classden312} for that symbol's section, coordinates
and order. A conjugated symbol uses the corresponding conjugate
place. Define
\[
 B_{\rm data}=
   \sum_{p\notin\mathcal S}\log p\sum_{w\mid p}\omega_w
      \sum_{i=1}^q(H^{S_0}_{i,w}+H^{eS}_{i,w}).
\]
Then \(B_{\rm out}\le B_{\rm data}\), and \(B_{\rm data}\) is
finite and determined before evaluating any theta derivative.
\end{proposition}
\begin{proof}
In fixed rational class frames expand the successive determinant
twists and substitutions in \eqref{eq:sub81}. Replacing each
monomial in its matrix of differential symbols by a new variable
gives a polynomial linear in the variables of that edge.
On composing the \(m\) maps, every resulting term contains one
such variable for each edge. The finite path sum, normal-coordinate
factor, restriction, and small class contraction act on these
expressions by algebraic linear maps. They determine the coefficients
\(\beta\) in the marked arithmetic normalization and include the
actual input and test values and the mass factors. The local
estimates for the specified CM sections bound the valuation of
each evaluated term below by \(v_w(\beta)-\sum_j\Gamma_w(J_j)\).
Thus the valuation of the period on any generator is at least
\(-H^{D'}_{i,w}\). This extends by linearity to any integral
linear combination of generators and after localization at \(w\).
Expanding a determinant bounds \(b_{D',w}\) below by
\(-\sum_i H^{D'}_{i,w}\). Taking the negative parts on the
outside primes proves the inequality. In
Theorem~\ref{thm:classden312} one can use the common integer
denominators in Theorem~\ref{thm:periodden312} for every edge.
They depend only on the finite translations, linear substitutions,
fixed elliptic model and degrees. Those denominators and the
finitely many nonzero coefficients \(\beta\) involve only finitely
many negative valuations. This proves finiteness.
\end{proof}

\begin{corollary}[An upper bound for the centered count]\leavevmode
\label{cor:centerheight451}
In the hypotheses of Corollary~\ref{cor:gramsignal450}, let
\(E_{\rm omit}=2\sum_{p\notin\mathcal S}U_p\log p\), where
the sum is restricted to the unramified signal primes \(p\ge5\)
in Theorem~\ref{thm:centerwords451}. Then on the same fields
and intervals
\[
 \begin{split}
 \Delta_u\ \le{}&
   \frac1{q[J:\mathbb Q]}\sum_\sigma\log\det G_\sigma
       -A_{\rm sig}(x)\\
 &+\frac{2L_{\mathcal S}+B_{\rm data}}q+
          2\log D+\log m_G+E_{\rm omit}
          +o_{\psi,\varepsilon,A}(b/x^A).
 \end{split}
\]
If \(\mathcal S\) omits an active prime, its contribution to
\(E_{\rm omit}\) is at most \(2hma\log p/(p-1)\).
\end{corollary}
\begin{proof}
Split the prime sum of Theorem~\ref{thm:centerwords451}
at \(\mathcal S\), insert Corollary~\ref{cor:gramsignal450}
on that set, and apply Proposition~\ref{prop:outsidewordcost451}
to the outside denominators. In the omitted sum discard the
nonnegative thresholds and bound the geometric series by \(1/(p-1)\).
\end{proof}

A proof of \eqref{eq:gap} by this inequality requires projected
columns of adequate rank and a uniform upper estimate for the
displayed excess over \(A_{\rm sig}(x)\), with the finite
costs included. The explicit main term and the signs have been
computed here at the original word weights.


\subsection{The ramified prime in the polynomial-conductor comparison}
\label{sec:ramifiedcenter460}

Choose the polynomial-conductor intervals of
Theorem~\ref{thm:conductor232}. At their point \(x\), write \(j_x\)
for the dyadic index in that theorem. The field and its character
modulus are
\[
 E_x=K_\psi C_{j_x}(\sqrt{-q},\sqrt{-d_0}),\qquad
 N_x=2^{j_x+2}Fqd_0\le x^{D_0},
\]
where \(D_0\) is fixed. We use the \(K_d=\QQ(\sqrt{-d_0})\)
orientation and the actual word weights of
Corollary~\ref{cor:count111}. In the construction of
Theorem~\ref{thm:conductor232}, both \(q,d_0\) exceed \(x^{401}\)
throughout these intervals. Indeed, in its notation they exceed
\(H^{401W}\), and \(x\le V^{3/2}\le H^W\). This lower bound
has a fixed exponent. We evaluate their contribution explicitly
before estimating it.

Use additive notation in \(G=\operatorname{Gal}(E_x/\QQ)\).
Write \(c_E\) for complex conjugation, \(I_p\) for inertia and
\(F_p\) for a lift of arithmetic Frobenius modulo inertia.
Let \(\iota_d\) be the nonidentity element of \(I_{d_0}\).
The field \(K_\psi C_{j_x}(\sqrt{-q})\) is unramified at \(d_0\),
and \(K_d\) has ramification index two there. Thus
\(I_{d_0}=\{0,\iota_d\}\) and \(\iota_d\) acts trivially on the
first field. For a translate \(u\), put
\[
 W_{d_0}(u)=(\log d_0)\sum_{i=1}^m\sum_{t\in T_{\rm hi}}
       \sum_{k\ge1}
       [\,a/d_0^{2k}-e_{i,t+c_E+2kF_{d_0}+\iota_d}\,]_+ .
\]
The weights in this formula are those for \(u\). Even multiples
of \(F_{d_0}\) are independent of the lift.

\begin{lemma}[The ramified terms in the selected quadratic signs]\leavevmode
\label{lem:ramcenter460}
In the portion of the prime count with trivial sign on \(K_q\)
and nontrivial sign on \(K_d\), the contribution from ramified
primes is exactly \(W_{d_0}(u)\). In particular
\[
           0\le W_{d_0}(u)
              \le \frac{hma\log d_0}{d_0^2-1}.
\]
\end{lemma}
\begin{proof}
The character of \(K_d\) is unramified and split at \(2\) because
\(-d_0\equiv1\bmod8\). For odd primes of the fixed conductor
\(F\), and for \(q\), the congruences of
Theorem~\ref{thm:long232} give \(-d_0\equiv1\) at the prime
for sufficiently large intervals. Those primes have trivial
\(K_d\)-sign in every inertia coset. At \(d_0\) the character
of \(K_q\) is unramified with value
\((d_0/q)=(-1/q)=-1\). Here the primitive quadratic character
of discriminant \(-q\) is the Legendre character modulo \(q\).
The trivial \(K_q\)-sign therefore selects the even powers
\(d_0^{2k}\). In the coset \(2kF_{d_0}+I_{d_0}\), the even
multiple has trivial sign on both quadratic fields. Its sum
with \(\iota_d\) is the unique member with the required signs.
Formula~\eqref{eq:frob} gives it coefficient
\(2/|I_{d_0}|=1\). The weight at \(t\) and this group element
is the one displayed above. This exhausts the ramified primes.
Discard the nonnegative thresholds and sum the geometric series
in \(d_0^{-2}\) to obtain the upper bound.
\end{proof}

\begin{theorem}[The centered identity on the polynomial-conductor family]\leavevmode
\label{thm:polycenter460}
On these fields, uniformly over the translates and for each
fixed \(A>0\), the centered count satisfies
\[
 \Delta_u
   = 2\!\!\sum_{\substack{p\ge5\ {\rm unramified\ in}\ E_x\\
               p\ {\rm split\ in}\ K_q,\ {\rm inert\ in}\ K_d}}
                U_p\log p-A_{\rm sig}(x)+W_{d_0}(u)
                         +o_{\psi,\varepsilon,A}(b/x^A).
\]
For the sparse sets of Theorem~\ref{thm:rangenew50},
\[
 \frac1{|A_0|}\sum_{u\in A_0} W_{d_0}(u)
          =(1+o_{\psi,\varepsilon}(1))
                          \frac{hma\log d_0}{d_0^2}.
\]
\end{theorem}
\begin{proof}
For primes unramified in \(E_x\) the exact count in the
selected signs is \(2U_p\log p\) by
\eqref{eq:cp73}. The primes \(2,3\) split in \(K_d\) on these
intervals. In the other signs contributing to the count, the
group element acts nontrivially on both quadratic fields. Its
target \(t+c_E+g\) lies in \(\Phi_0\). At every such target
\(a/e_{i,t+c_E+g}\ge X'/2\) for large \(x\) by the word
weights. The character expansion uses primitive characters and
the full inertia average \eqref{eq:frob}, also at \(q,d_0\).

We check the large-argument error on the chosen family.
Every real primitive character of \(E_x\) without an auxiliary
quadratic factor has conductor dividing \(8F\), so these form
a fixed finite set. With such a factor, the zero-free rectangle
of Theorem~\ref{thm:long232} rules out the possible exceptional
real zero near one. The zero-free region and zero count used in
Lemma~\ref{lem:far14} therefore apply with
\(\log(2N_x)=O_{\psi,\varepsilon}(\log x)\). In detail, at
\(y\ge X'/2\), truncation at height \(\exp(\sqrt{\log y})\)
bounds the normalized contribution of the low zeros by a
constant times
\[
 \log^3(2N_x)\,
 \exp\!\left(-c\min\left\{\sqrt{\log y},
                       \frac{\log y}{\log(2N_x)}\right\}\right).
\]
The tail is \(O(\log(2N_x(3+\exp(\sqrt{\log y})))/
\exp(\sqrt{\log y}))\), using the \(1/|\rho(1-\rho)|\)
weights as in that lemma. The elementary terms of the explicit
formula obey the same required estimate after division by \(y\).
Since \(\log y\ge x^{\varepsilon/3}-\log2\), the total normalized
error is \(O(\exp(-c' x^{\varepsilon/8}))\). Multiplication by
the word weight and summation in the character expansion give
error \(O(hmna\exp(-c'x^{\varepsilon/8}))\). The quantities
\(h,m,n,a/b\) have polynomial growth in \(x\), so this is
\(o(b/x^A)\).

Character orthogonality and the logarithms of the word weights
now give the same complementary main term as the computation
in Theorem~\ref{thm:centerwords451}. More explicitly, subtracting
the part with both signs nontrivial from the main term of one
edge and one high place gives
\[
 \frac{a}{2}\left(\log a-1-\gamma
     -\frac1n\sum_{\mu\in\Phi_1}\log e_{i,\mu}
     -D_{\xi_q}+D_{\xi_d}+D_{\xi_{qd}}\right).
\]
The replacement of the \(h_2\) distinguished weights by \(b\)
and the other weights by \(b/(nX')\) in these logarithms has
total error \(O(x^B X')=o(b/x^A)\) for a fixed \(B\).
Summing thus gives \(A_{\rm sig}(x)\). Including the ramified
part of the selected signs by Lemma~\ref{lem:ramcenter460}
proves the identity.

For the mean, keep \(k=1\) in \(W_{d_0}\). For each \(i,t\) the
target in \(\Phi_1\) is fixed as \(u\) varies. The proportion
of translates for which it is distinguished is at most
\(|T_2|/|A_0|\), by translating \(T_2\) in \(A_0\). This is
\(o(1)\) for the sets in question. Indeed, their distinct
integer representatives in the proof of
Theorem~\ref{thm:rangenew50} lie in
\([0,|A_0|/(8L))\), where \(L\) tends to infinity. At an
undistinguished target the weight equals
\(b/(nX')+O(1+2mH_0)=o(a/d_0^2)\) uniformly. This follows
from \(d_0\le x^{D_0}\), polynomial growth of \(mH_0\), and
\(X'=\exp(x^{\varepsilon/3})\). Consequently the \(k=1\)
term averaged over translates is at least
\((1-o(1))hma\log d_0/d_0^2\). The lemma gives the matching
upper bound.
\end{proof}

\begin{corollary}[The comparison and the cohomology on the same shortened field]\leavevmode
\label{cor:shortcenter460}
For the indicated polynomial-conductor fields use
\(-A_{\rm sig}(x)+W_{d_0}(u)\) in the centered upper bound of
Corollary~\ref{cor:allcost454}, including the refinements of
Corollaries~\ref{cor:terminalrow455} and~\ref{cor:splitcost456}.
The correction is \(O_{\psi,\varepsilon}(hmb\,x^{-801}\log x)\).
Corollaries~\ref{cor:dyadicdepth458} and~\ref{cor:dyadicorbit459}
hold there for the faithful characters of \(C_{j_x}\), with
\(f=2^{j_x-2}\).
\end{corollary}
\begin{proof}
The inequalities for the integral minors at unramified primes
use the same signed columns, row places, and common terminal
prime on the fields of Theorem~\ref{thm:conductor232}. At the
ramified coefficient places for the periods, the lower bound
from the different and the marked differential in
Theorem~\ref{thm:outsidedifferent453} already applies. It does
not estimate the ramified term of the prime count. Substitute
the identity just proved when converting the two period ideals
to that centered count. This adds precisely \(W_{d_0}(u)\).
The numerical upper bound follows from \(a\le2xb\),
\(d_0>x^{401}\), and \(\log d_0\ll\log x\). In particular it
is \(o(x^{5\varepsilon}\mathcal B)\) in \eqref{eq:gap}, since
\(\mathcal B\ge2nm^2b x^\varepsilon\) and \(h\le n\).
The asymptotic for its mean also shows why an error
\(o(b/x^A)\) for every fixed \(A\) requires retaining this
term on the polynomial family.

Finally write \(o=[K_\psi:\QQ]\). The inequalities for the
degree of \(K_\psi C_{j_x}\) give
\[
     x^{2+8\varepsilon}/o\ \le\ 2^{j_x}
                             \ <\ 2x^{2+8\varepsilon}.
\]
Thus \(j_x\) tends to infinity, the faithful conductor is
\(2^{j_x+2}<8x^{2+8\varepsilon}\), and the odd part of
\([E_x:\QQ]\) divides \(o\). These are exactly the degree and
conductor inputs of Corollary~\ref{cor:dyadicdepth458}.
The order computation at \(\ell\equiv5\bmod8\) gives
\(f=2^{j_x-2}\) with this index as well. The simultaneous
class group bound and the two-orbit argument of
Corollary~\ref{cor:dyadicorbit459} now hold for these characters
on the very same \(E_x\).
\end{proof}


\subsection{A sparse set supporting the residual cohomology on the common fields}
\label{sec:shiftedsupport465}

We use the polynomial-conductor fields \(E_x\) of
Theorem~\ref{thm:conductor232} on one interval \([V,V^{3/2}]\).
In particular the two auxiliary primes are fixed across this
interval. Keep \(P(x),b_0(x),R_{\ell,r}\) and the logarithmic
length \(\mathcal R(x)\) of Section~\ref{sec:missing}. For
\(x\) in the interval let \(t(x)=\lfloor x^{401}\rfloor\), with
extreme values \(t_-=\lfloor V^{401}\rfloor\) and
\(t_*=\lfloor V^{1203/2}\rfloor\). Define
\[
 \delta=1-\frac{1+\log\log 2}{\log 2},\qquad
 H(T)=(\log T)^\delta\sqrt{\log\log T}.
\]
All the arguments of \(H\) below tend to infinity. For large
coefficient primes the two possible inertia residues in
Corollary~\ref{cor:unboundedprimes428} impose the same divisibility
condition on \(2r\). We count a weight once even when several primes
give that condition.

\begin{lemma}[Weights divisible after doubling by a large shifted prime]\leavevmode
\label{lem:shiftdivisors465}
Put \(Y=V^{100}\). Let \(\mathcal A_V\) consist of the positive
integers \(r\le t_*\) for which \(\ell-1\mid 2r\) for some prime
\(\ell>Y\). Uniformly for \(t_-\le T\le t_*\),
\[
                  |\mathcal A_V\cap[1,T]|\asymp T/H(T).
\]
\end{lemma}
\begin{proof}
Let \(N(Z,y)\) count positive integers \(n\le Z\) for which some
prime \(p>y\) satisfies \(p-1\mid n\). Ford's theorem gives
\[
  N(Z,y)\asymp_c
       \frac{Z}{(\log y)^\delta\sqrt{\log\log y}}
         \qquad(3\le y\le Z^{1-c})
\]
for each fixed \(0<c<1\) \cite[Theorem~1(i)]{I-I-ford465}.
An integer counted by \(N(T,Y)\) belongs to \(\mathcal A_V\),
since any divisor of \(r\) divides \(2r\). In the other direction
\(r\mapsto 2r\) injects \(\mathcal A_V\cap[1,T]\) into the
integers counted by \(N(2T,Y)\). Both applications allow
\(c=1/2\) because \(V^{100}\le \sqrt{t_-}\) for large \(V\).
Moreover \(\log Y/\log T\) is bounded above and below by
positive constants independent of \(T,V\) in this range. This
gives both inequalities. The lower bound counts the weights
allowed by the divisibility condition alone.
\end{proof}

\begin{theorem}[One set of weights for all residual groups up to the cutoff]\leavevmode
\label{thm:simultaneousweights465}
For every sufficiently distant interval of the stated
polynomial-conductor family there is a single set
\(\mathcal E_V\subseteq[1,t_*]\cap\ZZ\) such that
\[
 \begin{aligned}
  |\mathcal E_V\cap[1,T]|&\ll_{\psi,\varepsilon}T/H(T)
                        &&(t_-\le T\le t_*),\\
  R_{\ell,r}(E_x)&=0
                        &&\bigl(V\le x\le V^{3/2},\
             1\le r<t(x),\ r\notin\mathcal E_V,\
             2<\ell\le P(x)\text{ prime}\bigr).
 \end{aligned}
\]
In particular the fraction of weights which can support any such
group tends to zero uniformly in \(x\).
\end{theorem}
\begin{proof}
Use \(B=2\) in Corollary~\ref{cor:cut232}. It provides one set
\(\mathcal D_V\) for the whole interval with
\[
  |\mathcal D_V|\ll_{\psi,\varepsilon}
               V^{401}/(\log(2V))^2.
\]
Away from it the full complementary group \(M_{\ell,r}(E_x)\)
is annihilated by \(\ell^{\lceil V^{203}\rceil}\) for every odd
prime up to \(V^{199-6\varepsilon}/(\log(2V))^{3/2}\).
This cutoff exceeds \(Y\) for large \(V\), since
\(0<\varepsilon<1/52\). Also \(b_0(x)\ge\lceil V^{203}\rceil\)
on the whole interval for large \(V\). At such a prime
\(\ell\le\min\{Y,P(x)\}\), the subgroup \(U_{\ell,r}\) of
\(M_{\ell,r}\) is killed by \(\ell^{b_0(x)}\) as well, so the
quotient defining \(R_{\ell,r}\) is zero.

Take \(\mathcal E_V=(\mathcal D_V\cup\mathcal A_V)\cap[1,t_*]\).
The finite exceptional set of primes in
Corollary~\ref{cor:unboundedprimes428} depends only on the fixed
data, so it lies below \(Y\) on sufficiently distant intervals.
If \(Y<\ell\le P(x)\) and \(R_{\ell,r}\ne0\), that corollary
forces \(\ell-1\mid2r\). Thus \(r\in\mathcal A_V\). This also
accounts for the two moving primes by the same corollary. The two
prime ranges give the asserted simultaneous vanishing. Finally
\(V^{401}\ll T\), \(\log T\asymp\log V\), and
\(H(T)=o((\log T)^2)\). Adding the count of \(\mathcal D_V\)
to the upper estimate in Lemma~\ref{lem:shiftdivisors465} proves
the uniform weight count.
\end{proof}

\begin{corollary}[The remaining local mass on the same interval]\leavevmode
\label{cor:supportmass465}
At an allowed pair \(x,r\) write
\[
        W_x(r)=\sum_{\substack{2<\ell\le P(x)\\\ell\ {\rm prime}}}
                   (\log\ell)\operatorname{length}_{\ZZ_\ell}
                         R_{\ell,r}(E_x).
\]
For a constant \(C\) depending only on the fixed data one has
uniformly on the interval
\[
 0\le \mathcal R(x)-
       \sum_{\substack{1\le r<t(x)\\r\in\mathcal E_V}} W_x(r)
                      \ll_{\psi,\varepsilon} V^C.
\]
In particular the sufficient local estimate of
Proposition~\ref{prop:res23} would follow from the estimate
\[
 \max_{\substack{1\le r<t(x)\\r\in\mathcal E_V}} W_x(r)
       =o_{\psi,\varepsilon}\!
          \left(\frac{\exp(x^\varepsilon/2)H(t(x))}{t(x)}\right)
\]
uniformly on arbitrarily distant intervals of this family, with
the maximum of the empty set taken to be zero.
\end{corollary}
\begin{proof}
For \(\ell\le P(x)\) the sum over \(\mathcal E_V\) includes all
nonzero terms by Theorem~\ref{thm:simultaneousweights465}.
For \(\ell>P(x)\) the residual group is a subgroup of the full
complementary group. The weighted total for the latter is
\(O_{\psi,\varepsilon}(V^C)\) by
Corollary~\ref{cor:unboundedprimes428}. This proves the first
formula, including its sign. For the last implication multiply
the proposed maximum by the count at \(T=t(x)\). The resulting
bound is \(o(\exp(x^\varepsilon/2))\), and the polynomial error
has the same property because \(x\ge V\). It follows that
\(\mathcal R(x)\le \tfrac14\exp(x^\varepsilon/2)\) uniformly on
such intervals once \(V\) is large enough.
\end{proof}

The individual values of \(W_x(r)\) at the retained weights
remain to be bounded. Lemma~\ref{lem:shiftdivisors465} also
explains the dependence on arithmetic information beyond the
two inertia residues. The set of weights satisfying those
residues for at least one large prime already has order
\(T/H(T)\), whether or not its strict groups are nonzero.


\subsection{Full strict orders and the primes of the projected period}
\label{sec:sameprimes517}

Fix \(\psi\) and \(0<\varepsilon<1/52\). Take one of the
polynomial-conductor intervals \([V,V^{3/2}]\) of
Theorem~\ref{thm:conductor232}. Write \(E_*\) for the upper
endpoint field, \(n_*=[E_*:\QQ]\), \(N_*\) for its cyclotomic
modulus, \(R=\lfloor V^{1203/2}\rfloor\), and \(Y=V^{100}\).
Here \(\mathcal S_{p,r}(L)\) denotes the full strict group over
\(L\) for the twist \(-r\), with both signs included.
All constants below may depend on \(\psi,\varepsilon\).

Keep the fixed field \(K=\QQ(\sqrt{-s})\) and the split test
prime \(u\) of Proposition~\ref{prop:quadratic345}. Let \(D(V)\)
be the integer of Theorem~\ref{thm:originalorder424}. For each
\(r\) use the integral special-value determinant
\(D_{r+1}(E_*)\) of Lemma~\ref{lem:int13}, formed on \(E_*\) with
its modulus \(N_*\). Define
\[
 B(V)=\prod_{1\le r<R}|D_{r+1}(E_*)|,\qquad
                 J(V)=D(V)^2 B(V).
\]
All factors are nonzero and are fixed before the two auxiliary
period primes and the terminal prime are chosen. Indeed \(D(V)\)
uses the fixed split tests, the weights below \(R\), and the
characters of \(E_*\). The other factors are determined by \(E_*\)
and the two fixed denominator-clearing primes in
Lemma~\ref{lem:int13}. In particular \(\log J(V)\ll V^C\)
for some \(C\). For \(D(V)\) this is
Theorem~\ref{thm:originalorder424}. For \(B(V)\) the lemma gives
\(O(n_*R^2(n_*+\log(2N_*R)))\).

\begin{lemma}[An integer controlling the full strict order at separated twists]
\leavevmode\label{lem:fullinteger517}
There is a fixed finite set of primes \(\mathcal B\) such that
for every sufficiently distant interval, every prime
\(p\notin\mathcal B\), and every number subfield \(L\subseteq E_*\),
\[
 \prod_{\substack{1\le r<R\\p-1\nmid 2r}}
       |\mathcal S_{p,r}(L)|\quad\text{divides}\quad p^{v_p(J(V))}.
\]
In particular the product of the full orders on any such \(L\),
over all primes \(p>2R+1\) and all \(1\le r<R\), divides \(J(V)\).
\end{lemma}
\begin{proof}
Take \(\mathcal B\) to contain the primes below eleven and the
prime divisors of \(2suF[K_\psi:\QQ]|\Cl(K)|\), with \(F\) the
conductor of \(\psi\). The degree \(n_*\) divides a power of
two times \([K_\psi:\QQ]\). Thus \(p\nmid n_*\) outside
\(\mathcal B\). First work on \(E_*\). At a prime also coprime
to \(N_*\), Theorem~\ref{thm:adjacentweights427}, applied at
\(x=V^{3/2}\), bounds the sum of the complementary lengths over
the displayed weights by \(2v_p(D(V))\). The only other prime
divisors of \(N_*\) outside \(\mathcal B\) are the quadratic
primes \(q,d_0\). Theorem~\ref{thm:movingprimes428} gives the
same bound there even after including all weights.

For the opposite sign at a weight \(r\), put \(n=r+1\) and
\(e_n=(1+(-1)^nc)/2\). By Lemma~\ref{lem:fit13} the dual of
this strict group embeds in the \(e_n\)-part of degree-two
\(p\)-adic cohomology. In that summand the integral special-value
element \(e_n A\Theta_n\) of Lemma~\ref{lem:int13} lies in the
zeroth Fitting ideal. Indeed \(A\) is an integer multiple of
the annihilator \(W\) in Lemma~\ref{lem:fit13}. The abelian
group ring \(e_n\ZZ_p[\Gal(E_*/\QQ)]\) is a product of
unramified discrete valuation rings since \(p\nmid n_*\).
Over each factor the valuation of a member of the zeroth
Fitting ideal is at least the module length, by diagonalizing
a presentation over that ring. Multiplying by the residue
degrees and adding bounds the \(\ZZ_p\)-length by the valuation
of the determinant of multiplication by \(e_n A\Theta_n\).
After extending scalars its eigenvalues are the character
values whose product is \(D_n(E_*)\). Thus the sum of the
critical lengths over all these weights is at most \(v_p(B(V))\).
Adding the two signs gives the asserted divisibility on \(E_*\).

Restriction from \(L\) to \(E_*\) maps strict groups to strict
groups. It preserves vanishing at \(p\) and unramifiedness away
from \(p\) by restriction to the local subgroups. Corestriction
after restriction on cohomology is multiplication by
\([E_*:L]\), a unit at \(p\), so restriction is injective.
The order at every weight on \(L\) therefore divides the order
on \(E_*\). For \(p>2R+1\), every weight in the interval satisfies
\(p-1\nmid 2r\). Also all primes in the fixed set are eventually
smaller. Multiplying the prime-power divisibilities then gives
the last assertion.
\end{proof}

For each odd prime \(p\le 2R+1\), let \(\mathcal E_{p,\tau}\)
be the finite ordinary class group test of
Theorem~\ref{thm:finiteweights515} at cutoff \(R\), with
\(\tau=\kappa_p^{-r}|_\Delta\) as there. Define a finite set of
prime and weight pairs by
\[
 \mathcal F_V=\{(p,r):\ 2<p\le2R+1\text{ prime},   1\le r<R,\ r\in\mathcal E_{p,\tau},                   (p\le Y\ \text{or}\ p-1\mid 2r)\}.
\]
Let \(\mathcal W_V\) be its projection onto the weights.

\begin{theorem}[A polynomial bound for both signs outside the finite tests]
\leavevmode\label{thm:fullsparse517}
One has
\[
   \sum_{p>2}\ \sum_{\substack{1\le r<R\\(p,r)\notin\mathcal F_V}}
             \log |\mathcal S_{p,r}(E_*)| \ll V^C
\]
for some \(C\), where the outer sum is over primes. This same
bound holds for the groups over every \(E_x\), even up to the
endpoint weight cutoff \(R\). Uniformly for
\(\lfloor V^{401}\rfloor\le T'\le R\),
\[
 |\mathcal W_V\cap[1,T']|
    \ll \frac{T'}{(\log T')^\delta\sqrt{\log\log T'}},
 \qquad
 \delta=1-\frac{1+\log\log 2}{\log 2}.
\]
\end{theorem}
\begin{proof}
The full critical-sign sum on \(E_*\) is polynomial by
Theorem~\ref{thm:critical13} at the upper endpoint. For the
complementary sign first take \(p\le2R+1\) with either
\(p\le Y\) or \(p-1\mid2r\). Outside \(\mathcal F_V\) its length
is at most \(H_\tau\) by Theorem~\ref{thm:finiteweights515}.
Uniformly over this entire polynomial prime range, the proof
of Corollary~\ref{cor:smallresidual515} gives
\[
  H_\tau\log p\ll
   (1+\log(2n_*R))\,n_*^2 p^2\log(2N_*n_*p).
\]
There are at most \((2R+1)R\) pairs here. Both the number of
pairs and their individual log orders are polynomial in \(V\).
At every remaining prime and weight we have \(p>Y\) and
\(p-1\nmid 2r\). All such primes lie outside the fixed set of
the lemma. Summing its primewise bound costs at most \(\log J(V)\).
This proves the first assertion. At odd primes restriction to
\(E_*\) is injective for every \(E_x\), since the index is a
power of two. The bounds on the whole of \(E_*\) thus dominate
all these sums term by term.

Weights in \(\mathcal W_V\) coming from primes at most \(Y\)
number \(O(V^{306+24\varepsilon}\log(2V))\) by
Corollary~\ref{cor:smallresidual515}. This is smaller than the
displayed bound even at its shortest prefix. Every remaining
weight satisfies \(p-1\mid2r\) for some \(p>V^{100}\). The
upper estimate of Lemma~\ref{lem:shiftdivisors465}, applied
to that prefix, gives exactly the second term needed.
\end{proof}

\begin{corollary}[Testing the exceptional local length at a finite precision]
\leavevmode\label{cor:exactbudget517}
Fix \(x\) on the interval and a real threshold \(H\ge0\).
Write \(b=b_0(x)\) as in Section~\ref{sec:missing}.
For each prime occurring in \(\mathcal F_V\) put
\[
   m_p=b+1+\lfloor H/\log p\rfloor,\qquad
   j_p=m_p+\lfloor\log_p n_*\rfloor .
\]
For a pair with \(r<t(x)\) let
\(C=Q_{r,m_p}(\Cl(\mathcal O_{F_{j_p}}[1/p])\{p\})\)
over \(E_x\) as in Theorem~\ref{thm:ordinarydigits514}.
On this finite group let \(J^\vee_{p,r}\) be the field average
dual to \(J_{p,r}\) in Section~\ref{sec:missing}, or zero when
that operator is zero. Set
\[
  W^{\rm fin}(x,H)=
      \sum_{\substack{(p,r)\in\mathcal F_V\\r<t(x)}}
                    \log|p^b(1-J^\vee_{p,r})C|.
\]
The logarithmic sum of the orders of the actual residual groups
on these pairs is at most \(H\) if and only if
\(W^{\rm fin}(x,H)\le H\). In that case the two sums are equal.
The total contribution to \(\mathcal R(x)\) from all other pairs
is \(O(V^C)\).
\end{corollary}
\begin{proof}
For large \(V\) we have \(2R+1<P(x)\) at every \(x\) in this
interval, so the residual group in each tested pair uses
precision \(b\). Write \(A\) for the full complementary dual at
this pair. The common-field theorem identifies \(C=A/p^{m_p}A\).
The dual field action in the untwisted ideal-class coordinates
is represented by a lift \(g\) acting on ideals, multiplied by
\(\kappa_p(g)^r\). This compensates the relation
\(\delta=\kappa_p(\delta)^{-r}\) for lifts fixing \(E_x\), and
is the action obtained by dualizing the coefficient twist in
Theorem~\ref{thm:loginterval319}. In particular the average
over the finite relative field group acts on the finite quotient.
Its denominator is a power of two. The dual average is the
average over the inverses of the same elements. Projection onto
\(B=(1-J^\vee_{p,r})A\) commutes with reduction, as it is a
direct summand. Moreover the dual of the quotient by the subgroup
killed by \(p^b\) is \(p^bB\), by the cyclic decomposition of a
finite abelian group.

The finite summand thus uses \(p^b(B/p^{m_p}B)\). If the
reduction shortens a cyclic factor of \(B\), that factor alone
contributes \((m_p-b)\log p>H\) even to the finite sum.
Otherwise the entire contribution agrees with its full value.
All summands are nonnegative and reducing them can only
decrease their orders. This proves both implications and
the equality assertion. Outside \(\mathcal F_V\) each residual
order is at most the full order in Theorem~\ref{thm:fullsparse517},
which gives the last estimate.
\end{proof}

\begin{corollary}[Full vanishing at the three primes of the signed periods]
\leavevmode\label{cor:sameprimes517}
At each \(x\) on every sufficiently distant interval, the two
auxiliary primes \(g_1,g_2\) and the terminal prime \(T\) can be
chosen for the original fields, levels and words so that
\[
     \mathcal S_{p,r}(L)=0\quad
       (p\in\{g_1,g_2,T\},\ L\subseteq E_*,\ 1\le r<R).
\]
For the selected words the terminal window still contains
\(Y_x^{1-o(1)}\) eligible choices of \(T\), where \(Y_x\) is its
scale in Theorem~\ref{thm:common111}. At these same terminal
primes the signed congruences of Theorem~\ref{thm:connectedjets516}
hold at every precision on the period fields for which the
primes have been selected.
\end{corollary}
\begin{proof}
Exclude the prime divisors of \(J(V)\) before any of the three
choices. There are at most \(\log J(V)/\log2\ll V^C\) of them.
The first two primes are selected in \((U,2U]\), with
\(U=(k b_x)^8\) and \(b_x=\exp(x^\varepsilon)\), by
Proposition~\ref{prop:ray50} and Theorem~\ref{thm:levels88}.
There are \(U^{1-o(1)}\) candidates for the first prime. After
fixing it, the disjointness condition on the two sets
\(\mathcal B_r(g_i)\) rules out at most
\(O(r^3(1+U^{3/4}))\) integers for a polynomially bounded
rank parameter \(r\), by the proof of Theorem~\ref{thm:levels88}.
The extra polynomial exclusion leaves eligible choices for both.
Use the same words and admissible factors as in
Theorem~\ref{thm:common111} after fixing the auxiliary pair.
Its terminal candidate count has
\(\log Y_x\gg m x^\varepsilon\) for the prescribed chain length
\(m\). Subtracting \(O(V^C)\) again leaves \(Y_x^{1-o(1)}\)
candidates. The selected primes all exceed \(2R+1\), uniformly
since \(x\ge V\) and \(R\) is polynomial in \(V\).
The lemma now gives full vanishing on every stated subfield
and weight. These exclusions preserve the period congruences
and splitting conditions used in Lemma~\ref{lem:terminalblock496},
so the connected-coordinate calculation applies to the very
same terminal choices. The period choices themselves may
depend on \(x\).
\end{proof}

The remaining exceptional orders in Corollary~\ref{cor:exactbudget517}
need a uniform upper bound. On the period side the required rank
and archimedean height concern the signed columns after projection
at every translate in Section~\ref{sec:missing}.

\clearpage
\section{References}\label{nav:bibliography-I}
\begingroup\def\section*#1{}\begin{thebibliography}{99}\small
\bibitem{I-I-I-nist} NIST, \emph{Digital Library of Mathematical Functions}, Chapter 25, Section 25.15.\enspace\href{https://dlmf.nist.gov/25.15}{Dirichlet \(L\)-functions}.
\bibitem{I-I-I-elkies} N. D. Elkies, \emph{The asymptotic formula for primes in arithmetic progressions}, Math 229, Harvard University, course notes.\enspace\href{https://people.math.harvard.edu/~elkies/M229.19/pnt_q.pdf}{Dirichlet functions and their zeros}.
\bibitem{I-I-I-alpoge} L. Alp\"oge, \emph{Remark}, 20 November 2026. Pages 1 and 3 state the prime-sum conclusion and describe the theta-lift strategy.\enspace\href{https://alpo.ge/remark.pdf}{Original note}.
\bibitem{I-I-I-milne} J. S. Milne, \emph{Class Field Theory}, version 4.03, 2020, Chapter V, Theorem 3.27.\enspace\href{https://www.jmilne.org/math/CourseNotes/CFT.pdf}{Course notes}.
\bibitem{I-I-I-tz} J. Thorner and A. Zaman, \emph{Refinements to the prime number theorem for arithmetic progressions}, Corollary 1.4. \enspace\href{https://arxiv.org/abs/2108.10878}{Authors' paper}.
\bibitem{I-I-I-mont} H. L. Montgomery, \emph{Zeros of \(L\)-functions}, Invent. Math. \textbf{8} (1969), 346--354, Theorem 1. \enspace\href{https://link.springer.com/article/10.1007/BF01404638}{Original density theorem}.
\bibitem{I-I-I-behrend48} F. A. Behrend, \emph{On sets of integers which contain no three terms in arithmetical progression}, Proc. Nat. Acad. Sci. U.S.A. \textbf{32} (1946), 331--332. \href{https://doi.org/10.1073/pnas.32.12.331}{Original paper}.
\bibitem{I-I-I-tz50} J. Thorner and A. Zaman, \emph{A unified and improved Chebotarev density theorem}, Algebra Number Theory \textbf{13} (2019). \href{https://msp.org/ant/2019/13-5/ant-v13-n5-p02-p.pdf}{Theorem 1.4 and Theorem 3.1}.
\bibitem{I-I-I-tzleast50} J. Thorner and A. Zaman, \emph{An explicit bound for the least prime ideal in the Chebotarev density theorem}, Algebra Number Theory \textbf{11} (2017). \href{https://arxiv.org/abs/1604.01750}{Theorem 1.1}.
\bibitem{I-I-II-gross14} B. H. Gross, \emph{Algebraic modular forms}, Israel J. Math. \textbf{113} (1999), 61--93, Prop.~4.3. \href{https://doi.org/10.1007/BF02780173}{Original paper}.
\bibitem{I-I-ford465} K. Ford, \emph{Integers divisible by a large shifted prime}, Acta Arith. (2017), Theorem~1(i). \href{https://doi.org/10.4064/aa8495-2-2017}{Published paper}. \href{https://arxiv.org/pdf/1604.00281}{Author's preprint}.
\end{thebibliography}\endgroup
\end{document}
